\documentclass[12pt]{book}%

%\usepackage{fontspec}[greek]
%\usepackage{fourier-otf} % pour les arcs (double brève qui compte comme une longue)
%\usepackage{arcs} % pour les arcs (double brève qui compte comme une longue)
\usepackage{graphicx} % pour l'image de la coupe d'Onésime
\usepackage{fontspec}%
\usepackage{calc} % \ratio
\usepackage{fancyhdr}% en-têtes personnalisés
\usepackage[strict]{changepage}%pour modifier les marges des footnotes
\pagestyle{fancy}
\fancyfoot[LE,RO]{\thepage}
\fancyfoot[C]{}
\fancyhead[LE,RO]{}
%\fancyhead[RE]{\leftmark}
%\fancyhead[LO]{\rightmark}
\newcommand{\enTeteG}[1]{\fancyhead[LE]{#1}}% En-tête principale (pages paires de gauche)
\newcommand{\enTeteD}[1]{\fancyhead[RO]{#1}}% En-tête secondaire (pages impaires de droite)
\renewcommand{\headrulewidth}[0]{0pt}%filet de en-tête haut supprimé
\renewcommand{\footrulewidth}[0]{0.05mm}%filet de bas page
%\usepackage{totpages}
%\cfoot{{\thepage} / \ref{TotPages}}
%\usepackage{xunicode}%
%\usepackage{yfonts}%
\usepackage{hyperref}% pour le courriel
\usepackage{amsmath}% pour \overset
\usepackage{cancel}% pour \inRas
\usepackage{quoting}%
%\usepackage{csquotes} % Pour remplacer les guillemets anglais par des français
%\MakeOuterQuote{"}
\usepackage{ulem}% pour souligner en pointillé
\usepackage{comment}%
\usepackage{xcolor}%
\usepackage{fourier}% pour \widearc de la synizèze
%\usepackage{ragged2e}% pour la justification
%scale = {hscale, vscale}
%\usepackage[scale={0.75,0.8}]{geometry}%\usepackage[total={12cm,20cm}%Harmattan
%\usepackage[total={12cm,22cm}, top=3.5cm, right=4.5cm, footnotesep=1cm, includefoot]{geometry}% pour une page impaire
\usepackage[twoside, bindingoffset=1cm,total={12cm,22.5cm}, top=3.5cm, right=4.5cm, footnotesep=1cm, includefoot]{geometry}
%\usepackage{xpatch}
%\setlength{\hoffset}{-18pt}        
%\setlength{\oddsidemargin}{20pt} % Marge gauche sur pages impaires ATTENTION écrit 2 fois
%\setlength{\evensidemargin}{9pt} % Marge gauche sur pages paires
%\setlength{\marginparwidth}{54pt} % Largeur de note dans la marge
%\setlength{\textwidth}{481pt} % Largeur de la zone de texte (17cm)
%\setlength{\voffset}{-18pt} % Bon pour DOS
%\setlength{\marginparsep}{7pt} % Séparation de la marge
%\setlength{\topmargin}{0pt} % Pas de marge en haut
%\setlength{\headheight}{13pt} % Haut de page
%\setlength{\headsep}{10pt} % Entre le haut de page et le texte
%\setlength{\footskip}{27pt} % Bas de page + séparation
%\setlength{\textheight}{708pt} % Hauteur de la zone de texte (25cm)


%\newfontfamily\policeGr{Liberation Serif}[
%  Script=Greek,
%  NFSSFamily=liberationserif,
%
%]
%\newfontfamily\policeGr{Liberation Serif}%

%\newfontfamily{\policeFr}{Garamond Premier Pro}%
%\newfontfamily{\policeLat}{Book Antiqua}%
%\newfontfamily{\policeAn}{Comic Sans MS}%
%\newfontfamily{\policeAll}{Fraktur}%
%\newfontfamily{\policeIt}{BodoniFLF}%
%\newfontfamily{\policeEsp}{Great Moravia}%
% ...
%\setmainfont{Liberation Serif}%
%\setmainfont{Liberation Serif}%
%\newfontfamily{\policeFr}{Garamond Premier Pro}%
%\newfontfamily{\policeFr}{EB Garamond}%
\IfFontExistsTF{GFS Bodoni}{\newfontfamily{\policeFr}{GFS Bodoni}\newfontfamily{\policeGr}{GFS Bodoni}\setmainfont{GFS Bodoni}}{\newfontfamily{\policeFr}{GFS Didot}\newfontfamily{\policeGr}{GFS Didot}\setmainfont{GFS Didot}}%
%\newfontfamily{\policeLat}{Book Antiqua}%
\newfontfamily{\policeLat}{GFS Neohellenic}%
%\newfontfamily{\policeAn}{Comic Sans MS}%
\IfFontExistsTF{Courier 10 Pitch}{\newfontfamily{\policeAn}{Courier 10 Pitch}}{\newfontfamily{\policeAn}{FreeMono}}%
\newfontfamily{\policeEsp}{Abyssinica SIL}%
%\newfontfamily{\policeAll}{Fraktur}%
%\newfontfamily{\policeAll}{fraktur-bt.ttf}%
\IfFontExistsTF{Bertholdr Mainzer Fraktur}{\newfontfamily{\policeAll}{Bertholdr Mainzer Fraktur}}{\usepackage{yfonts}}%
%\newfontfamily{\policeIt}{BodoniFLF}%
\newfontfamily{\policeIt}{GFS Artemisia}%
%\newfontfamily{\policeEsp}{Great Moravia}%
%\newfontfamily\policeGr{GFSBodoni.otf}[BoldFont = GFSBodoniBold.otf, ItalicFont = GFSBodoniIt.otf, BoldItalicFont = GFSBodoniBoldIt.otf]%
%\setmainfont{GFSBodoni.otf}[BoldFont = GFSBodoniBold.otf, ItalicFont = GFSBodoniIt.otf, BoldItalicFont = GFSBodoniBoldIt.otf]

%\setdefaultlanguage[variant=ancient]{greek}%
\usepackage{polyglossia}%
\setdefaultlanguage{french}%
%\setotherlanguage[variant=ancient]{greek}%
%\setmainlanguage[variant=ancient]{greek} 
%\setotherlanguage[variant=ancient]{greek} 
%\setmainfont{Libertinus Serif} 
%\newfontfamily{\greekfont}{Liberation Serif}[
%  Script=Greek,
%  NFSSFamily=liberationserif,
%]

\usepackage[maxlevel=3]{csquotes}%
%\MakeOuterQuote{"}% pour les guillemets français
%\usepackage{verse}%
%\usepackage}[squarebr]{gmverse}% pour le rejet avec un crochet droit
\usepackage[parapparatus]{reledmac}%
\usepackage[parledgroup, nomaxlines]{reledpar}% 
%\usepackage[parledgroup, nomaxlines, nopbinverse]{reledpar}% nopbinverse empêche de casser un groupe
%\usepackage[parledgroup]{reledpar}% nomaxlines permettait d'avoir un nombre de lignes différents sur des pages qui se font face.
%\parledgrouptrue
%\setgoalfraction{0.9}% sept 2022
%\makeatletter
%\xpatchcmd{\mp@append@Xnotes}{\@parboxrestore}{\@parboxrestore\leftskip-1in}{}{}
%\makeatother
% Pour que les commentaires apparaissent sur les deux pages
\AtBeginDocument{%

%\Xbeforeinserting[B]{\espace{-2cm}}% code exécuté avant l'appel à la note
%\Xafternote[B]{\espace{20cm}}
%\Xbhookgroup[A,C]{\espace{-2cm}}
%\Xprenotes{2cm}
%\setlength{\skip\footins}{-2cm}
%\Xhangindent[B]{3cm}
%\sethangingsymbol{\protect\hfill}
%\versehangrightsquare% le rejet d'un vers se fait avec un crochet droit
\rightnoteupfalse 
\Xmaxhnotes[A,B]{0.50\textheight}%
\Xmaxhnotes[C]{0.22\textheight}%
%\Xmaxhnotes{0.9\textheight}%
%\maxhnotesX{0.45\textheight}%
\setRlineflag{}%
%\newcommand{\decoupeVers}[1]{\linenumannotation{#1}}%
\newcommand{\decoupeVers}[1]{}%
%\setlength{\oddsidemargin}{100pt} % Marge gauche sur pages impaires ATTENTION écrit 2 fois
% La Largeur du commentaire
%\Xwidth{\columnwidth+\marginparsep+\ledrsnotewidth}
% \Xgroupbyline[A,C]
\Xgroupbyline[A,C]
% \Xonlyside[A]{R}% imprime les testimonia page de droite
%\Xafterrule[A]{-1em}% supprime un peu d'esp vertical
%\Xwidth{\columnwidth+\marginparsep+\ledrsnotewidth}
\Xwidth{14cm}
%\Xsethangindent{20pt}%décale les footnotes
%\Xhangindent{-3cm}
%\Xparindent{-10cm}
%\sethangindentX{20pt}
%\Xhangindent = 20pt%
}%


\newcommand\Petit{\fontsize{9pt}{11pt}\selectfont} % x 1,2 pour e second argument qui est l'espacement interligne
\newcommand{\Taille}[1]{\fontsize{#1pt}{12pt}\selectfont}
\newcommand{\TaillePt}[1]{\fontsize{#1}{15pt}\selectfont}
\newlength{\LargAuto}%
\newlength{\LargTheor}%
\newlength{\Vrai}%
\newlength{\TP}%
\setlength{\TP}{12pt}%
\newcommand{\Automatique}[1]{%
\setlength{\LargTheor}{-72pt}%
\addtolength{\LargTheor}{\textwidth}%
\settowidth{\LargAuto}{#1}%
%\divide \LargTheor by \LargAuto%  \real{1.} *
\setlength{\Vrai}{\TP * \ratio{\LargTheor}{\LargAuto}}%
%\edef\TP{12* \ratio{\LargTheor}{\LargAuto}}%
{\TaillePt{\Vrai}#1}%
}%
\newlength{\largtext}%
\newcommand{\Strophe}[1]{\ledleftnote{%
\settowidth{\textelargeur}{\footnotesize{\francais{#1}}}%
%\espace{-\textelargeur}
\begin{minipage}{\textelargeur}\footnotesize{\francais{#1}}\end{minipage}}}%
\newcommand{\renverse}[1]{\rotatebox[origin=c]{180}{#1}}
\newcommand{\italique}[1]{\textit{#1}}%
\newcommand{\citComm}[1]{\begin{quoting}[vskip=0pt] #1 \end{quoting}}% vskip = \topsep + \partopsep par défaut
\newcommand{\passeLigne}[0]{\skipnumbering}%
\newcommand{\fin}[0]{\linebreak\passeLigne\protect\hfill $[$ }% Pour le rejet 
\newcommand{\ligne}[1]{\setline{#1}\hidenumbering}%
\newcommand{\espace}[1]{\hspace{#1}}%
\newcommand{\cesure}[0]{ [...] }%
\newcommand{\nieme}[1]{#1\italique{e}}%
\newcommand{\lettreBreve}[1]{$\breve{\mbox{#1}}$}
\newcommand{\synizese}[1]{$\widearc{\mbox{#1}}$}
\newcommand{\inRas}[2]{\if\relax\detokenize{#2}\relax\text{\sout{#1}}\else$\overset{\text{#2}}{\text{\sout{#1}}}$\fi}
\newcommand{\auDessus}[2]{$\overset{\text{#2}}{\text{#1}}$}
%\newcommand{\etc}[0]{κτλ}%
%\newcommand{\cit}[0]{\espace{1cm}} % indent. pour une citation
\newcommand{\expo}[1]{\textsuperscript{#1}}%
\newcommand{\alinea}[0]{ \ \ }%
\newcommand{\Fr}[0]{ Ar. \italique{fragm.} (Kassel-Austin) }%
\newcommand{\Syn}[0]{ Συναγωγή }%
\newcommand{\SynB}[0]{ Συναγωγή\expo{B} }% la deuxième version pour la lettre A
\newcommand{\ClemAlex}{ Clemens Alexandrinus }
\newcommand{\Eustathe}{ Eustathius }
\newcommand{\AelDion}{ Aelius Dionysius }
\newcommand{\EtymMagnum}{ Etymologicum magnum }
\newcommand{\EtymGen}{ Etymologicum genuinum }
\newcommand{\EtymSym}{ Etymologicum Symeonis }
\newcommand{\Apostol}{ Michael Apostolius }
\newcommand{\Macaire}{ Macarius } 
\newcommand{\Phryn}{ Phrynichus \italique{Praep. soph.} }
\newcommand{\Stobee}{ Joannes Stobaeus }
\newcommand{\EpimAlphab}{ Epimerismi alphabetici in Homerum (ed. Cramer) }
\newcommand{\GregCor}{ Gregorius Corinthius }
\newcommand{\Antiatt}{ Antiatticista (ed. Bekker) }
\newcommand{\Harp}{ Harpocration }% Autre édition de référence : Dindorf
\newcommand{\Rpc}[0]{R\expo{pc}}%
\newcommand{\Rac}[0]{R\expo{ac}}%
\newcommand{\Rdeux}[0]{R\expo{2}}%
\newcommand{\Mon}[0]{ M }%
\newcommand{\PsiUn}[0]{PSI-1194}
\newcommand{\PsiDeux}[0]{POxy-3839}
\newcommand{\PsiTrois}[0]{POxy-3840}
\newcommand{\PsiQuatre}[0]{POxy-4935}
\newcommand{\PsiCinq}[0]{POxy-1176}%P50

\newcommand{\francais}[1]{{\policeFr #1}}%
\newcommand{\latin}[1]{{\policeLat #1}}%
\IfFontExistsTF{Bertholdr Mainzer Fraktur}{\newcommand{\allemand}[1]{{\policeAll #1}}}{\newcommand{\allemand}[1]{{\textfrak{#1}}}}%
%\newcommand{\allemand}[1]{{\textfrak{#1}}}%
%\newcommand{\grec}[1]{{\policeGr #1}}%
\newcommand{\grec}[1]{#1}%
\newcommand{\anglais}[1]{{\policeAn #1}}%
\newcommand{\italien}[1]{{\policeIt #1}}%
\newcommand{\espagnol}[1]{{\policeEsp #1}}%

\newcommand{\Nauck}[0]{N\expo{2}}%
\newcommand*{\auteur}[1]{\textsc{#1}*}%
\newcommand{\cf}[0]{\italique{cf.}\xspace}%
\newcommand{\chap}[1]{\^{#1}}%
%\newcommand{\chap}[1]{$\hat{#1}$}


%\marginparwidth = 1.2\marginparwidth
%\ledrsnotewidth = 1.2\marginparwidth
\newlength{\larg}%
\newcommand{\decalage}[1]{%
\settowidth{\larg}{#1}%
\espace{\larg}%
}
\newcommand{\decalageN}[1]{%
\settowidth{\larg}{#1}%
\espace{-\larg}%
}
%\newcommand{\apparat}[2]{\ledrightnote{\espace{-1.2\marginparwidth}\begin{minipage}{1.4\marginparwidth}\footnotesize{\textbf{#1} #2}\end{minipage}}}%
\newcommand{\apparat}[2]{\ledrightnote{\espace{-2cm}\begin{minipage}{4.5cm}\footnotesize{\textbf{#1} #2}\end{minipage}}}%
%\newcommand{\lecon}[1]{\dotuline{#1}} % la leçon retenue de l'apparat critiqe
\newcommand{\lecon}[1]{\textbf{#1}}
%\newcommand{\douteux}[1]{\dotuline{#1}} % pour les lettres douteuses : un point en-dessous
\newcommand{\douteux}[1]{$\underset{\cdot}{\text{#1}}$}
\newcommand{\difference}[1]{\dotuline{#1}}


%\setsidenotesep{ sep }%
\newcommand{\espApp}[0]{ || }%\espace{10pt}%
\newcommand{\espTest}[0]{ || }%\espace{10pt}%
\newcommand{\ligneApp}[1]{ \ \ \textbf{#1} }%
%% 3 arguments : les numéros de ligne, le lemme et la note elle-même



%\newcommand{\commentaire}[3]{\edtext{}{\Bfootnote[nonum, nosep]{\espace{-2cm}\begin{minipage}{16cm}\espace{-1cm}\textbf{#1~~#2} #3\end{minipage}}}}%
\newcommand{\commentaire}[3]{\edtext{}{\Bfootnote[nonum,nosep]{\begin{adjustwidth}{-1.5cm}{-2cm}\espace{-1.5cm}\textbf{#1}\decalageN{\textbf{#1}}\espace{0pt}\espace{1.5cm}\textbf{#2} #3\end{adjustwidth}}}}% Avec une * : adjustwidth* pour la symétrie des pages
%\renewcommand{\footfudgefiddle}{68}%Note that you must call it before 
% \Xarrangement[C]{paragraph}% entrâine trop d'erreurs, cf. documentation
\newcommand{\debutNote}[1]{\edtext{}{\Cfootnote[nonum, nosep]{\espace{-1cm}\espace{-15pt}\textbf{\textbf{#1}}}}}%
%\newcommand{\debutNote}[1]{\edtext{}{\Cfootnote[nonum, nosep]{\espace{-1cm}\espace{-15pt}\textbf{\dotuline{\textbf{#1}}}}}}%
\newcommand{\note}[2]{\edtext{}{\Cfootnote[nonum, nosep]{\espace{-15pt}\textbf{#1} #2}}}%
\newcommand{\testimonia}[2]{\edtext{}{\Afootnote[nonum, nosep]{\espace{-15pt}\textbf{#1} #2}}}%
%\newcommand{\testimonia}[2]{}
\newcommand{\marqueur}[1]{\textsuperscript{#1}}%
%\newcommand{\cat}[0]{$""_"^"$}%
\newcommand{\cat}[0]{\texorpdfstring{\textsubscript{\textasciicircum}}{cat}}%
%\fnpos{critical-familiar}%
%\mpfnpos{familiar-critical}%

\newcommand{\eff}[1]{%
{\color{orange}#1}%
}%

\newlength{\textlarg}%
\newcommand{\typeVers}[1]{%
\settowidth{\textlarg}{\footnotesize{#1}}%
\hspace{-0.8cm}{\footnotesize{#1}}\hspace{-\textlarg}\hspace{0.8cm}}%
%\newlength{\largtext}%
%\newcommand{\strophe}[1]{%
%\settowidth{\largtext}{~~\footnotesize{\francais{#1}}}%
%~~\footnotesize{\francais{#1}}%\hspace{-\largtext}%
%}%
\newcommand{\strophe}[1]{\footnotesize{~\francais{#1}}}%
%\newcommand{\strophe}[1]{%
%\settowidth{\largtext}{\footnotesize{\francais{#1}}}%
%	\ledleftnote{%
		%\hspace{-0.5cm}\footnotesize{\francais{#1}}\hspace{\largtext}\hspace{0.5cm%}%
%		\hspace{-0.5cm}\footnotesize{\francais{#1}}\hspace{\largtext}\hspace{0.5cm}%
%	}%
%}%
%\newcommand{\personnageGrec}[1]{\ledleftnote{\textbf{\large{#1}}\espace{-1.5cm}}}%
\newlength{\textelargeur}%
\newcommand{\personnageGrec}[1]{\ledleftnote{%
%\settowidth{\textelargeur}{\large{\textbf{#1}}}%
\settowidth{\textelargeur}{\large{\textbf{Στεφανόπωλις }}}%
%\espace{-\textelargeur}
\begin{minipage}{\textelargeur}\begin{flushright}\large{\textbf{#1}}\end{flushright}\end{minipage}}}%
%\newcommand{\Strophe}[1]{\footnotesize{\francais{#1}}}%
%\newcommand{\Strophe}[1]{\ledleftnote{%
%\settowidth{\textelargeur}{\footnotesize{\francais{#1}}}%
%\espace{-\textelargeur}
%\begin{minipage}{\textelargeur}\footnotesize{\francais{#1}}\end{minipage}}}%
%\newcommand{\personnageFr}[1]{\ledleftnote{\textbf{\large{\francais{#1}}}\espace{-1.5cm}}}%
\newlength{\texto}%
\newcommand{\personnageFr}[1]{\ledleftnote{%
%\settowidth{\texto}{\large{\textbf{#1}}}%
\settowidth{\texto}{\large{\textbf{Marchande d}}}%
%\espace{-\textelargeur}
\begin{minipage}{\texto}\begin{flushright}\large{\textbf{#1}}\end{flushright}\end{minipage}}}%
%% Les différents types de vers
\newcommand{\trimIamb}[1]{%
\hangindent=20pt%
\espace{0pt}% la valeur précédente -20
#1%
}%
\newcommand{\trimIambFr}[1]{%
\hangindent=20pt%
\espace{0pt}% la valeur précédente -20
\francais{#1%
}}%
\newcommand{\trimIambCat}[1]{%
\hangindent=53pt%
\espace{33pt}% la valeur précédente -20
#1%
}%
\newcommand{\trimIambCatFr}[1]{%
\hangindent=53pt%
\espace{33pt}% la valeur précédente -20
\francais{#1%
}}%
\newcommand{\tetramIamb}[1]{%
\hangindent=10pt%
\espace{-10pt}% la valeur précédente -20
#1%
}%
\newcommand{\tetramIambFr}[1]{%
\hangindent=20pt%
\espace{0pt}% la valeur précédente -20
\francais{#1%
}}%
\newcommand{\tetramTroch}[1]{%
\hangindent=20pt%
\espace{0pt}% la valeur précédente -20
#1%
}%
\newcommand{\tetramTrochFr}[1]{%
\hangindent=20pt%
\espace{0pt}% la valeur précédente -20
\francais{#1%
}}%
\newcommand{\tetramIambCat}[1]{%
\hangindent=10pt%
\espace{-10pt}% la valeur précédente -20
#1%
}%
\newcommand{\tetramIambCatFr}[1]{%
\hangindent=20pt%
\espace{0pt}% la valeur précédente -20
\francais{#1%
}}%
\newcommand{\dimAnap}[1]{%
\hangindent=70pt%
\espace{50pt}% la valeur précédente -20
#1%
}%
\newcommand{\dimAnapFr}[1]{%
\hangindent=70pt%
\espace{50pt}% la valeur précédente -20
\francais{#1%
}}%
\newcommand{\anap}[1]{%
\hangindent=70pt%
\espace{50pt}% la valeur précédente -20
#1%
}%
\newcommand{\anapFr}[1]{%
\hangindent=70pt%
\espace{50pt}% la valeur précédente -20
\francais{#1%
}}%
% 2 dochmiaques
\newcommand{\diDoch}[1]{%
\hangindent=60pt%
\espace{40pt}% la valeur précédente -20
#1%
}%
\newcommand{\diDochFr}[1]{%
\hangindent=60pt%
\espace{40pt}% la valeur précédente -20
\francais{#1%
}}%
\newcommand{\tetramAnap}[1]{%
\hangindent=20pt%
\espace{0pt}% la valeur précédente -20
#1%
}%
\newcommand{\tetramAnapFr}[1]{%
\hangindent=20pt%
\espace{0pt}% la valeur précédente -20
\francais{#1%
}}%
\newcommand{\chantAgath}[1]{%
\hangindent=90pt%
\espace{70pt}% la valeur précédente -20
#1%
}%
\newcommand{\chantAgathFr}[1]{%
\hangindent=90pt%
\espace{70pt}% la valeur précédente -20
\francais{#1%
}}%
%% Pour les "vers" 295 à 311
\newcommand{\prose}[1]{%
\hangindent=30pt%
\espace{10pt}% la valeur précédente -20
#1%
}%
\newcommand{\proseFr}[1]{%
\hangindent=30pt%
\espace{10pt}% la valeur précédente -20
\francais{#1%
}}%
%% pour les vers courts de structure changeante dans un choeur
\newcommand{\chantChoeur}[1]{%
%{\centering
%	#1\par
%}%
\hangindent=50pt%
\espace{30pt}% la valeur précédente -20
#1%
}%
\newcommand{\chantChoeurFr}[1]{%
%{\centering
%	\francais{#1}\par
%}%
\hangindent=50pt%
\espace{30pt}% la valeur précédente -20
\francais{#1%
}}%
%% pour les vers 1015 à 1055
\newcommand{\paratragiques}[1]{%
%{\centering
%	#1\par
%}%
\hangindent=50pt%
\espace{30pt}% la valeur précédente -20
#1%
}%
\newcommand{\paratragiquesFr}[1]{%
%{\centering
%	\francais{#1}\par
%}%
\hangindent=30pt%
\espace{10pt}% la valeur précédente -20
\francais{#1%
}}%
% pour indiquer que les notes marginales s'alignent au niveau de la première ligne
%\leftnoteupfalse
%\rightnoteupfalse



\newenvironment{groupe}{%
\pstart%
\ledgroup%
}{%
\pend %
}%

\newcommand{\LSJ}[0]{LSJ}%

%% pour voir tous les numéros de ligne (version brouillonne)
\linenumincrement{1}%
\firstlinenum{1}%
%% Lignes à imprimer
\gdef\linenumberlist{2,6,10,15,20,24,30,36,40,45,50,55,60,65,69,74,80,87,91,94,102,105,110,115,120,125,131,135,140,145,150,155,160,165,170,174,180,185,190,196,201,205,215,220,225,229,234,239,245,250,261,265,270,276,280,285,290,295,300,305,310,315,320,325,330,335,340,345,350,355,360,365,370,375,379,385,390,395,400,405,410,415,420,425,430,435,440,445,450,455,460,465,470,475,480,485,490,495,499,501,505,510,515,521,526,530,535,541,545,550,556,561,565,572,575,580,585,591,596,600,609,614,620,624,629,635,640,650,656,660,665,670,675,680,685,690,695,700,708,713,720,725,729,735,741,745,750,755,761,766,770,775,780,786,790,795,800,805,810,815,820,825,831,835,840,845,850,856,860,865,870,875,880,884,890,894,900,905,915,921,930,935,940,946,950,955,960,965,970,975,980,985,990,996,1000,1010,1016,1020,1025,1030,1035,1040,1045,1050,1055,1060,1066,1071,1076,1100,1106,1116,1120,1124,1130,1141,1147,1150,1156,1161,1166,1171,1175,1180,1185,1190,1195,1200,1205,1211,1215,1221,1224,1230}



\usepackage{adjustbox}
\usepackage{array}
\setmonofont{FreeMono}
\setlength{\headheight}{16pt}
\renewcommand{\enTeteG}[1]{\fancyhead[LE]{\small #1}}
\renewcommand{\enTeteD}[1]{\fancyhead[RO]{\small #1}}

% In Spread 2, critical apparatus and testimonia are already displayed on Spread 1 Left Page:
\renewcommand{\apparat}[2]{}
\renewcommand{\testimonia}[2]{}
\makeatletter
\renewcommand*{\led@err@LeftOnRightPage}{}
\renewcommand*{\led@err@RightOnLeftPage}{}
\makeatother
\AtBeginDocument{%
  \Xmaxhnotes[A,B]{0.82\textheight}%
  \Xmaxhnotes[C]{0.55\textheight}%
}
\renewcommand{\commentaire}[3]{\edtext{}{\Bfootnote[nonum,nosep]{\begin{adjustwidth}{-1.5cm}{-2cm}\fontsize{8.0pt}{9.4pt}\selectfont\espace{-1.5cm}\textbf{#1}\decalageN{\textbf{#1}}\espace{0pt}\espace{1.5cm}\textbf{#2} #3\end{adjustwidth}}}}%
\renewcommand{\debutNote}[1]{\edtext{}{\Cfootnote[nonum, nosep]{\fontsize{7.8pt}{9.2pt}\selectfont\espace{-1cm}\espace{-15pt}\textbf{#1}}}}%
\renewcommand{\note}[2]{\edtext{}{\Cfootnote[nonum, nosep]{\fontsize{7.8pt}{9.2pt}\selectfont\espace{-15pt}\textbf{#1} #2}}}%

% Row macro for Spread 1 Left Page (Diplomatic-yet-corrected Ravennas + inline apparatus)
\newcommand{\ligneRav}[4]{%
  \noindent\hbox to 16.5cm{%
    \hbox to 0.55cm{\hfill\tiny\texttt{#1}\hspace{0.1cm}}%
    \hbox to 1.25cm{\hfill\scriptsize #2\hspace{0.15cm}}%
    \hbox to 8.1cm{\adjustbox{max width=8.0cm}{#3}\hfill}%
    \hspace{0.2cm}%
    \parbox[t]{6.4cm}{\raggedright\fontsize{6.7pt}{7.9pt}\selectfont #4}%
  }\par\vspace{0.35pt}%
}

% Row macro for Spread 2 in the Strict 4-Page-Per-Folio Edition (Established Greek / French)
\newcommand{\ligneEst}[3]{%
  \noindent\hbox to 16.5cm{%
    \hbox to 2.5cm{\hfill\scriptsize #1\hspace{0.2cm}}%
    \hbox to 13.2cm{\adjustbox{max width=13.1cm}{#2}\hfill}%
    \hbox to 0.8cm{\hfill\scriptsize\texttt{#3}}%
  }\par\vspace{0.3pt}%
}

\begin{document}
\thispagestyle{empty}
\noindent\begin{adjustbox}{max totalheight=19.5cm, max width=\textwidth, keepaspectratio}\begin{minipage}{\textwidth}\begin{center}
\vspace*{0.4cm}
{\Huge\textbf{ΑΡΙΣΤΟΦΑΝΟΥΣ ΘΕΣΜΟΦΟΡΙΑΖΟΥΣΑΙ}}\par\vspace{0.5cm}
{\LARGE\textbf{ARISTOPHANE --- LES THESMOPHORIES}}\par\vspace{0.8cm}
{\Large\textit{Nouvelle édition philologique et paléographique sur quatre pages}}\par\vspace{0.4cm}
{\large Tome I : Folios 01 à 15 (\texttt{ff. 162v–169v}, \textit{Hypothesis} \& vv. 1–634)}\par\vspace{0.5cm}
\hrule height 0.8pt\vspace{0.6cm}
{\normalsize
\textbf{Architecture des doubles pages (30 folios du \textit{Codex Ravennas} 429, \texttt{ff. 162v–177r}, vv. 1–1231) :}\par\vspace{0.3cm}
\begin{itemize}
  \item \textbf{Première double page — Page de gauche (pair) :} Transcription diplomatique corrigée du \textit{Codex Ravennas} 429 (\textbf{R}), avec \textit{notae personarum} tachygraphiques, \textit{paragraphoi}, \textit{dikola}, grattages (\textit{in rasura}), ajouts marginaux, les quatre \textit{parepigraphai} scéniques, soulignement pointillé des émendations (\dotuline{texte}), \textbf{apparat critique dans le prolongement de la ligne}, et en bas de page : (1) \textbf{\textit{loci paralleli}} et (2) \textbf{suite (en commentaire) de l'apparat critique détaillé}.
  \item \textbf{Première double page — Page de droite (impair) :} Fac-similé haute résolution en pleine page du folio correspondant du \textit{Codex Ravennas} 429 (Biblioteca Classense, Ravenne).
  \item \textbf{Seconde double page — Page de gauche (pair) :} Texte grec établi proprement, avec colométrie et attribution complète des répliques.
  \item \textbf{Seconde double page — Page de droite (impair) :} Traduction française en regard.
  \item \textbf{Bas de la seconde double page (gauche et droite) :} Long commentaire philologique, historique, métrique et scénique.
\end{itemize}
}\vspace{0.6cm}
\hrule height 0.8pt\vspace{0.5cm}
{\Large\textbf{Édition, traduction et commentaire par Louis HAUSEUX (\textit{Bibi})}}\par
\end{center}\end{minipage}\end{adjustbox}
\newpage

\clearpage\ifodd\value{page}\hbox{}\newpage\fi
\enTeteG{ΠΡΟΛΟΓΟΣ --- Codex Ravennas 429, f. 162v (Classense 164v) --- vv. 1–42}
\begin{adjustwidth}{-1.5cm}{-3.0cm}
\begin{adjustbox}{max totalheight=21.0cm, max width=16.5cm, keepaspectratio}
\begin{minipage}{16.5cm}
{\fontsize{7.2pt}{8.5pt}\selectfont
\noindent\textbf{Folio 01 --- Codex Ravennas 429, \texttt{f. 162v} (\texttt{Classense 164v}) : ΠΡΟΛΟΓΟΣ, vv. 1–42}\hfill \textit{Témoins :} $\text{R}^1, \Rpc, \Rdeux$ (\Mon)\par\vspace{2pt}
\hrule height 0.4pt\vspace{2.5pt}
\noindent\textbf{\scriptsize [$\Rdeux$ — \italique{Hypothesis} :]} {\scriptsize Εὐριπίδης δράματά τινα γράψας κατὰ γυναικῶν κινδυνεύει ὑπὸ τῶν γυναικῶν ἐν θεσμοφορίοις ἀποθανεῖν· διὸ πείθει Ἀγάθωνα τὸν τραγῳδοποιὸν γυναικείαν στολὴν ἀναλαβόντα εἰσελθεῖν εἰς τὰ θεσμοφόρια καὶ ὑπὲρ αὐτοῦ ἀπολογήσασθαι. τοῦ δὲ ἀντειπόντος, Μνησίλοχος ὁ κηδεστὴς αὐτοῦ ξυρηθεὶς καὶ γυναικείαν στολὴν ἀναλαβὼν εἰσέρχεται· καὶ φωραθεὶς ὑπὸ Κλεισθένους... προλογίζει δὲ Μνησίλοχος κηδεστὴς Εὐριπίδου.}\par\vspace{1pt}
\noindent\textbf{\scriptsize [$\Rdeux$ — \italique{Τὰ τοῦ δράματος πρόσωπα} :]} {\scriptsize Μνησίλοχος κηδεστὴς Εὐριπίδου $\cdot$ Εὐριπίδης $\cdot$ Θεράπων Ἀγάθωνος $\cdot$ Ἀγάθων $\cdot$ Κῆρυξ $\cdot$ Χορὸς ἐκ Θεσμοφοριαζουσῶν $\cdot$ Γυνή τις $\cdot$ Γυνὴ ἑτέρα $\cdot$ Κλεισθένης $\cdot$ Πρύτανις $\cdot$ Τοξότης.}\par\vspace{1.5pt}
\centerline{\footnotesize\textbf{$\ast\;\cdot\;\ast$ ἀριστοφάνους θεσμοφοριάζουσαι $\ast\;\cdot\;\ast$}}\vspace{1.5pt}
\hrule height 0.3pt\vspace{2pt}
}
{\fontsize{8.1pt}{9.5pt}\selectfont
\ligneRav{1}{μ\auDessus{ν}{Η}}{ὦ Ζεῦ \difference{Χελιδὼν} \difference{ἆρά ποτε} φανήσεται;}{\textbf{1} ἄρα πότε R : \lecon{ἆρά ποτε} Biset}
\ligneRav{2}{}{ἀπολεῖ μ᾽ ἀλοῶν \difference{ἅνθρωπος} ἐξ ἑωθινοῦ·}{\textbf{2} ἄνθρωπος R : \lecon{ἅνθρωπος} Boissonade}
\ligneRav{3}{}{παρὰ σοῦ πυθέσθαι ποῖ μ᾽ ἄγεις ωὖριπίδη;}{\textbf{3} τι R : \lecon{τε} Σ (δυνατόν), Suid. οι 150}
\ligneRav{4}{}{οἷόν \difference{τε}, πρὶν τὸν σπλῆνα κομιδῇ μ᾽ ἐκβαλεῖν·}{}
\ligneRav{5}{\auDessus{εὐ}{ρ}}{ἀλλ᾽ οὐκ ἀκούειν δεῖ σε πάνθ᾽ \difference{ὅσ᾽} αὐτίκα}{\textbf{5} ὅσ R : \lecon{ὅσ᾽} \Mon}
\ligneRav{}{—}{ὄψει παρεστώς. : πῶς λέγεις; αὖθις φράσον·}{}
\ligneRav{}{—}{οὐ δεῖ μ᾽ ἀκούειν; : οὐχ ἅ γ᾽ ἂν μέλλῃς ὁρᾶν·}{}
\ligneRav{}{μ\auDessus{ν}{Η}}{οὐδ᾽ ἆρ᾽ ὁρᾶν δεῖ μ᾽; : οὐχ ἅ γ᾽ ἂν ἀκούειν δέῃ·}{}
\ligneRav{}{μ\auDessus{ν}{Η}}{πῶς μοι παραινεῖς; δεξιῶς μέντοι λέγεις·}{}
\ligneRav{10}{}{οὐ φῂς σὺ χρῆναί μ᾽ οὔτ᾽ ἀκούειν οὔθ᾽ ὁρᾶν·}{\textbf{10} \lecon{ὁρᾶν.} R : ὁρᾶν; Gelenius}
\ligneRav{}{\auDessus{εὐ}{ρ}}{χωρὶς γὰρ αὐτοῖν ἑκατέρου 'στὶν ἡ φύσις·}{}
\ligneRav{}{μ\auDessus{ν}{Η}}{τοῦ μήτ᾽ ἀκούειν μήθ᾽ ὁρᾶν; : εὖ ἴσθ᾽ ὅτι·}{}
\ligneRav{}{μ\auDessus{ν}{Η}}{πῶς χωρίς; : οὕτω ταῦτα διεκρίθη τότε·}{}
\ligneRav{}{}{αἰθὴρ γὰρ ὅτε τὰ πρῶτα διεχωρίζετο}{}
\ligneRav{15}{}{καὶ ζᾦ᾽ ἐν αὑτῷ ξυνετέκνου κινούμενα,}{}
\ligneRav{16}{}{ᾧ μὲν βλέπειν χρὴ \difference{πρῶτ᾽} ἐμηχανήσατο}{\textbf{16} πρῶτα R : \lecon{πρῶτ᾽} Grynaeus}
\ligneRav{}{}{ὀφθαλμὸν ἀντίμιμον ἡλίου τροχῷ,}{}
\ligneRav{18}{}{ἀκοὴν δὲ χοάνης ὦτα διετετρήνατο·}{\textbf{18} \lecon{ἀκοὴν δὲ χοάνης} R, Willems : ἀκοὴς δὲ χοάνην Biset : ἀκοῇ δὲ χοάνην Boissonade}
\ligneRav{}{μ\auDessus{ν}{Η}}{διὰ τὴν χοάνην οὖν μήτ᾽ ἀκούω μήθ᾽ ὁρῶ;}{}
\ligneRav{20}{}{νὴ τὸν Δί᾽ ἥδομαί γε τουτὶ προσμαθών·}{}
\ligneRav{21}{}{οἷόν \difference{γέ} που ᾽στιν αἱ σοφαὶ ξυνουσίαι·}{\textbf{21} τε R : \lecon{γε} Σ}
\ligneRav{}{\auDessus{εὐ}{ρ}}{πόλλ᾽ ἂν μάθοις τοιαῦτα παρ᾽ ἐμοῦ. : πῶς ἂν οὖν}{}
\ligneRav{(23)}{}{$[\![\text{\sout{πρὸς τοῖς ἀγαθοῖς τούτοισιν ἐξεύροιμ᾽ ὅπως}}]\!]$}{\textbf{v. 23} πρὸς τοῖς ἀγαθοῖς τούτοισιν ἐξεύροιμ᾽ ὅπως R : \lecon{del.} Willems}
\ligneRav{24}{}{ἔτι \difference{προσμάθοιμι} χωλὸς εἶναι τὼ σκέλει·}{\textbf{24} προσμάθοι· μὴ R : \lecon{προσμάθοιμι} Ellebodius}
\ligneRav{25}{\auDessus{εὐ}{ρ}}{βάδιζε δευρὶ καὶ πρόσεχε τὸν νοῦν. : ἰδού·}{}
\ligneRav{}{\auDessus{εὐ}{ρ}}{ὁρᾷς τὸ θύριον τοῦτο; : νὴ τὸν Ἡρακλέα}{}
\ligneRav{27}{—}{οἶμαί γε : σίγα \difference{νυν} : σιωπῶ τὸ θύριον· \hfill\textbf{[α \Rpc]}}{\textbf{27} νῦν R : \lecon{νυν} Boissonade\espApp θύριον R : \lecon{θύριον;} Dindorf}
\ligneRav{28}{\auDessus{εὐ}{ρ}}{ἄκου᾽ : \difference{ἀκούω} καὶ σιωπῶ τὸ θύριον· \hfill\textbf{[β \Rpc]}}{\textbf{28} ἀκούσω R : \lecon{ἀκούω} Brunck \espApp vers. 28 post 30 transp. \Rac : \lecon{corr.} \Rpc}
\ligneRav{29}{\auDessus{εὐ}{ρ}}{ἐνταῦθ᾽ Ἀγάθων ὁ κλεινὸς οἰκῶν τυγχάνει \hfill\textbf{[γ \Rpc]}}{}
\ligneRav{30}{—}{ὁ τραγῳδοποιός : ποῖος οὗτος \difference{Ἁγάθων}· \hfill\textbf{[δ \Rpc]}}{\textbf{30-1} Ἀγάθων (bis) R : \lecon{Ἁγάθων} et \lecon{ἀγαθῶν} Bibi}
\ligneRav{31}{}{ἔστιν τις \difference{ἀγαθῶν}; μῶν ὁ μέλας ὁ καρτερός·}{\textbf{31} ἔστιν τις ἀγαθῶν \lecon{(Μνη.) continuat} R : (Εὐρ.) \Mon\espApp ἔστιν τις ἀγαθῶν \lecon{(Μνη.) continuat} R : (Εὐρ.) \Mon \espApp Ἀγάθων R : \lecon{ἀγαθῶν;} Bibi}
\ligneRav{32}{\auDessus{εὐ}{ρ}}{οὔκ, ἀλλ᾽ ἕτερός τις· $[\![\text{\sout{μ\auDessus{ν}{Η}}}]\!]$ οὐχ \difference{ἑόρακας} πώποτε;}{\textbf{32} (Μνη.) οὐχ ... πώποτε \Rac : \lecon{(Εὐρ.) continuat} \Rpc, Σ \espApp ἑόρακα R : \lecon{ἑόρακας} Bentley \ligneApp{v. 33} omm. \Rac, \Mon : in marg. \Rpc}
\ligneRav{33}{μ\auDessus{ν}{Η}}{$[\![\text{in marg. \Rpc{} : }$μῶν ὁ δασυπώγων : οὐχ ἑόρακας πώποτε·$]\!]$}{\textbf{33} omm. $\text{R}^{\text{ac}}, \text{M}$ : \lecon{in marg.} \Rpc{}}
\ligneRav{34}{μ\auDessus{ν}{Η}}{μὰ τὸν Δί᾽ οὔτοι γ᾽ ὥστε \difference{κἀμέ γ᾽} (\inRas{καὶ}{κἀ}μέ γ᾽) εἰδέναι·}{\textbf{34} καὶ μεγ᾽ \Rac \ (ὶ eras. \Rpc) : καὶ μεγ᾽ \Mon : \lecon{κἀμέ γ᾽} edd. omn.}
\ligneRav{35}{\auDessus{εὐ}{ρ}}{καὶ μὴν \difference{†}βε\inRas{α}{β}ίνηκας\difference{†} σύ γ᾽, ἀλλ᾽ οὐκ οἶσθ᾽ ἴσως·}{\textbf{35} βε\inRas{α}{β}ίνηκας \Rpc{} (ex βεαίνηκας \Rac{} ut vid.) : \lecon{†βεβίνηκας†} \lecon{Bibi}}
\ligneRav{}{}{ἀλλ᾽ ἐκποδὼν πτήξωμεν, ὡς ἐξέρχεται}{}
\ligneRav{}{}{θεράπων τις αὐτοῦ πῦρ ἔχων καὶ μυρρίνας}{}
\ligneRav{}{}{προθυσόμενος ἔοικε τῆς ποιήσεως·}{}
\ligneRav{39}{θερ/}{\hspace{1.1em}εὔφημος πᾶς ἔστω \difference{λαός},}{\textbf{39} λέως R, Suid. ε 3792 codd. AF : \lecon{λαός} Σ, Suid. codd. GIVM}
\ligneRav{40}{}{\hspace{1.1em}στόμα συγκλῄσας· ἐπιδημεῖ γὰρ}{}
\ligneRav{}{}{\hspace{1.1em}θίασος Μουσῶν ἔνδον μελάθρων}{}
\ligneRav{}{}{\hspace{1.1em}τῶν δεσποσύνων μελοποιῶν·}{}
}
\vspace{3pt}\noindent\rule{4.5cm}{0.4pt}\par\vspace{1.5pt}
{\fontsize{7.0pt}{8.3pt}\selectfont
\noindent\textbf{Loci paralleli (vv. 1–42) ---} \textbf{1} \Fr 617 πυθοῦ χελιδὼν πηνίκ᾽ ἄττα φαίνεται \espTest \textbf{2} (ἀλοῶν) Suidas α 1330 ; Photius α 1029 $=$ \SynB α 984 \espTest \textbf{18-9} (χοάνη = χώνη) Moeris χ 1 ~ Hesychius χ 576 χοάνη \espTest \textbf{18-9} (χοάνη = χώνη) Moeris χ 1 ; cf. Hesychius χ 576 χοάνη \espTest \textbf{24} \Macaire VI 90 ; \Apostol II 94 \espTest \textbf{33} (δασυπώγων) Pollux II 88 \espTest \textbf{36} (ἀλλ᾽ -- πτήξωμεν) Suidas π 3024 \espTest \textbf{39-41} Suidas ε 3791 (λεὼς codd. AF : λαὸς codd. GIVM, συγκλείσας, ἐνδημεῖ pro ἐπιδημεῖ ) \espTest \textbf{40-2} (ἐπιδημεῖ -- μελοποιῶν) Suidas θ 378\par
}
\vspace{2.5pt}\noindent\rule{4.5cm}{0.4pt}\par\vspace{1.5pt}
{\fontsize{6.9pt}{8.2pt}\selectfont
\noindent\textbf{Suite (en commentaire) de l'apparat critique \& notes paléographiques (\texttt{f. 162v}) ---} En haut de la page, \Rdeux{} (Main~B) a copié en semi-onciale de glossaire les six lignes de l'\italique{Hypothesis} ($\text{Εὐριπίδης δράματά τινα... προλογίζει δὲ Μνησίλοχος κηδεστὴς Εὐριπίδου}$) et les deux colonnes des \italique{Dramatis personae}, source unique du nom de \textbf{$\text{Μνησίλοχος}$} défendu par Louis Hauseux (« Le dossier Mnésiloque », \texttt{Introduction.tex}, ll. 325–357). --- Au-dessous, le titre de première main encadré d'ornements en semi-onciale d'apparat : ⁜ ἀριστοφάνους θεσμοφοριάζουσαι ⁜. --- \textbf{Transposition corrigée par lettres-chiffres (vv. 27–30) :} \Rac{} avait copié le v. 28 après le v. 30 ; \Rpc{} a rétabli l'ordre exact en inscrivant à l'extrémité droite des quatre lignes les lettres numérales \textbf{\textbf{α, γ, δ, β}}. --- \textbf{Sigle tachygraphique gratté et scholie au v. 32 :} Comme le relève Louis Hauseux (\texttt{vv 1-641.tex}, ll. 777–787), le manuscrit comportait à l'intérieur du vers, devant οὐχ ἑόρακα πώποτε, le sigle tachygraphique μ\auDessus{ν}{Η} ($\overset{\text{Η}}{\text{μν}}$ = $\text{Μνησίλοχος}$), qui a été ensuite \textbf{gratté (\italique{ut videtur})} ; l'erreur de ce sigle est en outre prouvée par la scholie $\Sigma^{\text{R}}$ du v. 32 (ἐπειδὴ οὐ πάλαι ἤρξατο διδάσκειν, ἀλλὰ πρὸ τριῶν ἐτῶν πρὸ τούτου), qui ne se comprend qu'avec la 2\ieme{} personne ἑόρακας prononcée par Euripide. --- \textbf{Omission de vers réparée en marge (v. 33) :} Le \textbf{v. 33} ($\text{μῶν ὁ δασυπώγων; --- οὐχ ἑόρακας πώποτε}$) avait été omis par saut du même au même dans la colonne par \Rac{} (et reste \textbf{omis dans \Mon{}}) ; il est rajouté sur deux lignes dans la marge droite par \Rpc{}. --- \textbf{Grattage au v. 34 et correction \latin{in rasura} au v. 35 :} Au v.~34, \Rac{} avait écrit καὶ μεγ' (copié tel quel par \Mon{}) ; \Rpc{} a gratté l'iota (ὶ eras.) pour amorcer la crase κἀμέ γ'. Au \textbf{v.~35}, le second $\beta$ de \textbf{βεβίνηκας} est corrigé \latin{in rasura} sur un $\alpha$ gratté (\textbf{βεαίνηκας} \Rac{}), ce qui, joint au silence complet de $\Sigma^{\text{R}}$, motive les \latin{cruces} \lecon{†βεβίνηκας†} de \lecon{Bibi}. --- \textbf{v. 1 :} ἄρα πότε R : \lecon{ἆρά ποτε} Biset \espApp χελιδὼν R : \lecon{Χελιδὼν} \lecon{Bibi} ap. Clements (dans l'\italique{Index nominum} et le commentaire : référence au devin \italique{Khelidôn}, πυθοῦ Χελιδόνος). --- \textbf{v. 3 :} τι R : \lecon{τε} $\Sigma$ (δυνατόν), Souda $\text{οι}\ 150$. --- \textbf{v. 10 :} \lecon{ὁρᾶν.} R (ponctuation haute maintenue par L. Hauseux) : ὁρᾶν; Gelenius. --- \textbf{v. 18 :} \lecon{ἀκοὴν δὲ χοάνης} R, défendu par Alphonse Willems et adopté par L. Hauseux : ἀκοὴς δὲ χοάνην Biset : ἀκοῇ δὲ χοάνην Boissonade. --- \textbf{v. 21 :} τε R : \lecon{γε} $\Sigma$. --- \textbf{v. 23 :} πρὸς τοῖς ἀγαθοῖς τούτοισιν ἐξεύροιμ' ὅπως R : \textbf{vers athétisé (\italique{del.}) par Alphonse Willems} et L. Hauseux comme interpolation issue d'une glose syntaxique. --- \textbf{v. 24 :} προσμάθοι· μὴ R (médivision avec \italique{dikolon} fautif) : \lecon{προσμάθοιμι} Ellebodius. --- \textbf{v. 28 :} ἀκούσω R : \lecon{ἀκούω} Brunck \espApp vers. 28 post 30 transp. \Rac{} : \lecon{corr.} \Rpc{} (litt. $\alpha, \gamma, \delta, \beta$ in mg.). --- \textbf{vv. 30–31 :} Ἀγάθων (bis) R : \lecon{Ἁγάθων} (v. 30, avec crase $\text{ὁ Ἀγάθων}$) et \lecon{ἀγαθῶν} (v. 31, génitif pluriel en calembour) \lecon{Bibi} \espApp Au \textbf{v. 31}, ἔστιν τις ἀγαθῶν \lecon{(Μνη.) continuat} R (maintenu par \textbf{Bibi}) : (Εὐρ.) \Mon{}. --- \textbf{v. 32 :} (Μνη.) (μ\auDessus{ν}{Η}) οὐχ ... πώποτε \Rac{} (sigle gratté \italique{ut vid.}) : \lecon{(Εὐρ.) continuat} $\text{R}^{\text{pc}}, \Sigma$ \espApp ἑόρακα R : \lecon{ἑόρακας} Bentley. --- \textbf{v. 35 :} βε\inRas{α}{β}ίνηκας \Rpc{} (ex βεαίνηκας \Rac{} ut vid.) : \lecon{†βεβίνηκας†} \lecon{Bibi}. --- \textbf{v. 39 :} λέως R, Souda $\varepsilon\ 3792$ codd. AF : \lecon{λαός} $\Sigma$, Souda codd. GIVM.\par
}
\end{minipage}
\end{adjustbox}
\end{adjustwidth}
\newpage
\enTeteD{Codex Ravennas 429 (Biblioteca Classense) --- Fac-similé du f. 162v (vv. 1–42)}
\begin{adjustwidth}{-1.5cm}{-3.0cm}
\begin{center}
\includegraphics[width=16.5cm,height=20.5cm,keepaspectratio]{/workspaces/Thesmophories/dossier_ravennas_thesmophories/folios_individuels_ff162v-177r/01_f162v_classense164v_vv0001-0042.jpg}\par
\vspace{2pt}
{\scriptsize\textbf{Codex Ravennas 429, f. 162v (Classense 164v)} --- ΠΡΟΛΟΓΟΣ, \textbf{vv. 1–42} ($\text{R}^1, \Rpc, \Rdeux$ (\Mon)).}
\end{center}
\end{adjustwidth}
\newpage
\clearpage\ifodd\value{page}\hbox{}\newpage\fi
\enTeteG{ΘΕΣΜΟΦΟΡΙΑΖΟΥΣΑΙ --- Texte grec établi (f. 162v, vv. 1–42)}
\enTeteD{LES THESMOPHORIES --- Traduction française (f. 162v, vv. 1–42)}
\begin{pages}
\begin{Leftside}
\beginnumbering
%% v. 1-4
\begin{groupe}%
\enTeteG{ΠΡΟΛΟΓΟΣ --- Prologue  vv. 1-294}%
%\enTeteD{ΠΡΟΛΟΓΟΣ --- Prologue  vv. 1-294}%
\personnageGrec{Μνησίλοχος}%
\typeVers{$3$ ia.}\trimIamb{ὦ Ζεῦ χελιδὼν ἆρά ποτε φανήσεται;\apparat{1}{ἄρα πότε R : \lecon{ἆρά ποτε} Biset \ligneApp{2} ἄνθρωπος R : \lecon{ἅνθρωπος} Boissonade \ligneApp{3} τι R : \lecon{τε} Σ (δυνατόν), Suid. οι 150}\\%\testimonia{1}{\Fr 617 πυθοῦ χελιδὼν πηνίκ᾽ ἄττα φαίνεται}\\
ἀπολεῖ μ᾽ ἀλοῶν ἅνθρωπος ἐξ ἑωθινοῦ.%\testimonia{2}{(ἀλοῶν) Suidas α 1330 ; Photius α 1029 $=$ \SynB α 984}%
\commentaire{1-10}{ὦ Ζεῦ ... οὐχ ἅ γ᾽ ἂν ἀκούειν δέῃ.}{Nous sommes au tout début de la pièce, il fait encore nuit, un personnage entre d'un pas pressé ; il est immédiatement reconnaissable à son masque bien connu des spectateurs : c'est Euripide\marqueur{1}. Il est suivi d'un vieillard qui traîne la patte et crie sa détresse -- \italique{topos} comique. (On apprendra plus loin qu'il est un parent d'Euripide ; la tradition l'identifie à Mnésiloque, le père de la première femme d'Euripide, Chœrile\marqueur{2}.) Réplique fracassante d'Euripide, et bien selon ses façons -- les fameuses γνῶμαι d'Euripide fort goûtées des spectateurs : ἀλλ᾽ οὐκ ἀκούειν δεῖ σε πάνθ᾽ ὅσ᾽ αὐτίκα ὄψει παρεστώς, « il n'est pas besoin que tu entendes tout ce que tu vas bientôt voir ».
Hors d'haleine, Mnésiloque n'est pas en état de suivre Euripide dans ses aphorismes : il s'est arrêté à la première partie : οὐ δεῖ μ᾽ ἀκούειν, « il ne faut pas que j'entende... » Euripide reprend : οὐχ ἅ γ᾽ ἂν μέλῃς ὁρᾶν, « pas certes ce que tu es sur le point de voir ».\\
C'est là que nous arrivons à l'étrange réponse de Mnésiloque au vers 7b : οὐδ ἆρ ὁρᾶν δεῖ μ᾽, « je ne dois pas non plus voir ? ». Ne comprenant pas pourquoi  Mnésiloque passe d'une défense d'entendre à une défense de voir, l'on est incité à mieux scruter les vers précédents. 
Le scholiaste nous met sur la piste : les constructions tortueuses d'Euripide sont sujettes à l'amphibologie\marqueur{3}.
Mnésiloque commet une erreur (volontaire) d'interprétation sur la dernière réplique d'Euripide : dans `οὐχ ἅ γ ἂν μέλῃς ὁρᾶν', il feint d'entendre `ἄγαν' à la place de `ἅ γ᾽ ἂν'. Aristophane ayant mis à dessein ces particules dans un ordre inhabituel\marqueur{4}.
Il est remarquable de noter que, des onze pièces conservées de notre poète, le cas ne se rencontre qu'une seule autre fois... et c'est encore un pastiche d'un Euripide postiche (celui des \italique{Grenouilles}, cette fois, vers 960) :%
\citComm{ἐξ \textbf{ὧν γ᾽ ἂν} ἐξηλεγχόμην.}%
Mnésiloque croit comprendre que le voir est lui aussi soumis au μηδὲν ἄγαν grec : `οὐκ ἄγαν μέλῃς ὁρᾶν', « tu ne dois pas trop voir »... d'où sa réaction du vers 7b.\\
}%
\debutNote{1-10}%
\note{1}{À quoi pouvait ressembler ce masque d'Euripide ? On peut supposer qu'il présentait une barbe très touffue (cf. le vers 190) ainsi que de multiples taches sur son visage : Ἐλέγετο δὲ καὶ βαθὺν πώγωνα θρέψαι καὶ ἐπὶ τῆς ὄψεως φακοὺς ἐσχηκέναι. (Tiré d'une des \italique{Vies} d'Euripide, rassemblées par Richard Kannicht pour son édition des fragments d'Euripide, cinquième volume des \italique{Tragicorum Graecorum Fragmenta} (2004). Ce détail se trouve à la page 47.)}%
\note{2}{\italique{Confer} l'introduction sur la vie d'Euripide et de Mnésiloque et le vers 291. %\eff{Il n'y a aucune raison de mettre en doute l'existence historique de Mnésiloque ainsi que le lien de parenté qu'il entretenait avec lui. S'il est vrai que les indices internes au texte d'Aristophane sont ténus L'arbre généalogique d'Euripide chez Kannicht concorde avec l'usage d'appeler son premier fils du nom du grand-père paternel et second de celui maternel. Cf. Willems Démosthène ?}
}%
\note{3}{ὁ μὲν τραγικώτερον καὶ ὑψηλότερον φράζει, ὁ δὲ ταπεινότερον ἀκούει, « l'un s'exprime dans un style plus tragique et élevé, quand l'autre comprend de façon plus basse ».}%
\note{4}{Le véritable Euripide s'exprimait ainsi : (\italique{Oreste}, vers 81) : Ἑλένη, τί σοι λέγοιμ᾽ \textbf{ἂν ἅ γε} παροῦσ᾽ ὁρᾷς;}\\
παρὰ σοῦ πυθέσθαι ποῖ μ᾽ ἄγεις ωὖριπίδη;%
\commentaire{1-4}{ὦ Ζεῦ ... πρὶν τὸν σπλῆνα κομιδῇ μ᾽ ἐκβαλεῖν.}{Ce quatrain est analogue à celui des vers 71-74. À deux reprises, Mnésiloque adresse une prière plaintive à Zeus avant de questionner Euripide. Ressemblance saisissante aussi avec les premiers vers des \italique{Nuées} (vv. 1-5) :%
\citComm{Στρεψιάδης --- ἰοὺ ἰού·\\
ὦ Ζεῦ βασιλεῦ τὸ χρῆμα τῶν νυκτῶν ὅσον·\\
ἀπέραντον. οὐδέποθ᾽ ἡμέρα γενήσεται;\\
καὶ μὴν πάλαι γ᾽ ἀλεκτρυόνος ἤκουσ᾽ ἐγώ·\\
οἱ δ᾽ οἰκέται ῥέγκουσιν· ἀλλ᾽ οὐκ ἂν πρὸ τοῦ.\\
Strepsiade --- Iou ! iou !\\
Ah ! Zeus roi, quel fardeau que ces nuits !\\
C'est interminable... Il ne fera donc jamais jour ?\\
Pourtant voilà longtemps que j'ai entendu le coq,\\
et mes serviteurs ronflent. Ah ! que non pas autrefois !}
Cette analogie permet de saisir le sens imagé de l'hirondelle (χελιδών) de notre passage : elle est mise à la place du coq des \italique{Nuées} pour ce qu'elle a de semblable : en chantant dès avant l'aube, elle annonce le jour bientôt renaissant\marqueur{1}. Que l'hirondelle était réputée chez les Grecs chanter\marqueur{2} au matin, cela est appuyé par l'adjectif ὀρθρογόη, \italique{aux lamentations matinales}, qu'Hésiode (\italique{Opera et Dies}, v. 568) applique à l'hirondelle\marqueur{3} :%
\citComm{ὀρθρογόη Πανδιονὶς ὦρτο χελιδὼν\\
La fille matinale et gémissante de Pandion, l'hirondelle, réveille...}%
Ce strident réveille-matin avait donné lieu à un expression proverbiale dont Aristophane fait mention ailleurs\marqueur{4}:%
\citComm{πυθοῦ χελιδὼν πηνίκ᾽ ἄττα φαίνεται\\
Renseigne-toi à quelle heure l'hirondelle apparaît.}%
Ici, ce n'est pas une hirondelle mais Euripide qui a jeté hors du lit Mnésiloque bien trop tôt.. Ce sera donc lui qui sera interrogé : \textbf{παρὰ σοῦ} πυθέσθαι (v. 3).}%
\debutNote{1-4}%
\note{1}{Et non pas le printemps -- qui n'est pas de saison ici --, comme on peut le lire drôlement dans nombre de commentaires.}%
\note{2}{Chanter ?.. ou plutôt babiller : l'hirondelle était proverbialement brocardée pour son caquet, cf. notamment Taillardat §§520-522 et les notes d'Alphonse Willems à propos du vers 682 des \italique{Grenouilles}.}%
\note{3}{Voir aussi Hésychius (ο 1189) : ὀρθ<ρ>ογόη· ἡ χελιδών.}%
\note{4}{Fragment 617 Kassel-Austin. Diversement recensé par les parémiographes (souvent sous la forme πυθοῦ Χελιδόνος) : Hésychius π 4311 ; Photius π 1519 ; Suidas π 3136 ; \Apostol XV, 14 ; \&c.}\\
οἷόν τε, πρὶν τὸν σπλῆνα κομιδῇ μ᾽ ἐκβαλεῖν.%
\commentaire{2}{ἀλοῶν.}{Ἀλοᾶν est un terme agricole signifiant \italique{battre le grain}, selon une technique bien définie : sur l'aire à battre (ἡ ἅλως), on fait tourner des bœufs, attachés à un piquet, qui piétinent et qui broient la céréale jonchant le sol. De ce fait, ἁλοᾶν est devenu un synonyme expressif de \grec{τύπτειν} avec le sens de \italique{broyer de coups}. (Note de Taillardat\marqueur{1})\\
Ici, ἁλοᾶν a le sens moral de \italique{tourmenter}. Conformité du grec et du français -- ou du latin, plutôt : l'on possède la même image en avec notre \italique{tribulation}, le latin \latin{tribulare} signifie « battre le blé, presser avec la herse (\latin{tribulum}) ».}%
\debutNote{2}%
\note{1}{Je me référerai souvent – et avec force citations ! – au magistral ouvrage que Jean Taillardat a consacré à la langue bigarrée et haute en couleurs d'Aristophane : \italique{Les images d'Aristophane} édité aux Belles Lettres (seconde édition de 1965). Je cite ici le §138 ; voir aussi les références du §606 : la métaphore était usuelle et pas de l'invention de notre poète : \italique{Grenouilles} 149, Sophocle fr. 19 \Nauck [$=$ fr. 20 Radt (dub.)] ; Hérondas 2, 34 ; Babrius, \italique{Le lion amoureux}, 98.15 ; καταλοᾶν Eschine 2, 140 ; Xénoph. \italique{Cyrop.} 7, 1, 31 ; Lucien \italique{Icarom.} 15 ; συναλοᾶν Hippon. 14AD\expo{3}.4 [$=$ fr. 92.4 Masson $=$ fr 95.4 Degani] ; Théocr. 22.128.}%
\\\passeLigne}%
\end{groupe}%
%% v. 5-6a
\begin{groupe}%
\personnageGrec{Εὐριπίδης}%
\trimIamb{ἀλλ᾽ οὐκ ἀκούειν δεῖ σε πάνθ᾽ ὅσ᾽ αὐτίκα\apparat{5}{ὅσ R : \lecon{ὅσ᾽} \Mon}\\
ὄψει παρεστώς.%
\commentaire{3}{πρὶν τὸν σπλῆνα κομιδῇ μ᾽ ἐκβαλεῖν.}{Les dérangements (palpitations) dus à la course – d'autant plus aggravés par l'âge avancé de Mnésiloque – étaient imputés par les Anciens à des troubles de la rate (σπλήν, \anglais{spleen} en anglais)\marqueur{1}. En français, celui qui court ventre à terre court aussi comme un « dé-raté ».%\eff{ventre à terre, rendre les boyaux}
}%
\debutNote{3}%
\note{1}{Voir Pline l'Ancien, \italique{Histoire naturelle} 11,205 ; Plaute, \italique{Mercator}, v. 124.}\\\passeLigne}%
\end{groupe}%
%% v. 6b-7a
\begin{groupe}%
\personnageGrec{Μνησίλοχος}%
\trimIamb{\ligne{6}\decalage{ὄψει παρεστώς.}πῶς λέγεις; αὖθις φράσον.\\
%\lednopb%
οὐ δεῖ μ᾽ ἀκούειν;\\\passeLigne}%
\end{groupe}%
%% v. 7b
\begin{groupe}%
\personnageGrec{Εὐριπίδης}%
\trimIamb{\decalage{οὐ δεῖ μ᾽ ἀκούειν;}\ligne{7}οὐχ ἅ γ᾽ ἂν μέλλῃς ὁρᾶν.\\\passeLigne}%
\end{groupe}%
%% v. 8a
\begin{groupe}%
\personnageGrec{Μνησίλοχος}%
\trimIamb{οὐδ᾽ ἆρ᾽ ὁρᾶν δεῖ μ᾽;\\\passeLigne}%
\end{groupe}%
%% v. 8b
\begin{groupe}%
\personnageGrec{Εὐριπίδης}%
\trimIamb{\decalage{οὐδ᾽ ἆρ᾽ ὁρᾶν δεῖ μ᾽;}\ligne{8}οὐχ ἅ γ᾽ ἂν ἀκούειν δέῃ.\\\passeLigne}%
\end{groupe}%
%% v. 9-10
\begin{groupe}%
\personnageGrec{Μνησίλοχος}%
\trimIamb{πῶς μοι παραινεῖς; δεξιῶς μέντοι λέγεις.\\
οὐ φῂς σὺ χρῆναί μ᾽ οὔτ᾽ ἀκούειν οὔθ᾽ ὁρᾶν.\apparat{10}{\lecon{ὁρᾶν.} R : ὁρᾶν; Gelenius}\\\passeLigne}%
\end{groupe}%
%% v. 11
\begin{groupe}%
\personnageGrec{Εὐριπίδης}%
\trimIamb{χωρὶς γὰρ αὐτοῖν ἑκατέρου 'στὶν ἡ φύσις.\\\passeLigne}%
\end{groupe}%
%% v. 12a
\begin{groupe}%
\personnageGrec{Μνησίλοχος}%
\trimIamb{τοῦ μήτ᾽ ἀκούειν μήθ᾽ ὁρᾶν;\\\passeLigne}%
\end{groupe}%
%% v. 12b
\begin{groupe}%
\personnageGrec{Εὐριπίδης}%
\trimIamb{\ligne{12}\decalage{τοῦ μήτ᾽ ἀκούειν μήθ᾽ ὁρᾶν;}εὖ ἴσθ᾽ ὅτι.\\\passeLigne}%
\end{groupe}%
%% v. 13a
\begin{groupe}%
\personnageGrec{Μνησίλοχος}%
\trimIamb{πῶς χωρίς;\\\passeLigne}%
\end{groupe}%
%% v. 13b-18
\begin{groupe}%
\personnageGrec{Εὐριπίδης}%
\trimIamb{\ligne{13}\decalage{πῶς χωρίς;}οὕτω ταῦτα διεκρίθη τότε.\\
αἰθὴρ γὰρ ὅτε τὰ πρῶτα διεχωρίζετο\\
καὶ ζᾦ᾽ ἐν αὑτῷ ξυνετέκνου κινούμενα,\\
ᾧ μὲν βλέπειν χρὴ πρῶτ᾽ ἐμηχανήσατο\apparat{16}{πρῶτα R : \lecon{πρῶτ᾽} Grynaeus}\\
ὀφθαλμὸν ἀντίμιμον ἡλίου τροχῷ,\\
ἀκοὴν δὲ χοάνης ὦτα διετετρήνατο.%\testimonia{18-9}{(χοάνη = χώνη) Moeris χ 1 ~ Hesychius χ 576 χοάνη}
\apparat{18}{\lecon{ἀκοὴν δὲ χοάνης} R, Willems : ἀκοὴς δὲ χοάνην Biset : ἀκοῇ δὲ χοάνην Boissonade}%
\commentaire{13-8}{}{Aristophane nous peint Euripide comme un de ces faiseurs de cosmogonies (tel le Socrate des \italique{Nuées} avec son Δῖνος). Le véritable Euripide avait effectivement un penchant pour ce genre de nuées\marqueur{1}. En l'occurrence, le pastiche de cosmogonie qu'Aristophane met dans la bouche d'Euripide semble inspiré d'Empédocle\marqueur{2}.}%
\debutNote{13-8}%
\note{1}{Voir par exemple le fragment 5 de \italique{Mélanippe la sage} dans l'édition de Jouan-van Looy aux Belles Lettres, tome VIII, \nieme{2} partie ($=$ fragment 484 de l'édition Kannicht).}%
\note{2}{Cf. notamment ce passage d'Eusèbe de Césarée (\italique{Préparation évangélique}, XV, 37 ; fragment A49 Diels-Kranz) : Ἐμπεδοκλῆς τὸν μὲν αἰθέρα πρῶτον διακριθῆναι et la création de l'œil par Aphrodite (B84.9 D.-K.) reconstitué par Blass (1883) : χοάνῃσι δίαντα τετρήατο θεσπεσίῃσιν.}%
\commentaire{16-8}{πρῶτα ἐμηχανήσατο ... ἀκοὴν δὲ χοάνης ὦτα διετετρήνατο.}{Pour l'explication de ce passage diversement retouché au gré des éditeurs, je ne peux que renvoyer au brillant exercice d'exégèse d'Alphonse Willems. (Ses notes sur les \italique{Thesmophories} sont reproduites en annexe, celle relative à ce passage en est la première. À propos d'Alphonse Willems, \italique{confer} l'Introduction.) Euripide recourt à une image tirée de la fonderie -- certainement plus parlante aux spectateurs de l'époque qu'à nous : l'oreille (ἀκοή) est comparée aux ὦτα (littéralement les \italique{oreilles} et, par analogie, les \italique{anses}, \italique{poignées} ou, comme ici, les \italique{godets} ou \italique{entonnoirs} d'un moule, χοάνη). Willems construit ainsi la phrase, en accord avec la sobre indication du scholiaste (λείπει ὡς) : διετετρήνατο δὲ ἀκοὴν (ὡς) χοάνης ὦτα.}%
\\\passeLigne}%
\end{groupe}%
%% v. 19-21
\begin{groupe}%
\personnageGrec{Μνησίλοχος}%
\trimIamb{διὰ τὴν χοάνην οὖν μήτ᾽ ἀκούω μήθ᾽ ὁρῶ;
\commentaire{19}{διὰ τὴν χοάνην οὖν μήτ᾽ ἀκούω μήθ᾽ ὁρῶ;}{Mnésiloque ne goûte guère aux subtilités des poètereaux et tord à plaisir leurs propos. Ses saillies fort désobligeantes se repèrent aisément puisqu'elles sont toujours superbement ignorées. C'en est ici le premier exemple ; Euripide le-faiseur-de-cosmogonies euripidise en lui racontant des histoires de fonderie et de moule... conclusion de Mnésiloque pour expliquer son accablement : « C'est donc à cause de ton moule que je suis moulu ? », διὰ τὴν χοάνην οὖν μήτ᾽ ἀκούω μήθ᾽ ὁρῶ\marqueur{1}. C'est bien le sens qu'il faut donner au tour `μήτ᾽ ἀκούω μήθ᾽ ὁρῶ' et qu'on peut rapprocher du latin \citComm{\latin{Nec aures me credo habere nec tango},\\
Je crois n'avoir plus ni oreilles ni main,} qui signifie justement qu'on ne se sent plus la force de faire quoi que ce soit\marqueur{2}.\\
%il est évident que Mnésiloque se livre ici à un détournement, instantanément intelligible de l'Athénien du \nieme{V} siècle av. J.-C... mais qui nous est bien plus difficile à saisir.
%On peut toutefois proposer une hypothèse.
%Le tour `μήτ᾽ ἀκούω μήθ᾽ ὁρῶ'\marqueur{1} nous semble à rapprocher de l'expression proverbiale latine : \latin{Nec aures me credo habere nec tango}. 
La langue d'Aristophane, comme nous l'allons voir en chemin, est truffée d'échos aux proverbes grecs ; on ne peut que regretter le manque d'intérêt criant des commentateurs d'Aristophane pour ce domaine. Nous nous efforcerons au contraire, et selon nos moyens, de souligner chacun des emprunts d'Aristophane à cette langue proverbiale. Et ils sont nombreux ! Le cas présent est malheureusement moins probant puisque je n'ai pas pu trouver d'autres occurrences dans la littérature grecque. (Surtout : le jeu de mots exact en grec, qui tourne autour de χοάνη, nous échappe.)
%Revenons au proverbe latin\marqueur{2}, \latin{Nec aures me credo habere nec tango}, «~je crois n'avoir plus ni oreilles ni main~» ; il est employé pour signifier qu'on ne se sent plus la force de faire quoi que ce soit. Si donc Mnésiloque est complètement \italique{moulu}, c'est d'après les explications d'Euripide, la faute à ce \italique{moule} (διὰ τὴν χοάνην). Le jeu de mots exact en grec, qui tourne autour de χοάνη, nous échappe.
%\eff{Nοῦς ὁρᾷ, καὶ νοῦς ἀκούει: ?Érasme 2500 ; n° 1971 du Dictionnaire des sentences latines et grecques.}
%L'explication qui nous semble la plus simple est un jeu de mots sur les différents sens de διά + accusatif.\\
%Le premier : \italique{à cause de} ; si donc Mnésiloque ne peut ni voir ni entendre, c'est \italique{à cause de} cette fichue χοάνη... conclusion apparemment logique du long exposé justificatif d'Euripide.\\
%Second sens possible de διά + accusatif : \italique{à travers}. La χοάνη est à prendre dans un sens littéral (de \italique{tamis} ou de \italique{filtre} ?), opaque à la lumière et au son ; le vers suivant est d'une insolente ironie pour Euripide (il fallait au moins un Euripide pour parvenir à la démonstration qu'on ne peut ni voir ni entendre \italique{à travers} un tel objet !)
}%
\debutNote{19}%
\note{1}{Ou bien : « Je ne vois plus Claire ; j'ai perdu Louis. » (traduction plus libre.)}%
\note{2}{Je renvoie le lecteur à l'admirable \italique{Dictionnaire des sentences latines et grecques} de Renzo Tosi (traduit de l'italien par Rebecca Lenoir ; proverbe n° 1971, cité par saint Jérôme notamment). Érasme avait recensé le proverbe dans ses \italique{Adages}, n° 2500. (Je cite les \italique{Adages} ainsi que leur traduction française selon l'édition de Jean-Christophe Saladin aux Belles Lettres (2019).) Toujours dans le §1971, Tosi cite aussi la sentence médiévale \latin{Nec audio nec video}, qui se superpose parfaitement au grec μήτ᾽ ἀκούω μήθ᾽ ὁρῶ.}\\
νὴ τὸν Δί᾽ ἥδομαί γε τουτὶ προσμαθών.\\
\lednopb%
οἷόν γέ που 'στιν αἱ σοφαὶ ξυνουσίαι.\apparat{21}{τε R : \lecon{γε} Σ}%
\commentaire{21}{οἷόν τέ που 'στιν αἱ σοφαὶ ξυνουσίαι.}{Le scholiaste note à ce vers une allusion à un mot célèbre : σοφοὶ τύραννοι τῶν σοφῶν ξυνουσίᾳ, qu'Aristophane dans ses \italique{Héros}\marqueur{1} mais aussi Platon\marqueur{2} et Anthistène\marqueur{3} attribuent à Euripide. Cela n'empêche pas moins notre scholiaste, bien sûr de lui, d'accorder la paternité de cette formule à Sophocle\marqueur{4}. Toujours est-il qu'Euripide n'en passait certainement pas moins aux yeux d'Aristophane et de ses spectateurs pour le père putatif de la citation. Sur le sens de γέ που, on lira avec intérêt la seconde note de Willems reproduite ci-après en annexe.}%
\debutNote{21}%
\note{1}{Fragment 323 de l'édition Kassel-Austin.}%
\note{2}{\italique{Théagès}, 125b et la \italique{République}, VIII, 568b. Platon devait avoir un point de vue sur Euripide assez proche de celui d'Aristophane : il chasse joyeusement de sa République idéale Euripide, cet homme qui occupe une place prestigieuse parmi les semi-habiles de sa cité.}%
\note{3}{Fragment 59 de l'édition Decleva-Caizzi}%
\note{4}{(Dans son Αἴας Λοκρός.) Fr. 14 de l'édition de Stefan Radt, \italique{Tragicorum Graecorum Fragmenta}, vol. 4, édition corrigée et augmentée de 1999. Il figure aussi parmi les fragments incertains d'Euripide (n° 888b) dans l'édition Kannicht des \italique{Tragicorum Graecorum Fragmenta}, vol 5. % \eff{chez Jouan - van Looy ?}
Le ton de la scholie est d'une telle assurance doublée d'une forte érudition qu'on serait tenté de l'en croire. Preuve que la scholie est de Didyme, selon le critère de Pierre Boudreaux ? (\italique{Le texte d'Aristophane et ses commentateurs} (1920). Boudreaux, brillant helléniste fauché bien jeune par la Grande guerre, réussit le tour de force d'écrire un bouquin -- publié de façon posthume -- passionnant quoique traitant de sujets à priori fort arides.)
Après tout, Euripide est coutumier du fait : Aulu-Gelle (\italique{Nuits attiques}, XIII, 18) nous rapporte entre autres exemples que le vers 193 des \italique{Bacchantes}, γέρων γέροντα παιδαγωγήσω σ᾽ ἐγώ, se retrouve aussi à l'identique chez Sophocle (fr. 695 Radt) ; de même, Aristote (Περὶ ποιητικῆς, XIII, 1458b22) cite une reprise d'Eschyle par Euripide.
À rebours de leur querelle dans les \italique{Grenouilles}, Aristote cite un vers iambique d'Eschyle jugé plat en raison du recours à un mot courant tandis qu'Euripide le reprit en l'agrémentant d'un mot emprunté et rare : οἷον τὸ αὐτὸ ποιήσαντος ἰαμβεῖον Αἰσχύλου καὶ Εὐριπίδου, ἓν δὲ μόνον ὄνομα μεταθέντος, ἀντὶ κυρίου εἰωθότος γλῶτταν, τὸ μὲν φαίνεται καλὸν τὸ δ᾽ εὐτελές. Αἰσχύλος μὲν γὰρ ἐν τῷ Φιλοκτήτῃ ἐποίησε·\textbf{φαγέδαιναν ἥ μου σάρκας ἐσθίει ποδός,} ὁ δὲ ἀντὶ τοῦ `ἐσθίει' τὸ `θοινᾶται' μετέθηκεν.
}%
\\\passeLigne}%
\end{groupe}%
%% v. 22a
\begin{groupe}%
\personnageGrec{Εὐριπίδης}%
\trimIamb{πόλλ᾽ ἂν μάθοις τοιαῦτα παρ᾽ ἐμοῦ.\\\passeLigne}%
\end{groupe}%
%% v. 22b-24
\begin{groupe}%
\personnageGrec{Μνησίλοχος}%
\trimIamb{\ligne{23}\decalage{πόλλ᾽ ἂν μάθοις τοιαῦτα παρ᾽ ἐμοῦ.}πῶς ἂν οὖν\apparat{v. 23}{πρὸς τοῖς ἀγαθοῖς τούτοισιν ἐξεύροιμ᾽ ὅπως R : \lecon{del.} Willems \ligneApp{24} προσμάθοι· μὴ R : \lecon{προσμάθοιμι} Ellebodius}\\
ἔτι προσμάθοιμι χωλὸς εἶναι τὼ σκέλει;%\testimonia{24}{\Macaire VI 90 ; \Apostol II 94}%\apparat{24}{προσμάθοι· μὴ R : \lecon{προσμάθοιμι} Ellebodius}%
\commentaire{22-4}{πῶς ἂν οὖν ... προσμάθοιμι χωλὸς εἶναι τὼ σκέλει;}{Aristophane joue sur un détournement de proverbe équivalent à notre «~Qui fréquente le boiteux apprend à boiter~» : Παρὰ χωλὸν οἰκῶν κἂν ἐπισκάζειν μαθήσῃ\marqueur{1}.
Le boiteux -- ou plutôt le \italique{faiseur de boiteux}, χωλοποιός\marqueur{2} -- c'est bien-sûr Euripide et sa fâcheuse prédilection pour ce genre de personnages\marqueur{3}.
La pointe de la plaisanterie de Mnésiloque est rejetée à la toute fin ; quand il précise vouloir être boiteux τὼ σκέλει, \italique{des deux jambes}. 
Le scholiaste a raison : %
\citComm{χωλὸς εἶναι· ἵνα μὴ περιπατήσῃ μετ᾽ αὐτοῦ.}%
Si Mnésiloque venait à boiter des deux jambes, il ne serait plus forcé de suivre Euripide. Le grec aime à préciser que l'on est rompu \italique{des deux jambes} pour souligner une fatigue intense\marqueur{4}.
(Voir là-dessus la \nieme{3} note d'Alphonse Willems, reproduite en annexe, qui justifie la suppression du vers 23 -- une glose explicative du suivant.)
Mais Euripide ne s'en laisse pas compter et lui ordonne de lui emboîter le pas (βάδιζε δευρὶ).\\ 
%De simples allusions suffisaient au public d'Aristophane, mais une moderne traduction devrait peut-être insister davantage et le rendre ainsi : «~Euripide, toi qui as une prédilection pour les éclopés, tu pourrais me faire deux belles jambes en coupant court à tous ces va-et-vient~», ou bien : «~m'apprendre à te désemboîter le pas~».\\
%L'autre explication est celle du scholiaste : si Mnésiloque venait à boiter des deux jambes, il ne serait plus forcé de suivre Euripide.\\
%Ma préférence va évidemment à la première hypothèse... que je ne suis malheureusement pas parvenu à étayer : ni pour montrer que Mnésiloque soit boiteux, ni que "boiter des deux jambes" pût signifier pour un Grec "marcher droit" ; je m'en remets donc à toi, ô lecteur : puisses-tu m'apporter quelque éclaircissement !\\
}%
\debutNote{22-4}%
\note{1}{C'est le proverbe n° 1674 du \italique{Dictionnaire des sentences latines \& grecques}. Recensé par Macaire au \nieme{XIV} siècle dans son anthologie Ῥοδωνιά (VI, 90 de l'édition de Leutsch \& Schneidewim du \italique{Corpus Paroemiographorum Graecorum}, vol. II) ainsi que sous d'autres formes proches par Michel Apostolius, Plutarque et un scholiaste de Pindare.}%
\note{2}{\italique{Les Grenouilles}, vers 846.}%
\note{3}{Ou autres aussi peu reluisants. Cf. les \italique{Acharniens} v. 411 sq. : Télèphe, Bellérophon, Philoctète le mendiant ; voir aussi la \italique{Paix} v. 147 ; les \italique{Oiseaux} v. 1293 et le fragment 53 de l'édition Kassel-Austin des \italique{Poetae Comici Graeci}, vol. III 2 (1984).}%
\note{4}{Voir par exemple la \italique{Paix} et la douloureuse ascension de l'orchestre à la scène de Trygée : ἔγωγέ τοι πεπόνηκα κομιδῇ τὼ σκέλει (v. 820) et encore : ἤλγουν τὼ σκέλει μακρὰν ὁδὸν διεληλυθώς(v. 825). 
Nous ne sommes pas dans la chanson de Nadaud où souhaiter claudiquer des deux côtés, c'est souhaiter ne plus claudiquer :%
\citComm{Si je marche plus droit qu'un autre\\
C'est que je boite des deux pieds.}}\\\passeLigne}%
\end{groupe}%
%% v. 25a
\begin{groupe}%
\personnageGrec{Εὐριπίδης}%
\trimIamb{\ligne{25}βάδιζε δευρὶ καὶ πρόσεχε τὸν νοῦν.\\\passeLigne}%
\end{groupe}%
%% v. 25b
\begin{groupe}%
\personnageGrec{Μνησίλοχος}%
\trimIamb{\ligne{25}\decalage{βάδιζε δευρὶ καὶ πρόσεχε τὸν νοῦν.}ἰδού.\\\passeLigne}%
\end{groupe}%
%% v. 26a
\begin{groupe}%
\personnageGrec{Εὐριπίδης}%
\trimIamb{ὁρᾷς τὸ θύριον τοῦτο;\\\passeLigne}%
\end{groupe}%
%% v. 26b-27a
\begin{groupe}%
\personnageGrec{Μνησίλοχος}%
\trimIamb{\ligne{26}\decalage{ὁρᾷς τὸ θύριον τοῦτο;}νὴ τὸν Ἡρακλέα\\
οἶμαί γε.\\\passeLigne}%
\end{groupe}%
%% v. 27b
\begin{groupe}%
\personnageGrec{Εὐριπίδης}%
\trimIamb{\ligne{27}\decalage{οἶμαί γε.}σίγα νυν.\apparat{27}{νῦν R : \lecon{νυν} Boissonade}\\\passeLigne}%
\end{groupe}%
%% v. 27c
\begin{groupe}%
\personnageGrec{Μνησίλοχος}%
\trimIamb{\ligne{27}\decalage{οἶμαί γε.σίγα νυν.}σιωπῶ τὸ θύριον.%\apparat{27}{θύριον R : \lecon{θύριον;} Dindorf}%
%\commentaire{27}{σιωπῶ τὸ θύριον}{\eff{Lire σιω πῶ τὸ θύριον; ? du verbe σείω ? puis a-krouw ?}}
\commentaire{27-8}{σίγα νυν --- σιωπῶ τὸ θύριον --- ἄκου᾽ --- ἀκούω καὶ σιωπῶ τὸ θύριον.}{Le trait d'esprit de Mnésiloque nous échappe.}%
\\\passeLigne}%
\end{groupe}%
%% v. 28a
\begin{groupe}%
\personnageGrec{Εὐριπίδης}%
\trimIamb{ἄκου᾽.\\\passeLigne}%
\end{groupe}%
%% v. 28b
\begin{groupe}%
\personnageGrec{Μνησίλοχος}%
\trimIamb{\ligne{28}\decalage{ἄκου᾽.}ἀκούω καὶ σιωπῶ τὸ θύριον.\apparat{28}{ἀκούσω R : \lecon{ἀκούω} Brunck \espApp vers. 28 post 30 transp. \Rac : \lecon{corr.} \Rpc }\\\passeLigne}%
\end{groupe}%
%% v. 29-30a
\begin{groupe}%
\personnageGrec{Εὐριπίδης}%
\trimIamb{ἐνταῦθ᾽ Ἀγάθων ὁ κλεινὸς οἰκῶν τυγχάνει
\commentaire{29}{Ἀγάθων.}{Voici enfin la réponse à la question de Mnésiloque (v. 4) sur la destination de leur équipée : la demeure d'Agathon, le fameux poète tragique\marqueur{1} – sauf pour Mnésiloque –, rendu célèbre par sa victoire au concours tragique des Lénéennes de 416\marqueur{2}. Platon immortalisa dans son \italique{Banquet} la réception qu'il fit chez lui le lendemain de cette victoire. Le rôle d'Agathon dans les \italique{Thesmophories} est semblable à celui d'Euripide dans les \italique{Acharniens} (v. 410-489) : il fournit le héros en matériel de travestissement. Euripide confirmera cette parenté aux vers 173-174\marqueur{3}. Après la mort d'Agathon, Aristophane lui rendit un hommage posthume dans les \italique{Grenouilles} aux vers 83-85.}%
\debutNote{29}%
\note{1}{On n'a conservé d'Agathon que des fragments et quelques témoignages. Ils ont été regroupés par Snell \& Kannicht au premier volume des \italique{Tragicorum Graecorum Fragmenta} (1986).}
\note{2}{Athénée, \italique{Deipnosophistes}, V, 216f-217a.}%
\note{3}{Aristophane n'est d'ailleurs pas le seul à rapprocher Euripide \& Agathon : ainsi fait Aristote dans sa \italique{Poétique} (1456a25-32) : Agathon introduisit des interludes (ἐμβόλιμα) chantés par le chœur sans rapport avec l'action, suivant la voie tracée par Euripide.% : %
%\citComm{καὶ τὸν χορὸν δὲ ἕνα δεῖ ὑπολαμβάνειν τῶν ὑποκριτῶν, καὶ μόριον εἶναι τοῦ ὅλου καὶ συναγωνίζεσθαι μὴ ὥσπερ Εὐριπίδῃ ἀλλ᾽ ὥσπερ Σοφοκλεῖ. τοῖς δὲ λοιποῖς τὰ ᾀδόμενα <οὐδὲν> μᾶλλον τοῦ μύθου ἢ ἄλλης τραγῳδίας ἐστίν: διὸ ἐμβόλιμα ᾁδουσιν πρώτου ἄρξαντος Ἀγάθωνος τοῦ τοιούτου. καίτοι τί διαφέρει ἢ ἐμβόλιμα ᾁδειν ἢ εἰ ῥῆσιν ἐξ ἄλλου εἰς ἄλλο ἁρμόττοι ἢ ἐπεισόδιον ὅλον;\\
%Le chœur doit être considéré comme l'un des acteurs ; il doit faire partie de l'ensemble et participer à l'action, non comme chez Euripide mais comme chez Sophocle. Chez tous les autres, les parties chantées n'ont pas plus de rapport avec l'histoire qu'avec une autre tragédie ; aussi chantent-ils des interludes, pratique dont l'origine remonte à Agathon. Et pourtant quelle différence y a-t-il entre chanter des interludes et adapter à une pièce une tirade ou un épisode tout entier tirés d'un autre ?}
}\\
ὁ τραγῳδοποιός.\\\passeLigne}%
\end{groupe}%
%% v. 30b
\begin{groupe}%
\personnageGrec{Μνησίλοχος}%
\trimIamb{\ligne{30}\decalage{ὁ τραγῳδοποιός.}ποῖος οὗτος Ἁγάθων;\apparat{30-1}{Ἀγάθων (bis) R : \lecon{Ἁγάθων} et \lecon{ἀγαθῶν} Bibi \ligneApp{31} ἔστιν τις ἀγαθῶν \lecon{(Μνη.) continuat} R : (Εὐρ.) \Mon}\\% \espApp Ἀγάθων R : \lecon{ἀγαθῶν;} Bibi}\\
ἔστιν τις ἀγαθῶν; μῶν ὁ μέλας ὁ καρτερός;%\apparat{31}{ἔστιν τις ἀγαθῶν \lecon{(Μνη.) continuat} R : (Εὐρ.) \Mon \espApp Ἀγάθων R : \lecon{ἀγαθῶν;} Bibi}%
\commentaire{30-3}{ποῖος οὗτος Ἁγάθων; ... οὐχ ἑόρακας πώποτε.}{La répartition des répliques de ce passage fluctue d'une édition à l'autre. %, hésitations du manuscrit de Ravenne lui-même : il comporte au niveau de ces vers de nombreuses corrections. 
Tâchons d'y voir plus clair !\\
À l'instar du manuscrit de Ravenne, nous laissons à Mnésiloque le début du vers 31 : ἔστιν τις Ἀγάθων; 
La clef de compréhension est la suivante : lisons :%
%La première erreur de répartition fut introduite par le scribe du manuscrit de Munich (copie du \italique{Ravennas} du \nieme{XV} siècle) et reprise depuis par la plupart des éditions : il attribua à Euripide le début du vers 31 : ἔστιν τις Ἀγάθων ---, avant que Mnésiloque n'enchaîne en demandant des précisions : μῶν ὁ μέλας ὁ καρτερός;\\
%Laisser le vers entier ἔστιν τις Ἀγάθων μῶν ὁ μέλας ὁ καρτερός; dans la bouche de Mnésiloque relevait d'une construction douteuse... mais il y a une clef de compréhension beaucoup plus simple et qui dispense de découper arbitrairement les vers (d'autant qu'il n'est pas vraisemblable qu'Euripide se laisse ainsi couper le sifflet par Mnésiloque) : il suffit de lire :%
\citComm{ἔστιν τις ἀγαθῶν;}%
Mnésiloque joue sur le nom d'Agathon ; \latin{nomen omen}, demande-t-il : Agathon est-il bien l'\italique{honnête homme} que suggère son nom ? Or, pour une nature rustre comme celle de Mnésiloque, un bon-homme se doit d'être un solide gaillard, à la peau bronzée par le travail en extérieur et non pas au teint fromage-blanc comme ces désœuvrés de fémoïdes ! 
Il continue sur sa lancée : μῶν ὁ μέλας ὁ καρτερός; « n'est-ce pas le noiraud, le costaud ? »
(Agathon est en réalité tout le contraire d'un rustaud ; Aristophane donne ainsi facilement à rire à son public aux dépens de l'original bien connu.)\\
Le vers suivant est écrit ainsi dans le manuscrit de Ravenne :%
\citComm{οὔκ ἀλλ᾽ ἕτερός τις --- οὐχ \textbf{ἑόρακα} πώποτε}%
Le signe mis pour signifier un changement d'interlocuteur\marqueur{1} ayant par la suite été gratté (\latin{ut videtur}).
L'entièreté du vers est effectivement à restituer à Euripide : une scholie écrite juste au-dessus du vers l'atteste : %que telle était ainsi une ancienne leçon : 
\citComm{\textbf{ἐπειδὴ} οὐ πάλαι ἤρξατο διδάσκειν, ἀλλὰ τρισὶν πρὸ τούτων ἔτεσιν.\\
\textbf{Puisqu}'il n'a pas commencé récemment de faire jouer des pièces, mais depuis trois ans. [sept, en réalité, et selon la façon grecque de compter]}
Cette scholie est évidemment là pour expliquer l'étonnement d'Euripide : οὐχ ἑόρακα\textbf{ς} πώποτε; « Tu ne l'as donc jamais vu ?! [alors que cela fait trois année qu'il monte des pièces] » ; le scholiaste ne l'aurait pas ainsi rédigée (en commençant par ἐπειδὴ...) si, au vers 32, le `οὐχ ἑόρακα πώποτε' était prononcé par Mnésiloque, « je ne l'ai jamais vu ».\\
Cette interprétation, outre qu'elle est anciennement corroborée par la scholie, est aussi celle qui permet la lecture la plus drôle : οὐχ ἑόρακας πώποτε; --- μῶν ὁ δασυπώγων; --- οὐχ ἑόρακας πώποτε, « Tu ne l'as jamais vu ?! --- N'est-il pas barbu ? --- Tu ne l'as jamais vu !.. », bien meilleure par le rythme et la répétition déconvenue d'Euripide que l'échange premièrement transcrit sur le manuscrit de Ravenne où Mnésiloque commençait par avouer son ignorance : (Μνη.) οὐχ ἑόρακα πώποτε. μῶν ὁ δασυπώγων; --- (Εὐρ.) οὐχ ἑόρακας πώποτε.%
}%
\debutNote{30-3}%
\note{1}{Non un tiret, \italique{paragraphos}, ni un double point ` : ', dont l'usage est courant dans notre manuscrit lorsque le changement d'interlocuteur intervient au milieu d'un vers, mais en l'occurrence une indication de personnage : μ\auDessus{ν}{Η}.}\\\passeLigne}%
\end{groupe}%
%% v. 32
\begin{groupe}%
\personnageGrec{Εὐριπίδης}%
\trimIamb{οὔκ, ἀλλ᾽ ἕτερός τις· οὐχ ἑόρακας πώποτε;\apparat{32}{(Μνη.) οὐχ ... πώποτε \Rac : \lecon{(Εὐρ.) continuat} \Rpc, Σ \espApp ἑόρακα R : \lecon{ἑόρακας} Bentley \ligneApp{v. 33} omm. \Rac, \Mon : in marg. \Rpc}% \ligneApp{33} (Μ) ($=$ (Μνη.) ?) \Rpc}
\\\passeLigne}%
\end{groupe}%
%% v. 33a
\begin{groupe}%
\personnageGrec{Μνησίλοχος}%
\trimIamb{μῶν ὁ δασυπώγων;%\testimonia{33}{(δασυπώγων) Pollux II 88}
\\\passeLigne}%
\end{groupe}%
%% v. 33b
\begin{groupe}%
\personnageGrec{Εὐριπίδης}%
\trimIamb{\ligne{33}\decalage{μῶν ὁ δασυπώγων;}οὐχ ἑόρακας πώποτε.\\\passeLigne}%
\end{groupe}%
%% v. 34
\begin{groupe}%
\personnageGrec{Μνησίλοχος}%
\trimIamb{μὰ τὸν Δί᾽ οὔτοι γ᾽ ὥστε κἀμέ γ᾽ εἰδέναι.\apparat{34}{καὶ μεγ᾽ \Rac \ (ὶ eras. \Rpc) : καὶ μεγ᾽ \Mon : \lecon{κἀμέ γ᾽} edd. omn.}\\\passeLigne}%
\end{groupe}%
%% v. 35-38
\begin{groupe}%
\personnageGrec{Εὐριπίδης}%
\trimIamb{καὶ μὴν †βεβίνηκας† σύ γ᾽, ἀλλ᾽ οὐκ οἶσθ᾽ ἴσως.
\commentaire{35}{καὶ μὴν βεβίνηκας σύ γ᾽, ἀλλ᾽ οὐκ οἶσθ᾽ ἴσως.}{Le ton est inhabituel -- voire plutôt choquant -- dans la bouche d'Euripide.
Qu'à cela ne tienne : puisqu'il semble d'autre part qu'Agathon tourne le dos lorsqu'il accorde ses faveurs\marqueur{1}, nos commentateurs en concluent qu'il est logique que Mnésiloque ne le connaisse pas -- sauf éventuellement au sens biblique... de Sodome, s'entend. 
Et c'est le grave et docte Euripide qui colporterait ce genre de ragots. %nous renseigne de la sorte. 
Y a-t-il besoin de commenter ce genre de commentaires ?.. je précise : ces commentaires \italique{modernes} ; car le βεβίνηκας du \italique{Ravennas} n'a pas laissé derrière lui la moindre scholie ancienne ; chose étonnante lorsqu'on sait que les anciens commentateurs étaient presqu'autant portés au scabreux que nos modernes. %Doublement étonnant même, car la forme βεβίνηκας méritait que les grammairiens s'y arrêtassent aussi : ce serait la seule attestation du verbe βινεῖν au parfait.
Il faut croire que le texte original d'Aristophane, l'\allemand{Urtext} des \italique{Thesmophories}, comportait une tout autre leçon : une banalité sur laquelle aucun commentateur ancien ne se sera arrêté.
Las !.. on ne détruit que ce que l'on remplace ; à défaut d'avoir une proposition alternative solidement étayée, il faut nous résoudre à \latin{crux}-désespérer...
%Ce †  †  
%Chaussons plutôt une bonne paire de lunettes et allons voir le texte d'un peu plus près. Von Velsen \marqueur{\italique{Aristophanis Thesmophoriazusae} (1883). Von Velsen a établi définitivement dans cette édition que le second manuscrit par lequel nous sont connues les \italique{Thesmophories}, le \italique{Monacensis gr. 492} (G dans son apparat critique, \Mon dans le nôtre), n'était qu'une copie du \nieme{XV} siècle du \italique{Ravennas}, apportant quelques émendations bien vues. Son édition compare exhaustivement le \italique{Ravennas} et le \italique{Monacensis} ; ce qui est d'un grand intérêt pour déterminer si une correction du \italique{Ravennas} est ancienne ou plus récente : il suffit d'aller comparer avec la leçon du manuscrit de Münich. Son apparat critique est le plus précis et permet d'avoir la meilleure idée possible du manuscrit de Ravenne -- avec celui de l'édition Prato ; en partie une reprise du premier mais revu et amélioré en de nombreux points. Pour qui voudrait pousser le scrupule plus loin, et puisque le manuscrit de Ravenne n'a toujours pas été numérisé -- ô scandale ! le manuscrit sans qui nous ne connaîtrions que neuf des onze pièces survivantes d'Aristophane !.. --, Van Leeuwen édita en fac-similé le manuscrit de Ravenne : \italique{Comoediae Undecim cum Scholiis : Codex Ravennas 137, 4, A} [$=$ \italique{Rav.} 429 (nouvelle cote)], \italique{phototypice editus} in \italique{Codices Graeci et Latini photographice depicti}, t. IX (1904), pp. XVIII-XIX.} remarqua que le `βεβίνηκας' du \italique{Ravennas} était la correction d'une première leçon : `βεαίνηκας' (l'alpha α ayant été gratté pour le transformer en bêta β dont la forme est très semblable dans notre manuscrit). Le scribe avait-il simplement mal formé son bêta ou bien avait-il omis cette lettre lors de la première rédaction ? 
%Si l'on fait la seconde hypothèse, il faut croire que la leçon recopiée par notre scribe était peu ou prou ΒΕΒΑΙΝΗΚΑC 
%Casanier sédentaire : pourquoi ne se croisent-ils pas ? car aucun des deux n'est un marcheur !
%Casanier -- même il ne quitte pas son lit !
%καὶ μὴ ἀναβαίνω
%καὶ συ γε (nombreuses références)
%
%    ΛΥ.                 Ἀλλ’, ὦ μέλ’, ὄψει τοι σφόδρ’ αὐτὰς Ἀττικάς,
%ἅπαντα δρώσας τοῦ δέοντος ὕστερον.
%Ἀλλ’ οὐδὲ Παράλων οὐδεμία γυνὴ πάρα,
%οὐδ’ ἐκ Σαλαμῖνος.
%ΚΛ.                                    Ἀλλ’ ἐκεῖναί γ’ οἶδ’ ὅτι   (59)
%ἐπὶ τῶν κελήτων διαβεβήκασ’ ὄρθριαι.   (60)
%ΛΥ.                 Οὐδ’ ἃς προσεδόκων κἀλογιζόμην ἐγὼ
%πρώτας παρέσεσθαι δεῦρο τὰς Ἀχαρνέων @1
%γυναῖκας, οὐχ ἥκουσιν.
%ΚΛ.                                       Ἡ γοῦν Θεογένους   (63)
%ὡς δεῦρ’ ἰοῦσα τἀκάτειον ᾔρετο.
%
% W.-Moellendorff, Zum Lexikon des Photios, SPAW 1907,2-14 (= Kleine Schriften, IV,. 528-541).
%Lysistrata 966 καὶ μὴ βινῶν τοὺς ὄρθρους;
%
%βινεῖν· παρὰ Σόλωνι τὸ βίᾳ μίγνυσθαι. τὸ δὲ κατὰ νόμον ὀπύειν + retrouver passage euripide. Cf. : Héraclides
% Χο.        εἰκὸς θεῶν ἱκτῆρας αἰδεῖσθαι, ξένε,
%καὶ μὴ βιαίῳ χερὶ δαιμόνων
%ἀπολείπειν [σ’] ἕδη· @1
%πότνια γὰρ Δίκα τάδ’ οὐ πείσεται.
%Κη.        ἔκπεμπέ νυν γῆς τούσδε τοὺς Εὐρυσθέως,    (105)
%
%
% Ποῖος γὰρ <ἔτ’> ἂν νέφρος ἀντίσχοι,
%ποία ψυχή, ποῖοι δ’ ὄρχεις,
%ποία δ’ ὀσφύς, ποῖος δ’ ὄρρος
%κατατεινόμενος   (965)
%καὶ μὴ βινῶν τοὺς ὄρθρους;
%
%Même problème de concentration qu'avec la belette / le calme ?
%Ecclesiazusae
%Νεις.        αἰαῖ, τί ποτε πείσομαι;
%οὐχ ἥκει μοὐταῖρος
%μόνη δ’ αὐτοῦ λείπομ’· ἡ   (912)
%γάρ μοι μήτηρ ἄλλῃ βέβηκεν.
%
%Ma ton Dia out'(e) oig(w) wste kai (e)me g' eidenai
%-- kai mhn bebhnai kai su g(e) allh\& ()
%kai mh βεβιακαs
%kai mhn dis enika
%Semble confirmer le commentaire du scholiaste au vers 258\marqueur{1} :  Profitons-en pour préciser que c'est bien le fait qu'Agathon se travestisse qui aurait pu l'empêcher d'être connu de Mnésiloque – sauf au sens biblique de Sodome – et non pas, comme l'ont soutenu nombre de commentateurs qui ont une vision bien réductrice de ce genre de rencontres, le simple fait qu'il tourne le dos lorsqu'il accorde ses faveurs.\eff{Le ton est inhabituel dans la bouche d'Euripide.... si bien que l'on se demande s'il ne vaudrait pas mieux lire bainw, piétiner, et le public d'interpréter différemment. Quoi qu'il en soit, ce vers semble confirmer}
}%
%anaphore allitération ἡνίκ’ ἂν σὺ νι-   (1210) κᾷς λέγων τὰς δίκας.
%\debutNote{35}%
%\note{1}{Commentaire paraissant de prime abord tout-à-fait fantaisiste et émanant d'une imagination délurée. Mais il semble corroboré par le }
\debutNote{35}%
\note{1}{Voir le vers 258 et sa scholie ; Agathon se livrerait au travestissement dans l'interlope \italique{Athens by night} pour y être éprouvé à la manière d'une femme -- mais le commentaire déluré du scholiaste n'est peut-être qu'une mauvaise interprétation du présent vers 35 adjointe à une forte dose d'imagination lubrique. Ce qui est sûr, c'est qu'au vers 205, Agathon dit craindre s'il était reconnu d'être lynché pour avoir chapardé nuitamment les œuvres de la Cypris : les Athéniennes voient en lui un concurrent irrégulier.}%
\\
ἀλλ᾽ ἐκποδὼν πτήξωμεν, ὡς ἐξέρχεται\\%\testimonia{36}{(ἀλλ᾽ -- πτήξωμεν) Suidas π 3024}\\
θεράπων τις αὐτοῦ πῦρ ἔχων καὶ μυρρίνας\commentaire{37}{μυρρίνας}{désigne les \italique{rameaux de myrte} que l'on brûlait lors des sacrifices (cf. les \italique{Guêpes}, v. 860 ; les \italique{Oiseaux} au v. 43).}\\
προθυσόμενος ἔοικε τῆς ποιήσεως.\\\passeLigne}%
\end{groupe}%
%% v. 39-41
\begin{groupe}%
\personnageGrec{Θεράπων}%
\typeVers{dim. anap.}\dimAnap{εὔφημος πᾶς ἔστω λαός,%\testimonia{39-41}{Suidas ε 3791 (λεὼς codd. AF : λαὸς codd. GIVM, συγκλείσας, ἐνδημεῖ pro ἐπιδημεῖ )}
\apparat{39}{λέως R, Suid. ε 3792 codd. AF : \lecon{λαός} Σ, Suid. codd. GIVM}\\
στόμα συγκλῄσας· ἐπιδημεῖ γὰρ%\testimonia{40-2}{(ἐπιδημεῖ -- μελοποιῶν) Suidas θ 378}%
\commentaire{39-62}{Passage en dimètres anapestiques.}{(Ceux aux vers 42 et 62 sont catalectiques.) Ce mètre est d'usage courant chez Aristophane pour les prières (cf. les \italique{Guêpes} v. 323-333 et la \italique{Paix} v. 947-1015). Le serviteur d'Agathon est tout empreint de la préciosité du maître, comme, du reste, le serviteur de l'Euripide des \italique{Acharniens} (v. 396-400) l'était.}%
\commentaire{39-40}{εὔφημος πᾶς ἔστω λέως / στόμα συγκλῄσας·}{Ce genre de formule est classique au début d'une prière ou d'un sacrifice. On peut néanmoins constater ici un peu plus d'emphase avec le tour tragique\marqueur{1} et l'utilisation de l'adjectif εὔφημος plutôt que le classique verbe εὐφημεῖν\marqueur{2} ou substantif εὐφημία\marqueur{3}.}%
\debutNote{39-40}%
\note{1}{C'est ce que souligne le scholiaste : παρὰ τοῖς τραγικοῖς οὕτω λέγεται, ``<εὔφημος> πᾶς ἔστω λαός''.}%
\note{2}{\italique{Acharniens} v. 241 ; \italique{Cavaliers} v. 1316 εὐφημεῖν χρὴ καὶ στόμα κλείεν ; \italique{Nuées} v. 263 ; \italique{Paix} v. 434 et 1316 ; \italique{Grenouilles} v. 354.}%
\note{3}{\italique{Guêpes} v. 868 ; \italique{Oiseaux} v. 959.}%
\commentaire{40}{ἐπιδημεῖ.}{\italique{Séjourner}, tel me paraît bien le sens de ἐπιδημεῖ, \latin{peregre advenit et adest}. (Xénophon, \italique{Mémor.}, I, 2, 61.) Cf. ἐπιχωρτάζει, \latin{venit ibique commoratur}. Mελοποιῶν a le sens du futur (voir note [de Willems] sur les \italique{Grenouilles} au v. 148, t. III p. 109-110). (Note de Willems.)}\\
θίασος Μουσῶν ἔνδον μελάθρων}%
\end{groupe}%
%% v. 42-45a (part 1)
\begin{groupe}\typeVers{(catalect.)}\dimAnap{τῶν δεσποσύνων μελοποιῶν.
\commentaire{41-2}{θίασος Μουσῶν ἔνδον μελάθρων / τῶν δεσποσύνων μελοποιῶν.}{Pour bien goûter ces vers, il faut savoir ce qu'est un \italique{thiase} : une bruyante troupe de gens célébrant des rites (bachiques la plupart du temps) et parcourant les rues. Cf. les \italique{Bacchantes} d'Euripide où des femmes munies de tambourins interprètent des danses scandées de cris rituels (évohé) avec une fougue poussée jusqu'à la frénésie. Par extension, ce mot très prisé des tragiques désigne une foule extrêmement bruyante ; Euripide l'emploie à propos d'une bande de centaures (\italique{Iphigénie à Aulis}, v. 1059) ; une troupe endiablée de jeunes danseuses (\italique{Iphigénie en Tauride}, v. 1146) ou encore une meute de guerriers (\italique{Phaéton}, v. 796). 
Pour décrire cette multitude de Muses, au lieu de recourir à la banale image de l'\italique{essaim}\marqueur{1}, à dessein Aristophane préfère évoquer un thiase : 
comme le remarque Sommerstein, l'inversion est totale : un lieu d'un cortège célébrant une divinité, c'est un thiase de Muses qui vient honorer Agathon.}
\debutNote{41-2}
\note{1}{Le Socrate des \italique{Nuées} (vers 297) parle d'un « essaim de déesses », μέγα γάρ τι θεῶν κινεῖται σμῆνος. Voir Jean Taillardat, \italique{Les images d'Aristophane}, §662.}\\\passeLigne}%
\end{groupe}%
\endnumbering
\end{Leftside}
\begin{Rightside}
\beginnumbering
%% v. 1-4
\begin{groupe}%
\personnageFr{Mnésiloque}%
\trimIambFr{Ô Zeus, est-ce qu'à la fin une hirondelle paraîtra ?\testimonia{1}{\Fr 617 πυθοῦ χελιδὼν πηνίκ᾽ ἄττα φαίνεται}\\
Cet homme me tuera de me faire tribuler dès le matin.\testimonia{2}{(ἀλοῶν) Suidas α 1330 ; Photius α 1029 $=$ \SynB α 984}\\
Y a-t-il moyen, avant que perdre tout de bon la rate,\\
de savoir de toi où tu me mènes, Euripide ?\\\passeLigne}%
\end{groupe}%
%% v. 5-6a
\begin{groupe}%
\personnageFr{Euripide}%
\trimIambFr{Eh ! il n'est pas besoin que tu entendes tout ce que\\
tout à l'heure tu verras de tes yeux.\\\passeLigne}%
\end{groupe}%
%% v. 6b-7a
\begin{groupe}%
\personnageFr{Mnésiloque}%
\trimIambFr{\ligne{6}\espace{2.5cm}Comment dis-tu ? Répète.\\
Il n'est pas besoin que j'entende.\\\passeLigne}%
\end{groupe}%
%% v. 7b
\begin{groupe}%
\personnageFr{Euripide}%
\trimIambFr{\espace{2.7cm}\ligne{7}Non, ce que tu vas voir.\\\passeLigne}%
\end{groupe}%
%% v. 8a
\begin{groupe}%
\personnageFr{Mnésiloque}%
\trimIambFr{Ni donc que je voie...\\\passeLigne}%
\end{groupe}%
%% v. 8b
\begin{groupe}%
\personnageFr{Euripide}%
\trimIambFr{\espace{3.1cm}\ligne{8}Non, ce qu'il faut que tu entendes.\\\passeLigne}%
\end{groupe}%
%% v. 9-10
\begin{groupe}%
\personnageFr{Mnésiloque}%
\trimIambFr{Singulier conseil ! Ma foi, c'est finement dit.\\
Tu prétends que je ne dois ni entendre ni voir.\\\passeLigne}%
\end{groupe}%
%% v. 11
\begin{groupe}%
\personnageFr{Euripide}%
\trimIambFr{C'est que distincte est l'essence de chacune de ces \fin choses.\\\passeLigne}%
\end{groupe}%
%% v. 12a
\begin{groupe}%
\personnageFr{Mnésiloque}%
\trimIambFr{Ne pas entendre et ne pas voir ?\\\passeLigne}%
\end{groupe}%
%% v. 12b
\begin{groupe}%
\personnageFr{Euripide}%
\trimIambFr{\ligne{12}\espace{4.4cm}Tiens-le pour assuré.\\\passeLigne}%
\end{groupe}%
%% v. 13a
\begin{groupe}%
\personnageFr{Mnésiloque}%
\trimIambFr{Comment distincte ?\\\passeLigne}%
\end{groupe}%
%% v. 13b-18
\begin{groupe}%
\personnageFr{Euripide}%
\trimIambFr{\ligne{13}\espace{1.8cm}Voici comment elles furent départies à l'origine.\\
L'Éther, quand d'abord il se dégagea
et procréa \\ en soi-même des animaux doués de mouvement,\\
commença par fabriquer pour qui doit voir\\
un œil analogique de l'orbe solaire\\
et fora une oreille à l'image des godets d'un moule.\testimonia{18-9}{(χοάνη = χώνη) Moeris χ 1 ~ Hesychius χ 576 χοάνη}\\\passeLigne}%
\end{groupe}%
%% v. 19-21
\begin{groupe}%
\personnageFr{Mnésiloque}%
\trimIambFr{Et à cause de ce moule je ne dois ni entendre ni voir ?\\
Par Zeus, je suis ravi d'apprendre encore cela.
Quelle chose tout de même, que les doctes accointances !\\\passeLigne}%
\end{groupe}%
%% v. 22a
\begin{groupe}%
\personnageFr{Euripide}%
\trimIambFr{Auprès de moi, tu en apprendras bien d'autres du \fin même genre.\\\passeLigne}%
\end{groupe}%
%% v. 22b-24
\begin{groupe}%
\personnageFr{Mnésiloque}%
\trimIambFr{\ligne{22}\espace{5.3cm}Ah ! que ne puis-je\\
en outre apprendre aussi à être perclus des jambes ?\testimonia{24}{\Macaire VI 90 ; \Apostol II 94}\\\passeLigne}%
\end{groupe}%
%% v. 25a
\begin{groupe}%
\personnageFr{Euripide}%
\trimIambFr{Avance ici et prête-moi ton attention.\\\passeLigne}%
\end{groupe}%
%% v. 25b
\begin{groupe}%
\personnageFr{Mnésiloque}%
\trimIambFr{\ligne{25}\espace{5.2cm}J'y suis.\\\passeLigne}%
\end{groupe}%
%% v. 26a
\begin{groupe}%
\personnageFr{Euripide}%
\trimIambFr{Tu vois cette petite porte ?\\\passeLigne}%
\end{groupe}%
%% v. 26b-27a
\begin{groupe}%
\personnageFr{Mnésiloque}%
\trimIambFr{\ligne{26}\espace{3.4cm}Oui, par Héraclès,\\
à ce que je crois, du moins.\\\passeLigne}%
\end{groupe}%
%% v. 27b
\begin{groupe}%
\personnageFr{Euripide}%
\trimIambFr{\ligne{27}\espace{1.4cm}Maintenant, motus.\\\passeLigne}%
\end{groupe}%
%% v. 27c
\begin{groupe}%
\personnageFr{Mnésiloque}%
\trimIambFr{\ligne{27}\espace{3.4cm}Je ne dis mot de la petite porte.\\\passeLigne}%
\end{groupe}%
%% v. 28a
\begin{groupe}%
\personnageFr{Euripide}%
\trimIambFr{Écoute !\\\passeLigne}%
\end{groupe}%
%% v. 28b
\begin{groupe}%
\personnageFr{Mnésiloque}%
\trimIambFr{\ligne{28}\espace{1.4cm}J'écoute et ne dis mot de la petite porte.\\\passeLigne}%
\end{groupe}%
%% v. 29-30a
\begin{groupe}%
\personnageFr{Euripide}%
\trimIambFr{C'est là que se trouve demeurer l'illustre Agathon,\\
le poète tragique.\\\passeLigne}%
\end{groupe}%
%% v. 30b
\begin{groupe}%
\personnageFr{Mnésiloque}%
\trimIambFr{\ligne{30}\espace{2.4cm}Qui est-ce, cet Agathon ?
Est-ce l'un de ces honnêtes hommes... celui au teint hâlé, plein d'allant ?\\\passeLigne}%
\end{groupe}%
%% v. 32
\begin{groupe}%
\personnageFr{Euripide}%
\trimIambFr{Non, c'est un autre. Tu ne l'as jamais vu ?\\\passeLigne}%
\end{groupe}%
%% v. 33a
\begin{groupe}%
\personnageFr{Mnésiloque}%
\trimIambFr{L'homme à la barbe touffue ?\testimonia{33}{(δασυπώγων) Pollux II 88}\\\passeLigne}%
\end{groupe}%
%% v. 33b
\begin{groupe}%
\personnageFr{Euripide}%
\trimIambFr{\ligne{33}\espace{3cm}Tu ne l'as jamais vu !..\\\passeLigne}%
\end{groupe}%
%% v. 34
\begin{groupe}%
\personnageFr{Mnésiloque}%
\trimIambFr{Non, par Zeus, pas du moins que je sache.\\\passeLigne}%
\end{groupe}%
%% v. 35-38
\begin{groupe}%
\personnageFr{Euripide}%
\trimIambFr{À coup sûr tu l'as enfourché, mais peut-être ne t'en \fin souvient-il pas.\\
Mais allons nous tapir à l'écart. Car sort de chez lui\testimonia{36}{(ἀλλ᾽ -- πτήξωμεν) Suidas π 3024}\\
un serviteur avec du feu et des myrtes.\\
Il va sacrifier, ce semble, en faveur du succès de sa \fin poésie.\\\passeLigne}%
\end{groupe}%
%% v. 39-45a
\begin{groupe}%
\personnageFr{Serviteur d'Agathon}%
\dimAnapFr{Que le peuple entier se recueille\testimonia{39-41}{Suidas ε 3791 (λεὼς codd. AF : λαὸς codd. GIVM, συγκλείσας, ἐνδημεῖ pro ἐπιδημεῖ )}\\
et ferme la bouche, car est venu séjourner\\
un thiase de muses dans la demeure}%
\end{groupe}%
%% v. 42-45a (part 1)
\begin{groupe}\dimAnapFr{de mon maître pour composer des chants.\testimonia{40-2}{(ἐπιδημεῖ -- μελοποιῶν) Suidas θ 378}\\\passeLigne}%
\end{groupe}%
\endnumbering
\end{Rightside}
\end{pages}
\Pages
\clearpage\ifodd\value{page}\hbox{}\newpage\fi
\enTeteG{ΠΡΟΛΟΓΟΣ --- Codex Ravennas 429, f. 163r (Classense 165r) --- vv. 43–79}
\begin{adjustwidth}{-1.5cm}{-3.0cm}
\begin{adjustbox}{max totalheight=21.0cm, max width=16.5cm, keepaspectratio}
\begin{minipage}{16.5cm}
{\fontsize{7.2pt}{8.5pt}\selectfont
\noindent\textbf{Folio 02 --- Codex Ravennas 429, \texttt{f. 163r} (\texttt{Classense 165r}) : ΠΡΟΛΟΓΟΣ, vv. 43–79}\hfill \textit{Témoins :} $\text{R}^1, \Rpc, \Rdeux$ (\Mon)\par\vspace{2pt}
\hrule height 0.4pt\vspace{2.5pt}
}
{\fontsize{8.1pt}{9.5pt}\selectfont
\ligneRav{43}{θερ/}{\hspace{1.1em}\auDessus{ἐχέτω}{\tiny ἐπεχέτω} δὲ πνοὰς νήνεμος αἰθήρ,}{\textbf{43} \lecon{ἐχέτω} R (avec la glose interne ἐπεχέτω)}
\ligneRav{44}{}{\hspace{1.1em}κῦμα δὲ πόντου μὴ κελαδείτω}{\textbf{44} \lecon{δὲ} R : τε Austin, cf. v.46, 47}
\ligneRav{45}{—}{\hspace{1.1em}γλαυκόν — : βομβάξ· : σίγα· τί λέγει\inRas{ς}{};}{\textbf{45} λέγει\inRas{ς}{} \Rpc}
\ligneRav{}{θερ/}{\hspace{1.1em}πτηνῶν τε γένη κατακοιμάσθω,}{}
\ligneRav{}{}{\hspace{1.1em}θηρῶν τ᾽ ἀγρίων πόδες ὑλοδρόμων}{}
\ligneRav{}{—}{\hspace{1.1em}μὴ λυέσθων. : βομβαλοβομβάξ·}{}
\ligneRav{}{θερ/}{\hspace{1.1em}μέλλει γὰρ ὁ καλλιεπὴς Ἀγάθων}{}
\ligneRav{50}{—}{\hspace{1.1em}\difference{πρανοῦν προτέρους} — : μῶν βινεῖσθαι;}{\textbf{50} πράμος R : πρόμος Grynaeus : \lecon{πρανοῦν} Bibi ap. Hesych. π 3188 \espApp ἡμέτερος R : \lecon{προτέρους} Bibi, cf. \italique{Nubes} v.562}
\ligneRav{}{θερ/}{\hspace{1.1em}τίς ὁ φωνήσας; : νήνεμος αἰθήρ·}{}
\ligneRav{}{θερ/}{\hspace{1.1em}δρυόχους, τιθέναι δράματος ἀρχάς·}{}
\ligneRav{53}{}{\hspace{1.1em}κάμπτει δὲ νέας \difference{ἁψῖδας} ἐπῶν,}{\textbf{53} ἀσπίδας R : \lecon{ἁψῖδας} Suid. δ 1547, λ 181}
\ligneRav{}{}{\hspace{1.1em}τὰ δὲ τορνεύει, τὰ δὲ κολλομελεῖ,}{}
\ligneRav{55}{}{\hspace{1.1em}καὶ γνωμοτυπεῖ κἀντονομάζει}{}
\ligneRav{56}{}{\hspace{1.1em}καὶ κηροχυτεῖ καὶ \difference{γογγύλλει}}{\textbf{56} γογγύλίζει R, Suid. γ 361 : \lecon{γογγύλλει} Porson : γογγυλίει Hermann (ad \italique{Nub.} 992)}
\ligneRav{}{—}{\hspace{1.1em}καὶ χοανεύει— : καὶ λαικάζει·}{}
\ligneRav{58}{θερ/}{\hspace{1.1em}τίς \difference{ἀγροιώτας} πελάθει θριγκοῖς;}{\textbf{58} ἀγριώτας R : \lecon{ἀγροιώτας} anonym. Parisinus}
\ligneRav{}{μ\auDessus{ν}{Η}}{\hspace{1.1em}ὃς ἕτοιμος σοῦ τοῦ τε ποιητοῦ}{}
\ligneRav{60}{}{\hspace{1.1em}τοῦ καλλιεποῦς \expo{κατὰ} τοῦ θριγκοῦ}{\textbf{60} κατὰ om. \Rac : add. \Rpc}
\ligneRav{61}{}{\hspace{1.1em}\difference{συγγογγύλας} καὶ συστρέψας}{\textbf{61} γογγυλίσας R : \lecon{συγγογγύλας} Enger}
\ligneRav{}{}{\hspace{1.1em}τουτὶ τὸ πέος χοανεῦσαι·}{}
\ligneRav{}{θερ/}{ἦ που νέος γ᾽ ὢν ἦσθ᾽ ὑβριστὴς ὦ γέρον·}{}
\ligneRav{}{\auDessus{εὐ}{ρ}}{ὦ δαιμόνιε τοῦτον μὲν ἔα χαίρειν, σὺ δὲ}{}
\ligneRav{65}{}{Ἀγάθωνά μοι δεῦρ᾽ ἐκκάλεσον πάσῃ τέχνῃ·}{}
\ligneRav{}{θερ/}{μηδὲν ἱκέτευ᾽· αὐτὸς γὰρ ἔξεισιν τάχα·}{}
\ligneRav{}{}{καὶ γὰρ μελοποιεῖν ἄρχεται· χειμῶνος οὖν}{}
\ligneRav{}{}{ὄντος κατακάμπτειν τὰς στροφὰς οὐ ῥᾴδιον,}{}
\ligneRav{69}{}{ἢν μὴ προίῃ θύρασι πρὸς τὸν ἥλιον·}{\textbf{69} \lecon{θύρασι} R : θύραζε Zanetti}
\ligneRav{70}{μ\auDessus{ν}{Η}}{τί οὖν ἐγὼ δρῶ; : περίμεν᾽, ὡς ἐξερχεται·}{}
\ligneRav{71}{μ\auDessus{ν}{Η}}{ὦ Ζεῦ τί δρᾶσαι διανοεῖ με τήμερον;}{\textbf{71-4} (Μνη.) (71) et (Εὐρ.) (72-4) R : \lecon{(Μνη.) (71-4)} Bibi ap. Σ : alii alia}
\ligneRav{}{}{νὴ τοὺς θεοὺς ἐγὼ πυθέσθαι βούλομαι}{}
\ligneRav{}{}{τί τὸ πρᾶγμα τουτί. τί στένεις; τί δυσφορεῖς;}{}
\ligneRav{74}{}{οὐ χρῆν σε κρύπτειν ὄντα κηδεστὴν ἐμόν·}{\textbf{74} \lecon{χρῆν} R, Suid. κ 1487 codd. FV : χρή Suid. codd. AGM \espApp \lecon{ἐμόν} R : ἐμέ Valckenaer ap. Dobree (p. 237)}
\ligneRav{75}{\auDessus{εὐ}{ρ}}{ἔστιν κακόν μοι μέγα τι προπεφυραμένον·}{}
\ligneRav{}{μ\auDessus{ν}{Η}}{ποῖόν τι; : τῇδε θἠμέρᾳ κριθήσεται}{}
\ligneRav{}{}{εἴτ᾽ ἔστ᾽ ἔτι ζῶν εἴτ᾽ ἀπόλωλ᾽ Εὐριπίδης·}{}
\ligneRav{}{μ\auDessus{ν}{Η}}{καὶ πῶς; ἐπεὶ νῦν γ᾽ οὔτε τὰ δικαστήρια}{}
\ligneRav{}{}{μέλλει δικάζειν οὔτε βουλῆς ἐσθ᾽ ἕδρα,}{}
}
\vspace{3pt}\noindent\rule{4.5cm}{0.4pt}\par\vspace{1.5pt}
{\fontsize{7.0pt}{8.3pt}\selectfont
\noindent\textbf{Loci paralleli (vv. 43–79) ---} \textbf{43-5} (ἐχέτω -- γλαυκόν) Suidas ε 3792 ; Zonaras 941 ἐχέτω \espTest \textbf{44} (κελαδείτω) Suidas κ 1278 \espTest \textbf{45} (βομβάξ) Hesychius β 787 βομβάξ ; Suidas ε 2807 ; \Eustathe 1718.9 \espTest \textbf{48} (βομβαλοβομβάξ) Suidas ε 2807 ; \Eustathe 790.13, 1718.13 \espTest \textbf{49} (καλλιεπὴς) Philostrate \italique{Vit. soph.} I 9.1 [$=$ II 12.8 Kayser] \espTest \textbf{50} (πραῖνοι) Hesychius π 3188 †πραῖνοι· πρηνίζειν, καταστρέφειν \espTest \textbf{50} (πραῖνοι) Hesychius π 3188 πραῖνοι· πρηνίζει, καταστρέφει \espTest \textbf{52-7} Suidas δ 1547 (sed τορεύει pro τορνεύει, χιανεύει pro χοανεύει) \espTest \textbf{52} Zonaras 569 δρύοχοι \espTest \textbf{53} (κάμπτει) Hesychius κ 617 κάμπτειν \espTest \textbf{54} (τορνεύει) Hesychius τ 1168 τορνεύει \espTest \textbf{56} (γογγύλλει) Hesychius γ 764 γογγύλλειν \espTest Suidas γ 361 $=$ Zonaras 449 γογγυλίζει \espTest \textbf{57} (χοανεύει) \Eustathe 1153.40 \espTest \textbf{58} (θριγκοῖς) Hesychius θ 750 θριγκοῖς· φραγμοῖς \espTest \textbf{62} (χοανεῦσαι $=$ χωνεῦσαι) Moeris χ 21 ; Hesychius χ 575 χοανεῦσαι \espTest \textbf{74} Suidas κ 1487, cf. Moeris κ 38\par
}
\vspace{2.5pt}\noindent\rule{4.5cm}{0.4pt}\par\vspace{1.5pt}
{\fontsize{6.9pt}{8.2pt}\selectfont
\noindent\textbf{Suite (en commentaire) de l'apparat critique \& notes paléographiques (\texttt{f. 163r}) ---} \textbf{v. 43 :} Glose interlinéaire/interne dans R : ἐπεχέτω glosant ἐχέτω. --- \textbf{v. 45 :} Grattage du sigma final dans λέγει\inRas{ς}{} par \Rpc{}. --- \textbf{v. 50 :} Lieu célèbre de corruption paléographique (πράμος ἡμέτερος R), magistralement résolu par Louis Hauseux (\textbf{Bibi}). --- \textbf{v. 60 :} La préposition κατὰ, omise par \Rac{}, a été rajoutée par \Rpc{} au-dessus de la ligne. --- \textbf{v. 43 :} \lecon{ἐχέτω} R (avec la glose interne ἐπεχέτω). --- \textbf{v. 44 :} \lecon{δὲ} R (maintenu par \textbf{Bibi}) : τε Austin (cf. vv. 46, 47). --- \textbf{v. 50 :} πράμος R (πρόμος Grynaeus) : \lecon{πρανοῦν} \lecon{Bibi} ap. Hésychius $\pi\ 3188$ (et πραῖνοι dans l'\italique{Index}) \espApp ἡμέτερος R : \lecon{προτέρους} \lecon{Bibi}, cf. \italique{Nuées} v. 562. \textbf{Explication paléographique (Commentaire ad v. 50) :} Le texte primitif πρανοῦν προτέρους (ou πρανοι ν) a subi : 1° une médivision avec confusion minuscule $\nu \to \mu$ et $\iota \to \varsigma$ donnant πράμος ; 2° la syllabe προ- de προτέρους, écrite au moyen du signe tachygraphique en forme de Ʒ inversé (semblable à un $\varepsilon$, cf. Gardthausen II p. 340), a été rattachée au $\nu$ précédent et lue <η>μ- + -ετερος $\to$ ἡμέτερος ! --- \textbf{v. 53 :} ἀσπίδας R (« boucliers », faute par attraction d'un mot militaire banal) : \lecon{ἁψῖδας} Souda $\delta\ 1547, \lambda\ 181$ (« cintres, voûtes d'hymnes »). --- \textbf{v. 60 :} κατὰ om. \Rac{} : \textbf{add.} \Rpc{}. --- \textbf{v. 69 :} \lecon{θύρασι} R (maintenu par \textbf{Bibi}) : θύραζε Zanetti. --- \textbf{vv. 71–74 :} (Μνη.) (71) et (Εὐρ.) (72–74) R : \lecon{(Μνη.) (71–74)} \lecon{Bibi} d'après $\Sigma^{\text{R}}$ (toute la tirade appartient à Mnésiloque).\par
}
\end{minipage}
\end{adjustbox}
\end{adjustwidth}
\newpage
\enTeteD{Codex Ravennas 429 (Biblioteca Classense) --- Fac-similé du f. 163r (vv. 43–79)}
\begin{adjustwidth}{-1.5cm}{-3.0cm}
\begin{center}
\includegraphics[width=16.5cm,height=20.5cm,keepaspectratio]{/workspaces/Thesmophories/dossier_ravennas_thesmophories/folios_individuels_ff162v-177r/02_f163r_classense165r_vv0043-0079.jpg}\par
\vspace{2pt}
{\scriptsize\textbf{Codex Ravennas 429, f. 163r (Classense 165r)} --- ΠΡΟΛΟΓΟΣ, \textbf{vv. 43–79} ($\text{R}^1, \Rpc, \Rdeux$ (\Mon)).}
\end{center}
\end{adjustwidth}
\newpage
\clearpage\ifodd\value{page}\hbox{}\newpage\fi
\enTeteG{ΘΕΣΜΟΦΟΡΙΑΖΟΥΣΑΙ --- Texte grec établi (f. 163r, vv. 43–79)}
\enTeteD{LES THESMOPHORIES --- Traduction française (f. 163r, vv. 43–79)}
\begin{pages}
\begin{Leftside}
\beginnumbering
%% v. 43+ (part 2)
\begin{groupe}
\dimAnap{\ligne{43}ἐχέτω δὲ πνοὰς νήνεμος αἰθήρ,%\testimonia{43-5}{(ἐχέτω -- γλαυκόν) Suidas ε 3792 ; Zonaras 941 ἐχέτω}%
\commentaire{43-7}{ἐχέτω δὲ πνοὰς νήνεμος αἰθήρ ...}{De même que dans les \italique{Bacchantes} d'Euripide (v. 1084), à la voix de Dionysos « l'éther se tut, le ravin herbeux étouffa le murmure des feuilles, et l'on n'eût point entendu le cri des bêtes sauvages ».%
\citComm{σίγησε δ᾽ αἰθήρ, σῖγα δ᾽ ὕλιμος νάπη\\
φύλλ᾽ εἶχε, θηρῶν δ᾽ οὐκ ἂν ἤκουσας βοήν.}%
Et aussi dans Théocrite (II, 38), tout se tait à l'arrivée d'Hécate. (Note de Willems.)}%
\commentaire{43}{ἐχέτω}{ est mis pour ἐπεχέτω (ici : \italique{retenir}) comme l'indique la scholie intérieure au texte du \italique{Ravennas}.}\\
κῦμα δὲ πόντου μὴ κελαδείτω\apparat{44}{\lecon{δὲ} R : τε Austin, cf. v.46, 47}\\%\testimonia{44}{(κελαδείτω) Suidas κ 1278}\apparat{44}{\lecon{δὲ} R : τε Austin, cf. v.46, 47}\\
γλαυκόν—\\\passeLigne}%
\end{groupe}%
%% v. 45b
\begin{groupe}%
\personnageGrec{Μνησίλοχος}%
\dimAnap{\ligne{45}\decalage{γλαυκόν—}βομβάξ.%\testimonia{45}{(βομβάξ) Hesychius β 787 βομβάξ ; Suidas ε 2807 ; \Eustathe 1718.9}%
\commentaire{45}{βομβάξ.}{%Forgée sur βόμβος, \italique{bruit sourd}, \italique{grondement}, \italique{bourdonnement}, βομβάξ est une interjection créée pour railler un style pompeux « ran plan plan » et avec redoublement (βομβαλοβομβάξ, v. 48). (Note de Chantraine.)\\
%Forgée sur βόμβος, \italique{bourdonnement}, βομβάξ est une interjection créée pour ridiculiser la ruche à muses dont
%L'interjection, qui moque les \italique{bourdons} 
%Il pourrait être plaisant de la traduire en citant la célèbre chanson de Dalida « Paroles \& paroles \& paroles » et au redoublement : « Et encore des paroles que tu sèmes dans l'Éther ! » 
Forgée sur le modèle d'ὄμπαξ -- interjection pour réclamer le silence « paix ! (πάξ) », « assez ! » --, βομβάξ fait entendre le \italique{bourdonnement}, βόμβος (ὁ βομβυλιός désigne tout type d'insecte bourdonnant). Plaisante manière pour Aristophane de ridiculiser l'image détournée du serviteur : le bruyant thiase et sa pompe de muses venues honorer la demeure d'Agathon sont ramenées à la plus commune métaphore de la \italique{ruche}.
}\\\passeLigne}%
\end{groupe}%
%% v. 45c
\begin{groupe}%
\personnageGrec{Εὐριπίδης}%
\dimAnap{\ligne{45}\decalage{γλαυκόν—βομβάξ.}σίγα. τί λέγει;\apparat{45}{λέγει\inRas{ς}{} \Rpc }\\\passeLigne}%
\end{groupe}%
%% v. 46-48a
\begin{groupe}%
\personnageGrec{Θεράπων}%
\dimAnap{πτηνῶν τε γένη κατακοιμάσθω,\\
θηρῶν τ᾽ ἀγρίων πόδες ὑλοδρόμων\\
μὴ λυέσθων.\\\passeLigne}%
\end{groupe}%
%% v. 48b
\begin{groupe}%
%%
\personnageGrec{Μνησίλοχος}%
\dimAnap{\ligne{48}\decalage{μὴ λυέσθων.}βομβαλοβομβάξ.%\testimonia{48}{(βομβαλοβομβάξ) Suidas ε 2807 ; \Eustathe 790.13, 1718.13}
\\\passeLigne}%
\end{groupe}%
%% v. 49-50a
\begin{groupe}%
%%
\personnageGrec{Θεράπων}%
\dimAnap{μέλλει γὰρ ὁ καλλιεπὴς Ἀγάθων%\testimonia{49}{(καλλιεπὴς) Philostrate \italique{Vit. soph.} I 9.1 [$=$ II 12.8 Kayser]}
\\
πρανοῦν προτέρους—%\testimonia{50}{(πραῖνοι) Hesychius π 3188 †πραῖνοι· πρηνίζειν, καταστρέφειν}
\apparat{50}{πράμος R : πρόμος Grynaeus : \lecon{πρανοῦν} Bibi ap. Hesych. π 3188 \espApp ἡμέτερος R : \lecon{προτέρους} Bibi, cf. \italique{Nubes} v.562}%
\commentaire{49-50}{μέλλει γὰρ ὁ καλλιεπὴς Ἀγάθων / πράμος ἡμέτερος.}{%Deux hapax dans cette courte réplique du serviteur d'Agathon : καλλιεπὴς et πράμος ; préciosité de style qui permet un détournement obscène de la part de Mnésiloque (vers 50).\\
%Si le sens de καλλιεπὴς se laisse saisir sans problème, ce n'est pas le cas de πράμος, ni autrement attesté ni même relevé par les grammairiens. Pour rentrer dans le mètre anapestique, la scansion est nécessairement $\cup \cup $, ce qui empêche tout rapprochement étymologique pertinent. 
%L'on reste dubitatif devant le †πράμος† du vers 50. Ce mot, 
πράμος n'est ni autrement attesté %-- si tant est que le texte ait été correctement transmis -- 
ni même relevé par les grammairiens et doit nécessairement être scandé $\cup \cup $ pour rentrer dans le mètre anapestique, empêchant tout rapprochement étymologique pertinent. 
Autant d'indices qui font soupçonner une erreur textuelle plutôt qu'un hapax agathonien.
πρόμος, \italique{le chef}, fut proposé par Grynaeus et quasi-unanimement repris depuis. 
%Pourquoi cette erreur aurait-elle été commise ? ne chicanons pas paléographiquement : même s'il est rare qu'un O se transforme en A, l'erreur serait tout-à-fait possible\marqueur{1}. Cette émendation contredit surtout le principe de la \italique{lectio difficilior} : un mot certes rare, πρόμος, mais bien attesté aurait été transformé en un hapax... 
La correction est très légère mais n'est pas satisfaisante : 
%Plus important encore : 
il manque au terme πρόμος une acception \italique{sensu obsceno} pour expliquer l'exclamation de Mnésiloque : μῶν βινεῖσθαι;\marqueur{1} 
De plus, %d'un point de vue grammatical, 
l'on s'attendrait plutôt à trouver un infinitif réclamé par μέλλει. % plutôt qu'un nom.
%Heureusement pour nous, 
Une glose d'Hésychius va nous permettre de débrouiller l'écheveau. %n'a pas manqué de relever le vocable d'Agathon, rare et à double sens :%
Cette glose est elle-même légèrement corrompue (π 3188) :%
\citComm{†πραῖνοι†· πρηνίζειν, καταστρέφειν.}%
Sans nous attarder pour l'instant sur les hypothèses de corruptions, constatons que Hésychius atteste l'existence du verbe πρᾳνόω (ou πρανόω\marqueur{2}).
Pour le sens, πρηνίζειν signifie \italique{précipiter}, \italique{renverser}, \italique{jeter} (en avant).
Le serviteur d'Agathon construit ainsi sa phrase :%
\citComm{μέλλει γὰρ Ἁγάθων πρανοῦν δρυόχους\marqueur{3}.\\
Agathon s'apprête à jeter les fondations\marqueur{3}.}%
Seulement, πρηνίζειν peut désigner aussi l'action de s'accroupir tête baissée\marqueur{4}. (πρανοῦν, comme accroupir, fait entendre la \italique{croupe}, πρανής.) C'est bien-sûr dans ce sens détourné que le prend Mnésiloque -- image d'un Agathon se jetant à son travail dans la position du penseur de Rodin ? L'occasion est trop belle : Mnésiloque s'imagine Agathon en position κύβδα (cf. le vers 489) et tournant le dos à d'éventuels assaillants --- μῶν βινεῖσθαι; demande-t-il. Le véritable Agathon devait effectivement travailler dans une posture qui prêtait à rire aux mauvais esprits : c'est sur son lit et dans cette même posture qu'il apparaît sur scène quelques vers plus loin ; cette position étrange lui vaudra d'être comparé par Mnésiloque à Cyrène \italique{aux douze postures} (δωδεκαμήχανος) (v. 98).\\
Reste à expliquer la présence dans le \italique{Ravennas} de †ἡμέτερος† après †πράμος†. %C'est là que la corruption de la glose d'Hésychius se révèle cruciale. Tout au long de notre commentaire des \italique{Thesmophories}, nous allons montrer combien le lexique d'Hésychius est important pour l'exégèse d'Aristophane : nombre d'entrées sont tirées des meilleurs commentaires d'Aristophane. C'en est ici, croyons-nous, un exemple. 
Hésychius glose †πραῖνοι† par deux infinitifs, πρηνίζειν et καταστρέφειν, preuve que le mot relevé en vedette devait originellement être lui aussi un infinitif. Le rétablissement de πρανοῦν se fait alors sans difficulté. % (la chute d'un ν est courante ainsi que l'iotacisme des υ en ι). 
C'est la même mécoupure du verbe πρανοῦν, lu †πρανοι \ ν† qui pourrait expliquer le †πράμος† et le début de †ἡμέτερος† du \italique{Ravennas} selon le \italique{stemma} suivant :  %†πραῖνοι† et †πράμος† pourraient donc tout-à-fait provenir d'un même archétype :
%$$\overset{\overset{\text{\allemand{Urtext}}}{\text{πρηνίζειν, καταστρέφειν}}}{\text{πρανοῦν (???)}}\xrightarrow{\text{mécoupure}}\text{πρανου ν(???)}\xrightarrow{\text{confusion υ/ι}}$$
%$$\overset{\text{\allemand{Urtext} : \ \ \ }\text{πρηνίζειν, καταστρέφειν}}{\text{πρανοῦν (???????)}}\xrightarrow{\text{mécoupure}}\overset{\text{πρηνίζειν, καταστρέφειν}}{\text{πρανου \ \ ν...}}\xrightarrow{\text{confusion υ/ι}}\overset{\text{Hésychius π 3188}}{\text{πραῖνοι· πρηνίζειν, καταστρέφειν}}$$
\begin{align*}
\overset{\overset{\text{\footnotesize{\allemand{Urtext}}}}{\text{πρηνίζειν, καταστρέφειν}}}{\text{πρανοῦν [...]}}\xrightarrow{\text{mécoup. \ } + \text{\ conf. \  υ/ι}}\overset{\overset{\text{\footnotesize{Archétype}}}{\text{πρηνίζειν, καταστρέφειν}}}{\text{πρανοι \ \ ν[...]}}&\xrightarrow{\text{\ \ \ \ \ }}\overset{\text{\footnotesize{Hésychius π 3188}}}{\text{πραῖνοι· πρηνίζειν, καταστρέφειν}}\\
& \xrightarrow{\text{confus. \ ν/μ et ι/ς}}\overset{\text{\footnotesize{\italique{Ravennas}}}}{\text{πραμος <η>μ[ετερος]}}
\end{align*}
%Cette mauvaise séparation des mots couplée à la confusion entre les μ et les ν pourrait-elle aussi expliquer l'étrange présence du †ἡμέτερος† ? reconstruit à partir du dernier ν lu en μ de πρανοῦν ; mais comment expliquer alors la fin : †ετερος†. Il faut savoir que d'anciens manuscrits comportaient de nombreuses abréviations\marqueur{Cf. par exemple la très ingénieuse reconstruction du vers 80 par Nauck.} ; l'une d'elle consistait à abréger προ- par une sorte de Ʒ renversé qui ressemblait à un epsilon : \renverse{Ʒ}\marqueur{Voir  p. 340.}.
Rien de plus semblable qu'un ι et un ς, ainsi qu'un μ et un ν  -- dans les manuscrits écrits en minuscules à l'instar du \italique{Ravennas}.
Il s'agit maintenant de comprendre la présence de †ετερος†. Sans tenter d'explications, je constate qu'au vers 562 des \italique{Nuées}, les deux principaux manuscrits des pièces d'Aristophane, notre \italique{Ravennas} et le \italique{Venetus}, portent deux leçons différentes : ἑτέρας \italique{versus} προτέρας\marqueur{5}. 
Je suppose ici la même corruption :%. Voici le texte d'Agathon restitué (sans l'interruption de Mnésiloque) :%
\citComm{μέλλει γὰρ ὁ καλλιεπὴς Ἀγάθων πρανοῦν προτέρους δρυόχους, τιθέναι δράματος ἀρχάς. κάμπτει δὲ νέας ἁψῖδας ἐπῶν κτλ.\\
Agathon au beau langage s'apprête à jeter les premiers fondements, à disposer les prémices d'un drame. Il tourne de nouvelles roues de mots etc.}%
%Nous disions qu'il manquait au terme πρόμος une acception obscène. % pour être une correction acceptable. Il faut croire que Sommerstein a pris un peu trop au pied de la lettre la plaisanterie dans 
%Le rapprochement πρόμος -- πόρνος en arguant de la plaisanteries de \italique{Nuées} (v. 394) sur l'indissociabilité -- certes pas au niveau de la prononciation comme nous l'allons voir -- entre βροντή et πορδή. En passant, profitons-en aussi pour exégéter ce jeu de mots mécompris. %Il y avait en grec un proverbe, βὴξ ἀντὶ πορδῆς, «~une toux à la place d'un pet~», qui servait à désigner quelqu'un tentant vainement de camoufler une chose honteuse. (Suidas (β 261), les parémiographes ainsi qu'Érasme, au n°563 de ses \italique{Adages}, l'ont relevé.) Il devait certainement aussi exister l'adage cousin βροντὴ ἀντὶ πορδῆς ; ainsi s'explique notre passage tout comme une anecdote amusante rapportée par Strabon (14, 5, 14) et par Macaire (4, 11) : à Tarse, au temps de la lutte entre Antoine et Octave, les partisans d'un certain Boëthus, mauvais poète et ami d'Antoine, banni de la cité suite à ses exactions, couvrirent les murs d'inscriptions injurieuses contre le vieil Athénodore, qui avait chassé Boëthus, dont celle-ci : `ἔργα νέων, βουλαὶ δὲ μέσων, πορδαὶ δὲ γερόντων', « Aux jeunes l'action, aux adultes le conseil, au vieux le pet », en substituant πορδαί aux \italique{prières} (εὐχαί) du vers original d'Hésiode, entré dans la langue proverbiale (Tosi n° 1081). Athénodore avait pris la chose en riant et avait fait inscrire à son tour: « au vieux, le tonnerre vengeur » (βρονταί), faisant évidemment allusion au pouvoir qu'il détenait. Enfin, il existait aussi l'expression «~péter au nez de quelqu'un~», Aristophane, dans un de ses éclairs de génie dont il est coutumier, a fait se rencontrer ces deux images (\italique{Nuées}, vers 293 sq.) : καὶ βούλομαι ἀνταποπαρδεῖν πρὸς τὰς βροντάς. Ce n'est pas autrement que s'explique ce passage aussi. (Noter le préfixe ἀντ- de l'hapax ἀνταποπέρδω, rappel de l'ἀντὶ du proverbe.)
%Il y avait en grec un proverbe, βὴξ ἀντὶ πορδῆς, « une toux à la place d’un pet », qu’on évoquait lorsque quelqu'un tentait vainement de camoufler une chose honteuse. Divers parémiographes\marqueur{7} l'ont recensé.%
%\citComm{Βὴξ ἀντὶ πορδῆς : ἐπὶ τῶν ἐν ἀπορίᾳ προσποιουμένων τι πράττειν. παρ’ ὅσον οἱ πέρδοντες λανθάνειν πειρώμενοι προσποιοῦνται βήττειν.\\
%Une toux à la place d’un pet : en référence à ceux qui, dans une situation difficile, feignent de faire quelque chose ; puisque lorsque les gens pètent, ils feignent de tousser pour essayer de se cacher.}%
%On peut inférer l’existence de l’adage cousin, βροντὴ ἀντὶ πορδῆς, d’après un passage des \italique{Nuées} que nous allons voir tout à l’heure ainsi qu’une anecdote amusante rapportée par Strabon (14, 5, 14)\marqueur{8}. 
%À Tarse, au temps de la lutte entre Antoine et Octave, les partisans d’un certain Boëthus, filou et mauvais poète de surcroît, furent bannis de la cité par le vieil Athénodore.%
%\citComm{οἱ δὲ πρῶτον μὲν κατετοιχογράφησαν αὐτοῦ τοιαῦτα\\
%‘ἔργα νέων, βουλαὶ δὲ μέσων, \textbf{πορδαὶ} δὲ γερόντων’.\\
%ἐπεὶ δ᾽ ἐκεῖνος ἐν παιδιᾶς μέρει δεξάμενος ἐκέλευσε παρεπιγράψαι ‘\textbf{βρονταὶ} δὲ γερόντων’.\\
%Avant de sortir, les bannis couvrirent les murs d’inscriptions injurieuses dans le genre de celle-ci :\\
%« Aux jeunes l’action ; aux adultes le conseil ; aux vieux le PET ».\\
%Athénodore, prenant la chose en riant, avait donné ordre qu’on écrivît à côté de l’injure : « Aux vieux le TONNERRE VENGEUR ».}%
%Les partisans de Boëthus avaient détourné un mot célèbre d’Hésiode en substituant πορδαί aux prières (εὐχαί) du vers original\marqueur{9}. Athénodore y avait répondu avec esprit : %prenant à rebours le proverbe βροντὴ ἀντὶ πορδῆς, il fit souligner que ces πορδαί ne lui étaient reprochées par les partisans de Boëthus que pour cacher leur ressentiment vis-à-vis de sa 
%il n'a pas commis de pets mais une décision \italique{foudroyante}, βρονταὶ renvoyant au pouvoir détenu par Athénodore. %Gardons dans un coin de notre tête ce proverbe – comme les contemporains d’Aristophane l’avaient et pouvaient l’identifier à la volée. 
%Socrate affirme à Strepsiade que ses Nuées sont les seules divinités qui soient et que tout le reste n’est que du vent (v. 365) :%
%\citComm{Σωκράτης --- αὗται γάρ τοι μόναι εἰσὶ θεαί, τἄλλα δὲ πάντ᾽ ἐστὶ φλύαρος.}%
%Strepsiade se montre sceptique : et Zeus l’Olympien ? n’est-il pas dieu ? Socrate lui explique qu’il n’est pas de Διός mais seulement un Δῖνος, un tourbillon, qui avec les Nuées explique la pluie\marqueur{10} comme le tonnerre. Il se sert pour cela d’une image : quand le ventre est plein, il se met à tonner. De même lorsque les Nuées gorgées d’eau se ruent l’une sur l’autre. Conclusion de Strepsiade (v. 394) :%
%\citComm{Στρεψιάδης --- ταῦτ᾽ ἄρα καὶ τὠνόματ᾽ ἀλλήλοιν βροντὴ καὶ πορδὴ ὁμοίω.\\
%Strepsiade --- C’est donc pour cela que les mots ‘toux’ et ‘pêt’ sont associés.}%
%C’est bien en regard du proverbe identifié plus haut que Strepiade rapproche les mots βροντή et πορδή et non d’une quelconque homophonie comme on l’a soutenu absurdement jusqu’ici. 
%Enfin, il existait aussi l’expression « péter au nez de quelqu’un » pour signifier mépriser la personne\marqueur{11}, Aristophane, dans un de ses éclairs de génie dont il est coutumier, a fait se rencontrer ces deux images : entendant l’arrivée tonitruantes des Nuées, Strepsiade veut les « honorer » d’un tonnerre de flatulences (v. 293) :%
%\citComm{Στρεψιάδης --- καὶ σέβομαί γ᾽ ὦ πολυτίμητοι καὶ βούλομαι ἀνταποπαρδεῖν πρὸς τὰς βροντάς.\\
%Strepsiade --- Je vous honore, Nuées très vénérées, et je veux répondre par une pétarade aux coups de tonnerre.}%
%On comprend bien ce que cet hommage recèle d’irrévérencieux ; Socrate de gourmander aussitôt (v. 296) :%
%\citComm{Σωκράτης --- οὐ μὴ σκώψει μηδὲ ποιήσεις ἅπερ οἱ τρυγοδαίμονες οὗτοι.\\
%Socrate --- Garde de plaisanter ni de faire comme ces pitoyables auteurs comiques.}%
%%Noter la présence d'une synérèse\marqueur{3} entre πρανοῖ et ἡμέτερος.
}%
\debutNote{49-50}%
%\note{1}{Aux vers 117 ou 1174, le manuscrit de Ravenne substitue un `ο' à l'`α' de l'\allemand{Urtext} d'Aristophane ; la correction de Grynaeus consiste à supposer l'erreur inverse.}%
\note{1}{Sommerstein voit en πρόμος une légère déformation suggérant à l'oreille πόρνος... « Notre chef de file --- Celui qu'on enfile ? » Mais que c'est loin en grec !.. Certains ont voulu justifier le rapprochement πρόμος -- πόρνος en arguant de la plaisanterie de \italique{Nuées} (v. 394) sur l'indissociabilité entre βροντή et πορδή -- certes pas au niveau de la prononciation. Nous avons profité des annexes auxquelles nous renvoyons le lecteur pour exégéter ce jeu de mots qui roule sur la langue proverbiale.}%
%\note{3}{Ou d'une \italique{synalèphe} pour reprendre le vocabulaire de Koster dans son \italique{Traité de métrique grecque}. D'autres parlent de \italique{synizèse} ; on croise même parfois le mot \italique{crase}... Que les grammairiens se mettent d'accord pour nous proposer un jargon qui soit sans ambiguïté.}%
\note{2}{Même genre de variantes orthographiques entre πρᾷος et πρᾶος.}%
\note{3}{Les δρύοχοι sont les \italique{chantiers}, ou poutres sur lesquelles s'établit la quille d'un vaisseau en construction. On rapprochera la métaphore de \italique{Guêpes} 30 : τὴν τρόπιν τοῦ πράγματος, « la base de l'affaire ». (Note de Taillardat §758 n. 1 ; voir aussi §673.)}%
\note{4}{Cf. Hésychius π 3199, π 3269 πρηνιχθέντα· ἐπὶ πρόσωπον καὶ ἐπὶ στόμα πεσόντα ~ Photius π 1138 et L'Etymologicum genuinum : πρηνίξαι πρηνιχθῆναι καὶ πεπρηνίχθαι· γνὺξ γάρ ἐστιν κυρίως τὸ εἰς γόνυ πεσεῖν ... τὸ δὲ πρηνιχθῆναι ἐπὶ τὸ πρόσωπον [πεσεῖν] et d'autres références dans l'abondant et récent apparat des témoins de l'édition Theodoridis de Photius.}%
\note{5}{Les écrits anciens comportaient de nombreuses abréviations. (Cf. par exemple la très ingénieuse reconstruction du vers 80 par Nauck.) Par exemple, προσ était souvent abrégé d'une sorte de Ʒ renversé qui ressemblait à un epsilon : \renverse{Ʒ}. %Faut-il supposer l'existence de la même abréviation pour προ qui expliquerait sa confusion avec ε 
Voir G. Bady, \italique{Paléographie grecque : éléments de cours}, ou bien V. Gardthausen, \italique{Griechische palaeographie}\expo{2}, II (1913), p. 340.}%
%\note{7}{Comme Diogénianus (III, 62), Érasme (\italique{Adages}, n° 563) ou Suidas (β 261). Nous citons la \italique{Souda}.}%
%\note{8}{Le parémiographe Macaire (IV, 11) a recensé la formule des partisans de Boëthus}%
%\note{9}{Vers qui depuis longtemps était rentré dans la langue proverbiale. Voir Tosi, \italique{Dictionnaire des sentences latines et grecques}, §1081.}%
%\note{10}{‘ἀλλὰ τίς ὕει ;’ (v. 368). En grec, il pleut, il tonne, il neige, etc., ne sont pas proprement des impersonnels : on sous-entend le sujet, Zeus (exprimé dans Homère). [Note de Willems.] Cette explication sera évidemment comprise de travers par Strepsiade qui pense que dans la lutte des dieux, le tourbillon a chassé Zeus (vv. 828, 1471). Strepsiade finit par apostasier son nouveau dieu qu’il ne considère plus qu’être une « vile poterie » (χυτρεοῦν, v. 1474 ; le δῖνος pouvant désigner un vase à boire).}%
%\note{11}{Voir Érasme (\italique{Adages} n° 676, \latin{Oppedere et oppedere contra tonitrua}) et les exemples qu’il cite : \italique{Ploutos}, v. 618 ; Horace, \italique{Satires}, I, 9, 69-70. On trouve aussi cet adage remanié dans le genre des actions vaines : « péter au nez d’un sourd », παρὰ κωφὸν ἀποπαρδεῖν. (Voir l’érudite note qu’a consacrée Tosi à ce dernier, §1898, \latin{non canimus surdis}.)}%
\\\passeLigne}%
\end{groupe}%
%% v. 50b
\begin{groupe}%
%%
\personnageGrec{Μνησίλοχος}%
\dimAnap{\ligne{50}\decalage{πραῖνοι ἡμέτερος—}μῶν βινεῖσθαι;\\\passeLigne}%
\end{groupe}%
%% v. 51a
\begin{groupe}%
\personnageGrec{Θεράπων}%
\dimAnap{τίς ὁ φωνήσας;\\\passeLigne}%
\end{groupe}%
%% v. 51b
\begin{groupe}%
\personnageGrec{Μνησίλοχος}%
\dimAnap{\ligne{51}\decalage{τίς ὁ φωνήσας;}νήνεμος αἰθήρ.\\\passeLigne}%
\end{groupe}%
%% v. 52-57a
\begin{groupe}%
\personnageGrec{Θεράπων}%
\dimAnap{δρυόχους, τιθέναι δράματος ἀρχάς.%\testimonia{52-7}{Suidas δ 1547 (sed τορεύει pro τορνεύει, χιανεύει pro χοανεύει)}\testimonia{52}{Zonaras 569 δρύοχοι}
\\
κάμπτει δὲ νέας ἁψῖδας ἐπῶν,%\testimonia{53}{(κάμπτει) Hesychius κ 617 κάμπτειν}
\apparat{53}{ἀσπίδας R : \lecon{ἁψῖδας} Suid. δ 1547, λ 181}\\
τὰ δὲ τορνεύει, τὰ δὲ κολλομελεῖ,%\testimonia{54}{(τορνεύει) Hesychius τ 1168 τορνεύει}
\\
καὶ γνωμοτυπεῖ κἀντονομάζει\\
καὶ κηροχυτεῖ καὶ γογγύλλει%\testimonia{56}{(γογγύλλει) Hesychius γ 764 γογγύλλειν \espTest Suidas γ 361 $=$ Zonaras 449 γογγυλίζει}
\apparat{56}{γογγύλίζει R, Suid. γ 361 : \lecon{γογγύλλει} Porson : γογγυλίει Hermann (ad \italique{Nub.} 992)}\\
καὶ χοανεύει—%\testimonia{57}{(χοανεύει) \Eustathe 1153.40}%
\commentaire{52-7}{τιθέναι δράματος ἀρχάς... καὶ χοανεύει.}{Agathon le tragique n'est pas un génie bâtisseur, un τέκτων ; il ne crée pas : il fabrique, il bricole. Il est donc tout-à-tour charpentier de marine, menuisier, modeleur : « il va poser des chantiers, fondements d'un drame. etc. » Ces métaphores sont évidemment burlesques ; elles avilissent délibérément l'art d'Agathon, ravalé au rang de métier et ne peuvent que tourner le poète en dérision...\\
Agathon est donc à mille lieues du mystère de l'inspiration, décrit par Platon (dans l'\italique{Ion}) : κοῦφον γὰρ χρῆμα ποιητής ἐστιν καὶ πτηνὸν καὶ ἱερόν, « le poète est chose légère, ailée, sacrée ». De fait, Platon connaît bien les images du menuisier et du modeleur\marqueur{1} mais jamais il ne les applique au poète (ἀποτορνεύειν est dit de Lysias, \italique{Phèdre} 243e : ἀκριβῶς ἕκαστα τῶν ὀνομάτων ἀποτετόρνευται et reste péjoratif. La même observation vaut pour κολλᾶν, en parlant d'un auteur, \italique{Ménexène} 263b ; \italique{Phèdre} 278e). À noter que ces images semblent banales ; ainsi, dans le \italique{Timée} (81b), on trouve : οἷον ἐκ δρυόχων, « comme depuis les longerons de la carcasse », c'est-à-dire \italique{entièrement}. κάμπτειν : \italique{courber}, \italique{plier} désigne en musique les courbes capricieuses d'un air ou les inflexions de la voix. (Notes de Taillardat\marqueur{2}.)}%
\debutNote{52-7}%
\note{1}{Voir Pierre Louis, \italique{Les Métaphores de Platon}, Belles Lettres, 1945, p. 98, 170, 204-5.}%
\note{2}{\italique{Les images d'Aristophane}, §758 et §784.}\\\passeLigne}%
\end{groupe}%
%% v. 57b
\begin{groupe}%
\personnageGrec{Μνησίλοχος}%
\dimAnap{\ligne{57}\decalage{καὶ χοανεύει—}καὶ λαικάζει.
\commentaire{56-7}{καὶ κηροχυτεῖ καὶ γογγύλίζει / καὶ χοανεύει — καὶ λαικάζει.}{La métaphore du serviteur est tirée de la technique de la \italique{cire perdue}\marqueur{1} ; un procédé de moulage pour obtenir une statue en métal à partir d'un modèle en cire : on commence par modeler à la cire la statue (κηροχυτεῖν), modèle qu'on recouvre d'une gangue d'argile qui fera office de moule. C'est à cette opération que renvoie probablement le verbe γογγυλίζειν (corr. γογγύλλειν), associant la gangue à une \italique{rave ronde}, γογγυλίς. % ; comme on parlerait d'une \italique{coque} en français ? (gangue et γογγυλίς ne semblent pas avoir de lien étymologique.)
Il faut laisser à ce moule d'argile des cheminées, certaines pour insérer des entonnoirs et y couler le métal en fusion (χοανεύειν) qui va prendre la place de la cire, et d'autres cheminées pour permettre à la cire fondue d'être évacuée.
Cet esprit tordu de Mnésiloque y voit la description d'une scène érotique.
Il interrompt d'une obscénité: « καὶ λαικάζει ».
Ce verbe a le sens de \latin{fellare}\marqueur{2}.% ; si le lecteur ne comprend pas pourquoi, alors tant mieux pour lui\marqueur{3}.
%L'on comprend en quoi ce pétrissage de cire puis d'introduction du moule ait pu lui suggérer cette image. Nous revenons dans la note suivante sur l'interprétation scabreuse par Mnésiloque du verbe γογγύλλειν\marqueur{3}.
%On comprend alors l'interprétation détournée que fait Mnésiloque des trois verbes précédents. Pour γογγύλλειν, les commentateurs \& traducteurs donnent le sens d'\italique{arrondir} ; mais si l'on garde \latin{fellare} en tête, il semble bien préférable de lui donner le sens de \italique{rendre ramassé}, \italique{compact}, \italique{serré}, autrement attesté\marqueur{2}.
}%
\debutNote{56-7}%
\note{1}{Voir sur cette technique l'article de C. C. Mattusch « The Berlin Foundry Cup: The Casting of Greek Bronze Statuary in the Early Fifth Century B. C. », \italique{American Journal of Archaeology}, 84 n° 4 (1980), pp. 435-444, ou son livre \italique{Greek Bronze Statuary: From the Beginnings Through the Fifth Century B.C.} (1988).}%
\note{2}{\italique{Confer} par exemple Chantraine, \italique{Dictionnaire étymologique de la langue grecque}. γογγύλλειν peut avoir le sens de \italique{rendre compact}, \italique{serré}. Cf. la \nieme{3} note §166 de Taillardat : \italique{ramassé}, \italique{compact}, \italique{serré} sont les sens que Taillardat donne à l'adjectif γογγύλος aux vers 25-28 de la \italique{Paix}.}%
%\note{3}{γογγύλλειν peut avoir le sens de \italique{rendre compact}, \italique{serré}. Cf. la \nieme{3} note §166 de Taillardat : \italique{ramassé}, \italique{compact}, \italique{serré} sont les sens que Taillardat donne à l'adjectif γογγύλος aux vers 25-28 de la \italique{Paix}.}%
%\note{2}{Cf. Hésychius (γ 764) : γογγύλλειν· συστρέφειν et la \nieme{3} note §166 de Taillardat : \italique{ramassé}, \italique{compact}, \italique{serré} sont les sens que Taillardat donne à l'adjectif γογγύλος aux vers 25-28 de la \italique{Paix}.}%
\\\passeLigne}%
\end{groupe}%
%% v. 58
\begin{groupe}%
\personnageGrec{Θεράπων}%
\dimAnap{τίς ἀγροιώτας πελάθει θριγκοῖς;%\testimonia{58}{(θριγκοῖς) Hesychius θ 750 θριγκοῖς· φραγμοῖς}
\apparat{58}{ἀγριώτας R : \lecon{ἀγροιώτας} anonym. Parisinus}\\\passeLigne}%
\end{groupe}%
%% v. 59-61
\begin{groupe}%
\personnageGrec{Μνησίλοχος}%
\dimAnap{ὃς ἕτοιμος σοῦ τοῦ τε ποιητοῦ\\
τοῦ καλλιεποῦς κατὰ τοῦ θριγκοῦ\apparat{60}{κατὰ om. \Rac : add. \Rpc}\\
συγγογγύλας καὶ συστρέψας\apparat{61}{γογγυλίσας R : %συγγογγυλίσας Brunck : 
\lecon{συγγογγύλας} Enger}%
\commentaire{61}{γογγυλίσας καὶ συστρέψας.}{Vers retouché \italique{gratia metri} d'après un passage de \italique{Lysistrata} (v. 975), jouée probablement l'année précédente (cf. Introduction).
Allons-y voir d'un peu plus près pour comprendre précisément l'image de Mnésiloque : Kinésias est dans un état lamentable après que Myrrhine l'ait fortement échaudé... et finalement laissé raide. Le demi-chœur des hommes en vient à souhaiter un traitement de faveur pour cette s[célérate] (v. 972-979) : %
\citComm{μιαρὰ δῆτ' ὦ Ζεῦ Ζεῦ· |
εἴθ' αὐτὴν ὥσπερ τοὺς θωμοὺς |
μεγάλῳ τυφῷ καὶ πρηστῆρι |
\textbf{ξυστρέψας καὶ ξυγγογγύλας} |
οἴχοιο φέρων, εἶτα μεθείης, |
ἡ δὲ φέροιτ' αὖ πάλιν ἐς τὴν γῆν, |
κᾆτ' ἐξαίφνης |
περὶ τὴν ψωλὴν περιβαίη.\\
Une gueuse ! ô Zeus, Zeus, puisse-t-elle, comme les tas de blé dans un coup de vent et une tourmente, roulée par toi [Kinésias] et ramassée en boule, être portée dans l'air, puis, lâchée, retomber à terre et du coup enfourcher ce priape.\\
\textbf{975} ξυγγογγυλίσας R, Γ, Suid. : \lecon{ξυγγογγύλας} Cobet}%
L'image est tirée de la technique du foulage et du vannage du grain. Voici ce qu'en dit le Daremberg \& Saglio\marqueur{1} :%
\citComm{Foulage (ἀλόησις) --- Les Grecs ne battaient les céréales ni en grange, ni en plein air, car ils n'avaient point de fléau pour faire sortir le grain de l'épi. Ils employaient la vieille méthode égyptienne du dépiquage ou foulage, qui n'est plus guère en usage que chez les Provençaux et les Languedociens. En juin, c’est-à-dire, à l'apparition d’Orion\marqueur{2}, on portait les gerbes sur une aire, qualifiée de μεγάλη\marqueur{3} , ἐϋκτίμενη\marqueur{4}, ἐϋτροχάλος\marqueur{5}, ἱερά\marqueur{6}, et construite, plus ou moins, d'après le modèle de l'aire sacrée d'Éleusis\marqueur{7}. Des ouvriers, ἐπαλώστες\marqueur{8}, dressent les gerbes a côté les unes des autres en les inclinant un peu vers le centre et on fait entrer les animaux : bœufs, chevaux, mulets, parfois même des ânes\marqueur{9} ; ils sont accouplés deux à deux et marchent en cercle, d’abord au pas, puis au trot, sous la direction d'un homme, placé au centre de l'aire, qui tient en main les guides, δεσμά\marqueur{10}. D'autres ouvriers, armés de fourches en bois, τρίναξ\marqueur{11}, régularisent la foulaison, ὁμαλιεῖται ὁ ἁλοητός, en retournant la paille et en la ramenant sous les pieds des animaux, pour faciliter le manège, δῖνος.\\
Nettoyage du grain (λίκμησις)  — Les Grecs ne connaissaient point le vannage proprement dit, opération de nos batteurs en grange, qui nettoye simplement le grain des glumes ou glumelles et que l'on exécute près de la porte d’une grange, dans un léger courant d'air, mais à l'abri du vent. Au contraire, les Grecs attendaient que le vent fût assez fort\marqueur{12} pour emporter les fétus de paille et leur permît de pratiquer le \italique{ventage}, qui peut se faire de deux façons :  1°  Avec une corbeille, λικμός, πλόκανον [\latin{Ventilabrum}] [...]    2°   Avec une pelle en bois, πτύον [\latin{Pala}]. C'est notre \italique{nettoyage à la roue} ; on jette circulairement contre le vent et, aussi loin que possible, à la hauteur d’un mètre au moins, une pelletée de grains qui retombent sur le sol, tandis que la paille et les balles sont emportées à l'extrémité de l'aire.}%
Ces longues mais utiles précisions ayant été apportées, le passage de \italique{Lysistrata} est maintenant clair : Myrrhine doit être traitée comme un fétu de paille qu'on retourne et compactifie en boule (ξυστρέψας καὶ ξυγγογγύλας\marqueur{13}) avant de le jeter (εἶτα μεθείης) pour qu'il jonche l'aire vide ; ici ἡ ψωλή $=$ τὸ κενόν\marqueur{8}, mais avec un double sens obscène : le [phallos] \italique{nu}, c'est-à-dire décalotté.\\
Revenons aux \italique{Thesmophories}. Le serviteur demande qui est le paysan (τίς ἀγροιώτας) qui se permet d'interrompre de la sorte son poème. Mnésiloque lui répond donc par une image agricole certainement populaire : il se prépare à le retourner et le mettre en boule (συγγογγύλας καὶ συστρέψας v. 61) aussi rustiquement que le chœur des vieillards voulait faire avec Myrrhine.%
}%
\debutNote{61}%
\note{1}{\italique{Dictionnaire des antiquités grecques et romaines}, IV-2, pp. 906-907. Article « \latin{Rustica Res} ».}%
\note{2}{Hésiode, \italique{Travaux et les jours}, v. 596.}%
\note{3}{\italique{Iliade}, XIII, v. 588.}%
\note{4}{\italique{Iliade}, XX, v. 496.}%
\note{5}{Hésiode \italique{Travaux et les jours}, v. 597.}%
\note{6}{\italique{Iliade}, V, v. 499.}%
\note{7}{Pausanias, I, 38.}%
\note{8}{Xénophon, \italique{Économiq.}, XVIII.}%
\note{9}{\italique{Anthol. palat.} IX, 301.}%
\note{10}{\italique{Anthol. palat.} VI, 104.}%
\note{11}{\italique{Anthol. palat.} VI, 95.}%
\note{12}{\italique{Iliade}, V, vv. 499 sq. ; XIII, vv. 588 sq. ; cf. \italique{Anth. palat.}, VI, 53.}%
\note{13}{Soit dit en passant, c'est certainement l'un de ces deux passages que glose Hésychius (γ 764) : γογγύλλειν· συστρέφειν.}%
}%
\end{groupe}%
%% v. 62
\begin{groupe}%
\typeVers{(catalect.)}\dimAnap{τουτὶ τὸ πέος χοανεῦσαι.%\testimonia{62}{(χοανεῦσαι $=$ χωνεῦσαι) Moeris χ 21 ; Hesychius χ 575 χοανεῦσαι}
\\\passeLigne}%
\end{groupe}%
%% v. 63
\begin{groupe}%
\personnageGrec{Θεράπων}%
\typeVers{3 ia.}\trimIamb{ἦ που νέος γ᾽ ὢν ἦσθ᾽ ὑβριστὴς ὦ γέρον.\\\passeLigne}%
\end{groupe}%
%% v. 64-65
\begin{groupe}%
\personnageGrec{Εὐριπίδης}%
\trimIamb{ὦ δαιμόνιε τοῦτον μὲν ἔα χαίρειν, σὺ δὲ\\
Ἀγάθωνά μοι δεῦρ᾽ ἐκκάλεσον πάσῃ τέχνῃ.
%\commentaire{65}{πάσῃ τέχνῃ.}{\eff{\italique{Bon gré, mal gré} (traduction de Willems.)}}
\\\passeLigne}%
\end{groupe}%
%% v. 66-69
\begin{groupe}%
\personnageGrec{Θεράπων}%
\trimIamb{μηδὲν ἱκέτευ᾽· αὐτὸς γὰρ ἔξεισιν τάχα.\\
καὶ γὰρ μελοποιεῖν ἄρχεται· χειμῶνος οὖν\\
ὄντος κατακάμπτειν τὰς στροφὰς οὐ ῥᾴδιον,
\commentaire{67-8}{χειμῶνος οὖν / ὄντος κατακάμπτειν τὰς στροφὰς οὐ ῥᾴδιον.}{Les Thesmophories avaient lieu du 11 au 13 du mois de Pyanepsion, jours qui tombent en général pendant notre seconde moitié du mois d'octobre.
%Jeu de mots sur στροφὰς qui, littéralement, sont les \italique{tours}.
Le serviteur continue de filer sa métaphore : la cire est plus facile à modeler lorsqu'elle est exposée au Soleil.
}\\
ἢν μὴ προίῃ θύρασι πρὸς τὸν ἥλιον.\apparat{69}{\lecon{θύρασι} R : θύραζε Zanetti}\\\passeLigne}%
\end{groupe}%
%% v. 70a
\begin{groupe}%
\personnageGrec{Μνησίλοχος}%
\trimIamb{τί οὖν ἐγὼ δρῶ;%
\commentaire{70-4}{τί οὖν ἐγὼ δρῶ; ... ὄντα κηδεστὴν ἐμόν.}{La répartition des ce passage diffère d'un éditeur à un autre ; chacun voulant apporter son grain de sel. Inscrivons-nous, nous aussi, dans cette grande tradition philologique et apportons notre pierre à l'édifice : une nouvelle répartition des répliques qui, espérons-le, clôturera un débat qui n'a que trop duré. La répartition des deux répliques du vers 70 ne pose aucun problème : il suffit de songer aux caractères des trois personnages en présence : le ton empoté du « τί οὖν ἐγὼ δρῶ; » ne peut convenir qu'au personnage pour l'instant passif et qui subit l'action : Mnésiloque\marqueur{1} (Euripide ruisselle de défauts mais garde toujours une certaine dignité ; jamais il ne s'abaisserait à lancer un mot aussi trivial et décontenancé – ses plaintes ont davantage de lyrisme !) La réponse : « περίμεν᾽, ὡς ἐξερχεται » ne peut être attribuée qu'à Euripide : c'est lui qui a besoin de Mnésiloque ; de plus, le serviteur ne délivre ni conseil ni ordre : il se contente de répondre le plus sèchement possible aux questions (cette répartition est d'ailleurs celle du manuscrit de Ravenne). Les vers 72 à 74 sont évidemment prononcés par Mnésiloque qui cherche à percer la détresse d'Euripide et nous révèle par la même occasion son lien de parenté avec le poète. Reste le vers 71 : il est plus que tentant de l'attribuer à Mnésiloque et faire ainsi des quatre vers 71-74 l'exact pendant des quatre vers 1-4 : une invocation à Zeus suivie d'une question : là pour savoir où Euripide le mène et ici pour savoir ce qui le tourmente. Problème : la présence d'un \italique{signe} (σημεῖον), probablement un \italique{paragraphos}, sur d'anciens manuscrits a été interprétée comme indiquant un changement d'interlocuteurs entre le vers 71 et les vers 72-74. Mais tel n'est pas le cas ici, et c'est ce que nous indiquait une ancienne scholie, très précieuse et incomprise\marqueur{2} : 
\citComm{τὸ σημεῖον· οὐ μόνον κηδεστὴς ὁ ἐκδιδούς, ἀλλὰ καὶ ὁ λαμβάνων.\\
\francais{Le signe [typographique] : non seulement le Parent est celui qui abandonne [la parole], mais aussi celui qui la [re]prend.}}%
La scholie ne précise pas à quoi servait ce σημεῖον ; serait-ce une pause ? quoi qu'il en soit, le protagoniste devait certainement faire quelque jeu de théâtre après l'invocation à Zeus du vers 71, identique à celui déjà fait après le premier vers et montrer ainsi aux spectateurs le parallélisme entre les deux quatrains.}%
\debutNote{70-4}%
\note{1}{Réplique analogue de Mnésiloque au vers 925 : ἐγὼ δ᾽ ὁ κακοδαίμων τί δρῶ;}
\note{2}{Incomprise déjà de Suidas – lexicographe hypothétique du \nieme{IX} ou \nieme{X} siècle, à qui l'on attribue le lexique de la \italique{Souda}. Ce dernier fourmille en références et citations d'Aristophane : la \italique{Souda} est un outil de recherche sur Aristophane indispensable. 
Il en existe une version numérique : \href{https://www.cs.uky.edu/~raphael/sol/sol-html/}{https://www.cs.uky.edu/~raphael/sol/sol-html/} – incomprise donc de Suidas ce qui tend à prouver son ancienneté : Suidas (κ 1487) biffe le lemme `τὸ σημεῖον', pourtant indispensable à la bonne compréhension et égare ainsi les différents éditeurs des scholies qui l'appliquent au vers 74, voire qui la retouchent pour finir de la rendre tout-à-fait inintelligible.}\\\passeLigne}%
\end{groupe}%
%% v. 70b
\begin{groupe}%
\personnageGrec{Εὐριπίδης}%
\trimIamb{\ligne{70}\decalage{τί οὖν ἐγὼ δρῶ;}περίμεν᾽, ὡς ἐξερχεται.\\\passeLigne}%
\end{groupe}%
%% v. 71
\begin{groupe}%
\personnageGrec{Μνησίλοχος}\trimIamb{ὦ Ζεῦ τί δρᾶσαι διανοεῖ με τήμερον;\apparat{71-4}{(Μνη.) (71) et (Εὐρ.) (72-4) R : \lecon{(Μνη.) (71-4)} Bibi ap. Σ : alii alia}%
\commentaire{71}{ὦ Ζεῦ τί δρᾶσαι διανοεῖ με τήμερον;}{Formule courante ; se reporter, par exemple, aux \italique{Oiseaux}, vv. 58 \& 62.}}%
\end{groupe}%
%% v. 72-74
\begin{groupe}%
\trimIamb{νὴ τοὺς θεοὺς ἐγὼ πυθέσθαι βούλομαι\\
τί τὸ πρᾶγμα τουτί. τί στένεις; τί δυσφορεῖς;\\
οὐ χρῆν σε κρύπτειν ὄντα κηδεστὴν ἐμόν.%\testimonia{74}{Suidas κ 1487, cf. Moeris κ 38}
\apparat{74}{\lecon{χρῆν} R, Suid. κ 1487 codd. FV : χρή Suid. codd. AGM \espApp \lecon{ἐμόν} R : ἐμέ Valckenaer ap. Dobree (p. 237)}\\\passeLigne}%
\end{groupe}%
%% v. 75
\begin{groupe}%
\personnageGrec{Εὐριπίδης}%
\trimIamb{ἔστιν κακόν μοι μέγα τι προπεφυραμένον.
\commentaire{75}{προπεφυραμένον}{Pour comprendre cette image, l'on se rapportera à un passage des \italique{Oiseaux} (v. 462-463) et le commentaire qu'en fait Taillardat (§755) :
\citComm{καὶ προπεφύραται λόγος εἷς μοι, / ὃν διαμάττειν οὐ κωλύει\\
\francais{j'ai sous la main un discours déjà bien délayé et rien n'empêche que j'achève de le pétrir.}}%
« Φυρᾶν signifie proprement \italique{délayer} de la farine dans un liquide\marqueur{1}. » Est donc προ-πεφυραμένον ce qui a déjà été délayé à l'avance mais point encore pétri ; autrement dit : le complot est ourdi mais rien n'est encore joué. L'avantage de ce genre de métaphore boulangère est qu'elle se transpose assez aisément en français\marqueur{2}.}%
\debutNote{75}%
\note{1}{Cf. Hippocrate \italique{Du régime} (\italique{Περί διαίτης}) 2,40 et Thucydide 3,49.}%
\note{2}{Comme l'épisode raconté par Démosthène dans les \italique{Cavaliers} (v. 55-57) : il s'est fait voler par Cléon les pains qu'il avait distribués... alors que c'est lui qui avait mis la main à la pâte ! Le verbe \italique{pétrir} et le \italique{pétrin} ont donné nombre d'expressions en français.}\\\passeLigne}%
\end{groupe}%
%% v. 76a
\begin{groupe}%
\personnageGrec{Μνησίλοχος}%
\trimIamb{ποῖόν τι;\\\passeLigne}%
\end{groupe}%
%% v. 76b-77
\begin{groupe}%
\personnageGrec{Εὐριπίδης}%
\trimIamb{\ligne{76}\decalage{ποῖόν τι;}τῇδε θἠμέρᾳ κριθήσεται\\
εἴτ᾽ ἔστ᾽ ἔτι ζῶν εἴτ᾽ ἀπόλωλ᾽ Εὐριπίδης.
\commentaire{76-7}{τῇδε θἠμέρᾳ κριθήσεται / εἴτ᾽ ἔστ᾽ ἔτι ζῶν εἴτ᾽ ἀπόλωλ᾽ Εὐριπίδης.}{Le style est paratragique (avec l'expression poétique `τῇδε θἠμέρᾳ' plutôt que la forme plus courante τήμερον). Apprécier surtout à sa juste valeur l'allitération ! Lyrisme ? emphase ? Les `ttttttttt' saccadent-ils le vers ou lui donnent-ils du rythme ? sans doute un peu des deux. (Euripide est souvent à la frontière du sublime et du ridicule, frontière qu'il passe et repasse allègrement).}\\\passeLigne}%
\end{groupe}%
%% v. 78-80 (part 1)
\begin{groupe}\personnageGrec{Μνησίλοχος}
\trimIamb{καὶ πῶς; ἐπεὶ νῦν γ᾽ οὔτε τὰ δικαστήρια\\
μέλλει δικάζειν οὔτε βουλῆς ἐσθ᾽ ἕδρα,\\\passeLigne}%
\end{groupe}%
\endnumbering
\end{Leftside}
\begin{Rightside}
\beginnumbering
%% v. 43+ (part 2)
\begin{groupe}
\dimAnapFr{\ligne{43}Que retienne ses haleines l'éther paisible,\testimonia{43-5}{(ἐχέτω -- γλαυκόν) Suidas ε 3792 ; Zonaras 941 ἐχέτω}\\
que le flot de la mer point ne retentisse\testimonia{44}{(κελαδείτω) Suidas κ 1278}\\
glauque...\\\passeLigne}%
\end{groupe}%
%% v. 45b
\begin{groupe}%
\personnageFr{Mnésiloque}%
\dimAnapFr{\ligne{45}\espace{1.5cm}Bombax !\testimonia{45}{(βομβάξ) Hesychius β 787 βομβάξ ; Suidas ε 2807 ; \Eustathe 1718.9}\\\passeLigne}%
\end{groupe}%
%% v. 45c
\begin{groupe}%
\personnageFr{Euripide}%
\dimAnapFr{\ligne{45}\espace{2.8cm}Chut ! que dit-il ?\\\passeLigne}%
\end{groupe}%
%% v. 46-48a
\begin{groupe}%
\personnageFr{Serviteur}%
\dimAnapFr{Que la gent ailée s'assoupisse,\\
que les pieds des bêtes sauvages vaguant par les bois
point ne se meuvent...\\\passeLigne}%
\end{groupe}%
%% v. 48b
\begin{groupe}%
\personnageFr{Mnésiloque}%
\dimAnapFr{\ligne{48}\espace{1.9cm}Bombalo bombax !\testimonia{48}{(βομβαλοβομβάξ) Suidas ε 2807 ; \Eustathe 790.13, 1718.13}\\\passeLigne}%
\end{groupe}%
%% v. 49-50a
\begin{groupe}%
\personnageFr{Serviteur}%
\dimAnapFr{... Car il s'apprête l'homme au beau langage, Agathon,\testimonia{49}{(καλλιεπὴς) Philostrate \italique{Vit. soph.} I 9.1 [$=$ II 12.8 Kayser]}
à aller au devant du gaillard...\testimonia{50}{(πραῖνοι) Hesychius π 3188 †πραῖνοι· πρηνίζειν, καταστρέφειν}\\\passeLigne}%
\end{groupe}%
%% v. 50b
\begin{groupe}%
\personnageFr{Mnésiloque}%
\dimAnapFr{\ligne{50}\espace{2.6cm}Et se faire enfiler ?\\\passeLigne}%
\end{groupe}%
%% v. 51a
\begin{groupe}%
\personnageFr{Serviteur}%
\dimAnapFr{Qui a parlé?\\\passeLigne}%
\end{groupe}%
%% v. 51b
\begin{groupe}%
\personnageFr{Mnésiloque}%
\dimAnapFr{\ligne{51}\espace{2.3cm}L'éther paisible.\\\passeLigne}%
\end{groupe}%
%% v. 52-57a
\begin{groupe}%
\personnageFr{Serviteur}%
\dimAnapFr{... du gaillard d'avant ; il va composer les \fin débuts d'un drame.\testimonia{52-7}{Suidas δ 1547 (sed τορεύει pro τορνεύει, χιανεύει pro χοανεύει)}\testimonia{52}{Zonaras 569 δρύοχοι}\\
Il cambre, en jantes nouvelles, des vers.\testimonia{53}{(κάμπτει) Hesychius κ 617 κάμπτειν}
Il met les uns sur le tour, il assemble les autres.\testimonia{54}{(τορνεύει) Hesychius τ 1168 τορνεύει}\\
Il frappe des sentences, fait s'opposer des mots,\\
malaxe la cire, l'entoure d'un moule,\testimonia{56}{(γογγύλλει) Hesychius γ 764 γογγύλλειν \espTest Suidas γ 361 $=$ Zonaras 449 γογγυλίζει}\\
puis introduit le godet...\testimonia{57}{(χοανεύει) \Eustathe 1153.40}\\\passeLigne}%
\end{groupe}%
%% v. 57b
\begin{groupe}%
\personnageFr{Mnésiloque}%
\dimAnapFr{\ligne{57}\espace{2.2cm}Et avale.\\\passeLigne}%
\end{groupe}%
%% v. 58
\begin{groupe}%
\personnageFr{Serviteur}%
\dimAnapFr{Quel est le manant qui approche cette \fin enceinte ?\testimonia{58}{(θριγκοῖς) Hesychius θ 750 θριγκοῖς· φραγμοῖς}\\\passeLigne}%
\end{groupe}%
%% v. 59-61
\begin{groupe}%
\personnageFr{Mnésiloque}%
\dimAnapFr{Quelqu'un qui s'apprête à te mettre en boule et te retourner, toi et ton poète au beau langage,
et à introduire par votre enceinte}%
\end{groupe}%
%% v. 62
\begin{groupe}%
\dimAnapFr{cet outil-ci.\testimonia{62}{(χοανεῦσαι $=$ χωνεῦσαι) Moeris χ 21 ; Hesychius χ 575 χοανεῦσαι}\\\passeLigne}%
\end{groupe}%
%% v. 63
\begin{groupe}%
\personnageFr{Serviteur}%
\trimIambFr{Dans ta jeunesse tu devais être un ribaud, vieillard..\\\passeLigne}%
\end{groupe}%
%% v. 64-65
\begin{groupe}%
\personnageFr{Euripide}%
\trimIambFr{(À Mnésiloque :) Diantre d'homme, laisse-le tranquille. (Au serviteur :) Toi,
va m'appeler Agathon bon gré, mal \fin gré.\\\passeLigne}%
\end{groupe}%
%% v. 66-69
\begin{groupe}%
\personnageFr{Serviteur}%
\trimIambFr{Inutile d'insister. Lui-même ne tardera pas à sortir,\\
car il commence un chant lyrique. Et de fait, en hiver,\\
il n'est pas facile de contourner ses strophes,\\
sans aller se mettre dehors au soleil.\\\passeLigne}%
\end{groupe}%
%% v. 70a
\begin{groupe}%
\personnageFr{Mnésiloque}%
\trimIambFr{Et moi, que vais-je faire ?\\\passeLigne}%
\end{groupe}%
%% v. 70b
\begin{groupe}%
\personnageFr{Euripide}%
\trimIambFr{\ligne{70}\espace{2.6cm}Demeure, puisqu'il arrive.\\\passeLigne}%
\end{groupe}%
%% v. 71
\begin{groupe}%
\personnageFr{Mnésiloque}%
\trimIambFr{Ô Zeus, que prétends-tu faire de moi aujourd'hui ?}%
\end{groupe}%
%% v. 72-74
\begin{groupe}%
\trimIambFr{Par les dieux, je tiens à savoir de quoi il retourne.\\
Qu'as-tu à geindre ? à te contrarier ?
Il n'y a pas de mystère à faire avec moi, puisque tu es mon gendre.\testimonia{74}{Suidas κ 1487, cf. Moeris κ 38}\\\passeLigne}%
\end{groupe}%
%% v. 75
\begin{groupe}%
\personnageFr{Euripide}%
\trimIambFr{Un grand malheur, manigancé d'avance, me menace.\\\passeLigne}%
\end{groupe}%
%% v. 76a
\begin{groupe}%
\personnageFr{Mnésiloque}%
\trimIambFr{De quelle sorte ?\\\passeLigne}%
\end{groupe}%
%% v. 76b-77
\begin{groupe}%
\personnageFr{Euripide}%
\trimIambFr{\ligne{76}\espace{1.7cm}Ce jour décidera\\
si Euripide vit encore ou s'il est mort.\\\passeLigne}%
\end{groupe}%
%% v. 78-80 (part 1)
\begin{groupe}\personnageFr{Mnésiloque}
\trimIambFr{Comment cela ? Car ni les tribunaux
ne sont dans le cas de juger, ni la Boulé ne siège,\\\passeLigne}%
\end{groupe}%
\endnumbering
\end{Rightside}
\end{pages}
\Pages
\clearpage\ifodd\value{page}\hbox{}\newpage\fi
\enTeteG{ΠΡΟΛΟΓΟΣ --- Codex Ravennas 429, f. 163v (Classense 165v) --- vv. 80–122}
\begin{adjustwidth}{-1.5cm}{-3.0cm}
\begin{adjustbox}{max totalheight=21.0cm, max width=16.5cm, keepaspectratio}
\begin{minipage}{16.5cm}
{\fontsize{7.2pt}{8.5pt}\selectfont
\noindent\textbf{Folio 03 --- Codex Ravennas 429, \texttt{f. 163v} (\texttt{Classense 165v}) : ΠΡΟΛΟΓΟΣ, vv. 80–122}\hfill \textit{Témoins :} $\text{R}^1, \Rpc, \Rdeux$ (\Mon)\par\vspace{2pt}
\hrule height 0.4pt\vspace{2.5pt}
}
{\fontsize{8.1pt}{9.5pt}\selectfont
\ligneRav{80}{μ\auDessus{ν}{Η}}{\difference{ἐπείπερ ἐστὶ} Θεσμοφορίων ἡ μέση·}{\textbf{80} ἐπεὶ τρίτη 'στὶ R : \lecon{ἐπείπερ ἐστὶ} Nauck (\italique{Mél.} (1866), II pp.468-71)}
\ligneRav{}{\auDessus{εὐ}{ρ}}{τοῦτ᾽ αὐτὸ γὰρ τοι κἀπολεῖν με προσδοκῶ·}{}
\ligneRav{}{}{αἱ γὰρ γυναῖκες ἐπιβεβουλεύκασί μοι}{}
\ligneRav{83}{}{κἀν Θεσμοφόροιν μέλλουσι περί \difference{μου} τήμερον}{\textbf{83} ἐμοῦ R : \lecon{μου} Biset}
\ligneRav{}{—}{ἐκκλησιάζειν ἐπ᾽ ὀλέθρῳ. : τιὴ τί δή;}{}
\ligneRav{85}{\auDessus{εὐ}{ρ}}{ὁτιὴ τραγῳδῶ καὶ κακῶς αὐτὰς λέγω·}{}
\ligneRav{86}{μ\auDessus{ν}{Η}}{νὴ τὸν Ποσειδῶ καὶ \difference{δικαίαν} ἂν πάθοις·}{\textbf{86} δί᾽ και᾽ R : δίκαιά γ᾽ Grynaeus : \lecon{δικαίαν} Wilamowitz}
\ligneRav{}{}{ἀτὰρ τίν᾽ ἐκ ταύτης σὺ μηχανὴν ἔχεις;}{}
\ligneRav{}{\auDessus{εὐ}{ρ}}{Ἀγάθωνα πεῖσαι τὸν τραγῳδοδιδάσκαλον}{}
\ligneRav{}{—}{ἐς Θεσμοφόροιν ἐλθεῖν. : τί δράσοντ᾽; εἰπέ μοι·}{}
\ligneRav{90}{\auDessus{εὐ}{ρ}}{ἐκκλησιάσοντ᾽ ἐν ταῖς \difference{γυναιξὶ} \difference{χἂν} δέῃ}{\textbf{90} γυναιξίν R : \lecon{γυναιξὶ} \Mon \espApp κἂν R : \lecon{χἂν} Markland (1763) ad Eur. \italique{Suppl.} 364}
\ligneRav{91}{—}{λέξονθ᾽ ὑπὲρ ἐμοῦ. : πότερα φανερὸν ἢ λάθρᾳ;}{\textbf{91} \lecon{φανερὸν} R : φανερῶς Meineke}
\ligneRav{}{\auDessus{εὐ}{ρ}}{λάθρᾳ, στολὴν γυναικὸς ἠμφιεσμένον·}{}
\ligneRav{}{μ\auDessus{ν}{Η}}{τὸ πρᾶγμα κομψὸν καὶ σφόδρ᾽ ἐκ τοῦ σοῦ τρόπου·}{}
\ligneRav{}{}{τοῦ γὰρ τεχνάζειν ἡμέτερος ὁ πυραμοῦς·}{}
\ligneRav{95}{\auDessus{εὐ}{ρ}}{σίγα. : τί δ᾽ ἔστιν; : Ἁγάθων ἐξέρχεται·}{\textbf{95} \lecon{τί δ᾽} R : τί Kaibel}
\ligneRav{96}{μ\auDessus{ν}{Η}}{καὶ ποῖός ἐστιν; : οὗτος οὑκκυκλούμενος·}{\textbf{96} ποῖο ἐστὶν οὗτος — οὑκκυκλούμενος R : \lecon{ποῖος ἐστὶν; — οὗτος οὑκκυκλούμενος} Tyrwhitt (p. 65)\espApp ποῖο ἐστὶν οὗτος — οὑκκυκλούμενος R : ποῖος ἐστὶν οὗτος — οὑκκυκλούμενος \Mon : ποῦ ᾽σθ᾽; — ὅπου ᾽στίν; οὗτος οὑκκυκλούμενος Meineke : \lecon{ποῖος ἐστὶν; — οὗτος οὑκκυκλούμενος} Tyrwhitt (p. 65)}
\ligneRav{97}{μ\auDessus{ν}{Η}}{ἀλλ᾽ ἦ τυφλὸς μέν εἰμ᾽· ἐγὼ γὰρ οὐχ ὁρῶ}{\textbf{97} ἢ R : \lecon{ἦ} Biset}
\ligneRav{}{}{ἄνδρ᾽ οὐδέν᾽ ἐνθάδ᾽ ὄντα, Κυρήνην δ᾽ ὁρῶ·}{}
\ligneRav{99}{\auDessus{εὐ}{ρ}}{σίγα· μελῳδεῖν νῦν παρασκευάζεται·}{\textbf{99} ἂν R : \lecon{νῦν} Rogers : alii alia}
\ligneRav{}{}{\hspace{2.2em}\textbf{[παρεπιγραφή R\expo{1} :]} \lecon{μινυρισμός}}{\textbf{post 99} \lecon{μινυρισμός} (παρεπιγραφή) R, $\Sigma^{\text{R}}$}
\ligneRav{100}{μ\auDessus{ν}{Η}}{Μύρμηκος ἀτραπούς, ἢ τί διαμινύρεται;}{\textbf{100} μύρμ- R : \lecon{Μύρμ-} Bibi ap Hesych. μ 1904}
\ligneRav{101}{ἀγάθ/}{\hspace{1.1em}ἱερὰν Χθονίαιν δε-}{\textbf{101} χθονίαις R : \lecon{χθονίαιν} Meineke}
\ligneRav{}{}{\hspace{1.1em}ξάμεναι λαμπάδα κοῦραι}{}
\ligneRav{103}{}{\hspace{1.1em}ξὺν ἐλευθέρᾳ πατρί-}{\textbf{103} \lecon{πατρίδι} R : πραπίδι Wecklein}
\ligneRav{}{}{\hspace{1.1em} δι χορεύσασθε βοάν·}{}
\ligneRav{}{}{\hspace{1.1em}τίνι δαιμόνων ὁ κῶμος;}{}
\ligneRav{105}{}{\hspace{1.1em}λέγε \difference{νιν} εὐπείστως δὲ τοὐμὸν}{\textbf{105} νυν R : \lecon{νιν} Dindorf \espApp εὐπίστως R, Suid. σ 178, 183 : \lecon{εὐπείστως} Reiske}
\ligneRav{106}{}{\hspace{1.1em}δαίμονας \difference{ἔχει} σεβίσαι·}{\textbf{106} ἔχεις R : \lecon{ἔχει} Suid.}
\ligneRav{107}{}{\hspace{1.1em}ἄγε νῦν \difference{ὄλβιζε} Μοῦσα}{\textbf{107} ὅπλιζε R : \lecon{ὄλβιζε} Bentley \espApp \lecon{Μοῦσα} R : μούσᾳ Bergk}
\ligneRav{}{}{\hspace{1.1em}χρυσέων ῥύτορα τόξων}{}
\ligneRav{}{}{\hspace{1.1em}Φοῖβον, ὃς ἱδρύσατο χώρας}{}
\ligneRav{110}{}{\hspace{1.1em}γύαλα Σιμουντίδι γᾷ·}{}
\ligneRav{}{}{\hspace{1.1em}χαῖρε καλλίσταις ἀοιδαῖς}{}
\ligneRav{}{}{\hspace{1.1em}Φοῖβ᾽ ἐν εὐμούσοισι τιμαῖς}{}
\ligneRav{}{}{\hspace{1.1em}γέρας ἱερὸν προφέρων·}{}
\ligneRav{}{}{\hspace{1.1em}τάν τ᾽ ἐν ὄρεσι δρυογόνοισι}{}
\ligneRav{115}{}{\hspace{1.1em}κόραν \difference{ἀείσατ᾽} Ἄρτεμιν ἀγροτέραν·}{\textbf{115} ἀεισαντ᾽ R : \lecon{ἀείσατ᾽} Zanetti (\italique{laudate} Divus)}
\ligneRav{}{}{\hspace{1.1em}ἕπομαι κλῄζουσα σεμνὰν}{}
\ligneRav{117}{}{\hspace{1.1em}γόνον ὀλβίζουσα Λατοῦς}{\textbf{117} σεμνὸν R : \lecon{σεμνὰν} Dindorf}
\ligneRav{}{}{\hspace{1.1em}Ἄρτεμιν ἀπειρολεχῆ·}{}
\ligneRav{}{}{\hspace{1.1em}Λατώ τε κρούματά τ' Ἀσιάδος ποδὶ}{}
\ligneRav{120}{}{\hspace{1.1em}παράρυθμ' εὔρυθμα Φρυγίων}{}
\ligneRav{122}{}{\hspace{1.1em}\difference{διὰ νεύματα} Χαρίτων·}{\textbf{122} διανεύματα R : δινεύματα Bentley : \lecon{διὰ νεύματα} Fritzsche}
}
\vspace{3pt}\noindent\rule{4.5cm}{0.4pt}\par\vspace{1.5pt}
{\fontsize{7.0pt}{8.3pt}\selectfont
\noindent\textbf{Loci paralleli (vv. 80–122) ---} \textbf{80} Hesychius τ 1440 τρίτη Θεσμοφορίων \espTest \textbf{93-4} Suidas τ 439 \espTest \textbf{93} Suidas κ 2025 \espTest \textbf{94} (ἡμέτερος ὁ πυραμοῦς) Suidas π 3192, cf. Tzetzes \italique{Chil.} XIII 281-288 \espTest \textbf{96} (οὑκκυκλούμενος) Zonaras 678 ἐκκυκλεῖ \espTest \textbf{100} Hesychius μ 1904 Μύρμηκος ἀτραπούς ; Photius μ 606 Μύρμηκος ἀτραπός ; Suidas μ 1445 \espTest \textbf{104-6} (εὐπίστως -- σεβίσαι) Suidas σ 183 \espTest \textbf{115-6} (Ἄρτεμιν ἀγροτέραν) \Eustathe 361.37 \espTest \textbf{120} \EtymMagnum 153.21 Ἀσιάδος κρούματα· τῆς κιθάρας. οὔτως Αριστοφάνης εἶπε, παρῳδῶν τὸ ἐξ Ἐρεχθέως Εὐριπίδου ; \SynB α 2226 $=$ Photius α 2957 $=$ Suidas 4147 $=$ \EtymGen (AB) α 1271 ed. Lasserre-Livadaras\par
}
\vspace{2.5pt}\noindent\rule{4.5cm}{0.4pt}\par\vspace{1.5pt}
{\fontsize{6.9pt}{8.2pt}\selectfont
\noindent\textbf{Suite (en commentaire) de l'apparat critique \& notes paléographiques (\texttt{f. 163v}) ---} \textbf{Première \italique{parepigraphè} (entre vv. 99 et 100) :} $\text{R}^1$ copie dans la colonne le mot scénique \lecon{μινυρισμός} (« fredonnement d'Agathon »), accompagné de la précieuse scholie $\Sigma^{\text{R}}$ : μεταξὺ τῶν δυοῖν ἀξιοῦσί τινες γράφειν 'μινυρισμός', ὡς πολλὰ τοιαῦτα παρεπιγράφεται. --- \textbf{Erreurs d'attribution de \Rdeux{} dans la monodie d'Agathon (vv. 101–129) :} Prenant le dialogue lyrique fictif d'Agathon (qui chante successivement le rôle de la coryphée troyenne et celui du chœur) pour un vrai dialogue dramatique, \Rdeux{} inscrit à tort en marge des \italique{notae personarum} faisant alterner αγαθ/ et χορ/. --- \textbf{v. 80 :} ἐπεὶ τρίτη 'στὶ R : \lecon{ἐπείπερ ἐστὶ} Nauck (\italique{Mélanges gréco-romains}, 1866, II, pp. 468–471), adopté et expliqué paléographiquement par L. Hauseux : dans l'archétype en onciale, l'abréviation \textbf{Π' = περ} dans ΕΠΕΙΠ'ΕϹΤΙ a été lue \textbf{Γ' = τρίτη} (ΕΠΕΙΓ'ΕϹΤΙ), sous l'influence d'une glose savante sur le troisième jour des Thesmophories (la \italique{Nesteia}). --- \textbf{post v. 99 :} \italique{parepigraphè} \lecon{μινυρισμός} R ($\Sigma^{\text{R}}$ μεταξὺ τῶν δυοῖν ἀξιοῦσί τινες γράφειν 'μινυρισμός', ὡς πολλὰ τοιαῦτα παρεπιγράφεται). --- \textbf{v. 100 :} μύρμ- R : \lecon{Μύρμ-} (Μύρμηκος ἀτραπούς) \lecon{Bibi} d'après Hésychius $\mu\ 1904$. --- \textbf{v. 107 :} ὅπλιζε R (métathèse de consonnes) : \lecon{ὄλβιζε} Bentley \espApp \lecon{Μοῦσα} R (μούσᾳ Bergk).\par
}
\end{minipage}
\end{adjustbox}
\end{adjustwidth}
\newpage
\enTeteD{Codex Ravennas 429 (Biblioteca Classense) --- Fac-similé du f. 163v (vv. 80–122)}
\begin{adjustwidth}{-1.5cm}{-3.0cm}
\begin{center}
\includegraphics[width=16.5cm,height=20.5cm,keepaspectratio]{/workspaces/Thesmophories/dossier_ravennas_thesmophories/folios_individuels_ff162v-177r/03_f163v_classense165v_vv0080-0122.jpg}\par
\vspace{2pt}
{\scriptsize\textbf{Codex Ravennas 429, f. 163v (Classense 165v)} --- ΠΡΟΛΟΓΟΣ, \textbf{vv. 80–122} ($\text{R}^1, \Rpc, \Rdeux$ (\Mon)).}
\end{center}
\end{adjustwidth}
\newpage
\clearpage\ifodd\value{page}\hbox{}\newpage\fi
\enTeteG{ΘΕΣΜΟΦΟΡΙΑΖΟΥΣΑΙ --- Texte grec établi (f. 163v, vv. 80–122)}
\enTeteD{LES THESMOPHORIES --- Traduction française (f. 163v, vv. 80–122)}
\begin{pages}
\begin{Leftside}
\beginnumbering
%% v. 80+ (part 2)
\begin{groupe}
\trimIamb{\ligne{80}ἐπείπερ ἐστὶ Θεσμοφορίων ἡ μέση.%\testimonia{80}{Hesychius τ 1440 τρίτη Θεσμοφορίων}
\apparat{80}{ἐπεὶ τρίτη 'στὶ R : \lecon{ἐπείπερ ἐστὶ} Nauck (\italique{Mél.} (1866), II pp.468-71)}%
\commentaire{80}{ἐπεὶ τρίτη 'στὶ Θεσμοφορίων ἡ μέση.}{Le manuscrit de Ravenne donne au vers 80 : « puisque c'est le troisième [jour] des Thesmophories, celui du milieu. » Problème : si les Thesmophories duraient bien trois jours, celui du milieu ne peut en aucune façon être le troisième. Cette contradiction eut la fâcheuse conséquence de déchaîner des torrents de commentaires (à commencer par celui du scholiaste) ; on a inventé, en contradiction avec tous les témoignages\marqueur{1}, des Thesmophories de quatre jours à grand renfort d'érudition et de savantes notes, et ce sous le prétexte qu'il ait pu y avoir avant ces trois jours d'autres festivités ailleurs (surprenant, non ? mais il suffit de prendre un calendrier des fêtes grecques pour constater qu'avec ce genre de raisonnements, on pourrait tout aussi bien \italique{prouver} que les Thesmophories duraient 365 jours par an ; au reste, il faudrait encore montrer que le milieu d'un ensemble de quatre éléments en est le \nieme{3} ; pour finir – \italique{last but not least} –, faire préciser à Aristophane quel jour tombe ἡ μέση serait un détail inepte tant cela était connu de tous les Athéniens). Le texte avait été pourtant brillamment rétabli par August Nauck\marqueur{2} et intégré avec explications dès l'édition aux Belles Lettres de Victor Coulon ; cela n'empêche pas la faute d'apparaître dans nombre d'éditions modernes. Venons-en donc à la bonne correction : l'apparition du τρίτη, \italique{troisième}, est due à la mauvaise lecture d'un Π' (abréviation de περί ou περ) en Γ' (\italique{troisième} lettre de l'alphabet) dans un ancien papyrus\marqueur{3} – cette abréviation est fréquente, et source d'erreurs aussi dans l'autre sens ; ainsi a-t-il fallu en maintes occasions rétablir un `τρίτη' à partir d'un gamma\marqueur{4}. Bref, voici l'histoire de notre début de \nieme{80} vers, très ingénieusement reconstituée : $$\overset{\text{\footnotesize{\allemand{Urtext}}}}{\text{ἐπείπερ ἐστὶ}}\xrightarrow{\text{onciales}}\text{ΕΠΕΙΠ'ΕCΤΙ}\xrightarrow{\text{corruption}}\text{ΕΠΕΙΓ'ΕCΤΙ}\xrightarrow{\text{interprétation}}\overset{\text{\footnotesize{\italique{Ravennas}}}}{\text{ἐπεὶ τρίτη ᾽στὶ.}}$$}%
\debutNote{80}%
\note{1}{Outre le scholiaste, qui \italique{sait} que les Thesmophories ne duraient que trois jours – sinon il ne prendrait pas la peine de tant ergoter -- et la glose dubitative d'Hésychius (τ 1440), citons : Alciphron 2,37.2 ; 3,39 ; et la « \anglais{speculative restauration} » du papyrus n° 3540.1-2, \italique{The Oxyrhynchus Papyri}, L (1983), ed. E. W. Handley}%
\note{2}{in \italique{Mélanges gréco-romains}, St.-Pétersburg (1866), t. II, pp. 468-71 ; voir aussi pp. 344-5 pour des témoignages sur les \italique{Thesmophories}.}%
\note{3}{La corruption doit être ancienne puisqu'elle – conjointe au bavardage du scholiaste – rendait déjà perplexe Hésychius (τ 1440): τρίτη Θεσμοφορίων· ζητεῖται πῶς ἅμα μὲν λέγει τρίτην Θεσμοφορίων εἶναι, ἅμα δὲ μέσην, τεττάρων οὐσῶν ἡμερῶν.}%
\note{4}{Pour rester avec notre cher Aristophane, un `τρεῖς' a dû être restitué au vers 1507 des \italique{Guêpes} là où tous les manuscrits comportent un `γε'.}\\\passeLigne}%
\end{groupe}%
%% v. 81-84a
\begin{groupe}%
\personnageGrec{Εὐριπίδης}%
\trimIamb{τοῦτ᾽ αὐτὸ γὰρ τοι κἀπολεῖν με προσδοκῶ.\\
αἱ γὰρ γυναῖκες ἐπιβεβουλεύκασί μοι\\
κἀν Θεσμοφόροιν μέλλουσι περί μου τήμερον\apparat{83}{ἐμοῦ R : \lecon{μου} Biset}%
\commentaire{83}{κἀν Θεσμοφόροιν.}{Le scholiaste précise que c'est une Ἀττικὴ σύνταξις καὶ δυϊκὴ, ἀντὶ τοῦ `ἐν τῷ τῶν θεσμοφόρων ἱερῷ' ; les \italique{deux} déesses thesmophores étant Déméter et sa fille Perséphone.}\\
ἐκκλησιάζειν ἐπ᾽ ὀλέθρῳ.\\\passeLigne}%
\end{groupe}%
%% v. 84b
\begin{groupe}%
\personnageGrec{Μνησίλοχος}%
\trimIamb{\ligne{84}\decalage{ἐκκλησιάζειν ἐπ᾽ ὀλέθρῳ.}τιὴ τί δή;\\\passeLigne}%
\end{groupe}%
%% v. 85
\begin{groupe}%
\personnageGrec{Εὐριπίδης}%
\trimIamb{ὁτιὴ τραγῳδῶ καὶ κακῶς αὐτὰς λέγω.\\\passeLigne}%
\end{groupe}%
%% v. 86-87
\begin{groupe}%
\personnageGrec{Μνησίλοχος}%
\trimIamb{νὴ τὸν Ποσειδῶ καὶ δικαίαν ἂν πάθοις.\apparat{86}{δί᾽ και᾽ R : δίκαιά γ᾽ Grynaeus : \lecon{δικαίαν} Wilamowitz}%
\commentaire{86-7}{νὴ τὸν Ποσειδῶ καὶ δί᾽ και᾽ ἂν πάθοις / ἀτὰρ τίν᾽ ἐκ ταύτης σὺ μηχανὴν ἔχεις;}{Ne voyant pas à quel nom féminin renvoyait `ἐκ ταύτης', on a retouché le texte du vers 87 à l'envi (ἐκ τούτων Bergler, εἰς ταύτας Sandbach, \&c.). 
C'est à mon avis aller un peu vite en besogne alors qu'il y a justement un problème textuel au vers précédent : il manque une syllabe, probablement au niveau du fautif `δί᾽ και᾽'. Je fais l'hypothèse que c'est justement à cet endroit que se trouvait ce nom féminin, repris par le `ἐκ ταύτης' du vers suivant. De quoi Euripide peut-il se mettre \italique{hors} en déployant quelque machine ? Réponse : tel le Strepsiade des \italique{Nuées}, c'est à la justice qu'il veut échapper. Au vu du texte conservé, la correction est évidente : il faut substituer au `δί᾽ και᾽' un `δικαίαν' ; la chute de l'`αν' s'expliquant par la courante erreur de «~saut du même au même~» (`δικαίαν' étant directement suivi de la particule `ἂν'). 
Lisons donc :
\citComm{νὴ τὸν Ποσειδῶ καὶ δικαίαν ἂν πάθοις.\\
ἀτὰρ τίν᾽ ἐκ ταύτης σὺ μηχανὴν ἔχεις;\\
\francais{Sainte Mer ! ce serait en justice que tu pâtirais.\\
Mais pour y échapper, de quelle machine disposes-tu ?}}
Cette correction est d'autant plus intéressante que ἡ δικαία, mot rare, est un synonyme poétique de δίκη\marqueur{1} ; Mnésiloque retourne-t-il contre Euripide une de ses propres sentences\marqueur{2} ?\\
Wilamowitz\marqueur{3} avait déjà proposé cette émendation sans qu'elle soit jamais reprise dans les éditions publiées depuis -- à une notable exception : celle de Prato.%encore une fois, on regrettera que les erreurs soient continuellement recopiées quand les meilleures restitutions sont superbement ignorées !.. manque de travail des différents éditeurs d'Aristophane ? ou manque de discernement ? Aristophane n'a pas eu la chance d'avoir beaucoup de commentateurs du calibre d'Alphonse Willems.
}%
\debutNote{86-7}%
\note{1}{Le mot est attesté dans l'\EtymMagnum à l'entrée Ἀθηνᾶ : le glosateur y recense Ἀθηναία, qui est à Ἀθήνη ce que δικαία est à δίκη (Ἀθήνη Ἀθηναία, ὡς δίκη δικαία).}%
\note{2}{Tout au long de la pièce, Euripide est pris au piège de ses propres γνῶμαι (aux vers 194, 275-276, etc). Est-ce aussi le cas ici ? Quoi qu'il en soit, on trouve des expressions très proches aux vers 1379 des \italique{Nuées} : %
\citComm{Φειδιππίδης -- νὴ τὸν Δί᾽ ἐν δίκῃ γ᾽ ἄν,}%
ou 403 de \italique{Lysistrata} :%
\citComm{Πρόβουλος -- νὴ τὸν Ποσειδῶ τὸν ἁλυκὸν δίκαιά [δικαία ?] γε.}%
}%
\note{3}{Au commentaire du vers 681 de \italique{La Folie d'Héraclès} : second volume de son \italique{Euripides Herakles erklärt von Ulrich von Wilamowitz-Moellendorff}. Wilamowitz cite \italique{en passant} le vers 86 des \italique{Thesmophories} ainsi corrigé sans même préciser la correction -- dont il semble bien être l'auteur -- tant elle devait lui sembler aller de soi.}\\
%
%\lednopb%
ἀτὰρ τίν᾽ ἐκ ταύτης σὺ μηχανὴν ἔχεις;
\commentaire{87}{μηχανὴν.}{Au sens imagé, une μηχανή est un \italique{expédient}, un \italique{tour}, un \italique{stratagème}. Certes, Euripide est un roué, un filou, un retors... mais Euripide est avant tout un homme de théâtre ! Il dispose de \italique{machineries} dont il use et abuse\marqueur{1}. Ce sont ces deux sens de μηχανή, le propre et le figuré, qui étaient instantanément saisis du spectateur du \nieme{V} siècle av. J.-C. Conformité du langage français avec le grec : le mot « truc », qui désignait au \nieme{XVII} siècle une machine de théâtre, s'est approprié le même sens imagé. % ; seul hic : le sens originel de \italique{machine théâtrale} s'est davantage estompé en français qu'en grec.
}%
\debutNote{87}%
\note{1}{Euripide recourut à la μηχανή massivement dans ses œuvres tardives. Aristote, dans sa \italique{Poétique} (XV, 1454a37), le lui reproche pour sa \italique{Médée} : Euripide privilégia un dénouement ἀπὸ μηχανῆς plutôt que par l'histoire elle-même : φανερὸν οὖν ὅτι καὶ τὰς λύσεις τῶν μύθων ἐξ αὐτοῦ δεῖ τοῦ μύθου συμβαίνειν, καὶ μὴ ὥσπερ ἐν τῇ Μηδείᾳ ἀπὸ μηχανῆς. Aristophane parodie en maints endroits cet artifice de théâtre. Pour un dossier plus complet, aller voir l'article «~Machines et \italique{deus ex machina} dans les spectacles comiques~» de Leonardo Fiorentini, in \italique{L'appareil scénique dans les spectacles de l'Antiquité} (2013).}\\\passeLigne}%
\end{groupe}%
%% v. 88-89a
\begin{groupe}%
\personnageGrec{Εὐριπίδης}%
\trimIamb{Ἀγάθωνα πεῖσαι τὸν τραγῳδοδιδάσκαλον\\
ἐς Θεσμοφόροιν ἐλθεῖν.\\\passeLigne}%
\end{groupe}%
%% v. 89b
\begin{groupe}%
\personnageGrec{Μνησίλοχος}%
\trimIamb{\ligne{89}\decalage{ἐς Θεσμοφόροιν ἐλθεῖν.}τί δράσοντ᾽; εἰπέ μοι.\\\passeLigne}%
\end{groupe}%
%% v. 90-91a
\begin{groupe}%
\personnageGrec{Εὐριπίδης}%
\trimIamb{ἐκκλησιάσοντ᾽ ἐν ταῖς γυναιξὶ χἂν δέῃ\apparat{90}{γυναιξίν R : \lecon{γυναιξὶ} \Mon \espApp κἂν R : \lecon{χἂν} Markland (1763) ad Eur. \italique{Suppl.} 364}\\
λέξονθ᾽ ὑπὲρ ἐμοῦ.
\commentaire{90-1}{ἐκκλησιάσοντ᾽ ἐν ταῖς γυναιξὶν κἂν δέῃ / λέξονθ᾽ ὑπὲρ ἐμοῦ.}{Je reprends la correction de Markland : substituer χἂν ($=$ καὶ ἃ ἂν) à κἂν ($=$ καὶ ἐὰν) qui présente un sens bien meilleur : Euripide ne dit certes pas « qu'au besoin » ou « s'il le faut » (κἂν) Agathon interviendra mais bien plutôt qu'il lui faut intervenir autant qu'il le pourra (χἂν). Euripide a sa vie en jeu – si l'on peut s'exprimer ainsi – et sait bien le courroux de la gent féminine ; l'inaction aurait forcément une issue fatale... heureuse pour Athènes et les Athéniennes mais un peu plus fâcheuse pour sa personne. 
Précisons que ἐν a ici le sens technique que donne le Bailly (A, III.10) : « lorsqu'il s'agit d'un orateur, d'un accusé, ἐν correspond au français \italique{devant} » (cf. les \italique{Acharniens} v. 630 ; les \italique{Guêpes} v. 1185 ; le fragment 161 K.-A.) ; il faut donner un sens actif au `ἐκκλησιάσοντ᾽ ἐν ταῖς γυναιξὶν' : « pour intervenir à l'assemblée devant les femmes ».}\\\passeLigne}%
\end{groupe}%
%% v. 91b
\begin{groupe}%
\personnageGrec{Μνησίλοχος}%
\trimIamb{\ligne{91}\decalage{λέξονθ᾽ ὑπὲρ ἐμοῦ.}πότερα φανερὸν ἢ λάθρᾳ;\apparat{91}{\lecon{φανερὸν} R : φανερῶς Meineke}\\\passeLigne}%
\end{groupe}%
%% v. 92
\begin{groupe}%
\personnageGrec{Εὐριπίδης}%
\trimIamb{λάθρᾳ, στολὴν γυναικὸς ἠμφιεσμένον.\\\passeLigne}%
\end{groupe}%
%% v. 93-94
\begin{groupe}%
\personnageGrec{Μνησίλοχος}%
\trimIamb{τὸ πρᾶγμα κομψὸν καὶ σφόδρ᾽ ἐκ τοῦ σοῦ τρόπου·%\testimonia{93-4}{Suidas τ 439}\testimonia{93}{Suidas κ 2025}
\\
%\lednopb%
τοῦ γὰρ τεχνάζειν ἡμέτερος ὁ πυραμοῦς.%\testimonia{94}{(ἡμέτερος ὁ πυραμοῦς) Suidas π 3192, cf. Tzetzes \italique{Chil.} XIII 281-288}%
\commentaire{94}{ἡμέτερος ὁ πυραμοῦς.}{L'expression proverbiale signifie : « à nous la palme de la victoire ». En effet, on offrait un gâteau de miel et de froment (πυραμοῦς), à titre de récompense, à celui qui avait le mieux résisté au sommeil dans une παννυχίς (scholiaste des \italique{Cavaliers}, au v. 277 où l'expression apparaît à l'identique ; Athénée 647b-c, 668c). (Note de Taillardat, §847.)}\\\passeLigne}%
\end{groupe}%
%% v. 95a
\begin{groupe}%
\personnageGrec{Εὐριπίδης}%
\trimIamb{σίγα.\\\passeLigne}%
\end{groupe}%
%% v. 95b
\begin{groupe}%
\personnageGrec{Μνησίλοχος}%
\trimIamb{\ligne{95}\decalage{σίγα.}τί δ᾽ ἔστιν;\apparat{95}{\lecon{τί δ᾽} R : τί Kaibel}\\\passeLigne}%
\end{groupe}%
%% v. 95c
\begin{groupe}%
\personnageGrec{Εὐριπίδης}%
\trimIamb{\ligne{95}\decalage{σίγα.τί δ᾽ ἔστιν;}Ἁγάθων ἐξέρχεται.\\\passeLigne}%
\end{groupe}%
%% v. 96a
\begin{groupe}%
\personnageGrec{Μνησίλοχος}%
\trimIamb{καὶ ποῖός ἐστιν;\apparat{96}{ποῖο ἐστὶν οὗτος — οὑκκυκλούμενος R : \lecon{ποῖος ἐστὶν; — οὗτος οὑκκυκλούμενος} Tyrwhitt (p. 65)}%\apparat{96}{ποῖο ἐστὶν οὗτος — οὑκκυκλούμενος R : ποῖος ἐστὶν οὗτος — οὑκκυκλούμενος \Mon : ποῦ ᾽σθ᾽; — ὅπου ᾽στίν; οὗτος οὑκκυκλούμενος Meineke : \lecon{ποῖος ἐστὶν; — οὗτος οὑκκυκλούμενος} Tyrwhitt (p. 65)}
\\\passeLigne}%
\end{groupe}%
%% v. 96b
\begin{groupe}%
\personnageGrec{Εὐριπίδης}%
\trimIamb{\ligne{96}\decalage{καὶ ποῖός ἐστιν;}οὗτος οὑκκυκλούμενος.%\testimonia{96}{(οὑκκυκλούμενος) Zonaras 678 ἐκκυκλεῖ}%
\commentaire{96}{οὑκκυκλούμενος}{($=$ ὁ ἐκκυκλούμενος). Agathon apparaît sur l'eccyclème, « machine de théâtre munie de roues qui, en tournant, faisait voir aux spectateurs ce qui était censé se passer dans une maison. » (Chantraine.) 
Le scholiaste l'avait noté :%
\citComm{ἐπὶ ἐκκυκλήματος γὰρ φαίνεται.}%
%\eff{Houston ! we have a textual problem !}
}\\\passeLigne}%
\end{groupe}%
%% v. 97-98
\begin{groupe}%
\personnageGrec{Μνησίλοχος}%
\trimIamb{ἀλλ᾽ ἦ τυφλὸς μέν εἰμ᾽· ἐγὼ γὰρ οὐχ ὁρῶ\apparat{97}{ἢ R : \lecon{ἦ} Biset}\\
ἄνδρ᾽ οὐδέν᾽ ἐνθάδ᾽ ὄντα, Κυρήνην δ᾽ ὁρῶ.
\commentaire{98}{Κυρήνην δ᾽ ὁρῶ.}{Cyrène était une célèbre courtisane évoquée aussi aux vers 1327-1328 des \italique{Grenouilles} : 
\citComm{ἀνὰ τὸ δωδεκαμήχανον / Κυρήνης μελοποιῶν\\
\francais{$[$Euripide$]$ compose ses chants à la mode de Cyrène aux douze postures.}}%
Le scholiaste y est plus loquace que le nôtre et nous apprend le trait distinctif de cette Cyrène :%
\citComm{Κυρήνη τις ἑταίρα ἐπίσημος, δωδεκαμήχανος ἐπικαλουμένη διὰ τὸ τοσαῦτα σχήματα ἐν τῇ συνουσίᾳ ποιεῖν.}%
Taillardat (§734) rapproche pertinemment ce sobriquet du titre d'un ``traité'' de Paxamos : Δωδεκάτεχνον, à propos duquel la \italique{Souda} (à l'entrée Πάξαμος) dit : ἔστι δὲ περὶ αἰσχρῶν σχημάτων. Que moque donc Aristophane en rapprochant, non cette fois la muse du poète comme dans les \italique{Grenouilles}, mais Agathon lui-même de Cyrène ? Bien-sûr son air efféminé ! bien-sûr son travestissement !.. mais pourquoi Cyrène en particulier ? Réponse : pour ce qu'elle a de remarquable (ἐπίσημος) qui lui valut son surnom ; il ne me paraît pas douteux qu'Agathon, qui vient de se faire rouler par l'eccyclème, apparaît affalé sur son lit dans une position suggestive, ou du moins qui peut prêter à la confusion de Mnésiloque. (Cf. la raillerie -- μῶν βινεῖσθαι; -- au vers 50.)}\\\passeLigne}%
\end{groupe}%
%% v. 99
\begin{groupe}%
\personnageGrec{Εὐριπίδης}%
\trimIamb{σίγα· μελῳδεῖν νῦν παρασκευάζεται.\apparat{99}{ἂν R : \lecon{νῦν} Rogers : alii alia}%
\commentaire{99}{μινυρισμός.}{Une scholie du \italique{Ravennas} semble avoir gardé la trace d'une ancienne didascalie : μινυρισμός était indiqué marginalement :
\citComm{μεταξὺ τῶν δυοῖν ἀξιοῦσί τινες γράφειν `μινυρισμός', ὡς πολλὰ τοιαῦτα παρεπιγράφεται.}%
}\\\passeLigne}%
\end{groupe}%
%% v. 100
\begin{groupe}%
\personnageGrec{Μνησίλοχος}%
\trimIamb{Μύρμηκος ἀτραπούς, ἢ τί διαμινύρεται;%\testimonia{100}{Hesychius μ 1904 Μύρμηκος ἀτραπούς ; Photius μ 606 Μύρμηκος ἀτραπός ; Suidas μ 1445 }
\apparat{100}{μύρμ- R : \lecon{Μύρμ-} Bibi ap Hesych. μ 1904}%
\commentaire{100}{μύρμηκος ἀτραπούς.}{Willems, en accord avec le scholiaste, se propose d'expliquer ainsi l'expression μύρμηκος ἀτραπούς : « par allusion à la marche saccadée et tortueuse des fourmis. Image familière aux Anciens. Dans l'\italique{Anthologie}, XI, 78, il est question des galeries obliques et droites, λόξα καὶ ὀρθά, d'une fourmilière. Voir aussi Plutarque, \italique{de Solertia animal.}, II, p. 968a. Plaute a dit : \latin{formicinum gradum movere}. » Taillardat (§784) abonde en son sens avec d'autres exemples à l'appui. %Rapprocher les airs contournés d'Agathon de sentiers tortus était certainement le but d'Aristophane.
Cependant, pour suggérer l'image de \italique{sentiers de fourmis}, n'aurait-il pas fallu, en français comme en grec, un pluriel (μυρμήκ\textbf{ων} ἀτραπούς) ?\\
Hésychius ne le comprenait pas autrement et voyait dans `μύρμηκος' un nom propre (μ 1904) :
\citComm{Μύρμηκος ἀτραπούς· Ἀθήνησιν ἐν Σκαμβωνιδῶν ἔστι Μύρμηκος ἀτραπός, ἀπὸ ἥρωος Μύρμηκος ὀνομαζομένη.\\
\francais{Dans [le dème] athénien des Scambônides [σκαμβός signifie \italique{courbe}, \italique{tortu}] se trouve un sentier de Myrmex, qui tire son nom du héros Myrmex.}}
Que l'image soit tirée de la fourmilière ou d'un nom propre, la critique n'en reste pas moins d'ordre musical\marqueur{1}.
%Il y a donc plus vraisemblablement une évocation d'un véritable lieu, connu des Athéniens, plutôt qu'une image animalière.\\
%Voici pour le sens littéral. Venons-en au côté facétieux du commentaire de Mnésiloque !\\
% L'image des lignes brisées de la fourmilière, μυρμηκιά, servit en musique pour décrire des trilles ou des arpèges comme l'atteste un fragment de Phérécrate\marqueur{2}.\\
Philoxène de Cythère, un représentant de cette «~nouvelle musique~» dénoncée par Aristophane, reçut d'ailleurs le surnom de Μύρμηξ (voir \italique{Souda} φ 393, à l'entrée Φιλόξενος). Le persiflage de Mnésiloque ne serait-il pas une allusion audit Philoxène, alias Myrmex ?
%\citComm{Μύρμηκος ἀτραπούς\\
%C'est la voie de Myrmex}
}%
\debutNote{100}%
%\note{1}{Relevé par nombre de parémiographes ainsi que Suidas à l'entrée Ἀγαθῶν ἀγαθίδες.}
\note{1}{L'image des lignes brisées de la fourmilière, μυρμηκιά, servit en musique pour décrire des trilles ou des arpèges. On trouve l'expression `ἐκτραπέλους μυρμηκιάς' chez le Pseudo-Plutarque, \italique{De Musica}, 30 ; il s'agit du fragment 155(v.23) de Phérécrate dans l'édition Kassel-Austin. À noter que le même Phérécrate a composé une comédie intitulée les μυρμηκάνθρωποι.}%
%\note{2}{Voir l'article de la \italique{Souda} (φ 393) qui lui est consacré, à l'entrée Φιλόξενος.}%
\\%
\commentaire{101}{Χθονίαιν}{ est un génitif se rapportant à ἱεράν. Cf. \italique{Ploutos}, 937, et \italique{Caval.}, 301. (Note de Willems.)}\passeLigne}%
\end{groupe}%
%% v. 101-129
\begin{groupe}%
\personnageGrec{Ἀγάθων}%
\typeVers{io. min\cat + io. min.}\chantAgath{ἱερὰν Χθονίαιν δε-\apparat{101}{χθονίαις R : \lecon{χθονίαιν} Meineke}%
\commentaire{101-29}{Agathon chante un fragment de tragédie.}{Sujet : un coryphée et un chœur de jeunes filles troyennes célèbrent en strophes alternées Léto et ses deux rejetons. Il est très évident qu'Agathon chante seul. Les arguments fournis par M. P. Mazon\marqueur{1} sont à cet égard sans réplique. (Note de Willems.)\\
Suidas (χ 403) nous confirme qu'Agathon avait l'habitude de chanter de la sorte : 
\citComm{Χορικός· ὁ τοῦ χοροῦ ἐξάρχων. τοιοῦτος ἦν Ἀγάθων, ὃς ὑποκριτικὰ μέλη ἐποίει καὶ αὐτὸς ὑποκρίνεται\\
\francais{Chorique : celui qui mène le chœur. Tel Agathon, qui avait l'habitude de créer des chants dramatiques et de les interpréter lui-même.}}%
L'on pas manqué de remarquer en quoi les ruptures de rythmes (παράρυθμ' εὔρυθμα) ou le contexte ionique de cette parodie d'Agathon en font un morceau efféminé et suave, ce qui est tout-à-fait juste quoiqu'il nous manque la musique pour juger sur pièces\marqueur{2} ; mais beaucoup trop de commentateurs se sont arrêté à ces considérations\marqueur{3}. 
(Opinion au demeurant est bien pratique puisqu'elle dispense d'avoir à porter trop d'attention sur le texte.) Ce fut pas le cas du regretté Alain Blanchard avait excellemment analysé la parodie d'Agathon dans son article : \italique{La parodie d'Agathon et la structure des }``Bucoliques''\italique{ de Virgile}\marqueur{4}. Blanchard décela dans ce chant une structure concentrique, en cela analogue aux \italique{Bucoliques} de Virgile ou encore aux cantates de Bach ; le point important étant que « la reconnaissance de ce schéma concentrique » est d'une grande « valeur exégétique ». Laissons-lui la parole : « la coryphée invite le chœur à célébrer Apollon, Artémis et leur mère Léto, et le chœur, docilement, obéit à cette injonction. En réalité, l'étude de composition, une fois admise l'existence d'une structure concentrique conduit à plus de subtilité.
Le parallélisme entre le premier élément (vv. 101-103) et le neuvième (vv. 126-129) réside dans le thème de la lumière (lumière des torches portées par le chœur, lumière qui brille en réponse dans les yeux divins), thème qui précède un appel à la célébration d'un dieu – l'appel étant plus solennel, plus brusque et plutôt à titre de rappel dans le neuvième. Le parallélisme entre le deuxième élément (104-106) et le huitième (123-125) réside dans le thème du σέβεσθαι, introduit en mineur à la fin du deuxième élément et qui éclate sur le mode majeur au début du huitième.
Le parallélisme entre le troisième élément (107-110) et le septième (120-122) réside dans le caractère national des divinités honorées : caractère national sensible surtout dans le troisième élément (Apollon constructeur des remparts de Troie), mais présent, plus discrètement, dans le septième avec les Charites phrygiennes.
Enfin, le thème de la majesté divine (majesté d'Apollon comme dispensateur de la gloire) entre le quatrième élément (111-113) et le sixième (117-119). Le cinquième (114-116), central, est un appel à honorer Artémis, divinité féminine. Dans la première partie de la parodie, les trois interventions du coryphée concernent des divinités armées de l'arc qu'il faut honorer avec un esprit libre et d'une voix forte (103), mais le coryphée s'adresse à un chœur de jeunes filles qui, pour ainsi dire, retourne complètement la nature de ces divinités dont on aperçoit alors le caractère ambigu. Déjà, aux appels à un culte viril qu'on leur faisait, elles avaient répondu, avec une soumission apparente bien féminine, par une profession de piété plus neutre (σεβίσαι, 106). Puis, avec elles, Apollon cesse d'être le dieu archer pour être le coryphée des Muses (112). Qu'on leur propose d'honorer Artémis, il n'est plus question avec elle de chasse – occupation virile – mais de maternité : maternité glorieuse de Léto (118), maternité refusée par elle-même par la déesse Vierge (119). Le ton est donné pour la seconde partie du poème où le chœur, à l'appel du coryphée, honorera, après Artémis, sa mère Léto et la cithare. Dans un cadre aussi favorable, le chœur peut même se permettre une concession – nous parlerions plutôt d'un appel spontané et aventureux du cœur – avec l'expression ἄρσενι βοᾷ, \italique{mâle accent}, en 125 qui rappelle en la précisant l'invitation faite en 103b par le coryphée. Mais qu'on invite de nouveau à honorer une divinité masculine comme Apollon : c'est la maternité de Léto qui est chantée (130). »\\
Fera-t-on la même erreur à propos de ce chant que les premiers auditeurs du sonnet d'Oronte en prétendant : « Il est bon » ? C'est en tout cas ce qu'il nous semble à l'analyse de ce petit bijou de malice au subtil renversement : d'un banal hommage au fondateur d'une cité, l'on passe doucement et insidieusement à un morceau d'érotisme sur fond de maternité frustrée, le tout échafaudé sur une structure solidement ancrée\marqueur{5}. Rien d'étonnant donc à ce que Mnésiloque en soit tout émoustillé ; rien d'étonnant non plus à ce que, par la bouche de Dionysos, Aristophane rende hommage à Agathon après sa mort : ἀγαθὸς ποιητὴς καὶ ποθεινὸς τοῖς φίλοις, « un bon poète apprécié de ses amis » (\italique{Grenouilles}, v. 84). Pour finir, l'élaboration complexe du morceau ainsi mise au jour par Blanchard fait voler en éclat la thèse fort répandue d'une improvisation de la part d'Agathon ; d'ailleurs l'on sait par Aristote qu'Agathon avait eu l'ambition dans une de ses tétralogies de traiter intégralement du sac de Troie, ambition qui l'avait mené à l'échec : (\italique{Poétique}, XVIII, 56a16-19) :
\citComm{ὅσοι πέρσιν Ἰλίου ὅλην ἐποίησαν καὶ μὴ κατὰ μέρος ὥσπερ Εὐριπίδης, <ἢ> Νιόβην καὶ μὴ ὥσπερ Αἰσχύλος, ἢ ἐκπίπτουσιν ἢ κακῶς ἀγωνίζονται, ἐπεὶ καὶ Ἀγάθων ἐξέπεσεν ἐν τούτῳ μόνῳ.\\
\francais{Tous les poètes qui veulent traiter la ruine de Troie entièrement et non pas par parties comme Euripide, ou la légende de Niobé tout entière au lieu de faire comme Eschyle, tantôt échouent complètement, tantôt font mauvaise figure dans les concours. Ainsi, Agathon échoua sur ce seul point.}}%
Je me plais à imaginer Aristophane, révolté devant cet échec injuste, remettre sur scène ce qu'il considérait comme la perle (érotique) du cycle troyen d'Agathon (c'est en tout cas ce sens de la justice qui l'avait poussé à vouloir faire rejouer ses \italique{Nuées}, mésestimées par les mauvais jurés athéniens).}%
\debutNote{101-29}%
\note{1}{in \italique{Essai sur la composition des comédies d'Aristophane} (1904), p. 127-128.}%
\note{2}{À ce propos, cf. les \italique{Grenouilles} v. 1299 où Eschyle dit s'être écarté du pré des Muses de Phrynichos – Muses par trop efféminées et ionisantes ; de même, l'\italique{Assemblée des femmes} (v. 883).%, quand une vieille veut alpaguer les beaux jeunes hommes qu'elle est en droit de s'approprier par force de loi, elle se lance alors dans une ariette au goût ionien (v. 883).
}%
\note{3}{%L'internationale du commentaire s'est donné le mot : « \allemand{Vor allem aber wird die Musik selbst durch Karikatur verspottet} » (Rau, p. 106) ; « \italien{più agli aspetti ritmico-musicali dell'inno che a quelli stilistici} » (Prato, p. 177) ; « \anglais{are due to the metre and the music rather than the words} » (Austin \& Olson p. 87) ; légèrement plus nuancé chez Rosella Saetta Cottone : «~une partie essentielle de la parodie devait reposer sur les modulations vocales par lesquelles l'acteur imitait la voix féminine~» ; \&c... %en dépouillant minutieusement tous les glosateurs d'Aristophane, on devrait pouvoir traduire cette phrase dans une bonne partie des idiomes de notre babélienne Terre. Cette opinion au demeurant est bien pratique puisqu'elle dispense d'avoir à porter trop d'attention sur le texte.
Le style d'Agathon est certes efféminé mais beaucoup moins caricatural qu'on a voulu le faire croire. Cf. le fragment 178 (Γηρυτάδης) de l'édition Kassel-Austin : Aristophane y emploie l'adjectif Ἀγαθώνειος dans un sens loin d'être péjoratif d'après les commentaires anciens (Zénobius I, 2 ; Hésych. α 281 ; Phot. α 83 ; Suid. α 125 ; \&c.) :%
\citComm{Ἀγαθώνιος αὔλησις: ἡ μαλακὴ καὶ ἐκλελυμένη: ἢ ἡ μήτε χαλαρὰ, μήτε πικρὰ, ἀλλ' εὔκρατος καὶ ἡδίστη.\\
Un air de flûte agathonien : doux et détendu ; ni mol ni strident, mais tempéré et très doux.}%
}%
\note{4}{\italique{Dans l'ouvroir du poète}, \italique{Hellenica} (2008), pp. 35-45. J'eus la chance de recevoir d'Alain Blanchard peu de temps avant sa mort de vifs encouragements, source d'énergie nécessaire à l'abattage de la somme de travail que représente ce commentaire des \italique{Thesmophories}. %Quel regret pour moi que ce grand helléniste n'ait pas consacré son temps de recherche à Aristophane plutôt qu'à son «~cher Ménandre~».
}%
\note{5}{Cf. pour plus de détails l'article d'A. Blanchard avec les chiffres à la clef.}%
}%
\end{groupe}%
%% v. 101-129
\begin{groupe}%
\typeVers{2 ion. min.}\chantAgath{ξάμεναι λαμπάδα κοῦραι}%
\end{groupe}%
%% v. 101-129
\begin{groupe}%
\decoupeVers{a}\typeVers{io. min\cat + bacch.}\chantAgath{ξὺν ἐλευθέρᾳ πατρί-\apparat{103}{\lecon{πατρίδι} R : πραπίδι Wecklein}}%
\end{groupe}%
%% v. 101-129
\begin{groupe}%
\decoupeVers{b}\typeVers{io. min. + io. min\cat}\chantAgath{\ligne{102}\alinea δι χορεύσασθε βοάν.\\}%
\end{groupe}%
%% v. 101-129
\begin{groupe}%
\typeVers{2 io. min. anacl.}\chantAgath{\ligne{104}τίνι δαιμόνων ὁ κῶμος;}%
\end{groupe}%
%% v. 101-129
\begin{groupe}%
\typeVers{2 troch.}\chantAgath{λέγε νιν εὐπείστως δὲ τοὐμὸν%\testimonia{104-6}{(εὐπίστως -- σεβίσαι) Suidas σ 183}
\apparat{105}{νυν R : \lecon{νιν} Dindorf \espApp εὐπίστως R, Suid. σ 178, 183 : \lecon{εὐπείστως} Reiske \ligneApp{106} ἔχεις R : \lecon{ἔχει} Suid. \ligneApp{107} ὅπλιζε R : \lecon{ὄλβιζε} Bentley \espApp \lecon{Μοῦσα} R : μούσᾳ Bergk}}%
\end{groupe}%
%% v. 101-129
\begin{groupe}%
\typeVers{troch. + io. min\cat}\chantAgath{δαίμονας ἔχει σεβίσαι.\\}%
\end{groupe}%
%% v. 101-129
\begin{groupe}%
\typeVers{ion. min. + troch.}\chantAgath{\ligne{107}ἄγε νῦν ὄλβιζε Μοῦσα}%
\end{groupe}%
%% v. 101-129
\begin{groupe}%
\typeVers{troch. + io. min.}\chantAgath{χρυσέων ῥύτορα τόξων}%
\end{groupe}%
%% v. 101-129
\begin{groupe}%
\typeVers{2 io. min.}\chantAgath{Φοῖβον, ὃς ἱδρύσατο χώρας}%
\end{groupe}%
%% v. 101-129
\begin{groupe}%
\typeVers{io. min. + io. min\cat}\chantAgath{γύαλα Σιμουντίδι γᾷ.\\}%
\end{groupe}%
%% v. 101-129
\begin{groupe}%
\typeVers{2 troch.}\chantAgath{χαῖρε καλλίσταις ἀοιδαῖς\\
Φοῖβ᾽ ἐν εὐμούσοισι τιμαῖς}%
\end{groupe}%
%% v. 101-129
\begin{groupe}%
\typeVers{io. min. + io. min\cat}\chantAgath{\ligne{112}γέρας ἱερὸν προφέρων.\\}%
\end{groupe}%
%% v. 101-129
\begin{groupe}%
\typeVers{2 troch.}\chantAgath{τάν τ᾽ ἐν ὄρεσι δρυογόνοισι}%
\end{groupe}%
%% v. 101-129
\begin{groupe}%
\typeVers{2 io. m. + io. m\cat}\chantAgath{κόραν ἀείσατ᾽ Ἄρτεμιν ἀγροτέραν.%\testimonia{115-6}{(Ἄρτεμιν ἀγροτέραν) \Eustathe 361.37}
\apparat{115}{ἀεισαντ᾽ R : \lecon{ἀείσατ᾽} Zanetti (\italique{laudate} Divus) \ligneApp{117} σεμνὸν R : \lecon{σεμνὰν} Dindorf}\\}%
\end{groupe}%
%% v. 101-129
\begin{groupe}%
\typeVers{io. min. + troch.}\chantAgath{\ligne{116}ἕπομαι κλῄζουσα σεμνὰν}%
\end{groupe}%
%% v. 101-129
\begin{groupe}%
\typeVers{io. min. + troch.}\chantAgath{γόνον ὀλβίζουσα Λατοῦς}%
\end{groupe}%
%% v. 101-129
\begin{groupe}%
\typeVers{io. min. + io. min\cat}\chantAgath{\ligne{118}Ἄρτεμιν ἀπειρολεχῆ.\\}%
\end{groupe}%
%% v. 101-129
\begin{groupe}%
\typeVers{io. m+io. m\cat+io. m}\chantAgath{Λατώ τε κρούματά τ' Ἀσιάδος ποδὶ%\testimonia{120}{\EtymMagnum 153.21 Ἀσιάδος κρούματα· τῆς κιθάρας. οὔτως Αριστοφάνης εἶπε, παρῳδῶν τὸ ἐξ Ἐρεχθέως Εὐριπίδου ; \SynB α 2226 $=$ Photius α 2957 $=$ Suidas 4147 $=$ \EtymGen (AB) α 1271 ed. Lasserre-Livadaras}
}%
\end{groupe}%
%% v. 101-129
\begin{groupe}%
\typeVers{2 io. min.}\chantAgath{παράρυθμ' εὔρυθμα Φρυγίων}%
\end{groupe}%
%% v. 101-129
\begin{groupe}%
\typeVers{io. min. + io. min\cat}\chantAgath{διὰ νεύματα Χαρίτων.\apparat{122}{διανεύματα R : δινεύματα Bentley : \lecon{διὰ νεύματα} Fritzsche}\\}%
\end{groupe}%
\endnumbering
\end{Leftside}
\begin{Rightside}
\beginnumbering
%% v. 80+ (part 2)
\begin{groupe}
\trimIambFr{\ligne{80}puisque nous sommes au jour intermédiaire des Thesmophories.\testimonia{80}{Hesychius τ 1440 τρίτη Θεσμοφορίων}%\commentaire{80}{}{Les difficultés que soulève ce passage me semblent avoir été résolues avec autant de clarté que de savoir par Guignaut, dans une note sur CREUZER, Religions de l'Antiquité, t. III, 3, 1150 et suiv.}%
\\\passeLigne}%
\end{groupe}%
%% v. 81-84a
\begin{groupe}%
\personnageFr{Euripide}%
\trimIambFr{C'est précisément ce qui fait que je m'attends à périr.\\
Car les femmes ont monté un coup contre moi.
Elles doivent, dans le temple des Thesmophories, à mon sujet aujourd'hui,
tenir assemblée afin de me perdre.\\\passeLigne}%
\end{groupe}%
%% v. 84b
\begin{groupe}%
\personnageFr{Mnésiloque}%
\trimIambFr{\ligne{84}\espace{4cm}Pourquoi cela ?\\\passeLigne}%
\end{groupe}%
%% v. 85
\begin{groupe}%
\personnageFr{Euripide}%
\trimIambFr{Parce que je les théâtralise et parle mal d'elles.\\\passeLigne}%
\end{groupe}%
%% v. 86-87
\begin{groupe}%
\personnageFr{Mnésiloque}%
\trimIambFr{Par Poseidôn, tu serais châtié en justice.\\
%
Mais, pour y échapper, de quelle machine disposes-tu ?\\\passeLigne}%
\end{groupe}%
%% v. 88-89a
\begin{groupe}%
\personnageFr{Euripide}%
\trimIambFr{De persuader à Agathon, le poète tragique,\\
de se rendre aux Thesmophories.\\\passeLigne}%
\end{groupe}%
%% v. 89b
\begin{groupe}%
\personnageFr{Mnésiloque}%
\trimIambFr{\ligne{89}\espace{3.5cm}Pour quoi faire, dis-moi ?\\\passeLigne}%
\end{groupe}%
%% v. 90-91a
\begin{groupe}%
\personnageFr{Euripide}%
\trimIambFr{Pour intervenir à l'assemblée devant les femmes et, autant qu'il est nécessaire,
me défendre.\\\passeLigne}%
\end{groupe}%
%% v. 91b
\begin{groupe}%
\personnageFr{Mnésiloque}%
\trimIambFr{\ligne{91}\espace{2.8cm}Ouvertement ou en tapinois ?\\\passeLigne}%
\end{groupe}%
%% v. 92
\begin{groupe}%
\personnageFr{Euripide}%
\trimIambFr{En tapinois, revêtu d'un costume féminin.\\\passeLigne}%
\end{groupe}%
%% v. 93-94
\begin{groupe}%
\personnageFr{Mnésiloque}%
\trimIambFr{C'est ingénieux et bien dans tes façons.\testimonia{93-4}{Suidas τ 439}\testimonia{93}{Suidas κ 2025}\\
En fait d'artifices, à nous le gâteau.\testimonia{94}{(ἡμέτερος ὁ πυραμοῦς) Suidas π 3192, cf. Tzetzes \italique{Chil.} XIII 281-288}\\\passeLigne}%
\end{groupe}%
%% v. 95a
\begin{groupe}%
\personnageFr{Euripide}%
\trimIambFr{Silence.\\\passeLigne}%
\end{groupe}%
%% v. 95b
\begin{groupe}%
\personnageFr{Mnésiloque}%
\trimIambFr{\ligne{95}\espace{1cm}Qu'y a-t-il ?\\\passeLigne}%
\end{groupe}%
%% v. 95c
\begin{groupe}%
\personnageFr{Euripide}%
\trimIambFr{\ligne{95}\espace{2.7cm}Agathon sort.\\\passeLigne}%
\end{groupe}%
%% v. 96a
\begin{groupe}%
\personnageFr{Mnésiloque}%
\trimIambFr{Quel homme est-ce ?\\\passeLigne}%
\end{groupe}%
%% v. 96b
\begin{groupe}%
\personnageFr{Euripide}%
\trimIambFr{\ligne{96}\espace{2cm}Celui-là qu'on roule dehors.\testimonia{96}{(οὑκκυκλούμενος) Zonaras 678 ἐκκυκλεῖ}\\\passeLigne}%
\end{groupe}%
%% v. 97-98
\begin{groupe}%
\personnageFr{Mnésiloque}%
\trimIambFr{Est-ce donc que je suis aveugle ? Je ne vois pas qu'il y ait
un homme ici. C'est Cyrène que je vois.\\\passeLigne}%
\end{groupe}%
%% v. 99
\begin{groupe}%
\personnageFr{Euripide}%
\trimIambFr{Chut ! Il se dispose à chanter.\\\passeLigne}%
\end{groupe}%
%% v. 100
\begin{groupe}%
\personnageFr{Mnésiloque}%
\trimIambFr{Des sentes de Fourmi, ou que sont ces fredons ?\testimonia{100}{Hesychius μ 1904 Μύρμηκος ἀτραπούς ; Photius μ 606 Μύρμηκος ἀτραπός ; Suidas μ 1445 }\\\passeLigne}%
\end{groupe}%
%% v. 101-129
\begin{groupe}%
\personnageFr{Agathon}%
\chantAgathFr{1° — Des deux Chtoniennes, recevez}%
\end{groupe}%
%% v. 101-129
\begin{groupe}%
\chantAgathFr{la torche consacrée, vierges,}%
\end{groupe}%
%% v. 101-129
\begin{groupe}%
\chantAgathFr{\ligne{101}et avec votre libre patrie}%
\end{groupe}%
%% v. 101-129
\begin{groupe}%
\chantAgathFr{entonnez un chœur de danse.\\}%
\end{groupe}%
%% v. 101-129
\begin{groupe}%
\chantAgathFr{\ligne{104}2° — À quelle divinité ce cômos ?\testimonia{104-6}{(εὐπίστως -- σεβίσαι) Suidas σ 183}}%
\end{groupe}%
%% v. 101-129
\begin{groupe}%
\chantAgathFr{Parle. Je me laisse aisément induire}%
\end{groupe}%
%% v. 101-129
\begin{groupe}%
\chantAgathFr{à honorer les dieux.\\}%
\end{groupe}%
%% v. 101-129
\begin{groupe}%
\chantAgathFr{3° — Allons, Muse, glorifie}%
\end{groupe}%
%% v. 101-129
\begin{groupe}%
\chantAgathFr{l'archer à l'arc d'or,}%
\end{groupe}%
%% v. 101-129
\begin{groupe}%
\chantAgathFr{\ligne{109}Phébus, qui fonda les murs}%
\end{groupe}%
%% v. 101-129
\begin{groupe}%
\chantAgathFr{de cette cité dans la contrée du Simois.\\}%
\end{groupe}%
%% v. 101-129
\begin{groupe}%
\chantAgathFr{4° — Sois salué de nos plus beaux chants,\\
Phébus, qui dans les fêtes des Muses}%
\end{groupe}%
%% v. 101-129
\begin{groupe}%
\chantAgathFr{\ligne{112}emportes le prix sacré.\\}%
\end{groupe}%
%% v. 101-129
\begin{groupe}%
\chantAgathFr{\ligne{114}5° — Et chantez la Vierge qui se plaît dans les}%
\end{groupe}%
%% v. 101-129
\begin{groupe}%
\chantAgathFr{montagnes dryogènes, Artémis-chasseresse.\\}%
\end{groupe}%
%% v. 101-129
\begin{groupe}%
\chantAgathFr{6° — À ta suite je célèbre et glorifie }%
\end{groupe}%
%% v. 101-129
\begin{groupe}%
\chantAgathFr{l'auguste rejeton de Léto,}%
\end{groupe}%
%% v. 101-129
\begin{groupe}%
\chantAgathFr{\ligne{118}Artémis ignorant l'hymen.\testimonia{115-6}{(Ἄρτεμιν ἀγροτέραν) \Eustathe 361.37}\\}%
\end{groupe}%
%% v. 101-129
\begin{groupe}%
\chantAgathFr{7° — Ainsi que Léto et les percussions de }%
\end{groupe}%
%% v. 101-129
\begin{groupe}%
\chantAgathFr{la cithare d'Asie, aux beaux rythmes contrariés marqués par le pied}%
\end{groupe}%
%% v. 101-129
\begin{groupe}%
\chantAgathFr{sous la direction des Charites phrygiennes.\\}%
\end{groupe}%
\endnumbering
\end{Rightside}
\end{pages}
\Pages
\clearpage\ifodd\value{page}\hbox{}\newpage\fi
\enTeteG{ΠΡΟΛΟΓΟΣ --- Codex Ravennas 429, f. 164r (Classense 166r) --- vv. 123–164}
\begin{adjustwidth}{-1.5cm}{-3.0cm}
\begin{adjustbox}{max totalheight=21.0cm, max width=16.5cm, keepaspectratio}
\begin{minipage}{16.5cm}
{\fontsize{7.2pt}{8.5pt}\selectfont
\noindent\textbf{Folio 04 --- Codex Ravennas 429, \texttt{f. 164r} (\texttt{Classense 166r}) : ΠΡΟΛΟΓΟΣ, vv. 123–164}\hfill \textit{Témoins :} $\text{R}^1, \Rpc, \Rdeux$, \PsiUn\par\vspace{2pt}
\hrule height 0.4pt\vspace{2.5pt}
}
{\fontsize{8.1pt}{9.5pt}\selectfont
\ligneRav{123}{ἀγάθ/}{\hspace{1.1em}σέβομαι Λατώ τ' ἄνασσαν}{}
\ligneRav{}{}{\hspace{1.1em}κίθαρίν τε ματέρ᾽ ὕμνων}{}
\ligneRav{125}{}{\hspace{1.1em}ἄρσενι βοᾷ δοκιμῳ·}{\textbf{125} \lecon{δοκιμῳ} R, Σ\expo{1}, Suid. α 4016 : δόκιμον Σ\expo{2} : δοκίμων Schöne}
\ligneRav{126}{}{\hspace{1.1em}τᾷ \difference{φάος} ἔσσυτο δαιμονίοις}{\textbf{126} φῶς R : \lecon{φάος} Burges}
\ligneRav{127}{}{\hspace{1.1em}ὄμμασιν \difference{ἁμετέρας} τε δι᾽ αἰφνι-}{\textbf{127} ἡμετέρας R : ὑμετέρας Nietzsche : \lecon{ἁμετέρας} Meineke}
\ligneRav{127}{}{\hspace{1.1em} δίου ὀπός. ὧν χάριν}{\textbf{127} ἡμετέρας R : ὑμετέρας Nietzsche : \lecon{ἁμετέρας} Meineke}
\ligneRav{128}{}{\hspace{1.1em}ἄνακτ᾽ ἄγαλλε \difference{Φοῖβον}·}{\textbf{128} \lecon{ἄγαλλε} R : ἀγάλλετε Austin (1987) \espApp Φοῖβον τιμᾷ R : \lecon{Φοῖβον} Dindorf}
\ligneRav{}{}{\hspace{1.1em}χαῖρ᾽ ὄλβιε παῖ Λατοῦς·}{}
\ligneRav{}{}{\hspace{2.2em}\textbf{[παρεπιγραφή R\expo{1} :]} \lecon{ὀλολύζει ὁ γέρων}}{\textbf{post 129} \lecon{ὀλολύζει ὁ γέρων} (παρεπιγραφή) R}
\ligneRav{130}{μ\auDessus{ν}{Η}}{ὡς ἡδὺ τὸ μέλος ὦ \difference{πότνιαι} Γενετυλλίδες}{\textbf{130} πότνια R : \lecon{πότνιαι} Suid. γ 141}
\ligneRav{}{}{καὶ θηλυδριῶδες καὶ κατεγλωττισμένον}{}
\ligneRav{}{}{καὶ μανδαλωτόν, ὥστ᾽ ἐμοῦ γ᾽ ἀκροωμένου}{}
\ligneRav{}{}{ὑπὸ τὴν ἕδραν αὐτὴν ὑπῆλθε γάργαλος·}{}
\ligneRav{134}{}{καί σ᾽ ὦ \difference{νεανίσκ᾽ εἴ τις} εἶ, κατ᾽ Αἰσχύλον}{\textbf{134-5} νεανίσχ᾽ εἴ τις R : \lecon{νεανίσκ᾽ εἴ τις} \Mon : νεανίσχ᾽ ἥτις Gelenius \espApp \textbf{135} Λυκουργίας R : \lecon{Λυκουργείας} Dindorf}
\ligneRav{135}{}{ἐκ τῆς Λυκουργείας ἐρέσθαι βούλομαι·}{}
\ligneRav{}{}{ποδαπὸς ὁ γύννις; τίς πάτρα; τίς ἡ στολή;}{}
\ligneRav{}{}{τίς ἡ τάραξις τοῦ βίου; τί βάρβιτος}{}
\ligneRav{}{}{λαλεῖ κροκωτῷ; τί δὲ λύρα κεκρυφάλῳ;}{}
\ligneRav{139}{}{τί λήκυθος καὶ στρόφιον; ὡς οὐ ξύμφορον·}{\textbf{139} \lecon{οὐ ξύμ-} R, Suid. β 190 V : ἀξυμ- Suid. M : ὀξύμ- Suid. A,G,I : ξύμ- Suid. F}
\ligneRav{140}{}{τίς δαὶ \difference{κατρόπτου} καὶ ξίφους κοινωνία;}{\textbf{140} \lecon{κατόπτρου} R, Suid. β 110 : \lecon{κατρόπτου} Dover ad \italique{Ran.} 47}
\ligneRav{141}{}{\difference{σύ τ᾽} αὐτὸς ὦ παῖ; πότερον ὡς ἀνὴρ τρέφει;}{\textbf{141} \lecon{σύ τ᾽} \PsiUn, Σ : τίς δ᾽ R}
\ligneRav{142}{}{καὶ ποῦ \difference{πέος}; ποῦ χλαῖνα; ποῦ Λακωνικαί;}{\textbf{142} \lecon{πέος} \PsiUn, R : τὸ πέος Suid. λ 63, π 987}
\ligneRav{}{}{ἀλλ᾽ ὡς γυνὴ δῆτ᾽ εἶτα ποῦ τὰ τιτθία;}{}
\ligneRav{}{}{τί φῄς; τί σιγᾷς; ἀλλὰ δῆτ᾽ ἐκ τοῦ μέλους}{}
\ligneRav{145}{}{ζητῶ σ᾽, ἐπειδή γ᾽ αὐτὸς οὐ βούλει φράσαι;}{}
\ligneRav{}{ἀγάθ/}{ὦ πρέσβυ πρέσβυ, τοῦ φθόνου μὲν τὸν ψόγον}{}
\ligneRav{}{}{ἤκουσα, τὴν δ᾽ ἄλγησιν οὐ παρεσχόμην‎}{}
\ligneRav{148}{}{ἐγὼ δὲ τὴν ἐσθῆθ᾽ ἅμα \difference{τῇ γνώμῃ} φορῶ·}{\textbf{148} γνώμη R : \lecon{τῇ γνώμῃ} Meineke}
\ligneRav{}{}{χρὴ γὰρ ποιητὴν ἄνδρα πρὸς τὰ δράματα}{}
\ligneRav{150}{}{ἃ δεῖ ποιεῖν πρὸς ταῦτα τοὺς τρόπους ἔχειν·}{}
\ligneRav{}{}{αὐτίκα γυναικεῖ᾽ ἢν ποιῇ τις δράματα,}{}
\ligneRav{}{}{μετουσίαν δεῖ τῶν τρόπων τὸ σῶμ᾽ ἔχειν·}{}
\ligneRav{}{μ\auDessus{ν}{Η}}{οὐκοῦν κελητίζεις, ὅταν Φαίδραν ποιῇς;}{}
\ligneRav{154}{ἀγάθ/}{ἀνδρεῖα δ᾽ \difference{ἄν} ποιῇ τις, ἐν τῷ σώματι}{\textbf{154} \lecon{ἄν} \PsiUn, R : ἢν Dindorf}
\ligneRav{155}{}{ἔνεσθ᾽ ὑπάρχον τοῦθ᾽. ἃ δ᾽ οὐ κεκτήμεθα,}{}
\ligneRav{}{}{μίμησις ἤδη ταῦτα συνθηρεύεται·}{}
\ligneRav{}{μ\auDessus{ν}{Η}}{ὅταν σατύρους τοίνυν ποιῇς, καλεῖν ἐμέ,}{}
\ligneRav{}{}{ἵνα συμποιῶ σοὔπισθεν ἐστυκὼς ἐγώ·}{}
\ligneRav{}{ἀγάθ/}{ἄλλως τ᾽ ἄμουσόν ἐστι ποιητὴν ἰδεῖν}{}
\ligneRav{160}{}{ἀγρεῖον ὄντα καὶ δασύν σκέψαι δ᾽ ὅτι}{}
\ligneRav{}{}{Ἴβυκος ἐκεῖνος κἀνακρέων ὁ Τήιος}{}
\ligneRav{162}{}{κἀλκαῖος, οἳ περὶ ἁρμονίαν \difference{ἐκεχήνεσαν},}{\textbf{162} \lecon{κἀλκαῖος} R, Suid. : (κ)ἀχαιός Σ \espApp ἐχύμησαν R : ἐχύμισαν Suid. : \lecon{ἐκεχήνεσαν} Bibi}
\ligneRav{163}{}{ἐμιτροφόρουν τε καὶ \difference{διέκειντ'} Ἰωνικῶς,}{\textbf{163} διεκίνων R : διεκίνουν Suid. ε 989, ι 495 : διεκλῶντ᾽ Toup : \lecon{διέκειντ'} Bibi}
\ligneRav{}{}{καὶ Φρύνιχος, τοῦτον γὰρ οὖν ἀκήκοας,}{}
}
\vspace{3pt}\noindent\rule{4.5cm}{0.4pt}\par\vspace{1.5pt}
{\fontsize{7.0pt}{8.3pt}\selectfont
\noindent\textbf{Loci paralleli (vv. 123–164) ---} \textbf{125} (ἄρσενι βοᾷ δοκιμῳ) Suidas α 4016 $=$ Scholium \espTest \textbf{128} (ἄγαλλε) Photius α 164 $=$ \SynB α 145 \espTest \textbf{130-2} (ἡδὺ -- μανδαλωτόν) Suidas γ 141 \espTest \textbf{130} (Γενετυλλίδες) Hesychius γ 343 Γενετυλλίς \espTest \textbf{131-2} (θηλυδριῶδες -- μανδαλωτόν) Suidas μ 134 \espTest \textbf{131} (θηλυδριῶδες καὶ κατεγλωττισμένον) Suidas κ 912 \espTest Hesychius θ 491 θηλυδριῶδες \espTest (κατεγλωττισμένον) Pollux II.109 \espTest \textbf{132} Hesychius μ 224 μανδαλωτόν $\simeq$ Photius μ 86, cf. Photius γ 115 \espTest \textbf{132-3} (ὥστ᾽ -- γάργαλος) Suidas γ 83 (sed τοῦ μέλους post ἀκροωμένου addito) \espTest \textbf{133} (γάργαλος $=$ γαργαλισμός) Moeris γ 23, cf. \Phryn 56.9 \espTest \textbf{134} Hesychius η 948 ἥτις· ὅστις (?) \espTest \textbf{136} (ποδαπὸς -- πάτρα) Suidas γ 504 \espTest (γύννις) Hesychius γ 1015 γύννις, cf. \Phryn 60.19 \espTest \textbf{136-40} (τίς ἡ στολή -- κοινωνία) Suidas β 110 (δὲ pro δαὶ) \espTest \textbf{137} (βάρβιτος) Hesychius β 211 βάρβιτος, cf. Photius β 62 \espTest \textbf{141-2} (πότερον -- πέος) Suidas π 987 \espTest \textbf{142} (ποῦ τὸ πέος -- Λακωνικαί) Suidas λ 63, π 987 (τὸ πέος pro πέος)\espTest Moeris χ 19 χλαῖνα \espTest (Λακωνικαί) Hesychius λ 225 Λακωνικαί ; Photius λ 57 Λακωνικαί ; scholium \italique{Vespae} 1158-60 \espTest \textbf{142} (ποῦ τὸ πέος -- Λακωνικαί) Suidas λ 63 \espTest Moeris χ 19 χλαῖνα \espTest (Λακωνικαί) Hesychius λ 225 Λακωνικαί ; Photius λ 57 Λακωνικαί ; scholium \italique{Vespae} 1158-60 \espTest \textbf{146-7} Suidas α 1097 (μὲν omisso et δὲ pro δ') \espTest \textbf{149-50} \Apostol XVIII 41h \espTest \textbf{159-61} Suidas α 1633 (κ' Ἀνακρ.) \espTest \textbf{161-2} Suidas ε 4049 (Ἀνακρέων, Ἀλκαῖος, οἵπερ ἁρμονίαν ἐχύμισαν) \espTest \textbf{162-3} (κ' Ἀλκαῖος -- Ἰωνικῶς) Suidas ε 989 (sed οἵπερ ἁρμονίαν ἐχύμισαν) $\simeq$ Zonaras 707 ἐμιτροφόρουν \espTest \textbf{162} Zonaras 942 οἵπερ ἁρμονίαν ἐχύμισαν \espTest (ἐχύμισαν) Hesychius ε 7681 ἐχύμισαν· ἤρτυσαν \espTest \textbf{162-3} (κ' Ἀλκαῖος -- Ἰωνικῶς) Suidas ε 989 $\simeq$ Zonaras 707 ἐμιτροφόρουν \espTest \textbf{162} Zonaras 942 οἵπερ ἁρμονίαν ἐχύμισαν \espTest (ἐχύμισαν) Hesychius ε 7681 ἐχύμισαν \espTest \textbf{163} Suidas ι 495 ἐμιτροφόρουν τε καὶ διεκίνουν Ἰωνικῶς\par
}
\vspace{2.5pt}\noindent\rule{4.5cm}{0.4pt}\par\vspace{1.5pt}
{\fontsize{6.9pt}{8.2pt}\selectfont
\noindent\textbf{Suite (en commentaire) de l'apparat critique \& notes paléographiques (\texttt{f. 164r}) ---} \textbf{Témoin papyrologique :} \textbf{\PsiUn{} (\italique{PSI} XI 1194 + XIV p. XV)} couvre les \textbf{vv. 133–156}. --- \textbf{Deuxième \italique{parepigraphè} (entre vv. 129 et 130) :} $\text{R}^1$ copie sur une ligne propre \lecon{ὀλολύζει ὁ γέρων}. --- \textbf{Confrontation avec \PsiUn{} (v. 141) :} Au v. 141, R porte la leçon fautive τίς δ' αὐτὸς ὦ παῖ· πότερον ὡς ἀνὴρ τρέφει;, tandis que le papyrus \PsiUn{} et la scholie $\Sigma^{\text{R}}$ conservent l'authentique \lecon{σύ τ'}. --- \textbf{Double émendation paléographique de \textbf{Bibi} aux vv. 162–163 :} ἐχύμησαν ... διεκίνων R $\to$ \lecon{ἐκεχήνεσαν ... διέκειντ'} \textbf{Bibi}. --- \textbf{v. 125 :} \lecon{δοκίμῳ} R, $\Sigma^1$, Souda $\alpha\ 4016$ : δόκιμον $\Sigma^2$ : δοκίμων Schöne. --- \textbf{v. 127 :} ἡμετέρας R (ὑμετέρας Nietzsche) : \lecon{ἁμετέρας} Meineke (dorisme lyrique). --- \textbf{v. 128 :} \lecon{ἄγαλλε} R (ἀγάλλετε Austin 1987) \espApp Φοῖβον τιμᾷ R (τιμᾷ est une glose intruse) : \lecon{Φοῖβον} Dindorf. --- \textbf{vv. 134–135 :} νεανίσχ' εἴ τις R (aspiration fautive par anticipation de l'esprit rude) : \lecon{νεανίσκ' εἴ τις} \Mon{} (νεανίσχ' ἥτις Gelenius) \espApp \textbf{v. 135 :} Λυκουργίας R (iotacisme) : \lecon{Λυκουργείας} Dindorf. --- \textbf{v. 141 :} \lecon{σύ τ'} $\Psi^1, \Sigma$ : τίς δ' R. --- \textbf{v. 162 :} \lecon{κἀλκαῖος} R, Souda ((κ)ἀχαιός $\Sigma$) \espApp ἐχύμησαν R (ἐχύμισαν Souda) : \lecon{ἐκεχήνεσαν} \lecon{Bibi}. \textbf{Explication paléographique (Commentaire ad v. 162) :} Dans ἐκεχήνεσαν (plus-que-parfait de $\text{χάσκω}$, « ils restaient bouche bée », cf. \italique{Ach.} 10), le saut visuel εκεχ- $\to$ εχ- (haplographie), joint à l'iotacisme $\eta \to \upsilon$ et à la confusion minuscule constante $\nu \to \mu$, a produit le monstre ἐχύμησαν de R ! --- \textbf{v. 163 :} διεκίνων R (διεκίνουν Souda $\varepsilon\ 989, \iota\ 495$ ; διεκλῶντ' Toup) : \lecon{διέκειντ'} \lecon{Bibi} (ou δίκανον ap. Hésychius $\delta\ 1805$).\par
}
\end{minipage}
\end{adjustbox}
\end{adjustwidth}
\newpage
\enTeteD{Codex Ravennas 429 (Biblioteca Classense) --- Fac-similé du f. 164r (vv. 123–164)}
\begin{adjustwidth}{-1.5cm}{-3.0cm}
\begin{center}
\includegraphics[width=16.5cm,height=20.5cm,keepaspectratio]{/workspaces/Thesmophories/dossier_ravennas_thesmophories/folios_individuels_ff162v-177r/04_f164r_classense166r_vv0123-0164.jpg}\par
\vspace{2pt}
{\scriptsize\textbf{Codex Ravennas 429, f. 164r (Classense 166r)} --- ΠΡΟΛΟΓΟΣ, \textbf{vv. 123–164} ($\text{R}^1, \Rpc, \Rdeux$, \PsiUn).}
\end{center}
\end{adjustwidth}
\newpage
\clearpage\ifodd\value{page}\hbox{}\newpage\fi
\enTeteG{ΘΕΣΜΟΦΟΡΙΑΖΟΥΣΑΙ --- Texte grec établi (f. 164r, vv. 123–164)}
\enTeteD{LES THESMOPHORIES --- Traduction française (f. 164r, vv. 123–164)}
\begin{pages}
\begin{Leftside}
\beginnumbering
%% v. 101-129
\begin{groupe}%
\typeVers{io. min. + troch.}\chantAgath{\ligne{123}σέβομαι Λατώ τ' ἄνασσαν}%
\end{groupe}%
%% v. 101-129
\begin{groupe}%
\typeVers{2 io. min. anacl.}\chantAgath{κίθαρίν τε ματέρ᾽ ὕμνων}%
\end{groupe}%
%% v. 101-129
\begin{groupe}%
\typeVers{troch. + io. min\cat}\chantAgath{ἄρσενι βοᾷ δοκιμῳ.%\testimonia{125}{(ἄρσενι βοᾷ δοκιμῳ) Suidas α 4016 $=$ Scholium}
\apparat{125}{\lecon{δοκιμῳ} R, Σ\expo{1}, Suid. α 4016 : δόκιμον Σ\expo{2} : δοκίμων Schöne \ligneApp{126} φῶς R : \lecon{φάος} Burges}\\}%
\end{groupe}%
%% v. 101-129
\begin{groupe}%
\typeVers{4 dact\cat}\chantAgath{\ligne{126}τᾷ φάος ἔσσυτο δαιμονίοις}%
\end{groupe}%
%% v. 101-129
\begin{groupe}%
\typeVers{6 dact. acatal.}\chantAgath{ὄμμασιν ἁμετέρας τε δι᾽ αἰφνι-\\
\alinea δίου ὀπός. ὧν χάριν\apparat{127}{ἡμετέρας R : ὑμετέρας Nietzsche : \lecon{ἁμετέρας} Meineke \ligneApp{128} \lecon{ἄγαλλε} R : ἀγάλλετε Austin (1987) \espApp Φοῖβον τιμᾷ R : \lecon{Φοῖβον} Dindorf}}%
\end{groupe}%
%% v. 101-129
\begin{groupe}%
\typeVers{2 iamb\cat}\chantAgath{\ligne{128}ἄνακτ᾽ ἄγαλλε Φοῖβον.%\testimonia{128}{(ἄγαλλε) Photius α 164 $=$ \SynB α 145}
\\}%
\end{groupe}%
%% v. 101-129
\begin{groupe}%
\typeVers{telesil. scazon}\chantAgath{\ligne{129}χαῖρ᾽ ὄλβιε παῖ Λατοῦς.%\apparat{post v. 129}{parepigraphe ὀλολύζει R : ὀλολύζει ὁ γέρων Suid. ο 194}
\\\passeLigne}%
\end{groupe}%
%% v. 130-145
\begin{groupe}%
\personnageGrec{Μνησίλοχος}%
\typeVers{3 ia.}\trimIamb{ὡς ἡδὺ τὸ μέλος ὦ πότνιαι Γενετυλλίδες%\testimonia{130-2}{(ἡδὺ -- μανδαλωτόν) Suidas γ 141}\testimonia{130}{(Γενετυλλίδες) Hesychius γ 343 Γενετυλλίς}
\apparat{130}{πότνια R : \lecon{πότνιαι} Suid. γ 141}%
\commentaire{130}{ὦ πότνιαι Γενετυλλίδες.}{Déesses de la génération, les Génétyllides sont aussi citées par Aristophane au vers 52 des \italique{Nuées} et au second vers de \italique{Lysistrata} ; de ce dernier passage on peut inférer que les femmes étaient réputées célébrer des rites licencieux en leur honneur.}\\
καὶ θηλυδριῶδες καὶ κατεγλωττισμένον%\testimonia{131-2}{(θηλυδριῶδες -- μανδαλωτόν) Suidas μ 134}\testimonia{131}{(θηλυδριῶδες καὶ κατεγλωττισμένον) Suidas κ 912 \espTest Hesychius θ 491 θηλυδριῶδες \espTest (κατεγλωττισμένον) Pollux II.109}%
\commentaire{131}{θηλυδριῶδες.}{Semble être un néologisme d'Aristophane. Chantraine\marqueur{1} suppose un relais *θηλύδριον qui expliquerait aussi θηλυδρίας, bien attesté. Aristophane a-t-il porté son choix sur θηλυδριώδης pour faire entendre θηριώδης et toute l'animalité du spasme de Mnésiloque ? (cf. κατεγλωττισμένον et μανδαλωτόν.)}%
\debutNote{131}%
\note{1}{\italique{Dictionnaire étymologique de la langue grecque}, \italique{sub voce} θηλή.}\\
καὶ μανδαλωτόν, ὥστ᾽ ἐμοῦ γ᾽ ἀκροωμένου%\testimonia{132}{Hesychius μ 224 μανδαλωτόν $\simeq$ Photius μ 86, cf. Photius γ 115}\testimonia{132-3}{(ὥστ᾽ -- γάργαλος) Suidas γ 83 (sed τοῦ μέλους post ἀκροωμένου addito)}%
\commentaire{131-2}{καὶ κατεγλωττισμένον / καὶ μανδαλωτόν}{: ces mots renvoient tous deux à ce que les anglophones appellent – ne boudons pas l'hommage qui nous est ainsi rendu ! – \italique{french kiss}. À propos de μανδαλωτόν, je ne résiste pas à citer Taillardat (§195) qui met au jour la bigarrure de la langue aristophanienne : « le substantif, τὸ φίλημα, \italique{le baiser}, est déjà connu d'Eschyle ; il se trouve chez Sophocle, chez Euripide, chez les comiques Platon et Nicophon. En revanche, il n'est pas attesté dans les œuvres conservées d'Aristophane, bien que le poëte emploie φιλεῖν comme synonyme de κυνεῖν [...] Aristophane préfère des mots imagés et comme ``techniques'' : τὸ περιπεταστόν (\italique{Acharniens} 1201), τὸ ἐπιμανδαλωτόν, \italique{ib.}, τὸ μανδαλωτόν (\italique{Thesm.} 132), les deux derniers termes évoquant l'idée d'une porte fermée ; ὁ μάνδαλος, comme βάλανος, désigne la clavette assurant la barre de la porte. La clavette, c'est ici la langue. »}\\
ὑπὸ τὴν ἕδραν αὐτὴν ὑπῆλθε γάργαλος.%\testimonia{133}{(γάργαλος $=$ γαργαλισμός) Moeris γ 23, cf. \Phryn 56.9}%
%\commentaire{133}{ὑπὸ τὴν ἕδραν αὐτὴν ὑπῆλθε γάργαλος.}{Rabelais, \italique{Gargantua}, XIII :\citComm{\francais{[sa soye] me causoit au fondement une volupté bien grande.}}}
\\
καί σ᾽ ὦ νεανίσκ᾽ εἴ τις εἶ, κατ᾽ Αἰσχύλον%\testimonia{134}{Hesychius η 948 ἥτις· ὅστις (?)}%
\apparat{134-5}{νεανίσχ᾽ εἴ τις R : \lecon{νεανίσκ᾽ εἴ τις} \Mon : νεανίσχ᾽ ἥτις Gelenius \espApp \textbf{135} Λυκουργίας R : \lecon{Λυκουργείας} Dindorf}%
\commentaire{134}{νεανίσχ᾽ εἴ τις εἶ.}{Le manuscrit de Ravenne comporte une coquille : le κ de νεανίσκ᾽ est aspiré alors qu'il est suivi de εἴ. Plusieurs solutions ont été apportées : lire ἥτις (justifiant le féminin par Hésychius (η 948) glosant ἥτις· ὅστις). D'autres se sont contentés de substituer un κ au χ. J'ai une préférence pour cette dernière correction\marqueur{1} qui me semble être une pique contre Agathon : « jeune homme... si tu en es un ! » car Agathon serait âgé au moment de la représentation de 35 ans ou environ ; à croire que son mode de vie le conserve bien... mais tout-de-même ! cette allusion permanente à la jeunesse d'Agathon – comme dans le \italique{Banquet} de Platon censé se passer en 416, cinq ans plus tôt – me paraît plutôt être une allusion au fait qu'Agathon pèche par certain attribut pour être tout-à-fait considéré comme un homme\marqueur{2}.}%
\debutNote{134}%
\note{1}{Signalons qu'une coquille analogue (οὐχ᾽ εἴρηχ᾽ R en place d'un οὐκ εἴρηχ᾽) se retrouve au vers 556. De plus, les anciens manuscrits sur lesquels on a copié les \italique{Thesophories} n'étaient pas systématiquement élidés ; ce genre de fautes s'explique donc aisément. Cf. par exemple v. 148 ἐσθητ' ἅμα ; v. 279 θραττ' έπου ; v. 1178 ἔρχετ' ὡς. }%
\note{2}{καὶ ποῦ πέος; demande Mnésiloque au vers 142.}%
\\
ἐκ τῆς Λυκουργείας ἐρέσθαι βούλομαι.\\
\commentaire{134-5}{κατ᾽ Αἰσχύλον / ἐκ τῆς Λυκουργείας.}{D'après le scholiaste, la \italique{Lycurgie} est le nom collectif d'une tétralogie d'Eschyle, composée des \italique{Édoniens} (fr.\marqueur{1} 57-67), des \italique{Bassarides} (fr. 23-25), des \italique{Jouvenceaux} (Νεανίσκοι, fr. 146-9) et du drame satyrique \italique{Lycurgue} (fr. 124-6). Elle compte les mésaventures du roi Lycurgue essayant de contrer l'invasion du culte dionysiaque en Thrace. Laissons la parole à François Jouan\marqueur{2} : « La légende de Lycurgue était déjà connue d'Homère. L'\italique{Iliade} rappelait sommairement que Lycurgue avait été aveuglé par Zeus pour avoir pourchassé les Nourrices de Dionysos et forcé le petit dieu à se jeter à la mer\marqueur{3}. La version eschyléenne est très différente, et comporte en particulier une localisation précise : le héros est roi des Édoniens et vit en Thrace, dans la région du Pangée. On admet en général que cette tétralogie [...] datait de la fin de la vie d'Eschyle, entre 466 et 459. Or dans cette période, des forces athéniennes commandées par Cimon guerroyaient précisément contre les Édoniens, autour d'Eion et d'Ennéahodoi. [...] Les \italique{Édoniens} sont la pièce la moins mal connue, en particulier du fait de l'existence de fragments d'un \italique{Lycurgue} du poète latin Naevius, qui semble avoir suivi d'assez près le drame d'Eschyle. Dans le prologue, un messager annonçait au roi thrace l'arrivée d'un personnage entouré d'un bruyant cortège de musiciens et de danseurs – Dionysos et son thiase. Le roi ordonnait de les arrêter et de les faire comparaître devant lui. Dans la parodos, le chœur, sans doute composé de sujets de Lycurgue, décrivait le cortège dionysiaque. Un garde rendait compte de l'arrestation des trublions et de leur docilité. Le roi interrogeait lui-même Dionysos, comme le Penthée des \italique{Bacchantes}. Aussi mécontent des réponses de Dionysos que le prince thébain, il l'emprisonnait avec sa suite dans son palais, où il devait trouver, lui disait-il ironiquement «~un logement gratuit [ἄδασμος, fr. 63]~». Mais la défaite du dieu n'était que provisoire. Bientôt les liens des captifs tombaient d'eux-mêmes, Dionysos apparaissait dans sa gloire sur le toit du palais, qui vacillait sur sa base et se nimbait de flammes : autre motif inclus dans les \italique{Bacchantes}\marqueur{5}. » \\
C'est donc bien à cette première tragédie des \italique{Édoniens} que semble se référer Mnésiloque, interrogeant à la manière de Lycurgue, associant ainsi Agathon à Dionysos (métaphore qui avait commencée d'être filée au vers 41 quand le serviteur d'Agathon parlait du thiase de Muses réunies chez son maître). Rappelons que, pour les Grecs de l'époque, Dionysos était un dieu jeune : « Les particularités de son physique et de son costume sont détaillées aussi bien par Eschyle\marqueur{6} que par Euripide. Il est grand et mince avec de longues jambes, une chevelure blonde, abondante et parfumée, de belles boucles. Il est beau, avec un teint vermeil et des traits de visage plutôt féminins, l'air souriant. Il est vêtu d'une robe asiatique de fin tissu, tombant jusqu'à ses pieds et tient en main le thyrse. »\marqueur{7} 
(Ajoutons à propos de cette \italique{crocote} au cœur des péripéties de Mnésiloque qu'elle est aussi associée à Dionysos dans les \italique{Grenouilles} (v.46) ; par Cratinos, fr. 40 de l'édition de Kassel-Austin et, comme on l'a vu, Eschyle, fr. 59.) 
Voilà qui explique un peu les associations d'idées de Mnésiloque !}%
\debutNote{134-5}%
\note{1}{Les fragments d'Eschyle sont cités selon l'édition de S. Radt, \italique{Tragicorum Graecorum Fragmenta}, vol. 3 (1985).}%
\note{2}{On lira avec intérêt deux de ses articles : « Dionysos chez Eschyle » (in \italique{Kernos} 5, 1992, p. 71-86) et « Une introduction au mythe des \italique{Bacchantes} » (\italique{Kentron} 14, 1998). C'est le premier que nous citons. }%
\note{3}{\italique{Iliade}, VI, 130-140.}%
\note{4}{Eur. \italique{Bacch}., 586-599 ; 622-624 ; 633.}%
\note{5}{Cf. Hérodote, II, 145.}%
\note{6}{Voir les fragments 59, 60, 62 des \italique{Édoniens}.}%
\note{7}{F. Jouan, \italique{Une introduction au mythe des} Bacchantes, p. 12.}%
ποδαπὸς ὁ γύννις; τίς πάτρα; τίς ἡ στολή;%\testimonia{136}{(ποδαπὸς -- πάτρα) Suidas γ 504 \espTest (γύννις) Hesychius γ 1015 γύννις, cf. \Phryn 60.19}\testimonia{136-40}{(τίς ἡ στολή -- κοινωνία) Suidas β 110 (δὲ pro δαὶ)}
}
\end{groupe}%
%% v. 130-145
\begin{groupe}%
\trimIamb{τίς ἡ τάραξις τοῦ βίου; τί βάρβιτος%\testimonia{137}{(βάρβιτος) Hesychius β 211 βάρβιτος, cf. Photius β 62}
\\
λαλεῖ κροκωτῷ; τί δὲ λύρα κεκρυφάλῳ;\\
%\commentaire{138}{τί δὲ λύρα κεκρυφάλῳ;}{Mnésiloque questionne Agathon sur les contrastes de son accoutrement. Il relève au vers précédent l'étrange vis-àv-is du βάρβιτος et du κροκωτός, attributs aussi portés par le Dionysos d'Eschyle ; il est donc probable qu'Aristophane soit toujours en train de pasticher la \italique{Lycurgie} (encore que la tournure avec le verbe λαλεῖν fasse davantage penser au style d'Aristophane qu'à celui d'Eschyle).\\
%Le vers 138,τί δὲ λύρα κεκρυφάλῳ, semble une redite du précédent, puisque λύρα est un synonyme de βάρβιτος et que κεκρύφαλος est un autre des atours d'une toilette soignée. En fait l'apparente répétition permet à Aristophane de jouer avec un proverbe en le détournant. (procédé identique au vers 140.)\\
%Aulu-Gelle nous a conservé la version latine de ce \italique{vieil adage} :
%\citComm{\latin{Nihil graculo cum fidibus}\\
%\francais{Une lyre et un choucas, ça fait deux}}%
%\eff{Equites 1021 : τί γάρ ἐστ᾽ Ἐρεχθεῖ καὶ κολοιοῖς καὶ κυνί;}
%Laissons la parole à Érasme, à qui nous devons grâce à ses précieux \italique{Adages}\marqueur{1} cette découverte :
%\citComm{\latin{Nihil graculo cum fidibus, id est nihil stolidis et imperitis cum bonis literis. Nam graculus avis est ineptae molestaeque garrulitatis atque obstrepera, adeo ut a voce qua sonat, nomen inditum putent auctore Quintiliano primo Institutionum libro. Tum sui generis conventus amat, quo mutuo garritu strepitus fiat odiosior. Porro cithara silentium postulat et aures attentas. A. Gellius Noctium Atticarum extrema lucubratione sane quam eleganter hoc adagium torquet in quosdam pinguiore ingenio homines stolidaque loquacitate praeditos, profanos, ἀμούσος, id est, a musarum sacris alienos, qui politiores literas elegantioremque doctrinam vel ridere possunt vel etiam contemnere, intelligere nequaquam possunt, quos ea demum juvant, quae ipsi scilicet didicerunt, sordida, muta, indocta. Hi pari importunitate et alienae doctrinae obturbant et suam inscitiam passim venditant et inculcant neminem prae se contemnentes. Hoc igitur hominum genus Gellius a suis commentariis, quibus nihil fieri potest neque tersius neque eruditius, quasi fuste quodam abigens, Vetus est, inquit, adagium, nihil graculo cum fidibus.}\\
%\francais{Une lyre et un choucas, ça fait deux. C'est-à-dire rien de commun entre les imbéciles ou les têtes creuses et les belles-lettres. En effet, le choucas est un oiseau qui casse les oreilles par son croassement ridicule et désagréable ; on pense d'ailleurs qu'il tient son nom du cri qu'il fait entendre, voyez Quintilien dans le livre I [6, 37] de l'\italique{Institution oratoire}. En outre, il aime faire bande avec ceux de son espèce [cf. \italique{Adages}, 123] : leur échange de criailleries en devient d'autant plus insupportable. À l'opposé, la cithare exige du silence et des oreilles attentives. À la toute fin de sa préface des \italique{Nuits attiques} (19-21), Aulu-Gelle avec beaucoup d'à-propos tourne cet adage contre ces hommes à l'esprit particulièrement épais, dotés d'un creux bagout, prêts à tout salir, \italique{amousoï}, c'est-à-dire étrangers aux Muses [cf. \italique{Adages}, 1518], qui peuvent se gausser d'une culture particulièrement raffinée ou d'un vaste savoir, et même en faire fi. Ils n'y peuvent rien comprendre, et ne trouvent de charme qu'aux grossièretés, aux inepties sans nom qu'ils ont eux-mêmes apprises. Tout aussi fielleux pour déranger qui se cultive, ils font étalage en tout lieu de leur ignorance et l'imposent à toute force, méprisant le monde entier sauf eux-mêmes. C'est ce genre d'hommes qu'Aulu-Gelle dans ses commentaires — d'une élégance et d'une érudition hors pair — malmène à coups de bâton pour ainsi dire : « Il existe un vieil adage : rien de commun entre un choucas et une lyre. »}}
%Si les \latin{fidibus} latines correspondent à la λύρα grecque ; le \latin{graculo}, \italique{choucas}, pourrait se traduire en grec par le mot κολοιός ou bien κόκκυξ. C'est le mot κόκκυξ qui devait être utilisé pour le proverbe, comme on peut l'inférer de la glose d'Hésychius (κ 3298) suivante :  
%\citComm{κόκκυξ· τὸ ζῷον. ἢ τὸ ὄρνεον [ ... ] καὶ λόφος, καὶ \textbf{περικεφαλαία}.\\
%coucou : l'animal [un poisson]. Ou l'oiseau. Ou une crête. Ou un couvre-chef.}
%Le mot κόκκυξ pouvait avoir le sens de περικεφαλαία ; Aristophane use donc de κεκρυφάλῳ comme "permutant méthaphorique" -- ainsi parlait Taillardat -- pour κόκκυγι.\\
%Il résulte que le \italique{vieil adage} dont parle Aulu-Gelle devait se rencontrer chez les Grecs, peu ou prou sous la forme :
%\citComm{τί λύρα λαλεῖ κόκκυγι;}}%
%\debutNote{138}%
%\note{1}{Je cite les \italique{Adages} ainsi que leur traduction française selon l'édition de Jean-Christophe Saladin aux Belles Lettres (2019). Pour l'adage qui nous concerne, il s'agit du \nieme{337}.}\\
\ligne{139}τί λήκυθος καὶ στρόφιον; ὡς οὐ ξύμφορον.\apparat{139}{\lecon{οὐ ξύμ-} R, Suid. β 190 V : ἀξυμ- Suid. M : ὀξύμ- Suid. A,G,I : ξύμ- Suid. F \ligneApp{140} \lecon{κατόπτρου} R, Suid. β 110 : \lecon{κατρόπτου} Dover ad \italique{Ran.} 47 \ligneApp{141} \lecon{σύ τ᾽} \PsiUn, Σ : τίς δ᾽ R}% \ligneApp{142} \lecon{πέος} \PsiUn, R : τὸ πέος Suid. λ 63, π 987}%
\commentaire{139-44}{Fragments d'un ancien papyrus.}{L'on possède quelques fragments d'un ancien papyrus (\PsiUn \ dans notre apparat) du premier ou second siècle ap. J.-C. pour les vers 133-144 ; 145-156 ; 237-245 ; 273-288 ; 594-596 ; 804-809. Il nous apprend en outre le \latin{numerus versuum} d'une ancienne édition des \italique{Thesmophories} : la pièce comportait alors 1201 vers (contre 1231 pour celle de Brunck sur laquelle se basent tous les éditeurs modernes). La papyrus est reproduit en annexe. Cf. \italique{Papiri della Società Italiana} XI, n° 1194 et  \italique{Papiri della Società Italiana} XIV, p. XV.}%
\commentaire{139}{τί λήκυθος καὶ στρόφιον;}{Dans sa liste des attributs masculins et féminins qui jurent ensemble chez Agathon, Mnésiloque oppose ici le \italique{lécythe} et le \italique{strophion}. Si le cas du strophion ne pose pas de problème : il s'agit du \latin{strophium} ou \latin{fascia pectoralis}\marqueur{1} : bandelette que les femmes enroulaient au niveau de leur poitrine pour la soutenir -- bref : un soutien-gorge --, en revanche, le lécythe a étonnamment égaré les commentateurs: ils lui ont donné son sens (courant) de \italique{fiole à huile}. D'aucuns ont même déployé force énergie pour tenter de montrer en quoi cet accessoire était typiquement masculin (!?) 
Ils eussent plus vite et mieux fait d'ouvrir n'importe quel lexique\marqueur{2} : ils y auraient appris que λήκυθος peut aussi désigner la \italique{pomme d'Adam}. On comprend mieux l'effarement de Mnésiloque face à cette bandelette toute féminine entourant la poitrine désespérément plate d'Agathon, surmontée d'une proéminente et peu ragoûtante pomme d'Adam ; admettez que cela est d'assez mauvais goût (ὡς οὐ ξύμφορον)\marqueur{3}.}%
\debutNote{139}%
\note{1}{Je renvoie ici aux articles correspondants de la monumentale somme qu'est le \italique{Dictionnaire des antiquités grecques et romaines} de Daremberg et Saglio}%
\note{2}{Les dictionnaires classiques : Bailly, \LSJ, \&c. Hésychius (λ 857): τὸ μεταξὺ τοῦ λαυκανίου καὶ αὐχένος ήχῶδες. Cf. le verbe ληκυθίζειν = λαρυγγίζειν = φαρυγγίζειν, \italique{donner de la voix}.}%
\note{3}{Faut-il en conclure que la présence d'une pomme d'Adam était caractéristique du masque comique représentant Agathon ?}\\
τίς δαὶ κατρόπτου καὶ ξίφους κοινωνία;\\
\commentaire{140}{κατόπτρου.}{La variante κάτροπτον (avec métathèse du ρ) est la seule attestée par l'épigraphie pour les inscriptions contemporaines de la représentation des \italique{Thesmophories}\marqueur{1}. Ce n'est qu'à partir de la fin du \nieme{IV} siècle que la forme κάτοπτρον, plus conforme à l'étymologie, semble prendre le dessus.}%
\debutNote{140}%
\note{1}{Voir par exemple les \italique{Inscriptiones Graecae} XXII.1471.47, 1544.58.
Le lecteur pourra aisément vérifier ces dires par lui-même en papillonnant sur les \italique{Inscriptiones Graecae} à l'adresse suivante : \href{https://epigraphy.packhum.org/}{https://epigraphy.packhum.org/}}%
σύ τ᾽ αὐτὸς ὦ παῖ; πότερον ὡς ἀνὴρ τρέφει;%\testimonia{141-2}{(πότερον -- πέος) Suidas π 987}
\\
\commentaire{140}{τίς δαὶ κατόπτρου καὶ ξίφους κοινωνία;}{Le sel de ce vers ne se peut goûter qu'à condition d'avoir en tête le proverbe grec\marqueur{1} :
\citComm{τίς γὰρ κατόπτρῳ καὶ τυφλῷ κοινωνία;\\
Qu'y a-t-il de commun entre un aveugle et un miroir ?}
Ce proverbe a été répertorié par les parémiographes\marqueur{2} dans sa version concise (κάτοπτρον ὁ τυφλός) avec d'autres associations tout aussi inutiles. Ainsi : φαλακρὸς κτένα , \italique{un chauve un peigne}.
L'on trouvait bien évidemment des proverbes équivalents dans notre vieux parler ; \italique{À l'aveugle ne duit} [$=$ convient]\italique{ peinture / Couleur, miroir ne figure}\marqueur{3}.
Aristophane joue par paronomase entre ξίφους et τυφλοῦ.
%\eff{Pour l'épée, Agathon l'a-t-elle car elle fait partie de l'attirail de Dionysos ? mais alors, et le reste ?}
Au vers 138, il devait sans doute se livrer à un détournement du même genre : on recense de nombreux proverbes mettant en relation une lyre avec un âne (ex. : τί γὰρ κοινὸν ... λύρᾳ καὶ ὄνῳ, voir Tosi n° 1934), animal lourd et rétif aux muses ; ou une lyre avec un choucas (\latin{Nihil graculo cum fidibus}, \italique{Adages} d'Érasme n° 337), oiseau criard.
}%
\debutNote{140}%
\note{1}{\nieme{3} vers d'un fragment comique adespote cité par Stobée (IV, 30, 6a) ; c'est le \nieme{903} fragment des \italique{Poetae Comici Graeci, Adespota}, VIII (1995).}%
\note{2}{\italique{Corpus paroemiographorum Graecorum}, II, p. 774 (\italique{Mantissae Proverbiorum} 100).}%
\note{3}{\italique{Le livre des proverbes français} de Le Roux de Lincy, Série n°V, à l'entrée \italique{Aveugle}. À l'initiative du génial Pierre Boutang, le livre a été réédité en 1996 chez Hachette, augmenté d'une sienne préface dans laquelle -- entre autres merveilles -- Boutang questionne avec bonheur l'étymologie du \italique{pro-verbe}.}%
%\commentaire{140}{κατόπτρου καὶ ξίφους}{Sur un fameux cratère\marqueur{1} (le « cratère de Würzburg », Martin-von-Wagner-Museum der Universitàt Wurzburg, H 5697 ) représentant la scène d’égorgement de la ‘petite’ de Mika, on voit un miroir suspenduest-ce le miroir que l'on voit sur le cratère de Würzburg ? Si oui, il resterait à élucider pourquoi il est représenté lors de la scène beaucoup plus lointaine de la parodie du \italique{Télèphe}.\eff{Pour ce qui est du rapport des miroirs et des épées – pourquoi Agathon aurait-il une épée ? – la réponse se trouvait peut-être dans les pièces perdues d'Eschyle... ou plutôt d'Agathon. Sinon, alors voir la justification de Dover, Grenouilles 47 pour changer κατόπτρου en κατρόπτου et si c'est au prétexte comme je le soupçonne d'une pureté étymologique, pousser un coup de gueule !!}}%
καὶ ποῦ πέος; ποῦ χλαῖνα; ποῦ Λακωνικαί;%\testimonia{142}{(ποῦ τὸ πέος -- Λακωνικαί) Suidas λ 63, π 987 (τὸ πέος pro πέος)\espTest Moeris χ 19 χλαῖνα \espTest (Λακωνικαί) Hesychius λ 225 Λακωνικαί ; Photius λ 57 Λακωνικαί ; scholium \italique{Vespae} 1158-60}
\\
ἀλλ᾽ ὡς γυνὴ δῆτ᾽ εἶτα ποῦ τὰ τιτθία;\\
τί φῄς; τί σιγᾷς; ἀλλὰ δῆτ᾽ ἐκ τοῦ μέλους\\%
%\commentaire{144}{ἐκ τοῦ μέλους}{« d'après ton chant », serait-on tenté de traduire ; mais le μέλος (litt. : \italique{membre}, \italique{partie}) peut désigner aussi une partie du corps. Y a-t-il un jeu de mots ?}
ζητῶ σ᾽, ἐπειδή γ᾽ αὐτὸς οὐ βούλει φράσαι;\\\passeLigne}%
\end{groupe}%
%% v. 146-152
\begin{groupe}%
%
\personnageGrec{Ἀγάθων}%
\trimIamb{ὦ πρέσβυ πρέσβυ, τοῦ φθόνου μὲν τὸν ψόγον%\testimonia{146-7}{Suidas α 1097 (μὲν omisso et δὲ pro δ')}
\\
ἤκουσα, τὴν δ᾽ ἄλγησιν οὐ παρεσχόμην‎\\
ἐγὼ δὲ τὴν ἐσθῆθ᾽ ἅμα τῇ γνώμῃ φορῶ.\apparat{148}{γνώμη R : \lecon{τῇ γνώμῃ} Meineke}\\
χρὴ γὰρ ποιητὴν ἄνδρα πρὸς τὰ δράματα%\testimonia{149-50}{\Apostol XVIII 41h}
\\
ἃ δεῖ ποιεῖν πρὸς ταῦτα τοὺς τρόπους ἔχειν.\\
αὐτίκα γυναικεῖ᾽ ἢν ποιῇ τις δράματα,\\
μετουσίαν δεῖ τῶν τρόπων τὸ σῶμ᾽ ἔχειν.\\\passeLigne}%
\end{groupe}%
%% v. 153
\begin{groupe}%
\personnageGrec{Μνησίλοχος}%
\trimIamb{οὐκοῦν κελητίζεις, ὅταν Φαίδραν ποιῇς;
\commentaire{153}{κελητίζεις, ὅταν Φαίδραν ποιῇς.}{Allusion à la \italique{pendula Venus} des Latins. Cf. \italique{Guêpes}, 500 ; \italique{Lysistrata}, 60 et 773. Le poète Machon disait καθιππάζειν (Athénée, XIII, p. 581d). (Note d'Alphonse Willems. Phèdre était éprise d'Hippolyte, dont le nom signifie littéralement le «~délieur de chevaux~» ; c'est-à-dire celui dont les chevaux courent à bride abattue.)}\\\passeLigne}%
\end{groupe}%
%% v. 154-156
\begin{groupe}%
\personnageGrec{Ἀγάθων}%
\trimIamb{ἀνδρεῖα δ᾽ ἄν ποιῇ τις, ἐν τῷ σώματι\apparat{154}{\lecon{ἄν} \PsiUn, R : ἢν Dindorf}\\
ἔνεσθ᾽ ὑπάρχον τοῦθ᾽. ἃ δ᾽ οὐ κεκτήμεθα,\\
μίμησις ἤδη ταῦτα συνθηρεύεται.\\\passeLigne}%
\end{groupe}%
%% v. 157-158
\begin{groupe}%
\personnageGrec{Μνησίλοχος}%
\trimIamb{ὅταν σατύρους τοίνυν ποιῇς, καλεῖν ἐμέ,
\commentaire{157}{ὅταν σατύρους τοίνυν ποιῇς.}{Lors des Grandes Dionysies, chaque poète tragique en lice devait présenter une tétralogie composée de trois tragédies et d'un drame satyrique.}\\
ἵνα συμποιῶ σοὔπισθεν ἐστυκὼς ἐγώ.
\commentaire{158}{συμποιῶ σοὔπισθεν ἐστυκὼς}{ ; ἐστυκὼς de στύω, \italique{bander}, alors qu'on attendait ἑστηκώς, \italique{se tenir debout} ; σοὔπισθεν est mis ici pour σοί ὄπισθέν (le datif σοί réclamé par le συμ- de συμποιῶ).}\\\passeLigne}%
\end{groupe}%
%% v. 159-167
\begin{groupe}%
%
\personnageGrec{Ἀγάθων}%
\trimIamb{ἄλλως τ᾽ ἄμουσόν ἐστι ποιητὴν ἰδεῖν%\testimonia{159-61}{Suidas α 1633 (κ' Ἀνακρ.)}
\\
ἀγρεῖον ὄντα καὶ δασύν σκέψαι δ᾽ ὅτι}%
\end{groupe}%
%% v. 159-167 (part 1)
\begin{groupe}\trimIamb{Ἴβυκος ἐκεῖνος κἀνακρέων ὁ Τήιος
\commentaire{161-3}{Ἴβυκος ἐκεῖνος κἀνακρέων ὁ Τήιος / κἀλκαῖος, οἳ περὶ ἁρμονίαν ἐχύμησαν / ἐμιτροφόρουν τε καὶ διεκίνων Ἰωνικῶς,}{Les trois poètes lyriques Ibycos de Rhégium, Anacréon de Téos et Alcée de Lesbos sont aussi rapprochés par le scholiaste de Pindare (\italique{Isthmiques}, 2.1b\marqueur{1}) pour d'autres raisons que d'ordre harmonique : ce sont trois auteurs pédérastiques ; Cicéron les regroupe aussi sur le critère de la passion amoureuse (\italique{Tusculanes}, IV, 33) :
\citComm{\latin{quid denique homines doctissimi et summi poëtae de se ipsis et carminibus edant et cantibus ? Fortis vir in sua re publica cognitus quae de iuvenum amore scribit Alcaeus ! nam Anacreontis quidem tota poësis est amatoria. Maxume vero omnium flagrasse amore Reginum Ibycum apparet ex scriptis.}\\
\francais{Avec quelle licence de très savants hommes et de grands poètes n'ont-ils pas chanté leurs galanteries ? Alcée, ce fameux guerrier de Mitylène, que n'a-t-il pas écrit de ces inclinations à la Grecque ? Anacréon respire-t-il autre chose que l'amour ? On voit la passion encore poussée bien plus loin dans les poésies d'Ibycus.}}
Le spectateur du \nieme{V} siècle ne s'y trompait pas et savait à quoi s'en tenir avec ces trois bougres-là et comment interpréter correctement les euphémismes agathoniens.
Le commentaire ancien le plus pertinent sur notre passage n'est pas écrit en marge du \italique{Ravennas} -- comme c'est hélas bien souvent le cas pour les \italique{Thesmophories} -- mais se trouve dans la \italique{Souda} (ε 989 -- et quasiment à l'identique dans le lexique de Zonaras, p. 707) :
\citComm{ἀντὶ τοῦ μίτραις\marqueur{2} τουτέστι ζώναις\marqueur{2} διαζωσάμενοι ἁβρῶς περιεπάτουν.\\
Signifie : ceints de mitres -- c'est-à-dire de ceintures -- ils vivaient de façon luxurieuse.}
Ce commentaire est doublement intéressant : d'abord, il précise que la μίτρα dont il est question est à prendre au sens de \italique{ceinture} et non de bandeau. Ensuite parce qu'il glose trois des mots du vers 163 : ἐμιτροφόρουν par μίτραις διαζωσάμενοι ; Ἰωνικῶς par ἁβρῶς\marqueur{3} ; et -- plus intéressant -- le \italique{contra metrum} †διεκίνων† (†διεκίνουν† dans la \italique{Souda}) par περιεπάτουν (ici avec le sens figuré de \italique{se conduire}, \italique{vivre de telle} ou \italique{telle manière}). 
Cette glose est donc incompatible avec l'émendation de Toup, très largement reprise dans les éditions modernes. Toup proposa la correction suivante : διεκλῶντ' (de διακλάω , \italique{briser en deux}, d'où \italique{briser une vie par la mollesse} (?!)).
Il est bien préférable -- tant du point de vue paléographique que du sens -- de lire διέκειντ', de διάκειμαι, \italique{être disposé}, synonyme de περιπατέω, et qui se construit fréquemment avec un adverbe (ici : διάκειμαι Ἰωνικῶς, \italique{avoir des dispositions à la vie ionienne}, c'est-à-dire à la mollesse et la luxure\marqueur{3}).\\
Les sous-entendus d'Agathon sont à peine voilés : les trois poètes pédérastes portent la ceinture à la mode ionienne, c'est-à-dire lâche, voire franchement dénouée...
La métaphore est la suivante : tels qu'ils portaient la ceinture, béante, telles étaient leurs \italique{harmonies}. Le vers 162, οἳ περὶ ἁρμονίαν †ἐχύμησαν†\marqueur{4}, se laisse maintenant comprendre. Il faut lire :
\citComm{οἳ περὶ ἁρμονίαν ἐκεχήνεσαν}
La correction est très légère : une nouvelle faute de iotacisme (entre η et υ), de confusion μ/ν (quasiment identiques sur notre \italique{Ravennas}, comme sur beaucoup de manuscrits écrits en minuscules) et d'une haplographie bien explicable :  $\text{εκεχ} \xrightarrow{\text{saut du même au m. }} \text{ εχ}$.\\
χάσκειν signifie littéralement \italique{avoir la bouche grande ouverte}\marqueur{5}. En bâillant de la sorte aux corneilles, l'on pouvait émettre de légers sons, comparables aux coquilles grillant sur des chardons\marqueur{6} ou à cette vieille lyre de Cratinos, le rival hors-d'âge d'Aristophane, chevrotant ses derniers et mauvais accords (\italique{Cavaliers}, v. 533)\marqueur{7} :
\citComm{τῶν θ' ἁρμονιῶν διαχασκουσῶν.}
(Noter la collocation entre (δια-)χάσκω et ἁρμονία.)  C'est la même expression qu'on retrouve dans la bouche d'Agathon : περὶ ἁρμονίαν ἐκεχήνεσαν, « béer une harmonie », c'est-à-dire bêler -- ou chevroter, telle est l'image en français, point très éloignée du grec -- quelques paroles en musique\marqueur{8}. Le sel de la métaphore réside dans le double sens de περιχάσκειν. Le verbe est justement attesté au sens d'\italique{ouvrir une ceinture}\marqueur{9}.\\
Il est remarquable de noter que la tradition millénaire du texte d'Aristophane ait conservé la leçon ἐχύμησαν avec un η jusqu'au scribe du \italique{Ravennas} ; et ce en dépit de manuscrits anciens fautifs portant ἐχύμισαν et de commentaires afférents (ἔγχυμον ἐποίησαν) -- ayant pénétré jusqu'aux marges du \italique{Ravennas} -- suggérant la présence d'un verbe *χυμίζω. Preuve encore une fois que l'œuvre conservée d'Aristophane est, comme l'avait déjà noté Willems, l'une des plus fidèlement transmises avec l'\italique{Iliade} et l'\italique{Odyssée}\marqueur{10}.
}
\debutNote{161-3}
\note{1}{Au tout début de la célèbre Isthmique où Pindare affirme qu'autrefois la Muse ne se faisait pas mercenaire mais chantait la beauté. Le scholiaste de noter que cela s'applique aux trois poètes lyriques : ταῦτα δὲ τείνει καὶ εἰς τοὺς περὶ Ἀλκαῖον καὶ Ἴβυκον καὶ Ἀνακρέοντα... ainsi qu'à tous les poètes apparemment fort portés sur la pédérastie : καὶ εἴ τινες τῶν πρὸ αὐτοῦ δοκοῦσι περὶ τὰ παιδικὰ ἠσχολῆσθαι· οὗτοι γὰρ παλαιότεροι Πινδάρου. Et d'appuyer le trait avec cette anecdote : on interrogeait Anacréon sur la raison qu'il avait d'écrire des hymnes à de jeunes adolescents plutôt qu'aux dieux. Anacréon répondit que c'était justement ces derniers leurs dieux : Ἀνακρέοντα γοῦν ἐρωτηθέντα, φασί, διατί οὐκ εἰς θεοὺς ἀλλ εἰς παῖδας γράφεις τοὺς ὕμνους; εἰπεῖν, ὅτι οὗτοι ἡμῶν θεοί εἰσιν.}
\note{2}{Zonaras -- ou le \italique{pseudo}-Zonaras pour les esprits chagrins -- comporte un accusatif en place de l'ablatif : ἀντὶ τοῦ μίτρας ἤτοι ζώνας διαζωσάμενοι.}
\note{3}{Voir Érasme, \italique{Adages} (n° 469, \latin{Risus Ionicus}).}
\note{4}{La rêverie du scholiaste reprise par Suidas, ἐχύμισαν· ἀντὶ τοῦ ἔγχυμον ἐποίησαν, n'est évidemment pas à prendre au sérieux. Pas plus que l'invention du verbe χυμίζω... à moins qu'on ne m'explique à quelle sauce (χυμός) tourner le passage pour qu'il ait un rapport avec les inflexions de voix (\italique{harmonies}). Ces divagations auront probablement troublé davantage encore les scribes, déjà en butte à un texte corrompu.}
\note{5}{Bien souvent : \italique{être béat d'admiration}. Voir Taillardat §§472 et 578.}
\note{6}{ἀνέχασκον εἷς ἕκαστος ἐμφερέστατα ὀπτωμέναις κόγχαισιν ἐπὶ τῶν ἀνθράκων. Aristophane, \italique{fr.} 67 K.-A.}
\note{7}{C'est bien ainsi qu'il faut interpréter τῶν θ' ἁρμονιῶν διαχασκουσῶν : « émettre des sons, ou accords » et non « aux joints béants » (en parlant de la lyre), comme on peut le lire.}
\note{8}{ἁρμονία a le sens d'arrangement de paroles pour un morceau de musique dans Platon, \italique{Theæt.} 175e-176a.}
\note{9}{Voir LSJ, I. b. Attesté par Héliodore le médecin (I\italique{er} ou \nieme{II} s. ap. J.-C). De longs fragments nous en ont été conservés par Oribase (\italique{Collectiones medicae} XLVIII, 58.5) : cf. les \italique{Œuvres d'Oribase}, t. 4, p. 320 (1862), de Bussemaker et Daremberg. À propos d'un patient à qui l'on applique un bandage, appelé ceinture : τῇ ζώνῃ \textbf{περιχάσκων} ζώννυται ὁ πάσχων, καὶ τότε ὁ εὐθὺς τελαμὼν ἀπὸ τῆς ὀσφύος ἄγεται ἐπὶ τὴν ἕδραν καὶ τὸν περίναιον, « On entoure le malade avec la ceinture, en l'\italique{ouvrant} largement ; puis on amène la bande droite des lombes au siège et au périnée » (trad. Daremberg).}
\note{10}{Voir l'introduction à ses premières notes sur les \italique{Cavaliers}. L'une des raisons en est que le Comique, au contraire des Tragiques, ne fut que peu rejoué. Les premières éditions critiques de ses œuvres furent sans doute beaucoup plus aisées à constituer que celles d'un Euripide, dont beaucoup plus de versions différentes des pièces circulaient. Victor Coulon, l'éditeur aux Belles Lettres d'Aristophane, abonda dans le sens de Willems : à l’opposé du courant majoritaire de son temps, il s'est astreint à conserver le plus possible les leçons des manuscrits contre les conjectures des philologues. Ce « conservatisme », qu'il défend pour le cas d'Aristophane dans son \italique{Essai sur la méthode de la critique conjecturale appliquée au texte d’Aristophane} (1933) lui valut d’être encore considéré à la fin des années 90 par Sommerstein comme le meilleur éditeur d’Aristophane.}
\\
κἀλκαῖος, οἳ περὶ ἁρμονίαν ἐκεχήνεσαν,
\apparat{162}{\lecon{κἀλκαῖος} R, Suid. : (κ)ἀχαιός Σ \espApp ἐχύμησαν R : ἐχύμισαν Suid. : \lecon{ἐκεχήνεσαν} Bibi \ligneApp{163} διεκίνων R : διεκίνουν Suid. ε 989, ι 495 : διεκλῶντ᾽ Toup : \lecon{διέκειντ'} Bibi}
\\
ἐμιτροφόρουν τε καὶ διέκειντ' Ἰωνικῶς,
\\
καὶ Φρύνιχος, τοῦτον γὰρ οὖν ἀκήκοας,
\commentaire{164}{Φρύνιχος, τοῦτον γὰρ οὖν ἀκήκοας.}{Le poète Phrynicos, qui vivait au temps des guerres médiques portait naturellement le riche costume de son temps. Cf. les \italique{Cavaliers}, v. 1331-1332 :
\citComm{Ἀλλαντοπώλης — ὅδ᾽ ἐκεῖνος ὁρᾶν τεττιγοφόρας, ἀρχαίῳ σχήματι λαμπρός,\\
οὐ χοιρινῶν ὄζων ἀλλὰ σπονδῶν, σμύρνῃ κατάλειπτος.\\
\francais{Le charcutier — Le voilà qui apparait, la cigale à la chevelure, dans tout l'éclat du costume antique, fleurant, non les coquilles mais les libations de paix, et frotté avec de la myrrhe.}}
Et Thucydide (I, 6) :
\citComm{καὶ οἱ πρεσβύτεροι αὐτοῖς τῶν εὐδαιμόνων \{διὰ τὸ ἁβροδίαιτον\} οὐ πολὺς χρόνος ἐπειδὴ χιτῶνάς τε λινοῦς ἐπαύσαντο φοροῦντες καὶ χρυσῶν τεττίγων ἐνέρσει κρωβύλον ἀναδούμενοι τῶν ἐν τῇ κεφαλῇ τριχῶν.\\
Les vieillards athéniens de la classe aisée portaient encore, il n'y a pas longtemps, des tuniques (longues) de lin, et piquaient des cigales d'or dans leurs cheveux, ramassés et noués sur le sommet de la tête.}
(Note de Willems.)\\
Le scholiaste explique l'incise d'Agathon (τοῦτον γὰρ οὖν ἀκήκοας) ainsi : ἐπεὶ Ἀθηναῖος ἦν.}\\\passeLigne}%
\end{groupe}%
\endnumbering
\end{Leftside}
\begin{Rightside}
\beginnumbering
%% v. 101-129
\begin{groupe}%
\chantAgathFr{\ligne{123}8° — J'honore Léto la souveraine}%
\end{groupe}%
%% v. 101-129
\begin{groupe}%
\chantAgathFr{et la cithare mère des hymnes,\testimonia{120}{\EtymMagnum 153.21 Ἀσιάδος κρούματα· τῆς κιθάρας. οὔτως Αριστοφάνης εἶπε, παρῳδῶν τὸ ἐξ Ἐρεχθέως Εὐριπίδου ; \SynB α 2226 $=$ Photius α 2957 $=$ Suidas 4147 $=$ \EtymGen (AB) α 1271 ed. Lasserre-Livadaras}}%
\end{groupe}%
%% v. 101-129
\begin{groupe}%
\chantAgathFr{à la voix mâle et éprouvée.\\}%
\end{groupe}%
%% v. 101-129
\begin{groupe}%
\chantAgathFr{9° — À ces accents, une lumière jaillit}%
\end{groupe}%
%% v. 101-129
\begin{groupe}%
\chantAgathFr{dans les yeux des dieux\testimonia{125}{(ἄρσενι βοᾷ δοκιμῳ) Suidas α 4016 $=$ Scholium}\\
\ligne{126}et traverse notre soudain regard.}%
\end{groupe}%
%% v. 101-129
\begin{groupe}%
\chantAgathFr{Aussi bien, glorifie le seigneur Phébus.\testimonia{128}{(ἄγαλλε) Photius α 164 $=$ \SynB α 145}\\}%
\end{groupe}%
%% v. 101-129
\begin{groupe}%
\chantAgathFr{10° — Salut, heureux fils de Léto.\\\passeLigne}%
\end{groupe}%
\begin{groupe}%
\personnageFr{Mnésiloque}%
\trimIambFr{Quel chant suave, ô Génétyllides vénérées\testimonia{130-2}{(ἡδὺ -- μανδαλωτόν) Suidas γ 141}\testimonia{130}{(Γενετυλλίδες) Hesychius γ 343 Γενετυλλίς}\\
et qui sent la féminie et les coups de langue\testimonia{131-2}{(θηλυδριῶδες -- μανδαλωτόν) Suidas μ 134}\testimonia{131}{(θηλυδριῶδες καὶ κατεγλωττισμένον) Suidas κ 912 \espTest Hesychius θ 491 θηλυδριῶδες \espTest (κατεγλωττισμένον) Pollux II.109}\\
et les baisers lascifs. En l'écoutant\testimonia{132}{Hesychius μ 224 μανδαλωτόν $\simeq$ Photius μ 86, cf. Photius γ 115}\testimonia{132-3}{(ὥστ᾽ -- γάργαλος) Suidas γ 83 (sed τοῦ μέλους post ἀκροωμένου addito)}
j'ai ressenti un prurit jusque sous le fondement.\testimonia{133}{(γάργαλος $=$ γαργαλισμός) Moeris γ 23, cf. \Phryn 56.9}
Jouvenceau, si tu en es un, à la façon d'Eschyle\testimonia{134}{Hesychius η 948 ἥτις· ὅστις (?)}
dans sa \italique{Lycurgie}, je veux t'interroger.
D'où es-tu, ô femmelette ? Quelle est ta patrie ? Qu'est cet accoutrement,\testimonia{136}{(ποδαπὸς -- πάτρα) Suidas γ 504 \espTest (γύννις) Hesychius γ 1015 γύννις, cf. \Phryn 60.19}\testimonia{136-40}{(τίς ἡ στολή -- κοινωνία) Suidas β 110 (δὲ pro δαὶ)}}
\end{groupe}%
%% v. 130-145
\begin{groupe}%
\trimIambFr{cette incohérence dans tes façons ? Un barbitos,\testimonia{137}{(βάρβιτος) Hesychius β 211 βάρβιτος, cf. Photius β 62}
qu'a-t-il à dire à une crocote, une lyre à une résille ?
Quoi ! une pomme d'Adam avec un strophion ? Quelle  disconvenance !
Le rapport d'un miroir à une épée ?\\
Et toi-même, ô jeunesse, qui es-tu ? Un homme ?\testimonia{141-2}{(πότερον -- πέος) Suidas π 987}
Mais où est ton cas ? où ton manteau ? où tes laconiennes ?\testimonia{142}{(ποῦ τὸ πέος -- Λακωνικαί) Suidas λ 63, π 987 (τὸ πέος pro πέος)\espTest Moeris χ 19 χλαῖνα \espTest (Λακωνικαί) Hesychius λ 225 Λακωνικαί ; Photius λ 57 Λακωνικαί ; scholium \italique{Vespae} 1158-60}\\
Alors tu es femme ? Mais où sont tes seins ?
Que dis-tu ? Pourquoi te taire ? Alors d'après ton chant,
je te déchiffre puisque tu ne veux pas t'expliquer toi-même.\\\passeLigne}%
\end{groupe}%
%% v. 146-152
\begin{groupe}%
\personnageFr{Agathon}%
\trimIambFr{Vieillard, vieillard, ces reproches dictés par l'envie,\testimonia{146-7}{Suidas α 1097 (μὲν omisso et δὲ pro δ')}\\
je les ai entendus, mais n'en ai point senti l'atteinte.\\
Je porte la mise congruente à mes pensées.\\
Car c'est une loi que le poète, suivant les drames\testimonia{149-50}{\Apostol XVIII 41h}\\
qu'il lui faut faire, modèle sur eux ses mœurs.\\
Et par exemple, si l'on compose des drames féminins,\\
il faut que le corps participe aux mœurs de la femme..\\\passeLigne}%
\end{groupe}%
%% v. 153
\begin{groupe}%
\personnageFr{Mnésiloque}%
\trimIambFr{Ainsi tu chevauches quand tu composes une Phèdre ?\\\passeLigne}%
\end{groupe}%
%% v. 154-156
\begin{groupe}%
\personnageFr{Agathon}%
\trimIambFr{Compose-t-on des drames mâles, dans le corps
réside ce qui est requis. Mais ce que nous ne possédons pas,\\
l'imitation est là qui l'attrape.\\\passeLigne}%
\end{groupe}%
%% v. 157-158
\begin{groupe}%
\personnageFr{Mnésiloque}%
\trimIambFr{Eh bien, quand tu feras des drames satyriques, appelle-moi,
et je coopérerai avec toi, de derrière, avec raideur.\\\passeLigne}%
\end{groupe}%
%% v. 159-167
\begin{groupe}%
\personnageFr{Agathon}%
\trimIambFr{D'ailleurs, c'est chose choquante à voir \testimonia{159-61}{Suidas α 1633 (κ' Ἀνακρ.)}\\qu'un poète
grossier et hirsute. Avise-toi que}%
\end{groupe}%
%% v. 159-167 (part 1)
\begin{groupe}\trimIambFr{le fameux Ibycos, et Anacréon de Téos\testimonia{161-2}{Suidas ε 4049 (Ἀνακρέων, Ἀλκαῖος, οἵπερ ἁρμονίαν ἐχύμισαν)}\\
et Alcée, accordaient leur voix,\testimonia{162-3}{(κ' Ἀλκαῖος -- Ἰωνικῶς) Suidas ε 989 (sed οἵπερ ἁρμονίαν ἐχύμισαν) $\simeq$ Zonaras 707 ἐμιτροφόρουν}\testimonia{162}{Zonaras 942 οἵπερ ἁρμονίαν ἐχύμισαν \espTest (ἐχύμισαν) Hesychius ε 7681 ἐχύμισαν· ἤρτυσαν}\\
leur mitre et leur vit à l'ionienne. Et Phrynichos,\testimonia{163}{Suidas ι 495 ἐμιτροφόρουν τε καὶ διεκίνουν Ἰωνικῶς}\\
— celui-là, en effet, tu l'as entendu, —\\\passeLigne}%
\end{groupe}%
\endnumbering
\end{Rightside}
\end{pages}
\Pages
\clearpage\ifodd\value{page}\hbox{}\newpage\fi
\enTeteG{ΠΡΟΛΟΓΟΣ --- Codex Ravennas 429, f. 164v (Classense 166v) --- vv. 165–207}
\begin{adjustwidth}{-1.5cm}{-3.0cm}
\begin{adjustbox}{max totalheight=21.0cm, max width=16.5cm, keepaspectratio}
\begin{minipage}{16.5cm}
{\fontsize{7.2pt}{8.5pt}\selectfont
\noindent\textbf{Folio 05 --- Codex Ravennas 429, \texttt{f. 164v} (\texttt{Classense 166v}) : ΠΡΟΛΟΓΟΣ, vv. 165–207}\hfill \textit{Témoins :} $\text{R}^1, \Rpc, \Rdeux$ (\Mon)\par\vspace{2pt}
\hrule height 0.4pt\vspace{2.5pt}
}
{\fontsize{8.1pt}{9.5pt}\selectfont
\ligneRav{165}{ἀγάθ/}{αὐτός τε καλὸς ἦν καὶ καλῶς ἠμπέσχετο}{\textbf{165} \lecon{ἠμπέσχετο} R : ἠμπισχετο Elmsley ad Eur. \italique{Med.} 1128 (\italique{induebatur} Divus)}
\ligneRav{}{}{διὰ τοῦτ᾽ ἄρ᾽ αὐτοῦ καὶ κάλ᾽ ἦν τὰ δράματα·}{}
\ligneRav{}{}{ὅμοια γὰρ ποιεῖν ἀνάγκη τῇ φύσει·}{}
\ligneRav{168}{μ\auDessus{ν}{Η}}{ταῦτ᾽ ἄρ᾽ ὁ \difference{Φιλοκλέης} αἰσχρὸς ὢν αἰσχρῶς ποιεῖ,}{\textbf{168-9} Φιλοκλῆς, Ξενοκλῆς R : \lecon{Φιλοκλέης}, \lecon{Ξενοκλέης} Bentley \espApp \textbf{169} \lecon{ὁ δ᾽ αὖ} R : ὁ δὲ Brunck}
\ligneRav{}{}{ὁ δ᾽ αὖ Ξενοκλέης ὢν κακὸς κακῶς ποιεῖ,}{}
\ligneRav{170}{}{ὁ δ᾽ αὖ Θέογνις ψυχρὸς ὢν ψυχρῶς ποιεῖ·}{}
\ligneRav{}{ἀγάθ/}{ἅπασ᾽ ἀνάγκη ταῦτα γάρ τοι γνοὺς ἐγὼ}{}
\ligneRav{}{—}{ἐμαυτὸν ἐθεράπευσα. : πῶς πρὸς τῶν θεῶν;}{}
\ligneRav{173}{\auDessus{εὐ}{ρ}}{παῦσαι βαΰζων καὶ γὰρ ἐγὼ τοιοῦτος ἦν}{\textbf{173} ἦν R, Suid. β 192 : \lecon{ἦ} Dindorf (1825)}
\ligneRav{}{}{ὢν τηλικοῦτος, ἡνίκ᾽ ἠρχόμην ποιεῖν·}{}
\ligneRav{175}{μ\auDessus{ν}{Η}}{μὰ τὸν Δί᾽ οὐ ζηλῶ σε τῆς παιδεύσεως·}{}
\ligneRav{176}{\auDessus{εὐ}{ρ}}{ἀλλ᾽ ὧνπερ οὕνεκ᾽ ἦλθον, ἔα μ᾽ εἰπεῖν. : λέγε·}{\textbf{176} \lecon{(Ἀγάθ.)} λέγε R, Fritzsche : (Μνη.) Küster}
\ligneRav{}{\auDessus{εὐ}{ρ}}{Ἀγάθων, σοφοῦ πρὸς ἀνδρός, ὅστις ἐν βραχεῖ}{}
\ligneRav{178}{}{πολλοὺς καλῶς \difference{οἷός} τε συντέμνειν λόγους·}{\textbf{178} οἷόν R : \lecon{οἷός} Suid. σ 792, οι 154, Stob. III,35.3 \espApp \lecon{συντέμνειν} R, Suid. οι 154, Stob. : συντεμεῖν Suid. σ 792 \espApp \textbf{179} κοινῇ R, Suid. κ 2556, σ 792 : \lecon{καινῇ} Biset ap. Σ}
\ligneRav{179}{}{ἐγὼ δὲ \difference{καινῇ} ξυμφορᾷ πεπληγμένος}{\textbf{179} κοινῇ R, Souda $\kappa\ 2556, \sigma\ 792$ (confusion minuscule $\alpha \to \text{ο}$) : \lecon{καινῇ} Biset d'après $\Sigma^{\text{R}}$}
\ligneRav{180}{—}{ἱκέτης ἀφῖγμαι πρὸς σέ. : τοῦ χρείαν ἔχων;}{}
\ligneRav{}{\auDessus{εὐ}{ρ}}{μέλλουσί μ᾽ αἱ γυναῖκες ἀπολεῖν τήμερον}{}
\ligneRav{}{}{τοῖς Θεσμοφορίοις, ὅτι κακῶς αὐτὰς λέγω·}{}
\ligneRav{}{ἀγάθ/}{τίς οὖν παρ᾽ ἡμῶν ἐστιν ὠφέλειά σοι;}{}
\ligneRav{184}{\auDessus{εὐ}{ρ}}{\difference{ἡ πᾶσ᾽·} ἐὰν γὰρ ἐγκαθεζόμενος λάθρᾳ}{\textbf{184} ἡ πᾶσα R : \lecon{ἡ πᾶσ᾽·} Küster : ἅπασ' Scaliger}
\ligneRav{185}{}{ἐν ταῖς γυναιξίν, ὡς δοκῶν εἶναι γυνή,}{}
\ligneRav{}{}{ὑπεραποκρίνῃ μου, σαφῶς σώσεις ἐμέ·}{}
\ligneRav{}{}{μόνος γὰρ ἂν λέξειας ἀξίως ἐμοῦ·}{}
\ligneRav{}{ἀγάθ/}{ἔπειτα πῶς οὐκ αὐτὸς ἀπολογεῖ παρών;}{}
\ligneRav{}{\auDessus{εὐ}{ρ}}{ἐγὼ φράσω σοι. πρῶτα μὲν γιγνώσκομαι}{}
\ligneRav{190}{}{ἔπειτα πολιός εἰμι καὶ πώγων᾽ ἔχω,}{}
\ligneRav{}{}{σὺ δ᾽ εὐπρόσωπος λευκὸς ἐξυρημένος}{}
\ligneRav{}{}{γυναικόφωνος ἁπαλὸς εὐπρεπὴς ἰδεῖν·}{}
\ligneRav{}{ἀγάθ/}{Εὐριπίδη— : τί ἔστιν; : ἐποίησάς ποτε,}{}
\ligneRav{}{}{`χαίρεις ὁρῶν φῶς, πατέρα δ᾽ οὐ χαίρειν δοκεῖς;'}{}
\ligneRav{195}{\auDessus{εὐ}{ρ}}{ἔγωγε. : μή νυν ἐλπίσῃς τὸ σὸν κακὸν}{}
\ligneRav{196}{}{ἡμᾶς ὑφέξειν. καὶ \difference{γὰρ ἂν} μαινοίμεθ᾽ ἂν·}{\textbf{196} γὰρ R : \lecon{γὰρ ἂν} Suid. υ 709}
\ligneRav{}{}{ἀλλ᾽ αὐτὸς ὅ γε σόν ἐστιν οἰκείως φέρε·}{}
\ligneRav{}{}{τὰς συμφορὰς γὰρ οὐχὶ τοῖς τεχνάσμασιν}{}
\ligneRav{}{}{φέρειν δίκαιον ἀλλὰ τοῖς παθήμασιν·}{}
\ligneRav{200}{μ\auDessus{ν}{Η}}{καὶ μὴν σύ γ᾽ ὦ κατάπυγον εὐρύπρωκτος εἶ}{}
\ligneRav{}{}{οὐ τοῖς λόγοισιν ἀλλὰ τοῖς παθήμασιν·}{}
\ligneRav{}{\auDessus{εὐ}{ρ}}{τί δ᾽ ἔστιν ὅτι δέδοικας ἐλθεῖν αὐτόσε;}{}
\ligneRav{}{ἀγάθ/}{κάκιον ἀπολοίμην ἂν ἢ σύ. : πῶς; : ὅπως;}{}
\ligneRav{204}{}{δοκῶν γυναικῶν ἔργα νυκτερείσια}{\textbf{204} \lecon{νυκτερείσια} R : νυκτερήσια Bothe}
\ligneRav{205}{}{κλέπτειν ὑφαρπάζειν τε θήλειαν Κύπριν·}{}
\ligneRav{206}{μ\auDessus{ν}{Η}}{ἰδού γε κλέπτειν· νὴ Δία βινεῖσθαι μὲν οὖν·}{\textbf{206} \lecon{νὴ Δία} R : μà Δία Fritzsche}
\ligneRav{}{}{ἀτὰρ ἡ πρόφασίς γε νὴ Δί᾽ εἰκότως ἔχει·}{}
}
\vspace{3pt}\noindent\rule{4.5cm}{0.4pt}\par\vspace{1.5pt}
{\fontsize{7.0pt}{8.3pt}\selectfont
\noindent\textbf{Loci paralleli (vv. 165–207) ---} \textbf{168} (ὁ Φιλοκλέης αἰσχρὸς $=$ κορυδός) scholium \italique{Aves} 1295a.α \espTest \textbf{173} (βαύ̈ζων -- ἦν) Suidas β 192 $\simeq$ Zonaras 379 (sed γὰρ καὶ pro καὶ γὰρ) \espTest \textbf{175} (οὐ -- παιδεύσεως) Suidas ζ 63 \espTest \textbf{176} Scholium v. 500 (sed ἀλλ' εἰπεῖν βούλομαι pro ἔα μ' εἰπεῖν) \espTest \textbf{177-80} (σοφοῦ -- πρὸς σέ) Suidas σ 792 \espTest \textbf{177-8} (σοφοῦ -- λόγους) Suidas οι 154 \espTest \textbf{179} Suidas κ 2556 \espTest \textbf{192} (γυναικόφωνος) Pollux II 111 \espTest \textbf{193} (ἐποίησάς ποτε) \Eustathe 4.9 \espTest \textbf{195-7} (μή νυν -- φέρε) Suidas υ 709 (νῦν pro νυν et ἀλλ' omisso) \espTest \textbf{198-201} (τὰς συμφορὰς -- παθήμασιν) Suidas τ 440 (τεχνάσμασι et παθήμασι pro -σιν) \espTest \textbf{198-201} (τὰς συμφορὰς -- παθήμασιν) Suidas τ 440 (-σι pro -σιν) %(τεχνάσμασι et παθήμασι pro -σιν)\par
}
\vspace{2.5pt}\noindent\rule{4.5cm}{0.4pt}\par\vspace{1.5pt}
{\fontsize{6.9pt}{8.2pt}\selectfont
\noindent\textbf{Suite (en commentaire) de l'apparat critique \& notes paléographiques (\texttt{f. 164v}) ---} \textbf{Contexte :} Fin de la tirade d'Agathon (vv. 165–172), requête d'Euripide priant Agathon d'aller plaider sa cause aux Thesmophories et refus prudent d'Agathon (vv. 173–207). --- Page très régulière de 43 lignes. On relève les formes contractes non métriques de noms propres tragiques (Φιλοκλῆς, Ξενοκλῆς R aux vv. 168–169 pour -κλέης), la confusion minuscule $\alpha/\text{ο}$ dans κοινῇ (v. 179 R et Souda) pour καινῇ (Biset d'après $\Sigma^{\text{R}}$), et la graphie iotacisante νυκτερείσια (v. 204 R, maintenue par L. Hauseux). --- \textbf{v. 179 :} κοινῇ R, Souda $\kappa\ 2556, \sigma\ 792$ (confusion minuscule $\alpha \to \text{ο}$) : \lecon{καινῇ} Biset d'après $\Sigma^{\text{R}}$.\par
}
\end{minipage}
\end{adjustbox}
\end{adjustwidth}
\newpage
\enTeteD{Codex Ravennas 429 (Biblioteca Classense) --- Fac-similé du f. 164v (vv. 165–207)}
\begin{adjustwidth}{-1.5cm}{-3.0cm}
\begin{center}
\includegraphics[width=16.5cm,height=20.5cm,keepaspectratio]{/workspaces/Thesmophories/dossier_ravennas_thesmophories/folios_individuels_ff162v-177r/05_f164v_classense166v_vv0165-0207.jpg}\par
\vspace{2pt}
{\scriptsize\textbf{Codex Ravennas 429, f. 164v (Classense 166v)} --- ΠΡΟΛΟΓΟΣ, \textbf{vv. 165–207} ($\text{R}^1, \Rpc, \Rdeux$ (\Mon)).}
\end{center}
\end{adjustwidth}
\newpage
\clearpage\ifodd\value{page}\hbox{}\newpage\fi
\enTeteG{ΘΕΣΜΟΦΟΡΙΑΖΟΥΣΑΙ --- Texte grec établi (f. 164v, vv. 165–207)}
\enTeteD{LES THESMOPHORIES --- Traduction française (f. 164v, vv. 165–207)}
\begin{pages}
\begin{Leftside}
\beginnumbering
%% v. 165+ (part 2)
\begin{groupe}
\trimIamb{\ligne{165}αὐτός τε καλὸς ἦν καὶ καλῶς ἠμπέσχετο\\%\apparat{165}{\lecon{ἠμπέσχετο} R : ἠμπισχετο Elmsley ad Eur. \italique{Med.} 1128 (\italique{induebatur} Divus)}\\
διὰ τοῦτ᾽ ἄρ᾽ αὐτοῦ καὶ κάλ᾽ ἦν τὰ δράματα.\\
ὅμοια γὰρ ποιεῖν ἀνάγκη τῇ φύσει.\\\passeLigne}%
\end{groupe}%
%% v. 168-170
\begin{groupe}%
\personnageGrec{Μνησίλοχος}%
\trimIamb{ταῦτ᾽ ἄρ᾽ ὁ Φιλοκλέης αἰσχρὸς ὢν αἰσχρῶς ποιεῖ,%\testimonia{168}{(ὁ Φιλοκλέης αἰσχρὸς $=$ κορυδός) scholium \italique{Aves} 1295a.α}
\apparat{168-9}{Φιλοκλῆς, Ξενοκλῆς R : \lecon{Φιλοκλέης}, \lecon{Ξενοκλέης} Bentley \espApp \textbf{169} \lecon{ὁ δ᾽ αὖ} R : ὁ δὲ Brunck}\\
%
ὁ δ᾽ αὖ Ξενοκλέης ὢν κακὸς κακῶς ποιεῖ,\\
%
%\lednopb
ὁ δ᾽ αὖ Θέογνις ψυχρὸς ὢν ψυχρῶς ποιεῖ.
\commentaire{168-70}{Φιλοκλῆς, Ξενοκλῆς, Θέογνις.}{Trois poètereaux tragiques contemporains d'Aristophane.\\
À propos du vilain Philoclès, deux scholies nous ont été transmises dans le manuscrit de Ravenne\marqueur{1} :%
\citComm{Φιλοκλέης: τραγῳδίας ποιητὴς, αἰσχρὸς τὸ πρόσωπον καὶ σιμὸς καὶ μικρὸς τὸ σῶμα\marqueur{2}.\\
ἐκ τούτου δηλονότι καὶ ἐν τοῖς Ὄρνισιν (v. 1295) οὐκ ἄλλως κορυδῷ αὐτὸν ὡ[μοίου] ἀλλ᾽ ἢ ὅτι μικρὸς καὶ αἰσχρός. καὶ Τηλεκλείδης Ἡσιόδοις\marqueur{3}·
`ἀλλ᾽ ἡ τάλαινα Φιλοκλέα βδε\douteux{λ}\douteux{λ}...\douteux{λ}.οθεν οὖν·
εἰ δ᾽ ἐστὶν Αἰσχύλου φρόνημ᾽ ἔχων'.}%
Le scholiaste ne se trompe pas lorsqu'il prétend que la comparaison avec l'\italique{alouette huppée}, κορυδός, était d'ordre physique\marqueur{4} : on connaît par les parémiographes\marqueur{5} le proverbe grec «~être plus hideux que Corydée~» :%
\citComm{Κορυδέως εἰδεχθέστερος· ἐπὶ δυσμορφίᾳ ἐκωμῳδεῖτο καὶ οὗτος καὶ οἱ παῖδες αὐτοῦ ὧν ἕνα τὸν Ἀρχέστρατον φασιν.}%
La pauvre alouette huppée n'est pas bien lotie par la tradition populaire : en plus d'être laide, même dans la force de l'âge elle est chétive (plus que l'aigle vieillissant : ἀετοῦ γῆρας, κορύδου νεότης), et n'a pas pour autant le ramage du rossignol (ἐν ἀμούσοις καὶ κόρυδος φθέγγεται\marqueur{6})... « aux pays des aveugles, les borgnes sont rois ».\\
(Voir le commentaire au vers 441 pour le cas Xénoclès.)\\
Jean Taillardat a consacré la note que voici à Théognis\marqueur{7} :\\
Théognis le Tragique est ennuyeux et \italique{froid}, ψυχρός, au point qu'on le surnomme la \italique{Neige} (Χιών, Σ \italique{Acharniens} 11). La représentation d'une de ses pièces suffit à déchaîner les rigueurs de l'hiver, même sous le ciel déjà peu clément de la Thrace (\italique{Ach.} 138-140). Pris de pitié pour un froid rimailleur, Pisthétaire lui offre un vêtement de peau (\italique{Oiseaux} 935) : 
\citComm{ἔχε τὴν σπολάδα· πάντως δέ μοι ῥιγῶν δοκεῖς.\\
\francais{Prends ce manteau ; tu me parais complètement gelé.}}
Cette plaisanterie est continuée aux vers 946 et 955 (ταυταγὶ | τὰ κρυερὰ, \italique{ces vers glacés} ; cf. νιφοβόλους ἀναβολάς (\italique{ib.} 1385), « des préludes battus de neige ». Cet emploi de ψυχρός est banal\marqueur{8}.)}
\debutNote{168-70}
\note{1}{Manuscrit qui possède l'unique exemplaire des \italique{Thesmophories} qui nous soit parvenu -- \italique{unicodex} dit-on. Le manuscrit de Munich (\nieme{XV} siècle) n'étant qu'une recopie du premier (voir l'édition de von Velsen des \italique{Aristophanis Thesmophoriazusae} (1883) pour une comparaison complète du \italique{Ravennas} (R) et du \italique{Monacensis} (G)).}%
\note{2}{Suidas (φ 379) reprend cette scholie qu'il joint à d'autres éléments tirés du scholiaste aux \italique{Oiseaux} (v. 281).}%
\note{3}{fr. 15 in \italique{Poetae Comici Graeci}, VII, edd. Kassel et Austin.}%
\note{4}{Ce qui explique le surnom du \italique{Coryde} dont il est question dans les \italique{Deipnosophistes} d'Athénée (VI, 241a \&c.) «~qui aimait à se faire persifler~».}%
\note{5}{Zénobius IV, 59 ; Diogenianus V, 57 ; Michel Apostolius IX, 92 ; \&c. à retrouver dans le \italique{Corpus paroemiographorum graecorum}, ed. Leutsch. Cf. aussi Hésychius (σ 3681). Kassel \& Austin l'ont inclus (n° 827) aux \italique{Comica Adespota} de leur édition des \italique{Poetae Comici Graeci}, VIII.}%
\note{6}{Voir ce qu'en dit Érasme, qui a recensé ce proverbe dans ses \italique{Adages} (n° 1192).}%
\note{7}{\italique{Les images d'Aristophane} (1965), §759.}%
\note{8}{Ψυχρός en parlant d'un \italique{poète} : Alexis 179 [$=$ 184 Kassel-Austin $=$ Athénée, III, 123e], Machon (\italique{ap.} Athénée 580a), d'un \italique{orateur} : Plat. \italique{Euthyd.} 284e ; Théophil. 4, 4, \&c. Ψυχρός se dit aussi d'une tragédie manquée (\italique{Thesm.} 848), d'une plaisanterie sans sel (Eupol. 244 [$=$ 261 Kassel-Austin]).}\\\passeLigne}%
\end{groupe}%
%% v. 171-172a
\begin{groupe}%
\personnageGrec{Ἀγάθων}%
\trimIamb{ἅπασ᾽ ἀνάγκη ταῦτα γάρ τοι γνοὺς ἐγὼ\\
ἐμαυτὸν ἐθεράπευσα.\\\passeLigne}%
\end{groupe}%
%% v. 172b
\begin{groupe}%
\personnageGrec{Μνησίλοχος}%
\trimIamb{\ligne{172}\decalage{ἐμαυτὸν ἐθεράπευσα.}πῶς πρὸς τῶν θεῶν;\\\passeLigne}%
\end{groupe}%
%% v. 173-174
\begin{groupe}%
\personnageGrec{Εὐριπίδης}%
\trimIamb{παῦσαι βαΰζων καὶ γὰρ ἐγὼ τοιοῦτος ἦν%\testimonia{173}{(βαύ̈ζων -- ἦν) Suidas β 192 $\simeq$ Zonaras 379 (sed γὰρ καὶ pro καὶ γὰρ)}
\apparat{173}{ἦν R, Suid. β 192 : \lecon{ἦ} Dindorf (1825)}%
\commentaire{172-3}{πῶς πρὸς τῶν θεῶν; — παῦσαι βαΰζων.}{Goûtons à l'allitération en `t', `p' et surtout en `o' aboyée par Mnésiloque : π\textbf{ῶ}ς πρ\textbf{ὸ}ς τ\textbf{ῶ}ν θε\textbf{ῶ}ν ; c'est cette dissonance que lui reproche Euripide – pour la forme, car la question de Mnésiloque est aussi tout-à-fait désobligeante et ironique... il est certainement à l'affût de quelque double sens au «~ἐμαυτὸν ἐθεράπευσα~» d'Agathon. 
Euripide interrompt donc : παῦσαι βαΰζων, « cesse d'aboyer ! » -- quoique le verbe «~aboyer~» ne rende pas le registre tragique du verbe βαΰζω, formé par onomatopée sur le cri βαῢ βαΰ du chien (αὗ αὗ au vers 903 des \italique{Guêpes}).}\\
ὢν τηλικοῦτος, ἡνίκ᾽ ἠρχόμην ποιεῖν.\\\passeLigne}%
\end{groupe}%
%% v. 175
\begin{groupe}%
\personnageGrec{Μνησίλοχος}%
\trimIamb{μὰ τὸν Δί᾽ οὐ ζηλῶ σε τῆς παιδεύσεως.%\testimonia{175}{(οὐ -- παιδεύσεως) Suidas ζ 63}%
\commentaire{173-5}{καὶ γὰρ ἐγὼ τοιοῦτος ἦν ... — οὐ ζηλῶ σε τῆς παιδεύσεως.}{Euripide se reconnaît – non sans quelque condescendance – une parenté artistique\marqueur{1} quand Mnésiloque répond comme si Agathon était héritier des mœurs d'Euripide.}%
\debutNote{173-5}%
\note{1}{Confirmée par Aristote. (Voir le commentaire \italique{infra} au vers 29, n. 3.)}\\\passeLigne}%
\end{groupe}%
%% v. 176a
\begin{groupe}%
\personnageGrec{Εὐριπίδης}%
\trimIamb{ἀλλ᾽ ὧνπερ οὕνεκ᾽ ἦλθον, ἔα μ᾽ εἰπεῖν.%\testimonia{176}{Scholium v. 500 (sed ἀλλ' εἰπεῖν βούλομαι pro ἔα μ' εἰπεῖν)}
\\\passeLigne}%
\end{groupe}%
%% v. 176b
\begin{groupe}%
\personnageGrec{Ἀγάθων}%
\trimIamb{\ligne{176}\decalage{ἀλλ᾽ ὧνπερ οὕνεκ᾽ ἦλθον, ἔα μ᾽ εἰπεῖν.}λέγε.
\apparat{176}{\lecon{(Ἀγάθ.)} λέγε R, Fritzsche : (Μνη.) Küster}%
\commentaire{176}{λέγε.}{\latin{Nam neque Euripides neque Agatho Mnesilochi jocos solent curare. 
Apparet vero, orationem sic cohaerere},%
\citComm{\grec{παῦσαι βαΰζων — άλλ' ... ἔα μ᾽ εἰπεῖν}\\
\latin{desine me maledictis insequi, sed potius ea me patere loqui.}}%
\latin{Jam necessario is, quis rogatus est, respondere debet, estque certo verbum λέγε dic igitur Agathonis.} (Note de Fritzsche.)}\\\passeLigne}%
\end{groupe}%
%% v. 177-180a
\begin{groupe}%
%
\personnageGrec{Εὐριπίδης}%
\trimIamb{Ἀγάθων, σοφοῦ πρὸς ἀνδρός, ὅστις ἐν βραχεῖ%\testimonia{177-80}{(σοφοῦ -- πρὸς σέ) Suidas σ 792}\testimonia{177-8}{(σοφοῦ -- λόγους) Suidas οι 154}
\\
πολλοὺς καλῶς οἷός τε συντέμνειν λόγους.
\apparat{178}{οἷόν R : \lecon{οἷός} Suid. σ 792, οι 154, Stob. III,35.3 \espApp \lecon{συντέμνειν} R, Suid. οι 154, Stob. : συντεμεῖν Suid. σ 792 \espApp \textbf{179} κοινῇ R, Suid. κ 2556, σ 792 : \lecon{καινῇ} Biset ap. Σ}%
\commentaire{177-8}{Ἀγάθων σοφοῦ πρὸς ἀνδρός ὅστις ἐν βραχεῖ / πολλοὺς καλῶς οἷός τε συντέμνειν λόγους.}{Citation quasi-textuelle de l'\italique{Éole} (Αἰόλος) d'Euripide rapportée par Stobée (3, 35, 3)\marqueur{1} :\citComm{παῖδες, σοφοῦ πρὸς ἀνδρός ὅστις ἐν βραχεῖ,\\
πολλοὺς λόγους οἷός τε συντέμνειν καλῶς}%
P.-L. Courier, \italique{Pamphlet des pamphlets} :
\citComm{\francais{C'est un mérite non commun, ni facile, de clore en peu de mots beaucoup de sens.}%
(Note de Willems.)}}%
\debutNote{177-8}
\note{1}{À retrouver dans l'édition de C. Wachsmuth \& O. Hense, \italique{Ioannis Stobaei Anthologium}, vol. 3, 1904, p. 688 l. 5-7. C'est le fragment 21 de l'édition Jouan-van Looy aux Belles Lettres, t. VIII, 1\italique{ère} partie ; il correspond respectivement aux fragments 28 et 33 des éditions Kannicht et Mette.}%
\\
%\lednopb%
ἐγὼ δὲ καινῇ ξυμφορᾷ πεπληγμένος%\testimonia{179}{Suidas κ 2556}%
\commentaire{179}{ἐγὼ δὲ κοινῇ ξυμφορᾷ πεπληγμένος.}{L'émendation de Biset : substituer καινῇ à κοινῇ est reprise par quasiment tous les éditeurs. Elle a l'avantage de concorder avec la scholie : μεγάλῃ καὶ θαυμαστῇ (καινός désigne au propre ce qui est \italique{nouveau}, différent de ce qui s'est produit jusqu'à présent, et par suite, ce qui n'était pas attendu, \italique{étrange}, \italique{extraordinaire} – avec une évolution toute identique à notre \italique{in-ouï}). 
Nous sommes devant l'un des cas\marqueur{1} où la scholie rend compte d'un texte en meilleur état que celui auquel elle est attenante. (À noter qu'Euripide emploie l'expression καινάς συμφοράς aux vers 1226-7 de l'\italique{Oreste}.)\\
Cette glose s'est retrouvée dans la \italique{Souda} avec la mauvaise vedette (à l'article Κοινολογία). Pis : Suidas -- ou l'un des rédacteurs de la \italique{Souda} si l'on préfère --, voyant le problème aura ajouté aux deux termes dont le sens est incompatible avec κοινῇ un tiers, qui, lui, l'est, ἀνθρωπίνῃ.
Il écrit : κοινῇ δὲ ἀντὶ τοῦ ἀνθρωπίνῃ, μεγάλῃ καὶ θαυμαστῇ... Ce sont là des méthodes bien peu rigoureuses qui doivent nous le faire tenir en suspicion\marqueur{2}.%
%La \italique{Souda} va à l'encontre de cette émendation naturelle : à l'article Κοινολογία, elle définit ainsi κοινῇ· ἀνθρωπίνῃ, μεγάλῃ καὶ θαυμαστῇ -- reprenant le commentaire du \italique{Ravennas} et ajoutant un tiers terme, θαυμαστῇ, compatible avec une définition de κοινῇ.
%Suidas -- ou l'un des rédacteurs de la \italique{Souda} si l'on préfère --, disposant de la scholie et du texte fautif, aura ajouté le sens de κοινῇ par lui connu : ἀνθρωπίνῃ\marqueur{2}. 
%Pour ne pas laisser de doute quant au sens de l'épithète de ξυμφορᾷ, reportons-nous au style du véritable Euripide (\italique{Alceste}, 404-5) : ὥστ᾽ ἐγὼ / καὶ σφὼ \textbf{βαρείᾳ} συμφορᾷ πεπλήγμεθα et plus loin (v. 856) καίπερ \textbf{βαρείᾳ} συμφορᾷ πεπληγμένος. 
}%
\debutNote{179}%
\note{1}{Voir pour un autre exemple le commentaire au vers 1174.}%
\note{2}{Voir \italique{infra} le commentaire au vers 735 pour une erreur analogue de Suidas.}\\
ἱκέτης ἀφῖγμαι πρὸς σέ.\\\passeLigne}%
\end{groupe}%
%% v. 180b
\begin{groupe}%
\personnageGrec{Ἀγάθων}%
\trimIamb{\ligne{180}\decalage{ἱκέτης ἀφῖγμαι πρὸς σέ.}τοῦ χρείαν ἔχων;
\commentaire{180}{τοῦ χρείαν ἔχων; }{Euripide, \italique{Suppliantes}, 115 : Τί χρῆμα θηρῶν καὶ τίνος χρείαν ἔχων;}\\\passeLigne}%
\end{groupe}%
%% v. 181-182
\begin{groupe}%
\personnageGrec{Εὐριπίδης}%
\trimIamb{μέλλουσί μ᾽ αἱ γυναῖκες ἀπολεῖν τήμερον\\
%
τοῖς Θεσμοφορίοις, ὅτι κακῶς αὐτὰς λέγω.\\\passeLigne}%
\end{groupe}%
%% v. 183
\begin{groupe}%
\personnageGrec{Ἀγάθων}%
\trimIamb{τίς οὖν παρ᾽ ἡμῶν ἐστιν ὠφέλειά σοι;\\\passeLigne}%
\end{groupe}%
%% v. 184-187
\begin{groupe}%
\personnageGrec{Εὐριπίδης}%
\trimIamb{ἡ πᾶσ᾽· ἐὰν γὰρ ἐγκαθεζόμενος λάθρᾳ\apparat{184}{ἡ πᾶσα R : \lecon{ἡ πᾶσ᾽·} Küster : ἅπασ' Scaliger}\\
ἐν ταῖς γυναιξίν, ὡς δοκῶν εἶναι γυνή,\\
ὑπεραποκρίνῃ μου, σαφῶς σώσεις ἐμέ.%
\commentaire{186}{σαφῶς σώσεις ἐμέ.}{Le sigmatisme tant reproché à Euripide ! toujours avec le verbe σῴζειν, \italique{sauver}, le véritable Euripide avait sifflé au vers 476 de sa \italique{Médée} : ἔσῳσά σ' ὡς ἴσασιν Ἑλλήνων ὅσοι. «~Grand merci de nous avoir sauvé des sigmas d'Euripide~» avait persiflé Platon le Comique\marqueur{1} en reprenant le début de vers : %
\citComm{εὖ γέ σοι γενοιθ᾽ ἡμᾶς ὅτι\\
ἔσωσας ἐκ τῶν σίγμα τῶν Εὐριπίδου.}%
}%
\debutNote{186}%
\note{1}{Fragment 29(30) de l'édition Kassel-Austin, \italique{Poetae Comici Graeci}, vol. VII (1989) ; chez Meineke, \italique{Fragmenta Comicorum Graecorum}, il se trouve au tome II, p. 626.}\\
μόνος γὰρ ἂν λέξειας ἀξίως ἐμοῦ.\\\passeLigne}%
\end{groupe}%
%% v. 188
\begin{groupe}%
\personnageGrec{Ἀγάθων}%
\trimIamb{ἔπειτα πῶς οὐκ αὐτὸς ἀπολογεῖ παρών;\\\passeLigne}%
\end{groupe}%
%% v. 189-192
\begin{groupe}%
\personnageGrec{Εὐριπίδης}%
\trimIamb{ἐγὼ φράσω σοι. πρῶτα μὲν γιγνώσκομαι\\
ἔπειτα πολιός εἰμι καὶ πώγων᾽ ἔχω,\\
σὺ δ᾽ εὐπρόσωπος λευκὸς ἐξυρημένος\\
γυναικόφωνος ἁπαλὸς εὐπρεπὴς ἰδεῖν.%\testimonia{192}{(γυναικόφωνος) Pollux II 111}
\\\passeLigne}%
\end{groupe}%
%% v. 193a
\begin{groupe}%
\personnageGrec{Ἀγάθων}%
\trimIamb{Εὐριπίδη—\\\passeLigne}%
\end{groupe}%
%% v. 193b
\begin{groupe}%
\personnageGrec{Εὐριπίδης}%
\trimIamb{\ligne{193}\decalage{Εὐριπίδη—}τί ἔστιν;\\\passeLigne}%
\end{groupe}%
%% v. 193c-194
\begin{groupe}%
\personnageGrec{Ἀγάθων}%
\trimIamb{\ligne{193}\decalage{Εὐριπίδη—τί ἔστιν;}ἐποίησάς ποτε,%\testimonia{193}{(ἐποίησάς ποτε) \Eustathe 4.9}
\\
`χαίρεις ὁρῶν φῶς, πατέρα δ᾽ οὐ χαίρειν δοκεῖς;'
\commentaire{193}{χαίρεις ὁρῶν φῶς, πατέρα δ᾽ οὐ χαίρειν δοκεῖς;}{Citation tirée de l'\italique{Alceste} d'Euripide, v. 691 : Phérès envoie paître son fils Admète qui lui demande de mourir à sa place. C'est finalement Alceste qui se sacrifiera pour son mari ; tout comme Mnésiloque, parent par \italique{alliance} d'Euripide, ne va pas tarder à accepter ce qui ressemble fort à un suicide.}\\\passeLigne}%
\end{groupe}%
%% v. 195a
\begin{groupe}%
\personnageGrec{Εὐριπίδης}%
\trimIamb{ἔγωγε.\\\passeLigne}%
\end{groupe}%
%% v. 195b-199
\begin{groupe}%
\personnageGrec{Ἀγάθων}%
\trimIamb{\ligne{195}\decalage{ἔγωγε.}μή νυν ἐλπίσῃς τὸ σὸν κακὸν%\testimonia{195-7}{(μή νυν -- φέρε) Suidas υ 709 (νῦν pro νυν et ἀλλ' omisso)}
\\
ἡμᾶς ὑφέξειν. καὶ γὰρ ἂν μαινοίμεθ᾽ ἂν.\apparat{196}{γὰρ R : \lecon{γὰρ ἂν} Suid. υ 709}\\
ἀλλ᾽ αὐτὸς ὅ γε σόν ἐστιν οἰκείως φέρε.\\
τὰς συμφορὰς γὰρ οὐχὶ τοῖς τεχνάσμασιν%\testimonia{198-201}{(τὰς συμφορὰς -- παθήμασιν) Suidas τ 440 (τεχνάσμασι et παθήμασι pro -σιν)}
\\
φέρειν δίκαιον ἀλλὰ τοῖς παθήμασιν.\\\passeLigne}%
\end{groupe}%
%% v. 200-201
\begin{groupe}%
\personnageGrec{Μνησίλοχος}%
\trimIamb{καὶ μὴν σύ γ᾽ ὦ κατάπυγον εὐρύπρωκτος εἶ\\
οὐ τοῖς λόγοισιν ἀλλὰ τοῖς παθήμασιν.\commentaire{200-1}{εὐρύπρωκτος εἶ / οὐ τοῖς λόγοισιν ἀλλὰ τοῖς παθήμασιν.}{Mnésiloque reprend l'opposition τοῖς τεχνάσμασιν / τοῖς παθήμασιν d'Agathon en substituant τοῖς λόγοισιν à τοῖς τεχνάσμασιν.\\
Le spectateur attend donc l'opposition classique τοῖς λόγοις / τοῖς ἔργοις : \italique{en théorie} / \italique{en pratique}... sauf qu'ici les \italique{pratiques} d'Agathon sont jugées par trop passives.}\\\passeLigne}%
\end{groupe}%
%% v. 202
\begin{groupe}%
\personnageGrec{Εὐριπίδης}%
\trimIamb{τί δ᾽ ἔστιν ὅτι δέδοικας ἐλθεῖν αὐτόσε;\\\passeLigne}%
\end{groupe}%
%% v. 203a
\begin{groupe}%
\personnageGrec{Ἀγάθων}%
\trimIamb{κάκιον ἀπολοίμην ἂν ἢ σύ.\\\passeLigne}%
\end{groupe}%
%% v. 203b
\begin{groupe}%
\personnageGrec{Εὐριπίδης}%
\trimIamb{\ligne{203}\decalage{κάκιον ἀπολοίμην ἂν ἢ σύ.}πῶς;\\\passeLigne}%
\end{groupe}%
%% v. 203c-205
\begin{groupe}%
\personnageGrec{Ἀγάθων}%
\trimIamb{\ligne{203}\decalage{κάκιον ἀπολοίμην ἂν ἢ σύ.πῶς;}ὅπως;\\
δοκῶν γυναικῶν ἔργα νυκτερείσια\apparat{204}{\lecon{νυκτερείσια} R : νυκτερήσια Bothe}%
\commentaire{204}{νυκτερείσια.}{La forme canonique est νυκτερήσια, \italique{nocturne} (attestée tardivement par Lucien et Sextus Empiricus), forgée en miroir de ἡμερήσιος. %(\italique{nocturne}, attesté tardivement par Lucien et Sextus Empiricus). S'il n'est pas douteux qu'Aristophane emploie cet adjectif en ce sens – ou bien le forge en miroir de ἡμερήσιος (Eschyle, \italique{Agamemnon}, 22) –, 
Gardons-nous pour autant de retoucher la graphie du codex de Ravenne ; la réplique perdrait de son sel aristophanesque : % si l'on ne conservait pas la prononciation transmise par le codex de Ravenne et certainement voulue par Aristophane lui-même : 
νυκτ-ερείσια fait entendre le substantif du verbe ἐρείδειν, « \latin{quod in comoedia est, subagitare mulierem} » (Fritzsche).}\\
κλέπτειν ὑφαρπάζειν τε θήλειαν Κύπριν.%
\commentaire{205}{ὑφαρπάζειν τε θήλειαν Κύπριν.}{ L'expression ὑφαρπάζειν Κύπριν se retrouve au \nieme{722} vers de l'\italique{Assemblée des femmes}.
Les Athéniennes ont toujours une bonne raison d'en vouloir aux poètes tragiques : lorsqu'ils ne médisent pas d'elles, ils se livrent à une concurrence déloyale.%
\citComm{Mais les jeunes femmes d'Athènes\\
Privées d'leurs époux, d'leurs galants,\\
Accumulèrent la rancune patiemment...\\
Puis un jour, ivres de colère,\\
Elles s'armèrent de bâtons,\\
Et, farouches, elles immolèrent Agathon !}%
%Remarque en passant : le fait que l'action se déroule au \italique{jour du milieu} des Thesmophories, νηστεία, et qu'il fût théoriquement un jour d'abstinence sexuelle n'empêchera en rien les médisances... au contraire !
}\\\passeLigne}%
\end{groupe}%
%% v. 206-207
\begin{groupe}%
\personnageGrec{Μνησίλοχος}%
\trimIamb{ἰδού γε κλέπτειν· νὴ Δία βινεῖσθαι μὲν οὖν.\apparat{206}{\lecon{νὴ Δία} R : μà Δία Fritzsche}\\
ἀτὰρ ἡ πρόφασίς γε νὴ Δί᾽ εἰκότως ἔχει.
\commentaire{207}{ἀτὰρ ἡ πρόφασίς γε νὴ Δί᾽ εἰκότως ἔχει.}{Molière, \italique{Misanthrope}, IV, 3 :
\citComm{\francais{Ah! le détour est bon et l'excuse admirable.}}%
\textbf{La défaite}: prétexte fallacieux pour se défaire de ce qui est importun, embarrassant. (J'ai laissé le mot soigneusement choisi par Willems.)}\\\passeLigne}% 
\end{groupe}%
\endnumbering
\end{Leftside}
\begin{Rightside}
\beginnumbering
%% v. 165+ (part 2)
\begin{groupe}
\trimIambFr{\ligne{165}beau lui-même, il s'habillait de beaux vêtements
et c'est aussi ce qui fait que ses drames étaient beaux.
Car c'est une nécessité que l'on compose suivant sa nature.\\\passeLigne}%
\end{groupe}%
%% v. 168-170
\begin{groupe}%
\personnageFr{Mnésiloque}%
\trimIambFr{\small{Voilà aussi pourquoi Philoclès, qui est laid, compose de laides choses ;\testimonia{168}{(ὁ Φιλοκλέης αἰσχρὸς $=$ κορυδός) scholium \italique{Aves} 1295a.α}
%
pourquoi Xénoclès, qui est mauvais, en compose de  mauvaises ;
%
pourquoi Théognis, qui est froid, \fin compose froidement.}\\\passeLigne}%
\end{groupe}%
%% v. 171-172a
\begin{groupe}%
\personnageFr{Agathon}%
\trimIambFr{De toute rigueur. C'est, en effet, pour avoir su cela\\
que j'ai soigné ma personne.\\\passeLigne}%
\end{groupe}%
%% v. 172b
\begin{groupe}%
\personnageFr{Mnésiloque}%
\trimIambFr{\ligne{172}\espace{3.2cm}Comment, au nom des dieux ?\\\passeLigne}%
\end{groupe}%
%% v. 173-174
\begin{groupe}%
\personnageFr{Euripide}%
\trimIambFr{Cesse de japper. Car tel j'étais aussi moi-même\testimonia{173}{(βαύ̈ζων -- ἦν) Suidas β 192 $\simeq$ Zonaras 379 (sed γὰρ καὶ pro καὶ γὰρ)}\\
à son âge, quand je commençais à produire.\\\passeLigne}%
\end{groupe}%
%% v. 175
\begin{groupe}%
\personnageFr{Mnésiloque}%
\trimIambFr{Par Zeus ! je ne t'envie pas ton éducation.\testimonia{175}{(οὐ -- παιδεύσεως) Suidas ζ 63}\\\passeLigne}%
\end{groupe}%
%% v. 176a
\begin{groupe}%
\personnageFr{Euripide}%
\trimIambFr{Mais laisse-moi dire ce qui m'amène.\testimonia{176}{Scholium v. 500 (sed ἀλλ' εἰπεῖν βούλομαι pro ἔα μ' εἰπεῖν)}\\\passeLigne}%
\end{groupe}%
%% v. 176b
\begin{groupe}%
\personnageFr{Agathon}%
\trimIambFr{\ligne{176}\espace{5.5cm}Parle.\\\passeLigne}%
\end{groupe}%
%% v. 177-180a
\begin{groupe}%
\personnageFr{Euripide}%
\trimIambFr{Agathon, « c'est le propre d'un sage de savoir résumer en bref
de longs discours ».\testimonia{177-80}{(σοφοῦ -- πρὸς σέ) Suidas σ 792}\testimonia{177-8}{(σοφοῦ -- λόγους) Suidas οι 154}\\
Frappé d'une calamité inouïe,\testimonia{179}{Suidas κ 2556}\\
je viens à toi en suppliant.\\\passeLigne}%
\end{groupe}%
%% v. 180b
\begin{groupe}%
\personnageFr{Agathon}%
\trimIambFr{\ligne{180}\espace{3.5cm}Que te faut-il ?\\\passeLigne}%
\end{groupe}%
%% v. 181-182
\begin{groupe}%
\personnageFr{Euripide}%
\trimIambFr{Les femmes pensent consommer ma perte, aujourd'hui,\\
%
aux Thesmophories, parce que je dis du mal d'elles.\\\passeLigne}%
\end{groupe}%
%% v. 183
\begin{groupe}%
\personnageFr{Agathon}%
\trimIambFr{Et que pouvons-nous pour toi ?\\\passeLigne}%
\end{groupe}%
%% v. 184-187
\begin{groupe}%
\personnageFr{Euripide}%
\trimIambFr{Tout. Si siégeant en tapinois
parmi les femmes,\\ en feignant d'être l'une d'elles,
tu prends ma défense,\\ tu ne manqueras pas à me sauver.\\
Car seul tu parleras d'une manière digne de moi.\\\passeLigne}%
\end{groupe}%
%% v. 188
\begin{groupe}%
\personnageFr{Agathon}%
\trimIambFr{Alors, que ne vas-tu te défendre en personne ?\\\passeLigne}%
\end{groupe}%
%% v. 189-192
\begin{groupe}%
\personnageFr{Euripide}%
\trimIambFr{Je te le dirai. D'abord on me connaît.\\
Ensuite je suis tout chenu, et j'ai de la barbe.\\
Toi, tu es avenant de figure, blanc, rasé,\\
doué d'une voix de femme, mignon, plaisant à voir.\testimonia{192}{(γυναικόφωνος) Pollux II 111}\\\passeLigne}%
\end{groupe}%
%% v. 193a
\begin{groupe}%
\personnageFr{Agathon}%
\trimIambFr{Euripide...\\\passeLigne}%
\end{groupe}%
%% v. 193b
\begin{groupe}%
\personnageFr{Euripide}%
\trimIambFr{\ligne{193}\espace{1.5cm}Eh bien ?\\\passeLigne}%
\end{groupe}%
%% v. 193c-194
\begin{groupe}%
\personnageFr{Agathon}%
\trimIambFr{\ligne{193}\espace{3cm}Tu écrivis un jour : « Tu prends plaisir\testimonia{193}{(ἐποίησάς ποτε) \Eustathe 4.9}
à voir cette lumière, penses-tu que ton père y prenne \fin moins plaisir ? »\\\passeLigne}%
\end{groupe}%
%% v. 195a
\begin{groupe}%
\personnageFr{Euripide}%
\trimIambFr{Sans doute.\\\passeLigne}%
\end{groupe}%
%% v. 195b-199
\begin{groupe}%
\personnageFr{Agathon}%
\trimIambFr{\ligne{195}\espace{1.2cm}N'espère donc pas que je prenne sur mon compte ton infortune.\testimonia{195-7}{(μή νυν -- φέρε) Suidas υ 709 (νῦν pro νυν et ἀλλ' omisso)}
Il faudrait que je fusse fou.\\
Charge-toi personnellement de tes propres affaires.\\
Car les revers, il convient de les subir à force, non d'artifices,\testimonia{198-201}{(τὰς συμφορὰς -- παθήμασιν) Suidas τ 440 (-σι pro -σιν) %(τεχνάσμασι et παθήμασι pro -σιν)
}
mais d'endurance.\\\passeLigne}%
\end{groupe}%
%% v. 200-201
\begin{groupe}%
\personnageFr{Mnésiloque}%
\trimIambFr{Au fait toi-même, bougre de cul, tu as l'anus béant à force
non de discours mais d'endurance.\\\passeLigne}%
\end{groupe}%
%% v. 202
\begin{groupe}%
\personnageFr{Euripide}%
\trimIambFr{Mais qu'est-ce qui te fait craindre d'aller là ?\\\passeLigne}%
\end{groupe}%
%% v. 203a
\begin{groupe}%
\personnageFr{Agathon}%
\trimIambFr{Je finirais plus misérablement que toi.\\\passeLigne}%
\end{groupe}%
%% v. 203b
\begin{groupe}%
\personnageFr{Euripide}%
\trimIambFr{\ligne{203}\espace{4cm}Pourquoi ?\\\passeLigne}%
\end{groupe}%
%% v. 203c-205
\begin{groupe}%
\personnageFr{Agathon}%
\trimIambFr{\ligne{203}\espace{5cm}Pourquoi ?\\
Je passe pour ravir les travaux nocturnes des femmes,\\
et dérober la Cypris féminine.\\\passeLigne}%
\end{groupe}%
%% v. 206-207
\begin{groupe}%
\personnageFr{Mnésiloque}%
\trimIambFr{Mon cul, ravir ! Eh par Zeus! dis de te faire enfiler.\\
Mais, par Zeus, la défaite est bonne.%
%\commentaire{207}{La défaite}{: prétexte fallacieux pour se défaire de ce qui est importun, embarrassant. (J'ai laissé le mot choisi avec soin par Willems.)}
\\\passeLigne}%
\end{groupe}%
\endnumbering
\end{Rightside}
\end{pages}
\Pages
\clearpage\ifodd\value{page}\hbox{}\newpage\fi
\enTeteG{ΠΡΟΛΟΓΟΣ --- Codex Ravennas 429, f. 165r (Classense 167r) --- vv. 208–251}
\begin{adjustwidth}{-1.5cm}{-3.0cm}
\begin{adjustbox}{max totalheight=21.0cm, max width=16.5cm, keepaspectratio}
\begin{minipage}{16.5cm}
{\fontsize{7.2pt}{8.5pt}\selectfont
\noindent\textbf{Folio 06 --- Codex Ravennas 429, \texttt{f. 165r} (\texttt{Classense 167r}) : ΠΡΟΛΟΓΟΣ, vv. 208–251}\hfill \textit{Témoins :} $\text{R}^1, \Rpc, \Rdeux$, \PsiUn\par\vspace{2pt}
\hrule height 0.4pt\vspace{2.5pt}
}
{\fontsize{8.1pt}{9.5pt}\selectfont
\ligneRav{208}{\auDessus{εὐ}{ρ}}{τί \difference{οὖν}; ποιήσεις \difference{ταῦτα}; : μὴ δόκει γε σύ·}{\textbf{208} οὖν et ταῦτα R : \lecon{οὖν;} et \lecon{ταῦτα;} Ellebodius}
\ligneRav{209}{\auDessus{εὐ}{ρ}}{ὦ τρισκακοδαίμον, ὡς ἀπόλωλ᾽ Εὐριπίδης·}{\textbf{209} \lecon{τρισκακόδαιμον} R : τρισκακοδαίμων \Mon \espApp \lecon{(Εὐρ.)} R : (Μνη.) Elmsley ad \italique{Ach}. 475}
\ligneRav{210}{μ\auDessus{ν}{Η}}{ὦ φίλτατ᾽ ὦ κηδεστὰ μὴ σαυτὸν προδῷς·}{}
\ligneRav{}{\auDessus{εὐ}{ρ}}{πῶς οὖν ποιήσω δῆτα; : τοῦτον μὲν μακρὰ}{}
\ligneRav{}{}{κλάειν κέλευ᾽, ἐμοὶ δ᾽ ὅ τι βούλει χρῶ λαβών·}{}
\ligneRav{}{\auDessus{εὐ}{ρ}}{ἄγε νυν ἐπειδὴ σαυτὸν ἐπιδίδως ἐμοί,}{}
\ligneRav{}{—}{ἀπόδυθι τουτὶ θοἰμάτιον. : καὶ δὴ χαμαί·}{}
\ligneRav{215}{—}{ἀτὰρ τί μέλλεις δρᾶν \difference{μ᾽}; : ἀποξυρεῖν ταδί,}{\textbf{215} με R : \lecon{μ᾽} edd. omnes}
\ligneRav{216}{—}{τὰ κάτω δ᾽ ἀφεύειν. : ἀλλὰ πρᾶττ᾽, \difference{εἴ σοι} δοκεῖ·}{\textbf{216} \lecon{ἀφεύειν} R : ἀφαύειν Suid. α 4565 \espApp είσω R : \lecon{εἴ σοι} Bonini}
\ligneRav{217}{}{ἢ μὴ \difference{διδόναι γ' ἑαυτὸν} ὤφελόν ποτε·}{\textbf{217} διδόναι γ' αὐτὸν R : διδόναι γ' ἑμαυτὸν Scaliger : \lecon{διδόναι γ' ἑαυτὸν} Bibi auctore Willems}
\ligneRav{}{\auDessus{εὐ}{ρ}}{Ἀγάθων σὺ μέντοι ξυροφορεῖς ἑκάστοτε,}{}
\ligneRav{}{—}{χρῆσόν τι νῦν ἡμῖν ξυρόν. : αὐτὸς λάμβανε}{}
\ligneRav{220}{—}{ἐντεῦθεν ἐκ τῆς \difference{ξυροδόκης}. : γενναῖος εἶ·}{\textbf{220} ξυροδίκης R : \lecon{ξυροδόκης} Σ, Pollux II.32, X.140}
\ligneRav{}{}{κάθιξε· φύσα τὴν γνάθον τὴν δεξιάν·}{}
\ligneRav{222}{μ\auDessus{ν}{Η}}{\difference{οἴμοι}. : τί κέκραγας; ἐμβαλῶ σοι πάτταλον,}{\textbf{222} ὤιμοι R : \lecon{οἴμοι} Dindorf}
\ligneRav{223}{—}{ἢν μὴ σιωπᾷς. : \difference{ἀτταταῖ ἰατταταῖ}·}{\textbf{223} ἀττατα· ατταται R : \lecon{ἀτταταῖ ἰατταταῖ} Biset}
\ligneRav{}{\auDessus{εὐ}{ρ}}{οὗτος σὺ ποῖ θεῖς; : εἰς τὸ τῶν σεμνῶν θεῶν·}{}
\ligneRav{225}{}{οὐ γὰρ μὰ τὴν Δήμητρά γ᾽ ἐνταυθοῖ μενῶ}{\textbf{225} \lecon{Δήμητρά γ᾽} R : Δήμητρ᾽ ἔτ᾽ Porson, cf. ad \italique{Eq}. 698 : Δημήτερ᾽ Fritzsche}
\ligneRav{226}{—}{τεμνόμενος. : \difference{οὔκουν} καταγέλαστος δῆτ᾽ ἔσει}{\textbf{226} οὐκ οὖν R : \lecon{οὔκουν} Brunck}
\ligneRav{}{}{τὴν ἡμίκραιραν τὴν ἑτέραν ψιλὴν ἔχων;}{}
\ligneRav{}{μ\auDessus{ν}{Η}}{ὀλίγον μέλει μοι. : μηδαμῶς πρὸς τῶν θεῶν}{}
\ligneRav{}{—}{προδῷς με· χώρει δεῦρο. : κακοδαίμων ἐγώ·}{}
\ligneRav{230}{\auDessus{εὐ}{ρ}}{ἔχ᾽ ἀτρέμα σαυτὸν κἀνάκυπτε· ποῖ στρέφει;}{}
\ligneRav{}{μ\auDessus{ν}{Η}}{μυμῦ. : τί μύζεις; πάντα πεποίηται καλῶς·}{}
\ligneRav{232}{μ\auDessus{ν}{Η}}{οἴμοι κακοδαίμων, ψιλὸς αὖ στρατεύσομαι·}{\textbf{232} \lecon{αὖ} R, Suid. ψ 105 : οὖν von Velsen}
\ligneRav{}{\auDessus{εὐ}{ρ}}{μὴ φροντίσῃ· ὡς εὐπρεπὴς φανεῖ πάνυ·}{}
\ligneRav{234}{—}{βούλει \difference{θεάσασθ᾽ αὑτόν}; : εἰ δοκεῖ, φέρε·}{\textbf{234} θεάσασθαι σεαυτόν R : \lecon{θεάσασθ᾽ αὑτόν} Fritzsche : θεᾶσθαι σαυτόν Porson}
\ligneRav{235}{\auDessus{εὐ}{ρ}}{ὁρᾷς σεαυτόν; : οὐ μὰ Δί᾽ ἀλλὰ \difference{Κλεισθένη}·}{\textbf{235} Κλεισθένην R, Suid. κ 1756 : \lecon{Κλεισθένη} Brunck}
\ligneRav{}{\auDessus{εὐ}{ρ}}{ἀνίστασ᾽, ἵν᾽ ἀφεύσω σε, κἀγκύψας ἔχε·}{}
\ligneRav{}{μ\auDessus{ν}{Η}}{οἴμοι κακοδαίμων δελφάκιον γενήσομαι·}{}
\ligneRav{}{\auDessus{εὐ}{ρ}}{ἐνεγκάτω τις ἔνδοθεν δᾷδ᾽ ἢ λύχνον·}{}
\ligneRav{}{}{ἐπίκυπτε· τὴν κέρκον φυλάττου νυν ἄκραν·}{}
\ligneRav{240}{μ\auDessus{ν}{Η}}{ἐμοὶ μελήσει νὴ Δία, πλήν γ᾽ ὅτι κάομαι·}{}
\ligneRav{}{}{οἴμοι τάλας. ὕδωρ ὕδωρ ὦ γείτονες}{}
\ligneRav{242}{}{τὸν \difference{[ἕ]τερον} πρὶν ἀντιλαβέσθαι φλόγα·}{\textbf{242} \lecon{[ἓ]τερον} \PsiUn \ (restituit Medaglia) : ἓτερον om. R}
\ligneRav{}{\auDessus{εὐ}{ρ}}{θάρρει. : τί θαρρῶ καταπεπυρπολημένος;}{}
\ligneRav{}{\auDessus{εὐ}{ρ}}{ἀλλ᾽ οὐκ ἔτ᾽ οἰδὲν πρᾶγμά σοι· τὰ πλεῖστα γὰρ}{}
\ligneRav{245}{—}{ἀποπεπόνηκας. : φεῦ ἰοὺ τῆς ἀσβόλου·}{\textbf{245} \lecon{φεῦ} R, Suid. α 4132 : φῦ Dindorf, Wilamowitz \italique{Lysistr.} v.294}
\ligneRav{246}{}{αἰθὸς \difference{γεγένημαι} πάντα τὰ περὶ τὴν τράμιν·}{\textbf{246} γεγέννημαι R : \lecon{γεγένημαι} Suid. τ 903}
\ligneRav{247}{\auDessus{εὐ}{ρ}}{μὴ φροντίσῃς· ἕτερος γὰρ αὐτὰ σπογγιεῖ·}{\textbf{247} \lecon{σπογγιεῖ} R : σφογγιεῖ Elmsley ad \italique{Ach.} 463}
\ligneRav{248}{μ\auDessus{ν}{Η}}{\difference{οἰμώξετ᾽} ἆρ᾽ \difference{εἴ τις} τὸν ἐμὸν πρωκτὸν πλυνεῖ·}{\textbf{248} οἰμώζετ᾽ R : \lecon{οἰμώξετ᾽} Zanetti \espApp εἰς R : \lecon{εἴ τις} Brunck}
\ligneRav{}{\auDessus{εὐ}{ρ}}{Ἀγάθων, ἐπειδὴ σαυτὸν ἐπιδοῦναι φθονεῖς,}{}
\ligneRav{250}{}{ἀλλ᾽ ἱμάτιον γοῦν χρῆσον ἡμῖν \difference{τουτῳὶ}}{\textbf{250} τουτωί R : \lecon{τουτῳὶ} edd. plerique}
\ligneRav{}{}{καὶ στρόφιον· οὐ γὰρ ταῦτά γ᾽ ὡς οὐκ ἔστ᾽ ἐρεῖς·}{}
}
\vspace{3pt}\noindent\rule{4.5cm}{0.4pt}\par\vspace{1.5pt}
{\fontsize{7.0pt}{8.3pt}\selectfont
\noindent\textbf{Loci paralleli (vv. 208–251) ---} \textbf{215-6} (ἀποξυρεῖν -- ἀφεύειν) Suidas α 4565 (ἀφαύειν pro ἀφεύειν), cf. Scholium ταῦτα δὲ ἔλαβεν ἐκ τῶν Ἰδαίων Κρατίνου (fr. 90 Kassel-Austin) ; \ClemAlex, \italique{Στρώματα} VI 26.4 Ἀριστοφάνης δὲ ὁ κωμικὸς ἐν ταῖς πρώταις Θεσμοφοριαζούσαις τὰ ἐκ τῶν Κρατίνου Ἐμπιπραμένων μετήνεγκεν ἔπη. \espTest \textbf{219-20} (χρῆσόν -- ξυροδόκης) Pollux X 140 (νῦν omisso) \espTest \textbf{220} (ξυροδόκης) Pollux II 32 ; \Eustathe 130.43, 1399.38 ; cf. Scholium \espTest \textbf{223} (ἀτταταῖ ἰατταταῖ) Suidas ε 2807 \espTest \textbf{227} Suidas η 334 \espTest (ἡμίκραιραν) Hesychius η 496 ἡμίκραιρα ; \EtymGen (AB) 150 ed. Miller ; \Eustathe 710.50, 1127.35 ; \AelDion η 9 \espTest \textbf{231} \EtymGen (AB) 218 ed. Miller \espTest (μῦ μῦ) Suidas ε 2807 ; cf. \Eustathe 723.9 \espTest (τί -- καλῶς) Suidas μ 1383 ; \Apostol XVI 60c \espTest (μύζεις) \Eustathe 440.23 \espTest \textbf{232} Suidas ψ 105 \espTest \textbf{235} Suidas κ 1756 Κλεισθένην ὁρῶ \espTest \textbf{237} Suidas δ 205 (sed γέγονα pro γενήσομαι) \espTest \textbf{244-5} (τὰ πλεῖστα -- ἀσβόλου) Photius α 2946 $=$ \SynB α 2213 ; \AelDion α 186 \espTest \textbf{245} (φεῦ -- ἀσβόλου) Suidas α 4132 \espTest (ἄσβολος $=$ ἀσβόλη) Moeris α 17 $\simeq$ Hesychius α 7634 ; cf. Phrynicus \italique{Praep. soph} 28.1 (ἄσβολος fem.) \espTest \textbf{246} Suidas τ 903 (sed τὸν pro τὴν) \espTest (αἰθὸς γεγένημαι) \Eustathe 83.6, 992.60, 1314.15 \espTest (αἰθὸς) \Eustathe 1385.59 \espTest (τράμιν) Hesychius τ 1243 ; Photius τ 415 \espTest \textbf{248} (πλυνεῖ) Moeris π 79\par
}
\vspace{2.5pt}\noindent\rule{4.5cm}{0.4pt}\par\vspace{1.5pt}
{\fontsize{6.9pt}{8.2pt}\selectfont
\noindent\textbf{Suite (en commentaire) de l'apparat critique \& notes paléographiques (\texttt{f. 165r}) ---} \textbf{Témoin papyrologique :} \textbf{\PsiUn{} (\italique{PSI} XI 1194)} couvre les \textbf{vv. 237–245} et comble au \textbf{v. 242} l'omission de [ἓ]τερον dans R. --- \textbf{v. 242 :} R avait omis un crétique avant πρὶν ἀντιλαβέσθαι ; le papyrus \PsiUn{} (lu par Medaglia) restitue \lecon{[ἓ]τερον}, tombé par saut visuel dans l'ancêtre de R. --- \textbf{v. 250 :} Exemple typique de \textbf{iota muet adscrit} dans $\text{R}^1$ : τουτωί. --- \textbf{v. 217 :} διδόναι γ' αὐτὸν R (διδόναι γ' ἐμαυτὸν Scaliger) : \lecon{διδόναι γ' ἑαυτὸν} \lecon{Bibi} \italique{auctore Willems}. --- \textbf{v. 220 :} ξυροδίκης R (iotacisme/attraction de $\text{δίκη}$) : \lecon{ξυροδόκης} $\Sigma$, Pollux II.32, X.140 (« étui à rasoir »). --- \textbf{v. 234 :} θεάσασθαι σεαυτόν R (glose pronominale non contracte) : \lecon{θεάσασθ' αὑτόν} Fritzsche (θεᾶσθαι σαυτόν Porson). --- \textbf{v. 242 :} \lecon{[ἓ]τερον} \PsiUn{} (restituit Medaglia) : ἕτερον om. R. --- \textbf{v. 248 :} οἰμώζετ' R : \lecon{οἰμώξετ'} Zanetti \espApp εἰς R (haplographie de ΤΙ dans ΕΙΤΙϹ) : \lecon{εἴ τις} Brunck.\par
}
\end{minipage}
\end{adjustbox}
\end{adjustwidth}
\newpage
\enTeteD{Codex Ravennas 429 (Biblioteca Classense) --- Fac-similé du f. 165r (vv. 208–251)}
\begin{adjustwidth}{-1.5cm}{-3.0cm}
\begin{center}
\includegraphics[width=16.5cm,height=20.5cm,keepaspectratio]{/workspaces/Thesmophories/dossier_ravennas_thesmophories/folios_individuels_ff162v-177r/06_f165r_classense167r_vv0208-0251.jpg}\par
\vspace{2pt}
{\scriptsize\textbf{Codex Ravennas 429, f. 165r (Classense 167r)} --- ΠΡΟΛΟΓΟΣ, \textbf{vv. 208–251} ($\text{R}^1, \Rpc, \Rdeux$, \PsiUn).}
\end{center}
\end{adjustwidth}
\newpage
\clearpage\ifodd\value{page}\hbox{}\newpage\fi
\enTeteG{ΘΕΣΜΟΦΟΡΙΑΖΟΥΣΑΙ --- Texte grec établi (f. 165r, vv. 208–251)}
\enTeteD{LES THESMOPHORIES --- Traduction française (f. 165r, vv. 208–251)}
\begin{pages}
\begin{Leftside}
\beginnumbering
%% v. 208a
\begin{groupe}%
\personnageGrec{Εὐριπίδης}%
\trimIamb{τί οὖν; ποιήσεις ταῦτα;\apparat{208}{οὖν et ταῦτα R : \lecon{οὖν;} et \lecon{ταῦτα;} Ellebodius}\\\passeLigne}%
\end{groupe}%
%% v. 208b
\begin{groupe}%
\personnageGrec{Ἀγάθων}%
\trimIamb{\ligne{208}\decalage{τί οὖν; ποιήσεις ταῦτα;}μὴ δόκει γε σύ.\\\passeLigne}%
\end{groupe}%
%% v. 209
\begin{groupe}%
\personnageGrec{Εὐριπίδης}%
\trimIamb{ὦ τρισκακοδαίμον, ὡς ἀπόλωλ᾽ Εὐριπίδης.
\apparat{209}{\lecon{τρισκακόδαιμον} R : τρισκακοδαίμων \Mon \espApp \lecon{(Εὐρ.)} R : (Μνη.) Elmsley ad \italique{Ach}. 475}%
\commentaire{209}{ὦ τρισκακόδαιμον ὡς ἀπόλωλ᾽ Εὐριπίδης.}{La plupart des éditeurs modernes découpent cette réplique, attribuant la première partie à Euripide, puis Mnésiloque de lancer l'apostrophe : Εὐριπίδη κτλ. Cela les oblige à retoucher le texte pour faire de ὤ τρισκακοδαίμων le sujet de ἀπόλωλ(α) et changer Εὐριπίδης en Εὐριπίδη... 
Tout cela paraît bien tarabiscoté alors que le texte de notre unicodex se laisse comprendre facilement : « Ô triple malheureux ! ainsi périt Euripide », apostrophe qu'Euripide s'adresse à lui-même d'autant moins étonnante qu'il a déjà parlé de lui à la \nieme{3} personne (v. 77).
En outre, si l'attribution des répliques dans le manuscrit de Ravenne est souvent folklorique, la découpe des répliques, elle, est bien plus sûre. Philogues, réfrénez vos pulsions à vouloir tailler ainsi le vers.}\\\passeLigne}%
\end{groupe}%
%% v. 210
\begin{groupe}%
\personnageGrec{Μνησίλοχος}%
\trimIamb{ὦ φίλτατ᾽ ὦ κηδεστὰ μὴ σαυτὸν προδῷς.\\\passeLigne}%
\end{groupe}%
%% v. 211a
\begin{groupe}%
\personnageGrec{Εὐριπίδης}%
\trimIamb{πῶς οὖν ποιήσω δῆτα;\\\passeLigne}%
\end{groupe}%
%% v. 211b-212
\begin{groupe}%
\personnageGrec{Μνησίλοχος}%
\trimIamb{\ligne{211}\decalage{πῶς οὖν ποιήσω δῆτα;}τοῦτον μὲν μακρὰ\\
κλάειν κέλευ᾽, ἐμοὶ δ᾽ ὅ τι βούλει χρῶ λαβών.\\\passeLigne}%
\end{groupe}%
%% v. 213-214a
\begin{groupe}%
\personnageGrec{Εὐριπίδης}%
\trimIamb{ἄγε νυν ἐπειδὴ σαυτὸν ἐπιδίδως ἐμοί,\\
ἀπόδυθι τουτὶ θοἰμάτιον.\\\passeLigne}%
\end{groupe}%
%% v. 214b-215a
\begin{groupe}%
\personnageGrec{Μνησίλοχος}%
\trimIamb{\ligne{214}\decalage{ἀπόδυθι τουτὶ θοἰμάτιον.}καὶ δὴ χαμαί.\\
ἀτὰρ τί μέλλεις δρᾶν μ᾽;\apparat{215}{με R : \lecon{μ᾽} edd. omnes}\\\passeLigne}%
\end{groupe}%
%% v. 215b-216a
\begin{groupe}%
\personnageGrec{Εὐριπίδης}%
\trimIamb{\ligne{215}\decalage{ἀτὰρ τί μέλλεις δρᾶν μ᾽;}ἀποξυρεῖν ταδί,%\testimonia{215-6}{(ἀποξυρεῖν -- ἀφεύειν) Suidas α 4565 (ἀφαύειν pro ἀφεύειν), cf. Scholium ταῦτα δὲ ἔλαβεν ἐκ τῶν Ἰδαίων Κρατίνου (fr. 90 Kassel-Austin) ; \ClemAlex, \italique{Στρώματα} VI 26.4 Ἀριστοφάνης δὲ ὁ κωμικὸς ἐν ταῖς πρώταις Θεσμοφοριαζούσαις τὰ ἐκ τῶν Κρατίνου Ἐμπιπραμένων μετήνεγκεν ἔπη.}
\\
τὰ κάτω δ᾽ ἀφεύειν.%
\apparat{216}{\lecon{ἀφεύειν} R : ἀφαύειν Suid. α 4565 \espApp είσω R : \lecon{εἴ σοι} Bonini}\\\passeLigne}%
\end{groupe}%
%% v. 216b-217
\begin{groupe}%
\personnageGrec{Μνησίλοχος}%
\trimIamb{\ligne{216}\decalage{τὰ κάτω δ᾽ ἀφεύειν.}ἀλλὰ πρᾶττ᾽, εἴ σοι δοκεῖ.\\
ἢ μὴ διδόναι γ' ἑαυτὸν ὤφελόν ποτε.%
\apparat{217}{διδόναι γ' αὐτὸν R : διδόναι γ' ἑμαυτὸν Scaliger : %'πιδοῦναι 'μαυτὸν Dawes : 
\lecon{διδόναι γ' ἑαυτὸν} Bibi auctore Willems}%
\commentaire{217}{ἢ μὴ διδόναι γ' αὐτὸν ὤφελόν ποτε.}{Le vers pèche contre le mètre.
Willems (voir sa note 4 en annexe), en guerre perpétuelle contre les retoucheurs de texte, rendait justice à ce qu'il croyait être le texte original. (En fait légèrement modifié : γ' ἑμαυτὸν à la place de γ' αὐτὸν. Willems a-t-il été trompé par Dindorf ?) 
Il s'en justifiait par des exemples qui ne laissent nul doute de l'epsilonesque correction qu'il aurait lui-même proposée s'il avait eu le texte des manuscrits sous les yeux : substituer γ' ἑαυτὸν à γ' αὐτὸν.\\ 
% Cette correction est d'autant plus facile à accepter si l'on s'imagine une version en onciales, ΓΕΑΥΤΟΝ ; un scribe ayant malencontreusement lu γε αυτον avant d'élider le γε. 
Au reste, comme l'avait souligné Willems, « διδόναι ἑαυτὸν τινι est la formule ordinaire, et qui se lit partout : \italique{Cavaliers}., 740 ; Sophocle, \italique{Philoct}., 84 ; \italique{Trach}., 1117; Hérodote, VI, 108; Thucydide, II, 68; Démosthène, \italique{Sur la couronne}, 503d.~»}%
%\debutNote{217}%
%\note{1}{En fait légèrement modifié : γ' ἑμαυτὸν à la place de γ' αὐτὸν. Willems a-t-il été trompé par Dindorf ? (qui prétendait -- à tort -- que γ' ἑμαυτὸν était la lecture originale du manuscrit de Münich.)}%
\\\passeLigne}%
\end{groupe}%
%% v. 218-219a
\begin{groupe}%
\personnageGrec{Εὐριπίδης}%
\trimIamb{Ἀγάθων σὺ μέντοι ξυροφορεῖς ἑκάστοτε,\\
χρῆσόν τι νῦν ἡμῖν ξυρόν.%\testimonia{219-20}{(χρῆσόν -- ξυροδόκης) Pollux X 140 (νῦν omisso)}
\\\passeLigne}%
\end{groupe}%
%% v. 219b-220a
\begin{groupe}%
\personnageGrec{Ἀγάθων}%
\trimIamb{\ligne{219}\decalage{χρῆσόν τί νυν ἡμῖν ξυρόν.}αὐτὸς λάμβανε\\
ἐντεῦθεν ἐκ τῆς ξυροδόκης.%\testimonia{220}{(ξυροδόκης) Pollux II 32 ; \Eustathe 130.43, 1399.38 ; cf. Scholium}%
\apparat{220}{ξυροδίκης R : \lecon{ξυροδόκης} Σ, Pollux II.32, X.140}\\\passeLigne}%
\end{groupe}%
%% v. 220b-221
\begin{groupe}%
\personnageGrec{Εὐριπίδης}%
\trimIamb{\ligne{220}\decalage{ἐντεῦθεν ἐκ τῆς ξυροδόκης.}γενναῖος εἶ.\\
κάθιξε· φύσα τὴν γνάθον τὴν δεξιάν.\\\passeLigne}%
\end{groupe}%
%% v. 222a
\begin{groupe}%
\personnageGrec{Μνησίλοχος}%
\trimIamb{οἴμοι.
\apparat{222}{ὤιμοι R : \lecon{οἴμοι} Dindorf}\\\passeLigne}%
\end{groupe}%
%% v. 222b-223a
\begin{groupe}%
\personnageGrec{Εὐριπίδης}%
\trimIamb{\ligne{222}\decalage{οἴμοι.}τί κέκραγας; ἐμβαλῶ σοι πάτταλον,\\
ἢν μὴ σιωπᾷς.\\\passeLigne}%
\end{groupe}%
\begin{groupe}%
\personnageGrec{Μνησίλοχος}%
\trimIamb{\ligne{223}\decalage{ἢν μὴ σιωπᾷς.}ἀτταταῖ ἰατταταῖ.%\testimonia{223}{(ἀτταταῖ ἰατταταῖ) Suidas ε 2807}
\apparat{223}{ἀττατα· ατταται R : \lecon{ἀτταταῖ ἰατταταῖ} Biset}\\\passeLigne}%
\end{groupe}%
%% v. 224a
\begin{groupe}%
\personnageGrec{Εὐριπίδης}%
\trimIamb{οὗτος σὺ ποῖ θεῖς;\commentaire{224}{οὗτος σὺ ποῖ θεῖς;}{formule identique au vers 854 des \italique{Guêpes}.}\\\passeLigne}%
\end{groupe}%
%% v. 224b-226a
\begin{groupe}%
\personnageGrec{Μνησίλοχος}%
\trimIamb{\ligne{224}\decalage{οὗτος σὺ ποῖ θεῖς;}εἰς τὸ τῶν σεμνῶν θεῶν·
\commentaire{224}{εἰς τὸ τῶν σεμνῶν θεῶν·}{Le scholiaste :
\citComm{ἀντὶ τοῦ `εἰς <τὸ> τῶν ᾿Ερινύων'. ἄσυλον γὰρ εἶχον οἵ καταλαμβάνοντες τὰ ἱερὰ τῶν Ἐρινύων.}}\\
οὐ γὰρ μὰ τὴν Δήμητρά γ᾽ ἐνταυθοῖ μενῶ\apparat{225}{\lecon{Δήμητρά γ᾽} R : Δήμητρ᾽ ἔτ᾽ Porson, cf. ad \italique{Eq}. 698 : Δημήτερ᾽ Fritzsche}\\
τεμνόμενος.\\\passeLigne}%
\end{groupe}%
%% v. 226b-227
\begin{groupe}%
\personnageGrec{Εὐριπίδης}%
\trimIamb{\ligne{226}\decalage{τεμνόμενος.}οὔκουν καταγέλαστος δῆτ᾽ ἔσει\apparat{226}{οὐκ οὖν R : \lecon{οὔκουν} Brunck}\\
τὴν ἡμίκραιραν τὴν ἑτέραν ψιλὴν ἔχων;%\testimonia{227}{Suidas η 334 \espTest (ἡμίκραιραν) Hesychius η 496 ἡμίκραιρα ; \EtymGen (AB) 150 ed. Miller ; \Eustathe 710.50, 1127.35 ; \AelDion η 9}
\\\passeLigne}%
\end{groupe}%
%% v. 228a
\begin{groupe}%
\personnageGrec{Μνησίλοχος}%
\trimIamb{ὀλίγον μέλει μοι.\\\passeLigne}%
\end{groupe}%
%% v. 228b-229a
\begin{groupe}%
\personnageGrec{Εὐριπίδης}%
\trimIamb{\ligne{228}\decalage{ὀλίγον μέλει μοι.}μηδαμῶς πρὸς τῶν θεῶν\\
προδῷς με· χώρει δεῦρο.\\\passeLigne}%
\end{groupe}%
%% v. 229b
\begin{groupe}%
\personnageGrec{Μνησίλοχος}%
\trimIamb{\ligne{229}\decalage{προδῷς με· χώρει δεῦρο.}κακοδαίμων ἐγώ.\\\passeLigne}%
\end{groupe}%
%% v. 230
\begin{groupe}%
\personnageGrec{Εὐριπίδης}%
\trimIamb{ἔχ᾽ ἀτρέμα σαυτὸν κἀνάκυπτε· ποῖ στρέφει;
\commentaire{230}{ἔχ᾽ ἀτρέμα σαυτὸν κἀνάκυπτε·}{C'est la leçon du \italique{Ravennas}. 
M. van Leeuwen, d'après Dobree, corrige : ἔχ᾽ ἀτρέμας αὐτοῦ, \italique{quia ferri nequit accusativus} ; lequel αὐτοῦ était à sa place en \italique{Oiseaux}, 1200, mais pas du tout ici. 
Hérodote a pourtant dit, en parlant de l'armée athénienne : εἶχον ἀτρέμας σφέας αὐτοὺς ἵνα ἐτάχθησαν\marqueur{1}. Hérodote le pouvait, mais non Aristophane, qui ailleurs s'est borné à ἔχ' ἀτρέμα\marqueur{2}. Quoi ! voilà un comique qui, à quatre reprises (car on s'est donné la peine de compter), a écrit : \italique{reste coi}, et la cinquième fois il se permettrait d'écrire : \italique{tiens-toi coi}. 
À quoi servirait la statistique si l'on tolérait de pareilles licences ?
(Note vengeresse n°5 de Willems.)}%
\debutNote{230}%
\note{1}{Hérodote, IX, 54 (aussi III, 79). — De même Platon, \italique{Epist.}, 348 B, écrit : ἐκποδὼν ἑαυτὸν ἔσχεν et Euripide, \italique{Iph. Taur.}, 1129 : ἐκποδὼν τοῦδ' ἔχειν μιάσματος.}%
\note{2}{ll est vrai qu'une des quatre fois les bons manuscrits portent ἔχ' ἀτρεμί (\italique{Nuées}, 261). Était-ce assez dégoûtant ? Heureusement que M. van Leeuwen y a mis bon ordre.}\\\passeLigne}%
\end{groupe}%
%% v. 231a
\begin{groupe}%
\personnageGrec{Μνησίλοχος}%
\trimIamb{μυμῦ.%\testimonia{231}{\EtymGen (AB) 218 ed. Miller \espTest (μῦ μῦ) Suidas ε 2807 ; cf. \Eustathe 723.9 \espTest (τί -- καλῶς) Suidas μ 1383 ; \Apostol XVI 60c \espTest (μύζεις) \Eustathe 440.23}
\\\passeLigne}%
\end{groupe}%
%% v. 231b
\begin{groupe}%
\personnageGrec{Εὐριπίδης}%
\trimIamb{\ligne{231}\decalage{μυμῦ.}τί μύζεις; πάντα πεποίηται καλῶς.%\commentaire{231}{τί μύζεις;}{\eff{Érasme 702}}
\\\passeLigne}%
\end{groupe}%
%% v. 232
\begin{groupe}%
\personnageGrec{Μνησίλοχος}%
\trimIamb{οἴμοι κακοδαίμων, ψιλὸς αὖ στρατεύσομαι.%\testimonia{232}{Suidas ψ 105}
\apparat{232}{\lecon{αὖ} R, Suid. ψ 105 : οὖν von Velsen}%
\commentaire{232}{ψιλὸς αὖ στρατεύσομαι.}{ψιλός, comme \italique{levis} en latin, signifie à la fois \italique{lisse}, \italique{sans poil}, et \italique{léger}, \italique{armé à la légère}. (Note de Willems.)\\
% Il n'y a donc pas lieu de corriger le texte : le « αὖ » sert au jeu de mots.\\
Cette note sur ψιλός me sera l'occasion de digresser sur l'un des plus beaux échanges à bâtons rompus de la Tragédie grecque : chez Sophocle (\italique{Ajax}, v. 1120 sq.), entre Ménélas et Teucer :
\citComm{Μενέλαος — Ὁ τοξότης ἔοικεν οὐ σμικρὸν φρονεῖν.\\	
Τεύκρος — Οὐ γὰρ βάναυσον τὴν τέχνην ἐκτησάμην.\\
Μενέλαος — Μέγ᾽ ἄν τι κομπάσειας, ἀσπίδ᾽ εἰ λάβοις.\\
Τεύκρος — Κἂν \textbf{ψιλὸς} ἀρκέσαιμι σοί γ᾽ ὡπλισμένῳ.\\
Μενέλαος — Ἡ γλῶσσά σου τὸν θυμὸν ὡς δεινὸν τρέφει.\\
Τεύκρος — Ξὺν τῷ δικαίῳ γὰρ μέγ᾽ ἔξεστιν φρονεῖν.}%
Ménélas s'étonne de ce que Teucer se permette de lui adresser des reproches – à lui le roi de Sparte ! –, vil archer de son état --- Mais l'archerie est loin d'être méprisable lui rétorque Teucer --- Ménélas ironise : Que serait-ce s'il portait un bouclier --- La réplique de Teucer fuse : « Même \italique{ψιλὸς}, je tiendrais bon devant toi et ton [royal] armement ». Comprendre : même \italique{fantassin de l'armée légère} ; c'est-à-dire l'échelon le plus bas chez les militaires – encore plus bas qu'archer –, poste souvent réservé à des esclaves et donc le plus opposé à la royale position de Ménélas et le plus inégal en terme d'armements. C'en est trop pour Ménélas qui s'indigne de cette langue par trop pendue. C'est alors que Teucer lance la fameuse et cinglante réplique : « La fierté est permise, quand on a pour soi la justice ! »}\\\passeLigne}%
\end{groupe}%
%% v. 233-234a
\begin{groupe}%
\personnageGrec{Εὐριπίδης}%
\trimIamb{μὴ φροντίσῃ· ὡς εὐπρεπὴς φανεῖ πάνυ.\\
βούλει θεάσασθ᾽ αὑτόν;\apparat{234}{θεάσασθαι σεαυτόν R : \lecon{θεάσασθ᾽ αὑτόν} Fritzsche : θεᾶσθαι σαυτόν Porson}%
\commentaire{234}{βούλει θεάσασθαι σεαυτόν; — εἰ δοκεῖ, φέρε.}{Le vers péchant contre le mètre, la correction la plus convaincante me semble être celle de Fritzsche qui substitue αὑτόν au σεαυτόν, faisant l'hypothèse que σεαυτόν aura été suggéré à un scribe inattentif par la proximité avec ὁρᾷς σεαυτόν du vers suivant.}%
\\\passeLigne}%
\end{groupe}%
%% v. 234b
\begin{groupe}%
\personnageGrec{Μνησίλοχος}%
\trimIamb{\ligne{234}\decalage{βούλει θεάσασθ᾽ αὑτόν;}εἰ δοκεῖ, φέρε.\\\passeLigne}%
\end{groupe}%
%% v. 235a
\begin{groupe}%
\personnageGrec{Εὐριπίδης}%
\trimIamb{ὁρᾷς σεαυτόν;%
%\commentaire{234-5}{φέρε — ὁρᾷς σεαυτόν;}{Sur le τόπος comique du miroir, \cf Platon, \italique{Alcibiade}, I, 132.\eff{Remarquons que ce doit être à ce moment que Mnésiloque saisit le miroir que l'on retrouve quelques scènes plus loin sur le cratère de Würtzburg. Et le poignard serait-il le rasoir ?}}
\\\passeLigne}%
\end{groupe}%
%% v. 235b
\begin{groupe}%
\personnageGrec{Μνησίλοχος}%
\trimIamb{\ligne{235}\decalage{ὁρᾷς σεαυτόν;}οὐ μὰ Δί᾽ ἀλλὰ Κλεισθένη.%\testimonia{235}{Suidas κ 1756 Κλεισθένην ὁρῶ}
\apparat{235}{Κλεισθένην R, Suid. κ 1756 : \lecon{Κλεισθένη} Brunck}%
\commentaire{235}{Κλεισθένην.}{Clisthène, l'efféminé notoire, personnage qui apparaîtra dans la pièce au vers 571 -- et sera d'ailleurs reconnue comme femme -- avant de prêter allégeance au genre féminin. 
Cette tête de Turc d'Aristophane est raillée comme : une pseudo-femme (\italique{Nuées} 355, \italique{Oiseaux} 829-831, \italique{Thesmophories} 571) ; en tout cas femme en amour (\italique{Lysistrata} 621, 1091-1092, \italique{Thesmophories} 574, \italique{Grenouilles} 48) ; imberbe comme femme (\italique{Acharniens} 117-121, \italique{Grenouilles} 422-424) ; Clisthène ne contrebalance même pas ces innombrables défauts – qui ne sont pas tous du ressort de sa volonté – par une misogynie affectée : il trouve encore le moyen d'être traître à son sexe en s'alliant au genre féminin (\italique{Lysistrata} 621, \italique{Thesmophories} 575-576, 602).}\\\passeLigne}%
\end{groupe}%
%% v. 236
\begin{groupe}%
\personnageGrec{Εὐριπίδης}%
\trimIamb{ἀνίστασ᾽, ἵν᾽ ἀφεύσω σε, κἀγκύψας ἔχε.\\\passeLigne}%
\end{groupe}%
%% v. 237
\begin{groupe}%
\personnageGrec{Μνησίλοχος}%
\trimIamb{οἴμοι κακοδαίμων δελφάκιον γενήσομαι.%\testimonia{237}{Suidas δ 205 (sed γέγονα pro γενήσομαι)}%
\commentaire{237}{δελφάκιον}{ est un diminutif d'ἡ δέλφαξ, le porc \italique{mature}, opposé au χοῖρος, \italique{le porcelet}.
%Hésychius (δ 599) précise que le mot δελφάκιον, comme χοῖρος, il peut renvoyer au γυναικεῖον [αἰδοῖον]. Mais ne nous égarons pas ! 
N'en déplaise à certains commentateurs, il n'est fait aucune allusion ici au sexe féminin -- pas plus qu'au vers 1061 de \italique{Lysistrata}. Et d'abord, ce serait absurde : il n'est question que du côté pile ; côté épilé par flambage. 
L'image est ailleurs : c'est celle de de viande de cochon dont on retire les poils par le feu.
Cf. la \italique{Souda} (δ 205):%
\citComm{δελφάκιον : ὁ μικρὸς χοῖρος. Ἀριστοφάνης: οἴμοι κακοδαίμων, δελφάκιον γέγονα. τοῦτο εἴρηκεν ἀφευθεὶς τὸν πρωκτόν. ἀντὶ τοῦ φλογισθείς. μετὰ γὰρ τὸ τυθῆναι τὰ δελφάκια φλογίζονται, ἵνα ψιλωθῶσιν.}%
(On pourra aussi se reporter à Homère\marqueur{1} ou Eschyle\marqueur{2}.)}%
\debutNote{237}%
\note{1}{\italique{Iliade}, IX, 467-8 et \italique{Odyssée}, II, 300 ; XIV, 73-5}%
\note{2}{Fragment 310 de l'édition de S. Radt, \italique{Tragicorum Graecorum Fragmenta}, vol. 3 (1985).}%
\\\passeLigne}%
\end{groupe}%
%% v. 238-239
\begin{groupe}%
\personnageGrec{Εὐριπίδης}%
\trimIamb{ἐνεγκάτω τις ἔνδοθεν δᾷδ᾽ ἢ λύχνον.\\
ἐπίκυπτε· τὴν κέρκον φυλάττου νυν ἄκραν.
\commentaire{239}{τὴν κέρκον φυλάττου νυν ἄκραν.}{Euripide reprend l'image du δελφάκιον flambé ; il parle bien-sûr de l'\italique{olisbos} de Mnésiloque.}\\\passeLigne}%
\end{groupe}%
%% v. 240-242
\begin{groupe}%
\personnageGrec{Μνησίλοχος}%
\trimIamb{ἐμοὶ μελήσει νὴ Δία, πλήν γ᾽ ὅτι κάομαι.\\
οἴμοι τάλας. ὕδωρ ὕδωρ ὦ γείτονες%
\commentaire{241-2}{οἴμοι τάλας. ὕδωρ ὕδωρ ὦ γείτονες / πρὶν ἀντιλαβέσθαι πρωκτὸν ἓτερον τῆς φλογός.}{Mnésiloque est littéralement en train de prendre feu ; toute sa réplique est un pastiche d'un appel au secours, comme si sa propre maison s'embrasait\marqueur{1} et que celles, contiguës, de ses voisins allaient prendre feu à leur tour (sauf qu'il s'agit ici non d'οἶκός mais de πρωκτός)... il appelle donc à puiser de l'eau « ὕδωρ ! ὕδωρ ! » %, « À l'eau ! Non mais, Allô quoi ! » \\%j'suis un homme et j'ai plus de poils au c... »\\
Horace prodigue ce conseil dans ses \italique{Épîtres} (1, 18, 84) :
\citComm{\latin{Tua res agitur, paries cum priximus ardet}\\
\francais{C'est ton intérêt qui est en jeu, lorsque la maison de ton voisin brûle\marqueur{2}.}}
L'image de l'incendie qui se propage de maison en maison était si fréquente dans le monde latin et roman\marqueur{3} qu'on a peine à croire qu'elle ne le fût pas aussi dans le monde grec -- témoin ce passage d'Aristophane.}%
\debutNote{241-2}%
\note{1}{La scholie au vers 242 : (πρωκτὸν·) ὡς εἰ ἔλεγεν ``οἰκίαν''.}%
\note{2}{J'emprunte la traduction au magistral \italique{Dictionnaire des sentences latines et grecques} de Renzo Tosi (traduit de l'italien par Rebecca Lenoir). La consultation de ce dictionnaire me fut d'une grande aide pour découvrir, en ce passage comme en d'autres, des ponts entre la langue d'Aristophane et l'univers des proverbes.}%
\note{3}{Voir le \italique{Dictionnaire} sus-cité (§524).}\\
πρὶν ἀντιλαβέσθαι πρωκτὸν ἓτερον τῆς φλογός.\apparat{242}{\lecon{[ἓ]τερον} \PsiUn \ (restituit Medaglia) : ἓτερον om. R}%
\\\passeLigne}%
\end{groupe}%
%% v. 243a
\begin{groupe}%
\personnageGrec{Εὐριπίδης}%
\trimIamb{θάρρει.\\\passeLigne}%
\end{groupe}%
%% v. 243b
\begin{groupe}%
\personnageGrec{Μνησίλοχος}%
\trimIamb{\ligne{243}\decalage{θάρρει.}τί θαρρῶ καταπεπυρπολημένος;
\commentaire{243}{θάρρει — τί θαρρῶ;}{Plaute, \italique{Curculio}, IV, 3, 21 : \latin{bellator vale — quid valeam ?}}\\\passeLigne}%
\end{groupe}%
%% v. 244-245a
\begin{groupe}%
\personnageGrec{Εὐριπίδης}%
\trimIamb{ἀλλ᾽ οὐκ ἔτ᾽ οἰδὲν πρᾶγμά σοι· τὰ πλεῖστα γὰρ%\testimonia{244-5}{(τὰ πλεῖστα -- ἀσβόλου) Photius α 2946 $=$ \SynB α 2213 ; \AelDion α 186}
\\
ἀποπεπόνηκας.\\\passeLigne}%
\end{groupe}%
%% v. 245b-246
\begin{groupe}%
\personnageGrec{Μνησίλοχος}%
\trimIamb{\ligne{245}\decalage{ἀποπεπόνηκας.}φεῦ ἰοὺ τῆς ἀσβόλου.%\testimonia{245}{(φεῦ -- ἀσβόλου) Suidas α 4132 \espTest (ἄσβολος $=$ ἀσβόλη) Moeris α 17 $\simeq$ Hesychius α 7634 ; cf. Phrynicus \italique{Praep. soph} 28.1 (ἄσβολος fem.)}
\\%\apparat{245}{\lecon{φεῦ} R, Suid. α 4132 : φῦ Dindorf, Wilamowitz \italique{Lysistr.} v.294 \ligneApp{246} γεγέννημαι R : \lecon{γεγένημαι} Suid. τ 903}\\
αἰθὸς γεγένημαι πάντα τὰ περὶ τὴν τράμιν.%\testimonia{246}{Suidas τ 903 (sed τὸν pro τὴν) \espTest (αἰθὸς γεγένημαι) \Eustathe 83.6, 992.60, 1314.15 \espTest (αἰθὸς) \Eustathe 1385.59 \espTest (τράμιν) Hesychius τ 1243 ; Photius τ 415}
\apparat{246}{γεγέννημαι R : \lecon{γεγένημαι} Suid. τ 903}\\\passeLigne}%
\end{groupe}%
%% v. 247
\begin{groupe}%
\personnageGrec{Εὐριπίδης}%
\trimIamb{μὴ φροντίσῃς· ἕτερος γὰρ αὐτὰ σπογγιεῖ.\apparat{247}{\lecon{σπογγιεῖ} R : σφογγιεῖ Elmsley ad \italique{Ach.} 463}%
\commentaire{247}{μὴ φροντίσῃς· ἕτερος γὰρ αὐτὰ σπογγιεῖ.}{Les Grecs se servaient d'éponges pour la toilette et, dans un usage détourné, pouvaient même les utiliser en guise de torche-cul comme cela ressort d'une scène des \italique{Grenouilles} (vv. 481 sqq.). %: %
%\citComm{Διόνυσος — ἀλλ᾽ ὡρακιῶ.\\
%ἀλλ᾽ οἶσε πρὸς τὴν καρδίαν μου σφογγιάν.\\
%Ξανθίας — ἰδοὺ λαβέ, προσθοῦ.\\
%Διόνυσος — ποῦ 'στιν;
%Ξανθίας — ὦ χρυσοῖ θεοὶ\\
%ἐνταῦθ᾽ ἔχεις τὴν καρδίαν;\\
%Διόνυσος — δείσασα γὰρ\\
%ἐς τὴν κάτω μου κοιλίαν καθείρπυσεν.\\
%Ξανθίας — ὦ δειλότατε θεῶν σὺ κἀνθρώπων.\\
%Διόνυσος —                ἐγώ;\\
%πῶς δειλὸς ὅστις σπογγιὰν ᾔτησά σε;\\
%οὐκ ἂν ἕτερός γ’ αὔτ’ ἠργάσατ’ ἀνήρ.\\
%Ξανθίας —                ἀλλὰ τί;  \\
%Διόνυσος — κατέκειτ’ ἂν ὀσφραινόμενος, εἴπερ δειλὸς ἦν· \\
%ἐγὼ δ’ ἀνέστην καὶ προσέτ’ ἀπεψησάμην.\\
%Dionysos — Je me sens défaillir ; apporte-moi une éponge à m'appliquer sur le cœur. — Tiens, prends, applique-là. — Où est-elle ? — Dieux d'or ! c'est là que tu loges ton cœur ? — La peur l'a fait descendre dans mon bas-ventre. — }%
% Retournons à notre vers : 
Qui est cet ἕτερος ? Réponse : le voisin d'à côté : « La situation est bien celle qu'on attend : en cas d'incendie, les voisins, soucieux de leur propre intérêt, commençaient par faire la chaîne pour apporter de l'eau. Puis, le feu une fois éteint, il se trouvait presque toujours un voisin complaisant pour aider le propriétaire (ou le locataire) malchanceux à nettoyer sa maison et à y remettre de l'ordre. » (Taillardat\marqueur{1}.)%
%\eff{Ranae Ἐγώ;  (486)
%	Πῶς δειλὸς ὅστις σπογγιὰν ᾔτησά σε;
%	Οὐ τἂν ἕτερός γ᾽ αὔτ᾽ εἰργάσατ᾽ ἀνήρ.}
}%
\debutNote{247}%
\note{1}{Jean Taillardat, «~Aristophane, \italique{Thesmophories} v. 247, 289, 721 sqq. : essai d'interprétation du texte transmis~», in \italique{Revue des Études Grecques}, t. 85, fasc. 406-408, juillet-décembre 1972, p. XVII.
Le grand Taillardat avait donc compris le sens de ce passage avant même que le \nieme{242} vers, lacunaire, ne fût rétabli par la découverte d'un papyrus. 
%Il n'en demeure pas moins encore énigmatique : dans les \italique{Grenouilles}, Dionysos s'écrie au vers 487, indigné qu'on l'ait pu traiter de lâche : Πῶς δειλὸς ὅστις σπογγιὰν ᾔτησά σε; / Οὐ τἂν ἕτερός γ᾽ αὔτ᾽ εἰργάσατ᾽ ἀνήρ. L'éponge et le courage semblent proverbialement liés (cf. μελάμπυγος). 
% Pas énigmatique : Cf. Willems !!
}%
\\\passeLigne}%
\end{groupe}%
%% v. 248
\begin{groupe}%
\personnageGrec{Μνησίλοχος}%
\trimIamb{οἰμώξετ᾽ ἆρ᾽ εἴ τις τὸν ἐμὸν πρωκτὸν πλυνεῖ.%\testimonia{248}{(πλυνεῖ) Moeris π 79}
\apparat{248}{οἰμώζετ᾽ R : \lecon{οἰμώξετ᾽} Zanetti \espApp εἰς R : \lecon{εἴ τις} Brunck}%
\commentaire{248}{οἰμώξετ᾽ ἄρ᾽ εἴ τις τὸν ἐμὸν πρωκτὸν πλυνεῖ.}{πλύνειν signifie au propre \italique{nettoyer} (et fait donc écho au σπογγιεῖ du vers précédent) et par suite \italique{dire du mal} (\italique{passer un savon} en français ; l'on manque toujours de la délicatesse requise lorsqu'on s'occupe de la toilette d'autrui).
Le sel de la boutade de Mnésiloque nous échappe... 
%Contrairement à ce qu'ont affirmé nombre de commentateurs \& traducteurs, 
% Ce n'est pas au sens dérivé mais bien au sens propre qu'il faut prendre ici πλύνειν. %(On se demande bien pourquoi certains médisants s'aviseraient de dire du mal du derrière de Mnésiloque et plus encore pourquoi ce dernier, qui n'est pas un être d'une grande finesse, s'en offusquerait.)
Mnésiloque souhaite-t-il bien du courage à qui voudrait \italique{s'attaquer} -- au prime sens de πλύνειν : \italique{commencer de nettoyer} -- à son fondement ? %, dont la saleté égale celle des écuries d'Augias, et dont le \italique{travail} -- tel celui de l'hydre de Lerne -- serait à recommencer sitôt abattu !..
Ingrate tâche qui serait aussitôt à recommencer. Travail de Sisyphe qu'on pourrait rapprocher du proverbe grec πλίνθον πλύνειν, « laver une brique »\marqueur{}. Dans le même genre d'idée, Aristophane emploie au vers 604 des \italique{Guêpes} l'expression qui paraît proverbiale πρωκτὸς λουτροῦ περιγιγνόμενος, « le c.. l'emporte sur l'eau lustrale »\marqueur{2}.
}%
\debutNote{248}%
\note{1}{Voir Photius ο 360 (\italique{s.v} ὄνου πόκοι), ο 363 (ὄνου πόκαι) et l'apparat conséquent des \latin{testimonia} et autres \latin{loci paralleli} de l'édition Theodoridis. Le proverbe fut répertorié entre autres par Plutarque, \italique{Corpus paroemiographorum Graecorum}, I, περὶ τῶν ἀδυνάτων ; Diogenianus (VII, 50), Zenobius (VI, 48). Le scholiaste des \italique{Guêpes} (280b) le cite : <λίθον ἕψεις·> ὡς καὶ τὸ `\textbf{πλίνθον πλύνεις}' καὶ `χύτραν ποικίλλεις' καὶ `εἰς ὕδωρ γράφεις' καὶ `Αἰθίοπα λευκαίνεις' καὶ `κατὰ θαλάττης σπείρεις' καὶ τὰ τοιαῦτα. C'est d'ailleurs en écho à ce proverbe qu'Aristophane rapproche au vers 514 du \italique{Ploutos} les verbes πλινθουργεῖν et πλύνειν.}%
\note{2}{Bdélycléon lançait cette sentence à la barbe de Philocléon, saoul de discours et qui, quelque effort qu'il fît pour soigner sa glossolalite, n'en retomberait pas moins un jour dans son travers.
Le proverbe devait donc peu ou prou signifier qu'un naturel grossier, le πρωκτός, retournait inévitablement à sa fange ; en français : un \italique{indécrottable}. Il est repris tel quel par les parémiographes Diogenianus (III, 53) et Michel Apostolius (XIV, 78), Hésychius (π 4144), Photius (π 1450), Suidas (π 2951). La source principale doit être d'anciens commentaires au vers 604 des \italique{Guêpes}. (Cf. les scholies aux manuscrits d'Aristophane ; l'on y apprend que Καλλίστρατος lui-même avait signalé l'expression comme proverbiale.)% L'exégèse est de mon cru ; les anciens commentateurs n'ayant proposé que des explications pour le moins étranges..
}\\\passeLigne}%
\end{groupe}%
%% v. 249-251
\begin{groupe}%
\personnageGrec{Εὐριπίδης}%
\trimIamb{Ἀγάθων, ἐπειδὴ σαυτὸν ἐπιδοῦναι φθονεῖς,\\
ἀλλ᾽ ἱμάτιον γοῦν χρῆσον ἡμῖν τουτῳὶ\apparat{250}{τουτωί R : \lecon{τουτῳὶ} edd. plerique}\\
καὶ στρόφιον· οὐ γὰρ ταῦτά γ᾽ ὡς οὐκ ἔστ᾽ ἐρεῖς.\\\passeLigne}%
\end{groupe}%
\endnumbering
\end{Leftside}
\begin{Rightside}
\beginnumbering
%% v. 208a
\begin{groupe}%
\personnageFr{Euripide}%
\trimIambFr{Eh bien ? le feras-tu ?\\\passeLigne}%
\end{groupe}%
%% v. 208b
\begin{groupe}%
\personnageFr{Agathon}%
\trimIambFr{\ligne{208}\espace{3.5cm}N'y compte pas.\\\passeLigne}%
\end{groupe}%
%% v. 209
\begin{groupe}%
\personnageFr{Euripide}%
\trimIambFr{Ô triple malheureux, comment périt Euripide.\\\passeLigne}%
\end{groupe}%
%% v. 210
\begin{groupe}%
\personnageFr{Mnésiloque}%
\trimIambFr{Mon ami, mon gendre, n'abandonne point.\\\passeLigne}%
\end{groupe}%
%% v. 211a
\begin{groupe}%
\personnageFr{Euripide}%
\trimIambFr{Que puis-je faire ?\\\passeLigne}%
\end{groupe}%
%% v. 211b-212
\begin{groupe}%
\personnageFr{Mnésiloque}%
\trimIambFr{\ligne{211}\espace{3.4cm}Envoie-le se faire\\
pendre, lui, et dispose de moi comme tu voudras.\\\passeLigne}%
\end{groupe}%
%% v. 213-214a
\begin{groupe}%
\personnageFr{Euripide}%
\trimIambFr{Voyons, puisque tu te remets entre mes mains,\\
ôte ce manteau.\\\passeLigne}%
\end{groupe}%
%% v. 214b-215a
\begin{groupe}%
\personnageFr{Mnésiloque}%
\trimIambFr{\ligne{214}\espace{3.7cm}Le voilà par terre.\\
Mais que penses-tu faire de moi ?\\\passeLigne}%
\end{groupe}%
%% v. 215b-216a
\begin{groupe}%
\personnageFr{Euripide}%
\trimIambFr{\ligne{215}\espace{3.7cm}Te raser cela\testimonia{215-6}{(ἀποξυρεῖν -- ἀφεύειν) Suidas α 4565 (ἀφαύειν pro ἀφεύειν), cf. Scholium ταῦτα δὲ ἔλαβεν ἐκ τῶν Ἰδαίων Κρατίνου (fr. 90 Kassel-Austin) ; \ClemAlex, \italique{Στρώματα} VI 26.4 Ἀριστοφάνης δὲ ὁ κωμικὸς ἐν ταῖς πρώταις Θεσμοφοριαζούσαις τὰ ἐκ τῶν Κρατίνου Ἐμπιπραμένων μετήνεγκεν ἔπη.}\\
et te flamber par dessous.\\\passeLigne}%
\end{groupe}%
%% v. 216b-217
\begin{groupe}%
\personnageFr{Mnésiloque}%
\trimIambFr{\ligne{216}\espace{3cm}Eh bien, fais à ta guise ;\\
autrement je n'aurais jamais dû me mettre à ta merci.%\commentaire{217}{}{Cf. sur le texte de ce passage, v. 217, Bull. cité, note 4, ci-après.}.
\\\passeLigne}%
\end{groupe}%
%% v. 218-219a
\begin{groupe}%
\personnageFr{Euripide}%
\trimIambFr{Agathon, pour sûr, tu portes toujours sur toi des rasoirs.
Prête-moi donc un peu un rasoir.\testimonia{219-20}{(χρῆσόν -- ξυροδόκης) Pollux X 140 (νῦν omisso)}\\\passeLigne}%
\end{groupe}%
%% v. 219b-220a
\begin{groupe}%
\personnageFr{Agathon}%
\trimIambFr{\ligne{219}\espace{3.8cm}Prends-le toi-même,\\
là, dans l'étui.\testimonia{220}{(ξυροδόκης) Pollux II 32 ; \Eustathe 130.43, 1399.38 ; cf. Scholium}\\\passeLigne}%
\end{groupe}%
%% v. 220b-221
\begin{groupe}%
\personnageFr{Euripide}%
\trimIambFr{\ligne{220}\espace{4cm}Tu es galant homme. 
(À Mnésiloque.) Assieds-toi. Enfle la joue droite. (Il le rase.)\\\passeLigne}%
\end{groupe}%
%% v. 222a
\begin{groupe}%
\personnageFr{Mnésiloque}%
\trimIambFr{Aïe !\\\passeLigne}%
\end{groupe}%
%% v. 222b-223a
\begin{groupe}%
\personnageFr{Euripide}%
\trimIambFr{\ligne{222}\espace{1cm}Qu'as-tu à crier ? Je t'enfonce un bâillon\\
si tu ne te tais.\\\passeLigne}%
\end{groupe}%
%% v. 223b
\begin{groupe}%
\personnageFr{Mnésiloque}%
\trimIambFr{\ligne{223}\espace{2.2cm}Aïe, aïe, aïe ! aie, aïe, aïe !%\testimonia{223}{(ἀτταταῖ ἰατταταῖ) Suidas ε 2807}
\\\passeLigne}%
\end{groupe}%
%% v. 224a
\begin{groupe}%
\personnageFr{Euripide}%
\trimIambFr{Holà ! où cours-tu ?\\\passeLigne}%
\end{groupe}%
%% v. 224b-226a
\begin{groupe}%
\personnageFr{Mnésiloque}%
\trimIambFr{\ligne{224}\espace{2.6cm}Au temple des augustes\\
déesses. Car, par Déméter, je ne resterai pas là à me faire
taillader.\\\passeLigne}%
\end{groupe}%
%% v. 226b-227
\begin{groupe}%
\personnageFr{Euripide}%
\trimIambFr{\ligne{226}\espace{1.7cm}Mais tu seras grotesque\\
avec ta moitié de figure dégarnie de poil.\testimonia{227}{Suidas η 334 \espTest (ἡμίκραιραν) Hesychius η 496 ἡμίκραιρα ; \EtymGen (AB) 150 ed. Miller ; \Eustathe 710.50, 1127.35 ; \AelDion η 9}\\\passeLigne}%
\end{groupe}%
%% v. 228a
\begin{groupe}%
\personnageFr{Mnésiloque}%
\trimIambFr{Peu m'en chaut.\\\passeLigne}%
\end{groupe}%
%% v. 228b-229a
\begin{groupe}%
\personnageFr{Euripide}%
\trimIambFr{\ligne{228}\espace{2.7cm}Au nom des dieux,\\
ne m'abandonne pas. Arrive ici.\\\passeLigne}%
\end{groupe}%
%% v. 229b
\begin{groupe}%
\personnageFr{Mnésiloque}%
\trimIambFr{\ligne{229}\espace{3.6cm}Misère de moi !\\\passeLigne}%
\end{groupe}%
%% v. 230
\begin{groupe}%
\personnageFr{Euripide}%
\trimIambFr{Tiens-toi coi et lève la tête. Où te tournes-tu?\\\passeLigne}%
\end{groupe}%
%% v. 231a
\begin{groupe}%
\personnageFr{Mnésiloque}%
\trimIambFr{Mumû.\\\passeLigne}%
\end{groupe}%
%% v. 231b
\begin{groupe}%
\personnageFr{Euripide}%
\trimIambFr{\ligne{231}\espace{1.1cm}Pourquoi mumû ? Tout est fini, et selon les règles.\testimonia{231}{\EtymGen (AB) 218 ed. Miller \espTest (μῦ μῦ) Suidas ε 2807 ; cf. \Eustathe 723.9 \espTest (τί -- καλῶς) Suidas μ 1383 ; \Apostol XVI 60c \espTest (μύζεις) \Eustathe 440.23}\\\passeLigne}%
\end{groupe}%
%% v. 232
\begin{groupe}%
\personnageFr{Mnésiloque}%
\trimIambFr{Ah malheur ! Je vais resservir dans l'infanterie \fin légère.\testimonia{232}{Suidas ψ 105}\\\passeLigne}%
\end{groupe}%
%% v. 233-234a
\begin{groupe}%
\personnageFr{Euripide}%
\trimIambFr{Ne t'inquiète pas. Tu auras un air tout avenant.\\
Veux-tu te mirer ? (Il lui tend un miroir.)\\\passeLigne}%
\end{groupe}%
%% v. 234b
\begin{groupe}%
\personnageFr{Mnésiloque}%
\trimIambFr{\ligne{234}\espace{3.5cm}Si c'est ton plaisir. Passe.\\\passeLigne}%
\end{groupe}%
%% v. 235a
\begin{groupe}%
\personnageFr{Euripide}%
\trimIambFr{Te vois-tu ?\\\passeLigne}%
\end{groupe}%
%% v. 235b
\begin{groupe}%
\personnageFr{Mnésiloque}%
\trimIambFr{\ligne{235}\espace{2.4cm}Non, pas Zeus. Je vois Clisthène.\testimonia{235}{Suidas κ 1756 Κλεισθένην ὁρῶ}\\\passeLigne}%
\end{groupe}%
%% v. 236
\begin{groupe}%
\personnageFr{Euripide}%
\trimIambFr{Lève-toi, que je te flambe. Tiens-toi penché.\\\passeLigne}%
\end{groupe}%
%% v. 237
\begin{groupe}%
\personnageFr{Mnésiloque}%
\trimIambFr{Ah malheur ! Je vais être cochon de lait.\testimonia{237}{Suidas δ 205 (sed γέγονα pro γενήσομαι)}\\\passeLigne}%
\end{groupe}%
%% v. 238-239
\begin{groupe}%
\personnageFr{Euripide}%
\trimIambFr{Qu'on m'apporte de dedans une torche ou une lampe.\\
Courbe-toi, et gare au bout de la queue.\\\passeLigne}%
\end{groupe}%
%% v. 240-242
\begin{groupe}%
\personnageFr{Mnésiloque}%
\trimIambFr{J'y veillerai, par Zeus... N'empêche que je brûle.\\
Ah misère ! de l'eau, de l'eau, voisins,\\
avant que la fesse d'à côté ne s'enflamme.\\\passeLigne}%
\end{groupe}%
%% v. 243a
\begin{groupe}%
\personnageFr{Euripide}%
\trimIambFr{Du sang froid.\\\passeLigne}%
\end{groupe}%
%% v. 243b
\begin{groupe}%
\personnageFr{Mnésiloque}%
\trimIambFr{\ligne{243}\espace{1.4cm}Du sang froid ? quand je suis dévoré par les flammes !..\\\passeLigne}%
\end{groupe}%
%% v. 244-245a
\begin{groupe}%
\personnageFr{Euripide}%
\trimIambFr{Mais tu n'as plus à t'inquiéter. Car te voilà dépêtré\testimonia{244-5}{(τὰ πλεῖστα -- ἀσβόλου) Photius α 2946 $=$ \SynB α 2213 ; \AelDion α 186}\\
du principal.\\\passeLigne}%
\end{groupe}%
%% v. 245b-246
\begin{groupe}%
\personnageFr{Mnésiloque}%
\trimIambFr{\ligne{245}\espace{2.3cm}Phu ! Ah! cette suie !\testimonia{245}{(φεῦ -- ἀσβόλου) Suidas α 4132 \espTest (ἄσβολος $=$ ἀσβόλη) Moeris α 17 $\simeq$ Hesychius α 7634 ; cf. Phrynicus \italique{Praep. soph} 28.1 (ἄσβολος fem.)}\\
Je suis tout feu dans l'entrefesson.\testimonia{246}{Suidas τ 903 (sed τὸν pro τὴν) \espTest (αἰθὸς γεγένημαι) \Eustathe 83.6, 992.60, 1314.15 \espTest (αἰθὸς) \Eustathe 1385.59 \espTest (τράμιν) Hesychius τ 1243 ; Photius τ 415}\\\passeLigne}%
\end{groupe}%
%% v. 247
\begin{groupe}%
\personnageFr{Euripide}%
\trimIambFr{Ne t'inquiète pas. Le voisin d'à côté t'épongera cela.\\\passeLigne}%
\end{groupe}%
%% v. 248
\begin{groupe}%
\personnageFr{Mnésiloque}%
\trimIambFr{Il en coûtera à celui qui s'attaquera à mon derrière.\testimonia{248}{(πλυνεῖ) Moeris π 79}\\\passeLigne}%
\end{groupe}%
%% v. 249-251
\begin{groupe}%
\personnageFr{Euripide}%
\trimIambFr{Agathon, puisque tu refuses de te dévouer à moi,\\
prête-nous du moins un vêtement\\
et un strophion. Tu ne nieras pas que tu en possèdes..\\\passeLigne}%
\end{groupe}%
\endnumbering
\end{Rightside}
\end{pages}
\Pages
\clearpage\ifodd\value{page}\hbox{}\newpage\fi
\enTeteG{ΠΡΟΛΟΓΟΣ --- Codex Ravennas 429, f. 165v (Classense 167v) --- vv. 252–294}
\begin{adjustwidth}{-1.5cm}{-3.0cm}
\begin{adjustbox}{max totalheight=21.0cm, max width=16.5cm, keepaspectratio}
\begin{minipage}{16.5cm}
{\fontsize{7.2pt}{8.5pt}\selectfont
\noindent\textbf{Folio 07 --- Codex Ravennas 429, \texttt{f. 165v} (\texttt{Classense 167v}) : ΠΡΟΛΟΓΟΣ, vv. 252–294}\hfill \textit{Témoins :} $\text{R}^1, \Rpc, \Rdeux$, \PsiUn\par\vspace{2pt}
\hrule height 0.4pt\vspace{2.5pt}
}
{\fontsize{8.1pt}{9.5pt}\selectfont
\ligneRav{252}{ἀγάθ/}{λαμβάνετε καὶ χρῆσθ᾽· οὐ φθονῶ. : τί οὖν λάβω;}{\textbf{252} (Εὐρ.) R : \lecon{(Μνη.)} Bentley}
\ligneRav{253}{\auDessus{εὐ}{ρ}}{ὅ τι; τὸν κροκωτὸν πρῶτον ἐνδύου λαβών·}{\textbf{253} (Ἀγάθ.) R : \lecon{(Εὐρ.)} Dobree}
\ligneRav{254}{μ\auDessus{ν}{Η}}{νὴ τὴν Ἀφροδίτην ἡδύ γ᾽ ὄζει ποσθίου·}{\textbf{254} (Εὐρ.) R}
\ligneRav{255}{ἀγάθ/}{σύζωσον ἀνύσας. : αἶρε νῦν στρόφιον. : ἰδού·}{\textbf{255} (Ἀγάθ.) et dicolon ante αἶρε praefixit R\espApp (Εὐρ.) ἰδού — (Μνη.) ἴθι ... σκέλη R}
\ligneRav{256}{μ\auDessus{ν}{Η}}{ἴθι νυν κατάστειλόν με τὰ περὶ τὼ \difference{σκέλει}·}{\textbf{256} σκέλη R : \lecon{σκέλει} van Leeuwen}
\ligneRav{}{\auDessus{εὐ}{ρ}}{κεκρυφάλου δεῖ καὶ μίτρας. : ἡδὶ μὲν οὖν}{}
\ligneRav{}{}{κεφαλὴ περίθετος, ἣν ἐγὼ νύκτωρ φορῶ·}{}
\ligneRav{}{\auDessus{εὐ}{ρ}}{νὴ τὸν Δί᾽ ἀλλὰ κἀπιτηδεία πάνυ·}{}
\ligneRav{260}{μ\auDessus{ν}{Η}}{\difference{ἆρ᾽ ἁρμόσει} μοι; : νὴ Δί᾽ ἀλλ᾽ ἄριστ᾽ ἔχει·}{\textbf{260} ἠραρμοσηι R : \lecon{ἆρ᾽ ἁρμόσει} Küster\espApp (Ἀγάθ.) R : \lecon{(Εὐρ.)} Ellebodius}
\ligneRav{261}{—}{φέρ᾽ ἔγκυκλον· : \difference{τουτὶ} \difference{λάβ᾽} ἀπὸ τῆς κλινίδος·}{\textbf{261} τοῦτὶ R : \lecon{τουτὶ} Suid. ε 135, κ 1809 \espApp λάμβαν᾽ R, Suid. : \lecon{λάβ᾽} Daubuz, Küster}
\ligneRav{}{\auDessus{εὐ}{ρ}}{ὑποδημάτων δεῖ. : τἀμὰ ταυτὶ λάμβανε·}{}
\ligneRav{263}{μ\auDessus{ν}{Η}}{χὤπως μὲν ἀνδρείως· : Ἀγάθωνος γ᾽ ο\auDessus{ῦ}{ν} ἐστὶν·}{\textbf{263} μοι R : \lecon{μοι;} Küster\espApp \lecon{(Ἀγάθ.)} χαλαρὰ ... φορῶν R : (Μνη.) Brunck, edd. post. \espApp γ᾽ο\auDessus{ῦ}{ν} \Rpc \espApp χαίρεις R, Suid. A,G,M : \lecon{χαίροις} Suid. F,V}
\ligneRav{263}{μ\auDessus{ν}{Η}}{χὤπως μὲν ἀνδρείως· : Ἀγάθωνος γ᾽ ο\auDessus{ῦ}{ν} ἐστὶν·}{\textbf{263} μοι R : \lecon{μοι;} Küster\espApp \lecon{(Ἀγάθ.)} χαλαρὰ ... φορῶν R : (Μνη.) Brunck, edd. post. \espApp γ᾽ο\auDessus{ῦ}{ν} \Rpc \espApp χαίρεις R, Suid. A,G,M : \lecon{χαίροις} Suid. F,V}
\ligneRav{264}{}{εἴσω τις ὡς τάχιστά μ᾽ ἐσκυκλησάτω·}{\textbf{264} nulla nota mut. pers. R : \lecon{(Ἀγάθ.)} ἀλλ᾽ ... Bibi}
\ligneRav{266}{\auDessus{εὐ}{ρ}}{\difference{ἁνὴρ} μὲν ἡμῖν οὑτοσὶ καὶ δὴ γυνὴ}{\textbf{266} ἀνὴρ R : \lecon{ἁνὴρ} Fritzsche}
\ligneRav{267}{}{\difference{τό γ᾽} εἶδος· ἢν λαλῇς δ᾽, ὅπως τῷ φθέγματι}{\textbf{267} τόδ᾽ R : \lecon{τό γ᾽} Küster}
\ligneRav{}{—}{γυναικιεῖς εὖ καὶ πιθανῶς. : πειράσομαι·}{}
\ligneRav{}{\auDessus{εὐ}{ρ}}{βάδιζε τοίνυν. : μὰ τὸν Ἀπόλλω οὔκ, ἤν γε μὴ}{}
\ligneRav{270}{—}{ὀμόσῃς ἐμοί— : τί χρῆμα; : συσσώσειν ἐμὲ}{}
\ligneRav{}{}{πάσαις τέχναις, ἤν μοί τι περιπίπτῃ κακόν·}{}
\ligneRav{}{\auDessus{εὐ}{ρ}}{`ὄμνυμι τοίνυν αἰθέρ᾽ οἴκησιν Διός.'}{}
\ligneRav{273}{μ\auDessus{ν}{Η}}{τί μᾶλλον ἢ τὴν Ὑποκράτους ξυνοικίαν;}{\textbf{273} \lecon{Ὑποκράτους} R : Ἱπποκράτους Σ}
\ligneRav{}{\auDessus{εὐ}{ρ}}{ὄμνυμι τοίνυν πάντας ἄρδην τοὺς θεούς·}{}
\ligneRav{275}{μ\auDessus{ν}{Η}}{μέμνησο τοίνυν ταῦθ᾽, ὅτι ἡ φρὴν ὤμοσεν,}{}
\ligneRav{}{}{ἡ γλῶττα δ᾽ οὐκ ὀμώμοκ᾽· οὐδ᾽ ὥρκωσ᾽ ἐγώ·}{}
\ligneRav{}{}{\hspace{2.2em}\textbf{[παρεπιγραφή R\expo{1}, \PsiUn{} :]} \lecon{ὀλολύζουσι τε ἱερὸν ὠθεῖται}}{\textbf{post 276} \lecon{ὀλολύζουσι τε ἱερὸν ὠθεῖται} R, ολολυ[ \PsiUn, $\Sigma^{\text{R}}$}
\ligneRav{277}{\auDessus{εὐ}{ρ}}{ἔκσπευδε ταχέως· ὡς τὸ τῆς ἐκκλησίας}{\textbf{277} \lecon{ἔκσπευδε} R : εἄ· σπεῦδε Maas}
\ligneRav{278}{}{σημεῖον ἐν τῷ Θεσμοφορίῳ φαίνεται·}{\textbf{278} \lecon{Θεσμοφορίῳ} R, cf. \italique{IG} II\expo{2} 1177.24 ; 2498.12 : Θεσμοφόρῳ Suid. σ 325 : Θεσμοφορείῳ Scaliger}
\ligneRav{}{—}{ἐγὼ δ᾽ ἄπειμι. : δεῦρό νυν ὦ Θρᾷτθ᾽ ἕπου·}{}
\ligneRav{280}{}{ὦ Θρᾷττα θέασαι, καομένων τῶν λαμπάδων}{}
\ligneRav{}{}{ὅσον τὸ χρῆμ᾽ ἀνέρχεθ᾽ ὑπὸ τῆς λιγνύος·}{}
\ligneRav{282}{}{ἀλλ᾽, ὦ \difference{περικαλλεῖ} \difference{Θεσμοφόρω}, δέξασθέ με}{\textbf{282} Θεσμοφόρῳ R : \lecon{Θεσμοφόρω} Photius α 1386\espApp περὶ καλλῆ R : \lecon{περικαλλεῖ} Kaibel}
\ligneRav{283}{}{ἀγαθῇ τύχῃ \difference{καὶ δεῦρο καὶ πάλιν} οἴκαδε·}{\textbf{283} καὶ δεῦρο πάλιν R : \lecon{καὶ δεῦρο καὶ πάλιν} Bentley}
\ligneRav{}{}{ὦ Θρᾷττα τὴν κίστην κάθελε, κᾆτ᾽ ἔξελε}{}
\ligneRav{285}{}{\difference{τὰ πόπαν᾽}, ὅπως λαβοῦσα θύσω ταῖν θεαῖν·}{\textbf{285} \lecon{τὰ πόπαν᾽} \PsiUn : τὸ πόπανον R\espApp \lecon{τὰ πόπαν᾽} \PsiUn, Porson : τὸ πόπανον R}
\ligneRav{}{}{δέσποινα πολυτίμητε Δήμητερ φίλη}{}
\ligneRav{287}{}{καὶ \difference{Φερρέφαττα}, πολλὰ πολλάκις μέ σοι}{\textbf{287} \lecon{Φερρέφαττα} \PsiUn, Thiersch : Φερσέφαττα R}
\ligneRav{}{}{θύειν ἔχουσαν, εἰ δὲ μή, ἀλλὰ νῦν λαθεῖν}{}
\ligneRav{289}{}{καὶ τὴν θυγατέρα \difference{Χοιρίλ'} ἀνδρός μοι τυχεῖν}{\textbf{289} χοῖρον R : χοιρίον Σ : \lecon{Χοιρίλ'} Bibi}
\ligneRav{290}{}{πλουτοῦντος, ἄλλως τ᾽ ἠλιθίου κἀβελτέρου,}{}
\ligneRav{291}{}{καὶ \difference{πρὸς θυλάκιον} νοῦν ἔχειν μοι καὶ φρένας·}{\textbf{291} πρὸς θάληκον R : Ποσθαλίσκον Dindorf : \lecon{πρὸς θυλάκιον} Bibi}
\ligneRav{}{}{ποῦ ποῦ καθίζωμ᾽ ἐν καλῷ, τῶν ῥητόρων}{}
\ligneRav{}{}{ἵν᾽ ἐξακούω; σὺ δ᾽ ἄπιθ᾽ ὦ Θρᾷττ᾽ ἐκποδών·}{}
\ligneRav{}{}{δούλοις γὰρ οὐκ ἔξεστ᾽ ἀκούειν τῶν λόγων·}{}
}
\vspace{3pt}\noindent\rule{4.5cm}{0.4pt}\par\vspace{1.5pt}
{\fontsize{7.0pt}{8.3pt}\selectfont
\noindent\textbf{Loci paralleli (vv. 252–294) ---} \textbf{254} Suidas π 2114 \espTest (ποσθίου) Photius π 1113 πόσθιον· αἰδοῖον· ἢ (sic) τὴν Ἀφροδίτην ; Hesychius π 3096 \espTest \textbf{257-8} (ἡδὶ -- φορῶ) Suidas π 1165 (sed ἥδιον pro ἡδὶ) \espTest \textbf{258} Photius π 690 περίθετος κεφαλή ; \Eustathe 1280.63 \espTest (κεφαλὴ περίθετος) Pollux II 35 \espTest \textbf{261} Suidas ε 135 (sed τουτὶ pro τι) \espTest (ἔγκυκλον -- κλινίδος) Suidas κ 1809 \espTest \textbf{263} Suidas χ 7 χαλαρὰ γοῦν χαίρεις φορῶν ὑποδήματα (sed χαίροις codd. FV) \espTest \textbf{274} (ἄρδην) Hesychius α 7097 \espTest \textbf{275-6} cf. \Eustathe 443.37 \espTest \textbf{277-8} (ὡς -- φαίνεται) Suidas σ 325 (sed Θεσμοφόρῳ pro Θεσμοφορίῳ) \espTest \textbf{282} Photius α 1386 ; \AelDion α 113 \espTest \textbf{285} (πόπαν᾽) \Eustathe 437.4 \espTest \textbf{289} (χοῖρον) Euripides, testimonia 4, 18 Kannicht Χοιρίλη \espTest \textbf{291} (θάληκον) Hesychius θ 50 θάλλικα ; Pollux X 152 θυλάκιον [...] ἐπὶ τοῦ αὐτοῦ δὲ καὶ σακίον, ὅταν φῇ ἐν Θεσμοφοριαζούσαις \espTest \textbf{292} (τῶν ῥητόρων) Photius ρ 106 ῥήτορα· καὶ τὴν γυναῖκα. Ἀριστοφάνης\par
}
\vspace{2.5pt}\noindent\rule{4.5cm}{0.4pt}\par\vspace{1.5pt}
{\fontsize{6.9pt}{8.2pt}\selectfont
\noindent\textbf{Suite (en commentaire) de l'apparat critique \& notes paléographiques (\texttt{f. 165v}) ---} \textbf{Témoin papyrologique :} \textbf{\PsiUn{} (\italique{PSI} XI 1194)} couvre les \textbf{vv. 273–288}. --- \textbf{Troisième \italique{parepigraphè} (post v. 276) :} R écrit dans la colonne \lecon{ὀλολύζουσι τε ἱερὸν ὠθεῖται}, avec une \textbf{\italique{koronis}} dans la marge gauche et la scholie $\Sigma^{\text{R}}$ : παρεπιγραφή. ἐκκυκλεῖται ἐπὶ τὸ ἔξω τὸ θεσμοφόριον. Le papyrus \PsiUn{} confirme ολολυ[. --- \textbf{Abréviation supralinéaire au v. 263 :} γ' ο\auDessus{ῦ}{ν} \Rpc{}. --- \textbf{Deux conjectures majeures de \textbf{Bibi} (vv. 289 et 291) :} Au v. 289, \lecon{Χοιρίλ'} (nom de la fille de Mnésiloque et épouse d'Euripide, corrompu en χοῖρον dans R et χοιρίον dans $\Sigma^{\text{R}}$) ; au v. 291, \lecon{πρὸς θυλάκιον} (ou πρὸς θάλικα ap. Hésychius $\theta\ 50$) pour πρὸς θάληκον de R. --- \textbf{v. 260 :} ἠραρμοσηι R (iotacisme et médivision) : \lecon{ἆρ' ἁρμόσει} Küster \espApp (Ἀγάθ.) R : \lecon{(Εὐρ.)} Ellebodius. --- \textbf{v. 263 :} μοι R : \lecon{μοι;} Küster \espApp \textbf{(Ἀγάθ.) χαλαρὰ ... φορῶν} R (maintenu par \textbf{Bibi} contre (Μνη.) Brunck) \espApp \lecon{γ' ο\auDessus{ῦ}{ν}} \Rpc{} \espApp χαίρεις R, Souda AGM : \lecon{χαίροις} Souda FV. --- \textbf{v. 264 :} nulla nota mut. pers. R : \lecon{(Ἀγάθ.)} ἀλλ' ... \lecon{Bibi}. --- \textbf{v. 273 :} \lecon{Ὑποκράτους} R (maintenu par \textbf{Bibi} : calembour comique sur $\text{ὑπό}$, contre la normalisation Ἱπποκράτους de $\Sigma^{\text{R}}$). --- \textbf{v. 289 :} χοῖρον R (χοιρίον $\Sigma$) : \lecon{Χοιρίλ'} (ou Χοιρίν') \lecon{Bibi} (fille de Mnésiloque et épouse d'Euripide). --- \textbf{v. 291 :} πρὸς θάληκον R (Ποσθαλίσκον Dindorf) : \lecon{πρὸς θυλάκιον} (ou πρὸς θάλικα ap. Hésychius $\theta\ 50$) \lecon{Bibi} (par métathèse).\par
}
\end{minipage}
\end{adjustbox}
\end{adjustwidth}
\newpage
\enTeteD{Codex Ravennas 429 (Biblioteca Classense) --- Fac-similé du f. 165v (vv. 252–294)}
\begin{adjustwidth}{-1.5cm}{-3.0cm}
\begin{center}
\includegraphics[width=16.5cm,height=20.5cm,keepaspectratio]{/workspaces/Thesmophories/dossier_ravennas_thesmophories/folios_individuels_ff162v-177r/07_f165v_classense167v_vv0252-0294.jpg}\par
\vspace{2pt}
{\scriptsize\textbf{Codex Ravennas 429, f. 165v (Classense 167v)} --- ΠΡΟΛΟΓΟΣ, \textbf{vv. 252–294} ($\text{R}^1, \Rpc, \Rdeux$, \PsiUn).}
\end{center}
\end{adjustwidth}
\newpage
\clearpage\ifodd\value{page}\hbox{}\newpage\fi
\enTeteG{ΘΕΣΜΟΦΟΡΙΑΖΟΥΣΑΙ --- Texte grec établi (f. 165v, vv. 252–294)}
\enTeteD{LES THESMOPHORIES --- Traduction française (f. 165v, vv. 252–294)}
\begin{pages}
\begin{Leftside}
\beginnumbering
%% v. 252a
\begin{groupe}%
\personnageGrec{Ἀγάθων}%
\trimIamb{λαμβάνετε καὶ χρῆσθ᾽· οὐ φθονῶ.\\\passeLigne}%
\end{groupe}%
%% v. 252b
\begin{groupe}%
\personnageGrec{Μνησίλοχος}%
\trimIamb{\ligne{252}\decalage{λαμβάνετε καὶ χρῆσθ᾽· οὐ φθονῶ.}τί οὖν λάβω;\apparat{252}{(Εὐρ.) R : \lecon{(Μνη.)} Bentley}\\\passeLigne}%
\end{groupe}%
%% v. 253
\begin{groupe}%
\personnageGrec{Εὐριπίδης}%
\trimIamb{ὅ τι; τὸν κροκωτὸν πρῶτον ἐνδύου λαβών.\apparat{253}{(Ἀγάθ.) R : \lecon{(Εὐρ.)} Dobree}\\\passeLigne}%
\end{groupe}%
%%% v. 254
\begin{groupe}%
\personnageGrec{Μνησίλοχος}%
\trimIamb{νὴ τὴν Ἀφροδίτην ἡδύ γ᾽ ὄζει ποσθίου.%\testimonia{254}{Suidas π 2114 \espTest (ποσθίου) Photius π 1113 πόσθιον· αἰδοῖον· ἢ (sic) τὴν Ἀφροδίτην ; Hesychius π 3096}
\apparat{254}{(Εὐρ.) R}\\\passeLigne}%
\end{groupe}%
%% v. 255a
\begin{groupe}%
\personnageGrec{Ἀγάθων}%
\trimIamb{σύζωσον ἀνύσας.\\\passeLigne}%
\end{groupe}%
%% v. 255b
\begin{groupe}%
\personnageGrec{Εὐριπίδης}%
\trimIamb{\ligne{255}\decalage{σύζωσον ἀνύσας.}αἶρε νῦν στρόφιον.\apparat{255}{(Ἀγάθ.) et dicolon ante αἶρε praefixit R}%
\commentaire{255}{σύζωσον ἀνύσας --- αἶρε νῦν στρόφιον.}{La répartition de ce vers comme de tout ce passage est très changeante d'une édition à une autre.
J'ai essayé pour ma part de respecter le découpage des vers du manuscrit de Ravenne, sans tailler ni fusionner çà et là des répliques qui ne l'étaient pas originellement.
J'ai confié à Euripide les conseils d'habillement, Agathon continuant d'adopter une attitude froide : il se contente de tendre à Mnésiloque les vêtements pour ne pas que ce rustre aille puiser directement dans sa garde-robe.
La principale question qui se soit posée est celle de la première partie du vers 255 : σύζωσον ἀνύσας --- αἶρε νῦν στρόφιον.
Il comporte un \italique{dicôlon} entre ἀνύσας et αἶρε sans pour autant indiquer de changement de locuteur... Ce \italique{dicôlon} est-il là pour signifier une simple pause le temps que Mnésiloque passe la crocote ? ou y a-t-il véritablement changement de locuteur ?
`σύζωσον ἀνύσας' me semble être une réplique au ton sec, marquant la pointe d'agacement qu'Agathon aurait pu sentir à la remarque de Mnésiloque du vers 254. Aussi lui ai-je accordé cette réplique ; puis Euripide de demander le \italique{strophion} (αἶρε νῦν στρόφιον) et Agathon de le tendre à Mnésiloque (ἰδού).}\\\passeLigne}%
\end{groupe}%
%% v. 255c
\begin{groupe}%
\personnageGrec{Ἀγάθων}%
\trimIamb{\ligne{255}\decalage{σύζωσον ἀνύσας. αἶρε νῦν στρόφιον.}ἰδού.\apparat{255-6}{(Εὐρ.) ἰδού — (Μνη.) ἴθι ... σκέλη R}\\\passeLigne}%
\end{groupe}%
%% v. 256
\begin{groupe}%
\personnageGrec{Μνησίλοχος}%
\trimIamb{ἴθι νυν κατάστειλόν με τὰ περὶ τὼ σκέλει.\apparat{256}{σκέλη R : \lecon{σκέλει} van Leeuwen}\\\passeLigne}%
\end{groupe}%
%% v. 257a
\begin{groupe}%
\personnageGrec{Εὐριπίδης}%
\trimIamb{κεκρυφάλου δεῖ καὶ μίτρας.\\\passeLigne}%
\end{groupe}%
%% v. 257b-258
\begin{groupe}%
\personnageGrec{Ἀγάθων}%
\trimIamb{\ligne{257}\decalage{κεκρυφάλου δεῖ καὶ μίτρας.}ἡδὶ μὲν οὖν%\testimonia{257-8}{(ἡδὶ -- φορῶ) Suidas π 1165 (sed ἥδιον pro ἡδὶ)}
\\
κεφαλὴ περίθετος, ἣν ἐγὼ νύκτωρ φορῶ.%\testimonia{258}{Photius π 690 περίθετος κεφαλή ; \Eustathe 1280.63 \espTest (κεφαλὴ περίθετος) Pollux II 35}%
\commentaire{258}{κεφαλὴ περίθετος, ἣν ἐγὼ νύκτωρ φορῶ.}{De tout l'attirail nécessaire pour se métamorphoser en femme, il est étonnant que ce soit le plus évident à priori qui ait rencontré le plus de difficultés : c'est Fritzsche qui le premier s'avisa que la \grec{κεφαλὴ περίθετος} désignait une \italique{perruque}. Hé oui ! car une résille et un bandeau ne suffiraient à cacher l'absence volumique d'une belle crinière, d'où la suggestion avisée d'Agathon – en homme habitué à de tels travestissements. Ce sens est clairement attesté par Pollux (X, 170) :%
\citComm{\grec{πηνήκην δὲ καὶ προκόμιον, καὶ ὡς Ἀριστφάνης ἐν Γηρυτάδῃ περίθεσιν· καὶ περιθέτην δέ, ὡς ἐν Ἄμφιδος Ἀλκμαίωνι καὶ Μενάνδρου Ὀλυνθίᾳ.}}%
Les synonymes \grec{πηνήκη} et \grec{προκόμιον} ne souffrent aucune ambiguïté.\\
Le tacle d'Aristophane envers les mœurs agathoniennes – qui lui fait préciser : ἣν ἐγὼ νύκτωρ φορῶ – n'est pas clair... le scholiaste à l'imagination débridée -- qui n'était probablement pas plus renseigné que nous sur la question -- y voyait une allusion aux chasses masquées contre-nature de notre poètereau dans l'interlope \italique{Athens by night}. Cette rêverie ne trouve qu'un faible écho au vers 35 (et s'appuie peut-être sur une transmission fautive du texte, comme nous en avons fait l'hypothèse \italique{ad loc.}).}\\\passeLigne}%
\end{groupe}%
%% v. 259
\begin{groupe}%
\personnageGrec{Εὐριπίδης}%
\trimIamb{νὴ τὸν Δί᾽ ἀλλὰ κἀπιτηδεία πάνυ.\\\passeLigne}%
\end{groupe}%
%%% v. 260a
\begin{groupe}%
\personnageGrec{Μνησίλοχος}%
\trimIamb{ἆρ᾽ ἁρμόσει μοι;\apparat{260}{ἠραρμοσηι R : \lecon{ἆρ᾽ ἁρμόσει} Küster}\\\passeLigne}%
\end{groupe}%
%%% v. 260b-261a
\begin{groupe}%
\personnageGrec{Εὐριπίδης}%
\trimIamb{\ligne{261}\decalage{ἆρ᾽ ἁρμόσει μοι;}νὴ Δί᾽ ἀλλ᾽ ἄριστ᾽ ἔχει.\apparat{260}{(Ἀγάθ.) R : \lecon{(Εὐρ.)} Ellebodius}\\
φέρ᾽ ἔγκυκλον·%\testimonia{261}{Suidas ε 135 (sed τουτὶ pro τι) \espTest (ἔγκυκλον -- κλινίδος) Suidas κ 1809}%
\commentaire{261}{φέρ᾽ ἔγκυκλον.}{Se reporter à la note au vers 643 pour un développement sur l'\italique{encycle}. Ce mot désigne une sorte de \italique{jupe courte}.}\\\passeLigne}%
\end{groupe}%
%%% v. 261b
\begin{groupe}%
\personnageGrec{Ἀγάθων}%
\trimIamb{\ligne{261}\decalage{φέρ᾽ ἔγκυκλον·}τουτὶ λάβ᾽ ἀπὸ τῆς κλινίδος.\apparat{261}{τοῦτὶ R : \lecon{τουτὶ} Suid. ε 135, κ 1809 \espApp λάμβαν᾽ R, Suid. : \lecon{λάβ᾽} Daubuz, Küster}\\\passeLigne}%
\end{groupe}%
%%% v. 262a
\begin{groupe}%
\personnageGrec{Εὐριπίδης}%
\trimIamb{ὑποδημάτων δεῖ.\\\passeLigne}%
\end{groupe}%
%%% v. 262b
\begin{groupe}%
\personnageGrec{Ἀγάθων}%
\trimIamb{\ligne{262}\decalage{ὑποδημάτων δεῖ.}τἀμὰ ταυτὶ λάμβανε.\\\passeLigne}%
\end{groupe}%
%% v. 263
\begin{groupe}%
\personnageGrec{Μνησίλοχος}%
\trimIamb{ἆρ᾽ ἁρμόσει μοι;\apparat{263}{ μοι R : \lecon{μοι;} Küster}%
\commentaire{263}{ἆρ᾽ ἁρμόσει μοι Ἀγ. χαλαρὰ γ᾽ο\auDessus{ῦ}{ν} χαίρεις φορῶν.}{Mnésiloque, devant les chaussures qu'il trouve étroites --- « Mais m'iront-elles ? » --- Agathon « Tu préférerais (χαίροις Suid.) en chausser de larges... »\\
Rien de plus clair : Agathon « n'aurait eu garde de porter des souliers larges, car c'eût été faire preuve de d'inélégance et même de rusticité [...] Mnésiloque, en vieillard qui aime ses aises, est donc fondé à craindre que les souliers ne soient trop étroits pour lui. »\marqueur{1}. 
Seulement, Brunck eut l'idée saugrenue de confier l'intégralité du vers 263 à Mnésiloque alors que le \italique{Ravennas} indiquait clairement le passage de témoin à Agathon pour la seconde partie, χαλαρὰ γο\auDessus{ῦ}{ν} χαίρεις φορῶν. Ainsi, c'est Mnésiloque qui prononcerait « tu aimes donc à te chausser large »... allusion aux mœurs sodomotiques d'Agathon. En voilà bien assez pour que tous les éditeurs postérieurs suivent en meute Brunck.
Le scholiaste était le premier à avoir ouvert le bal des affreux, qui glosait : <χαλαρὰ·>χαῦνα. διαβάλλει δὲ πάλιν τὸν Ἀγάθωνα ὡς χαῦνον.
Personne ne s'est avisé que si Mnésiloque avait effectivement lancé le trait à sous-entendu obscène qu'on lui impute, alors jamais Agathon ne lui aurait répondu avec bienveillance `σὺ τοῦτο γίγνωσκ᾽', « à toi de t'en aviser ».
Mais cette engeance de fouines à ordures perd tout sens critique dès qu'elle se met à flairer du graveleux.
%Il fut malheureusement suivi par tous les éditeurs postérieurs, qui ne se sont pas rendus compte du contre-sens qu'ils faisaient. En aucun cas Mnésiloque ne peut reprocher à Agathon de se chausser large :  Le scholiaste est le premier à avoir ouvert le bal des affreux, qui glose <χαλαρὰ·>χαῦνα. διαβάλλει δὲ πάλιν τὸν Ἀγάθωνα ὡς χαῦνον. Il ne s'est pas avisé que jamais Agathon n'aurait répondu à Mnésiloque si ce dernier avait lancé le trait à sous-entendu obscène qu'on lui impute.% : « tu aimes à te chausser large ». %Il l'aurait ignoré comme à son accoutumée. 
%Mais cette engeance de fouines à ordures perd tout sens critique dès qu'elle se met à flairer du graveleux.\\
Willems l'avait noté : Mnésiloque se trouve désappointé devant l'étroitesse des élégants souliers d'Agathon ; cette paire lui ira-t-elle : ἆρ᾽ ἁρμόσει μοι; se demande-t-il.
À ces mots, Agathon ne se sent pas de joie : ce rustre en préférerait sans doute une bien large ; c'est bien Agathon qui prononce `χαλαρὰ γοῦν χαίροις φορῶν'.
Et c'est à ce moment seulement que Mnésiloque saisit la perche ; il taquine une nouvelle fois Agathon sur ses mœurs supposées :%
\citComm{Μνησίλοχος --- σὺ τοῦτο γίγνωσκ᾽.\\
Mnésiloque --- Ah ça ! tu en sais quelque chose...}%
C'en est trop pour la patience d'Agathon : excédé, il décide de se retirer superbement avec l'eccyclème (exactement comme autrefois l'Euripide des \italique{Acharniens}).%
%\eff{En fait la correction est la bonne : Agathon est moqué ailleurs (Gérutades 178) comme mh xalara, c'est-à-dire non serré. Évidemment, les allusions scabreuses au trou-du-cul d'Agathon ne sont que pure imagination des fouines à ordure. Zénob. I 2 (cf. Diogenianus I 7)}
}%
\debutNote{263}%
\note{1}{Se reporter encore une fois au grand Alphonse Willems ; la citation est tirée de sa \nieme{6} note, reproduite à la fin.}%
\\\passeLigne}%
\end{groupe}%
%% v. 263b
\begin{groupe}%
\personnageGrec{Ἀγάθων}%
\trimIamb{\ligne{262}\decalage{ἆρ᾽ ἁρμόσει μοι;}χαλαρὰ γοῦν χαίροις φορῶν.%\testimonia{263}{Suidas χ 7 χαλαρὰ γοῦν χαίρεις φορῶν ὑποδήματα (sed χαίροις codd. FV)}
\apparat{263}{\lecon{(Ἀγάθ.)} χαλαρὰ ... φορῶν R : (Μνη.) Brunck, edd. post. \espApp γ᾽ο\auDessus{ῦ}{ν} \Rpc \espApp χαίρεις R, Suid. A,G,M : \lecon{χαίροις} Suid. F,V \ligneApp{264} nulla nota mut. pers. R : \lecon{(Ἀγάθ.)} ἀλλ᾽ ... Bibi}\\\passeLigne}%
\end{groupe}%
%% v. 264
\begin{groupe}%
\personnageGrec{Μνησίλοχος}%
\trimIamb{\ligne{263}σὺ τοῦτο γίγνωσκ᾽\\\passeLigne}%
\end{groupe}%
%% v. 264b-265
\begin{groupe}%
\personnageGrec{Ἀγάθων}%
\trimIamb{\decalage{σὺ τοῦτο γίγνωσκ᾽}ἀλλ᾽ ἔχεις γὰρ ὧν δέει,\\
εἴσω τις ὡς τάχιστά μ᾽ ἐσκυκλησάτω.%
\commentaire{265}{εἴσω τις ὡς τάχιστά μ᾽ ἐσκυκλησάτω.}{De même que l'Euripide des \italique{Acharniens} (vers 409) demande à être roulé par l'eccyclème pour sortir de la scène : il ne lui agrée pas de descendre (καταβαίνειν, cf. \italique{infra} v. 1215) :%
\citComm{ἀλλ' ἐκκυκλήσομαι· καταβαίνειν δ᾽ οὐ σχολή.}%
}\\\passeLigne}%
\end{groupe}%
%% v. 266-268a
\begin{groupe}%
%
\personnageGrec{Εὐριπίδης}%
\trimIamb{ἁνὴρ μὲν ἡμῖν οὑτοσὶ καὶ δὴ γυνὴ\apparat{266}{ἀνὴρ R : \lecon{ἁνὴρ} Fritzsche \ligneApp{267} τόδ᾽ R : \lecon{τό γ᾽} Küster}\\
τό γ᾽ εἶδος· ἢν λαλῇς δ᾽, ὅπως τῷ φθέγματι\\%\apparat{267}{τόδ᾽ R : \lecon{τό γ᾽} Küster}\\
γυναικιεῖς εὖ καὶ πιθανῶς.\\\passeLigne}%
\end{groupe}%
%% v. 268b
\begin{groupe}%
\personnageGrec{Μνησίλοχος}%
\trimIamb{\ligne{268}\decalage{γυναικιεῖς εὖ καὶ πιθανῶς.}πειράσομαι.\\\passeLigne}%
\end{groupe}%
%% v. 269a
\begin{groupe}%
\personnageGrec{Εὐριπίδης}%
\trimIamb{βάδιζε τοίνυν.\\\passeLigne}%
\end{groupe}%
%% v. 269b-270a
\begin{groupe}%
\personnageGrec{Μνησίλοχος}%
\trimIamb{\ligne{269}\decalage{βάδιζε τοίνυν.}μὰ τὸν Ἀπόλλω οὔκ, ἤν γε μὴ\\
ὀμόσῃς ἐμοί—\\\passeLigne}%
\end{groupe}%
%% v. 270b
\begin{groupe}%
\personnageGrec{Εὐριπίδης}%
\trimIamb{\ligne{270}\decalage{ὀμόσῃς ἐμοί—}τί χρῆμα;\\\passeLigne}%
\end{groupe}%
%% v. 270c-271
\begin{groupe}%
\personnageGrec{Μνησίλοχος}%
\trimIamb{\ligne{270}\decalage{ὀμόσῃς ἐμοί—τί χρῆμα;}συσσώσειν ἐμὲ\\
πάσαις τέχναις, ἤν μοί τι περιπίπτῃ κακόν.\\\passeLigne}%
\end{groupe}%
%%% v. 272
\begin{groupe}%
\personnageGrec{Εὐριπίδης}%
\trimIamb{`ὄμνυμι τοίνυν αἰθέρ᾽ οἴκησιν Διός.'%
\commentaire{272}{ὄμνυμι τοίνυν αἰθέρ᾽ οἴκησιν Διός.}{Selon le scholiaste, ce vers est une citation de la \italique{Mélanippe [la sage]} d'Euripide\marqueur{1} (ἐκ τῆς Μελανίππης Εὐριπίδου) ; citation qui devait être fameuse à la façon dont Aristophane l'évoque encore aux vers 100 et 311 des \italique{Grenouilles}.}%
\debutNote{272}%
\note{1}{\nieme{8} fragment de l'édition Jouan-van Looy aux Belles Lettres ; fr. 487 de l'édition Kannicht ; 670 chez Mette.}\\\passeLigne}%
\end{groupe}%
%% v. 273
\begin{groupe}%
\personnageGrec{Μνησίλοχος}%
\trimIamb{τί μᾶλλον ἢ τὴν Ὑποκράτους ξυνοικίαν;\apparat{273}{\lecon{Ὑποκράτους} R : Ἱπποκράτους Σ}%
\commentaire{273}{τί μᾶλλον ἢ τὴν Ὑποκράτους ξυνοικίαν;}{Au Ciel de qui se moque Euripide ?! Jurer sur l'Éther ? Et pourquoi pas `τὴν Ὑποκράτους ξυνοικίαν' (c'est-à-dire « par une porcherie ») ?\\
« Les fils d'Hippocratès étaient connus de tout le monde à Athènes dès l'époque des \italique{Nuées} (v. 1001). Car leur père étant fils d'Ariphron, propre frère de Périclès, ils étaient les petits neveux du fameux homme d'État. On leur avait donné le surnom de porcs à cause de leur \italique{balourdise} (ὑομουσία, \italique{Cavaliers}, v. 986). Les textes sont unanimes sur ce point. Athénée (III, 96e) : %
\citComm{ἑξῆς ἐπεισηνέχθη, μητρόπολίς τις ὡς ἀληθῶς οὖσα καὶ μήτηρ τῶν Ἱπποκράτους υἱῶν, οὓς εἰς ὑωδίαν κωμῳδουμένους οἶδα.\\
On servit ensuite une vulve de truie, mère des fils d'Hippocratès, lesquels, je le sais, ont été raillés par les comiques à cause de leur cochonnerie\marqueur{1}.}%
Cela se trouve confirmé par le scholiaste d'Aristophane sur les \italique{Nuées}, v. 1001, et par Suidas. De même Photius (σ 845), au mot σῦς :%
\citComm{ὗς τοὺς Ἱπποκράτους υἱοὺς ἔλεγον καὶ τοὺς Παναιτίου καὶ Μέμνονος εἰς ὑηνίαν κωμῳδοῦντες.\\
On nommait porcs les fils d'Hippocratès ainsi que ceux de Panaitios et de Memnon. Les comiques les raillent à cause de leur cochonnerie.}%
(Autres textes dans Meineke, \italique{Fr. Com. Gr.}, t. II, p. 477. [Pour une bibliographie complète et mise à jour, voir les \italique{testimonia} de l'édition Theodoridis de Photius à l'adresse σ 845.])
Le logis de famille d'Hippocratès est donc une porcherie. Il est à croire que l'auteur ne se faisait pas faute de prononcer συῳκία. » (Note de Willems.) \\
Je conserve la leçon du manuscrit de Ravenne : \grec{Ὑποκράτους}, sans reprendre la correction communément admise \grec{Ἱπποκράτους}. En effet, le préfixe \grec{ὑ}- sert à construire des composés sur le mot \italique{porc} ; ainsi, pour se borner aux textes cités par Willems, on recense \grec{ὑωδία} (littéralement, la \italique{ressemblance à un porc}) et \grec{ὑηνία} ( littéralement \italique{nature de porc}). 
En prononçant : \grec{τὴν ὑ-ποκράτου-ς-(σ)υ-ῳκίαν}, Mnésiloque insiste doublement sur la porcinité de l'hippocratique maison.
Voilà aussi pourquoi les enfants qui l'habitaient sont plus souvent qualifiés d'« \grec{υ-ἱοί Ἱπποκράτους} » que de « παῖδες Ἱπποκράτους ».%\marqueur{2}.
}%
\debutNote{273}%
\note{1}{Comprendre : \italique{grossièreté}, \italique{lourdeur}. ὗς, le \italique{porc} est pour les Grecs le symbole de l'ignorance, partant, de la sottise. Voir à ce sujet le toujours bon Taillardat, §451. L'adjectif Ὑάντιος servait à désigner les Béotiens, peuples à l'éternelle réputation de lourdeur et d'ignorance, cf. l'adage n° 906 d'Érasme : ἡ βοιωτία ὗς.}%
%\note{2}{Le scholiaste des \italique{Nuées} (v. 1001) confirme ce jeu de mots : l'Injuste s'adresse à Phidippide : « Écoute-le [le Juste], jeune homme... et tu ressembleras aux fils d'Hippocratès », τοῖς Ἱπποκράτους υἱέσιν εἴξεις. Tous les spectateurs comprenant ὑσίν.}
\\\passeLigne}%
\end{groupe}%
%% v. 274
\begin{groupe}%
\personnageGrec{Εὐριπίδης}%
\trimIamb{ὄμνυμι τοίνυν πάντας ἄρδην τοὺς θεούς.%\testimonia{274}{(ἄρδην) Hesychius α 7097}
\\\passeLigne}%
\end{groupe}%
%% v. 275-276
\begin{groupe}%
\personnageGrec{Μνησίλοχος}%
\trimIamb{μέμνησο τοίνυν ταῦθ᾽, ὅτι ἡ φρὴν ὤμοσεν,%\testimonia{275-6}{cf. \Eustathe 443.37}
\\
ἡ γλῶττα δ᾽ οὐκ ὀμώμοκ᾽· οὐδ᾽ ὥρκωσ᾽ ἐγώ.
\commentaire{275-6}{ὅτι ἡ φρὴν ὤμοσεν / ἡ γλῶττα δ᾽ οὐκ ὀμώμοκ᾽· οὐδ᾽ ὥρκωσ᾽ ἐγώ.}{Mnésiloque se méfie avec raison d'Euripide-le-Retors, il a notamment en tête le vers 612 de son \italique{Hippolyte}\marqueur{1} : 
\citComm{ἡ γλῶσσ' ὀμώμοχ', ἡ δὲ φρὴν ἀνώμοτος.}%
(La citation d'Euripide comportait déjà 7 `o' (dont 4 d'affilée) ; avec Aristophane, le compteur explose !)}%
\debutNote{275-6}%
\note{1}{Cf. Érasme, \italique{Adages} n° 1441, \latin{Lingua juravit} et Tosi, \italique{Dictionnaire des sentences latines et grecques}, §2252.}%
\apparat{post 276}{parepigraphe ὀλολύζουσι τε ἱερὸν ὠθεῖται R : ολολυ[ \PsiUn}%
\commentaire{post 276}{ὀλολύζουσι τε ἱερὸν ὠθεῖται.}{Le vers 276 est suivi de l'indication scénique : ὀλολύζουσι τε ἱερὸν ὠθεῖται ainsi que d'une κορωνίς en marge, qui marque le changement de scène. 
Le scholiaste confirme : παρεπιγραφή. ἐκκυκλεῖται ἐπὶ τὸ ἔξω τὸ θεσμοφόριον et l'on trouve déjà ολολυ[ sur un fragment de papyrus datant du second siècle après Jésus-Christ.\\
La maison d'Agathon est donc roulée par l'eccyclème au fond pour se transformer en temple des Thesmophores ; est hissé le \italique{signal}, le σημεῖον du vers 278 tandis que le chœur des femmes fait retentir l'ὀλολυγή ; c'est à-dire l'espèce de hurlement sacré qu'Hérodote (IV, 189) faisait remonter aux Libyens.
(Cette ululation dont il est question déjà dans Homère (\italique{Iliade}, VI, 301), et qui retentissait dans les temples en signe tantôt de douleur tantôt de joie.)
}\\\passeLigne}%
\end{groupe}%
%% v. 277-279a
\begin{groupe}%
\personnageGrec{Εὐριπίδης}%
\trimIamb{ἔκσπευδε ταχέως· ὡς τὸ τῆς ἐκκλησίας%\testimonia{277-8}{(ὡς -- φαίνεται) Suidas σ 325 (sed Θεσμοφόρῳ pro Θεσμοφορίῳ)}
\\%\apparat{277}{\lecon{ἔκσπευδε} R : εἄ· σπεῦδε Maas}\\
σημεῖον ἐν τῷ Θεσμοφορίῳ φαίνεται.\apparat{278}{\lecon{Θεσμοφορίῳ} R, cf. \italique{IG} II\expo{2} 1177.24 ; 2498.12 : Θεσμοφόρῳ Suid. σ 325 : Θεσμοφορείῳ Scaliger \ligneApp{282} Θεσμοφόρῳ R : \lecon{Θεσμοφόρω} Photius α 1386}\\
ἐγὼ δ᾽ ἄπειμι.%\\
\\\passeLigne}%
\end{groupe}%
%% v. 279b-294
\begin{groupe}%
\personnageGrec{Μνησίλοχος}%
\trimIamb{\ligne{279}\decalage{ἐγὼ δ᾽ ἄπειμι.}δεῦρό νυν ὦ Θρᾷτθ᾽ ἕπου.\\
ὦ Θρᾷττα θέασαι, καομένων τῶν λαμπάδων
\commentaire{280}{Θρᾷττα}{est un nom courant pour les esclaves Thraces (comme dans les \italique{Acharniens}).}\\
ὅσον τὸ χρῆμ᾽ ἀνέρχεθ᾽ ὑπὸ τῆς λιγνύος.
\commentaire{281}{ὑπὸ τῆς λιγνύος.}{Sur le sens de ὑπὸ : \italique{avec}, \italique{accompagné de} (souvent en parlant d'un accompagnement musical. Cf. Photius, \italique{Bibliothèque} 320b : οἱ παλαιοὶ τὴν ὑπό ἀντὶ τῆς μετά πολλάκις ἐλάμβανον), voir la note d'Alphonse Willems n°7, écrite d'une juste colère (et reproduite à la fin).}\\
ἀλλ᾽, ὦ περικαλλεῖ Θεσμοφόρω, δέξασθέ με%\testimonia{282}{Photius α 1386 ; \AelDion α 113}
\apparat{282}{περὶ καλλῆ R : \lecon{περικαλλεῖ} Kaibel%, Hall \& Geldart
 \ligneApp{283} καὶ δεῦρο πάλιν R : \lecon{καὶ δεῦρο καὶ πάλιν} Bentley \ligneApp{285} \lecon{τὰ πόπαν᾽} \PsiUn : τὸ πόπανον R}\\
ἀγαθῇ τύχῃ καὶ δεῦρο καὶ πάλιν οἴκαδε.\\\passeLigne}
\end{groupe}%
%% v. 279b-294
\begin{groupe}%
\trimIamb{%\ledpb%
ὦ Θρᾷττα τὴν κίστην κάθελε, κᾆτ᾽ ἔξελε\\
τὰ πόπαν᾽, ὅπως λαβοῦσα θύσω ταῖν θεαῖν.%\testimonia{285}{(πόπαν᾽) \Eustathe 437.4}
\\%\apparat{285}{\lecon{τὰ πόπαν᾽} \PsiUn, Porson : τὸ πόπανον R}\\
δέσποινα πολυτίμητε Δήμητερ φίλη\\
καὶ Φερρέφαττα, πολλὰ πολλάκις μέ σοι\apparat{287}{\lecon{Φερρέφαττα} \PsiUn, Thiersch : Φερσέφαττα R}\\
θύειν ἔχουσαν, εἰ δὲ μή, ἀλλὰ νῦν λαθεῖν%
\commentaire{288}{εἰ δὲ μή, ἀλλὰ νῦν λαθεῖν.}{Mnésiloque, qui déclame sa prière à haute voix afin de mystifier l'assemblée des femmes, prononce en aparté cette réplique.}\\
καὶ τὴν θυγατέρα Χοιρίλ' ἀνδρός μοι τυχεῖν%\testimonia{289}{(χοῖρον) Euripides, testimonia 4, 18 Kannicht Χοιρίλη}
\apparat{289}{χοῖρον R : χοιρίον Σ : \lecon{Χοιρίλ'} Bibi}%
\commentaire{289}{καὶ τὴν θυγατέρα χοῖρον ἀνδρός μοι τυχεῖν.}{Les vers 288 à 291 sont corrompus en plusieurs endroits. Occasion trop belle pour les philologues de nous infliger leur saleté morale. Et ils s'en sont donné à cœur joie : Mnésiloque userait d'obscénités pour parler de sa fille (τὴν θυγατέρα χοῖρον ; χοῖρον, \italique{le porc}, allusion selon eux au sexe de la femme), quand il ne lui souhaiterait tout bonnement pas les mœurs les plus infamantes. % Il en vaet de son fils (πρὸς θάληκον (??)). 
Et ce au beau milieu d’une prière aux deux déesses parmi les plus révérées !.. Juste courroux d'Alphonse Willems\marqueur{1}. (Allez donc chercher dans tout Aristophane la moindre irrévérence à Déméter ou Coré.) On ne le répétera jamais assez : « Mais pour qui prennent-ils Aristophane ? » 
Las ! on ne détruit que ce que l'on remplace et le coup d'épée de Willems – qu'il avait pourtant aiguisée – porta dans l'eau car sa note était essentiellement aporétique. Prions les déesses thesmophores de conférer à la nôtre assez de force pour dissiper à tout jamais ces nuées. 
%Las ! ce coup d'épée de Willems – qu’il avait pourtant aiguisée – a porté dans l'eau. %Revenons donc à la charge pour tenter à notre tour de dissiper ces nuées !
%Il ne sera pas bien difficile de dissiper ces nuées : Mnésiloque nomme sa fille ; †χοῖρον† est évidemment une corruption. 
Le nom ou surnom de la fille de Mnésiloque, †χοῖρον†, est évidemment corrompu.
La preuve en est que scholiaste avait une autre leçon devant les yeux : χοιρίον\marqueur{2}. 
N'allons pas inventer un sobriquet obscène ; reportons-nous plutôt aux quelques éléments biographiques qui nous sont parvenus\marqueur{3}. Mnésiloque est le beau-père d'Euripide, et la fille qu'il lui a cédé en mariage s'appelle justement Χοιρίλη\marqueur{4} ; la correction est toute indiquée : la présence d'un nombre propre (rare) a naturellement tendance à induire des fautes de transmission du texte ; de plus, la coquille consistant à substituer un Ν à un Λ est monnaie courante. 
Mnésiloque profite donc de cette prière fictive pour tacler son gendre Euripide et souhaiter à sa fille Chœrile un mari riche et de savoir gérer à bon escient le porte-monnaie conjugal -- Euripide avait-il la réputation d'être mauvais économe ? (cf. la note ci-dessous.) }%
\debutNote{289}%
\note{1}{Tome II de son édition posthume d'Aristophane (1919), pp. 548-549. %Édition malheureusement désormais introuvable mais en accès libre chez les grandes bibliothèques américaines numériques – à condition de se faire passer pour américain.
Cette note sur ce passage des \italique{Thesmophories}, la \nieme{8} qu'il écrivit, est reproduite en annexe.}%
\note{2}{La scholie : ὡς τῆς θυγατρὸς αὐτοῦ κατ' ἐπίκλησιν οὕτω καλουμένης. οἷον βοίδιον ἢ χρυσίον ἢ μυρτ<ί>ον. ἴσως δὲ τὸ αἰδοῖον τὸ γυναικεῖον αἰνίττεται.}%
\note{3}{Voir le « Dossier Mnésiloque » de l'Introduction.}%
\note{4}{Certains manuscrits portent Χοιρίνη.}\\
πλουτοῦντος, ἄλλως τ᾽ ἠλιθίου κἀβελτέρου,\\
καὶ πρὸς θυλάκιον νοῦν ἔχειν μοι καὶ φρένας.%\testimonia{291}{(θάληκον) Hesychius θ 50 θάλλικα ; Pollux X 152 θυλάκιον [...] ἐπὶ τοῦ αὐτοῦ δὲ καὶ σακίον, ὅταν φῇ ἐν Θεσμοφοριαζούσαις}
\apparat{291}{πρὸς θάληκον R : Ποσθαλίσκον Dindorf : \lecon{πρὸς θυλάκιον} Bibi}%
\commentaire{291}{καὶ πρὸς θάληκον νοῦν ἔχειν μοι καὶ φρένας.}{« Pour ce qui du †θάληκον†, qu'elle ait esprit et intelligence. » Le vers 291 a fait l'objet de nombreux débats ; et pour cause : le mot θάληκον est inconnu au bataillon des lexiques.
% Que signifie donc ce `θάληκον', inconnu autrement ?
Une scholie en marge du manuscrit de Ravenne se référant probablement au \nieme{291} vers indique :%
\citComm{τὸν παιδαρίσκον. ἴσως δὲ παρὰ τὴν πόσθην αὐτὸ συνέθηκεν,\\
le jeune garçon. Peut-être à rapprocher du pénis (πόσθη).}%
Les commentateurs s'en sont remis à cette scholie pour en induire que πρὸς θάληκον était une transmission fautive d'un surnom de garçon, forgé sur le mot πόσθη, \italique{le membre}. %\marqueur{2}. 
La reconstruction encore la plus en vogue est celle de Dindorf : πρὸς θάληκον $ \to $ Ποσθαλίσκον, où Ποσθαλίσκον désignerait le fils (fictif) de Mnésiloque. 
À ce ποσθαλίσκον, Willems a opposé deux raisons\marqueur{1}, que je crois fortes et fondées : « il est plus que douteux que ce dernier mot [Ποσθαλίσκον] soit grec, les diminutifs se faisant en ισκος et non en αλισκος. (On cite, il est vrai, κωραλίσκος, παιδαρίσκος et μειρακίσκος. Ce serait concluant s'il n'y avait pas l'intermédiaire de κωράλιον, παιδάριον et μειράκιον\marqueur{2}. Si des diminutifs comme ποσθάλιον et ποσθαλίσκος avaient existé, ils n'auraient pu manquer d'être relevés par d'anciens lexicographes.) »
À l'école de Willems, j'ai aussi appris à me méfier du scholiaste\marqueur{3}, cet « ignare byzantin », et faire bien davantage confiance à l'« alexandrin » Hésychius. Autres lieux, autres temps et autres sciences !
Comme l'a déjà fait remarquer Alphonse Willems, Hésychius n'a jamais compulsé de ποσθαλίσκον ou autre composé de πόσθη pouvant expliquer notre πρὸς θάληκον. 
En revanche, l'on trouve chez Hésychius à l'adresse θ 50 :
\citComm{θάλλικα· σάκκου εἶδος\\
Thallika : sorte de sac.}%
Voilà qui est beaucoup plus proche du texte transmis dans le \italique{Ravennas}, θάληκον. 
%Substituer θάλικα à θάληκον semble à peine être une correction par rapport à ce qui était proposé jusqu'ici\marqueur{7}.
%(La glose d'Hésychius sur θαλικοποιοί laisse à penser que le sac en question était fait de cuir et qu'il ressemblait à un petit coffret.)\\
Pour le sens, j'ai reçu l'aide précieuse de Jean-Victor Vernhes qui « estime qu'il s'agit du sac où cette femme met des bijoux, la clef de sa maison, ses emplettes, sa carte de crédit (et bientôt, vous allez voir, celle de son mari), et qu'il symbolise d'une façon générale, ses affaires, ses intérêts, au profit desquels elle exploite son nigaud de conjoint. Elle finira par tout fourrer dans son sac ! »\marqueur{4}
% Comprendre ce θάλλικα comme une métonymie du portefeuille du ménage -- la \italique{cassette} --%\marqueur{8} (ou, pour rester plus proche de ce à quoi devait ressembler un véritable θάλιξ, une \italique{cassette}), portefeuille que la fille devra gérer avec économie et esprit contrairement à ces idiots d'hommes, me paraît être tout-à-fait dans le ton d'Aristophane ! 
Thème cher à Aristophane : confier aux femmes les clefs de la cité tout autant que celles du \italique{portefeuille du ménage} ; pour ces choses-là, elles ont beaucoup plus d'esprit et d'intelligence, νοῦν ἔχειν καὶ φρένας. Il y a de nombreux passages de l'\italique{Assemblée des Femmes} ou de \italique{Lysistrata} qui feraient écho à ce genre de prière.
Néanmoins, le mot θάλλικα (de θάλλιξ ?) est rare et circonscrit au lexique d'Hésychius. Il semble être une métathèse d'un mot beaucoup plus courant et qu'on trouve chez Aristophane : θυλάκιον. 
Si J.-V. Verhnes avait dû condenser son propos en grec, il aurait écrit quelque chose comme : σακίον, ἐν οἷσπερ τἀργύριον ταμιεύεται. Hé bien, c'est exactement la glose qu'on retrouve dans un passage de l'\italique{Onomasticon} de Pollux citant... les \italique{Thesmophories} d'Aristophane.
Pollux (X, 152) donne une liste de divers synonymes pour désigner une \italique{bourse}. Voici ce qu'il dit :%
\citComm{καὶ μιξάμενοι δ’ ἂν εἴποιμεν σκεύη, βαλάντια καὶ βαλαντίδια ὡς ἐν Αἰξὶν Εὔπολις. καὶ \textbf{θυλάκιον} δὲ καὶ θυλακίσκον·\\
Ἀριστοφάνης γοῦν ἐν Τριφάλητι τοῦτο ὑποδηλοῖ, ὅταν φῇ
\citComm{ἔπειτ' ἐπὶ τοὖψον ἧκε, τὴν σπυρίδα λαβών\\
καὶ θυλακίσκον καὶ τὸ μέγα βαλάντιον. }
$[$ ... $]$
ἐπὶ τοῦ αὐτοῦ δὲ καὶ σακίον, ὅταν φῇ ἐν Θεσμοφοριαζούσαις 
\citComm{\textbf{σακίον, ἐν οἷσπερ τἀργύριον ταμιεύεται.}}
Et les mêlant nous dirions, comme dans les \italique{Chèvres} d'Eupolis, les ustensiles, les grandes bourses et les petites bourses. Ainsi que θυλάκιον et θυλακίσκον. Aristophane l'illustre dans \italique{Triphalès}, lorsqu'il dit :
\citComm{Ensuite il arriva aux hors-d'œuvre, tenant la corbeille, un petit sac (θυλακίσκον) et une grande bourse.}
Ce mot désigne aussi une petite bourse, comme cela apparaît dans les \italique{Thesmophories} :
\citComm{\italique{une petite bourse, dans lesquelles on dépose l'argent.}}}%
Aristophane a écrit deux pièces intitulées \italique{Les Thesmophories}. Ce passage a été considéré comme un fragment des \italique{Thesmophories} perdues (fr. 343 de l'édition Kassel-Austin). En fait, ce petit bout de phrase explicatif, σακίον· ἐν οἷσπερ τἀργύριον ταμιεύεται, ressemble bien davantage à une glose d'un commentateur. Il arrive fréquemment à Pollux de fusionner des gloses avec le texte d'Aristophane lui-même\marqueur{5}. Ce n'est pas en réalité aux \italique{Thesmophories} perdues mais bien au vers 291 des \italique{Thesmophories} conservées que se rattache cette glose, en tant que définition de θυλάκιον :%
\citComm{$<$θυλάκιον·$>$ σακίον, ἐν οἷσπερ τἀργύριον ταμιεύεται.}%
Une métathèse ancienne aura donné la leçon fautive du \italique{Ravennas}, comportant †θάληκον† en place de θυλάκιον, tandis que cette métathèse adjointe à sa glose `σακίον' rejoignait le lexique d'Hésychius sous la forme : θάλλικα· σάκκου εἶδος. 
Pour séduisante que soit cette reconstruction, elle n'en demeure pas moins conjecturale ; le \italique{Ravennas} ne comportant pour les \italique{Thesmophories} que très peu de scholies, et parmi les moins intéressantes.
Finissons d'ailleurs cette note en montrant que notre conjecture est aussi compatible avec l'unique scholie conservée, τὸν παιδαρίσκον. ἴσως δὲ παρὰ τὴν πόσθην αὐτὸ συνέθηκεν. Le scholiaste ayant devant lui θάληκον -- qu'il lisait indifféremment θάλικον -- devait l'analyser comme un mot formé sur τό θάλος, \italique{la jeune pousse} et, dans un sens métaphorique, l'\italique{enfant}, suivi d'un suffixe en -ικος ; l'on comprend le quiproquo, facilité par l'évocation deux vers plus haut de la fille de Mnésiloque et d'un imaginaire sexe féminin (χοιρίον).
%Un point reste à élucider : comment expliquer la scholie : τὸν παιδαρίσκον. ἴσως δὲ παρὰ τὴν πόσθην αὐτὸ συνέθηκεν ? 
%Si l'on admet que le scholiaste avait devant lui θάλικα ou θάλικον, et qu'il l'analysait non comme un dérivé de θάλιξ, qu'il ignorait, mais comme un mot formé sur τό θάλος, \italique{la jeune pousse} et, dans un sens métaphorique, \italique{rejeton}, \italique{enfant}, suivi d'un suffixe en ικος, l'on peut comprendre le quiproquo, facilité par l'évocation deux vers plus haut de la fille de Mnésiloque. De toute façon l'hésitation du scholiaste (ἴσως δὲ) interdit de penser qu'il ait pu avoir devant lui une évocation aussi claire de la πόσθη que Ποσθαλίσκον.
%Lisons donc le vers 291 ainsi :
%\citComm{καὶ πρὸς θάλικα νοῦν ἔχειν μοι καὶ φρένας.\\
%\francais{et m'avoir l'esprit et l'intelligence appliqués à la cassette.}}
%Abordons une question en rapport avec ce passage, celle de notre anti-héros : Mnésiloque. Dans les indications de personnage du manuscrit de Ravenne, le protagoniste des \italique{Thesmophories} est indifféremment nommé Μνησίλοχος ou κηδεστής, le \italique{parent par alliance}. 
%Parent par alliance, il l'est assurément\marqueur{9} !.. mais est-il vraiment ce Mnésiloque que les \italique{Vies} d'Euripide\marqueur{10} présentent comme le père de la première femme d'Euripide, Χοιρίνη ou Χοιρίλη\marqueur{11} ? Si tel est le cas, alors il est certain qu'Aristophane ne pouvait se priver de cette prière factice de Mnésiloque à sa fille pour tacler Euripide ! 
%Faut-il supposer que l'acteur mâchait `χοιρίον' au vers 289 pour faire entendre Χοιρίν'\marqueur{12} ? Ou, au contraire, faut-il arguer de l'absence d'allusion à Euripide dans ce passage pour aller dans le même sens que la majorité des éditeurs modernes, pour qui la précision de l'identité du κηδεστής est un ajout postérieur à Aristophane ? La question, pour moi, n'est toujours pas tranchée ; qu'un lecteur éclairé m'éclaire !
}%
\debutNote{291}%
%\note{1}{Le manuscrit de Ravenne, qu'on abrège `R', comporte en fait : πεδαρίσκον.}%
%\note{2}{Comme le terme de caresse pour la jeune fille, χοιρίον, peut désigner – dans un autre contexte – le sexe féminin. Joie mauvaise des philologues qui ont ainsi pu laisser folâtrer leur imagination et plaquer leur saleté morale sur Aristophane. Juste courroux d'Alphonse Willems (\nieme{8} note) devant ces reconstructeurs de pièce qui nous inventent un personnage usant d'obscénités pour parler de sa propre fille (quand il ne lui souhaite pas les mœurs les plus infamantes), et ce au beau milieu d'une prière aux deux déesses parmi les plus révérées !! (Allez donc chercher dans tout Aristophane la moindre irrévérence à Déméter ou Coré.) On ne le répétera jamais assez : « Mais pour qui prennent-ils Aristophane ? » Las ! on ne détruit que ce que l'on remplace et le coup d'épée de Willems – qu'il avait aiguisée – avait porté dans l'eau car sa note était essentiellement aporétique. Prions les déesses thesmophores de conférer à la nôtre assez de force pour dissiper à tout jamais ces nuées !}%
\note{1}{Note 8, reproduite à la fin.}%
\note{2}{Chantraine, dans sa \italique{Formation des noms en grec ancien} aux pages 406 et suivantes, parle longuement du suffixe -ισκος sans même évoquer un possible -αλισκος ou -αρισκος ; Chantraine cite παιδαρίσκος et μειρακίσκος en expliquant que le Grec recourait à ce dernier diminutif parce que l'autre, μειράκιον, était devenu d'usage banal.}%
%\note{5}{Willems pense notamment à Hésychius. %, qui semble avoir fait un dépouillement complet d'Aristophane (ou, du moins, a dépouillé le commentaire que Didyme en a fait, cf. P. Boudreaux, \italique{Le texte d'Aristophane et ses commentateurs}), à une époque où le texte transmis était bien meilleur que le nôtre, Hésychius étant lui-même un érudit comme il ne peut plus en exister ; il avait probablement lu davantage de grec que nous n'en avons conservé. 
%Il pense peut-être aussi au lexique, beaucoup plus tardif (probablement composé au \nieme{X} siècle), de l'hypothétique Suidas : la \italique{Souda} ; lexique qui fait la part belle à Aristophane et qui est une mine précieuse de renseignements car il s'appuie sur des leçons du texte parfois meilleures, notamment pour notre pièce des \italique{Thesmophories} qui n'est transmise intégralement que par un unicodex, le manuscrit de Ravenne. 
%Suidas \& Hésychius n'ont effectivement pas manqué de relever le surnom πόσθων -- \italique{couillon}, traduira-t-on -- qui apparaît dans la \italique{Paix} (vers 1300). Ce surnom a aussi été relevé par le lexicographe Photius. De ποσθαλίσκον, par contre, point de trace !}%
\note{3}{Je parle toujours, par abus de langage et conformément à un usage trop bien établi, \italique{du} scholiaste. Il est bien évident qu'en réalité les commentaires en marge des manuscrits sont de différents auteurs, et ont été compulsés les uns à la suite des autres postérieurement. Certaines scholies font preuve d'une grande érudition et apportent parfois des leçons meilleures que celle du texte transmis, d'autres, par contre, – la plupart, hélas ! – ne sont, au mieux, d'aucun intérêt voire souvent sources d'erreurs pour qui y accorderait trop d'importance. Pour finir de gâter l'affaire, les \italique{Thesmophories} est la pièce d'Aristophane la moins pourvue en scholies.}%
\note{4}{La glose d'Hésychius sur θαλικοποιοί laisse à penser que le sac en question était fait de cuir et qu'il ressemblait à un petit coffret -- une \italique{cassette}.}%
%\note{7}{Le remplacement du double lambda `λλ' par un simple `λ' est commandé par les lois de scansion. Pour une telle graphie, cf. la glose d'Hésychius voisine : θαλιοποιοί· οἱ τὰ σκυτούμενα κιβώτια, καὶ τοὺς δερμάτινους ῥίσκους ἐργαζόμενοι.
%Aristophane prenant souvent des libertés avec le doublement des consonnes, et Hésychius n'étant pas bien sûr sur ce point (Hésychius, lui aussi, est conservé par un unique manuscrit bien postérieur : qu'on estime du \nieme{XIV} ou \nieme{XV} siècle), on ne peut même pas appeler ça une correction.
%Le iota et le êta qui se prononçaient de la même façon et sont allègrement permutés sans scrupules dans les manuscrits tardifs. Pour la confusion de la terminaison `α' et `ον', la faute est compréhensible pour un mot rare et certainement inconnu du scribe qui y aura vu un suffixe en -ικος. 
%(Une confusion analogue, entre un neutre, singulier en `ον', et pluriel en `α', avait déjà eu lieu au vers 285 : le manuscrit de Ravenne comporte τὸ πόπανον ; le sens poussait Porson en son temps à réclamer la correction τὰ πόπαν(α) ; cette émendation fut ensuite confirmée par la découverte d'un papyrus.)}%
%\note{8}{ὁ θύλακος désigne en grec le \italique{sac à provisions}.}%
\note{5}{Voir au commentaire du vers 1175 un autre exemple.}%
%\note{9}{Les vers 74 et 210 l'attestent.}%
%\note{10}{Voir, par exemple, la \italique{Souda} à l'entrée Εὐριπίδης.}%
%\note{11}{L'épigraphie fait état de plusieurs Χοιρίνη tandis que le nom Χοιρίλη ne semble pas attesté.}%
%\note{12}{Je n'exclue d'ailleurs pas que ce fût peut-être la leçon originale ; le iota de `χοιρίον' ayant été ajouté d'après la scholie. Voir le « dossier Mnésiloque » en Introduction.}
\\
ποῦ ποῦ καθίζωμ᾽ ἐν καλῷ, τῶν ῥητόρων%\testimonia{292}{(τῶν ῥητόρων) Photius ρ 106 ῥήτορα· καὶ τὴν γυναῖκα. Ἀριστοφάνης}
\\
ἵν᾽ ἐξακούω; σὺ δ᾽ ἄπιθ᾽ ὦ Θρᾷττ᾽ ἐκποδών.\\
δούλοις γὰρ οὐκ ἔξεστ᾽ ἀκούειν τῶν λόγων.\\\passeLigne}%
\end{groupe}%
\endnumbering
\end{Leftside}
\begin{Rightside}
\beginnumbering
%% v. 252a
\begin{groupe}%
\personnageFr{Agathon}%
\trimIambFr{Prenez et usez-en ; je ne m'y oppose pas.\\\passeLigne}%
\end{groupe}%
%% v. 252b
\begin{groupe}%
\personnageFr{Mnésiloque}%
\trimIambFr{\ligne{252}\espace{5cm}Que faut-il prendre ?\\\passeLigne}%
\end{groupe}%
%% v. 253
\begin{groupe}%
\personnageFr{Euripide}%
\trimIambFr{Quoi ? prends d'abord cette crocote et revêts-la.\\\passeLigne}%
\end{groupe}%
%%% v. 254
\begin{groupe}%
\personnageFr{Mnésiloque}%
\trimIambFr{Par Aphrodite, elle fleure bon le priape.\testimonia{254}{Suidas π 2114 \espTest (ποσθίου) Photius π 1113 πόσθιον· αἰδοῖον· ἢ (sic) τὴν Ἀφροδίτην ; Hesychius π 3096}\\\passeLigne}%
\end{groupe}%
%% v. 255a
\begin{groupe}%
\personnageFr{Agathon}%
\trimIambFr{Serre-t'en vite la taille.\\\passeLigne}%
\end{groupe}%
%% v. 255b
\begin{groupe}%
\personnageFr{Euripide}%
\trimIambFr{\ligne{255}\espace{2.5cm}Maintenant passe-nous un strophion.\\\passeLigne}%
\end{groupe}%
%% v. 255c
\begin{groupe}%
\personnageFr{Agathon}%
\trimIambFr{\ligne{255}\espace{5.5cm}Là.\\\passeLigne}%
\end{groupe}%
%% v. 256
\begin{groupe}%
\personnageFr{Mnésiloque}%
\trimIambFr{Maintenant arrange-moi les jambes.\\\passeLigne}%
\end{groupe}%
%% v. 257a
\begin{groupe}%
\personnageFr{Euripide}%
\trimIambFr{Une résille et un bandeau ?\testimonia{257-8}{(ἡδὶ -- φορῶ) Suidas π 1165 (sed ἥδιον pro ἡδὶ)}\\\passeLigne}%
\end{groupe}%
%% v. 257b-258
\begin{groupe}%
\personnageFr{Agathon}%
\trimIambFr{\ligne{257}\espace{4.2cm}Plutôt cette\\
perruque, que je porte la nuit.\testimonia{258}{Photius π 690 περίθετος κεφαλή ; \Eustathe 1280.63 \espTest (κεφαλὴ περίθετος) Pollux II 35}\\\passeLigne}%
\end{groupe}%
%% v. 259
\begin{groupe}%
\personnageFr{Euripide}%
\trimIambFr{Eh ! par Zeus, c'est même tout juste l'affaire.\\\passeLigne}%
\end{groupe}%
%%% v. 260a
\begin{groupe}%
\personnageFr{Mnésiloque}%
\trimIambFr{M'ira-t-elle ?\\\passeLigne}%
\end{groupe}%
%%% v. 260b-261a
\begin{groupe}%
\personnageFr{Euripide}%
\trimIambFr{\ligne{261}\espace{2.7cm}Mais c'est parfait, par Zeus.\\
Un encycle ?\\\passeLigne}%
\end{groupe}%
%%% v. 261b
\begin{groupe}%
\personnageFr{Agathon}%
\trimIambFr{\ligne{261}\espace{2.4cm}Prends celui-là, sur la couchette.\testimonia{261}{Suidas ε 135 (sed τουτὶ pro τι) \espTest (ἔγκυκλον -- κλινίδος) Suidas κ 1809}\\\passeLigne}%
\end{groupe}%
%%% v. 262a
\begin{groupe}%
\personnageFr{Euripide}%
\trimIambFr{Des chaussures ?\\\passeLigne}%
\end{groupe}%
%%% v. 262b
\begin{groupe}%
\personnageFr{Agathon}%
\trimIambFr{\ligne{262}\espace{2.6cm}Voici les miennes, prends-les.\\\passeLigne}%
\end{groupe}%
%% v. 263
\begin{groupe}%
\personnageFr{Mnésiloque}%
\trimIambFr{Est-ce qu'elles m'iront ? \\\passeLigne}%
\end{groupe}%
%% v. 263b
\begin{groupe}%
\personnageFr{Agathon}%
\trimIambFr{\espace{3.6cm}Tu préférerais en chausser de larges...\testimonia{263}{Suidas χ 7 χαλαρὰ γοῦν χαίρεις φορῶν ὑποδήματα (sed χαίροις codd. FV)}\\\passeLigne}%
\end{groupe}%
%% v. 264a
\begin{groupe}%
\personnageFr{Mnésiloque}%
\trimIambFr{Ça tu en sais quelque chose.\\\passeLigne}%
\end{groupe}%
%% v. 264b-265
\begin{groupe}%
\personnageFr{Agathon}%
\trimIambFr{\espace{3.1cm} Or çà, tu as ce qu'il te faut.\\
Qu'on me roule au plus vite à l'intérieur.\\\passeLigne}%
\end{groupe}%
%% v. 266-268a
\begin{groupe}%
\personnageFr{Euripide}%
\trimIambFr{Voilà notre homme devenu femme,
du moins en apparence. Si tu parles, tâche par ton timbre de voix
de simuler la femme, avec aisance et naturel.\\\passeLigne}%
\end{groupe}%
%% v. 268b
\begin{groupe}%
\personnageFr{Mnésiloque}%
\trimIambFr{\ligne{268}\espace{4.1cm}(Voix de fausset.) J'y tâcherai.\\\passeLigne}%
\end{groupe}%
%% v. 269a
\begin{groupe}%
\personnageFr{Euripide}%
\trimIambFr{Va maintenant.\\\passeLigne}%
\end{groupe}%
%% v. 269b-270a
\begin{groupe}%
\personnageFr{Mnésiloque}%
\trimIambFr{\ligne{269}\espace{2.3cm}Non, par Apollon, à moins que\\
tu ne me jures...\\\passeLigne}%
\end{groupe}%
%% v. 270b
\begin{groupe}%
\personnageFr{Euripide}%
\trimIambFr{\ligne{270}\espace{2.3cm}Quoi ?\\\passeLigne}%
\end{groupe}%
%% v. 270c-271
\begin{groupe}%
\personnageFr{Mnésiloque}%
\trimIambFr{\ligne{270}\espace{4cm}D'aider à me sauver
par toutes tes rubriques, s'il me tombe dessus quelque malheur.\\\passeLigne}%
\end{groupe}%
%%% v. 272
\begin{groupe}%
\personnageFr{Euripide}%
\trimIambFr{Eh bien, « je jure, par l'Éther, la demeure de Zeus ».\\\passeLigne}%
\end{groupe}%
%% v. 273
\begin{groupe}%
\personnageFr{Mnésiloque}%
\trimIambFr{Pourquoi pas aussi bien par le logis de famille d'Hippocratès ?\\\passeLigne}%
\end{groupe}%
%% v. 274
\begin{groupe}%
\personnageFr{Euripide}%
\trimIambFr{Eh bien! je jure par tous les dieux en bloc.\testimonia{274}{(ἄρδην) Hesychius α 7097}\\\passeLigne}%
\end{groupe}%
%% v. 275-276
\begin{groupe}%
\personnageFr{Mnésiloque}%
\trimIambFr{Au surplus, souviens-toi de ceci : c'est ton cœur qui a juré,\testimonia{275-6}{cf. \Eustathe 443.37}
et non ta langue ; je ne la tiens pas pour liée.\\\passeLigne}%
\end{groupe}%
%% v. 277-279a
\begin{groupe}%
\personnageFr{Euripide}%
\trimIambFr{Hâte-toi de partir. Car le signal de l'assemblée\ligne{276}
paraît sur le Thesmophorion. Moi, je m'en vais.\testimonia{277-8}{(ὡς -- φαίνεται) Suidas σ 325 (sed Θεσμοφόρῳ pro Θεσμοφορίῳ)}
(Tous sortent. Le décor du fond change et représente le Thesmophorion.\ligne{276}
Entrent Mnésiloque suivi d'une servante portant une ciste.)\\\passeLigne}%
\end{groupe}%
%% v. 279b-294
\begin{groupe}%
\personnageFr{Mnésiloque}%
\trimIambFr{\ligne{279}\espace{2.3cm}Par ici, Thratta, suis-moi.\\
Vois, Thratta, sous la lueur fumeuse des torches ardentes,
vois la foule qui monte.\\
— Thesmophores, merveilles de beauté, accueillez-moi\testimonia{282}{Photius α 1386 ; \AelDion α 113}\\
sous d'heureux auspices, ici et à mon retour au logis.\\\passeLigne}
\end{groupe}%
%% v. 279b-294
\begin{groupe}%
\trimIambFr{Thratta, dépose la ciste, et tires-en
le gâteau, pour que je l'offre aux deux déesses.\testimonia{285}{(πόπαν᾽) \Eustathe 437.4}
Maîtresse très honorée, chère Déméter,
et toi Perséphone, puissé-je mainte et mainte fois
être à même de te sacrifier. (À part :) À tout le moins, n'être pas découvert aujourd'hui.
(Fort :) Que ma fille Chœrile trouve un mari riche\testimonia{289}{(χοῖρον) Euripides, testimonia 4, 18 Kannicht Χοιρίλη}
fortuné, et au surplus benêt et stupide.
Et qu'elle me soit sensée et avisée sur le porte-monnaie.\testimonia{291}{(θάληκον) Hesychius θ 50 θάλλικα ; Pollux X 152 θυλάκιον [...] ἐπὶ τοῦ αὐτοῦ δὲ καὶ σακίον, ὅταν φῇ ἐν Θεσμοφοριαζούσαις}\\
Où donc, où m'asseoir en bonne place pour entendre les orateurs ?\testimonia{292}{(τῶν ῥητόρων) Photius ρ 106 ῥήτορα· καὶ τὴν γυναῖκα. Ἀριστοφάνης}
Toi, Thratta, retire-toi à l'écart,
car il n'est pas permis aux esclaves d'ouïr ce qui se dit ici. \\\passeLigne}%
\end{groupe}%
\endnumbering
\end{Rightside}
\end{pages}
\Pages
\clearpage\ifodd\value{page}\hbox{}\newpage\fi
\enTeteG{ΠΑΡΟΔΟΣ --- Codex Ravennas 429, f. 166r (Classense 168r) --- vv. 295–336}
\begin{adjustwidth}{-1.5cm}{-3.0cm}
\begin{adjustbox}{max totalheight=21.0cm, max width=16.5cm, keepaspectratio}
\begin{minipage}{16.5cm}
{\fontsize{7.2pt}{8.5pt}\selectfont
\noindent\textbf{Folio 08 --- Codex Ravennas 429, \texttt{f. 166r} (\texttt{Classense 168r}) : ΠΑΡΟΔΟΣ, vv. 295–336}\hfill \textit{Témoins :} $\text{R}^1, \Rpc, \Rdeux$, \PsiUn\par\vspace{2pt}
\hrule height 0.4pt\vspace{2.5pt}
}
{\fontsize{8.1pt}{9.5pt}\selectfont
\ligneRav{295}{\inRas{Κῆρυξ}{}}{\hspace{1.1em}\difference{εὐφημία} ἔστω,}{\textbf{295} εὐφημί᾽ R : \lecon{εὐφημία} Suid. ε 3797}
\ligneRav{}{}{\hspace{1.1em}εὐφημία ἔστω·}{}
\ligneRav{297}{}{\hspace{1.1em}εὔχεσθε ταῖν Θεσμοφόροιν $[\![\text{\sout{τῇ Δήμητρῃ καὶ τῇ Κόρῃ}}]\!]$}{\textbf{297} Θεσμοφόροιν τῇ Δήμητρῃ καὶ τῇ Κόρῃ R : \lecon{Θ.} Blaydes}
\ligneRav{}{}{\hspace{1.1em}καὶ τῷ Πλούτῳ καὶ τῇ Καλλιγενείᾳ}{}
\ligneRav{300}{}{\hspace{1.1em}καὶ τῇ Κουροτρόφῳ $[\![\text{\sout{τῇ Γῇ}}]\!]$}{}
\ligneRav{301}{}{\hspace{1.1em}καὶ τῷ Ἑρμῇ καὶ \difference{ταῖς Χ}άρισιν}{\textbf{301} Χάρισιν R : \lecon{ταῖς Χ.} Meineke}
\ligneRav{}{}{\hspace{1.1em}ἐκκλησίαν τήνδε καὶ σύνοδον τὴν νῦν}{}
\ligneRav{}{}{\hspace{1.1em}κάλλιστα καὶ ἄριστα ποῆσαι,}{}
\ligneRav{304}{}{\hspace{1.1em}πολυωφελῶς μὲν \difference{τῇ π}όλει τῇ Ἀθηναίων,}{\textbf{304} πόλει R : \lecon{τῇ π.} Meineke}
\ligneRav{304}{}{\hspace{1.1em}}{\textbf{304} πόλει R : \lecon{τῇ π.} Meineke}
\ligneRav{305}{}{\hspace{1.1em}τυχηρῶς δ᾽ ἡμῖν αὐταῖς·}{\textbf{305} \lecon{ἡμῖν} R : ὑμῖν Fritzsche}
\ligneRav{306}{}{\hspace{1.1em}\expo{καὶ}τὴν δρῶσαν καὶ ἀγορεύουσαν}{\textbf{306} \expo{καὶ}τὴν ἀγορεύουσαν \Rdeux}
\ligneRav{306}{}{\hspace{1.1em}\expo{καὶ}τὴν δρῶσαν καὶ ἀγορεύουσαν}{\textbf{306} \expo{καὶ}τὴν ἀγορεύουσαν \Rdeux}
\ligneRav{307}{}{\hspace{1.1em}τὰ βέλτιστα περὶ τὸν δῆμον \difference{τὸν Ἀθηναίων}}{\textbf{307} τῶν Ἀθηναίων R : \lecon{τὸν Ἀθηναίων} Grynaeus}
\ligneRav{307}{}{\hspace{1.1em}}{\textbf{307} τῶν Ἀθηναίων R : \lecon{τὸν Ἀθηναίων} Grynaeus}
\ligneRav{}{}{\hspace{1.1em}καὶ τὸν τῶν γυναικῶν}{}
\ligneRav{}{}{\hspace{1.1em}ταύτην νικᾶν·}{}
\ligneRav{310}{}{\hspace{1.1em}ταῦτ᾽ εὔχεσθε καὶ ὑμῖν αὐταῖς τἀγαθά·}{\textbf{310} \lecon{ὑμῖν} R : ἡμῖν \Mon}
\ligneRav{310}{}{\hspace{1.1em}}{\textbf{310} \lecon{ὑμῖν} R : ἡμῖν \Mon}
\ligneRav{}{}{\hspace{1.1em}ἰὴ παιὼν ἰὴ παιὼν ἰὴ παιών. χαίρωμεν·}{}
\ligneRav{312}{χορ/}{\hspace{1.1em}δεχόμεθα καὶ θεῶν γένος}{\textbf{312} δεχόμεσθα R : \lecon{-μεθα} Dindorf}
\ligneRav{313}{}{\hspace{1.1em}\difference{λιτόμεθα} ταῖσδ᾽ ἐπ᾽ εὐχαῖς}{\textbf{313} λιτόμεσθα R : \lecon{-μεθα} Bourdin\espApp λιτόμεσθα R : \lecon{λιτόμεθα} Bourdin}
\ligneRav{}{}{\hspace{1.1em}φανέντας ἐπιχαρῆναι·}{}
\ligneRav{315}{}{\hspace{1.1em}Ζεῦ μεγαλώνυμε χρυσολύρα τε}{}
\ligneRav{}{}{\hspace{1.1em}Δῆλον ὃς ἔχεις ἱεράν,}{}
\ligneRav{317}{}{\hspace{1.1em}καὶ σὺ \difference{παγκρατὲς} \difference{κόρα}}{\textbf{317} παμκρατὲς R : \lecon{παγκρατὲς} ed. pr. \espApp κόρᾳ R : \lecon{κόρα} \Mon}
\ligneRav{}{}{\hspace{1.1em}γλαυκῶπι χρυσόλογχε πόλιν}{}
\ligneRav{}{}{\hspace{1.1em}οἰκοῦσα περιμάχητον, ἐλθὲ δεῦρο·}{}
\ligneRav{320}{}{\hspace{1.1em}\difference{θηροφόνη} (θηροφ\inRas{ω}{ό}νε) παῖ Λατοῦς,}{\textbf{320} θηροφ\inRas{ω}{ό}νε παῖ-- \Rpc : \lecon{θηροφόνη} Hermann}
\ligneRav{}{}{\hspace{1.1em}Λατοῦς χρυσώπιδος ἔρνος,}{}
\ligneRav{}{}{\hspace{1.1em}σύ τε πόντιε σεμνὲ Πόσειδον}{}
\ligneRav{}{}{\hspace{1.1em}ἁλιμέδον}{}
\ligneRav{324}{}{\hspace{1.1em}προλιπὼν μυχὸν \difference{ἰχθυόεντα},}{\textbf{324} ἰχθυόεντ᾽ R, Suid. οι 181 : \lecon{ἰχθυόεντα} Wilamowitz}
\ligneRav{324}{}{\hspace{1.1em}οἰστροδόνητον}{\textbf{324} ἰχθυόεντ᾽ R, Suid. οι 181 : \lecon{ἰχθυόεντα} Wilamowitz}
\ligneRav{325}{}{\hspace{1.1em}Νηρέος \difference{εἰναλίου} τε κόραι}{\textbf{325} ἐναλίου R : ἐνάλιοί Fritzsche : \lecon{εἰναλίου} Bentley}
\ligneRav{326}{}{\hspace{1.1em}Νύμφαι τ᾽ ὀρείπλαγκτοι·}{\textbf{326} \lecon{ὀρείπλαγκτοι} R : ὀρίπλαγκτοι Austin}
\ligneRav{}{}{\hspace{1.1em}χρυσέα τε φόρμιγξ}{}
\ligneRav{}{}{\hspace{1.1em}ἰαχήσειεν ἐπ᾽ εὐχαῖς}{}
\ligneRav{329}{}{\hspace{1.1em}ἡμετέραις· τελέως δ᾽ ἐκκλησιάσαιμεν Ἀθηνῶν}{\textbf{329} Ἀθηναίων R : \lecon{Ἀθηνῶν} Reisig}
\ligneRav{330}{}{\hspace{1.1em}εὐγενεῖς γυναῖκες·}{}
\ligneRav{331}{κηρ/}{εὔχεσθε τοῖς θεοῖσι τοῖς Ὀλυμπίοις}{\textbf{331} (Κηρ.) R : \lecon{(Κορ.)} van Leeuwen}
\ligneRav{}{}{καὶ ταῖς Ὀλυμπίαισι καὶ τοῖς Πυθίοις}{}
\ligneRav{}{}{καὶ ταῖσι Πυθίαισι καὶ τοῖς Δηλίοις}{}
\ligneRav{334}{}{καὶ ταῖσι Δηλίαισι \difference{τοῖς τ᾽} ἄλλοις θεοῖς,}{\textbf{334} καὶ τοῖς R : \lecon{τοῖς τ᾽} Grynaeus}
\ligneRav{335}{}{εἴ τις ἐπιβουλεύει τι τῷ δήμῳ κακὸν}{}
\ligneRav{}{}{τῷ τῶν γυναικῶν ἢ 'πικηρυκεύεται}{}
}
\vspace{3pt}\noindent\rule{4.5cm}{0.4pt}\par\vspace{1.5pt}
{\fontsize{7.0pt}{8.3pt}\selectfont
\noindent\textbf{Loci paralleli (vv. 295–336) ---} \textbf{299} (Καλλιγενείᾳ $=$ Δημήτερι) Hesychius κ 472 Καλλιγένειαν ; cf. Photius κ 118 ($=$ fr. Aristophanes 331 Kassel-Austin) \espTest \textbf{300} (Κουροτρόφῳ $=$ Δημήτηρ) Hesychius κ 3865 κουροτρόφος ; scholium vet. in Hesiodus \italique{Opera et dies} 227-9 \espTest \textbf{322-4} (σύ τε -- οἰστροδόνητον) Suidas οι 181 (sed σεμνὲ et ἁλιμέδον omissis) \espTest \textbf{335-7} Satyrus \italique{Vit. Euripid.} ($=$ \italique{The Oxyrhynchus Papyri} IX 1176) fr. 39 col. XII 8-16 (Εὐριπιδίῳ pro Εὐριπίδῃ (?)) \espTest \textbf{336} ('πικηρυκεύεται) Photius ε 1580 ἐπικηρυκεύεται $=$ \Syn ε 682 $=$ Suidas ε 2375 $=$ \EtymMagnum 360.21 $=$ \EtymGen (B) ; cf. Hesychius ε 4853\par
}
\vspace{2.5pt}\noindent\rule{4.5cm}{0.4pt}\par\vspace{1.5pt}
{\fontsize{6.9pt}{8.2pt}\selectfont
\noindent\textbf{Suite (en commentaire) de l'apparat critique \& notes paléographiques (\texttt{f. 166r}) ---} \textbf{Témoin papyrologique :} \textbf{\PsiCinq{} (\italique{P.Oxy.} IX 1176, Satyrus)} cite les \textbf{vv. 335–337}. --- \textbf{Grattage de la rubrique d'interlocuteur au v. 295 :} \Rdeux{} a gratté Κῆρυξ (\inRas{Κῆρυξ}{} \Rdeux{} : \auDessus{κη}{ρυ} \Mon{}, tandis que $\Sigma^{\text{R}}$ explique γυνὴ μιμουμένη κήρυκα ; L. Hauseux édite Κορ.-Ἱέρεια). --- \textbf{Intrusions de gloses religieuses (vv. 297–298) :} Dans la prière en prose de la Héraut, R a incorporé dans le texte les deux gloses interlinéaires τῇ Δήμητρῃ καὶ τῇ Κόρῃ (glosant τοῖν Θεσμοφόροιν) et τῇ Γῇ (glosant τῇ Κουροτρόφῳ). --- \textbf{Ajout supralinéaire de \Rdeux{} (v. 306) :} \expo{καὶ}τὴν ἀγορεύουσαν. --- \textbf{Correction \italique{in rasura} au v. 320 :} θηροφ\inRas{ω}{ό}νε παῖ-- \Rpc{} ($\omega$ de \Rac{} gratté et corrigé en ό). --- \textbf{v. 295 :} \inRas{Κῆρυξ}{} \Rdeux{} : \auDessus{κη}{ρυ} \Mon{} : γυνὴ μιμουμένη κήρυκα $\Sigma$ : Κρίτυλλα Sommerstein : \lecon{(Κορ.-Ἱέρεια)} \lecon{Bibi} \espApp εὐφημί' R : \lecon{εὐφημία} Souda $\varepsilon\ 3797$. --- \textbf{v. 297 :} Θεσμοφόροιν τῇ Δήμητρῃ καὶ τῇ Κόρῃ R : \lecon{Θεσμοφόροιν} Blaydes (suppression de la glose). --- \textbf{v. 298 :} τῇ Κουροτρόφῳ τῇ Γῇ R : \lecon{τῇ Κουροτρόφῳ} Dobree ex $\Sigma^{\text{R}}$.\par
}
\end{minipage}
\end{adjustbox}
\end{adjustwidth}
\newpage
\enTeteD{Codex Ravennas 429 (Biblioteca Classense) --- Fac-similé du f. 166r (vv. 295–336)}
\begin{adjustwidth}{-1.5cm}{-3.0cm}
\begin{center}
\includegraphics[width=16.5cm,height=20.5cm,keepaspectratio]{/workspaces/Thesmophories/dossier_ravennas_thesmophories/folios_individuels_ff162v-177r/08_f166r_classense168r_vv0295-0336.jpg}\par
\vspace{2pt}
{\scriptsize\textbf{Codex Ravennas 429, f. 166r (Classense 168r)} --- ΠΑΡΟΔΟΣ, \textbf{vv. 295–336} ($\text{R}^1, \Rpc, \Rdeux$, \PsiUn).}
\end{center}
\end{adjustwidth}
\newpage
\clearpage\ifodd\value{page}\hbox{}\newpage\fi
\enTeteG{ΘΕΣΜΟΦΟΡΙΑΖΟΥΣΑΙ --- Texte grec établi (f. 166r, vv. 295–336)}
\enTeteD{LES THESMOPHORIES --- Traduction française (f. 166r, vv. 295–336)}
\begin{pages}
\begin{Leftside}
\beginnumbering
%% v. 295-311
\begin{groupe}%
%
\personnageGrec{Κορ.-Ἱέρεια}%
\typeVers{prose}\prose{εὐφημία ἔστω, |\apparat{295}{%\inRas{Κῆρυξ}{} \Rdeux : \auDessus{κη}{ρυ} \Mon : γυνὴ μιμουμένη κήρυκα Σ : Κρίτυλλα Sommerstein \espApp 
εὐφημί᾽ R : \lecon{εὐφημία} Suid. ε 3797 \ligneApp{297} Θεσμοφόροιν τῇ Δήμητρῃ καὶ τῇ Κόρῃ R : \lecon{Θ.} Blaydes \ligneApp{298} τῇ Κουροτρόφῳ τῇ Γῇ R : \lecon{τῇ Κ.} Dobree ex Σ \ligneApp{301} Χάρισιν R : \lecon{ταῖς Χ.} Meineke}%
\commentaire{295-311}{}{La coryphée-prêtresse commence la prière, en prose selon l'usage : (scholie) :%
\citComm{πεζῇ εὐφημίᾳ χρῶντα οἱ κωμικοὶ, ἐπειδὰν εὐχὴν ἢ ψήφισμα εἰσάγωσιν.}%
Cf. en \italique{Oiseaux} (vers 864-888) pour une prière analogue.}%
\ligne{296} εὐφημία ἔστω. |\\
\ligne{297}εὔχεσθε ταῖν Θεσμοφόροιν |\\
\ligne{299}καὶ τῷ Πλούτῳ καὶ τῇ Καλλιγενείᾳ |%\testimonia{299}{(Καλλιγενείᾳ $=$ Δημήτερι) Hesychius κ 472 Καλλιγένειαν ; cf. Photius κ 118 ($=$ fr. Aristophanes 331 Kassel-Austin)}
\\
καὶ τῇ Κουροτρόφῳ |%\testimonia{300}{(Κουροτρόφῳ $=$ Δημήτηρ) Hesychius κ 3865 κουροτρόφος ; scholium vet. in Hesiodus \italique{Opera et dies} 227-9}
\\
\ligne{301}καὶ τῷ Ἑρμῇ καὶ ταῖς Χάρισιν |\\
\ligne{302}ἐκκλησίαν τήνδε καὶ σύνοδον τὴν νῦν |\\
\ligne{303}κάλλιστα καὶ ἄριστα ποῆσαι, |\\
\ligne{304}πολυωφελῶς μὲν τῇ πόλει τῇ Ἀθηναίων, |\apparat{304}{πόλει R : \lecon{τῇ π.} Meineke}\\% \ligneApp{305} \lecon{ἡμῖν} R : ὑμῖν Fritzsche}\\
τυχηρῶς δ᾽ ἡμῖν αὐταῖς. |\\%\apparat{305}{\lecon{ἡμῖν} R : ὑμῖν Fritzsche}\\
\ligne{306}καὶ τὴν δρῶσαν καὶ ἀγορεύουσαν |\apparat{306}{\expo{καὶ}τὴν ἀγορεύουσαν \Rdeux}\\
\ligne{307}τὰ βέλτιστα περὶ τὸν δῆμον τὸν Ἀθηναίων |\apparat{307}{τῶν Ἀθηναίων R : \lecon{τὸν Ἀθηναίων} Grynaeus}\\
\ligne{308}καὶ τὸν τῶν γυναικῶν |\ligne{309} ταύτην νικᾶν. |
\commentaire{307-8}{τὸν δῆμον τῶν Ἀθηναίων καὶ τὸν τῶν γυναικῶν.}{Plaisanterie \italique{contre l'attente} (παρὰ προσδοκίαν) : « Le peuple des Athéniens et celui... des femmes », là où l'on attendait \grec{καὶ τὸν τῶν ξυμμάχων}, « et celui de ses alliés ».}\\
ταῦτ᾽ εὔχεσθε καὶ ὑμῖν αὐταῖς τἀγαθά. |\apparat{310}{\lecon{ὑμῖν} R : ἡμῖν \Mon}\\
ἰὴ παιὼν ἰὴ παιὼν ἰὴ παιών. χαίρωμεν.\\\passeLigne}%
\end{groupe}%
%% v. 312-330
\begin{groupe}%
%
\personnageGrec{Χορός}%
\typeVers{$2$ iamb.}\chantChoeur{\ligne{312}δεχόμεθα καὶ θεῶν γένος\apparat{312}{δεχόμεσθα R : \lecon{-μεθα} Dindorf \ligneApp{313} λιτόμεσθα R : \lecon{-μεθα} Bourdin}%
\commentaire{312-30}{Chant du chœur.}{Pour l'étude métrique de ce chant du chœur, l'on pourra consulter Lætitia Parker et son ouvrage de référence, \italique{The Songs of Aristophanes}\marqueur{1}. Elle parvient à faire apparaître comme une sorte de refrain scandé $- \cup \cup - \cup \cup - \cup \cup -$\marqueur{2} aux vers 315, 320 et 325.\\
Quoique la découverte soit extrêmement intéressante et encore à creuser, j'ai préféré pour ma part reprendre la colométrie de Prato, plus proche du manuscrit de Ravenne : «~\italien{Il Coro accoglie l'invito dell'Aralda e rivolge la sua preghiera alle divinità peculiari dell' ecclesia (maschile) escluse nella precedente invocazione in prosa e qui menzionate, non senza adattamenti, col consueto criterio selettivo, quasi per conferire un crisma di legalità maggiore alle decisioni finali.}~» (Prato.)}%
\debutNote{312-30}%
\note{1}{Lætitia P. E. Parker, \italique{The songs of Aristophanes}, Clarendon Press Oxford (1997), pp. 406-409.}%
\note{2}{Lætitia désigne ce vers comme un \italique{énoplien} (dactylo-épitrite prolongé) et Prato comme un \italique{alcméen}.}}%
\end{groupe}%
%% v. 312-330
\begin{groupe}%
\typeVers{$2$ iamb\cat}\chantChoeur{λιτόμεθα ταῖσδ᾽ ἐπ᾽ εὐχαῖς%\apparat{313}{λιτόμεσθα R : \lecon{λιτόμεθα} Bourdin}
}%
\end{groupe}%
%% v. 312-330
\begin{groupe}%
\typeVers{$2$ iamb\cat}\chantChoeur{φανέντας ἐπιχαρῆναι.}%
\end{groupe}%
%% v. 315
\begin{groupe}%
\typeVers{$2$ adon.}\chantChoeur{Ζεῦ μεγαλώνυμε χρυσολύρα τε}%
\end{groupe}%
%% v. 316
\begin{groupe}%
\typeVers{cr. $+$ chor.}\chantChoeur{Δῆλον ὃς ἔχεις ἱεράν,}%
\end{groupe}%
%% v. 317
\begin{groupe}%
\typeVers{lecyth.}\chantChoeur{καὶ σὺ παγκρατὲς κόρα\apparat{317}{παμκρατὲς R : \lecon{παγκρατὲς} ed. pr. \espApp κόρᾳ R : \lecon{κόρα} \Mon}}%
\end{groupe}%
%% v. 318
\begin{groupe}%
\typeVers{$2$ iamb.}\chantChoeur{γλαυκῶπι χρυσόλογχε πόλιν}%
\end{groupe}%
%% v. 319
\begin{groupe}%
\typeVers{$3$ iamb\cat}\chantChoeur{οἰκοῦσα περιμάχητον, ἐλθὲ δεῦρο.}%
\end{groupe}%
%% v. 320
\begin{groupe}%
\typeVers{alcm.}\chantChoeur{καὶ πολυώνυμε θηροφόνη\apparat{320}{θηροφ\inRas{ω}{ό}νε παῖ-- \Rpc : \lecon{θηροφόνη} Hermann}}%
\end{groupe}%
%% v. 321
\begin{groupe}%
\typeVers{paroem.}\chantChoeur{Λατοῦς χρυσώπιδος ἔρνος,}%
\end{groupe}%
%% v. 322
\begin{groupe}%
\typeVers{paroem.}\chantChoeur{σύ τε πόντιε σεμνὲ Πόσειδον%\testimonia{322-4}{(σύ τε -- οἰστροδόνητον) Suidas οι 181 (sed σεμνὲ et ἁλιμέδον omissis)}
}%
\end{groupe}%
%% v. 323
\begin{groupe}%
\typeVers{cret.}\chantChoeur{ἁλιμέδον}%
\end{groupe}%
%% v. 324a
\begin{groupe}%
\decoupeVers{a}\typeVers{paroem.}\chantChoeur{προλιπὼν μυχὸν ἰχθυόεντα,\apparat{324}{ἰχθυόεντ᾽ R, Suid. οι 181 : \lecon{ἰχθυόεντα} Wilamowitz \ligneApp{325} ἐναλίου R : ἐνάλιοί Fritzsche : \lecon{εἰναλίου} Bentley \ligneApp{326} \lecon{ὀρείπλαγκτοι} R : ὀρίπλαγκτοι Austin}}%
\end{groupe}%
%% v. 324b
\begin{groupe}%
\ligne{324}\decoupeVers{b}\typeVers{adon.}\chantChoeur{οἰστροδόνητον}%
\end{groupe}%
%% v. 325
\begin{groupe}%
\typeVers{alcm.}\chantChoeur{Νηρέος εἰναλίου τε κόραι}%
\end{groupe}%
%% v. 326
\begin{groupe}%
\typeVers{pros.}\chantChoeur{Νύμφαι τ᾽ ὀρείπλαγκτοι.}%
\end{groupe}%
%% v. 327
\begin{groupe}%
\typeVers{ith.}\chantChoeur{χρυσέα τε φόρμιγξ}%
\end{groupe}%
%% v. 328
\begin{groupe}%
\typeVers{ba. $+$ ad.}\chantChoeur{ἰαχήσειεν ἐπ᾽ εὐχαῖς}%
\end{groupe}%
%% v. 329
\begin{groupe}%
\typeVers{2 hem.}\chantChoeur{\small{ἡμετέραις· τελέως δ᾽ ἐκκλησιάσαιμεν Ἀθηνῶν}\apparat{329}{Ἀθηναίων R : \lecon{Ἀθηνῶν} Reisig}}%
\end{groupe}%
%% v. 330
\begin{groupe}%
\typeVers{ith.}\chantChoeur{εὐγενεῖς γυναῖκες.
\commentaire{330}{εὐγενεῖς γυναῖκες.}{Confirme, s'il était besoin, que seules les citoyennes athéniennes assistaient aux Thesmophories (cf. le vers 294 lorsque Mnésiloque renvoie Thratta).}\\\passeLigne}%
\end{groupe}%
%% v. 331-351 (part 1)
\begin{groupe}\personnageGrec{Κορυφαῖος}
\typeVers{3 ia.}\trimIamb{\ligne{331}εὔχεσθε τοῖς θεοῖσι τοῖς Ὀλυμπίοις\apparat{331}{(Κηρ.) R : \lecon{(Κορ.)} van Leeuwen}\\
καὶ ταῖς Ὀλυμπίαισι καὶ τοῖς Πυθίοις\\
καὶ ταῖσι Πυθίαισι καὶ τοῖς Δηλίοις\\
καὶ ταῖσι Δηλίαισι τοῖς τ᾽ ἄλλοις θεοῖς,\apparat{334}{καὶ τοῖς R : \lecon{τοῖς τ᾽} Grynaeus}\\
εἴ τις ἐπιβουλεύει τι τῷ δήμῳ κακὸν
\\
τῷ τῶν γυναικῶν ἢ 'πικηρυκεύεται
\\\passeLigne}%
\end{groupe}%
\endnumbering
\end{Leftside}
\begin{Rightside}
\beginnumbering
%% v. 295-311
\begin{groupe}%
\personnageFr{Coryphée}%
\proseFr{(La prêtresse :) Recueillez-vous, recueillez-vous.\\
Priez les Thesmophores,\\
et Ploutos et Calligénie,\testimonia{299}{(Καλλιγενείᾳ $=$ Δημήτερι) Hesychius κ 472 Καλλιγένειαν ; cf. Photius κ 118 ($=$ fr. Aristophanes 331 Kassel-Austin)}\\
et la Nourricière,\testimonia{300}{(Κουροτρόφῳ $=$ Δημήτηρ) Hesychius κ 3865 κουροτρόφος ; scholium vet. in Hesiodus \italique{Opera et dies} 227-9}\\
et Hermès et les Charites,\\
de rendre cette ecclésie et présente réunion\\
la meilleure et la plus avantageuse qu'il se pourra,\\
pour le plus grand profit de la cité des Athéniens\\
et pour notre propre bien.\\
Que celle qui aura agi et parlé\\
le mieux en faveur du peuple des Athéniens
et du peuple des femmes voie prévaloir ses avis.
Demandez aux dieux cette faveur, et pour vous-mêmes toutes les prospérités.
Io Péan ! Io Péan ! Soyons en Joie.\\\passeLigne}%
\end{groupe}%
%% v. 312-330
\begin{groupe}%
\personnageFr{Chœur}%
\chantChoeurFr{\ligne{312}Nous acceptons ! Et nous supplions}%
\end{groupe}%
%% v. 312-330
\begin{groupe}%
\chantChoeurFr{la gent des dieux d'apparaître}%
\end{groupe}%
%% v. 312-330
\begin{groupe}%
\chantChoeurFr{et de prendre en gré ces prières.}%
\end{groupe}%
%% v. 312-330
\begin{groupe}%
\chantChoeurFr{Zeus au grand renom, et toi, dieu à la lyre d'or,}%
\end{groupe}%
%% v. 312-330
\begin{groupe}%
\chantChoeurFr{qui tiens Délos la sainte,}%
\end{groupe}%
%% v. 312-330
\begin{groupe}%
\chantChoeurFr{et toi, vierge toute-puissante,}%
\end{groupe}%
%% v. 312-330
\begin{groupe}%
\chantChoeurFr{aux yeux glauques, à la lance d'or, qui habites }%
\end{groupe}%
%% v. 312-330
\begin{groupe}%
\chantChoeurFr{une cité enviée, viens ici. Et toi aussi, }%
\end{groupe}%
%% v. 312-330
\begin{groupe}%
\chantChoeurFr{vierge aux noms multiples, meurtrière de bêtes, }%
\end{groupe}%
%% v. 312-330
\begin{groupe}%
\chantChoeurFr{rejeton de Léto aux yeux d'or.}%
\end{groupe}%
%% v. 312-330
\begin{groupe}%
\chantChoeurFr{Et toi, auguste Poseidon, thalassien,\testimonia{322-4}{(σύ τε -- οἰστροδόνητον) Suidas οι 181 (sed σεμνὲ et ἁλιμέδον omissis)}}%
\end{groupe}%
%% v. 312-330
\begin{groupe}%
\chantChoeurFr{modérateur de la mer, }%
\end{groupe}%
%% v. 312-330
\begin{groupe}%
\chantChoeurFr{quitte le gouffre poissonneux }%
\end{groupe}%
%% v. 312-330
\begin{groupe}%
\chantChoeurFr{fouetté des vents.}%
\end{groupe}%
%% v. 312-330
\begin{groupe}%
\chantChoeurFr{Et vous, filles de Nérée le marin,}%
\end{groupe}%
%% v. 312-330
\begin{groupe}%
\chantChoeurFr{et vous, Nymphes des montagnes.}%
\end{groupe}%
%% v. 312-330
\begin{groupe}%
\chantChoeurFr{Que la phorminx d'or}%
\end{groupe}%
%% v. 312-330
\begin{groupe}%
\chantChoeurFr{accompagne nos prières, et puissions-nous}%
\end{groupe}%
%% v. 312-330
\begin{groupe}%
\chantChoeurFr{\ligne{329}mener à bonne fin notre ecclésie,}%
\end{groupe}%
%% v. 312-330
\begin{groupe}%
\chantChoeurFr{femmes de pure lignée athénienne.\\\passeLigne}%
\end{groupe}%
%% v. 331-351 (part 1)
\begin{groupe}\personnageFr{Coryphée}
\trimIambFr{\ligne{331}Invoquez les dieux et les déesses de l'Olympe,\\
les dieux et les déesses de Pytho,\\
les dieux et les déesses de Délos\\
et les autres dieux.\\
Si quelqu'un forme quelque dessein contre le peuple\testimonia{335-7}{Satyrus \italique{Vit. Euripid.} ($=$ \italique{The Oxyrhynchus Papyri} IX 1176) fr. 39 col. XII 8-16 (Εὐριπιδίῳ pro Εὐριπίδῃ (?))}\\
des femmes, fait des ouvertures\testimonia{336}{('πικηρυκεύεται) Photius ε 1580 ἐπικηρυκεύεται $=$ \Syn ε 682 $=$ Suidas ε 2375 $=$ \EtymMagnum 360.21 $=$ \EtymGen (B) ; cf. Hesychius ε 4853}\\\passeLigne}%
\end{groupe}%
\endnumbering
\end{Rightside}
\end{pages}
\Pages
\clearpage\ifodd\value{page}\hbox{}\newpage\fi
\enTeteG{ΠΑΡΟΔΟΣ --- Codex Ravennas 429, f. 166v (Classense 168v) --- vv. 337–379}
\begin{adjustwidth}{-1.5cm}{-3.0cm}
\begin{adjustbox}{max totalheight=21.0cm, max width=16.5cm, keepaspectratio}
\begin{minipage}{16.5cm}
{\fontsize{7.2pt}{8.5pt}\selectfont
\noindent\textbf{Folio 09 --- Codex Ravennas 429, \texttt{f. 166v} (\texttt{Classense 168v}) : ΠΑΡΟΔΟΣ, vv. 337–379}\hfill \textit{Témoins :} $\text{R}^1, \Rpc, \Rdeux$, \PsiCinq\par\vspace{2pt}
\hrule height 0.4pt\vspace{2.5pt}
}
{\fontsize{8.1pt}{9.5pt}\selectfont
\ligneRav{337}{χορ/}{†\difference{Εὐριπίδῃ Μήδοις τ᾽}† ἐπὶ βλάβῃ τινὶ}{\textbf{337} ευριπιδιω\douteux{ι}[. .]\douteux{υ}τ \PsiCinq : Εὐριπίδῃ Μήδοις R : \lecon{Εὐριπίδῃ Μήδοις τ᾽} Biset}
\ligneRav{}{}{τῇ τῶν γυναικῶν, ἢ τυραννεῖν ἐπινοεῖ}{}
\ligneRav{}{}{ἢ τὸν τύραννον συγκατάγειν, ἢ παιδίον}{}
\ligneRav{340}{}{ὑποβαλλομένης \difference{κατεῖπέ τις}, ἢ δούλη τινὸς}{\textbf{340} κατεῖπεν τις R : \lecon{κατεῖπέ τις} Brunck}
\ligneRav{341}{}{προαγωγὸς οὖσ᾽ ἐνετρύλλισεν τῷ δεσπότῃ}{\textbf{341} ἐνετρύλλισεν R, Suid. ε 1271 : \lecon{ἐνετρύλισεν} Dindorf}
\ligneRav{}{}{ἢ πεμπομένη τις ἀγγελίας ψευδεῖς φέρει,}{}
\ligneRav{}{}{ἢ μοιχὸς εἴ τις ἐξαπατᾷ ψευδῆ λέγων}{}
\ligneRav{344}{}{καὶ μὴ δίδωσιν \difference{ἃν} ὑπόσχηταί ποτε,}{\textbf{344} ἅ ἂν R : \lecon{ἃν} Faber}
\ligneRav{345}{}{ἢ δῶρά τις δίδωσι μοιχῷ γραῦς γυνή,}{}
\ligneRav{346}{}{ἢ καὶ δέχεται προδιδοῦσ᾽ \difference{ἑταίρᾳ} τὸν φίλον,}{\textbf{346} ἑταίρα R : \lecon{ἑταίρᾳ} Fritzsche ap Σ}
\ligneRav{347}{}{κεἴ τις κάπηλος ἢ καπηλὶς τοῦ χοὸς}{\textbf{347} \lecon{χοὸς} R, Suid. λ 394 : χοῶς Brunck}
\ligneRav{}{}{ἢ τῶν κοτυλῶν τὸ νόμισμα διαλυμαίνεται,}{}
\ligneRav{349}{}{κακῶς ἀπολέσθαι τοῦτον αὐτὸν \difference{κᾠκίαν}}{\textbf{349} κ' οἰκίαν R : \lecon{κᾠκίαν} Brunck}
\ligneRav{350}{}{ἀρᾶσθε, ταῖς δ᾽ ἄλλαισιν ὑμῖν τοὺς θεοὺς}{\textbf{350} \lecon{ὑμῖν} R : ἡμῖν \Mon}
\ligneRav{351}{}{εὔχεσθε πάσαις πολλὰ δοῦναι κἀγαθά·}{\textbf{351} \lecon{κἀγαθά} R : τἀγαθά Zanetti}
\ligneRav{352}{χορ/}{\hspace{1.1em}\difference{ξυνευχόμεσθα} τέλεα μὲν}{\textbf{352} ξυνευχόμεθα R : \lecon{ξυνευχόμεσθα} Daubruz}
\ligneRav{353}{}{\hspace{1.1em}πόλει τέλεα τε δήμῳ}{\textbf{353} \lecon{τε} R : δὲ \Mon}
\ligneRav{354}{}{\hspace{1.1em}τάδ᾽ εὔγματ᾽ \difference{ἀποτελεῖσθαι},}{\textbf{354} γενέσθαι R : ἐκγενέσθαι Dindorf : \lecon{ἀποτελεῖσθαι} Willems}
\ligneRav{355}{}{\hspace{1.1em}τὰ δ᾽ ἄρισθ᾽ ὅσαις προσήκει}{}
\ligneRav{}{}{\hspace{1.1em}νικᾶν λεγούσαις· ὁπόσαι δ᾽}{}
\ligneRav{357}{}{\hspace{1.1em}\difference{ἐξαπατῶσιν} παραβαίνουσί τε τοὺς}{\textbf{357} ἐξαπατῶσι R : \lecon{ἐξαπατῶσιν} Bentley}
\ligneRav{}{}{\hspace{1.1em}ὅρκους τοὺς νενομισμένους}{}
\ligneRav{360}{}{\hspace{1.1em}κερδῶν \difference{οὕνεκ᾽} ἐπὶ βλάβῃ,}{\textbf{360} ἕνεκ᾽ R : \lecon{οὕνεκ᾽} Bentley}
\ligneRav{}{}{\hspace{1.1em}ἢ ψηφίσματα καὶ νόμον}{}
\ligneRav{}{}{\hspace{1.1em}ζητοῦσ᾽ ἀντιμεθιστάναι,}{}
\ligneRav{}{}{\hspace{1.1em}τἀπόρρητά τε τοῖσιν ἐχ-}{}
\ligneRav{364}{}{\hspace{1.1em} θροῖς τοῖς ἡμετέροις \difference{λέγουσ᾽},}{\textbf{364} λέγουσιν R, Suid. α 3501 : \lecon{λέγουσ᾽} Bentley}
\ligneRav{365}{}{\hspace{1.1em}ἢ Μήδους ἐπάγουσι τῆς}{}
\ligneRav{}{}{\hspace{1.1em}χώρας οὕνεκ᾽ ἐπὶ βλάβῃ,}{}
\ligneRav{367}{}{\hspace{1.1em}\difference{ἀσεβοῦσ᾽ ἀδικοῦσί} τε τὴν πόλιν·}{\textbf{367} ἀσεβοῦσιν ἀδικοῦσίν R : \lecon{ἀσεβοῦσ᾽ ἀδικοῦσί} Bothe}
\ligneRav{}{}{\hspace{1.1em}ἀλλ᾽ ὦ παγκρατὲς}{}
\ligneRav{}{}{\hspace{1.1em}Ζεῦ ταῦτα κυρώσειας, ὥσθ᾽}{}
\ligneRav{370}{}{\hspace{1.1em}ἡμῖν θεοὺς παραστατεῖν}{}
\ligneRav{}{}{\hspace{1.1em}καίπερ γυναιξὶν οὔσαις·}{}
\ligneRav{372}{χορ/}{ἄκουε \difference{πᾶσ᾽}. ἔδοξε τῇ βουλῇ τάδε}{\textbf{372} πᾶς R : \lecon{πᾶσ᾽} Voss}
\ligneRav{}{}{τῇ τῶν γυναικῶν· Τιμόκλει᾽ ἐπεστάτει,}{}
\ligneRav{374}{}{Λύσιλλ᾽ ἐγραμμάτευεν, εἶπε Ζωστράτη·}{\textbf{374} Ἀρ[. . . .]\douteux{λ}(vel \douteux{a}, \douteux{δ}, \douteux{μ})ειa \PsiCinq, Ἀρ[χίκλ]ει᾽ suppl. Hunt : \lecon{Τιμόκλει᾽} R, Σ, Suid. ε 2616}
\ligneRav{375}{}{ἐκκλησίαν ποιεῖν ἕωθεν τῇ μέσῃ}{}
\ligneRav{376}{}{τῶν Θεσμοφορίων, ᾗ μάλισθ᾽ ἡμῖν \difference{σχολή},}{\textbf{376} σχολῇ R : \lecon{σχολή} Zanetti}
\ligneRav{}{}{καὶ χρηματίζειν πρῶτα περὶ Εὐριπίδου,}{}
\ligneRav{}{}{ὅ τι χρὴ παθεῖν ἐκεῖνον· ἀδικεῖν γὰρ δοκεῖ}{}
\ligneRav{}{}{ἡμῖν ἁπάσαις. τίς ἀγορεύειν βούλεται;}{}
}
\vspace{3pt}\noindent\rule{4.5cm}{0.4pt}\par\vspace{1.5pt}
{\fontsize{7.0pt}{8.3pt}\selectfont
\noindent\textbf{Loci paralleli (vv. 337–379) ---} \textbf{341} Suidas ε 1271 \espTest (προαγωγὸς) Hesychius π 3308 προαγωγός \espTest (ἐνετρύλλισεν) Zonaras 745 οὔ σ' ἐνετρύλλησεν \espTest \textbf{347-8} Suidas χ 394 (sed εἴ pro κεἴ) \espTest \textbf{363-4} (τἀπόρρητά -- λέγουσ᾽) Suidas α 3501 ἀπόρρητά τε τοῖσιν ἐχθροῖσι τοῖς ἡμετέροις λέγουσιν \espTest \textbf{369-71} (ὥσθ᾽ -- οὔσαις) Suidas π 444 \espTest \textbf{373-4} Satyrus \italique{Vit. Euripid.} ($=$ \italique{The Oxyrhynchus Papyri} IX 1176) fr. 39 col. XII 1 (vid. ad 335-7 ; sed Ἀρ[χίκλ]ει᾽ pro Τιμόκλει᾽, suppl. Hunt) \espTest (Τιμόκλει᾽ -- ἐγραμμάτευεν) Suidas ε 2616 (sed ἐπιστάτει pro ἐπεστάτει et ἐγραμμάτευε pro ἐγραμμάτευεν) \espTest \textbf{377-8} (καὶ χρηματίζειν -- παθεῖν) Suidas χ 481 ; Lexicon Vindobonense 190.10\par
}
\vspace{2.5pt}\noindent\rule{4.5cm}{0.4pt}\par\vspace{1.5pt}
{\fontsize{6.9pt}{8.2pt}\selectfont
\noindent\textbf{Suite (en commentaire) de l'apparat critique \& notes paléographiques (\texttt{f. 166v}) ---} \textbf{Témoin papyrologique :} \textbf{\PsiCinq{} (\italique{P.Oxy.} IX 1176, Satyrus)} cite les \textbf{vv. 335–337} et \textbf{vv. 373–375}. --- \textbf{v. 346 :} ἑταίρα R : \lecon{ἑταίρᾳ} Fritzsche ap. $\Sigma^{\text{R}}$. --- \textbf{v. 354 :} γενέσθαι R (ἐκγενέσθαι Dindorf) : \lecon{ἀποτελεῖσθαι} Alphonse Willems (adopté par L. Hauseux : γενέσθαι est une glose banale substituée au verbe rituel ἀποτελεῖσθαι). --- \textbf{v. 374 :} Ἀρ[....]λεια \PsiCinq{} (Ἀρ[χίκλ]ει' suppl. Hunt) : \lecon{Τιμόκλει'} R, $\Sigma^{\text{R}}$, Souda $\varepsilon\ 2616$ (maintenu par L. Hauseux).\par
}
\end{minipage}
\end{adjustbox}
\end{adjustwidth}
\newpage
\enTeteD{Codex Ravennas 429 (Biblioteca Classense) --- Fac-similé du f. 166v (vv. 337–379)}
\begin{adjustwidth}{-1.5cm}{-3.0cm}
\begin{center}
\includegraphics[width=16.5cm,height=20.5cm,keepaspectratio]{/workspaces/Thesmophories/dossier_ravennas_thesmophories/folios_individuels_ff162v-177r/09_f166v_classense168v_vv0337-0379.jpg}\par
\vspace{2pt}
{\scriptsize\textbf{Codex Ravennas 429, f. 166v (Classense 168v)} --- ΠΑΡΟΔΟΣ, \textbf{vv. 337–379} ($\text{R}^1, \Rpc, \Rdeux$, \PsiCinq).}
\end{center}
\end{adjustwidth}
\newpage
\clearpage\ifodd\value{page}\hbox{}\newpage\fi
\enTeteG{ΘΕΣΜΟΦΟΡΙΑΖΟΥΣΑΙ --- Texte grec établi (f. 166v, vv. 337–379)}
\enTeteD{LES THESMOPHORIES --- Traduction française (f. 166v, vv. 337–379)}
\begin{pages}
\begin{Leftside}
\beginnumbering
%% v. 337+ (part 2)
\begin{groupe}
\trimIamb{\ligne{337}†Εὐριπίδῃ Μήδοις τ᾽† ἐπὶ βλάβῃ τινὶ\apparat{337}{ευριπιδιω\douteux{ι}[. .]\douteux{υ}τ \PsiCinq : Εὐριπίδῃ Μήδοις R : \lecon{Εὐριπίδῃ Μήδοις τ᾽} Biset}% Biset : \lecon{Εὐριπίδῃ ᾧ ταῦτ᾽} Bibi}%
\commentaire{336-7}{ἢ 'πικηρυκεύεται / Εὐριπίδῃ Μήδοις ἐπὶ βλάβῃ τινὶ.}{Dans sa parodie d'ecclésie, Aristophane ne manque pas de reprendre les imprécations contre les Mèdes, qui étaient de règle comme on le sait par Isocrate (\italique{Panégyrique}, 157) :
\citComm{Καὶ πολλῶν μὲν οἱ πατέρες ἡμῶν μηδισμοῦ θάνατον κατέγνωσαν, ἐν δὲ τοῖς συλλόγοις ἔτι καὶ νῦν ἀρὰς ποιοῦνται, πρὶν ἄλλο τι χρηματίζειν, εἴ τις ἐπικηρυκεύεται Πέρσαις τῶν πολιτῶν.\\
\francais{Nos pères ont condamné à mort plusieurs citoyens pour leur adhésion au parti des Perses (μηδισμός). Encore aujourd'hui, dans nos assemblées, avant de traiter aucune affaire, on prononce des imprécations coutre celui des citoyens qui fait des ouvertures aux Mèdes.}}
L'ennemi héréditaire, pour la féminine assemblée, est remplacé -- on s'y attendait ! -- par Euripide.
Remplacé ? pas tout-à-fait ; car au côté d'Εὐριπίδῃ continuent de subsister étrangement les Μήδοις... détail qui fleure bon son interpolation : le spectateur athénien n'avait pas besoin qu'on lui souligne ainsi le trait comique...
% (Plutôt que d'une plaisanterie παρὰ προσδοκίαν, l'on devrait plus proprement parler d'une plaisanterie μετὰ προσδοκίας.)\\
%Interpolation, alors ? c'est ce que nous allons pouvoir essayer de montrer grâce à un papyrus d'Oxyrhynque !\\
Or nous possédons un papyrus d'Oxyrhynque\marqueur{1} du \nieme{II} siècle après Jésus-Christ\marqueur{2} qui comporte une citation des vers 335-337 des \italique{Thesmophories}. Malheureusement, le vers qui nous intéresse est lacunaire :%
\citComm{ευριπιδιω\douteux{ι}[. .]\douteux{υ}\marqueur{3}τ|επιβλαβη\douteux{ι}τι|νι}%
Hunt avait déjà remarqué en son temps que la lacune était bien insuffisante pour qu'ait pu s'y trouver Μήδοις τ᾽\marqueur{4}.%
%\eff{La lecture probable d'un υ\marqueur{4} avant le τ laisse à penser qu'il faut compléter ainsi la lacune : [τα]ῦτ᾽\marqueur{5}.\\
%Il suffit alors de faire la simple hypothèse de la chute d'un êta -- je la fais : ευριπιδ<η>ι pour que le passage devienne intelligible :%
%\citComm{ηεπικη|ρυκευετ[αι]ευρι|πιδ<η>ιω\douteux{ι}[τα]\douteux{υ}τ|επιβλαβη\douteux{ι}τι|νι [τῇ τῶν γυναικῶν]\\
%ἢ 'πικηρυκεύεται / Εὐριπίδῃ ᾧ ou bien ws ταῦτ᾽ ἐπὶ βλάβῃ τινὶ / τῇ τῶν γυναικῶν\\
%ou s'il envoie un ambassadeur à Euripide qui [utilise] ces informations au détriment des femmes}
%(Avec le redoublement Εὐριπίδῃ / ᾧ repris plus loin avec τινὶ / τῇ ; pastiche du style de la procédure ecclésiastique.)
}%
\debutNote{336-7}%
\note{1}{\italique{The Oxyrhynchus Papyri}, IX (1912), ed. Arthur S. Dunt. Il s'agit du papyrus n° 1176, fragment 39, colonne XII 8-16. Ce papyrus date du \nieme{II} siècle ap. J.-C. et comporte un fragment d'une \italique{Vie} d'Euripide par Satyros de Callatis. Je me réfère aussi à l'édition ainsi qu'au commentaire de ce fragment par Stefan Schorn, \italique{Satyros aus Kallatis} (2002).}%
\note{2}{Quand notre manuscrit de Ravenne est beaucoup plus tardif : \nieme{IX} ou \nieme{X} s. ap. J.-C.}%
\note{3}{Schorn indique que la présence d'un \douteux{c} est moins probable, et que ce n'est assurément pas un \douteux{ε}.}%
\note{4}{Ce qui n'a pas empêché les éditeurs depuis de prétendre dans leur apparat critique que ce papyrus confirmait l'addition du τ᾽ de Biset : Εὐριπίδῃ Μήδοις <τ᾽>... sauf que le τ identifié par Hunt ne pouvait en aucun cas se rattacher aux Μήδοις, absents du papyrus.}%
%\note{5}{On trouve l'expression ταῦτʼ ἐπικηρυκεύομαι au vers 1163, lorsqu'Euripide conclue la paix avec les femmes ; ταῦτα δι᾽ ἀγγέλων ἐπεκηρυκεύετο et ταῦτά οἱ πέμψας ὁ Σιτάλκης ἐπεκηρυκεύετο ou encore ταῦτα μὲν ἐπεκηρυκεύσατο τοῖσι Δηλίοισι chez Hérodote (I, 60 ; IV, 80 ; VI, 97).}
\\
τῇ τῶν γυναικῶν, ἢ τυραννεῖν ἐπινοεῖ\\
ἢ τὸν τύραννον συγκατάγειν, ἢ παιδίον
\commentaire{338-9}{ἢ τυραννεῖν ἐπινοεῖ ἢ τὸν τύραννον συγκατάγειν.}{Citation quasi-textuelle d'une ancienne loi, peut-être de Dracon, rapportée dans la \italique{Constitution d'Athènes} (XVI, 10) :
\citComm{ἦσαν δὲ καὶ τοῖς Ἀθηναίοις οἱ περὶ τῶν τυράννων νόμοι πρᾷοι κατ᾽ ἐκείνους τοὺς καιρούς, οἵ τ᾽ ἄλλοι καὶ δὴ καὶ ὁ μάλιστα καθήκων πρὸς τὴν τῆς τυραννίδος <κατάστασιν>. νόμος γὰρ αὐτοῖς ἦν ὅδε. `θέσμια τάδε Ἀθηναίων καὶ πάτρια· ἐάν τινες τυραννεῖν ἐπανιστῶνται ἐπὶ τυραννίδι, ἢ συγκαθιστῇ τὴν τυραννίδα, ἄτιμον εἶναι καὶ αὐτὸν καὶ γένος.'\\
\francais{Les lois des Athéniens sur la tyrannie étaient alors [à l'époque de Pisitrate] peu sévères, celle surtout qui visait les entreprises des tyrans et dont voici le texte : « Les lois athéniennes établies par nos pères portent que : quiconque aspire à la tyrannie ou forme un complot pour l'établir, sera frappé d'atimie dans sa personne et dans sa race. »}}%
Étant donnée l'époque troublée de la représentation des \italique{Thesmophories}, cette loi devait être sur toutes les bouches ; pour preuve, la même année que le représentation des \italique{Thesmophories}, en 410/09\marqueur{1}, Démophante proposa un décret inspiré de cette dernière garantissant l'impunité à l'assassin de toute personne complice de crime de lèse-démocratie (cette loi fut gravée sur une stèle fameuse à Athènes qu'Andocide cite dans son \italique{Discours sur les mystères}, 97) :%
\citComm{\grec{καὶ ἐάν τις τυραννεῖν ἐπαναστῇ ἢ τὸν τύραννον συγκαταστήσῃ}.}%
}%
\debutNote{338-9}%
\note{1}{C. Carey, Ifigeneia G., Brenda G.-W., \italique{Use and abuse of law in the Athenian courts}, p. 295.}\\
ὑποβαλλομένης κατεῖπέ τις, ἢ δούλη τινὸς\apparat{340}{κατεῖπεν τις R : \lecon{κατεῖπέ τις} Brunck % : del. τις Bentley 
\ligneApp{341} ἐνετρύλλισεν R, Suid. ε 1271% (sed ἐνετρύλλισε)
 : \lecon{ἐνετρύλισεν} Dindorf}\\
προαγωγὸς οὖσ᾽ ἐνετρύλλισεν τῷ δεσπότῃ%\testimonia{341}{Suidas ε 1271 \espTest (προαγωγὸς) Hesychius π 3308 προαγωγός \espTest (ἐνετρύλλισεν) Zonaras 745 οὔ σ' ἐνετρύλλησεν}
\\
ἢ πεμπομένη τις ἀγγελίας ψευδεῖς φέρει,}
\end{groupe}%
%% v. 331-351
\begin{groupe}%
%
\trimIamb{ἢ μοιχὸς εἴ τις ἐξαπατᾷ ψευδῆ λέγων\\
καὶ μὴ δίδωσιν ἃν ὑπόσχηταί ποτε,\apparat{344}{ἅ ἂν R : \lecon{ἃν} Faber}\\
ἢ δῶρά τις δίδωσι μοιχῷ γραῦς γυνή,\\
ἢ καὶ δέχεται προδιδοῦσ᾽ ἑταίρᾳ τὸν φίλον,\apparat{346}{ἑταίρα R : \lecon{ἑταίρᾳ} Fritzsche ap Σ \ligneApp{347} \lecon{χοὸς} R, Suid. λ 394 : χοῶς Brunck}%
\commentaire{346}{ἢ καὶ δέχεται προδιδοῦσ᾽ ἑταίρα τὸν φίλον.}{Je reprends la correction de Fritzsche qui substitue avec bonheur \grec{ἑταίρᾳ} à \grec{ἑταίρα}. Le scholiaste précise d'ailleurs qu'il y a différentes versions dans les manuscrits : certaines avec le iota souscrit, d'autres sans : 
\citComm{γρ<άφεται> ἔν τισιν ἑταίρα καὶ ἔστιν ἕτερος νοῦς [...] ἐὰν μὲν ἄνευ τοῦ ῑ γρ<άφηται> ἑταίρα [...] ἐὰν δὲ σὺν τῷ ῑ}
La leçon sans iota est incohérente : en quoi une hétaïre, dont c'est le métier de recevoir des dons en échange de faveurs \italique{trahirait-elle} (\grec{προδιδοῦσα}) son amant ce faisant ? 
(Willems, sentant bien la difficulté, s'éloigne en traduisant : « à l'insu de son amant ».) 
Seule donc la première version est satisfaisante : il est question au vers précédent d'une vieille faisant des cadeaux pour s'attirer les faveurs d'un galant (thème qu'Aristophane développe dans le \italique{Ploutos}, vv. 959 sqq.) ; après la lubricité, le lucre : cette fois la vieille reçoit des cadeaux pour servir d'entremetteuse en livrant (προδιδοῦσ᾽) le soupirant (τὸν φίλον) à une femme irrégulière (ἑταίρᾳ). (Rôle de vieille procureuse que tiendra d'ailleurs Euripide à la fin de la pièce.)% thème de la vieille procureuse qui sera repris à la fin de la pièce\eff{ (v. ....)}.
}\\
κεἴ τις κάπηλος ἢ καπηλὶς τοῦ χοὸς%\testimonia{347-8}{Suidas χ 394 (sed εἴ pro κεἴ)}
\\%\apparat{347}{\lecon{χοὸς} R, Suid. λ 394 : χοῶς Brunck}\\
ἢ τῶν κοτυλῶν τὸ νόμισμα διαλυμαίνεται,%
\commentaire{347-8}{κεἴ τις κάπηλος ἢ καπηλὶς τοῦ χοὸς ἢ τῶν κοτυλῶν τὸ νόμισμα διαλυμαίνεται.}{Les débitantes de mauvais aloi sont brocardées de la même façon par Aristophane dans le \italique{Ploutos}, v. 435-6 :
\citComm{Βλεψίδημος — ἆρ᾽ ἐστὶν ἡ καπηλὶς ἡ ᾽κ τῶν γειτόνων,\\
ἣ ταῖς κοτύλαις ἀεί με διαλυμαίνεται;}%
}\\
κακῶς ἀπολέσθαι τοῦτον αὐτὸν κᾠκίαν\apparat{349}{κ' οἰκίαν R : \lecon{κᾠκίαν} Brunck}\\% \ligneApp{350} \lecon{ὑμῖν} R : ἡμῖν \Mon}\\
ἀρᾶσθε, ταῖς δ᾽ ἄλλαισιν ὑμῖν τοὺς θεοὺς\\%\apparat{350}{\lecon{ὑμῖν} R : ἡμῖν \Mon}\\
εὔχεσθε πάσαις πολλὰ δοῦναι κἀγαθά.\apparat{351}{\lecon{κἀγαθά} R : τἀγαθά Zanetti}\\\passeLigne}%
\end{groupe}%
%% v. 352-354
\begin{groupe}%
%
\personnageGrec{Χορός}%
\typeVers{$2$ iamb.}\chantChoeur{ξυνευχόμεσθα τέλεα μὲν\apparat{352}{ξυνευχόμεθα R : \lecon{ξυνευχόμεσθα} Daubruz}%
\commentaire{352-71}{}{Je reprends encore une fois la colométrie de Prato pour ce chant du chœur.}}%
\end{groupe}%
%% v. 353
\begin{groupe}%
\typeVers{$2$ iamb\cat}\chantChoeur{πόλει τέλεα τε δήμῳ%\apparat{353}{\lecon{τε} R : δὲ \Mon}
}%
\end{groupe}%
%% v. 354
\begin{groupe}%
\typeVers{$2$ iamb\cat}\chantChoeur{τάδ᾽ εὔγματ᾽ ἀποτελεῖσθαι,\apparat{354}{γενέσθαι R : ἐκγενέσθαι Dindorf : \lecon{ἀποτελεῖσθαι} Willems}%
\commentaire{354}{τάδ᾽ εὔγματα γενέσθαι.}{J'adopte la correction de Willems qui substitue ἀποτελεῖσθαι à γενέσθαι (\nieme{9} note, reproduite en annexe).}}%
\end{groupe}%
%% v. 355
\begin{groupe}%
\typeVers{$2$ iamb\cat}\chantChoeur{τὰ δ᾽ ἄρισθ᾽ ὅσαις προσήκει}%
\end{groupe}%
%% v. 356
\begin{groupe}%
\typeVers{ia. $+$ chor.}\chantChoeur{νικᾶν λεγούσαις· ὁπόσαι δ᾽
\commentaire{356-67}{}{Sur tout ce passage, vers 356 à 367, pour le texte et l'interprétation, allez voir la \nieme{10} note d'Alphonse Willems qui fut le point de départ à l'excellent article de Taillardat \& Blanc reproduit ci-après.}}%
\end{groupe}%
%% v. 357-8
\begin{groupe}%
\ligne{357}\decoupeVers{8}\typeVers{$3$ chor.}\chantChoeur{ἐξαπατῶσιν παραβαίνουσί τε τοὺς\apparat{357}{ἐξαπατῶσι R : \lecon{ἐξαπατῶσιν} Bentley}}%
\end{groupe}%
%% v. 359
\begin{groupe}%
\ligne{359}\typeVers{glyc.}\chantChoeur{ὅρκους τοὺς νενομισμένους}%
\end{groupe}%
%% v. 360
\begin{groupe}%
\typeVers{glyc.}\chantChoeur{κερδῶν οὕνεκ᾽ ἐπὶ βλάβῃ,\apparat{360}{ἕνεκ᾽ R : \lecon{οὕνεκ᾽} Bentley}}%
\end{groupe}%
%% v. 361
\begin{groupe}%
\typeVers{glyc.}\chantChoeur{ἢ ψηφίσματα καὶ νόμον}%
\end{groupe}%
%% v. 362
\begin{groupe}%
\typeVers{glyc.}\chantChoeur{ζητοῦσ᾽ ἀντιμεθιστάναι,}%
\end{groupe}%
%% v. 363
\begin{groupe}%
\typeVers{glyc.}\chantChoeur{τἀπόρρητά τε τοῖσιν ἐχ-%\testimonia{363-4}{(τἀπόρρητά -- λέγουσ᾽) Suidas α 3501 ἀπόρρητά τε τοῖσιν ἐχθροῖσι τοῖς ἡμετέροις λέγουσιν}
}%
\end{groupe}%
%% v. 364
\begin{groupe}%
\typeVers{glyc.}\chantChoeur{\alinea θροῖς τοῖς ἡμετέροις λέγουσ᾽,\apparat{364}{λέγουσιν R, Suid. α 3501 : \lecon{λέγουσ᾽} Bentley}}%
\end{groupe}%
%% v. 365
\begin{groupe}%
\typeVers{glyc.}\chantChoeur{ἢ Μήδους ἐπάγουσι τῆς}%
\end{groupe}%
%% v. 366
\begin{groupe}%
\typeVers{glyc.}\chantChoeur{χώρας οὕνεκ᾽ ἐπὶ βλάβῃ,
\commentaire{365-6}{ἢ Μήδους ἐπάγουσι τῆς χώρας / οὕνεκ᾽ ἐπὶ βλάβῃ.}{Voilà deux vers bien malmenés et retouchés à l'envi par divers éditeurs d'Aristophane. On doit à Jean Taillardat et Alain Blanc d'avoir définitivement confirmé la leçon du manuscrit de Ravenne, tout en y relevant un jeu de mots jusqu'alors passé à l'as. La lecture de leur article\marqueur{1} est si rapicolante que je me ferais scrupule d'y taillarder, en voici donc un très copieux extrait : « Pendant plusieurs siècles, l'interprétation des vers 365-366 a été controversée. Les savants ont d'abord été d'accord pour chercher dans τῆς χώρας, v. 366, une désignation géographique (“l'Attique”) et pour comprendre, comme Andrea Divo, “\latin{quaecunque ... Persas inducunt in regionem nostram}”. Mais traduire τῆς χώρας οὕνεκα par “\latin{in regionem nostram}”\marqueur{2}, c'est prêter à οὕνεκα un sens qu'il n'a jamais. D'où diverses conjectures\marqueur{3} et, enfin, la \italique{crux philologorum} dont Van Leeuwen a marqué les vers 365-367. C'est à Willems\marqueur{4} que revient le mérite d'avoir fait progresser, d'une façon décisive, l'interprétation de ce texte : il a rendu leur juste sens aux leçons transmises, ἐπάγειν et οὕνεκα, en traduisant “toutes celles ... qui sollicitent les Mèdes en considération de notre pays, mais à notre dam”. Il commente ainsi sa version (t. II, p. 551-553) : “ Ici il s'agit de ceux qui sollicitent les Mèdes « dans l'intérêt du pays », τῆς χώρας οὕνεκα [ἐπὶ βλάβῃ]. « Dans l'intérêt du pays » est tout simplement un euphémisme équivalent à « contre les Péloponnésiens », ce qu'il fallait éviter de dire. On sait, en effet, que les oligarques intriguaient auprès des satrapes d'Asie Mineure pour détacher le Grand Roi de l'alliance de Sparte et le faire entrer dans celle d'Athènes (Thucydide VIII, 53). Mais l'on sait non moins bien que cette politique indignait Aristophane ». 
Que Μήδους ἐπάγουσι τῆς χώρας οὕνεκα puisse signifier « elles sollicitent l'alliance des Mèdes pour l'amour de notre pays », cela est sûr\marqueur{5} ; mais on ne voit pas qu'il y ait là « un euphémisme » lié à une sorte d'interdit. Mieux vaudrait dire que le chœur rapporte ironiquement les propos de toute femme « qui cherche à traiter avec lui (= l'ennemi) dans l'intérêt \italique{soi-disant d'Athènes} »\marqueur{5}. On peut admettre en effet que cette phrase transpose adroitement le prétexte donné par les intrigantes (Μήδους ἐπάγομεν τῆς χώρας οὕνεκα) : l'enjambement τῆς | χώρας οὕνεκα invitait le chœur à faire un sort au mot χώρας ainsi mis en relief, en sorte qu'une intonation particulière (ironique) des acteurs pouvait suffire à faire comprendre “pour l'amour de notre pays, disent-elles”.}%
\note{1}{«~Aristophane, ``Thesmophories'', v. 366 : τῆς χώρας οὕνεκα et l'ambition des oligarques dans l'hiver de 412-411~», in \italique{Revue de philologie, de littérature et d'histoire ancienne}, t. LX, 1986, fasc. 2, p. 183-186}%
\note{2}{Andreas Divus, \italique{Aristophanis comoediae undecim e graeco in latinum ad verbum translatae}, Basiliae, 1552. Brunck, en 1783, est même plus précis : « quaecunque ... Medos inducunt in regionem nostram, ad eam perdendam » (Brunck lit ἐπάγουσι τῆς χώρας ἐπὶ βλάβῃ, en supprimant οὕνεκα).}%
\note{3}{Par exemple : ἐπάγουσι τῆς | χώρας ἀσεβοῦσιν (Dobrée, suivi par Dindorf et Blaydes ; οὕνεκ᾽ ἐπὶ βλάβῃ disparaît) ; ἐπάγουσι τῆς | χώρας ἐπὶ βλάβῃ (Brunck) ; ἐπάγουσι χώρας | κερδῶν οὕνεκα ἐπὶ βλάβῃ (Reiske) ; ἐπάγουσι τῆς | ἔχθρας οὕνεκ᾽ ἐπὶ βλάβῃ (Hermann) ; ἐπάγουσι τῆς | ὥρας οὕνεκ᾽ ἐπὶ βλάβῃ (Fritzsche) ; ἐπάγουσι γῇ | κερδῶν οὕνεκ᾽ ἐπὶ βλάβῃ (van Velsen, suivi par Hall-Geldart et Rogers, ce dernier « with some hesitation ») ; τῆς | χώρας οὐ μεμελημένοι (C. Austin, voir ci-dessous, n. 5). Cet acharnement à modifier le texte transmis par le Ravennas (seul manuscrit des \italique{Thesmophories}) est d'autant plus étrange qu'il est métriquement sain : les vers 359-366 sont des glyconiens. }%
\note{4}{Il s'agit de la note 10 de Willems, reproduite en annexe.}%
\note{5}{Malgré C. Austin (\italique{P. C. Ph. S.}, 20, 1974, p. 1-2) qui refuse cette interprétation et corrige le texte (voir note 3).}%
\commentaire{}{}{\ \ \ Il y a cependant autre chose dans ce texte : un jeu de mots sur τῆς χώρας, plaisanterie dont personne ne s'est encore avisé. En effet, le parallélisme éclatant, et sûrement dû à Aristophane lui-même, de κερδῶν οὕνεκ᾽ ἐπὶ βλάβῃ (v. 360) et de τῆς χώρας οὕνεκ᾽ ἐπὶ βλάβῃ (v. 366) incite à chercher dans τῆς χώρας, autant que le nom de l'Attique, une variation de la notion exprimée par κερδῶν ; car le chœur peut faire allusion, ici comme là, à l'intérêt personnel et à l'avidité des gens qu'il maudit. Or le mot χώρα désigne, chez plusieurs auteurs, “la place” qu'on occupe dans la vie, par exemple “le rang”, qu'on a dans la hiérarchie sociale ou militaire\marqueur{7}. Dans les \italique{Thesmophories}, v. 366, c'est de rang dans la hiérarchie politique et sociale qu'il s'agit ; si bien que τῆς χώρας οὕνεκα, signifie à la fois “pour l'amour de notre pays (disent-elles)” et “pour se pousser aux (bonnes) places”\marqueur{8}, littéralement “en vue de leur place personnelle”\marqueur{9}. 
Plus que le cas général des σπουδαρχίδαι, “les sires de Brigue-charges”\marqueur{10}, Aristophane vise ici une circonstance particulière. Il pense, comme Willems l'a dit, aux intrigues ambitieuses d'Alcibiade\marqueur{11} et de Pisandre qui, pendant l'hiver 412-411, prétendaient arracher la Perse à son alliance spartiate et cherchaient à obtenir, pour Athènes, l'appui et l'or du Grand Roi\marqueur{12}, pour eux-mêmes, le pouvoir à Athènes. Et que fait Pisandre, à la fin de 412, quand le peuple le délègue auprès d'Alcibiade et de Tissapherne ? Avant de quitter l'Attique pour l'Asie Mineure, “il rend visite aux sociétés secrètes déjà établies à Athènes pour s'occuper des procès et des magistratures”, ξυνωμοσίας, αἵπερ ἐτύγχανον πρότερον ἐν τῇ πόλει οὖσαι ἐπὶ δίκαις καὶ ἀρχαῖς, ἁπάσας ἐπελθὼν... (Thuc. VIII, 54, 4). C'est dire que les comploteurs avaient ouvert la chasse aux places. Les oligarques se consultent ἐπὶ δίκαις καὶ ἀρχαῖς et tout le monde le sait à Athènes : deux mois avant les \italique{Thesmophories}, Aristophane a déjà dénoncé “ces gens qui s'agglomèrent et font bourre pour se pousser aux charges” τούς γε συνισταμένους τούτους καὶ τοὺς πιλοῦντας ἑαυτοὺς | ἐπὶ ταῖς ἀρχαῖσι (\italique{Lysistrata}, v. 577 sq.). Ces ambitieux sont concussionnaires et vénaux : “c'est pour être en mesure de voler que Pisandre et ceux qui ne pensent qu'aux charges créent à chaque instant grabuge et désordre. Qu'ils fassent donc, pour cela, tout ce qu'ils voudront ; mais cet argent-ci\marqueur{13}, il n'y a plus de danger qu'ils l'emportent”, 
\citComm{ἵνα γὰρ Πείσανδρος ἔχοι κλέπτειν χοἰ ταῖς ἀρχαῖς ἐπέχοντες,\\
ἀεί τινα κορκορυγὴν ἐκύκων. οἱ δ᾽ οὖν τοῦδ᾽ οὕνεκα δρώντων\\
ὅ τι βούλονται: τὸ γὰρ ἀργύριον τοῦτ᾽ οὐκέτι μὴ καθέλωσιν.\\
(\italique{Lysistrata}, v. 490-492).}%
Ἐπὶ (δίκαις καὶ) ἀρχαῖς, dit Thucydide ; ἐπὶ ταῖς ἀρχαῖσι, οἱ ταῖς ἀρχαῖς ἐπέχοντες, dit Aristophane. Peut-on rêver plus belles gloses au τῆς χώρας οὕνεκα des \italique{Thesmophories} ? Aux yeux d'Aristophane, tremper dans la grenouillère des combinaisons louches de l'hiver de 412-411, c'est donc prétendre à faire d'une pierre deux coups : devenir un personnage important à Athènes et s'enrichir en profitant de tous les avantages, licites ou non\marqueur{14}, attachés aux ἀρχαί ; c'est tout cela que résume τῆς χώρας οὕνεκ᾽ ἐπὶ βλάβῃ, car de telles ambitions ne réussissent qu'au mépris du bien public. }%
%\begin{comment}%
\note{6}{Willems, II, p. 552. C'est nous [Taillardat \& Blanc] qui soulignons soi-disant.}%
\note{7}{Xén., \italique{Anab.}, 5, 6, 13 (ἐν ἀνδραπόδων χώρᾳ ἐσόμεθα), \italique{Cyr.}, 2, 1, 18 (ἐν μισθοφόρου χώρᾳ εἶναι), Polybe, 1, 43, 1 (τῶν ἡγεμόνων τινὲς τῶν τὰς μεγίστας χώρας ἐχόντων ἐν τοῖς μισθοφόροις), Pausanias, 10, 1, 8 (ὀ δὲ χώραν ἐν τοῖς ἄρχουσιν ἔχων τὴν μεγίστην μάντις ἦν Τελλίας ὁ Ἠλεῖος). Avec le sens de « place dans la société », χώρα paraît déjà exister chez Théognis, v. 822. }%
\note{8}{On justifiera ci-dessous la traduction « aux \italique{bonnes} places ». }%
\note{9}{Tel est le sens de l'article, cf. Kühner-Gerth, \italique{Satzlehre}, I, p. 593, n°2. C'est l'emploi intentionnel du singulier τῆς χώρας (Kühner-Gerth, I, p. 14, n°4), au lieu du pluriel τῶν χωρῶν (plus fréquent dans ce cas, cf. Kühner-Gerth, I, p. 19, Anm. 3), qui permet le jeu de mots. Pour οὕνεκα, \cf \italique{Lysistrata}, v. 491 : τοῦδ᾽ οὕνεκα « en vue de cela (= en vue des magistratures et des vols qu'elles permettent) ». }%
\note{10}{\italique{Ach.}, 595 ; pour le fait, \cf \italique{Equ.}, 157-159, 164-165, et \italique{Aves}, 798-800.}%
\note{11}{Envisageant de provoquer une révolution à Athènes, il préparait alors son retour dans sa patrie (Thuc., 8, 47, 1-2).}%
\note{12}{Pour le rôle de l'or perse à ce moment, voir Thucydide, 8, 29, 1-2 ; 8, 45, 2-6; 8, 48, 2-3 ; 8, 53, 2.}%
\note{13}{Le trésor déposé à l'Acropole dont les Athéniennes barricadées interdisent l'entrée.}%
\note{14}{Voir J. Taillardat, \italique{Les images d'Aristophane}, pp. 408-423. Pour les avantages licites, \cf \italique{Acharniens}, v. 65 sqq. (indemnité de mission officielle et vie facile des ambassadeurs athéniens auprès du Grand Roi). }}%
\end{groupe}%
%% v. 367
\begin{groupe}%
\typeVers{reiz. $+$ ia.}\chantChoeur{ἀσεβοῦσ᾽ ἀδικοῦσί τε τὴν πόλιν.\apparat{367}{ἀσεβοῦσιν ἀδικοῦσίν R : \lecon{ἀσεβοῦσ᾽ ἀδικοῦσί} Bothe}}%
%\end{comment}%
\end{groupe}%
%% v. 368-371
\begin{groupe}%
\typeVers{dochm.}\chantChoeur{ἀλλ᾽ ὦ παγκρατὲς}%
\end{groupe}%
%% v. 369
\begin{groupe}%
%
\typeVers{$2$ iamb.}\chantChoeur{Ζεῦ ταῦτα κυρώσειας, ὥσθ᾽%\testimonia{369-71}{(ὥσθ᾽ -- οὔσαις) Suidas π 444 }
}%
\end{groupe}%
%% v. 370
\begin{groupe}%
%
\typeVers{$2$ iamb.}\chantChoeur{ἡμῖν θεοὺς παραστατεῖν}%
\end{groupe}%
%% v. 371
\begin{groupe}%
\typeVers{$2$ iamb\cat}\chantChoeur{καίπερ γυναιξὶν οὔσαις.\\\passeLigne}%
\end{groupe}%
%% v. 372-379
\begin{groupe}%
%
\personnageGrec{Κορυφαῖος}%
\typeVers{trim. iamb.}\trimIamb{\ligne{372}ἄκουε πᾶσ᾽. ἔδοξε τῇ βουλῇ τάδε\apparat{372}{%(Κηρ.) R : \lecon{(Κορ.)} van Leeuwen : (Κρίτ.) Sommerstein \espApp 
πᾶς R : \lecon{πᾶσ᾽} Voss \ligneApp{374} Ἀρ[. . . .]\douteux{λ}(vel \douteux{a}, \douteux{δ}, \douteux{μ})ειa \PsiCinq, Ἀρ[χίκλ]ει᾽ suppl. Hunt : \lecon{Τιμόκλει᾽} R, Σ, Suid. ε 2616 \ligneApp{376} σχολῇ R : \lecon{σχολή} Zanetti}\\
τῇ τῶν γυναικῶν· Τιμόκλει᾽ ἐπεστάτει,%\testimonia{373-4}{Satyrus \italique{Vit. Euripid.} ($=$ \italique{The Oxyrhynchus Papyri} IX 1176) fr. 39 col. XII 1 (vid. ad 335-7 ; sed Ἀρ[χίκλ]ει᾽ pro Τιμόκλει᾽, suppl. Hunt) \espTest (Τιμόκλει᾽ -- ἐγραμμάτευεν) Suidas ε 2616 (sed ἐπιστάτει pro ἐπεστάτει et ἐγραμμάτευε pro ἐγραμμάτευεν)}%
\commentaire{373}{Ἀρ[χίκλ]ει᾽ / Τιμόκλει᾽.}{Comme aux vers 336-337, le \italique{papyrus Oxyrhynchi} n° 1176\marqueur{1} diffère de la leçon du \italique{Ravennas} pour le nom de notre \italique{épistate}\marqueur{2} ; Hunt a suppléé aux quatre ou cinq lettres manquantes Ἀρ[. . . . .]ειa en proposant : Ἀρ[χίκλ]ει᾽. Pourquoi pas ? d'autres noms étaient cependant possibles.
Quoique le papyrus soit certainement plus fiable que notre \italique{Ravennas}, en l'absence de certitude sur la leçon originale du papyrus, je me bornerai à suivre celle de l'unicodex.}%
\debutNote{373}%
\note{1}{\italique{The Oxyrhynchus Papyri}, IX (1912), ed. Arthur S. Hunt. Il s'agit du papyrus n°1176, fragment 39, colonne XII 8-16 ; papyrus du \nieme{II} siècle ap. J.-C. comportant un fragment d'une \italique{Vie} d'Euripide par Satyros de Callatis. Voir aussi l'édition Schorn de Satyros, \italique{Satyros aus Kallatis} (2002).}%
\note{2}{On sait que c'est au niveau des noms propres que les textes s'altèrent le plus lors des copies successives, surtout si, comme probablement ici, il s'agit d'un personnage inconnu ailleurs.}\\
Λύσιλλ᾽ ἐγραμμάτευεν, εἶπε Ζωστράτη·\\
ἐκκλησίαν ποιεῖν ἕωθεν τῇ μέσῃ\\
τῶν Θεσμοφορίων, ᾗ μάλισθ᾽ ἡμῖν σχολή,%\apparat{376}{σχολῇ R : \lecon{σχολή} Zanetti}
\\
\commentaire{376-7}{ἐκκλησίαν ποιεῖν ἕωθεν τῇ μέσῃ / τῶν Θεσμοφορίων, ᾗ μάλισθ᾽ ἡμῖν σχολή.}{Remarquons au passage que les assemblées se tenaient à l'aube\marqueur{1}.< Le fait que le \italique{jour du milieu} soit celui du \italique{jeûne} (νηστεία) explique que les femmes aient tout le loisir pour tenir assemblée.}%
\debutNote{376-7}%
\note{1}{Coutume confirmée par d'autres passages d'Aristophane : \italique{Acharniens}, 19-20 ; \italique{Assemblée des femmes}, 20-21, 84-85, 739-741.}%
καὶ χρηματίζειν πρῶτα περὶ Εὐριπίδου,%\testimonia{377-8}{(καὶ χρηματίζειν -- παθεῖν) Suidas χ 481 ; Lexicon Vindobonense 190.10}
\\
ὅ τι χρὴ παθεῖν ἐκεῖνον· ἀδικεῖν γὰρ δοκεῖ\\
ἡμῖν ἁπάσαις. τίς ἀγορεύειν βούλεται;
\commentaire{373-9}{ἄκουε πᾶς ... τίς ἀγορεύειν βούλεται;}{L'assemblée à huis-clos féminin se donne les formes pour être en droit de juger Euripide : tout ce passage est un pastiche d'un \grec{προβούλευμα}, une \italique{résolution du Conseil}. En voici un exemple : le début du décret proposé par Démophante évoqué plus haut (Andocide, \italique{Discours sur les mystères}, 96) :
\citComm{ΝΟΜΟΣ\\
Ἔδοξε τῇ βουλῇ καὶ τῷ δήμῳ. Αἰαντὶς ἐπρυτάνευε, Κλειγένης ἐγραμμάτευε, Βοηθὸς ἐπεστάτει. Τάδε Δημόφαντος συνέγραψεν. Ἄρχει χρόνος τοῦδε τοῦ ψηφίσματος ἡ βουλὴ οἱ πεντακόσιοι λαχόντες τῷ κυάμῳ, οἷς Κλειγένης πρῶτος ἐγραμμάτευεν.\marqueur{1}}%
}%
\debutNote{373-9}%
\note{1}{Autres exemples : Thucydide, IV, 118 ; \italique{Inscriptiones Graecae} I\expo{3} 5, 9, 10, 12, 17 \&c. \&c.}\\\passeLigne}%
\end{groupe}%
\endnumbering
\end{Leftside}
\begin{Rightside}
\beginnumbering
%% v. 337+ (part 2)
\begin{groupe}
\trimIambFr{\ligne{337}à Euripide et aux Mèdes au détriment\\
des femmes, vise à être tyran
ou à ramener le tyran,\\
dénonce une femme qui a supposé un enfant; si une esclave
ayant servi d'entremetteuse l'est allé chuchoter au mari,\testimonia{341}{Suidas ε 1271 \espTest (προαγωγὸς) Hesychius π 3308 προαγωγός \espTest (ἐνετρύλλισεν) Zonaras 745 οὔ σ' ἐνετρύλλησεν}
ou si chargée d'un message, elle fait un rapport mensonger ;}
\end{groupe}%
%% v. 331-351
\begin{groupe}%
\trimIambFr{si un galant suborne par de belles paroles,\\
et ne tient pas les promesses qu'il a faites ;\\
si une vieille fait des présents à un galant,\\
ou en reçoit en livrant à une courtisane son galant ;\\
si un débitant ou une débitante\testimonia{347-8}{Suidas χ 394 (sed εἴ pro κεἴ)}\\
falsifie la mesure légale des chous ou des cotyles,\\
souhaitez qu'il périsse misérablement avec sa maison ;\\
mais à vous autres toutes priez les dieux\\
d'accorder des biens en abondance.\\\passeLigne}%
\end{groupe}%
%% v. 352-354
\begin{groupe}%
\personnageFr{Chœur}%
\chantChoeurFr{Nous unissons nos prières aux tiennes}%
\end{groupe}%
%% v. 352-354
\begin{groupe}%
\chantChoeurFr{pour le bien de la cité, pour le bien du peuple}%
\end{groupe}%
%% v. 352-354
\begin{groupe}%
\chantChoeurFr{afin que ces vœux s'accomplissent,}%
\end{groupe}%
%% v. 355-357
\begin{groupe}%
\chantChoeurFr{et que la victoire demeure, comme de raison,}%
\end{groupe}%
%% v. 352-354
\begin{groupe}%
\chantChoeurFr{à celles qui auront le mieux parlé. Mais }%
\end{groupe}%
%% v. 352-354
\begin{groupe}%
\chantChoeurFr{toutes celles qui en imposent ou transgressent}%
\end{groupe}%
%% v. 358-361
\begin{groupe}%
\chantChoeurFr{les serments établis,}%
\end{groupe}%
%% v. 352-354
\begin{groupe}%
\chantChoeurFr{dans leur propre intérêt, au détriment commun,}%
\end{groupe}%
%% v. 352-354
\begin{groupe}%
\chantChoeurFr{celles qui cherchent à bouleverser}%
\end{groupe}%
%% v. 362-364
\begin{groupe}%
\chantChoeurFr{les décrets et la loi,}%
\end{groupe}%
%% v. 352-354
\begin{groupe}%
\chantChoeurFr{révèlent nos secrets à nos ennemis,\testimonia{363-4}{(τἀπόρρητά -- λέγουσ᾽) Suidas α 3501 ἀπόρρητά τε τοῖσιν ἐχθροῖσι τοῖς ἡμετέροις λέγουσιν}}%
\end{groupe}%
%% v. 352-354
\begin{groupe}%
\chantChoeurFr{ou se positionnent }%
\end{groupe}%
%% v. 365-367
\begin{groupe}%
\chantChoeurFr{vis-à-vis des Mèdes pour le bien de notre empire...}%
\end{groupe}%
%% v. 352-354
\begin{groupe}%
\chantChoeurFr{et empirent notre position. Celles-là sont }%
\end{groupe}%
%% v. 352-354
\begin{groupe}%
\chantChoeurFr{coupables de lèse-majesté divine et de lèse-nation.}%
\end{groupe}%
%% v. 368-371
\begin{groupe}%
\chantChoeurFr{Ô tout-puissant}%
\end{groupe}%
%% v. 352-354
\begin{groupe}%
%
\chantChoeurFr{Zeus, ratifie ces vœux et fais\testimonia{369-71}{(ὥσθ᾽ -- οὔσαις) Suidas π 444 }}%
\end{groupe}%
%% v. 352-354
\begin{groupe}%
\chantChoeurFr{que les dieux nous assistent,}%
\end{groupe}%
%% v. 352-354
\begin{groupe}%
\chantChoeurFr{quoique nous soyons femmes.\\\passeLigne}%
\end{groupe}%
%% v. 372-379
\begin{groupe}%
%
\personnageFr{Coryphée}%
\trimIambFr{\ligne{373}Écoutez toutes. La Boulé
des femmes, Timocléa épistate,\testimonia{373-4}{Satyrus \italique{Vit. Euripid.} ($=$ \italique{The Oxyrhynchus Papyri} IX 1176) fr. 39 col. XII 1 (vid. ad 335-7 ; sed Ἀρ[χίκλ]ει᾽ pro Τιμόκλει᾽, suppl. Hunt) \espTest (Τιμόκλει᾽ -- ἐγραμμάτευεν) Suidas ε 2616 (sed ἐπιστάτει pro ἐπεστάτει et ἐγραμμάτευε pro ἐγραμμάτευεν)}
Lysilla greffière, Sostraté auteur de la motion, a décrété ce qui suit :\\
Une assemblée aura lieu le matin du jour intermédiaire
des Thesmophories, celui où nous sommes le plus de loisir ;
et on délibérera d'abord sur Euripide\testimonia{377-8}{(καὶ χρηματίζειν -- παθεῖν) Suidas χ 481 ; Lexicon Vindobonense 190.10}\\
et la peine qu'il doit subir. Car il est tenu coupable\\
par nous toutes. Qui veut prendre la parole ?\\\passeLigne}%
\end{groupe}%
\endnumbering
\end{Rightside}
\end{pages}
\Pages
\clearpage\ifodd\value{page}\hbox{}\newpage\fi
\enTeteG{ΕΠΕΙΣΟΔΙΟΝ Α' --- Codex Ravennas 429, f. 167r (Classense 169r) --- vv. 380–422}
\begin{adjustwidth}{-1.5cm}{-3.0cm}
\begin{adjustbox}{max totalheight=21.0cm, max width=16.5cm, keepaspectratio}
\begin{minipage}{16.5cm}
{\fontsize{7.2pt}{8.5pt}\selectfont
\noindent\textbf{Folio 10 --- Codex Ravennas 429, \texttt{f. 167r} (\texttt{Classense 169r}) : ΕΠΕΙΣΟΔΙΟΝ Α', vv. 380–422}\hfill \textit{Témoins :} $\text{R}^1, \Rpc, \Rdeux$ (\Mon)\par\vspace{2pt}
\hrule height 0.4pt\vspace{2.5pt}
}
{\fontsize{8.1pt}{9.5pt}\selectfont
\ligneRav{380}{\auDessus{γυ}{ν}}{ἐγώ. : περίθου νυν τόνδε πρῶτον πρὶν λέγειν·}{\textbf{380} (Γυνή τις) R : \lecon{(Μίκα)} Rogers, cf. v. 760\espApp \lecon{πρῶτον} R : πρότερον Suid. π 1166}
\ligneRav{}{}{σίγα σιώπα, πρόσεχε τὸν νοῦν· χρέμπτεται γὰρ ἤδη}{}
\ligneRav{}{}{ὅπερ ποιοῦσ᾽ οἱ ῥήτορες. μακρὰν ἔοικε λέξειν·}{}
\ligneRav{383}{καλλιλεξία}{φιλοτιμίᾳ μὲν οὐδεμιᾷ μὰ τὼ θεὼ}{\textbf{383} (\auDessus{γυ}{ν}) \Rac : (καλλιλεξία γυνή) \Rpc, \Mon : (\lecon{Μίκα}) Rogers}
\ligneRav{}{}{λέξουσ᾽ ἀνέστην ὦ γυναῖκες· ἀλλὰ γὰρ}{}
\ligneRav{385}{}{βαρέως φέρω τάλαινα πολὺν ἤδη χρόνον}{}
\ligneRav{386}{}{προπηλακιζομένας ὁρῶς \difference{ἡμᾶς} ὑπὸ}{\textbf{386} ὑμᾶς R : \lecon{ἡμᾶς} Cobet}
\ligneRav{}{}{Εὐριπίδου τοῦ τῆς λαχανοπωλητρίας}{}
\ligneRav{}{}{καὶ πολλὰ καὶ παντοῖ᾽ ἀκουούσας κακά·}{}
\ligneRav{389}{}{τί γὰρ οὗτος ἡμᾶς οὐκ ἐπισμῇ τῶν κακῶν;}{\textbf{389} \lecon{ἡμᾶς} R, Suid. ε 2592 : ἡμῖν Brunck}
\ligneRav{390}{}{ποῦ δ᾽ οὐχὶ διαβέβληχ᾽, ὅπουπερ \difference{ἔμβραχυ}}{\textbf{390} ἐν βραχὺ R : \lecon{ἔμβραχυ} Suid., Σ \italique{Gorg}. et \italique{Hipp.m}.}
\ligneRav{391}{}{εἰσὶν θεαταὶ καὶ τραγῳδοὶ καὶ χοροί,}{\textbf{391} \lecon{τραγῳδοὶ καὶ} R : τραγῳδικοὶ Schol. Plat. \italique{Theag} : τραγῳδία καὶ Suid. οι 128}
\ligneRav{392}{}{τὰς μοιχοτρόπους, τὰς ἀνδρεραστρίας καλῶν,}{\textbf{392} μυ- R : \lecon{μοι-} Suid. \espApp ἀνδρεα- R : \lecon{ἀνδρερα-} Suid.}
\ligneRav{393}{}{τὰς οἰνοπίπας, τὰς προδότιδας, τὰς λάλους,}{\textbf{393} οἰνοπότιδας R : \lecon{-πίπας} Suid. : -πίπους Σ}
\ligneRav{394}{}{τὰς οὐδὲν ὑγιές, τὰς μέγ᾽ ἀνδράσιν κακόν·}{\textbf{394} ἀνδράσι R : \lecon{-σιν} Daubuz}
\ligneRav{395}{}{ὥστ᾽ εὐθὺς εἰσιόντες ἀπὸ τῶν ἰκρίων}{}
\ligneRav{}{}{ὑποβλέπουσ᾽ ἡμᾶς σκοποῦνταί τ᾽ εὐθέως}{}
\ligneRav{}{}{μὴ μοιχὸς ἔνδον ᾖ τις ἀποκεκρυμμένος·}{}
\ligneRav{398}{}{δρᾶσαι δ᾽ ἔθ᾽ ἡμῖν οὐδὲν \difference{ὧνπερ} καὶ πρὸ τοῦ}{\textbf{398} ὥσπερ R : \lecon{ὧνπερ} Küster}
\ligneRav{}{}{ἔξεστι· τοιαῦθ᾽ οὗτος ἐδίδαξεν κακὰ}{}
\ligneRav{400}{}{τοὺς ἄνδρας ἡμῶν· ὥστ᾽ ἐάν \difference{τις ξυμπλέκῃ}}{\textbf{400} τις πλέκῃ R : \lecon{τις ξυμπλέκῃ} Willems : δή τις πλέκῃ Austin : alii alia}
\ligneRav{}{}{γυνὴ στέφανον, ἐρᾶν δοκεῖ· κἂν ἐκβάλῃ}{}
\ligneRav{}{}{σκεῦός τι κατὰ τὴν οἰκίαν πλανωμένη,}{}
\ligneRav{403}{}{\difference{ἁνὴρ} ἐρωτᾷ· `τῷ κατέαγεν ἡ χύτρα;}{\textbf{403} ἀνὴρ R : \lecon{ἁνὴρ} Brunck}
\ligneRav{}{}{οὐκ ἔσθ᾽ ὅπως οὐ τῷ Κορινθίῳ ξένῳ.'}{}
\ligneRav{405}{}{κάμνει κόρη τις, εὐθὺς \difference{ἁδελφὸς} λέγει·}{\textbf{405} ἀδελφὸς R : \lecon{ἁδελφὸς} Biset}
\ligneRav{}{}{`τὸ χρῶμα τοῦτό μ᾽ οὐκ ἀρέσκει τῆς κόρης.'}{}
\ligneRav{}{}{εἶἑν, γυνή τις ὑποβαλέσθαι βούλεται}{}
\ligneRav{}{}{ἀποροῦσα παίδων, οὐδὲ τοῦτ᾽ ἔστιν λαθεῖν·}{}
\ligneRav{409}{}{\difference{ἅνδρες} γὰρ ἤδη παρακάθηνται πλησίον·}{\textbf{409} ἄνδρες R : \lecon{ἅνδρες} Brunck.}
\ligneRav{410}{}{πρὸς τοὺς γέροντάς θ᾽ οἳ πρὸ τοῦ τὰς μείρακας}{}
\ligneRav{}{}{ἤγοντο, διαβέβληκεν, ὥστ᾽ οὐδεὶς γέρων}{}
\ligneRav{412}{}{γαμεῖν \difference{ἐθέλει} γυναῖκα διὰ τοὔπος τοδὶ}{\textbf{412} θέλει R : \lecon{ἐθέλει} Kappeyne van de Coppello \italique{Obs. crit.} I.2}
\ligneRav{}{}{`δέσποινα γὰρ γέροντι νυμφίῳ γυνή'·}{}
\ligneRav{}{}{εἶτα διὰ τοῦτον ταῖς γυναικωνίτισιν}{}
\ligneRav{415}{}{σφραγῖδας ἐπιβάλλουσιν ἤδη καὶ μοχλοὺς}{}
\ligneRav{}{}{τηροῦντες ἡμᾶς, καὶ προσέτι Μολοττικοὺς}{}
\ligneRav{417}{}{τρέφουσι μορμολυκεῖα τοῖς \difference{μοιχοῖς} κύνας·}{\textbf{417} μυχοῖς R : \lecon{μοιχοῖς} \Mon}
\ligneRav{}{}{καὶ ταῦτα μὲν ξυγγνώσθ᾽. ἃ δ᾽ ἦν ἡμῖν πρὸ τοῦ}{}
\ligneRav{419}{}{αὐταῖς ταμιεῦσαι καὶ προαιρούσαις λαβεῖν}{\textbf{419} ταμιεύεσθαι R : \lecon{ταμιεῦσαι καὶ} Reiske}
\ligneRav{420}{}{ἄλφιτον ἔλαιον οἶνον, οὐδὲ ταῦτ᾽ ἔτι}{}
\ligneRav{}{}{ἔξεστιν. οἱ γὰρ ἄνδρες ἤδη κλῄδια}{}
\ligneRav{}{}{αὐτοὶ φοροῦσι κρυπτὰ κακοηθέστατα}{}
}
\vspace{3pt}\noindent\rule{4.5cm}{0.4pt}\par\vspace{1.5pt}
{\fontsize{7.0pt}{8.3pt}\selectfont
\noindent\textbf{Loci paralleli (vv. 380–422) ---} \textbf{380} Suid. π 1166 περίθου νῦν τόνδε πρότερον \espTest \textbf{381} (χρέμπτεται) \Phryn 126.1 \espTest \textbf{383-4} (φιλοτιμίᾳ -- γυναῖκες) Suidas φ 431 \espTest \textbf{389} Suidas ε 2592 \espTest (ἐπισμῇ) Hesychius ε 5195 ἐπισμῇ \espTest \textbf{390-4} Suidas οι 128 (οἷς θέα pro εἰσὶν θεαταὶ et τραγῳδία pro τραγῳδοὶ) \espTest \textbf{390-1} scholium Plato \italique{Theages} 127c (πῶς pro ποῦ et τραγῳδικοὶ pro τραγῳδοὶ καὶ) $=$ scholium Plato \italique{Gorgias} 457b $=$ scholium Plato \italique{Hippias min.} 365d \espTest (ποῦ -- θεαταὶ) Suidas ε 955 (sed εἰσὶ pro εἰσὶν) \espTest \textbf{390-1} scholium Plato \italique{Theages} 127c $=$ scholium Plato \italique{Gorgias} 457b $=$ scholium Plato \italique{Hippias min.} 365d \espTest (ποῦ -- θεαταὶ) Suidas ε 955 (sed εἰσὶ pro εἰσὶν) \espTest \textbf{392} (ἀνδρεραστρίας) Pollux III 71 ; \Phryn 34.18 ; \Eustathe 380.15 \espTest \textbf{393} (οἰνοπίπας) scholium \italique{Aves} 307 \espTest \textbf{395-6} (ὥστ᾽ εὐθὺς -- ἡμᾶς) Suidas ι 275 \espTest \textbf{395} (ἰκρίων) Hesychius ι 501 ἴκρια \espTest \textbf{404} (τῷ Κορινθίῳ ξένῳ) Hesychius κ 3629 Κορίνθιος ξένος \espTest \textbf{413} Menandros \italique{Monostich.} 191 ($=$ Euripides fr. 804.3 ed. Kannicht) ; \Apostol V 93b \espTest \textbf{417} (μορμολυκεῖα) \EtymMagnum 590.53 \espTest \textbf{421-8} (οἱ γὰρ ἄνδρες -- ἐξαψαμένους) Suidas λ 54 (αὐτοὶ, τρεῖς et ᾡκότριψ Εὐριπίδης omissis, Λακωνίκα pro Λακωνίκ᾽, θριπηδέστατ᾽ pro θριπήδεστ᾽ et σφραγῖδα pro σφραγίδια) \espTest \textbf{421-7} (οἱ γὰρ ἄνδρες -- θριπήδεστ᾽) Suidas κ 159 (αὐτοὶ et ᾡκότριψ Εὐριπίδης omissis, Λακωνίκα pro Λακωνίκ᾽ et θριπηδέστατ᾽ pro θριπήδεστ᾽) \espTest \textbf{421-7} (οἱ γὰρ ἄνδρες -- θριπήδεστ᾽) \espTest \textbf{422} (κακοηθέστατα) \Eustathe 974.37, 1108.25, 1603.52\par
}
\vspace{2.5pt}\noindent\rule{4.5cm}{0.4pt}\par\vspace{1.5pt}
{\fontsize{6.9pt}{8.2pt}\selectfont
\noindent\textbf{Suite (en commentaire) de l'apparat critique \& notes paléographiques (\texttt{f. 167r}) ---} \textbf{Note de personnage descriptive au v. 383 :} Là où \Rac{} avait simplement abrégé (\auDessus{γυ}{ν}), \Rpc{} (suivi par \Mon{}) a développé en marge \lecon{(καλλιλεξία γυνή)} (« la femme au beau langage », glose descriptive passée dans la rubrique de locuteur ; Rogers l'identifie à \textbf{Mika} d'après le v. 760). --- \textbf{Iotacismes $\text{οι} \to \upsilon$ de $\text{R}^1$ :} μυχοτρόπους (v. 392) et μυχοῖς (v. 417) pour μοιχο-, redressés par la Souda et \Mon{}. --- \textbf{Restitution par \textbf{Bibi} du fragment de Sappho (fr. 70/125) dans $\Sigma^{\text{R}}$ ad v. 401 :} La scholie du \italique{Ravennas} au v. 401 conserve un précieux vers de Sappho sous la forme corrompue \lecon{αυταόρα αἱστεφανηπλόκουν} ; \textbf{Louis Hauseux (\textbf{Bibi}, \texttt{vv 1-641.tex}, ll. 3272–3294) restitue l'éolien αὐτὰ κόραις στεφαναπλόκεο} (ou -πλόκην), montrant que αυταόρα αἱ procède de la chute du $\kappa$ et d'une médivision de αὐτὰ κόραις. --- \textbf{v. 390 :} ἐν βραχὺ R : \lecon{ἔμβραχυ} Souda, $\Sigma$ ad Plat. \italique{Gorg.} et \italique{Hipp. min.} --- \textbf{v. 392 :} μυ- R : \lecon{μοι-} (μοιχοτρόπους) Souda \espApp ἀνδρεα- R (haplographie du second $\rho$) : \lecon{ἀνδρερα-} (ἀνδρεραστρίας) Souda. --- \textbf{v. 393 :} οἰνοπότιδας R (substitution d'un mot banal au composé comique rare) : \lecon{οἰνοπίπας} Souda (-πίπους $\Sigma$). --- \textbf{v. 400 :} τις πλέκῃ R : \lecon{τις ξυμπλέκῃ} Alphonse Willems (adopté par L. Hauseux ; δή τις πλέκῃ Austin). --- \textbf{$\Sigma^{\text{R}}$ ad v. 401 (Sappho fr. 70/125) :} αυταόρα αἱστεφανηπλόκουν $\Sigma^{\text{R}}$ : \lecon{αὐτὰ κόραις στεφαναπλόκεο} \lecon{Bibi}.\par
}
\end{minipage}
\end{adjustbox}
\end{adjustwidth}
\newpage
\enTeteD{Codex Ravennas 429 (Biblioteca Classense) --- Fac-similé du f. 167r (vv. 380–422)}
\begin{adjustwidth}{-1.5cm}{-3.0cm}
\begin{center}
\includegraphics[width=16.5cm,height=20.5cm,keepaspectratio]{/workspaces/Thesmophories/dossier_ravennas_thesmophories/folios_individuels_ff162v-177r/10_f167r_classense169r_vv0380-0422.jpg}\par
\vspace{2pt}
{\scriptsize\textbf{Codex Ravennas 429, f. 167r (Classense 169r)} --- ΕΠΕΙΣΟΔΙΟΝ Α', \textbf{vv. 380–422} ($\text{R}^1, \Rpc, \Rdeux$ (\Mon)).}
\end{center}
\end{adjustwidth}
\newpage
\clearpage\ifodd\value{page}\hbox{}\newpage\fi
\enTeteG{ΘΕΣΜΟΦΟΡΙΑΖΟΥΣΑΙ --- Texte grec établi (f. 167r, vv. 380–422)}
\enTeteD{LES THESMOPHORIES --- Traduction française (f. 167r, vv. 380–422)}
\begin{pages}
\begin{Leftside}
\beginnumbering
%% v. 380a
\begin{groupe}%
\personnageGrec{Μίκα}%
\trimIamb{ἐγώ.\apparat{380}{(Γυνή τις) R : \lecon{(Μίκα)} Rogers, cf. v. 760}\\\passeLigne}%
\end{groupe}%
%% v. 380b-382
\begin{groupe}%
\personnageGrec{Κορυφαῖος}%
\trimIamb{\ligne{380}\decalage{ἐγώ.}περίθου νυν τόνδε πρῶτον πρὶν λέγειν.%\testimonia{380}{Suid. π 1166 περίθου νῦν τόνδε πρότερον}
\apparat{380}{\lecon{πρῶτον} R : πρότερον Suid. π 1166}}%
\end{groupe}%
\begin{groupe}%
\typeVers{tetr. iamb\cat}\tetramIamb{σίγα σιώπα, πρόσεχε τὸν νοῦν· χρέμπτεται γὰρ ἤδη%\testimonia{381}{(χρέμπτεται) \Phryn 126.1}%
\commentaire{381}{χρέμπτεται.}{Elle \italique{tousse}, et non elle crache, comme traduisent les dictionnaires. L'ancien français avait un terme excellent, qui s'est perdu avec tant d'autres : « Il se lève, compose sa robe, sa barbe et ses yeux à la modestie, et ayant craché et toussi clair, il commença ». D'Aubigné, \italique{Avent. du baron de Foeneste}, I. IV, ch. 10. (Note de Willems.)}
%
\\
%
ὅπερ ποιοῦσ᾽ οἱ ῥήτορες. μακρὰν ἔοικε λέξειν.\\\passeLigne}%
\end{groupe}%
%% v. 383-432
\begin{groupe}%
\personnageGrec{Μίκα}%
\typeVers{trim. iamb.}\trimIamb{φιλοτιμίᾳ μὲν οὐδεμιᾷ μὰ τὼ θεὼ%\testimonia{383-4}{(φιλοτιμίᾳ -- γυναῖκες) Suidas φ 431}
\apparat{383}{(\auDessus{γυ}{ν}) \Rac : (καλλιλεξία γυνή) \Rpc, \Mon : (\lecon{Μίκα}) Rogers}\\
λέξουσ᾽ ἀνέστην ὦ γυναῖκες· ἀλλὰ γὰρ\\
βαρέως φέρω τάλαινα πολὺν ἤδη χρόνον\\
προπηλακιζομένας ὁρῶς ἡμᾶς ὑπὸ\apparat{386}{ὑμᾶς R : \lecon{ἡμᾶς} Cobet}\\
Εὐριπίδου τοῦ τῆς λαχανοπωλητρίας\\
καὶ πολλὰ καὶ παντοῖ᾽ ἀκουούσας κακά.\\
τί γὰρ οὗτος ἡμᾶς οὐκ ἐπισμῇ τῶν κακῶν;%\testimonia{389}{Suidas ε 2592 \espTest (ἐπισμῇ) Hesychius ε 5195 ἐπισμῇ}
%\apparat{389}{\lecon{ἡμᾶς} R, Suid. ε 2592 : ἡμῖν Brunck}
\\
ποῦ δ᾽ οὐχὶ διαβέβληχ᾽, ὅπουπερ ἔμβραχυ%\testimonia{390-4}{Suidas οι 128 (οἷς θέα pro εἰσὶν θεαταὶ et τραγῳδία pro τραγῳδοὶ)}\testimonia{390-1}{scholium Plato \italique{Theages} 127c (πῶς pro ποῦ et τραγῳδικοὶ pro τραγῳδοὶ καὶ) $=$ scholium Plato \italique{Gorgias} 457b $=$ scholium Plato \italique{Hippias min.} 365d \espTest (ποῦ -- θεαταὶ) Suidas ε 955 (sed εἰσὶ pro εἰσὶν)}
\apparat{390}{%\lecon{ποῦ} R, Suid. : πῶς Σ \italique{Theag}. \espApp 
ἐν βραχὺ R %, Σ \italique{Theag}.
: \lecon{ἔμβραχυ} Suid., Σ \italique{Gorg}. et \italique{Hipp.m}. %\ligneApp{391} \lecon{τραγῳδοὶ καὶ} R : τραγῳδικοὶ Σ Plat.\italique{Theag} : τραγῳδία καὶ Suid. οι 128 
\ligneApp{392} μυ- R : \lecon{μοι-} Suid. \espApp ἀνδρεα- R%, Suid. AV 
 : \lecon{ἀνδρερα-} Suid. %GFM 
\ligneApp{393} οἰνοπότιδας R : \lecon{-πίπας} Suid. : -πίπους Σ \ligneApp{394} ἀνδράσι R : \lecon{-σιν} Daubuz}\\
εἰσὶν θεαταὶ καὶ τραγῳδοὶ καὶ χοροί,\\
τὰς μοιχοτρόπους, τὰς ἀνδρεραστρίας καλῶν,%\testimonia{392}{(ἀνδρεραστρίας) Pollux III 71 ; \Phryn 34.18 ; \Eustathe 380.15}
\\
τὰς οἰνοπίπας, τὰς προδότιδας, τὰς λάλους,%\testimonia{393}{(οἰνοπίπας) scholium \italique{Aves} 307}%
\commentaire{393}{τὰς οἰνοπίπας,}{Encore une fois, c'est le scholiaste, disposant de meilleurs manuscrits, qui nous permet de revenir à la leçon originale : après avoir évoqué Symmaque, il rapproche non sans raison \grec{οἰνοπίπης} de \grec{παιδοπίπης} (on aurait pu rajouter : \grec{γυναικοπίπης}, \grec{παρθενοπίπης} ou encore le \grec{πυροπίπης} / πυρροπίπης des \italique{Cavaliers}, v. 407) ; c'est-à-dire que le suffixe ne renverrait pas aux verbes \grec{πιπίζειν} ou \grec{πιπίσκειν} mais -\grec{οπίπης} serait de la famille de \grec{ὀπιπεύειν}, \italique{guigner}, \italique{guetter}, \italique{épier} (cf. aussi Hésychius (ο 1006) : \grec{ὀπῖπα· ἐξαπατᾷ, ἀπατεών, ἢ ἀπατῶν}). 
Aristophane ne va pas manquer de surenchérir : la difficile relation des femmes au vin sera allègrement mise en scène dans la suite de la pièce.}}
\end{groupe}%
%% v. 383-432
\begin{groupe}%
\trimIamb{τὰς οὐδὲν ὑγιές, τὰς μέγ᾽ ἀνδράσιν κακόν·\commentaire{394}{τὰς οὐδὲν ὑγιές.}{Érasme a compilé l'expression dans ses \italique{Adages} (n° 738, \latin{nihil sanum}).}\\
ὥστ᾽ εὐθὺς εἰσιόντες ἀπὸ τῶν ἰκρίων%\testimonia{395-6}{(ὥστ᾽ εὐθὺς -- ἡμᾶς) Suidas ι 275}\testimonia{395}{(ἰκρίων) Hesychius ι 501 ἴκρια}
\\
ὑποβλέπουσ᾽ ἡμᾶς σκοποῦνταί τ᾽ εὐθέως\\
μὴ μοιχὸς ἔνδον ᾖ τις ἀποκεκρυμμένος.\\
δρᾶσαι δ᾽ ἔθ᾽ ἡμῖν οὐδὲν ὧνπερ καὶ πρὸ τοῦ\apparat{398}{ὥσπερ R : \lecon{ὧνπερ} Küster}\\
ἔξεστι· τοιαῦθ᾽ οὗτος ἐδίδαξεν κακὰ}
\end{groupe}%
%% v. 383-432
\begin{groupe}%
\trimIamb{τοὺς ἄνδρας ἡμῶν· ὥστ᾽ ἐάν τις ξυμπλέκῃ\apparat{400}{τις πλέκῃ R : \lecon{τις ξυμπλέκῃ} Willems : δή τις πλέκῃ Austin : alii alia}%
\commentaire{400}{τοὺς ἄνδρας ἡμῶν· ὥστ᾽ ἐάν τις πλέκῃ.}{Le vers 400 pèche contre le mètre. Willems y supplée en lisant <ξυμ>πλέκῃ ; il s'en justifie dans la \nieme{11} note, reproduite à la fin.}%
\commentaire{400-1}{ὥστ᾽ ἐάν τις πλέκῃ / γυνὴ στέφανον, ἐρᾶν δοκεῖ·}{Le scholiaste s'est montré court sur ce travers qu'Euripide impute aux femmes tressandières : l'on ne sait de laquelle de ses pièces (perdues) le reproche est tiré.
Toutefois, tresser une couronne pour son amant(e) est chose bien attestée ; et le scholiaste de citer un vers de Sappho\marqueur{1} :
\citComm{νεωτέρων καὶ ἐρωτικῶν τὸ στεφανοπλοκεῖν. πρὸς τὸ ἔθος, ὅτι ἐστεφανηπλόκουν αἱ παλαιαί.\\
Σαπφώ· αὐτὰ κόραις στεφαναπλόκεο\\
Sappho : «~toi qui te faisais tresser des couronnes par des jeunes filles~».}
Athénée (XV, 670d) confirme la coutume d'offrir des couronnes à son être aimé :
\citComm{στεφανοῦσιν δὲ τὰς τῶν ἐρωμένων θύρας ἤτοι τιμῆς χάριν καθαπερεί τινος θεοῦ τὰ πρόθυρα [στεφανοῦσιν αὐτῶν], ἢ οὐ τοῖς ἐρωμένοις ἀλλὰ τῷ Ἔρωτι ποιούμενοι τὴν τῶν στεφάνων ἀνάθεσιν τοῦ μὲν Ἔρωτος τὸν ἐρώμενον ἄγαλμα, τούτου δὲ ναὸν ὄντα τὴν οἴκησιν στεφανοῦσι [τὰ τῶν ἐρωμένων πρόθυρα]: διὰ ταῦτα δὲ καὶ θύουσιν ἔνιοι ἐπὶ ταῖς τῶν ἐρωμένων θύραις.\\
Quant aux couronnes et aux guirlandes dont les amants ornent les portes des personnes qu'ils aiment, ils les y mettent ou pour leur rendre des honneurs, comme couronnant le porche d'une divinité ; ou, c'est à l'amour qu'ils les offrent en les posant, et non aux personnes qui sont aimées ; car ces personnes sont pour les amants les images mêmes du dieu d'amour; et la maison où elles demeurent en étant le temple, c'en est aussi le porche qu'ils ornent de couronnes, en les mettant aux portes de ces personnes. Voilà aussi pourquoi on immole quelquefois des victimes aux portes des personnes qu'on aime.}}%
\debutNote{400-1}%
\note{1}{Citation corrompue ; le manuscrit de Ravenne donne : αυταόρα αἱστεφανηπλόκουν. La lecture proposée est de mon cru. Théodore Reinach, dubitatif, reprenait la conjecture de Hermann : αὐτὰρ ὀραῖαι στεφανηπλόκευν. Ce fragment de Sappho est le \nieme{70} de l'édition Reinach-Puech aux Belles Lettres (1937) et le \nieme{125} de l'édition d'Éva-Maria Voigt, \italique{Sappho et Alcaeus} (1971).}\\
γυνὴ στέφανον, ἐρᾶν δοκεῖ· κἂν ἐκβάλῃ\\
σκεῦός τι κατὰ τὴν οἰκίαν πλανωμένη,\\
ἁνὴρ ἐρωτᾷ· `τῷ κατέαγεν ἡ χύτρα;\apparat{403}{ἀνὴρ R : \lecon{ἁνὴρ} Brunck}\\
οὐκ ἔσθ᾽ ὅπως οὐ τῷ Κορινθίῳ ξένῳ.'%\testimonia{404}{(τῷ Κορινθίῳ ξένῳ) Hesychius κ 3629 Κορίνθιος ξένος}%
\commentaire{401-4}{κἂν ἐκβάλῃ ... τῷ Κορινθίῳ ξένῳ.}{\italique{L'étranger de Corinthe}, c'est Bellérophon, objet d'une coupable passion de la part de la princesse Sthénébée dans la pièce éponyme d'Euripide. Cette dernière, voyant ses avances repoussées par le chaste héros, l'accuse devant le roi son mari d'avoir essayé d'abuser d'elle. Le roi ordonne à Bellérophon de partir pour la Carie. Il y reçoit pour mission d'affronter la Chimère. C'est pendant le temps que Bellérophon est en Carie, et que Sthénébée le croit mort, qu'elle verse des sortes de libations à la mémoire du héros, comme l'on fait au jeu du cottabe avec les personnes aimées. Athénée (X, 427e) nous a conservé ce fragment\marqueur{1} :
\citComm{πεσὸν δέ νιν λέληθεν οὐδὲν ἐκ χερός,\\
ἀλλ᾽ εὐθὺς αὐδᾷ· `τῷ Κορινθίῳ ξένῳ'\\
<La Nourrice>\\
\francais{Rien ne tombe, en lui échappant des mains,\\
qu'elle ne s'écrie aussitôt : « Pour l'hôte de Corinthe ! »}}%
Les maris athéniens prenaient leur femme pour autant de Sthénébées.\\
Ce passage est à rapprocher d'un fragment d'Aristophane, rapporté par Diogène Laërte\marqueur{2} :
\citComm{Τὰ δὲ πεσόντα μὴ ἀναιρεῖσθαι, ὑπὲρ τοῦ ἐθίζεσθαι μὴ ἀκολάστως ἐσθίειν ἢ ὅτι ἐπὶ τελευτῇ τινος· καὶ Ἀριστοφάνης δὲ τῶν ἡρώων φησὶν εἶναι τὰ πίπτοντα, λέγων ἐν τοῖς Ἥρωσι,\\
Μηδὲ γεύεσθ' ἅττ' ἂν ἐντὸς τῆς τραπέζης καταπέσῃ.\\
$[$ Pythagore $]$ interdit de ramasser ce qui tombe de la table pour habituer à manger avec tempérance, ou bien encore parce que cela est destiné aux morts. Ce qui tombe de la table est pour les héros, suivant Aristophane ; car il dit dans les \italique{Héros} :
\francais{« Ne goûtez point à ce qui tombe de la table. »}}}%
\debutNote{401-4}%
\note{1}{fragment 3, \italique{Sthénébée}, de l'édition Jouan-van Looy, Belles Lettres, t. VIII, \nieme{3} partie $=$ fr. 664 de l'édition Kannicht $=$ fr. 883 de l'édition Mette. Traduction Jouan \& van Looy.}%
\note{2}{\italique{Vies}, VIII, 1 (Pythagore), 34. Correspond au fragment 320(305) de l'édition Kassel-Austin (\italique{Poetae Comici Graeci}) des fragments d'Aristophane.}}
\end{groupe}%
%% v. 383-432
\begin{groupe}%
\trimIamb{κάμνει κόρη τις, εὐθὺς ἁδελφὸς λέγει·\apparat{405}{ἀδελφὸς R : \lecon{ἁδελφὸς} Biset}\\
`τὸ χρῶμα τοῦτό μ᾽ οὐκ ἀρέσκει τῆς κόρης.'\\
\commentaire{405-6}{εὐθὺς ἁδελφὸς λέγει· / `τὸ χρῶμα τοῦτό μ᾽ οὐκ ἀρέσκει τῆς κόρης.'}{Il la soupçonne d'être enceinte. La pièce d'Euripide à laquelle se réfère Aristophane n'est pas identifiée ; pourtant il semble clair qu'il s'agit d'une citation (ne serait-ce que par l'absence de résolution au vers 406).}%
εἶἑν, γυνή τις ὑποβαλέσθαι βούλεται\\
ἀποροῦσα παίδων, οὐδὲ τοῦτ᾽ ἔστιν λαθεῖν.
\commentaire{406-7}{γυνή τις ὑποβαλέσθαι βούλεται / ἀποροῦσα παίδων, οὐδὲ τοῦτ᾽ ἔστιν λαθεῖν.}{Le thème de la \italique{supposition d'enfants}\marqueur{1} est récurrent dans les \italique{Thesmophories} (cf. vv. 340, 502-516, 564-565). À en croire Satyros dans sa \italique{Vita Euripidis}\marqueur{2}, c'est à Euripide qu'on doit son omniprésence dans la Comédie Nouvelle... et le biographe d'Euripide de donner du grain à moudre à ses contempteurs :
\citComm{πρὸς γ[υ]ναῖκα καὶ πατρὶ πρὸς υ[ἱὸ]ν καὶ θερά[πον]τι πρὸς δ[εσ]πότην, ἢ τ[ὰ] κατὰ τὰς π[ερι]πετείας, β[ια]σμοὺς παρθένων, \textbf{ὑποβολὰς παιδίων}, ἀναγνωρισμοὺς διά τε δακτυλίων καὶ διὰ δεραίων· ταῦτα γάρ ἐστι δήπου τὰ συνέχοντα τα τὴν νεωτέραν κωμῳδίαν, ἃ πρὸς ἄκρον ἤγα[γ]εν Εὐρπίδης, Ὁμήρου [ὄ]ντος ἀρχῆς\\
\francais{$[$le mari$]$ contre sa femme, le père contre son fils, le serviteur contre son maître ; ou lors des péripéties : violences de jeunes filles, \textbf{substitutions d'enfants}, reconnaissances par des anneaux ou des colliers ; hé bien c'est tous ces éléments qui composent la trame de la Nouvelle Comédie, éléments qu'Euripide avait conduit au sommet, Homère en étant le point de départ}}%
La tradition ne nous a rapporté qu'une seule supposition d'enfants chez Euripide : dans les \italique{Phéniciennes} (v. 28-31) :
\citComm{Πολύβου δέ νιν λαβόντες ἱπποβουκόλοι\\
φέρουσ᾽ ἐς οἴκους ἔς τε δεσποίνης χέρας\\
ἔθηκαν. ἣ δὲ τὸν ἐμὸν ὠδίνων πόνον\\
μαστοῖς ὑφεῖτο καὶ πόσιν πείθει τεκεῖν,}%
ainsi qu'une allusion dans l'\italique{Alceste} (v. 639) : Admète, désapprouvant l'attitude de son père qui a refusé de donner sa vie pour lui, suppose qu'il a été supposé.}%
\debutNote{406-7}%
\note{1}{Quand les femmes font passer pour leur des enfants – mâles le plus souvent ! – auprès de leur mari.}%
\note{2}{\italique{The Oxyrhynchus Papyri}, IX.1776, fr. 39, Col. vii, 1-24}\\
ἅνδρες γὰρ ἤδη παρακάθηνται πλησίον·\apparat{409}{ἄνδρες R : \lecon{ἅνδρες} Brunck.}\\
πρὸς τοὺς γέροντάς θ᾽ οἳ πρὸ τοῦ τὰς μείρακας\\
ἤγοντο, διαβέβληκεν, ὥστ᾽ οὐδεὶς γέρων\\
γαμεῖν ἐθέλει γυναῖκα διὰ τοὔπος τοδὶ\apparat{412}{θέλει R : \lecon{ἐθέλει} Kappeyne van de Coppello \italique{Obs. crit.} I.2}\\
`δέσποινα γὰρ γέροντι νυμφίῳ γυνή'.%\testimonia{413}{Menandros \italique{Monostich.} 191 ($=$ Euripides fr. 804.3 ed. Kannicht) ; \Apostol V 93b}%
\commentaire{413}{δέσποινα γὰρ γέροντι νυμφίῳ γυνή.}{Stobée (IV, 22a109) nous a conservé la fameuse citation d'Euripide ; elle est tirée du \italique{Phoenix}\marqueur{1} :
\citComm{μοχθηρόν ἐστιν ἀνδρὶ πρεσβύτῃ τέκνα\\
δίδωσιν†, ὅστις οὐκέθ᾽ ὡραῖος γαμεῖ·\\
δέςποινα γὰρ γέροντι νυμφίῳ γυνή\\
\francais{C'est un tracas pour un vieillard <qui est doté d'enfants> de se marier alors qu'il n'en a plus l'âge ; avec un vieil époux, une femme est la maîtresse.}}%
\italique{Quand on est jeune, on aime en fou ; Quand on est vieux, qui aime est fou.}}%
\debutNote{413}%
\note{1}{fr. 3, \italique{Phoenix} de l'édition Jouan-van Looy aux Belles-lettres, tome VIII, \nieme{3} partie. Je reprends leur traduction. Ce fragment correspond au n°804 de l'édition Kannicht et 1139 de celle de Mette.}\\
εἶτα διὰ τοῦτον ταῖς γυναικωνίτισιν\\
σφραγῖδας ἐπιβάλλουσιν ἤδη καὶ μοχλοὺς}
\end{groupe}%
%% v. 383-432 (part 1)
\begin{groupe}\trimIamb{τηροῦντες ἡμᾶς, καὶ προσέτι Μολοττικοὺς\\
τρέφουσι μορμολυκεῖα τοῖς μοιχοῖς κύνας.
\apparat{417}{
μυχοῖς R : \lecon{μοιχοῖς} \Mon}
\commentaire{416-7}{Μολοττικοὺς / τρέφουσι μορμολυκεῖα τοῖς μοιχοῖς κύνας.}{C'est le conseil et le vœu d'Hippolyte dans la pièce d'Euripide (vv. 645 et suiv. ; note de Willems) : 
\citComm{χρῆν δ᾽ ἐς γυναῖκα πρόσπολον μὲν οὐ περᾶν,\\
ἄφθογγα δ᾽ αὐταῖς συγκατοικίζειν δάκη\\
θηρῶν, ἵν᾽ εἶχον μήτε προσφωνεῖν τινα\\
μήτ᾽ ἐξ ἐκείνων φθέγμα δέξασθαι πάλιν.\\
\francais{Il faudrait que les femmes n'eussent point auprès d'elles de servantes,\\
mais qu'elles fussent servies par de muets animaux,\\
pour qu'elles n'eussent personne à qui parler,\\
ni qui pût à son tour leur adresser la parole.}}}\\
καὶ ταῦτα μὲν ξυγγνώσθ᾽. ἃ δ᾽ ἦν ἡμῖν πρὸ τοῦ\\
αὐταῖς ταμιεῦσαι καὶ προαιρούσαις λαβεῖν\apparat{419}{ταμιεύεσθαι R : \lecon{ταμιεῦσαι καὶ} Reiske}\\
ἄλφιτον ἔλαιον οἶνον, οὐδὲ ταῦτ᾽ ἔτι
\commentaire{418-20}{ἃ δ᾽ ἦν ἡμῖν πρὸ τοῦ / αὐταῖς ταμιεύεσθαι προαιρούσαις λαβεῖν / ἄλφιτον ἔλαιον οἶνον.}{Je renvoie à la courte mais incisive note (la \nieme{12}) de Willems reproduite à la fin : il y fait un sort à la « déplorable conjecture de Scaliger » (qui avait substitué λαθεῖν à λαβεῖν au vers 419, donnant au passage un sens qui fait contre-sens).
}\\
ἔξεστιν. οἱ γὰρ ἄνδρες ἤδη κλῄδια
\\
αὐτοὶ φοροῦσι κρυπτὰ κακοηθέστατα
\\\passeLigne}%
\end{groupe}%
\endnumbering
\end{Leftside}
\begin{Rightside}
\beginnumbering
%% v. 380a
\begin{groupe}%
\personnageFr{Mika}%
\trimIambFr{Moi.\\\passeLigne}%
\end{groupe}%
%% v. 380b-382
\begin{groupe}%
\personnageFr{Coryphée}%
\trimIambFr{\ligne{380}\espace{1cm}Alors commence par ceindre cette couronne, avant \fin que parler. Paix ! \testimonia{380}{Suid. π 1166 περίθου νῦν τόνδε πρότερον}}%
\end{groupe}%
\begin{groupe}%
\tetramIambFr{Silence! Fais attention ! Car voici qu'elle tousse\testimonia{381}{(χρέμπτεται) \Phryn 126.1}
comme font les orateurs. Elle paraît en avoir long à dire.\\\passeLigne}%
\end{groupe}%
%% v. 383-432
\begin{groupe}%
\personnageFr{Mika}%
\trimIambFr{Aucun motif d'ambition, non par les deux déesses,\testimonia{383-4}{(φιλοτιμίᾳ -- γυναῖκες) Suidas φ 431}\\
ne m'a fait lever pour prendre la parole, ô femmes. Mais
je suis excédée, depuis longtemps, hélas !\\
de nous voir traînées dans la boue par \\
Euripide, le fils de la vendeuse d'herbes,\\
et constamment en butte à des injures de tout genre.\testimonia{389}{Suidas ε 2592 \espTest (ἐπισμῇ) Hesychius ε 5195 ἐπισμῇ}\\
Car quelle est l'injure dont il ne nous salisse ?\testimonia{390-4}{Suidas οι 128 (οἷς θέα pro εἰσὶν θεαταὶ et τραγῳδία pro τραγῳδοὶ)}\testimonia{390-1}{scholium Plato \italique{Theages} 127c (πῶς pro ποῦ et τραγῳδικοὶ pro τραγῳδοὶ καὶ) $=$ scholium Plato \italique{Gorgias} 457b $=$ scholium Plato \italique{Hippias min.} 365d \espTest (ποῦ -- θεαταὶ) Suidas ε 955 (sed εἰσὶ pro εἰσὶν)}
Où s'est-il fait faute de nous calomnier, pour peu simplement\\
qu'il y eût des spectateurs, des tragédiens et des choeurs ?\\
nous traitant de galantes, passionnées d'hommes,\testimonia{392}{(ἀνδρεραστρίας) Pollux III 71 ; \Phryn 34.18 ; \Eustathe 380.15}\\
sujettes au vin, traitresses, bavardes,\testimonia{393}{(οἰνοπίπας) scholium \italique{Aves} 307}}%
\end{groupe}%
%% v. 383-432
\begin{groupe}%
\trimIambFr{rien qui vaille, grand fléau des hommes.\\
Si bien qu'à peine rentrés du théâtre,\testimonia{395-6}{(ὥστ᾽ εὐθὺς -- ἡμᾶς) Suidas ι 275}\testimonia{395}{(ἰκρίων) Hesychius ι 501 ἴκρια}\\
ils nous soupçonnent, et guettent aussitôt\\
s'il n'est pas quelque amant caché.
Et plus rien ne nous est permis de ce que nous faisions autrefois,\\
tant il a enseigné de vilaines choses}
\end{groupe}%
%% v. 383-432
\begin{groupe}%
\trimIambFr{à nos époux. Une femme tresse-t-elle
une couronne, on la croit amoureuse ; laisse-t-elle tomber\\
un ustensile en allant et venant par la maison,\\
le mari demande : « Pour qui brise-t-on la marmite ?\\
Nul doute que ce ne soit pour l'étranger de Corinthe ».\testimonia{404}{(τῷ Κορινθίῳ ξένῳ) Hesychius κ 3629 Κορίνθιος ξένος}}
\end{groupe}%
%% v. 383-432
\begin{groupe}%
\trimIambFr{Une fille est-elle souffrante, aussitôt son frère de dire :\\
« Ce teint ne me plaît guère chez la fillette ».\\
Mais passons. Une femme privée d'enfants veut\\
en supposer un ; même cela ne se peut celer,\\
car les maris se tiennent tout proche.
Et auprès des vieillards, qui auparavant épousaient les jeunesses,
il nous a calomniées, tellement que pas un homme d'âge\\
ne consent à prendre femme, en raison de ce vers :\\
« Pour un époux âgé, la femme est un tyran ».\testimonia{413}{Menandros \italique{Monostich.} 191 ($=$ Euripides fr. 804.3 ed. Kannicht) ; \Apostol V 93b}\\
En outre, à cause de lui, sur les appartements de femmes,\\
ils appliquent des cachets et des barres}
\end{groupe}%
%% v. 383-432 (part 1)
\begin{groupe}\trimIambFr{pour nous garder, et par là-dessus ils entretiennent des chiens molosses,
épouvantails pour les galants.\testimonia{417}{(μορμολυκεῖα) \EtymMagnum 590.53}
Encore lui passerait-on cela : Mais ce qu'on nous permettait autrefois,
de diriger nous-mêmes le ménage et de prendre dans le cellier
de la farine, de l'huile, du vin, cela même ne nous est plus
possible : Car les maris portent sur eux de petites clefs\testimonia{421-8}{(οἱ γὰρ ἄνδρες -- ἐξαψαμένους) Suidas λ 54 (αὐτοὶ, τρεῖς et ᾡκότριψ Εὐριπίδης omissis, Λακωνίκα pro Λακωνίκ᾽, θριπηδέστατ᾽ pro θριπήδεστ᾽ et σφραγῖδα pro σφραγίδια)}\testimonia{421-7}{(οἱ γὰρ ἄνδρες -- θριπήδεστ᾽) Suidas κ 159 (αὐτοὶ et ᾡκότριψ Εὐριπίδης omissis, Λακωνίκα pro Λακωνίκ᾽ et θριπηδέστατ᾽ pro θριπήδεστ᾽)}\testimonia{421-7}{(οἱ γὰρ ἄνδρες -- θριπήδεστ᾽)}
secrètes, les plus malicieuses du monde,\testimonia{422}{(κακοηθέστατα) \Eustathe 974.37, 1108.25, 1603.52}\\\passeLigne}%
\end{groupe}%
\endnumbering
\end{Rightside}
\end{pages}
\Pages
\clearpage\ifodd\value{page}\hbox{}\newpage\fi
\enTeteG{ΕΠΕΙΣΟΔΙΟΝ Α' --- Codex Ravennas 429, f. 167v (Classense 169v) --- vv. 423–465}
\begin{adjustwidth}{-1.5cm}{-3.0cm}
\begin{adjustbox}{max totalheight=21.0cm, max width=16.5cm, keepaspectratio}
\begin{minipage}{16.5cm}
{\fontsize{7.2pt}{8.5pt}\selectfont
\noindent\textbf{Folio 11 --- Codex Ravennas 429, \texttt{f. 167v} (\texttt{Classense 169v}) : ΕΠΕΙΣΟΔΙΟΝ Α', vv. 423–465}\hfill \textit{Témoins :} $\text{R}^1, \Rpc, \Rdeux$ (\Mon)\par\vspace{2pt}
\hrule height 0.4pt\vspace{2.5pt}
}
{\fontsize{8.1pt}{9.5pt}\selectfont
\ligneRav{423}{\auDessus{γυ}{ν}}{Λακωνίκ᾽ ἄττα, τρεῖς ἔχοντα γομφίους·}{}
\ligneRav{}{}{πρὸ τοῦ μὲν οὖν ἦν ἀλλ᾽ ὑποῖξαι τὴν θύραν}{}
\ligneRav{425}{}{ποιησαμέναισι δακτύλιον τριωβόλου,}{}
\ligneRav{426}{}{νῦν δ᾽ οὗτος αὐτοὺς \difference{ᾡκότριψ} Εὐριπίδης}{\textbf{426} ὡκότριψ R : \lecon{ᾡκότριψ} Bentley : οἱκότριψ Kühner-Blass}
\ligneRav{427}{}{\difference{ἐδίδαξε} θριπήδεστ᾽ ἔχειν σφραγίδια}{\textbf{427} ἐδίδαξεν R : \lecon{ἐδίδαξε} Suid. λ 64}
\ligneRav{}{}{ἐξαψαμένους. νῦν οὖν ἐμοὶ τούτῳ δοκεῖ}{}
\ligneRav{429}{}{ὄλεθρόν τιν᾽ ἡμᾶς κυρκανᾶν \difference{ἁμωσγέπως},}{\textbf{429} ἀμωσγέπως R, Suid. : \lecon{ἁμωσγέπως} Blaydes}
\ligneRav{430}{}{ἢ φαρμάκοισιν ἢ μιᾷ γέ τῳ τέχνῃ,}{}
\ligneRav{}{}{ὅπως ἀπολεῖται. ταῦτ᾽ ἐγὼ φανερῶς λέγω,}{}
\ligneRav{}{}{τὰ δ᾽ ἄλλα μετὰ τῆς γραμματέως συγγράψομαι·}{}
\ligneRav{433}{—}{\hspace{1.1em}ἐγὼ γὰρ αὐτὴ πρῶτον, ἵνα μάλιστ᾽ \difference{οὔπω} (οὐπώ\inRas{τε}{ποτε})}{\textbf{433} οὔπω τε \Rac : οὐπώποτε \Rpc : οὔπω Dindorf}
\ligneRav{}{}{\hspace{1.1em}πολυπλοκωτέρας γυναικὸς}{}
\ligneRav{}{}{\hspace{1.1em}οὐδὲ δεινότερον λεγούσης·}{}
\ligneRav{437}{}{\hspace{1.1em}πάντα γὰρ λέγει δίκαια,}{\textbf{437} \lecon{πάσας} εἰδέας \lecon{ἐξήτασεν} R (sed \lecon{ἰδέας} Suid. ε 22) : πᾶσαν ἰδέαν ἐξετάζει Hermann}
\ligneRav{437}{}{\hspace{1.1em}†\difference{πάσας} δ᾽ \difference{ἰδέας} ἐξήτασεν†,}{\textbf{437} \lecon{πάσας} εἰδέας \lecon{ἐξήτασεν} R (sed \lecon{ἰδέας} Suid. ε 22) : πᾶσαν ἰδέαν ἐξετάζει Hermann}
\ligneRav{438}{}{\hspace{1.1em}πάντα δ᾽ ἐβάστασεν φρενὶ πυκνῶς τε}{\textbf{438} ἐβάστασεν R : \lecon{-σε} Suid.}
\ligneRav{439}{}{\hspace{1.1em}ποικίλους λόγους \difference{ἀνηῦρεν}}{\textbf{439} ἀνεῦρεν R, Suid. : \lecon{ἀνηῦρεν} Meineke}
\ligneRav{439}{}{\hspace{1.1em}εὖ διεζητημένους·}{\textbf{439} ἀνεῦρεν R, Suid. : \lecon{ἀνηῦρεν} Meineke}
\ligneRav{440}{}{\hspace{1.1em}ὥστ᾽ ἂν εἰ λέγοι παρ᾽ αὐτῆς}{\textbf{440} \lecon{αὐτῆς} R : αὐτὴν Zanetti}
\ligneRav{441}{}{\hspace{1.1em}\difference{Ξενοκλέης} ὁ Καρκίνου, δο-}{\textbf{441} Ξενοκλῆς R : \lecon{Ξενοκλέης} Bentley}
\ligneRav{441}{}{\hspace{1.1em}κεῖν ἂν αὐτόν, ὡς ἐγᾦμαι,}{\textbf{441} Ξενοκλῆς R : \lecon{Ξενοκλέης} Bentley}
\ligneRav{}{}{\hspace{1.1em}πᾶσιν ὑμῖν}{}
\ligneRav{}{}{\hspace{1.1em}ἄντικρυς μηδὲν λέγειν·}{}
\ligneRav{443}{γυνὴ ἑτέρα}{ὀλίγων \difference{ἕνεκα καὐτὴ} παρῆλθον ῥημάτων·}{\textbf{443} γυνὴ ἑτέρα \Rpc \espApp ἕνεκ᾽ αὐτὴ R : \lecon{ἕνεκα καὐτὴ} Porson}
\ligneRav{}{}{τὰ μὲν γὰρ ἄλλ᾽ αὕτη κατηγόρηκεν εὖ·}{}
\ligneRav{445}{}{ἃ δ᾽ ἐγὼ πέπονθα, ταῦτα λέξαι βούλομαι·}{}
\ligneRav{446}{}{ἐμοὶ γὰρ \difference{ἁνὴρ} ἀπέθανεν μὲν \difference{ἔγκυος}}{\textbf{446} ἀνὴρ R : \lecon{ἁνὴρ} Brunck \espApp ἐν Κύπρῳ R : \lecon{ἔγκυος} Bibi}
\ligneRav{}{}{παιδάρια πέντε καταλιπών, ἁγὼ μόλις}{}
\ligneRav{}{}{στεφανηπλοκοῦσ᾽ ἔβοσκον ἐν ταῖς μυρρίναις·}{}
\ligneRav{}{}{τέως μὲν οὖν ἀλλ᾽ ἡμικάκως ἐβοσκόμην·}{}
\ligneRav{450}{}{νῦν δ᾽ οὗτος ἐν \difference{ταῖσιν} τραγῳδίαις ποιῶν}{\textbf{450} ταῖσι R : \lecon{ταῖσιν} Brunck}
\ligneRav{}{}{τοὺς ἄνδρας ἀναπέπεικεν οὐκ εἶναι θεούς·}{}
\ligneRav{}{}{ὥστ᾽ οὐκέτ᾽ ἐμπολῶμεν οὐδ᾽ εἰς ἥμισυ·}{}
\ligneRav{453}{}{νῦν οὖν \difference{ἁπάσαισιν} παραινῶ καὶ λέγω}{\textbf{453} ἁπάσαισι R : \lecon{ἁπάσαισιν} Gellius \italique{Ν.Α.} 15.20}
\ligneRav{}{}{τοῦτον κολάσαι τὸν ἄνδρα πολλῶν οὕνεκα·}{}
\ligneRav{455}{}{ἄγρια γὰρ ἡμᾶς ὦ γυναῖκες δρᾷ κακά,}{}
\ligneRav{456}{}{ἅτ᾽ ἐν ἀγρίοισι \difference{τοῖς} λαχάνοις αὐτὸς τραφείς·}{\textbf{456} \lecon{τοῖς} Plut., \italique{Mor.}, 853C, Gellius \italique{N.Α.}15.20 : om. R}
\ligneRav{}{}{ἀλλ᾽ εἰς ἀγορὰν ἄπειμι· δεῖ γὰρ ἀνδράσιν}{}
\ligneRav{}{}{πλέξαι στεφάνους συνθηματιαίους εἴκοσιν·}{}
\ligneRav{}{χορ/}{\hspace{1.1em}ἕτερον αὖ τι λῆμα τοῦτο}{}
\ligneRav{460}{}{\hspace{1.1em}κομψότερον ἔτ᾽ ἢ τὸ πρότερον}{}
\ligneRav{460}{}{\hspace{1.1em}ἀναπέφηνεν·}{}
\ligneRav{461}{}{\hspace{1.1em}οἷα κατεστωμύλατο}{\textbf{461} \lecon{κατεστωμύλατο} R : κἀστωμύλατο Dobree}
\ligneRav{462}{}{\hspace{1.1em}οὐκ \difference{ἄκαιρα}, φρένας ἔχουσα}{\textbf{462} ἄκερα R : \lecon{ἄκαιρα} Zanetti}
\ligneRav{463}{}{\hspace{1.1em}καὶ πολύπλοκον νόημ᾽, οὐδ᾽}{\textbf{463} \lecon{πολύπλοκον} R : τι πολύπλοκον Dobree : πολύστροφον Hermann}
\ligneRav{}{}{\hspace{1.1em}ἀσύνετ᾽ ἀλλὰ πιθανὰ πάντα·}{}
\ligneRav{465}{}{\hspace{1.1em}δεῖ δὲ ταύτης τῆς ὕβρεως ἡ-}{}
\ligneRav{465}{}{\hspace{1.1em}μῖν τὸν ἄνδρα}{}
\ligneRav{465}{}{\hspace{1.1em}περιφανῶς δοῦναι δίκην·}{}
}
\vspace{3pt}\noindent\rule{4.5cm}{0.4pt}\par\vspace{1.5pt}
{\fontsize{7.0pt}{8.3pt}\selectfont
\noindent\textbf{Loci paralleli (vv. 423–465) ---} \textbf{423} (γομφίους) \Phryn 61.1 \espTest \textbf{426} (οἰκότριψ $=$ οἰκοτραφής) Moeris ο 25 ; (ᾡκότριψ) Hesychius ο 265-266 ; \AelDion ο 8-9 \espTest \textbf{428-9} (ἐμοὶ -- ἁμωσγέπως) Suidas κ 2771 (τινα pro τιν᾽ ἡμᾶς) \espTest \textbf{429} (κυρκανᾶν) Photius κ 1254 κυρκανᾶν $=$ Hesychius κ 4691 κυρκανᾷ ; \Eustathe 1213.50 \espTest (ἁμωσγέπως) Hesychius α 4182 ἁμωσγέπως ; \Syn α 449 $=$ Suidas α 1645 \espTest \textbf{432} (μετὰ τῆς γραμματέως συγγράψομαι) \Eustathe 1484.53 \espTest (τῆς γραμματέως) Pollux IV 19 \espTest \textbf{437-9} (πάσας -- ἀνηῦρεν) Suidas ε 22 πάσας δ' ἰδέας ἐξήτασε πάντα δ' ἐβάστασε φρενί, πυκνῶς τε ποικίλως λόγους ἀνεῦρεν \espTest \textbf{438} (πυκνῶς) Hesychius π 4335 πυκνῶς \espTest \textbf{442} (ἄντικρυς) Photius α 2210 ἄντικρυς ; \AelDion α 107 \espTest \textbf{446} Hesychius ε 239 ἔγκυος· ἐγγαστρωμένη γυνή \espTest \textbf{447-9} (μόλις -- τέως μὲν οὖν ἀλλ᾽) scholium\expo{T} Plato \italique{Hipparch.} 229d (μυρίναις pro μυρρίναις et μέν οὐκ ἄν pro μὲν οὖν ἀλλ᾽) \espTest \textbf{448} (μυρρίνη $=$ μυρσίνη) Moeris μ 23 \espTest \textbf{449} Suidas τ 321 \espTest (ἡμικάκως) Pollux VI 161 \espTest \textbf{453-6} Aulus Gellius XV 20.7 \espTest \textbf{455-6} Plut. \italique{Moral.}, \italique{Compar. Arist. et Menandr.} (853c) \espTest \textbf{457-8} (δεῖ γὰρ -- εἴκοσιν) Suidas σ 1591 (sed ἀνδράσι et εἴκοσι pro ἀνδράσιν et εἴκοσιν) \espTest \textbf{458} Athenaeus XV 680c \espTest (στεφάνους συνθηματιαίους) Pollux VII 200 \espTest (συνθηματιαίους) Hesychius σ 2624 συνθηματιαῖοι \espTest \textbf{459-60} (ἕτερον -- ἀναπέφηνεν) Suidas λ 441 (ἀναπέφηνε pro ἀναπέφηνεν)\par
}
\vspace{2.5pt}\noindent\rule{4.5cm}{0.4pt}\par\vspace{1.5pt}
{\fontsize{6.9pt}{8.2pt}\selectfont
\noindent\textbf{Suite (en commentaire) de l'apparat critique \& notes paléographiques (\texttt{f. 167v}) ---} \textbf{Correction au v. 433 :} \Rac{} avait écrit οὔπω τε, corrigé par \Rpc{} en οὐπώποτε (Dindorf restitue οὔπω). --- \textbf{Note d'interlocuteur au v. 443 :} \Rpc{} inscrit γυνὴ ἑτέρα. --- \textbf{Conjecture remarquable de \textbf{Bibi} au v. 446 :} ἀνὴρ μὲν ἐν Κύπρῳ R $\to$ \lecon{ἁνὴρ μὲν ἔγκυος} (ἔγκυος se rapportant à la femme ou construit avec l'idée que son mari l'a laissée enceinte de son 5\ieme{} enfant avant de mourir, au lieu de la localisation géographique hors de propos ἐν Κύπρῳ). --- \textbf{v. 437 :} \lecon{πάσας} εἰδέας \lecon{ἐξήτασεν} R (mais \lecon{ἰδέας} Souda $\varepsilon\ 22$, sans iotacisme ; πᾶσαν ἰδέαν ἐξετάζει Hermann). --- \textbf{v. 443 :} \lecon{γυνὴ ἑτέρα} \Rpc{} \espApp ἕνεκ' αὐτὴ R (haplographie de $\kappa\alpha-$) : \lecon{ἕνεκα καὐτὴ} Porson. --- \textbf{v. 446 :} ἀνὴρ R : \lecon{ἁνὴρ} Brunck \espApp ἐν Κύπρῳ R : \lecon{ἔγκυος} \lecon{Bibi}. --- \textbf{v. 456 :} \lecon{τοῖς} Plutarque, \italique{Moralia} 853C, Aulu-Gelle 15.20 : om. R (devant λαχάνοις).\par
}
\end{minipage}
\end{adjustbox}
\end{adjustwidth}
\newpage
\enTeteD{Codex Ravennas 429 (Biblioteca Classense) --- Fac-similé du f. 167v (vv. 423–465)}
\begin{adjustwidth}{-1.5cm}{-3.0cm}
\begin{center}
\includegraphics[width=16.5cm,height=20.5cm,keepaspectratio]{/workspaces/Thesmophories/dossier_ravennas_thesmophories/folios_individuels_ff162v-177r/11_f167v_classense169v_vv0423-0465.jpg}\par
\vspace{2pt}
{\scriptsize\textbf{Codex Ravennas 429, f. 167v (Classense 169v)} --- ΕΠΕΙΣΟΔΙΟΝ Α', \textbf{vv. 423–465} ($\text{R}^1, \Rpc, \Rdeux$ (\Mon)).}
\end{center}
\end{adjustwidth}
\newpage
\clearpage\ifodd\value{page}\hbox{}\newpage\fi
\enTeteG{ΘΕΣΜΟΦΟΡΙΑΖΟΥΣΑΙ --- Texte grec établi (f. 167v, vv. 423–465)}
\enTeteD{LES THESMOPHORIES --- Traduction française (f. 167v, vv. 423–465)}
\begin{pages}
\begin{Leftside}
\beginnumbering
%% v. 423+ (part 2)
\begin{groupe}
\trimIamb{\ligne{423}Λακωνίκ᾽ ἄττα, τρεῖς ἔχοντα γομφίους.%\testimonia{423}{(γομφίους) \Phryn 61.1}
\\
πρὸ τοῦ μὲν οὖν ἦν ἀλλ᾽ ὑποῖξαι τὴν θύραν
\commentaire{424}{ὑποῖξαι τὴν θύραν.}{Il s'agit de la porte de l'appartement (cf. v. 414) et non de celle du cellier.}\\
ποιησαμέναισι δακτύλιον τριωβόλου,\\
νῦν δ᾽ οὗτος αὐτοὺς ᾡκότριψ Εὐριπίδης%\testimonia{426}{(οἰκότριψ $=$ οἰκοτραφής) Moeris ο 25 ; (ᾡκότριψ) Hesychius ο 265-266 ; \AelDion ο 8-9}
\apparat{426}{ὡκότριψ R : \lecon{ᾡκότριψ} Bentley : οἱκότριψ Kühner-Blass \ligneApp{427} ἐδίδαξεν R : \lecon{ἐδίδαξε} Suid. λ 64 \ligneApp{429} ἀμωσγέπως R, Suid. : \lecon{ἁμωσγέπως} Blaydes}%
\commentaire{426}{ᾡκότριψ Εὐριπίδης.}{Ce terme, \italique{tâte-poule}, dont un Grec aurait saisi le sens métaphorique (cf. les verbes βλιμάζειν, σκιμαλίζειν), rend bien le mot ᾡκότριψ. Cf. ἀστύτριψ, \italique{pilier de cité}, Θησειότριψ, \italique{pilier du Théseion} (fragm. du Polyides), πορνότριψ, etc. (Note de Willems.)}\\
ἐδίδαξε θριπήδεστ᾽ ἔχειν σφραγίδια\\
ἐξαψαμένους. νῦν οὖν ἐμοὶ τούτῳ δοκεῖ%\testimonia{428-9}{(ἐμοὶ -- ἁμωσγέπως) Suidas κ 2771 (τινα pro τιν᾽ ἡμᾶς)}
\\
ὄλεθρόν τιν᾽ ἡμᾶς κυρκανᾶν ἁμωσγέπως,%\testimonia{429}{(κυρκανᾶν) Photius κ 1254 κυρκανᾶν $=$ Hesychius κ 4691 κυρκανᾷ ; \Eustathe 1213.50 \espTest (ἁμωσγέπως) Hesychius α 4182 ἁμωσγέπως ; \Syn α 449 $=$ Suidas α 1645}%\apparat{429}{ἀμωσγέπως R, Suid. : \lecon{ἁμωσγέπως} Blaydes}%
\commentaire{429}{ὄλεθρόν τιν᾽ ἡμᾶς κυρκανᾶν ἁμωσγέπως.}{Κυρκανᾶν est un synonyme très rare de κυκᾶν et se dit proprement d'un cuisinier préparant une mixture\marqueur{1}. Il est pris ici au sens de \italique{comploter la perte} de quelqu'un. Il est clair que l'expression ὄλεθρόν κυρκανᾶν est faite sur le modèle de κυρκανᾶν φάρμακα (cf. \italique{Ploutos} 302, 309 : τὰ φάρμακ᾽ ἀνακυκῶσαν) [...] Pour le ton, le verbe κυρκανᾶν paraît correspondre à nos métaphores familières \italique{cuisiner}, \italique{mitonner une affaire}. (Note de Taillardat, §421.)}%
\debutNote{429}%
\note{1}{Cf. Épinicus 1,8 : συγκυρκανᾶν (préparation d'une boisson faite de farine, de miel et de vin mêlés).}\\
ἢ φαρμάκοισιν ἢ μιᾷ γέ τῳ τέχνῃ,\\
%
ὅπως ἀπολεῖται. ταῦτ᾽ ἐγὼ φανερῶς λέγω,\\
%
τὰ δ᾽ ἄλλα μετὰ τῆς γραμματέως συγγράψομαι.%\testimonia{432}{(μετὰ τῆς γραμματέως συγγράψομαι) \Eustathe 1484.53 \espTest (τῆς γραμματέως) Pollux IV 19}%
\commentaire{432}{τῆς γραμματέως.}{Aristophane faisait-il sourire sa mâle assemblée lorsqu'il parlait de \italique{la} greffièr\italique{e} ? c'est ce que suggère Pollux (IV, 19) :%
\citComm{Ἀριστοφάνης δὲ παίζει ἐν Θεσμοφοριαζούσαις, λέγων ἡ γραμματεύς.}%
}\\\passeLigne}%
\end{groupe}%
%% v. 433-442
\begin{groupe}%
%
\personnageGrec{Χορός\strophe{στροφή}}\typeVers{par.}\chantChoeur{οὔπω ταύτης ἤκουσα\apparat{433}{οὔπω τε \Rac : οὐπώποτε \Rpc : οὔπω Dindorf}%
\commentaire{433-42}{}{Ce chant du chœur trouvera son écho aux vers 520-530 : semblable réplique au point de vue métrique\marqueur{1} qui permettait aux deux passages d'être chantés sur la même mélodie. Là comme ici, le chœur s'exprime à la première personne.}%
\debutNote{433-42}%
\note{1}{À une exception près : au vers 437b où le texte est probablement corrompu.}}%
\end{groupe}%
%% v. 434-5
\begin{groupe}%
\decoupeVers{5}\typeVers{$2$ troch.}\chantChoeur{πολυπλοκωτέρας γυναικὸς}%
\end{groupe}%
%% v. 436
\begin{groupe}%
\ligne{436}\typeVers{$2$ troch.}\chantChoeur{οὐδὲ δεινότερον λεγούσης.\apparat{435}{\lecon{δεινότερον} R : δεινότερα Hermann}}%
\end{groupe}%
%% v. 437a
\begin{groupe}%
\decoupeVers{a}\typeVers{$2$ troch.}\chantChoeur{πάντα γὰρ λέγει δίκαια,}%
\end{groupe}%
%% v. 437b
\begin{groupe}%
\ligne{437}\decoupeVers{b}\typeVers{$2$ iamb.}\chantChoeur{†πάσας δ᾽ ἰδέας ἐξήτασεν†,%\testimonia{437-9}{(πάσας -- ἀνηῦρεν) Suidas ε 22 πάσας δ' ἰδέας ἐξήτασε πάντα δ' ἐβάστασε φρενί, πυκνῶς τε ποικίλως λόγους ἀνεῦρεν}
\apparat{437}{\lecon{πάσας} εἰδέας \lecon{ἐξήτασεν} R (sed \lecon{ἰδέας} Suid. ε 22) : πᾶσαν ἰδέαν ἐξετάζει Hermann \ligneApp{438} ἐβάστασεν R : \lecon{-σε} Suid. \ligneApp{439} %\lecon{ποικίλους} R : -λως Suid. \espApp 
ἀνεῦρεν R, Suid. : \lecon{ἀνηῦρεν} Meineke}}% Austin comme Fritzsche lisent ἀνεῦρεν dans R, mais semblent se tromper : Hall, Coulon & Prato contre eux
\end{groupe}%
%% v. 438
\begin{groupe}%
\typeVers{$2$ troch.}\chantChoeur{πάντα δ᾽ ἐβάστασεν φρενὶ πυκνῶς τε%\testimonia{438}{(πυκνῶς) Hesychius π 4335 πυκνῶς}
}%
\end{groupe}%
%% v. 439a
\begin{groupe}%
\decoupeVers{a}\typeVers{$2$ troch.}\chantChoeur{ποικίλους λόγους ἀνηῦρεν}%
\end{groupe}%
%% v. 439b
\begin{groupe}%
\ligne{439}\decoupeVers{b}\typeVers{$2$ troch\cat}\chantChoeur{εὖ διεζητημένους·}%
\end{groupe}%
%% v. 440
\begin{groupe}%
\typeVers{$2$ troch.}\chantChoeur{ὥστ᾽ ἂν εἰ λέγοι παρ᾽ αὐτῆς%\apparat{440}{\lecon{αὐτῆς} R : αὐτὴν Zanetti}
}%
\end{groupe}%
%% v. 441a
\begin{groupe}%
\decoupeVers{a}\typeVers{$2$ troch.}\chantChoeur{Ξενοκλέης ὁ Καρκίνου, δο-\apparat{441}{Ξενοκλῆς R : \lecon{Ξενοκλέης} Bentley}%
\commentaire{441}{Ξενοκλῆς ὁ Καρκίνου.}{Xénoclès, le plus célèbre des trois carcinites était un poète tragique non dépourvu de talent puisqu'en 415 il l'emporta sur Euripide\marqueur{1}. Mais on lui reprochait de chercher à captiver son auditoire par des moyens peu relevés. Platon le comique l'appelle δωδεκαμήχανος [l'homme] \italique{aux douze postures}, surnom qu'on donnait à la courtisane Cyrène\marqueur{2}. Aristophane le traite de μηχανοδίφης, inventeur de trucs. Il résulte d'un passage d'Athénée (IV, 134d), mis en rapport avec le nôtre, que Xénoclès cumulait avec la profession d'acteur et de poète tragique, celle d'orateur. Ce fut aussi le cas d'Eschine. Cela n'a rien d'étonnant, le métier d'acteur étant une excellente préparation à celui d'orateur. Il n'y a qu'à se souvenir de Démosthène et de ses cailloux. (Note de Willems ; voir aussi la note (t. I pp. 553-558) qu'il a consacrée au banquet de Philocléon conté par Xanthias (\italique{Guêpes}, vv. 1299-1323) où il est question des carcinites.)}%
\debutNote{441}%
\note{1}{Aristophane se montre toutefois sévère avec Xénoclès : il est méchant comme ses vers (v. 169) et son art comme sa personne ne méritent pas de survivre (\italique{Grenouilles}, v. 86).}}%
\note{2}{Voir la note au vers 235.}%
\end{groupe}%
%% v. 441b
\begin{groupe}%
\ligne{441}\decoupeVers{b}\typeVers{$2$ troch.}\chantChoeur{κεῖν ἂν αὐτόν, ὡς ἐγᾦμαι,}%
\end{groupe}%
%% v. 442a
\begin{groupe}%
\decoupeVers{a}\typeVers{troch.}\chantChoeur{πᾶσιν ὑμῖν\commentaire{442}{πᾶσιν ὑμῖν.}{Les spectateurs.}}%
\end{groupe}%
%% v. 442b
\begin{groupe}%
\ligne{442}\decoupeVers{b}\typeVers{$2$ troch\cat}\chantChoeur{ἄντικρυς μηδὲν λέγειν.%\testimonia{442}{(ἄντικρυς) Photius α 2210 ἄντικρυς ; \AelDion α 107}
\\\passeLigne}%
\end{groupe}%
%% v. 443-458
\begin{groupe}%
%
\personnageGrec{Στεφανόπωλις}%
\typeVers{trim. iamb.}\trimIamb{\ligne{443}ὀλίγων ἕνεκα καὐτὴ παρῆλθον ῥημάτων.\apparat{443}{γυνὴ ἑτέρα \Rpc %: Γυ. Β. Brunck : Κρίτυλλα Sommerstein : \lecon{Στεφανόπωλις} Rossella 
\espApp ἕνεκ᾽ αὐτὴ R : \lecon{ἕνεκα καὐτὴ} Porson}\\
τὰ μὲν γὰρ ἄλλ᾽ αὕτη κατηγόρηκεν εὖ·\\
ἃ δ᾽ ἐγὼ πέπονθα, ταῦτα λέξαι βούλομαι.\\
ἐμοὶ γὰρ ἁνὴρ ἀπέθανεν μὲν ἔγκυος%\testimonia{446}{Hesychius ε 239 ἔγκυος· ἐγγαστρωμένη γυνή}
\apparat{446}{ἀνὴρ R : \lecon{ἁνὴρ} Brunck \espApp ἐν Κύπρῳ R : \lecon{ἔγκυος} Bibi}
\commentaire{446}{ἐμοὶ γὰρ ἁνὴρ ἀπέθανεν μὲν ἐν Κύπρῳ.}{La précision du lieu de décès du mari (ἁνὴρ ἀπέθανεν ἐν Κύπρῳ) n'a pas manqué d'intriguer les modernes scholiastes.
Leur inclination naturelle étant portée au scabreux, ils ont pour une bonne partie voulu y voir un tour signifiant que la matrone athénienne, mère de cinq enfants, avait tué son mari au déduit avec ses appétits démesurés ; «~mourir de sa belle mort~», comme l'heureux \italique{Félix Faurenicator}. 
L'hypothèse est plaisante mais gratuite et donc à rejeter ! tous les lexiques consultés ne soufflent mot d'une telle expression qu'on aurait nécessairement relevée. 
D'autres se sont essayés à prendre ce vers littéralement. Là encore, ils se sont heurtés à un mur : les dernière expéditions sur l'île de Chypre feraient de notre procurette une septuagénaire au moins !.. et comme il ne semble pas y avoir de jeux de mots, il ne reste qu'une seule voie : une erreur de transcription du texte. 
(C'est d'autant plus vraisemblable que l'expression « ἀπέθανεν ἐν Κύπρῳ », si elle s'était véritablement trouvée dans le texte original, n'aurait pas manqué d'intriguer les commentateurs anciens qui, même sans y rien connaître, auraient ouvert les vannes de leur verve intarissable... au contraire : ce vers n'a étrangement pas laissé la moindre scholie ancienne derrière lui.)\\
L'intervention de la tresseuse de couronnes n'a donc rien de burlesque : elle essaie au contraire d'introduire du pathétique. Or, qu'y a-t-il de plus banal pour apitoyer son auditoire à peu de frais que de sortir le lieu-commun de la femme qui se retrouve veuve alors qu'elle est enceinte (ἔγκυος) ? 
\citComm{ἐμοὶ γὰρ ἁνὴρ ἀπέθανεν μὲν \textbf{ἔγκυος}\\
Mon mari mourut alors que j'étais enceinte \\
Me laissant cinq enfants à charge etc. etc.}}\\
παιδάρια πέντε καταλιπών, ἁγὼ μόλις%\testimonia{447-9}{(μόλις -- τέως μὲν οὖν ἀλλ᾽) scholium\expo{T} Plato \italique{Hipparch.} 229d (μυρίναις pro μυρρίναις et μέν οὐκ ἄν pro μὲν οὖν ἀλλ᾽)}
\\
στεφανηπλοκοῦσ᾽ ἔβοσκον ἐν ταῖς μυρρίναις.%\testimonia{448}{(μυρρίνη $=$ μυρσίνη) Moeris μ 23}
\\
%
τέως μὲν οὖν ἀλλ᾽ ἡμικάκως ἐβοσκόμην·%\testimonia{449}{Suidas τ 321 \espTest (ἡμικάκως) Pollux VI 161}
\\
νῦν δ᾽ οὗτος ἐν ταῖσιν τραγῳδίαις ποιῶν\apparat{450}{ταῖσι R : \lecon{ταῖσιν} Brunck}\\
%
τοὺς ἄνδρας ἀναπέπεικεν οὐκ εἶναι θεούς·
\commentaire{450-1}{ἐν ταῖσιν τραγῳδίαις ποιῶν / τοὺς ἄνδρας ἀναπέπεικεν οὐκ εἶναι θεούς·}{Petit florilège des \italique{\grec{γνῶμαι}} théoclastes d'Euripide :\\
\italique{Hécube}, v. 489-492 :%
\citComm{Ταλθύβιος — ὦ Ζεῦ, τί λέξω; πότερά σ᾽ ἀνθρώπους ὁρᾶν;\\
ἢ δόξαν ἄλλως τήνδε κεκτῆσθαι μάτην,\\
ψευδῆ, δοκοῦντας δαιμόνων εἶναι γένος\\
τύχην δὲ πάντα τἀν βροτοῖς ἐπισκοπεῖν;}%
\italique{Béllérophon}\marqueur{1} :%
\citComm{Βελλεροφόντης — φῆσίν τις εἶναι δῆτ᾽ ἐν οὐρανῷ θεούς;\\
οὐκ εἰσίν, οὐκ εἴσ᾽,}%
Prologue de \italique{Mélanippe la sage}\marqueur{2} (douteux) :%
\citComm{Ζεὺς ὅςτις ὁ Ζεύς, οὐ γὰρ οἶδα πλὴν λόγῳ\\
La véritable version euripidéenne étant peut-être\marqueur{3} :\\
Μελάνιππη — Ζεύς, ὡς λέλεκται τῆς ἀληθείας ὕπο,\\
Ἕλλην᾽ ἔτιχθ᾽, ὃς ἐξέφυσεν Αἴολον·}%
}%
\debutNote{450-1}%
\note{1}{Deux premières lignes du \nieme{10} fragment de l'édition Jouan-van Looy aux Belles Lettres ; fr. 286 de l'édition Kannicht ; fr. 385 Mette.}%
\note{2}{Fragment douteux car on recense différentes versions de ce début de prologue. Citation apocryphe ? retouche d'Euripide d'un vers ayant fait scandale ? quoi qu'il en soit, Jouan \& van Looy le mettent en tête de leur \italique{Mélanippe} sans le comptabiliser ; fr. 480 de l'édition Kannicht.}%
\note{3}{Premier fragment de l'édition Jouan-van Looy ; 481 chez Kannicht et 665a-c chez Mette.}%
}
\end{groupe}%
%% v. 443-458
\begin{groupe}%
%
\trimIamb{ὥστ᾽ οὐκέτ᾽ ἐμπολῶμεν οὐδ᾽ εἰς ἥμισυ.\\
%\commentaire{452}{ὥστ᾽ οὐκέτ᾽ ἐμπολῶμεν οὐδ᾽ εἰς ἥμισυ.}{Le reproche apparent de l'oratrice est que l'impiété d'Euripide a entraîné... la chute des ventes de couronnes ; accusation on-ne-peut-plus frivole ! et qui déchaîne pourtant l'enthousiasme du chœur (\grec{κομψότερον ἔτ᾽ ἢ τὸ πρότερον ... οἷα κατεστωμύλατο ... φρένας ἔχουσα \&c.}). En outrec cette crise couronnomique est en pleine contradiction avec la fin de son discours~: quand elle dit s'en aller car elle a quantité de commandes à honorer. Il y a donc certainement jeu de mots et double discours ; je mettrais ma main à couper qu'il se trouve au vers 452, à rapprocher du vers 456. Je laisse à d'autres le soin d'expliquer que la vendeuse de couronnes dit à mi-mots se prostituer~; ce sont d'ailleurs les mêmes dignes professeurs qui, après avoir estimé son âge à une soixantaine d'année, nous lancent avec le même ton que cette mère de cinq enfants ferait le tapin dans le marché aux myrthes (hypothèses qui, conjuguées, ne méritent qu'un grand éclat de rire, le même que celui du film \italique{La traversée de Paris}). Restons sérieux~: s'il y a quelques sous-entendus salaces, ils ne sont certainement pas là pour que l'oratrice se décrédibilise elle-même~; faut-il le rappeler~: la cible est Euripide !
%Le poète est ridiculisé ailleurs (v. 1044-1048 des \italique{Grenouilles}) comme cocu notoire. Qu'il soit ici aussi joué sur ce plan et sur sa consommation excessive de certains légumes réputés par les Grecs ne pas rendre très vaillants~; ces déboires entraînant chez le sujet l'agraire aigreur de ses euripidéennes tragédies... voici qui, à défaut d'une démonstration, a au moins le mérite de la vraisemblance et de justifier les éloges du chœur. Et si, comme je le crois, il n'y a pas là de sous-entendu obscène, alors il s'agirait simplement d'une explication de l'impiété d'Euripide du fait que ce ne soit pas un homme du ciel, mais de la terre (le verbe \grec{πολεῖν} signifie \italique{retourner la terre} ; attaque identique au vers 840 des \italique{Grenouilles})}%
%
νῦν οὖν ἁπάσαισιν παραινῶ καὶ λέγω%\testimonia{453-6}{Aulus Gellius XV 20.7}
\apparat{453}{ἁπάσαισι R : \lecon{ἁπάσαισιν} Gellius \italique{Ν.Α.} 15.20}\\
τοῦτον κολάσαι τὸν ἄνδρα πολλῶν οὕνεκα·\\
ἄγρια γὰρ ἡμᾶς ὦ γυναῖκες δρᾷ κακά,%\testimonia{455-6}{Plut. \italique{Moral.}, \italique{Compar. Arist. et Menandr.} (853c)}
\\
ἅτ᾽ ἐν ἀγρίοισι τοῖς λαχάνοις αὐτὸς τραφείς.\apparat{456}{\lecon{τοῖς} Plut., \italique{Mor.}, 853C, Gellius \italique{N.Α.}15.20 : om. R}
\commentaire{455-6}{Dossier sur la mère d'Euripide, la vendeuse d'herbes (λαχανοπωλήτρια).}{(\italique{Acharn.} vv. 478-479 ; \italique{Cavaliers} vv. 19-20 ; \italique{Thesm.} vv. 455-456, 910 ; \italique{Grenouilles} vv. 840, 946-947.)


Aristophane plaisante Euripide à maintes reprises sur la profession de sa mère, vendeuse d'herbes\marqueur{1}. Il fit ses choux gras d'une herbe en particulier : ἡ σκάνδιξ, \italique{le scandix peigne de Vénus}. Les commentateurs, anciens comme modernes, se sont interrogés depuis lors pour savoir s'il fallait prendre littéralement cette information, ou bien de façon imagée.
Sans parvenir à aucune conclusion convaincante ; la raison en est que, pour comprendre les attaques d'Aristophane, il fallait séparer le problème en deux :\\
--- quand il évoque le σκάνδιξ, Aristophane moque le style d'Euripide, tout en finasseries (nous allons voir que le verbe διασκανδικίζειν a le même double sens que \italique{épicer} en français ; le sens propre : on épice une galette ; le sens imagé : on épice un discours) ;\\
--- quand il évoque la mère d'Euripide, la verdurière, Aristophane fait entendre que ce dernier est un \italique{fils de la terre}, \latin{terrae filius} ; dont les origines sont obscures et l'arbre généalogique impossible à établir ; bref : Euripide est un bâtard !


Cette note risque fort de trop tirer en longueur ; le dossier sur la mère d'Euripide et son σκάνδιξ est en effet assez épais. En voici le plan :\\
--- je montre d'abord que le vers 910 des \italique{Thesmophories} ne fait pas partie de ce dossier et que les `ἰφύων', \italique{lavandes}\marqueur{2}, n'ont jamais apparu dans le texte d'Aristophane ;\\
--- il y a par contre une autre herbe que le scandix à rétablir au vers 20 des \italique{Cavaliers} : l'ἀπόκυνον (Hésychius α 6416) ;\\
--- le σκάνδιξ et le verbe διασκανδικίζειν renvoient au style et à la façon de parler ;\\
--- descendre d'une verdurière, c'est être un \italique{fils de la terre}, \latin{terrae filius} : un bâtard ;\\
--- c'est la bâtardise d'Euripide qui explique sa rancœur vis-à-vis des femmes ; cela ressort de \italique{Thesm.} vv. 455-456 grâce à l'identification d'un proverbe (αὐτὸς τραφείς...) analogue à notre \italique{tel père, tel fils} ;\\
--- la note se conclue par une tentative pseudétymologique sur le σκάνδιξ pour expliquer les images qu'il a générées.



Pour le vers 910 \italique{Thesmophories}, je renvoie le lecteur au commentaire \italique{ad locum}. 
Il n'y a aucune allusion aux herbes d'Euripide, qui ont envahi bien mal à propos jusqu'aux sources les plus fiables de ce vers.

Par contre, il faut en rétablir une aux vers 19-20 des \italique{Cavaliers}. La pièce débute par la traditionnelle scène des deux esclaves -- les généraux Nicias et Démosthène en l'occurrence, généraux qui avaient préparé la victoire de Sphactérie (victoire éclatante par la capture de la garnison spartiate). 
Mais ces deux généraux s'étaient fait voler la vedette par le braillard Cléon débarqué avec un contingent sur place au dernier moment et qui – coup du sort – l'avait emporté\marqueur{3}. Ils ressassent leur amertume au début de la pièce en se renvoyant la balle pour trouver une idée de salut : que faire ? — Dis-le ! — Non, toi dis-le\marqueur{4}... Nicias s'interroge : comment tourner cela précieusement à la manière d'Euripide (κομψευριπικῶς). Démosthène rétorque alors (vers 19-20) :
\citComm{Δημοσθένης --- μὴ 'μοί γε, μὴ 'μοί, μὴ διασκανδικίσῃς·\\
ἀλλ᾽ εὑρέ τιν᾽ †ἀπόκινον† ἀπὸ τοῦ δεσπότου.\\
Démosthène --- Non pas, non pas, ne saupoudre pas de scandix !\\
mais va trouver quelque .......  qui le décroche du maître.}%
Le « maître », c'est Δῆμος Πυκνίτης, le Populo-citoyen-roi-de-l'Assemblée, mené par le bout du nez par son « chien » : l'infâme Cléon. 
(Tout au long des \italique{Cavaliers}, Cléon est renvoyé à son canin statut de \italique{Chien de Dêmos}, animal flatteur et cajoleur – sûrement par reprise ironique de quelque propos du véritable Cléon, qui se prenait pour le chien de garde de la cité. Voir aussi les \italique{Guêpes}, vv. 852-853, avec le procès aux deux chiens.)


Que le lecteur m'accorde pour l'instant que le verbe διασκανδικίζω signifie « assaisonner à l'aide de σκάνδιξ » et, dans un sens imagé, « agrémenter », « enjoliver [un discours de circonlocutions] » ; le σκάνδιξ jouant un rôle analogue à notre \italique{piment}, et les multiples sens du verbe « pimenter » (Nous l'allons montrer tout-à-l'heure.)

L'ἀπόκινος des manuscrits, défini par le scholiaste comme une danse licencieuse et qui fait entendre par son étymologie le verbe ἀπο-κινέω, \italique{se retirer}, n'a rien à faire ici\marqueur{5}.

Comme toujours devant les difficultés, nous savons que nous pouvons compter sur notre bon vieil et fidèle ami : Hésychius. Voici ce que grand érudit a répondu à nos questions (α 6416) :%
\citComm{ἀπόκυνον: μάζα μεμαγμένη φαρμάκῳ, πρὸς ἀναίρεσιν κυνῶν. ἢ εἶδος βοτάνης\\
apokynon : galette préparée avec du poison [$=$ l'herbe dont il est question juste après], pour la destruction des chiens. \italique{Item} : sorte d'herbe.}


Hésychius nous explique que l'ἀπόκυνον, comme le σκάνδιξ, est une \italique{herbe sauvage}... à cette différence que celui-ci sert de condiment aux pauvres\marqueur{6} quand celui-là nous débarrasse des nuisibles : chiens selon Hésychius, Cléon pour Aristophane. 
Je propose donc la lecture suivante pour les vers 19-20 des \italique{Cavaliers} :
\citComm{Δημοσθένης --- μὴ 'μοί γε, μὴ 'μοί, μὴ διασκανδικίσῃς·\\
ἀλλ᾽ εὑρέ τιν᾽ ἀπόκυνον ἀπὸ τοῦ δεσπότου.\\
(ἀπόκινον codd. omn. : \lecon{ἀπόκυνον} Bibi ap. Hesych.)\\
Démosthène --- Arrête-moi ça : ne condiments pas !\\
mais trouve-moi plutôt quelque mort-aux-chiens qui en débarrasse son maître.}%
Ce passage des \italique{Cavaliers} est intéressant pour la bonne interprétation du verbe διασκανδικίζειν : Démosthène l'utilise d'abord comme synonyme de parler précieusement à la manière d'Euripide (κομψευριπικῶς) ; on dirait en français « épicer un discours ».  Mais au vers 20, il le sous-entend avec son sens premier : « Va me trouver de l'ἀπόκυνον [pour m'en épicer une galette] qui nous débarrassera le maître [de ce chien de Cléon]. »

Le verbe διασκανδικίζω n'est pas un hapax d'Aristophane ; il se trouve aussi dans un fragment de Téléclide, un autre poète comique contemporain d'Euripide. Ce fragment\marqueur{7} nous a été conservé par Athénée (\italique{Deipnosophistes}, VI, 56d.) :
\citComm{Ξυγγενέσθαι διὰ χρόνου †λιπαρείτω† με\\
δρυπεπέσι μάζαις καὶ διασκανδικίσαι.}
Tout ce que l'on peut inférer du contexte chez Athénée, c'est que l'adjectif δρυπεπής, \italique{bien mûr}, lorsqu'il est substantivé comme ici, désigne une \italique{olive}. 
Le sens que nous avons donné à διασκανδικίζω, « assaisonner à l'aide de σκάνδιξ », rendrait bien mieux compte aussi de ce second passage que celui indiqué dans les dictionnaires : « se nourrir de σκάνδιξ »\marqueur{8}.
Je traduis donc : « Il m'importune (λιπαρεῖ ?) depuis longtemps pour ajouter des olives aux galettes ainsi que les agrémenter de scandix ».

Nous n'avons pas encore parlé d'une citation non identifiée à l'entrée Σκάνδιξ de la \italique{Souda} (σ 536) qui montre clairement que le scandix, en un sens imagé, désigne un \italique{euphémisme} :
\citComm{καὶ μή μοι τὰ τοῦ σκανδικοπώλου παρατραγῴδει, σεμνὸν τὸ τῆς Πειρήνης ὕδωρ ἀποκαλοῦντος,\\
« et ne m'use pas de l'emphase du vendeur de scandix qui appelle `sacrée' la fontaine de Pirène ».}%
La cible est clairement Euripide : la citation peut faire allusion au vers 69 de sa \italique{Médée} ou aux vers 205-206 des \italique{Troyennes}. Il est toujours délicat de proposer des hypothèses sur des fragments aussi brefs... proposons tout-de-même notre hypothèse qui vaut ce qu'elle vaut – c'est-à-dire pas grand-chose : cette fontaine située sur l'Acrocorinthe délivrait effectivement une eau pure et abondante\marqueur{9} ; ce n'est donc pas sur la qualité de l'eau qu'Euripide enrobe de scandix la vérité... peut-être sur ce qu'on y pratiquait ? La fontaine de Pirène se trouvait juste en contrebas du temple d'Aphrodite dont certaines de ses préposées se livraient à une prostitution « sacrée »\marqueur{10}. Imagine-t-on un tragique contemporain, si peu tragique qu'il nous parlerait du « \italique{sacré} bois de... Boulogne » (!?) ce serait un peu fort de scandix !.. Cet exemple nous dessine très clairement le sens imagé du scandix : Euripide \italique{dore la pilule} en saupoudrant de scandix des choses basses ; il nous fait passer des vessies pour des lanternes.

Voici pour le style d'Euripide. Passons à sa généalogie. L'on connaît le \italique{topos} comique : chaque fois qu'un poète voulait discréditer l'un de ses contemporains, il l'accusait de n'être pas de pure race (εὐγενής) : soit qu'il fut Athénien d'importation récente\marqueur{11}, soit qu'on puisse émettre des doutes sur sa lignée officielle. Euripide n'y échappe pas. Dans les \italique{Grenouille} (vv. 946-947), alors qu'il expose sa théorie dramatique, il est interrompu :%
\citComm{Εὐριπίδης --- εἶτ’ οὐκ ἐλήρουν ὅ τι τύχοιμ’ οὐδ’ ἐμπεσὼν ἔφυρον,\\
ἀλλ’ οὑξιὼν πρώτιστα μέν μοι τὸ γένος εἶπ’ ἂν εὐθὺς\\
τοῦ δράματος.\\
Διόνυσος --- \decalage{τοῦ }κρεῖττον γὰρ ἦν σοι νὴ Δί’ ἢ τὸ σαυτοῦ.\\
Euripide --- Je ne débitais pas toutes les niaiseries qui me venaient, et je ne brouillais pas tout, sans crier gare ;
mais le premier qui entrait en scène exposait de prime-saut la filiation du drame.\\
Dionysos --- Cela valait mieux pour toi, par Zeus, que d'exposer la tienne.}%
L'attaque n'est pas neuve : quelques vers plus tôt, Euripide avait provoqué Eschyle ; la réplique ne s'était pas fait attendre (vv. 836-843) :%
\citComm{Εὐριπίδης --- ἐγᾦδα τοῦτον καὶ διέσκεμμαι πάλαι,\\
ἄνθρωπον ἀγριοποιὸν αὐθαδόστομον,\\
ἔχοντ᾽ ἀχάλινον ἀκρατὲς ἀπύλωτον στόμα,\\
ἀπεριλάλητον κομποφακελορρήμονα.\\
Αἰσχύλος --- `ἄληθες ὦ παῖ τῆς ἀρουραίας θεοῦ;'\\
σὺ δή με ταῦτ᾽ ὦ στωμυλιοσυλλεκτάδη\\
καὶ πτωχοποιὲ καὶ ῥακιοσυρραπτάδη;\\
ἀλλ᾽ οὔ τι χαίρων αὔτ᾽ ἐρεῖς.\\
Euripide --- Je le connais et l'ai de longtemps percé à jour,\\
cet homme aux héros sauvages, au verbe présomptueux,\\
à la langue sans frein, sans retenue, sans barrières,\\
à l'esprit rétif, ce fagoteur de mots ronflants.\\
Eschyle --- « Ah, par exemple ! fils de la déesse agreste », \\
c'est ainsi que tu parles de moi, raccoleur de fadaises,\\
faiseur de mendiants, ravaudeur de haillons.\\
Mais tu n'auras pas lieu de t'en applaudir.}
(Le scholiaste nous indique que la réaction cinglante d'Eschyle au vers 840 est une parodie d'un vers d'Euripide avec ἀρουραίας, \italique{agreste}, en place de θαλασσίας, \italique{marine}.)



Or on désigne comme fils de la (déesse) terre dans la langue proverbiale grecque (et latine) les personnes dont la généalogie est obscure. 
Dans la pièce d'Euripide qui porte son nom, Ion interroge Xouthos sur l'identité de sa mère ; comme Xouthos ignore la réponse, Ion s'exclame alors (v. 542) :%
\citComm{Ἴων --- γῆς ἄρ᾽ ἐκπέφυκα μητρός;\\
Ion --- J'ai donc eu la terre pour mère ?}
Érasme a recensé l'adage (n° 786) latin \latin{terrae filius} employé avec le même sens par Tertullien, Cicéron, \&c.
C'est bien le propre de l'αὐτό-χθων que de ne pas savoir à quand remonte l'installation de ses ancêtres dans le pays : sa filiation se perd dans la nuit des temps.

L'on comprend maintenant où veut en venir Aristophane avec ses légumes et autres herbes, étroitement associés à la terre : si nul Français ne s'offusquera qu'on dise de lui qu'il est né dans un chou ou une rose, en grec cela revient à émettre des doutes sur la naissance.
%Si en français dire de quelqu'un qu'il est né dans une rose ou dans un chou n'aurait rien d'insultant ; en grec, être élevé au milieu des légumes et autres herbes (λάχανα), étroitement associés à la terre, c'est être de naissance obscure... 
Ainsi s'explique notre passage des \italique{Thesmophories} (vv. 455-456).
Le procès d'Euripide a été ouvert et la parole est à l'accusation. Après un premier réquisitoire-fleuve, une seconde oratrice s'avance et prend la parole (vers 443-458).
Le point culminant de l'attaque se trouve dans le distique des vers 455-456 :%
\citComm{ἄγρια γὰρ ἡμᾶς ὦ γυναῖκες δρᾷ κακά,\\
ἅτ᾽ ἐν ἀγρίοισι τοῖς λαχάνοις αὐτὸς τραφείς.}%
Aristophane joue sur la répétition d'ἄγριος, qualifiant les mauvais comportements qu'Euripide impute à l'engeance féminine dans ses pièces\marqueur{12} et, au vers suivant, les λάχανα, les \italique{légumes} qui, lorsqu'ils sont ἄγρια, « sauvages », désignent des \italique{herbes sauvages}, de celles qui poussent sans qu'on les cultive (par opposition à ἥμερα). 

Si Euripide nourrit une telle haine vis-à-vis des femmes et les présente comme dépravées, c'est que sa propre mère était telle... Sur le fond, ce genre de psychanalyse de comptoir n'est pas d'une grande finesse, mais le coup sous la ceinture porte auprès des femmes et suffit pour déchaîner l'enthousiasme du chœur quelques vers plus loin\marqueur{13}.

Ces éloges concernent davantage la forme que le fond de la diatribe\marqueur{14}.

Le tour `αὐτὸς τραφείς' est une locution proverbiale qu'on retrouve ailleurs : le parémiographe Michel Apostolius\marqueur{15} a relevé dans l'\italique{Anthologie} de Stobée (IV, 25, 6), au chapitre consacré au respect filial, cette citation du poète Dikaiogène\marqueur{16} :%
\citComm{Αὐτὸς τραφεὶς δὲ τῶν φυτευσάντων ὕπο\\
καλῶς τὸν αὐτὸν ἔρανον\marqueur{17} αὐτοῖσιν νέμεις\\
Toi-même élevé [litt. : nourri] par tes parents [litt. : ceux qui t'ont semé],\\
tu leur partages comme de juste ton propre repas.}
La locution proverbiale complète était-elle `αὐτὸς τραφεὶς ὑπὸ τῶν φυτευσάντων' ? On comprendrait où veut en venir l'oratrice : si les parents d'Euripide, ses \italique{semeurs}, l'ont semé au milieu des légumes, il n'est pas douteux que cet homme de la terre y ait récolté tout le fiel qu'il déverse sur les femmes dans ses tragédies.

Dion Cassius (\italique{Histoire romaine}, XLVI, 4) met dans la bouche d'un orateur la même figure de rhétorique, semblable par la forme (ἐν $+$ datif $+$ αὐτὸς τραφείς) et par le ton. C'est une réponse à Cicéron et ses violentes \italique{Philippiques} contre Marc Antoine. Cicéron... qui est-ce, au juste ?
\citComm{ὁ πατὴρ αὐτῷ ὁ κναφεύς [ ... ] ἄνθρωπος ἀγαπητῶς ἔκ τε τούτων καὶ ἐκ τῶν πλυνῶν διατρεφόμενος [ ... ] \textbf{Ἐν οἷς αὐτὸς τραφεὶς} οὐκ ἀπεικότως τοὺς κρείσσονας αὐτοῦ καὶ πατεῖ καὶ πλύνει.\\
\francais{Son père était foulon, lui qui avait grand-peine à vivre de ce métier et de son lavoir. À la suite d'une telle éducation, non sans invraisemblance, ce personnage foule aux pieds ceux qui valent mieux que lui et leur lave la tête.}}
Ces clefs d'interprétation permettent d'expliquer le vers 478 des \italique{Acharniens}, alors que Dikaiopolis emprunte à Euripide tout un fatras de matériel puisé dans ses propres drames, voilà qu'il déclenche le courroux d'Euripide en lui demandant :
\citComm{Δικαιόπολις --- σκάνδικά μοι δὸς μητρόθεν δεδεγμένος.\\
Εὐριπίδης --- ἁνὴρ ὑβρίζει· κλῇε πηκτὰ δωμάτων.\\
Dikaiopolis --- Donne-moi de ce scandix que tu tiens de ta mère.\\
Euripide --- Il m'insulte. « Ferme sur lui les ais solides de la porte. »}
Aristophane fait coup double : l'on comprend bien pourquoi Dikaiopolis veut emprunter à Euripide son scandix : il a besoin de la faconde du tragédien pour la joute verbale qui va l'opposer au chœur des vieillards d'Acharnes\marqueur{18}. 
On pourrait dire que Dikaiopolis compte les « enfariner » (διασκανδικίζειν) avec l'art oratoire d'Euripide... mais en précisant que ce scandix lui vient de sa mère, Dikaiopolis fait allusion au sobriquet de fils de verdurière qui colle au poète.
Euripide le comprend bien qui chasse l'insolent : ἁνὴρ ὑβρίζει· κλῇε πηκτὰ δωμάτων.\\

Cette trop longue note touchant enfin à sa fin, nous nous permettrons quelque aventureuse conjecture sur l'origine de cette utilisation imagée du σκάνδιξ : il nous semble que si le grec a choisi cette plante précisément et non une autre, c'est non pas pour ses vertus (réelles ou supposées), mais bien plutôt pour des raisons de sonorité. Cette intuition m'est venue de la consultation de l'\latin{Etymologicum Genuium}\marqueur{19} à l'entrée σκάνδαλον, la \italique{chausse-trappe}, pris ici en un sens imagé : σκάνδαλον: τὸ σκάζειν περὶ τὴν ἀλήθειαν. 
Image saisissante et évocatrice s'il en est ! « boîter autour de la vérité », « tourner autour du pot-aux-roses ». 
Elle s'appuie sur un rapprochement pseudo-étymologique\marqueur{20} entre le mot σκάνδ-αλον et le verbe σκάζειν, boîter... le σκάνδ-ιξ ne pouvait-il pas suggérer la même image : celle de \italique{tordre}, \italique{rendre boiteuse} ou de \italique{camoufler} la vérité en la saupoudrant d'une herbe de vil prix ?


Voilà qui donne quelques nouveaux éléments d'appréciation sur le problème d'Euripide, son σκάνδιξ et sa maman la λαχανοπωλήτρια.

}%
\debutNote{455-6}%
\note{1}{\italique{Thesm.} vv. 387, 456 ; \italique{Acharn.} v. 478 ; \italique{Grenouilles} v. 840.}%
\note{2}{Dans un article passionnant sur quelques plantes apparaissant chez Aristophane, mêlant avec un rare bonheur de savantes études de textes et des connaissances botaniques très pointues, la regrettée Suzanne Amigues identifiait en fait l'ἴφυον avec la nielle des blés (\latin{Agrostemma githago L.}). (« Quelques légumes de disette chez Aristophane et Plutarque » paru en 1988 dans le \italique{Journal des Savants}, pp. 157-171.)}%
\note{3}{Pour plus de détails, le lecteur pourra utilement se reporter au quart livre de Thucydide, chapitres 28 et suivants ainsi qu'à la longue tirade du Démosthène des \italique{Cavaliers}, à partir du vers 40.}%
\note{4}{Aristophane fait-il allusion à un complot qui aurait été tramé contre Cléon ? contre ce va-t-en-guerre de Cléon qui empêchait toute négociation de paix au moment-même où Athènes était en position de force pour se sortir favorablement de cette guerre fratricide qui n'avait que trop duré.}%
\note{5}{Il faut donc rejeter la jolie trouvaille du grand Alphonse Willems : le « branle de sortie »...}%
\note{6}{Le scandix semble être une épice peu appréciée : Andocide parle ainsi des privations des premières années de la guerre du Péloponnèse : μηδὲ ἄγρια λάχανα καὶ σκάνδικας ἔτι φάγοιμεν, « puissions-nous ne plus manger des herbes sauvages ni de scandix ».}%
\note{7}{Fragment 40(38) dans l'édition de Kassel-Austin, \italique{Poetae Comici Graeci}, vol. VII.}%
\note{8}{« manger du cerfeuil » Bailly ; « \espagnol{dar de comer perifollo} » DGE ; « \anglais{feed on chervil} » LSJ ; « Kerbel essen » Pape.}%
\note{9}{Cf. ce qu'en disent Pausanias, Strabon et Suidas à l'entrée Πειρήνη ainsi que cette citation d'Antipater prise dans l'\italique{Anthologie grecque} : Λαί̈δ' ἔχω πολιῆτιν ἐϋζώνοιο Κορίνθου, Πειρήνης λευκῶν φαιδροτέρην λιβάδων, « Je tiens Laïs, citoyenne de Corinthe, bien ceinturée, plus brillante que les gouttes blanches de Pirène ».}%
\note{10}{La réputation des Corinthiennes était telle que les Grecs en ont inventé le verbe κορινθιάζειν.}%
\note{11}{Voir par exemple \italique{Grenouilles} vv. 717-733 qui est sans ambiguïté. L'on s'étonne de trouver dans les modernes traductions des \italique{Thesmophories} (v. 330) εὐγενεῖς γυναῖκες rendu par `nobles femmes', comme si la possibilité de célébrer les Thesmophores était conditionnée à une forme quelconque de noblesse. Il serait bon avant d'aborder Aristophane de se défaire de nos préjugés modernes, faute de quoi l'on s'expose à de tels contre-sens. Lire là-dessus l'indispensable article d'Alphonse Willems « Aristophane et la démocratie athénienne » reproduit en annexe et qu'on peut trouver dans son \italique{Aristophane} posthume (1919), t. III, pp. 445-473.}%
\note{12}{Aristophane, au vers 349 des \italique{Nuées}, qualifie ainsi d' ἄγριος le chevelu Hiéronyme, un poète dithyrambique, fils de Xénophane qui passait pour pédéraste ; et le scholiaste de noter qu'on appelait ἀγρίους (et débauchés) les pédérastes : ἀγρίους (καὶ κολλοποδιώκτας) ἐκάλουν τοὺς παιδεραστάς. De même, à l'entrée Ἀγρίους de la \italique{Souda}, on peut lire : Αἰσχίνης ἐν τῷ κατὰ Τιμάρχου τοὺς σφόδρα ἐπτοημένους περὶ τὰ παιδικὰ οὕτως ὠνόμασε. καὶ Μένανδρος δὲ ἄγριον εἶπε κυβευτὴν τὸν λίαν περὶ τὸ κυβεύειν ἐσπουδακότα, « Eschine dans son \italique{Contre Timarque} nomme ainsi les hommes en proie à une excessive attirance pour les enfants. Et Ménandre dit qu'un joueur est ἄγριον, celui qui est trop porté au jeu ». L'on voit que le terme ἄγριος, au sens moral, désigne un \italique{intempérant}.}%
\note{13}{κομψότερον ἔτ᾽ ἢ τὸ πρότερον, « plus intelligent encore que le premier [discours] » [...] οἷα κατεστωμύλατο, « comme elle a bien parlé » [...] φρένας ἔχουσα, « elle a de l'esprit » \&c.}%
\note{14}{L'oratrice précédente, au vers 387, avait déjà traité Euripide de fils de la λαχανοπωλήτρια. Plutarque citait ces deux vers comme exemple d'un jeu de mots qu'il goûtait fort peu (voir ses \italique{Œuvres morales}, 853C, \italique{Abrégé de la comparaison d'Aristophane et de Ménandre}). Je ne suis pas sûr que cet inénarrable raseur, le Cicéron des lettres grecques, ait bien saisi intégralement le détournement d'Aristophane.}%
\note{15}{\italique{Corpus paroemiographorum graecorum}, II, ed. Leutsch, Apostolii centuria V 30a, p. 341.}%
\note{16}{L'on ne possède pas grand-chose sur ce poète grec, auteur de tragédies et dithyrambes, qui vécut probablement au \nieme{IV} siècle av. J.-C. Tout ce que l'on sait de lui est regroupé dans les \italique{Tragicorum Graecorum Fragmenta}, vol. I, ed. Snell \& Kannicht. Il y est classé en tant que « poète mineur » n° 52.}%
\note{17}{L'on trouve chez Démosthène dans son \italique{Contre Neæra}, 8, l'expression ironique τὸν αὐτὸν ἔρανον ἀποδοῦναι, \italique{payer quelqu'un de la même monnaie}.}%
\note{18}{C'est la raison de sa visite chez Euripide (vers 416-417) : δεῖ γάρ με λέξαι τῷ χορῷ ῥῆσιν μακράν· /
αὕτη δὲ θάνατον, ἢν κακῶς λέξω, φέρει. 
« Car j'ai à débiter au chœur une longue tirade. /
Il m'en coûte la vie, si je m'en tire mal. »}%
\note{19}{Lexique du \nieme{IX} siècle souvent associé à d'autres, postérieurs mais plus complets : l'\latin{Etymologicum magnum} et l'\latin{Etymologicum Symeonis}.}%
\note{20}{Je me vois forcé de préciser \italique{pseudo}-étymologique, car mes velléités de rapprocher étymologiquement σκάνδαλον de σκάζειν furent foudroyées par Alexandre Perret, dont j'eus la chance d'assister à des cours de linguistique, cours qui ressemblaient davantage à d'enivrants voyages à travers les contrées et les parlers d'Europe (et ce jusqu'aux confins de l'Asie). Voici ce que m'a très aimablement répondu M. Perret : « Après recherches fructueuses, je pense avoir trouvé une réponse dans le \italique{Lexikon der Indogermanischen Verben} (LIV) mais il ne semble pas que les deux racines soient liées. Σκανδάλον semble en effet rattaché à un thème *skend- qui signifierait « bondir » (on a skándati en sanskrit pour « sauter »), mais le suffixe -αλον me dérange dans la mesure où ça semble pré-grec. Aucune étymologie n'est assez bien solide pour ce mot qui peut être un emprunt. Quelle motivation pour passer de « bondir » à « faire tomber » ? À la limite, un autre thème *sk'end- (avec un k- palatal donc invisible en grec mais pas en sanskrit) avec le sens parfois de « faire tomber » (skr. Chándayati: « tomber »), est en soit beaucoup plus plausible !! Mais pour σκάζω, on rapproche ça d'une racine *skeng- avec effectivement le sens de « boiteux, tors » (All. \italique{Hinken}). Ce qui m'étonne est qu'Hesychius recense à la fois σκάνδαλον: \italique{piège à souris} et σκάνδαλος: \italique{obstacle}, mais est-ce que les pièges à souris étaient des trappes où tomber, ou bien des claquoirs comme à notre époque ? Ça mérite investigation. En soit passer de *skend- à σκάζω avec un suffixe -yo de présent n'est pas choquant. »}%
%\note{1}{Pour ce genre de situation, le rire du soldat allemand dans la \italique{Traversée de Paris} est, ce me semble, la réaction la plus saine qui soit.}%
%\note{2}{Aristophane, au vers 349 des \italique{Nuées}, qualifie ainsi d' ἄγριος le chevelu Hiéronyme, un poète dithyrambique, fils de Xénophane qui passait pour pédéraste ; et le scholiaste de noter qu'on appelait ἀγρίους (et débauchés) les pédérastes : ἀγρίους (καὶ κολλοποδιώκτας) ἐκάλουν τοὺς παιδεραστάς. De même, à l'entrée Ἀγρίους de la \italique{Souda}, on peut lire : Αἰσχίνης ἐν τῷ κατὰ Τιμάρχου τοὺς σφόδρα ἐπτοημένους περὶ τὰ παιδικὰ οὕτως ὠνόμασε. καὶ Μένανδρος δὲ ἄγριον εἶπε κυβευτὴν τὸν λίαν περὶ τὸ κυβεύειν ἐσπουδακότα, « Eschine dans son \italique{Contre Timarque} nomme ainsi les hommes en proie à une excessive attirance pour les enfants. Et Ménandre dit qu'un joueur est ἄγριον, celui qui est trop porté au jeu ». L'on voit que le terme ἄγριος, au sens moral, désigne un \italique{intempérant}.}%
%\note{3}{L'oratrice précédente, au vers 387, avait déjà traité Euripide de fils de la λαχανοπωλήτρια.}%
%\note{4}{En rester là donnerait raison au grave Plutarque qui citait ces deux vers comme exemple d'un jeu de mots qu'il goûtait fort peu (voir ses \italique{Œuvres morales}, 853C, \italique{Abrégé de la comparaison d'Aristophane et de Ménandre}). Je ne suis pas sûr que cet inénarrable raseur, le Cicéron des lettres grecques, ait bien saisi intégralement le détournement d'Aristophane que nous allons voir.}%
%\note{5}{\italique{Corpus paroemiographorum graecorum}, II, ed. Leutsch, Apostolii centuria V 30a, p. 341.}%
%\note{6}{L'on ne possède pas grand-chose sur ce poète grec, auteur de tragédies et dithyrambes, qui vécut probablement au \nieme{IV} siècle av. J.-C. Tout ce que l'on sait de lui est regroupé dans les \italique{Tragicorum Graecorum Fragmenta}, vol. I, ed. Snell \& Kannicht. Il y est classé en tant que « poète mineur » n° 52.}%
%\note{7}{Stobée, livre IV, chapitre 25, 6.}%
%\note{8}{L'on trouve chez Démosthène dans son \italique{Contre Neæra}, 8, l'expression ironique τὸν αὐτὸν ἔρανον ἀποδοῦναι, \italique{payer quelqu'un de la même monnaie}.}%
%\note{9}{Voir la note sur le vers 910 et ces étranges plantes, ἰφύων, des \italique{nielles des blés} (Agrostemma githago), qui ont envahi bien mal à propos jusqu'aux sources les plus fiables.}%
%\note{10}{Pour un article passionnant sur quelques plantes apparaissant chez Aristophane, passionnant car il mêle avec un rare bonheur de savantes études de textes et des connaissances botaniques fort appréciables, le lecteur se reportera à l'article de Suzanne Amigues : « Quelques légumes de disette chez Aristophane et Plutarque » paru en 1988 dans le \italique{Journal des Savants} (pp. 157-171). Passionnant quand bien-même je suis obligé de m'inscrire en faux contre l'interprétation qu'elle donne du vers 910 des \italique{Thesmophories} (nous en reparlerons).}%
%\note{11}{Pour plus de détails, le lecteur pourra utilement se reporter au quart livre de Thucydide, chapitres 28 et suivants ainsi qu'à la longue tirade du Démosthène des \italique{Cavaliers}, à partir du vers 40.}%
%\note{12}{Aristophane fait-il allusion à un complot qui aurait été tramé contre Cléon ? contre ce va-t-en-guerre de Cléon qui empêchait toute négociation de paix au moment-même où Athènes était en position de force pour se sortir favorablement de cette guerre fratricide qui n'avait que trop duré.}%
%\note{13}{Le \italique{maître}, c'est Δῆμος Πυκνίτης, le Populo-citoyen-roi-de-l'Assemblée, mené par le bout du nez par la pire espèce rustique d'embabouineur qui soit : le Marchand-de-cuir Cléon.}%
%\note{14}{Bailly. Idem pour le Liddell-Scott : « feed on chervil (σκάνδιξ) ». Quant au Diccionario Griego-Español, il donne : « dar de comer perifollo ».}%
%\note{15}{Fragment 40(38) dans l'édition de Kassel-Austin, \italique{Poetae Comici Graeci}, vol. VII.}%
%\note{16}{\italique{Deipnosophistes}, VI, 56d.}%
%\note{17}{Je dois au grand Alponse Willems cette jolie trouvaille qui évoque tout en un : le « branle de sortie »... trouvaille qui porte malheureusement sur une transmission fautive du texte comme nous l'allons voir tout à l'heure.}%
%\note{18}{Et je n'ai encore jamais trouvé chez Aristophane d'analogue à notre éculée « sans transition ».}%
%\note{19}{Le scandix semble être une épice peu appréciée : Andocide parle ainsi des privations des premières années de la guerre du Péloponnèse : μηδὲ ἄγρια λάχανα καὶ σκάνδικας ἔτι φάγοιμεν, « puissions-nous ne plus manger des herbes sauvages ni de scandix ».}%
%\note{20}{J'accorderais ainsi plus de crédit aux biographies d'Euripide qui le présentent comme ayant eu des parents riches que celles qui en font – peut-être d'après des interprétations trop littérales des Comiques – un \italique{self-made-man} né dans l'indigence.}%
%\note{21}{Voir les \italique{Testimonia} des \italique{Vitae Euripidis} réunis par Kannicht au cinquième volume, consacré à Euripide, des \italique{Tragicorum Graecorum Fragmenta}.}%
%\note{22}{Pour le sens imagé de κάπηλος, voir par exemple ce qu'en dit Hésychius : κίβδης· κακοῦργος, κάπηλος, « frelateur : malfaiteur, \italique{boutiquier} » ; on trouve encore ceci : Φοινικελίκτην καὶ λόγων ἀλαζόνα: ἀπατηλὸν καὶ κάπηλον, « Phénichien et fort en gueule : un fourbe et un \italique{boutiquier} » ; les Phéniciens étant un peuple qui se livrait au trafic, la langue populaire leur rendait justice. Ainsi pour le français ; en ouvrant par exemple \italique{Le livre des proverbes français}, proverbes compulsés par Le Roux de Lincy, et allant visiter la ville chère à Rousseau \& Calvin, croquée comme il se doit par le blason populaire, on trouve : « Trois Juifs font un Balois, trois Balois font un Genevois » ou encore : « Les Genevois ont vertu de cent lieues de loing ». Chez Aristophane, ce sont plus particulièrement les débitants de vin qui sont accusés de gruger les mesures légales.}%
%\note{23}{Nous n'avons pas encore parlé d'une citation non identifiée à l'entrée Σκάνδιξ de la \italique{Souda} : καὶ μή μοι τὰ τοῦ σκανδικοπώλου παρατραγῴδει, σεμνὸν τὸ τῆς Πειρήνης ὕδωρ ἀποκαλοῦντος, « et ne m'use pas de l'emphase du vendeur de scandix qui appelle `sacrée' la fontaine de Pirène ». La cible est clairement Euripide : la citation peut faire allusion au vers 69 de sa \italique{Médée} ou aux vers 205-206 des \italique{Troyennes}. Il est toujours délicat de proposer des hypothèses sur des fragments aussi brefs... proposons tout-de-même notre hypothèse qui vaut ce qu'elle vaut – c'est-à-dire pas grand-chose : cette fontaine située sur l'Acrocorinthe délivrait effectivement une eau pure et abondante (cf. ce qu'en disent Pausanias, Strabon et Suidas à l'entrée Πειρήνη ainsi que cette citation d'Antipater prise dans l'\italique{Anthologie grecque} : Λαί̈δ' ἔχω πολιῆτιν ἐϋζώνοιο Κορίνθου, Πειρήνης λευκῶν φαιδροτέρην λιβάδων, « Je tiens Laïs, citoyenne de Corinthe, bien ceinturée, plus brillante que les gouttes blanches de Pirène ») ; ce n'est donc pas sur la qualité de l'eau qu'Euripide enrobe de scandix la vérité... peut-être sur ce qu'on y pratiquait ? La fontaine de Pirène se trouvait juste en contrebas du temple d'Aphrodite dont certaines de ses préposées se livraient à une prostitution « sacrée » (la réputation des Corinthiennes était telle que les Grecs en ont inventé le verbe κορινθιάζειν).}%
%\note{24}{« Entre l'artisan et le politique », in : \italique{Dieux, héros et médecins grecs} (2001), pp. 65-96.}%
%\note{25}{Tzetzes (dans une scholie citée par Fornaro) parle, à propos d'Euripide, de κάπηλα γράμματα et oppose ce poète κάπηλον et ἀγοραῖον à Eschyle. De même Lucien (\italique{De Historia conscribenda}, 44) emploie ces deux adjectifs pour caractériser un vocabulaire vulgaire (ὀνόματα ἀγοραῖα καὶ καπηλικά).}%
%\note{26}{Quintilien \italique{Inst. orat.} XII, 10, 58.}%
%\note{27}{\italique{Acharniens}, vers 469.}%
%\note{28}{\italique{Grenouilles}, vers 941-942.}%
%\note{29}{Voir Dioscoride, \italique{De Materia medica Libri V}, ed. M. Wellmann, Berlin 1907, vol. I, livre II 123, 138.}%
%\note{30}{Cf. lorsque Ion demande à Xouthos (\italique{Ion}, v. 542) :%
%\citComm{Ἴων -- γῆς ἄρ᾽ ἐκπέφυκα μητρός;\\
%J'ai donc eu la terre pour mère ?}%
%L'adage fut recensé par Érasme : (n° 786, \latin{terrae filius}).
%Est-ce là qu'Eschyle veut en venir dans les \italique{Grenouilles} (v. 840), lorsqu'il lance à Euripide la cinglante réponse : ἄληθες ὦ παῖ τῆς ἀρουραίας θεοῦ.}%
%\note{31}{Lexique du \nieme{IX} siècle souvent associé à d'autres, postérieurs mais plus complets : l'\latin{Etymologicum magnum} et l'\latin{Etymologicum Symeonis}.}%
%\note{32}{Je me vois forcé de préciser \italique{pseudo}-étymologique, car mes velléités de rapprocher étymologiquement σκάνδαλον de σκάζειν furent foudroyées par Alexandre Perret, dont j'eus la chance d'assister à des cours de linguistique, cours qui ressemblaient davantage à d'enivrants voyages à travers les contrées et les parlers d'Europe (et ce jusqu'aux confins de l'Asie). Voici ce que m'a très aimablement répondu M. Perret : « Après recherches fructueuses, je pense avoir trouvé une réponse dans le \italique{Lexikon der Indogermanischen Verben} (LIV) mais il ne semble pas que les deux racines soient liées. Σκανδάλον semble en effet rattaché à un thème *skend- qui signifierait « bondir » (on a skándati en sanskrit pour « sauter »), mais le suffixe -αλον me dérange dans la mesure où ça semble pré-grec. Aucune étymologie n'est assez bien solide pour ce mot qui peut être un emprunt. Quelle motivation pour passer de « bondir » à « faire tomber » ? À la limite, un autre thème *sk'end- (avec un k- palatal donc invisible en grec mais pas en sanskrit) avec le sens parfois de « faire tomber » (skr. Chándayati: « tomber »), est en soit beaucoup plus plausible !! Mais pour σκάζω, on rapproche ça d'une racine *skeng- avec effectivement le sens de « boiteux, tors » (All. \italique{Hinken}). Ce qui m'étonne est qu'Hesychius recense à la fois σκάνδαλον: \italique{piège à souris} et σκάνδαλος: \italique{obstacle}, mais est-ce que les pièges à souris étaient des trappes où tomber, ou bien des claquoirs comme à notre époque ? Ça mérite investigation. En soit passer de *skend- à σκάζω avec un suffixe -yo de présent n'est pas choquant. »}
\\
ἀλλ᾽ εἰς ἀγορὰν ἄπειμι· δεῖ γὰρ ἀνδράσιν%\testimonia{457-8}{(δεῖ γὰρ -- εἴκοσιν) Suidas σ 1591 (sed ἀνδράσι et εἴκοσι pro ἀνδράσιν et εἴκοσιν)}
\\
πλέξαι στεφάνους συνθηματιαίους εἴκοσιν.%\testimonia{458}{Athenaeus XV 680c \espTest (στεφάνους συνθηματιαίους) Pollux VII 200 \espTest (συνθηματιαίους) Hesychius σ 2624 συνθηματιαῖοι}
\commentaire{458}{στεφάνους συνθηματιαίους.}{Athénée (XV, 680c), dans le long chapitre consacré aux couronnes de ses \italique{Deipnosophistes}, parle de ces couronnes \italique{synthématiées} :%
\citComm{συνθηματιαῖοι στέφανοι, ἠργολαβημένοι καὶ ἐκδόσιμοι. Ἀριστοφάνης ἐν Θεσμοφοριαζούσαις.}%
Glose identique chez Hésychius (σ 2624) :%
\citComm{συνθηματιαῖοι: οἱ ἠργολαβημένοι καὶ ἐκδεδομένοι, ἐκδόσιμοι.}}\\\passeLigne}%
\end{groupe}%
%% v. 459-465
\begin{groupe}%
%
\personnageGrec{Χορός}%
\typeVers{$2$ toch.}\chantChoeur{ἕτερον αὖ τι λῆμα τοῦτο%\testimonia{459-60}{(ἕτερον -- ἀναπέφηνεν) Suidas λ 441 (ἀναπέφηνε pro ἀναπέφηνεν)}
}%
\end{groupe}%
%% v. 460a
\begin{groupe}%
\decoupeVers{a}\typeVers{$2$ troch.}\chantChoeur{κομψότερον ἔτ᾽ ἢ τὸ πρότερον}%
\end{groupe}%
%% v. 460b
\begin{groupe}%
\ligne{460}\decoupeVers{b}\typeVers{troch.}\chantChoeur{ἀναπέφηνεν.}%
\end{groupe}%
%% v. 461
\begin{groupe}%
\typeVers{choria. $+$ ia.}\chantChoeur{οἷα κατεστωμύλατο%\apparat{461}{\lecon{κατεστωμύλατο} R : κἀστωμύλατο Dobree}
}%
\end{groupe}%
%% v. 462
\begin{groupe}%
\typeVers{$2$ troch.}\chantChoeur{οὐκ ἄκαιρα, φρένας ἔχουσα\apparat{462}{ἄκερα R : \lecon{ἄκαιρα} Zanetti \ligneApp{463} \lecon{πολύπλοκον} R : τι πολύπλοκον Dobree : πολύστροφον Hermann}}%
\end{groupe}%
%% v. 463
\begin{groupe}%
\typeVers{$2$ troch.}\chantChoeur{καὶ πολύπλοκον νόημ᾽, οὐδ᾽%\apparat{463}{\lecon{πολύπλοκον} R : τι πολύπλοκον Dobree : πολύστροφον Hermann}%
\commentaire{463}{πολύπλοκον.}{Les diverses corrections qui ont été proposées pour faire rentrer de force ce vers dans le mètre trochaïque ne me paraissent pas indispensables : on peut simplement considérer qu'il n'y a pas de \italique{correptio attica} et scander :%
\citComm{καὶ πολύπ|λοκον νόημ᾽, οὐδ᾽\\
$ - \cup - \cup - \cup - - $}%
}%
}%
\end{groupe}%
%% v. 464
\begin{groupe}%
\typeVers{$2$ troch.}\chantChoeur{ἀσύνετ᾽ ἀλλὰ πιθανὰ πάντα.}%
\end{groupe}%
%% v. 465a
\begin{groupe}%
\decoupeVers{a}\typeVers{$2$ troch.}\chantChoeur{δεῖ δὲ ταύτης τῆς ὕβρεως ἡ-}%
\end{groupe}%
%% v. 465b
\begin{groupe}%
\ligne{465}\decoupeVers{b}\typeVers{troch.}\chantChoeur{μῖν τὸν ἄνδρα}%
\end{groupe}%
%% v. 465c
\begin{groupe}%
\ligne{465}\decoupeVers{c}\typeVers{$2$ troch\cat}\chantChoeur{περιφανῶς δοῦναι δίκην.\\\passeLigne}%
\end{groupe}%
\endnumbering
\end{Leftside}
\begin{Rightside}
\beginnumbering
%% v. 423+ (part 2)
\begin{groupe}
\trimIambFr{\ligne{423}je ne sais quelles laconiennes à trois dents.\testimonia{423}{(γομφίους) \Phryn 61.1}
Autrefois du moins il tenait à nous d'ouvrir furtivement la porte,
en nous faisant graver un anneau de trois oboles.
Aujourd'hui ce tâte-poule d'Euripide\testimonia{426}{(οἰκότριψ $=$ οἰκοτραφής) Moeris ο 25 ; (ᾡκότριψ) Hesychius ο 265-266 ; \AelDion ο 8-9}
leur a enseigné l'emploi de cachets vermoulus
qu'ils portent suspendus sur eux. Maintenant donc j'opine\testimonia{428-9}{(ἐμοὶ -- ἁμωσγέπως) Suidas κ 2771 (τινα pro τιν᾽ ἡμᾶς)}
que nous combinions de manière ou d'autre un procédé pour le détruire,\testimonia{429}{(κυρκανᾶν) Photius κ 1254 κυρκανᾶν $=$ Hesychius κ 4691 κυρκανᾷ ; \Eustathe 1213.50 \espTest (ἁμωσγέπως) Hesychius α 4182 ἁμωσγέπως ; \Syn α 449 $=$ Suidas α 1645}
soit par des drogues, soit par un moyen quelconque,
afin qu'il meure. Ceci, je le dis à pleine bouche.
Le reste, je le coucherai par écrit, avec l'aide de la greffière.\testimonia{432}{(μετὰ τῆς γραμματέως συγγράψομαι) \Eustathe 1484.53 \espTest (τῆς γραμματέως) Pollux IV 19}\\\passeLigne}%
\end{groupe}%
%% v. 433-442
\begin{groupe}%
\personnageFr{Chœur}%
\chantChoeurFr{Jamais je n'entendis}%
\end{groupe}%
%% v. 433-442
\begin{groupe}%
\chantChoeurFr{femme plus féconde en retours}%
\end{groupe}%
%% v. 433-442
\begin{groupe}%
\chantChoeurFr{ni qui sache mieux parler.}%
\end{groupe}%
%% v. 433-442
\begin{groupe}%
\chantChoeurFr{Tout ce qu'elle dit est juste.}%
\end{groupe}%
%% v. 433-442
\begin{groupe}%
\chantChoeurFr{Elle a tout envisagé sous tous égards,\testimonia{437-9}{(πάσας -- ἀνηῦρεν) Suidas ε 22 πάσας δ' ἰδέας ἐξήτασε πάντα δ' ἐβάστασε φρενί, πυκνῶς τε ποικίλως λόγους ἀνεῦρεν}}%
\end{groupe}%
%% v. 433-442
\begin{groupe}%
\chantChoeurFr{tout soupesé dans son esprit : et avec discrétion,\testimonia{438}{(πυκνῶς) Hesychius π 4335 πυκνῶς}}%
\end{groupe}%
%% v. 433-442
\begin{groupe}%
\chantChoeurFr{elle a trouvé un heureux}%
\end{groupe}%
%% v. 433-442
\begin{groupe}%
\chantChoeurFr{choix d'adroites raisons.}%
\end{groupe}%
%% v. 433-442
\begin{groupe}%
\chantChoeurFr{À tel point que si Xénoclès,}%
\end{groupe}%
%% v. 433-442
\begin{groupe}%
\chantChoeurFr{fils de Carkinos, parlait tout de suite après elle,}%
\end{groupe}%
%% v. 433-442
\begin{groupe}%
\chantChoeurFr{il vous paraîtrait je pense,}%
\end{groupe}%
%% v. 433-442
\begin{groupe}%
\ligne{443}\chantChoeurFr{à tous}%
\end{groupe}%
%% v. 433-442
\begin{groupe}%
\chantChoeurFr{parler sans rien dire.\testimonia{442}{(ἄντικρυς) Photius α 2210 ἄντικρυς ; \AelDion α 107}\\\passeLigne}%
\end{groupe}%
%% v. 443-458
\begin{groupe}%
\personnageFr{Marchande de couronnes}%
\trimIambFr{\ligne{443}Ce n'est que pour ajouter quelques mots que je m'avance \fin à mon tour.\\
Les autres chefs d'accusation ont été traités à souhait.
Mais ce que j'ai souffert personnellement, je tiens à le dire.
Mon mari mourut alors que j'étais enceinte,\testimonia{446}{Hesychius ε 239 ἔγκυος· ἐγγαστρωμένη γυνή}
me laissant cinq petits enfants, que je nourrissais à grand'peine\testimonia{447-9}{(μόλις -- τέως μὲν οὖν ἀλλ᾽) scholium\expo{T} Plato \italique{Hipparch.} 229d (μυρίναις pro μυρρίναις et μέν οὐκ ἄν pro μὲν οὖν ἀλλ᾽)}
en tressant des couronnes sur le marché aux myrtes.\testimonia{448}{(μυρρίνη $=$ μυρσίνη) Moeris μ 23}
En attendant, je gagnais ma vie vaille que vaille.\testimonia{449}{Suidas τ 321 \espTest (ἡμικάκως) Pollux VI 161}
Mais aujourd'hui cet homme, avec ses tragédies,
a persuadé aux gens qu'il n'y a pas de dieux ;}
\end{groupe}%
%% v. 443-458
\begin{groupe}%
\trimIambFr{de sorte que notre commerce est réduit environ de moitié.
C'est pourquoi je vous recommande à toutes et vous requiers
de le punir, cet homme, pour quantité de raisons.\testimonia{453-6}{Aulus Gellius XV 20.7}
Car il se porte contre nous, ô femmes, à de sauvages violences,
élevé qu'il fut parmi les herbes sauvages.\testimonia{455-6}{Plut. \italique{Moral.}, \italique{Compar. Arist. et Menandr.} (853c)}
Mais je m'en vas à l'agora.\testimonia{457-8}{(δεῖ γὰρ -- εἴκοσιν) Suidas σ 1591 (sed ἀνδράσι et εἴκοσι pro ἀνδράσιν et εἴκοσιν)}\\
Car j'ai vingt couronnes à tresser à prix convenu.\testimonia{458}{Athenaeus XV 680c \espTest (στεφάνους συνθηματιαίους) Pollux VII 200 \espTest (συνθηματιαίους) Hesychius σ 2624 συνθηματιαῖοι}\\\passeLigne}%
\end{groupe}%
%% v. 459-465
\begin{groupe}%
\personnageFr{Chœur}%
\chantChoeurFr{Cette autre a montré un tour d'esprit\testimonia{459-60}{(ἕτερον -- ἀναπέφηνεν) Suidas λ 441 (ἀναπέφηνε pro ἀναπέφηνεν)}}%
\end{groupe}%
%% v. 459-465
\begin{groupe}%
\chantChoeurFr{plus piquant encore}%
\end{groupe}%
%% v. 459-465
\begin{groupe}%
\chantChoeurFr{que la première.}%
\end{groupe}%
%% v. 459-465
\begin{groupe}%
\chantChoeurFr{Comme elle a débité,}%
\end{groupe}%
%% v. 459-465
\begin{groupe}%
\chantChoeurFr{non du hors-d'oœvre — elle a du sens}%
\end{groupe}%
%% v. 459-465
\begin{groupe}%
\chantChoeurFr{et un esprit délié — ni de l'abstrus, }%
\end{groupe}%
%% v. 459-465
\begin{groupe}%
\chantChoeurFr{mais toutes choses emportant conviction !}%
\end{groupe}%
%% v. 459-465
\begin{groupe}%
\chantChoeurFr{Oui, il faut que de cet outrage}%
\end{groupe}%
%% v. 459-465
\begin{groupe}%
\chantChoeurFr{l'auteur}%
\end{groupe}%
%% v. 459-465
\begin{groupe}%
\chantChoeurFr{subisse un châtiment exemplaire.\\\passeLigne}%
\end{groupe}%
\endnumbering
\end{Rightside}
\end{pages}
\Pages
\clearpage\ifodd\value{page}\hbox{}\newpage\fi
\enTeteG{ΕΠΕΙΣΟΔΙΟΝ Α' --- Codex Ravennas 429, f. 168r (Classense 170r) --- vv. 466–508}
\begin{adjustwidth}{-1.5cm}{-3.0cm}
\begin{adjustbox}{max totalheight=21.0cm, max width=16.5cm, keepaspectratio}
\begin{minipage}{16.5cm}
{\fontsize{7.2pt}{8.5pt}\selectfont
\noindent\textbf{Folio 12 --- Codex Ravennas 429, \texttt{f. 168r} (\texttt{Classense 170r}) : ΕΠΕΙΣΟΔΙΟΝ Α', vv. 466–508}\hfill \textit{Témoins :} $\text{R}^1, \Rpc, \Rdeux$ (\Mon)\par\vspace{2pt}
\hrule height 0.4pt\vspace{2.5pt}
}
{\fontsize{8.1pt}{9.5pt}\selectfont
\ligneRav{466}{μ\auDessus{ν}{Η}}{τὸ μὲν ὦ γυναῖκες ὀξυθυμεῖσθαι σφόδρα}{}
\ligneRav{467}{}{Εὐριπίδῃ, τοιαῦτ᾽ \difference{ἀκουούσας} κακά,}{\textbf{467} ἀκουούσαις R : \lecon{ἀκουούσας} Zanetti}
\ligneRav{}{}{οὐ θαυμάσιόν ἐστ᾽, οὐδ᾽ ἐπιζεῖν τὴν χολήν·}{}
\ligneRav{469}{}{καὐτὴ γὰρ ἔγωγ᾽, οὕτως \difference{ὀναίμην} τῶν τέκνων,}{\textbf{469} ὀνοίμην R : \lecon{ὀναίμην} Suid. ο 330}
\ligneRav{470}{}{μισῶ τὸν ἄνδρ᾽ ἐκεῖνον, εἰ μὴ μαίνομαι·}{}
\ligneRav{471}{}{ὅμως δ᾽ ἐν \difference{ἀλλήλαισι} χρὴ δοῦναι λόγον·}{\textbf{471} ἀλλήλοισι R : \lecon{ἀλλήλαισι} Brunck}
\ligneRav{472}{}{αὐταὶ γάρ ἐσμεν, κοὐδεμί᾽ ἔκφορος λόγου·}{\textbf{472} \lecon{ἔκφορος} R : ἐκφορὰ Valckenaer ad Eur. \italique{Hipp.} 294}
\ligneRav{473}{}{τί ταῦτ᾽ ἔχουσαι \difference{᾽κεῖνον} αἰτιώμεθα}{\textbf{473} κεῖνον R, Suid τ 683 : \lecon{᾽κεῖνον} Meineke}
\ligneRav{474}{}{βαρέως τε φέρομεν, \difference{εἰ δύ᾽} ἡμῶν ἢ τρία}{\textbf{474} ἢ δύ᾽ R : \lecon{εἰ δύ᾽} Grynaeus\espApp ἢ δύ᾽ R : \lecon{εἰ δύ᾽} Grynaeus (\italique{si} Divus)}
\ligneRav{475}{}{κακὰ ξυνειδὼς εἶπε δρώσας μυρία;}{}
\ligneRav{476}{}{ἐγὼ γὰρ αὐτὴ πρῶτον, ἵνα \difference{μἄλλην} λέγω,}{\textbf{476} μὴ ἄλλην R : \lecon{μἄλλην} Dindorf}
\ligneRav{477}{}{ξύνοιδ᾽ ἐμαυτῇ \difference{πολλά γ᾽· ἀλλ᾽} ἐκεῖνο δ᾽ οὖν}{\textbf{477} πολλ᾽ R : δεινὰ πολλ᾽ Bentley : \lecon{πολλά γ᾽· ἀλλ᾽} Bibi}
\ligneRav{478}{}{δεινότατον, ὅτε νύμφη μὲν ἦ τρεῖς ἡμέρας,}{\textbf{478} ἦν R : \lecon{ἦ} Dindorf\espApp ἦν R : \lecon{ἦ} Dindorf (1825)}
\ligneRav{479}{}{ὁ δ᾽ ἀνὴρ παρ᾽ ἐμοὶ \difference{καθηῦδεν}· ἦν δέ μοι φίλος,}{\textbf{479} καθεύδειν R : \lecon{καθηῦδεν} Scaliger}
\ligneRav{480}{}{ὅσπερ με διεκόρησεν \difference{οὖσαν ἑπτ}έτιν·}{\textbf{480} διεκόρευσεν R : \lecon{-ρησεν} Pollux cf. \italique{Pax} 59 \espApp ἑπτ. οὖσ. R : \lecon{οὖσαν ἑπτ.} Pollux}
\ligneRav{481}{}{οὖτος πόθῳ \difference{μου ᾽κνυεν} ἐλθὼν τὴν θύραν·}{\textbf{481} μ᾽ οὔκνυεν R : \lecon{μου ᾽κνυεν} Suid. ε 520}
\ligneRav{482}{}{\difference{κᾆτ᾽} εὐθὺς ἔγνων· εἶτα καταβαίνω λάθρᾳ·}{\textbf{482} κᾶτ᾽ R : \lecon{κᾆτ᾽} Küster}
\ligneRav{}{}{ὁ δ᾽ ἀνὴρ ἐρωτᾷ `ποῖ σὺ καταβαίνεις;' `ὅποι;}{}
\ligneRav{}{}{στρόφος μ᾽ ἔχει τὴν γαστέρ᾽ ὦνερ κὠδύνη·}{}
\ligneRav{485}{}{ἐς τὸν κοπρῶν᾽ οὖν ἔρχομαι.' `βάδιζέ νυν.'}{}
\ligneRav{486}{}{κᾆθ᾽ ὁ μὲν ἔτριβεν κεδρίδας \difference{ἄννηθον} σφάκον·}{\textbf{486} ἄνηθον R, Suid. : \lecon{ἄννηθον} Dindorf (1830)}
\ligneRav{}{}{ἐγὼ δὲ καταχέασα τοῦ στροφέως ὕδωρ}{}
\ligneRav{488}{}{ἐξῆλθον ὡς τὸν μοιχόν· εἶτ᾽ \difference{ἠρειδόμην}}{\textbf{488} ἐρειδόμην R : \lecon{ἠρειδόμην} Küster}
\ligneRav{489}{}{κ\inRas{υν}{ύβ}δ᾽ ἐχομένην τῆς δάφνης·}{\textbf{489} κυνδ᾽ \Rac : \lecon{κύβδ᾽} \Rpc, \Mon}
\ligneRav{490}{}{ταῦτ᾽ οὐδεπώποτ᾽ εἶφ᾽, ὁρᾶτ᾽, Εὐριπίδης·}{}
\ligneRav{}{}{οὐδ᾽ ὡς ὑπὸ τῶν δούλων τε κὠρεωκόμων}{}
\ligneRav{}{}{σποδούμεθ᾽, ἢν μὴ 'χωμεν ἕτερον, οὐ λέγει·}{}
\ligneRav{493}{}{ὑπὸ τῶν δούλων τε κὠρεωκόμων \difference{ληκώμεθα} (\inRas{ληκ}{κιν}ώμεθα)}{\textbf{493} \lecon{ληκώμεθα} \Rac (?), Suid. : κινώ- \Rpc}
\ligneRav{494}{}{τὴν νύχθ᾽, ἕωθεν \difference{σκόροδα δια}μασώμεθα,}{\textbf{494} σκορόδια R : \lecon{σκόροδα δια} Bentley\espApp post νύχθ' add. ὅλην Suid. (reminiscentia Pherecratis ? cf. Photius λ 254) \espApp σκορόδια R : \lecon{σκόροδα δια} Bentley}
\ligneRav{495}{}{ἵν᾽ \difference{ὀσφρόμενος} \difference{ἁνὴρ} \difference{ἀπὸ} τείχους εἰσιὼν}{\textbf{495} ὀσφραινό- R : \lecon{ὀσφρόμενος} Küster \espApp ἀνὴρ R : \lecon{ἁ-} Dawes \espApp ἀπὸ τοῦ R : \lecon{ἀπὸ} Biset\espApp ὀσφραινόμενος R : \lecon{ὀσφρόμενος} Bentley, Küster \espApp ἀνὴρ R : \lecon{ἁνὴρ} Dawes \espApp ἀπὸ τοῦ R : \lecon{ἀπὸ} Biset}
\ligneRav{}{}{μηδὲν κακὸν δρᾶν ὑποτοπῆται. ταῦθ᾽, ὁρᾷς,}{}
\ligneRav{}{}{οὐπώποτ᾽ εἶπεν. εἰ δὲ Φαίδραν λοιδορεῖ,}{}
\ligneRav{}{}{ἡμῖν τί τοῦτ᾽ ἔστ᾽; οὐδ᾽ ἐκεῖν᾽ εἴρηκέ πω,}{}
\ligneRav{500}{}{ὡς ἡ γυνὴ δεικνῦσα τἀνδρὶ τοὔγκυκλον}{\textbf{500} ὑπ᾽ αὐγὰς οἷόν ἐστιν ἐγκεκαλυμμένον R : ὑπαυγάσ᾽ κτλ. Hermann : \lecon{del} van Herwerden}
\ligneRav{(500)}{}{$[\![\text{\sout{ὑπ᾽ αὐγὰς οἷόν ἐστιν ἐγκεκαλυμμένον}}]\!]$}{\textbf{v. 500} ὑπ᾽ αὐγὰς οἷόν ἐστιν ἐγκεκαλυμμένον R : \lecon{del.} van Herwerden, Bibi (ex gl. ᾦόν)}
\ligneRav{501}{}{τὸν \difference{μοιχὸν} ἐξέπεμψεν, οὐκ εἴρηκέ πω·}{\textbf{501} μυχὸν R : \lecon{μοιχὸν} Gryaneus}
\ligneRav{502}{}{ἑτέραν δ᾽ ἐγᾦδ᾽ \difference{ἣ 'φασκεν} ὠδίνειν γυνὴ}{\textbf{502} ήφασκεν R : \lecon{ἣ 'φασκεν} Zanetti}
\ligneRav{}{}{δέχ᾽ ἡμέρας, ἔως ἐπρίατο παιδίον·}{}
\ligneRav{504}{}{ὁ δ᾽ ἀνὴρ περιήρχετ᾽ \difference{ὠκυτόκι᾽} ὠνούμενος·}{\textbf{504} \lecon{περιήρχετ᾽} R : περιέτρεχ᾽ Bergk \espApp ὠκυτόκειa R : \lecon{ὠκυτόκι᾽} Pollux II.7}
\ligneRav{505}{}{τὸ δ᾽ εἰσέφερε γραῦς ἐν χύτρᾳ τὸ παιδίον,}{\textbf{505} \lecon{τὸ δ᾽} R : τότ᾽ Biset : τῇδ᾽ Daubuz}
\ligneRav{506}{}{τὸ δ᾽ εἰσέφερε γραῦς ἐν χύτρᾳ τὸ παιδίον, ἵνα μὴ \expo{βοῴη} κηρίῳ βεβυσμένον·}{\textbf{506} βοῴη suppl. \Rpc}
\ligneRav{}{}{εἶθ᾽ ὡς ἔνευσεν ἡ φέρουσ᾽, εὐθὺς βοᾷ,}{}
\ligneRav{}{}{`ἄπελθ᾽ ἄπελθ᾽, ἤδη γὰρ ὦνέρ μοι δοκῶ}{}
}
\vspace{3pt}\noindent\rule{4.5cm}{0.4pt}\par\vspace{1.5pt}
{\fontsize{7.0pt}{8.3pt}\selectfont
\noindent\textbf{Loci paralleli (vv. 466–508) ---} \textbf{466} (ὀξυθυμεῖσθαι) Photius ο 380 ὀξυθυμεῖσθαι \espTest \textbf{468} (θαυμάσιον $=$ θαυμαστόν) Moeris θ 16 $\simeq$ Hesychius θ 146 θαυμάσιος \espTest \textbf{469} Suidas ο 330 (sed ἐγώ pro ἔγωγ᾽) ; cf. (ὀναίμην) Hesychius ο 852 ὀναίμην ; \Syn ο 162 ὀναίμην $=$ Photius ο 332 \espTest \textbf{470} Suidas μ 328 (τἄνδρ᾽ pro τὸν ἄνδρ᾽) \espTest \textbf{473-4} (τί ταῦτ᾽ -- φέρομεν) Suidas τ 683 \espTest \textbf{480} Pollux III 42 (οὗτος pro ὅσπερ et διεκόρησεν cod. A pro διεκόρευσεν) \espTest (διεκόρησεν) Hesychius δ 1558 διεκόρησε \espTest \textbf{480} Pollux III 42 (οὗτος pro ὅσπερ et διεκόρησεν pro διεκόρευσεν) \espTest (διεκόρησεν) Hesychius δ 1558 διεκόρησε \espTest \textbf{481} Suidas ε 520 ; Zonaras 669 ἔκνυεν; \Eustathe 1746.6 (μοὔκνυεν ἕλκων pro μου ᾽κνυεν ἐλθὼν) \espTest \textbf{484} \Eustathe 1746.45 \espTest (στρόφος) \Eustathe 503.20 \espTest \textbf{486} Suidas σ 1712 \espTest (κεδρίδας) Hesychius κ 1980 κεδρίδες ; Photius κ 1980 κ 517 \espTest \textbf{489} (Ἀγυιᾶ) Stephanus Byzantinus 22.15 ; cf. Moeris α 13 \espTest \textbf{493-4} Suidas λ 434 Ληκώμεθα: κινώμεθα. Ἀριστοφάνης: μάλισθ' ὅταν ὑπό του ληκώμεθα τὴν νύχθ' ὅλην. Post νύχθ' add ὅλην Suid. (reminiscentia Pherecratis ? cf. Photius λ 254) \espTest \textbf{493-4} Suidas λ 434 Ληκώμεθα: κινώμεθα. Ἀριστοφάνης: μάλισθ' ὅταν ὑπό του ληκώμεθα τὴν νύχθ' ὅλην \espTest \textbf{504} (ὠκυτόκι᾽) Pollux II 7 ; \Phryn 129.3 \espTest \textbf{506} (κηρίῳ βεβυσμένον) Hesychius κ 2549 κηρίῳ βύσασα ; Photius κ 675 κηρίῳ βῦσαι\par
}
\vspace{2.5pt}\noindent\rule{4.5cm}{0.4pt}\par\vspace{1.5pt}
{\fontsize{6.9pt}{8.2pt}\selectfont
\noindent\textbf{Suite (en commentaire) de l'apparat critique \& notes paléographiques (\texttt{f. 168r}) ---} \textbf{v. 477 :} Le vers de R (οὐδ' εἴρηκεν οὐδὲν πρὸς ἃ δρῶμεν· πολλ') est amputé d'une syllabe finale ; \textbf{\textbf{Bibi} restitue πολλά γ'· ἀλλ'} (chute de -ά γ'· ἀλλ' par saut du même au même entre les deux ΛΛ). --- \textbf{v. 489 :} \Rac{} avait écrit κυνδ' (confusion minuscule $\beta/\nu$), corrigé en \lecon{κύβδ'} par \Rpc{} (suivi par \Mon{}). --- \textbf{v. 493 :} Comme le précise Louis Hauseux (\texttt{vv 1-641.tex}, l. 3715), le copiste avait d'abord écrit la forme crue \lecon{ληκώμεθα} (attestée par la Souda $\sigma\ 1126$ et adoptée par L. Hauseux), puis a \textbf{gratté les trois premières lettres ληκ-} pour les remplacer par κιν- (\lecon{\inRas{ληκ}{κιν}ώμεθα} \Rpc{}) sous l'influence de la glose κινώμεθα. --- \textbf{v. 500 (genèse paléographique d'un vers interpolé) :} Le v. 500 (ὑπ' αὐγὰς οἷόν ἐστιν ἐγκεκαλυμμένον), athétisé par van Herwerden et L. Hauseux (\texttt{vv 1-641.tex}, ll. 3781–3806), est né d'une glose ou citation marginale [τὸν μοιχόν·] ὑπαυγάσας. ᾦόν ἐστιν ἐγκεκαλυμμένον, dans laquelle ΩΙΟΝ (avec iota muet adscrit = $\text{ᾦόν}$, « un œuf enveloppé ») a été lu par iotacisme et bêtacisme $\omega/\text{ο}$ \textbf{ΟΙΟΝ (οἷόν)}, entraînant la réécriture métrique ὑπ' αὐγὰς οἷόν ἐστιν.... --- \textbf{v. 506 :} L'optatif \lecon{βοῴη}, omis par \Rac{}, est suppléé par \Rpc{}. --- \textbf{v. 477 :} πολλ' R (δεινὰ πολλ' Bentley) : \lecon{πολλά γ'· ἀλλ'} \lecon{Bibi}. --- \textbf{v. 481 :} μ' οὔκνυεν R (médivision et diphtongaison) : \lecon{μου 'κνυεν} Souda $\varepsilon\ 520$ (cf. Hésychius $\kappa\ 3133$). --- \textbf{v. 501 :} μυχὸν R (iotacisme) : \lecon{μοιχὸν} Grynaeus.\par
}
\end{minipage}
\end{adjustbox}
\end{adjustwidth}
\newpage
\enTeteD{Codex Ravennas 429 (Biblioteca Classense) --- Fac-similé du f. 168r (vv. 466–508)}
\begin{adjustwidth}{-1.5cm}{-3.0cm}
\begin{center}
\includegraphics[width=16.5cm,height=20.5cm,keepaspectratio]{/workspaces/Thesmophories/dossier_ravennas_thesmophories/folios_individuels_ff162v-177r/12_f168r_classense170r_vv0466-0508.jpg}\par
\vspace{2pt}
{\scriptsize\textbf{Codex Ravennas 429, f. 168r (Classense 170r)} --- ΕΠΕΙΣΟΔΙΟΝ Α', \textbf{vv. 466–508} ($\text{R}^1, \Rpc, \Rdeux$ (\Mon)).}
\end{center}
\end{adjustwidth}
\newpage
\clearpage\ifodd\value{page}\hbox{}\newpage\fi
\enTeteG{ΘΕΣΜΟΦΟΡΙΑΖΟΥΣΑΙ --- Texte grec établi (f. 168r, vv. 466–508)}
\enTeteD{LES THESMOPHORIES --- Traduction française (f. 168r, vv. 466–508)}
\begin{pages}
\begin{Leftside}
\beginnumbering
%% v. 466-519
\begin{groupe}%
%
\personnageGrec{Μνησίλοχος}%
\trimIamb{τὸ μὲν ὦ γυναῖκες ὀξυθυμεῖσθαι σφόδρα%\testimonia{466}{(ὀξυθυμεῖσθαι) Photius ο 380 ὀξυθυμεῖσθαι}%
\commentaire{466 sqq.}{}{Scholiaste : \grec{ἐπιτηέδες ὀ κηδεστὴς ὑπὲρ τὰς ἄλλας γυναῖκας γυναικίζεται, ἵνα διὰ τούτου λαθῃ}... prudence qui va bientôt le quitter.}\\
Εὐριπίδῃ, τοιαῦτ᾽ ἀκουούσας κακά,\apparat{467}{ἀκουούσαις R : \lecon{ἀκουούσας} Zanetti}\\% (\italique{audientes} Divus)}\\
οὐ θαυμάσιόν ἐστ᾽, οὐδ᾽ ἐπιζεῖν τὴν χολήν.%\testimonia{468}{(θαυμάσιον $=$ θαυμαστόν) Moeris θ 16 $\simeq$ Hesychius θ 146 θαυμάσιος}
\\
καὐτὴ γὰρ ἔγωγ᾽, οὕτως ὀναίμην τῶν τέκνων,%\testimonia{469}{Suidas ο 330 (sed ἐγώ pro ἔγωγ᾽) ; cf. (ὀναίμην) Hesychius ο 852 ὀναίμην ; \Syn ο 162 ὀναίμην $=$ Photius ο 332}
\apparat{469}{ὀνοίμην R : \lecon{ὀναίμην} Suid. ο 330}\\
μισῶ τὸν ἄνδρ᾽ ἐκεῖνον, εἰ μὴ μαίνομαι.%\testimonia{470}{Suidas μ 328 (τἄνδρ᾽ pro τὸν ἄνδρ᾽)}
\\
ὅμως δ᾽ ἐν ἀλλήλαισι χρὴ δοῦναι λόγον·\apparat{471}{ἀλλήλοισι R : \lecon{ἀλλήλαισι} Brunck}%
%\commentaire{471}{ἐν ἀλλήλοισι χρὴ δοῦναι λόγον.}{Même expression dans Hérodote (I, 97) : « καὶ ἐδίδοσαν σφίσι λόγον ».}
\\
αὐταὶ γάρ ἐσμεν, κοὐδεμί᾽ ἔκφορος λόγου.\\%\apparat{472}{\lecon{ἔκφορος} R : ἐκφορὰ Valckenaer ad Eur. \italique{Hipp.} 294 \ligneApp{473} κεῖνον R, Suid τ 683 : \lecon{᾽κεῖνον} Meineke}\\
τί ταῦτ᾽ ἔχουσαι ᾽κεῖνον αἰτιώμεθα%\testimonia{473-4}{(τί ταῦτ᾽ -- φέρομεν) Suidas τ 683}
\apparat{473}{κεῖνον R, Suid τ 683 : \lecon{᾽κεῖνον} Meineke \ligneApp{474} ἢ δύ᾽ R : \lecon{εἰ δύ᾽} Grynaeus}% (\italique{si} Divus}%
\commentaire{473}{τί ταῦτ᾽ ἔχουσαι κεῖνον αἰτιώμεθα.}{Ce vers rappelle celui des \italique{Acharniens} (vers 514) :
\citComm{τί ταῦτα τοὺς Λάκωνας αἰτιώμεθα;}
Ils pourraient donc bien être tous deux une reprise parodique du \italique{Télèphe} d'Euripide, pièce maintes fois détournée ici et là.}\\
βαρέως τε φέρομεν, εἰ δύ᾽ ἡμῶν ἢ τρία\\%\apparat{474}{ἢ δύ᾽ R : \lecon{εἰ δύ᾽} Grynaeus (\italique{si} Divus)}\\
κακὰ ξυνειδὼς εἶπε δρώσας μυρία;\\
ἐγὼ γὰρ αὐτὴ πρῶτον, ἵνα μἄλλην λέγω,\apparat{476}{μὴ ἄλλην R : \lecon{μἄλλην} Dindorf %(1839) 
\ligneApp{477} πολλ᾽ R : δεινὰ πολλ᾽ Bentley : \lecon{πολλά γ᾽· ἀλλ᾽} Bibi \ligneApp{478} ἦν R : \lecon{ἦ} Dindorf}\\% (1825)}\\
ξύνοιδ᾽ ἐμαυτῇ πολλά γ᾽· ἀλλ᾽ ἐκεῖνο δ᾽ οὖν%\apparat{477}{πολλ᾽ R : δεινὰ πολλ᾽ Bentley : \lecon{πολλά γ᾽· ἀλλ᾽} Bibi \ligneApp{478} ἦν R : \lecon{ἦ} Dindorf (1825)}%
\commentaire{477}{ξύνοιδ᾽ ἐμαυτῇ πολλ᾽ ἐκεῖνο δ᾽ οὖν.}{Il manque un pied au vers 477 ; Bentley y suppléa en ajoutant <δεινὰ> πολλ᾽ pour introduire le superlatif δεινότατον du vers suivant. C'est une hypothèse (Dawes préféra πολλὰ δείν᾽). J'en propose une autre qui a l'avantage d'expliquer paléographiquement la lacune par un \italique{saut du même au }quasi-\italique{même} (il n'y a qu'une patte de mouche de différence entre un π et un γ) : \textbf{ΠΟΛΛ'} à la place de ΠΟΛΛΑ\textbf{Γ'ΑΛΛ'}.}\\
δεινότατον, ὅτε νύμφη μὲν ἦ τρεῖς ἡμέρας,}
\end{groupe}%
%% v. 466-519
\begin{groupe}%
%
\trimIamb{ὁ δ᾽ ἀνὴρ παρ᾽ ἐμοὶ καθηῦδεν· ἦν δέ μοι φίλος,\apparat{479}{καθεύδειν R : \lecon{καθηῦδεν} Scaliger % : καθεῦδεν Bentley, Brunck 
\ligneApp{480} διεκόρευσεν R %, Pollux III.42 codd.FS 
: \lecon{-ρησεν} Pollux %cod.A %, Fritzsche 
cf. \italique{Pax} 59
\espApp ἑπτ. οὖσ. R : \lecon{οὖσαν ἑπτ.} Pollux % \ligneApp{481} μ᾽ οὔκνυεν R : \lecon{μου ᾽κνυεν} Suid. ε 520 
\ligneApp{482} κᾶτ᾽ R : \lecon{κᾆτ᾽} Küster}\\% : κἀγὼ Lenting}\\
ὅσπερ με διεκόρησεν οὖσαν ἑπτέτιν.%\testimonia{480}{Pollux III 42 (οὗτος pro ὅσπερ et διεκόρησεν cod. A pro διεκόρευσεν) \espTest (διεκόρησεν) Hesychius δ 1558 διεκόρησε}%
\commentaire{480}{με διεκόρευσεν ἑπτέτιν οὖσαν.}{Mnésiloque n'aura pas mis longtemps à perdre toute retenue : il s'invente bien imprudemment une histoire repoussante qui n'a d'égale que celle de la Quartilla de Pétrone (\italique{Satyrico}, XXV) ou de la marquise de Brinvilliers. (Willems.)}\\
οὖτος πόθῳ μου ᾽κνυεν ἐλθὼν τὴν θύραν·%\testimonia{481}{Suidas ε 520 ; Zonaras 669 ἔκνυεν; \Eustathe 1746.6 (μοὔκνυεν ἕλκων pro μου ᾽κνυεν ἐλθὼν)}
\\
κᾆτ᾽ εὐθὺς ἔγνων· εἶτα καταβαίνω λάθρᾳ.\\
ὁ δ᾽ ἀνὴρ ἐρωτᾷ `ποῖ σὺ καταβαίνεις;' `ὅποι;\\
στρόφος μ᾽ ἔχει τὴν γαστέρ᾽ ὦνερ κὠδύνη·%\testimonia{484}{\Eustathe 1746.45 \espTest (στρόφος) \Eustathe 503.20}%
\commentaire{484}{στρόφος μ᾽ ἔχει τὴν γαστέρ᾽ ὦνερ κὠδύνη·}{\grec{στρόφος} est le terme «~propre~» pour désigner la \italique{colique} (cf. Hippocrate, \italique{Des maladies aiguës}, XXXVII ; \italique{Aphorismes}, V, 41).}\\
ἐς τὸν κοπρῶν᾽ οὖν ἔρχομαι.' `βάδιζέ νυν.'\\
κᾆθ᾽ ὁ μὲν ἔτριβεν κεδρίδας ἄννηθον σφάκον·%\testimonia{486}{Suidas σ 1712 \espTest (κεδρίδας) Hesychius κ 1980 κεδρίδες ; Photius κ 1980 κ 517}
\apparat{486}{%\lecon{ἔτριβεν} R : ἔτριβε Suid. σ 1712 \espApp 
ἄνηθον R, Suid. : \lecon{ἄννηθον} Dindorf (1830)}%
\commentaire{486}{κᾆθ᾽ ὁ μὲν ἔτριβεν κεδρίδας ἄννηθον σφάκον·}{Le niais mari se jette sur les graines de genévrier, d'aneth et de sauge car elles sont réputées préserver du mal que sa femme prétend avoir. Voir ce que dit Pline l'Ancien au sujet de ces plantes. %\marqueur{1}.
(Pline l'Ancien, \italique{Histoire naturelle}. D'après l'édition et la traduction d'Émile Littré, \italique{Histoire naturelle de Pline}, tome second (1877). XXIV, 36 pour le genévrier (ἡ κεδρίς) ; XX, 74 pour l'aneth (τό ἄννηθον) ; XXII, 71 pour la sauge (ὁ σφάκος).)
%À toutes fins utiles pour le lecteur – qui pourrait bien trouver un réel secours grâce à ces connaissances qui ne sont pas toutes théoriques –, nous reproduisons ce que Pline l'Ancien nous apprend des vertus de ces plantes :
%\citComm{Pour le genévrier (ἡ κεδρίς)\marqueur{1} :\\
%\latin{Juniperus vel ante caetera omnia excalfacit, extenuat, cedro assimilis. Et ejus duo genera : altera major, altera minor. Utraque accensa serpentes fugat. Semen \textbf{stomachi}, pectoris, lateris \textbf{doloribus utile}.}\\
%\francais{Le genévrier est plus que toute autre plante échauffant et atténuant, semblable en cela au cèdre. Il en est de deux espèces : l'un, plus grand que l'autre. Tous deux, brûlés, font fuir les serpents. La graine est bonne dans les \italique{douleurs d'estomac}, de poitrine et de côté.}\\
%Pour l'aneth (τό ἄννηθον)\marqueur{2} :\\
%\latin{Anethum quoque ructus movet , et tormina sedat : alvum sistit : epiphoris radices illinuntur ex aqua vel vino.}\\
%\francais{L'aneth (anethum graveolens, L.) aussi est carminatif, et calme les tranchées ; Il arrête les flux de ventre. On applique la racine dans de l'eau ou dans du vin sur les épiphoras.}\\
%Et enfin la sauge (ὁ σφάκος)\marqueur{3} :\\
%\latin{Est et silvestris elelisphacos} \grec{[le scholiaste du \italique{Ravennas} glose : σφάκον· ἐλελίσφακον. ὡς ἐπιτήδεια παῦσαι στρόφον]} \latin{dicta a Graecis, ab aliis \textbf{sphacos}. Ea est sativa lente levior, et folio minore, atque sicciore, et odoratiore. [...] \textbf{Bibitur cum absinthio ad dysenteriam}. Cum vino eadem commorantes menses trahit : abundantes sistit decocto ejus poto. Per se imposita herba vulnerum sanguinem cohibet.}\\
%\francais{La plante appelée par les Grecs élélisphacos ou \italique{sphacos} (salvia pomifera, L.) est une espèce de lentille sauvage plus légère que la lentille cultivée, à feuilles plus petites, plus sèches et plus odorantes. [...] On la fait boire avec l'absinthe pour la dysenterie. Avec le vin, elle fait venir les règles en retard; en décoction, elle arrête les règles trop abondantes ; appliquée seule, elle étanche le sang des plaies.}}
}%
%\debutNote{486}%
%\note{1}{Pline l'Ancien, \italique{Histoire naturelle}. D'après l'édition et la traduction d'Émile Littré, \italique{Histoire naturelle de Pline}, tome second (1877). XXIV, 36 pour le genévrier (ἡ κεδρίς) ; XX, 74 pour l'aneth (τό ἄννηθον) ; XXII, 71 pour la sauge (ὁ σφάκος).}%
%\note{1}{Pline l'Ancien, \italique{Histoire naturelle}, XXIV, 36. Je reprends l'édition et la traduction d'Émile Littré, \italique{Histoire naturelle de Pline}, tome second (1877).}%
%\note{2}{Pline l'Ancien, \italique{Histoire naturelle}, XX, 74.}%
%\note{3}{Pline l'Ancien, \italique{Histoire naturelle}, XXII, 71.}%
\\
ἐγὼ δὲ καταχέασα τοῦ στροφέως ὕδωρ
\commentaire{487}{ἐγὼ δὲ καταχέασα τοῦ στροφέως ὕδωρ.}{La perfide est rouée (στρόφις) : elle asperge d'eau les gonds de la porte pour ne pas qu'ils fassent de bruit. Le scholiaste l'avait compris ainsi : ἵνα ὀλισθηρὸς γενόμενος μὴ ψοφῇ. On a le même tour chez Plaute (\italique{Curculio}, I, 3, 1) :%
\citComm{\latin{Anus — Placide egredere, et sonitum prohibe forium, et \textbf{crepitum cardinum},\\
ne quod hic agimus, herus percipiat fieri, mea Planesium.\\
mane, \textbf{subfundam aquolam}.}%
}%
J'attire l'attention du lecteur sur le jeu de mots entre καταχέω / καταχέζω d'une part et στροφεύς / στρόφος d'autre part. Je n'en dis pas davantage car en voilà-z-assez !}\\
ἐξῆλθον ὡς τὸν μοιχόν· εἶτ᾽ ἠρειδόμην\apparat{488}{ἐρειδόμην R : \lecon{ἠρειδόμην} Küster \ligneApp{489} κυνδ᾽ \Rac : \lecon{κύβδ᾽} \Rpc, \Mon}\\
παρὰ τὸν Ἀγυιᾶ κύβδ᾽ ἐχομένη τῆς δάφνης.%\testimonia{489}{(Ἀγυιᾶ) Stephanus Byzantinus 22.15 ; cf. Moeris α 13}%\apparat{489}{κυνδ᾽ \Rac : \lecon{κύβδ᾽} \Rpc, \Mon}%
\commentaire{489}{παρὰ τὸν Ἀγυιᾶ κύβδ᾽ ἐχομένη τῆς δάφνης.}{L'Agyieus, surnom tiré d'ἀγυιά, la \italique{rue}, donné à Apollon, considéré comme protecteur des rues, devint par extension le nom des images et autels élevés en l'honneur du dieu, à Athènes, devant les portes des maisons (ἐν προθύροις). (Note d'Émile Saglio\marqueur{1}.) Mnésiloque parle bien d'un autel, ceint selon l'usage de laurier ; Hésychius (α 856) :
\citComm{ἀγυιεύς· ὁ πρὸ τῶν θυρῶν ἑστὼς βωμὸς ἐν σχήματι κίονος\\
\francais{Agyieus : l'autel placé devant les portes, en forme de pilier.}}%
Sachant ce qu'il y fait, Mnésiloque ajoute donc l'insulte au dieu à l'outrage de l'adultère. Au vu de la forme de l'ἀγυιεύς, on peut imaginer comment il prétend s'être tenu `παρὰ τὸν Ἀγυιᾶ κύβδα', mais pourquoi `ἐχομένη τῆς δάφνης', \italique{s'agrippant au laurier} ?! Quoique porté aux \italique{Realia}% par ma formation à l'école de Willems
, je me suis dispensé de reconstituer la scène ; cela ne m'empêchera pas d'affirmer : voilà une posture fort incommode !
Pourtant, cette image érotique d'un \italique{Penseur} de Rodin (κύβδα) se faisant levrer alors qu'il se tient à un arbre se retrouve ailleurs\marqueur{2} ; il y a un lièvre à lever. 
La réponse se trouve peut-être chez les parémiographes : ils ont recensé une des \italique{formules d'Éphèse} (Ἐφέσια γράμματα), incantations magiques pour conjurer le mauvais sort\marqueur{3}. :%
\citComm{Ἐφέσια γράμματα: ἐπῳδαί τινες ἦσαν, ἅσπερ οἱ φωνοῦντες ἐνίκων ἐν παντί. ταὐτὸ τῇ [lire : ταὐτότητι ?]·\
Δαφνίνην φορῶ βακτηρίαν.\\
Lettres éphésiennes : paroles magiques que ceux qui les prononçaient triomphaient en tout. En répétant :
« Je porte la branche de laurier ».}%
Ainsi, les Grecs \italique{touchaient du bois} (de laurier) pour éviter les malheurs ; cela pourrait expliquer l'attitude de notre libidineuse : s'accrochant aux branches et à de superstitieuses croyances, elle pense ainsi échapper au mauvais œil.% d'un flagrant délit d'adultère.
}%
\debutNote{489} %
\note{1}{Début de l'article \latin{Agyieus} de l'indispensable \italique{Dictionnaire des antiquités grecques et romaines} de Daremberg \& Saglio ; l'on peut y voir des représentations d'Agyeus.}%
\note{2}{Lire Théocrite, \italique{Idylles}, V, 115-119 et voir, par exemple, – ce qui n'est d'ailleurs pas indispensable ! – \italique{Antike Denkmäler}, Band II, 1908, Tafel 41, Figure 3.}%
\note{3}{\italique{Corpus paroemiographorum graecorum}, I, ed. E. L. a Leutsch, Diogeniani centuria IV 78, p. 244. On pourra consulter la \italique{Souda} ou Photius à l'entrée Ἀλεξιφάρμακα : ces charmes pouvaient concerner des époux. Aristophane met dans la bouche de Chrémyle qui prétend avoir quelque succès (\italique{Ploutos}, v. 213) : ὁ Φοῖβος αὐτὸς Πυθικὴν σείσας δάφνην.%Noter que ce fameux bâton de laurier que pouvaient prendre les convives lors de symposies avant d'entonner le chant était appelé αἴσακος, qui est aussi le nom d'un centaure..... Nonnios : postérieur ?
}\\
ταῦτ᾽ οὐδεπώποτ᾽ εἶφ᾽, ὁρᾶτ᾽, Εὐριπίδης·\\
οὐδ᾽ ὡς ὑπὸ τῶν δούλων τε κὠρεωκόμων\\
%
σποδούμεθ᾽, ἢν μὴ 'χωμεν ἕτερον, οὐ λέγει·}
\end{groupe}%
%% v. 466-519
\begin{groupe}%
%
\trimIamb{οὐδ᾽ ὡς ὅταν μάλισθ᾽ ὑπό του ληκώμεθα%\testimonia{493-4}{Suidas λ 434 Ληκώμεθα: κινώμεθα. Ἀριστοφάνης: μάλισθ' ὅταν ὑπό του ληκώμεθα τὴν νύχθ' ὅλην. Post νύχθ' add ὅλην Suid. (reminiscentia Pherecratis ? cf. Photius λ 254)}
\apparat{493}{\lecon{ληκώμεθα} \Rac (?), Suid. %λ 434 
: κινώ- \Rpc %, \Mon 
\ligneApp{494} σκορόδια R : \lecon{σκόροδα δια} Bentley \ligneApp{495} ὀσφραινό- R : \lecon{ὀσφρόμενος} Küster \espApp ἀνὴρ R : \lecon{ἁ-} Dawes \espApp ἀπὸ τοῦ R : \lecon{ἀπὸ} Biset}%
\commentaire{493}{μάλισθ᾽ ὑπό του ληκώμεθα.}{Le scribe ici encore a hésité entre \grec{ληκώμεθα}, qu'il écrivit probablement en premier, avant de gratter les trois premières lettres et de réécrire par dessus \grec{κινώμεθα}. N'étant pas un aficionado de la \italique{lectio difficilior}, il aura probablement préféré le verbe \grec{κινεῖν}, bien connu de lui, au fort rare \italique{ληκᾶν} (ou \grec{ληκᾶσθαι} dans la bouche d'une femme). Suidas (λ 434) ne s'y est pas trompé qui glose ληκώμεθα par κινώμεθα en citant ce passage d'Aristophane.}\\
τὴν νύχθ᾽, ἕωθεν σκόροδα διαμασώμεθα,%\apparat{494}{post νύχθ' add. ὅλην Suid. (reminiscentia Pherecratis ? cf. Photius λ 254) \espApp σκορόδια R : \lecon{σκόροδα δια} Bentley \ligneApp{495} ὀσφραινόμενος R : \lecon{ὀσφρόμενος} Bentley, Küster \espApp ἀνὴρ R : \lecon{ἁνὴρ} Dawes \espApp ἀπὸ τοῦ R : \lecon{ἀπὸ} Biset}%
\commentaire{494}{ἕωθεν σκόροδα διαμασώμεθα.}{(Pour ne pas que le mari cocu fleure le sacrifice à Vénus de sa femme.) 
À quand une thesmophorienne ecclésie pour châtier Aristophane qui enseigne aux Athéniens toutes les fourberies de leur femme ?.. à moins que le tour ne fût déjà connu (Athénée, XIII, 572c.) :%
\citComm{Διὰ ταῦθ' ὁ πόρνος οὗτος οὐδὲ τῶν πράσων\\
ἑκάστοτ' ἐπεδείπνει μεθ' ῆμῶν· τοῦτο δ' ἦν,\\
ἵνα μή τι λυπήσειε τὸν ἐραστὴν φιλῶν.}%
et encore avec de l'oignon (Xénophon, \italique{Symposium}, IV, 8.) :%
\citComm{Καὶ ὁ Χαρμίδης εἶπεν· `Ὦ ἄνδρες, ὁ Νικήρατος κρομμύων ὄζων ἐπιθυμεῖ οἴκαδε ἐλθεῖν, ἵν᾽ ἡ γυνὴ αὐτοῦ πιστεύῃ μηδὲ διανοηθῆναι μηδένα ἂν φιλῆσαι αὐτόν'\marqueur{2}.}%
}%
%\debutNote{494}%
%\note{1}{Athénée, XIII, 572c.}%
%\note{2}{Xénophon, \italique{Symposium}, IV, 8.}
\\
ἵν᾽ ὀσφρόμενος ἁνὴρ ἀπὸ τείχους εἰσιὼν%
\commentaire{495}{ἀπὸ τοῦ τείχους εἰσιὼν.}{Thucydide (VIII, 69) :%
\citComm{ἦσαν [δ᾽] Ἀθηναῖοι πάντες αἰεὶ οἱ μὲν ἐπὶ τείχει, οἱ δ᾽ ἐν τάξει, τῶν ἐν Δεκελείᾳ πολεμίων ἕνεκα ἐφ᾽ ὅπλοις.\\
\francais{Tous les Athéniens restaient continuellement, les uns sur les murailles, les autres à leurs postes, à cause des ennemis qui occupaient Décélie.}%
}%
}\\
μηδὲν κακὸν δρᾶν ὑποτοπῆται. ταῦθ᾽, ὁρᾷς,%
%\commentaire{496}{ὁρᾷς}{\eff{Pourquoi Mnésiloque passe-t-il du vouvoiement au tutoiement ?}}
\\
\ligne{497}οὐπώποτ᾽ εἶπεν. εἰ δὲ Φαίδραν λοιδορεῖ,\\
ἡμῖν τί τοῦτ᾽ ἔστ᾽; οὐδ᾽ ἐκεῖν᾽ εἴρηκέ πω,\\
ὡς ἡ γυνὴ δεικνῦσα τἀνδρὶ τοὔγκυκλον%
\commentaire{499}{τοὔγκυκλον.}{L'ἔγκυκλον était une sorte de jupe courte, portée par les femmes à l'intérieur. Cf. la note au vers 643.}%
\apparat{500}{ὑπ᾽ αὐγὰς οἷόν ἐστιν ἐγκεκαλυμμένον R : ὑπαυγάσ᾽ κτλ. Hermann %(1838) 
: \lecon{del} van Herwerden %, Willems 
\ligneApp{501} μυχὸν R : \lecon{μοιχὸν} Gryaneus}%
\commentaire{500}{ὑπ᾽ αὐγὰς οἷόν ἐστιν ἐγκεκαλυμμένον.}{Le manuscrit de Ravenne comporte :
\citComm{ὡς ἡ γυνὴ δεικνῦσα τἀνδρὶ τοὔγκυκλον\\
ὑπ᾽ αὐγὰς οἷόν ἐστιν ἐγκεκαλυμμένον\\
τὸν μυχὸν ἐξέπεμψεν, οὐκ εἴρηκέ πω.}
Gardons-nous, dans un premier temps, de biffer le \nieme{500} vers, ὑπ᾽ αὐγὰς οἷόν ἐστιν ἐγκεκαλυμμένον, à l'instar de Willems\marqueur{1}. Si effectivement interpolation il y a -- ce que nous allons essayer de démontrer --, celle-ci ne peut être que très ancienne ; en effet, une scholie érudite notait déjà l'abrègement violent (συστολή βίαιος) que devait subir ὑπ᾽ αὐγὰς pour rentrer dans le mètre iambique. 
%J'ai longtemps voulu conserver ὑπ᾽ αὐγὰς \italique{contra metrum} (quitte à le remplacer par des formes proches comme ὑπαυγάσ᾽ ou ὑπαυγάσας) car il me semblait donner des indications sur le stratagème employé par la pécheresse.
Willems parlait d'une «~pure cheville~» à propos de `οἷόν ἐστιν'. Le cas me paraît plus compliqué. %, lui, me laissait perplexe... (.)% ; le cas me semble plus grave : l'on doit faire violence au grec pour en proposer un sens acceptable.) 
Intéressons-nous y donc !\\
Après avoir retourné l'affaire dans tous les sens, je ne m'explique cet οἷόν ἐστιν que par la présence originelle de : ᾦον ἐστιν\marqueur{2}. Athénée (II, 57e) nous précise que ᾦον pouvait être utilisé pour ὑπερῷον, l'\italique{étage supérieur d'une maison}. Le scholiaste T de l'\italique{Iliade} (XVI, 184) emploie d'ailleurs le mot en ce sens -- était-ce le même que le scribe qui écrivit notre scholie devenue le vers 500\marqueur{3} ? Quoi qu'il en soit, l'adjectif ὑπερῷος signifie \italique{situé à l'étage supérieur}.
Si l'hypothèse est exacte, le sens se laisse alors clairement dégager : on comprend que l'amant (τὸν μοιχόν) regarde vers le bas (αὐγάζειν, au sens d'\italique{observer}, est préfixé de ὑπ- par décalque sur ὑφ-ορᾶν) la femme embobiner son cocu de mari, bien camouflé (ἐγκεκαλυμμένον). L'amant est donc à l'étage supérieur (ᾦον ἐστιν).
Une fois cette correction minime apportée, le \nieme{500} vers se peut lire ainsi : 
\citComm{[τὸν μοιχόν·] ὑπαυγάσας. ᾦον ἐστιν ἐγκεκαλυμμένον.}
Ainsi compris, il devient tout-à-fait clair que ce « vers » n'est effectivement qu'une glose du vers 501. Le malheur voulut qu'elle fît peu ou prou un trimètre iambique et soit alors très tôt intégrée au texte même d'Aristophane ; l'inintelligibilité du nouveau texte aidant à sa corruption, ΥΠΑΥΓΑΣΑΣ devint ΥΠΑΥΓΑΣ par haplographie et ΩΙΟΝ ΟΙΟΝ.
Lisons donc avec Willems : %
\citComm{ὡς ἡ γυνὴ δεικνῦσα τἀνδρὶ τοὔγκυκλον\\
τὸν μοιχὸν ἐξέπεμψεν, οὐκ εἴρηκέ πω.\\
\francais{Comment la femme qui exhibant l'encycle à son mari\\
fit échapper son amant, il n'en a pas soufflé mot.}}%
Le passage, ici comme ailleurs, n'est qu'allusif ; aussi fait-il probablement référence à une histoire bien connue des spectateurs. Victor Coulon\marqueur{4} l'a rapprochée de contes au motif semblable, couchés par écrit au \nieme{XI} siècle, et de source probablement arabe (parfois, la femme est aidée de sa mère dans son entreprise de mystification du cocu)~; les conteurs orientaux ayant eux-mêmes été familiarisés aux contes grecs. Ainsi, il nous apparaît plus que plausible qu'Aristophane rapporte ici non pas un fait-divers de la chronique scandaleuse du temps mais une historiette populaire ; ces «~révélations~» sur les Athéniennes ne seraient pas si compromettantes.\\% pour elles que ne l'étaient celles d'Euripide.\\
Voir le commentaire au vers 643 : l'encycle (ἔγκυκλον) est une courte jupe. Cette donnée finit d'anéantir les hypothèses loufoques %proposées par de savants universitaires 
où la femme cacherait l'amant aux yeux du mari en tendant l'encycle devant l'amant.
%Ce vers pèche contre le mètre. Willems, dans sa \nieme{13} note reproduite en annexe argue de cela et de sa construction des plus mal embouchées pour le biffer promptement. Pourtant, la présence de «~\grec{ὑπ᾽ αὐγὰς}~» paraît indispensable pour la révélation du stratagème : la femme emmène son mari à l'extérieur, \italique{au soleil} pour lui montrer son linge bien tissé\marqueur{1} et ce afin de laisser à son amant la possibilité de filer à l'anglaise\marqueur{2}. Je reprends la première conjecture de Hermann qui a l'avantage de ne changer aucune lettre : ὑπαυγάσ(αι)᾽ οἷόν ἐστιν. Le verbe \grec{ὑπαυγάζειν} n'est que peu attesté~: transitivement, Aristote\marqueur{3} l'emploie avec le sens de \italique{faire briller}, cependant Sophocle emploie bien le verbe αὐγάζειν avec le sens de \italique{voir} (au vers 217 de son \italique{Philoctète} ; de toute façon s'il fallait justifier par des exemples de chaque emploi d'un mot par Aristophane, nos traductions resteraient bien lacunaires !..) Ici, il est clair que le sens est : \italique{montrer \grec{ὑπ᾽ αὐγὰς}}~: \italique{montrer en plein jour}. Victor Coulon\marqueur{4} a rapproché cette histoire de contes au motif semblable, couchés par écrit au \nieme{XI} siècle, et de source probablement arabe (parfois, la femme est aidée de sa mère dans son entreprise de mystification du cocu)~: les conteurs orientaux ayant eux-mêmes été familiarisés aux contes grecs. Ainsi, il nous apparaît plus que plausible qu'Aristophane rapporte ici non pas un fait-divers de la chronique scandaleuse du temps mais une historiette populaire ; aussi ces «~révélations~» sur les Athéniennes ne sont peut-être pas aussi compromettantes pour elles que ne l'étaient celles d'Euripide. Passons à ἐγκεκαλυμμένον, littéralement \italique{bien affublé} d'où le sens dérivé de \italique{caché}. Ce terme, qui se rapporte au galant (τὸν μοιχὸν), est de prime abord étonnant : on comprend que l'amant veuille se cacher pour ne pas être reconnu en sortant, mais comment alors qu'il n'a que le bref laps de temps où la femme attire son mari à l'extérieur, comment ainsi pris de court trouverait-il le temps de bien se vêtir ? On l'imaginerait plutôt sortir précipitamment dans la même tenue – celle d'Adam – que lorsqu'il fut interrompu par l'arrivée inopportune du cocu. À la réflexion, c'est évident : le manteau (τοὔγκυκλον) prétendûment nouvellement tissé ne peut être, pour tromper le mari, que celui de l'amant ; de plus, il ne peut dans sa fuite dérober à la volée des habits sous peine d'éveiller les soupçons du cocu. Aussi, je mettrais ma main au feu que la version populaire – c'est-à-dire connue de tous les spectateurs – du conte faisait sortir l'amant précipitamment dans son plus simple appareil. Alors pourquoi cet ἐγκεκαλυμμένον ? Comme toujours : Hésychius à la rescousse ! Il glose comme exprès pour nous (ε 186) :
%\citComm{ἐγκαλύπτεσθαι· διατρέπεσθαι, αἰσχύνεσθαι.}%
%C'est bien cela : le cocufieur sort \italique{troublé}, \italique{honteux} comme un renard qu'une poule aurait pris... et nu comme un ver ! Aristophane use plaisamment à dessein de l'adjectif ἐγκεκαλυμμένον pour décrire cette situation ; il me semble que l'équivoque est plutôt bien rendue par l'expression : « partir la queue entre les jambes ».
}%
\debutNote{500}%
%\note{1}{De tout temps, la femme a pris plaisir à admirer de belles robes ; \cf Eur. \italique{Héc.} 1154 où les captives troyennes louent le tissage des femmes thraces \grec{ὑπ᾽ αὐγὰς τούσδε λεύσσουσαι πέπλους}. (Note de Victor Coulon, référence ci-dessous.)}%
%\note{2}{Et non pas l'invraisemblable interprétation de ces vers qu'on peut lire çà et là : la femme déploierait son \italique{encycle} pour camoufler au mari la fuite de l'amant caché derrière (sic) avant que ce dernier ne trouve encore le moyen de s'enfuir (re-sic). Il m'a été donné de lire un papier traitant de l'identification de ce vêtement qui donnait du crédit à cette interprétation. Et voilà comment une exégèse imbécile peut conduire à voir en l'ἔγκυκλον non une sorte de jupe, ce qu'il était (comme nous l'allons montrer à la note du vers 643), mais un \italien{mantello ampio, pesante}...}%
%\note{3}{Aristote, \italique{Ars Rhetorica}, 3, 1378}%
\note{1}{Willems argue de sa faute de mètre et de sa construction des plus mal embouchées. Le lecteur pourra utilement se reporter à \nieme{13} note ainsi que son introduction aux \italique{Notes sur la} Paix (t. II, pp. 93-99), toutes deux reproduites en annexe. Le premier à notre connaissance à avoir proposé la suppression du vers 500 est van Herwerden (\italique{Mnemosyne N. S.}, 24, 1896, p.293).}%
\note{2}{Rappelons que les iotas souscrits n'apparaissent qu'au cours du \nieme{XIII} siècle ; le manuscrit de Ravenne, du \nieme{IX} ou \nieme{X} siècle est adscrit tout comme les manuscrits qu'il transcrivait. La faute -- si faute il y a -- ne consiste par conséquent qu'en la confusion entre un Ω et un Ο.}%
\note{3}{Ce même scholiaste T connaissait bien notre comédie : il en cite des passages comme les vers 536-537 (voir les \latin{testimonia} \italique{ad locum}.).}%
\note{4}{«~Interprétation d'un passage d'Aristophane imité par les conteurs du Moyen-Âge~», in \italique{Annuaire de l'institut de philologie et d'histoire orientales}, tome II(1934) : \italique{Mélanges Bidez}, premier volume, p. 121-34.}\ligne{500}\\
τὸν μοιχὸν ἐξέπεμψεν, οὐκ εἴρηκέ πω.\\
\ligne{502}ἑτέραν δ᾽ ἐγᾦδ᾽ ἣ 'φασκεν ὠδίνειν γυνὴ\\%\apparat{502}{ήφασκεν R : \lecon{ἣ 'φασκεν} Zanetti}\\
δέχ᾽ ἡμέρας, ἔως ἐπρίατο παιδίον·}
\end{groupe}%
%% v. 520-530 (part 1)
\begin{groupe}\trimIamb{ὁ δ᾽ ἀνὴρ περιήρχετ᾽ ὠκυτόκι᾽ ὠνούμενος·
\apparat{504}{\lecon{περιήρχετ᾽} R : περιέτρεχ᾽ Bergk \espApp ὠκυτόκειa R : \lecon{ὠκυτόκι᾽} Pollux II.7 \ligneApp{505} \lecon{τὸ δ᾽} R : τότ᾽ Biset : τῇδ᾽ Daubuz \ligneApp{506} βοῴη suppl. \Rpc}\\
τὸ δ᾽ εἰσέφερε γραῦς ἐν χύτρᾳ τὸ παιδίον,\\
ἵνα μὴ βοῴη, κηρίῳ βεβυσμένον·
\\
εἶθ᾽ ὡς ἔνευσεν ἡ φέρουσ᾽, εὐθὺς βοᾷ,\\
`ἄπελθ᾽ ἄπελθ᾽, ἤδη γὰρ ὦνέρ μοι δοκῶ\\\passeLigne}%
\end{groupe}%
\endnumbering
\end{Leftside}
\begin{Rightside}
\beginnumbering
%% v. 466-519
\begin{groupe}%
\personnageFr{Mnésiloque}%
\trimIambFr{\ligne{466}Que vous soyez, ô femmes, vivement irritées\testimonia{466}{(ὀξυθυμεῖσθαι) Photius ο 380 ὀξυθυμεῖσθαι}\\
contre Euripide, qui si mal parla de vous,\\
n'est pas pour étonner, ni que votre bile déborde.\testimonia{468}{(θαυμάσιον $=$ θαυμαστόν) Moeris θ 16 $\simeq$ Hesychius θ 146 θαυμάσιος}\\
Car moi-même, — ainsi puissé-je être heureuse dans mes enfants, —\testimonia{469}{Suidas ο 330 (sed ἐγώ pro ἔγωγ᾽) ; cf. (ὀναίμην) Hesychius ο 852 ὀναίμην ; \Syn ο 162 ὀναίμην $=$ Photius ο 332}
je hais cet homme, ou je serais folle.\testimonia{470}{Suidas μ 328 (τἄνδρ᾽ pro τὸν ἄνδρ᾽)}\\
Néanmoins il convient que nous nous expliquions entre nous,
car nous sommes seules, et rien ne sera rapporté de ce qui se dit ici.
Pourquoi allons-nous l'accusant\testimonia{473-4}{(τί ταῦτ᾽ -- φέρομεν) Suidas τ 683}
et sommes-nous indignées de ce que, sachant deux ou trois
de nos méfaits, il les a divulgués, quand nous en commettons d'innombrables ?\\
Moi-même toute la première, pour ne rien dire des autres,
j'en ai sur la conscience un grand nombre, mais voici
le plus fort. J'étais mariée de trois jours,}
\end{groupe}%
%% v. 466-519
\begin{groupe}%
\trimIambFr{et mon mari dormait à mes côtés. Or j'avais un amant\\
qui m'avait déflorée à sept ans.\testimonia{480}{Pollux III 42 (οὗτος pro ὅσπερ et διεκόρησεν cod. A pro διεκόρευσεν) \espTest (διεκόρησεν) Hesychius δ 1558 διεκόρησε}\\
Pressé de désirs, il s'en vient gratter à la porte.\testimonia{481}{Suidas ε 520 ; Zonaras 669 ἔκνυεν; \Eustathe 1746.6 (μοὔκνυεν ἕλκων pro μου ᾽κνυεν ἐλθὼν)}\\
Je comprends sur-le-champ et descends à petit bruit.\\
Mon mari m'interroge : « Où descends-tu? — Où ?\\
Je suis prise de tranchées, mon homme, et j'ai mal.\testimonia{484}{\Eustathe 1746.45 \espTest (στρόφος) \Eustathe 503.20}\\
Aussi vais-je aux latrines. — Va donc ».\\
Et là-dessus il se met à broyer des grains de genévrier, de l'aneth, de la sauge.\testimonia{486}{Suidas σ 1712 \espTest (κεδρίδας) Hesychius κ 1980 κεδρίδες ; Photius κ 1980 κ 517}
Pour moi, ayant versé de l'eau sur les gonds,
\ligne{488}j'allai trouver mon galant et subir son assaut
\ligne{488}à côté de l'Aguieus, courbant l'échine et me tenant au laurier.\testimonia{489}{(Ἀγυιᾶ) Stephanus Byzantinus 22.15 ; cf. Moeris α 13}
Cela au grand jamais, voyez-vous, \fin Euripide n'en a parlé.\\
Non plus qu'il ne dit qu'à défaut d'autres,
nous nous faisons hocher par nos esclaves et par des muletiers ;}
\end{groupe}%
%% v. 466-519
\begin{groupe}%
\trimIambFr{ni qu'après avoir paillardé à cœur joie
la nuit, nous mâchons de l'ail le matin,\testimonia{493-4}{Suidas λ 434 Ληκώμεθα: κινώμεθα. Ἀριστοφάνης: μάλισθ' ὅταν ὑπό του ληκώμεθα τὴν νύχθ' ὅλην. Post νύχθ' add ὅλην Suid. (reminiscentia Pherecratis ? cf. Photius λ 254)}
afin que saisi par l'odeur, notre mari, à son retour du rempart,
n'ait aucun soupçon de notre conduite. Cela, voyez-vous,\\
il n'en a jamais parlé. Que s'il invective Phèdre,
que nous importe ? Il n'a pas raconté non plus le cas
de cette femme qui, faisant voir son encycle à son mari,\\
\ligne{501}fit partir son galant ; non, il ne l'a pas raconté.\\
Et j'en sais une autre qui prétendit être en travail
dix jours durant, jusqu'à ce qu'elle eût acheté un enfant.}
\end{groupe}%
%% v. 520-530 (part 1)
\begin{groupe}\trimIambFr{Le mari courait à la ronde faire emplette de drogues pour hâter la délivrance.\testimonia{504}{(ὠκυτόκι᾽) Pollux II 7 ; \Phryn 129.3}
L'enfant fut apporté par une vieille dans une marmite,
la bouche bourrée d'un rayon de miel pour qu'il ne criât point.\testimonia{506}{(κηρίῳ βεβυσμένον) Hesychius κ 2549 κηρίῳ βύσασα ; Photius κ 675 κηρίῳ βῦσαι}
\ligne{506}Puis, sur un signe de la porteuse, la femme aussitôt s'écria :\\
« Va-t'en, va-t'en, mon homme, je crois que je vais\\\passeLigne}%
\end{groupe}%
\endnumbering
\end{Rightside}
\end{pages}
\Pages
\clearpage\ifodd\value{page}\hbox{}\newpage\fi
\enTeteG{ΕΠΕΙΣΟΔΙΟΝ Α' --- Codex Ravennas 429, f. 168v (Classense 170v) --- vv. 509–550}
\begin{adjustwidth}{-1.5cm}{-3.0cm}
\begin{adjustbox}{max totalheight=21.0cm, max width=16.5cm, keepaspectratio}
\begin{minipage}{16.5cm}
{\fontsize{7.2pt}{8.5pt}\selectfont
\noindent\textbf{Folio 13 --- Codex Ravennas 429, \texttt{f. 168v} (\texttt{Classense 170v}) : ΕΠΕΙΣΟΔΙΟΝ Α', vv. 509–550}\hfill \textit{Témoins :} $\text{R}^1, \Rpc, \Rdeux$ (\Mon)\par\vspace{2pt}
\hrule height 0.4pt\vspace{2.5pt}
}
{\fontsize{8.1pt}{9.5pt}\selectfont
\ligneRav{509}{μ\auDessus{ν}{Η}}{τέξειν.' τὸ γὰρ ἦτρον τῆς χύτρας ἐλάκτισεν·}{}
\ligneRav{510}{}{χὠ μὲν γεγηθὼς \difference{ἀπέτρεχ᾽}, ἡ δ᾽ ἐξέσπασεν}{\textbf{510} ἔτρεχεν R : \lecon{ἀπέτρεχ᾽} van Herwerden}
\ligneRav{511}{}{ἐκ τοῦ στόματος τοῦ παιδίου, τὸ δ᾽ ἀνέκραγεν·}{\textbf{511} \lecon{τοῦ παιδίου} R : τὸ κηρίον Hirschig}
\ligneRav{512}{}{τὸ δ᾽ ἐξένευσε κἀνέκραγεν, \difference{ἥ ᾽φερεν} (\auDessus{εἰ}{η}\inRas{σ}{}έφερε)}{\textbf{512} \auDessus{εἰ}{η}\inRas{σ}{}έφερε \Rpc : \lecon{ἥ ᾽φερεν} Bentley}
\ligneRav{}{}{θεῖ μειδιῶσα πρὸς τὸν ἄνδρα καὶ λέγει,}{}
\ligneRav{514}{}{λέων λέων σοι γέγονεν, \difference{αὐτέκμαγμα} (αὖτ᾽ ἔ\inRas{γ}{κ}μαγμα) σόν,}{\textbf{514} αὖτ᾽ ἔκμαγμα (ex έγμαγμα) \Rpc, Pollux IX.131 : ἔκμαγμα Suid. λ 271, π 2114 : \lecon{αὐτέκμαγμα} Reiske}
\ligneRav{515}{}{τά τ᾽ ἄλλ᾽ ἁπαξάπαντα καὶ τὸ πόσθιον}{}
\ligneRav{516}{}{τῷ σῷ προσόμοιον, στρεβλὸν ὥσπερ κύτταρον.'}{\textbf{516} \lecon{κύτταρον} R, Suid. κ 2787 : κύτταρος Fritzsche}
\ligneRav{}{}{ταῦτ᾽ οὐ ποιοῦμεν τὰ κακά; νὴ τὴν Ἄρτεμιν}{}
\ligneRav{}{}{ἡμεῖς γε. κᾆτ᾽ Εὐριπίδῃ θυμούμεθα,}{}
\ligneRav{}{}{οὐδὲν παθοῦσαι μεῖζον ἢ δεδράκαμεν;}{}
\ligneRav{520}{—}{\hspace{1.1em}τουτὶ μέντοι θαυμαστόν,}{}
\ligneRav{521}{}{\hspace{1.1em}ὁπόθεν \difference{ηὑρέθη} τὸ χρῆμα,}{\textbf{521} εὑρέθη R : \lecon{ηὑρέθη} Meineke}
\ligneRav{}{}{\hspace{1.1em}χἤτις ἐξέθρεψε χώρα}{}
\ligneRav{}{}{\hspace{1.1em}τήνδε τὴν θρασεῖαν οὕτω·}{}
\ligneRav{}{}{\hspace{1.1em}τάδε γὰρ εἰπεῖν τὴν πανοῦργον}{}
\ligneRav{525}{}{\hspace{1.1em}κατὰ τὸ φανερὸν ὧδ᾽ ἀναιδῶς}{}
\ligneRav{526}{}{\hspace{1.1em}οὐκ ἂν ᾠόμην ἐν \difference{ἡμῖν}}{\textbf{526} ὑμῖν R : \lecon{ἡμῖν} Bonini}
\ligneRav{526}{}{\hspace{1.1em}οὐδὲ τολμῆσαί ποτ᾽ ἄν·}{\textbf{526} ὑμῖν R : \lecon{ἡμῖν} Bonini}
\ligneRav{527}{}{\hspace{1.1em}\difference{ἀλλὰ πᾶν} γένοιτ᾽ ἂν ἤδη·}{\textbf{527} ἀλλ᾽ ἅπαν R : \lecon{ἀλλὰ πᾶν} Cobet}
\ligneRav{}{}{\hspace{1.1em}τὴν παροιμίαν δ᾽ ἐπαινῶ}{}
\ligneRav{}{}{\hspace{1.1em}τὴν παλαιάν· ὑπὸ λίθῳ γὰρ}{}
\ligneRav{530}{}{\hspace{1.1em}παντί που χρὴ}{}
\ligneRav{530}{}{\hspace{1.1em}μὴ δάκῃ ῥητωρ ἀθρεῖν·}{}
\ligneRav{531}{χορ/}{ἀλλ᾽ οὐ γάρ ἐστι τῶν ἀναισχύτων φύσει γυναικῶν}{\textbf{531} (Κορ.) add. von Velsen}
\ligneRav{532}{}{οὐδὲν κάκιον εἰς ἅπαντα πλὴν ἄρ᾽ εἰ γυναῖκες·}{\textbf{532} \lecon{ἄρ᾽ εἰ} R : ἄλλαι Suid. α 1349}
\ligneRav{533}{γυνή τις}{οὔτοι μὰ τὴν \difference{Ἄγλαυρον} ὦ γυναῖκες εὖ φρονεῖτε,}{\textbf{533} (γυνὴ τις) \Rpc : \lecon{(Μίκα)} Rogers \espApp Ἄγραυλον R : \lecon{Ἄγλαυρον} Daubuz}
\ligneRav{}{}{ἀλλ᾽ ἢ πεφάρμαχθ᾽ ἢ κακόν τι μέγα πεπόνθατ᾽ ἄλλο,}{}
\ligneRav{535}{}{ταύτην ἐῶσαι τὴν φθόρον τοιαῦτα περιυβρίζειν}{}
\ligneRav{536}{}{ἡμᾶς ἁπάσας. εἰ μὲν οὖν τις ἔστιν· εἰ δὲ μή, ἡμεῖς}{\textbf{536} \lecon{μή ἡμεῖς} R : μἠμεῖς Bentley}
\ligneRav{537}{}{\difference{αὐταί} γε καὶ τὰ δουλάρια τέφραν ποθὲν λαβοῦσαι}{\textbf{537} αὐτοί R : \lecon{αὐταί} Bonini \espApp \lecon{γε} R : τε Σ T ad Hom.\italique{Il}.I.135}
\ligneRav{}{}{ταύτης ἀποψιλώσομεν τὸν χοῖρον, ἵνα διδαχθῇ}{}
\ligneRav{}{}{γυνὴ γυναῖκας οὖσα μὴ κακῶς λέγειν τὸ λοιπόν·}{}
\ligneRav{540}{μ\auDessus{ν}{Η}}{μὴ δῆτα τόν γε χοῖρον ὦ γυναῖκες. εἰ γὰρ οὔσης}{}
\ligneRav{}{}{παρρησίας κἀξὸν λέγειν ὅσαι πάρεσμεν ἀσταί,}{}
\ligneRav{542}{}{εἶτ᾽ εἶπον \difference{ἅ ᾽γίγνωσκον} ὑπὲρ Εὐριπίδου δίκαια,}{\textbf{542} ἅ γιγνωσκον R : \lecon{ἅ ᾽γίγνωσκον} Bentley}
\ligneRav{}{}{διὰ τοῦτο τιλλομένην με δεῖ δοῦναι δίκην ὑφ᾽ ὑμῶν;}{}
\ligneRav{}{\auDessus{γυ}{ν}}{οὐ γάρ σε δεῖ δοῦναι δίκην; ἥτις μόνη τέτληκας}{}
\ligneRav{545}{}{ὑπὲρ ἀνδρὸς ἀντειπεῖν, \difference{ὃς} ἡμᾶς πολλὰ κακὰ δέδρακεν}{\textbf{545} ὡς R : \lecon{ὃς} Grynaeus}
\ligneRav{546}{}{\difference{ἐπίτηδες} εὑρίσκων λόγους, ὅπου γυνὴ πονηρὰ}{\textbf{546} ἐξεπίτηδες R : \lecon{ἐπίτηδες} Scaliger}
\ligneRav{}{}{ἐγένετο, Μελανίππας ποιῶν Φαίδρας τε· Πηνελόπην δὲ}{}
\ligneRav{548}{}{οὐπώποτ᾽ \difference{ἐποίησ᾽}, ὅτι γυνὴ σώφρων ἔδοξεν εἶναι·}{\textbf{548} ἐποίησεν R : \lecon{ἐποίησ᾽} Scaliger}
\ligneRav{549}{μ\auDessus{ν}{Η}}{ἀλλ᾽ οὖν οὕνεκα τούτου γ᾽ \inRas{ἂρ}{ἂν} ὤφελες δοῦναι δίκην·}{\textbf{549} ἂν ex ἂρ corr. \Rpc}
\ligneRav{550}{}{τῶν νῦν γυναικῶν Πηνελόπην, Φαίδρας \difference{δ᾽ ἁπαξαπάσας}·}{\textbf{550} ἁπαξαπάσας R, M\expo{pc} : \lecon{δ᾽ ἁπαξαπάσας} M\expo{ac}}
}
\vspace{3pt}\noindent\rule{4.5cm}{0.4pt}\par\vspace{1.5pt}
{\fontsize{7.0pt}{8.3pt}\selectfont
\noindent\textbf{Loci paralleli (vv. 509–550) ---} \textbf{509} (τὸ γὰρ ἦτρον -- ἐλάκτισεν) Suidas η 637 ; Zonaras 1012 ἦτρον \espTest (ἦτρον) Moeris η 14 \espTest \textbf{514-6} Suid. λ 271 $\simeq$ Suid. π 2114 \espTest \textbf{514} Pollux IX 131 (λέων\expo{2} omisso et αὐτὸ ἔκμαγμα pro αὐτέκμαγμα) ; \Eustathe 1857.10, 462.40 ; cf. \AelDion ε 26 \espTest \textbf{515-6} (πόσθιον -- κύτταρον) Suidas κ 2787 \espTest \textbf{516} (στρεβλὸν) Hesychius σ 1977 στρεβλὸν \espTest{(κύτταρον)} Hesychius κ 4747 ; cf. Photius κ 1270 \espTest \textbf{527-30} Suidas α 1073 \espTest \textbf{528-30} \Stobee III 13.29 \espTest \textbf{531-2} Suidas α 1349 (sed ἄλλαι pro ἄρ᾽ εἰ) \espTest \textbf{536-7} (εἰ μὲν -- λαβοῦσαι) scholium\expo{bT} Homerus \italique{Il.} I 135-7 (ἔστι τις pro τις ἔστιν et αὐταί omisso)\par
}
\vspace{2.5pt}\noindent\rule{4.5cm}{0.4pt}\par\vspace{1.5pt}
{\fontsize{6.9pt}{8.2pt}\selectfont
\noindent\textbf{Suite (en commentaire) de l'apparat critique \& notes paléographiques (\texttt{f. 168v}) ---} \textbf{Correction complexe \italique{in rasura} et supralinéaire au v. 512 :} \Rac{} avait répété εἰσέφερε d'après le v. 509 ; \Rpc{} a gratté le $\sigma$ et tracé un $\eta$ au-dessus de εἰ (\lecon{\auDessus{εἰ}{η}\inRas{σ}{}έφερε}), visant ἣ 'φερεν (Bentley). --- \textbf{Correction au v. 514 :} αὖτ' ἔκμαγμα corrigé par \Rpc{} à partir de έγμαγμα. --- \textbf{Correction au v. 549 :} ἂν corrigé par \Rpc{} à partir de ἂρ. --- \textbf{v. 533 :} (γυνὴ τις) \Rpc{} : \lecon{(Μίκα)} Rogers \espApp Ἄγραυλον R (métathèse de liquides $\rho/\lambda$) : \lecon{Ἄγλαυρον} Daubuz. --- \textbf{v. 537 :} αὐτοί R : \lecon{αὐταί} Bonini \espApp \lecon{γε} R (τε $\Sigma$ T ad Hom. \italique{Il.} I.135).\par
}
\end{minipage}
\end{adjustbox}
\end{adjustwidth}
\newpage
\enTeteD{Codex Ravennas 429 (Biblioteca Classense) --- Fac-similé du f. 168v (vv. 509–550)}
\begin{adjustwidth}{-1.5cm}{-3.0cm}
\begin{center}
\includegraphics[width=16.5cm,height=20.5cm,keepaspectratio]{/workspaces/Thesmophories/dossier_ravennas_thesmophories/folios_individuels_ff162v-177r/13_f168v_classense170v_vv0509-0550.jpg}\par
\vspace{2pt}
{\scriptsize\textbf{Codex Ravennas 429, f. 168v (Classense 170v)} --- ΕΠΕΙΣΟΔΙΟΝ Α', \textbf{vv. 509–550} ($\text{R}^1, \Rpc, \Rdeux$ (\Mon)).}
\end{center}
\end{adjustwidth}
\newpage
\clearpage\ifodd\value{page}\hbox{}\newpage\fi
\enTeteG{ΘΕΣΜΟΦΟΡΙΑΖΟΥΣΑΙ --- Texte grec établi (f. 168v, vv. 509–550)}
\enTeteD{LES THESMOPHORIES --- Traduction française (f. 168v, vv. 509–550)}
\begin{pages}
\begin{Leftside}
\beginnumbering
%% v. 509+ (part 2)
\begin{groupe}
\trimIamb{\ligne{509}τέξειν.' τὸ γὰρ ἦτρον τῆς χύτρας ἐλάκτισεν·%\testimonia{509}{(τὸ γὰρ ἦτρον -- ἐλάκτισεν) Suidas η 637 ; Zonaras 1012 ἦτρον \espTest (ἦτρον) Moeris η 14}%
\commentaire{509}{τέξειν.}{Cf. sur le mot τέξειν la note 14 de Willems à lire dans les annexes.}
\commentaire{509}{τὸ γὰρ ἦτρον τῆς χύτρας ἐλάκτισεν·}{Il y a une plaisanterie \italique{contre l'attente} (παρὰ προσδοκίαν) : τῆς χύτρας est mis pour τῆς μητρός.}\\
χὠ μὲν γεγηθὼς ἀπέτρεχ᾽, ἡ δ᾽ ἐξέσπασεν\apparat{510}{ἔτρεχεν R : \lecon{ἀπέτρεχ᾽} van Herwerden}\\
ἐκ τοῦ στόματος τοῦ παιδίου, τὸ δ᾽ ἀνέκραγεν.\\%\apparat{511}{\lecon{τοῦ παιδίου} R : τὸ κηρίον Hirschig}\\
εἶθ᾽ ἡ μιαρὰ γραῦς, ἣ ᾽φερεν τὸ παιδίον,\apparat{512}{\auDessus{εἰ}{η}\inRas{σ}{}έφερε \Rpc : \lecon{ἥ ᾽φερεν} Bentley}\\
θεῖ μειδιῶσα πρὸς τὸν ἄνδρα καὶ λέγει,\\
`λέων λέων σοι γέγονεν, αὐτέκμαγμα σόν,%\testimonia{514-6}{Suid. λ 271 $\simeq$ Suid. π 2114}\testimonia{514}{Pollux IX 131 (λέων\expo{2} omisso et αὐτὸ ἔκμαγμα pro αὐτέκμαγμα) ; \Eustathe 1857.10, 462.40 ; cf. \AelDion ε 26}
\apparat{514}{αὖτ᾽ ἔκμαγμα (ex έγμαγμα) \Rpc, Pollux IX.131 : ἔκμαγμα Suid. λ 271, π 2114 : \lecon{αὐτέκμαγμα} Reiske}\\% \ligneApp{516} \lecon{κύτταρον} R, Suid. κ 2787 : κύτταρος Fritzsche}\\
τά τ᾽ ἄλλ᾽ ἁπαξάπαντα καὶ τὸ πόσθιον%\testimonia{515-6}{(πόσθιον -- κύτταρον) Suidas κ 2787}%
\commentaire{515}{ἁπαξάπαντα.}{La Fontaine, \italique{La chatte métamorphosée en femme} (II, 18) :%
\citComm{\francais{La croit femme en tout et partout.}}%
}\\
τῷ σῷ προσόμοιον, στρεβλὸν ὥσπερ κύτταρον.'%\testimonia{516}{(στρεβλὸν) Hesychius σ 1977 στρεβλὸν \espTest{(κύτταρον)} Hesychius κ 4747 ; cf. Photius κ 1270}%\apparat{516}{\lecon{κύτταρον} R, Suid. κ 2787 : κύτταρος Fritzsche}%
\commentaire{516}{κύτταρον,}{synonyme de κῶνος, \italique{pomme de pin}, que les Grecs appelaient aussi στρόβιλος (στρόμβος dans Nicandre. Hehn., p.292). (Note de Willems.)}\\
ταῦτ᾽ οὐ ποιοῦμεν τὰ κακά; νὴ τὴν Ἄρτεμιν\\
ἡμεῖς γε. κᾆτ᾽ Εὐριπίδῃ θυμούμεθα,\\
οὐδὲν παθοῦσαι μεῖζον ἢ δεδράκαμεν;
\commentaire{518-9}{θυμούμεθα, / οὐδὲν παθοῦσαι μεῖζον ἢ δεδράκαμεν;}{On sait grâce au scholiaste que ces vers sont tirés du \italique{Télèphe} : 
\citComm{\grec{παρὰ τὰ ἐκ Τηλέφου Εὐριπίδου· `εἰ δὴ θυμούμεθα παθόντες οὐδὲν μᾶλλον ἢ δεδρακότες'\marqueur{1}.}}
(fragment 20 du \italique{Télèphe} de l'édition Jouan-van Looy aux Belles Lettres, tome VIII, \nieme{3} partie. Correspond aux fragments 711 de l'édition Kannicht et 965 de celle de Mette.)}%
\debutNote{518-9}%
%\note{1}{fragment 20 du \italique{Télèphe} de l'édition Jouan-van Looy aux Belles Lettres, tome VIII, \nieme{3} partie. Correspond aux fragments 711 de l'édition Kannicht et 965 de celle de Mette.}
\\\passeLigne}%
\end{groupe}%
%% v. 520-530
\begin{groupe}%
\personnageGrec{Χορός\strophe{ἀντιστρ.}}%
\typeVers{par.}\chantChoeur{τουτὶ μέντοι θαυμαστόν,\commentaire{520-30}{}{Antistrophe du chant du chœur des vers 433-442.}}%
\end{groupe}%
%% v. 520-530
\begin{groupe}%
\typeVers{$2$ troch.}\chantChoeur{ὁπόθεν ηὑρέθη τὸ χρῆμα,\apparat{521}{εὑρέθη R : \lecon{ηὑρέθη} Meineke}}%
\end{groupe}%
%% v. 520-530
\begin{groupe}%
\typeVers{$2$ troch.}\chantChoeur{χἤτις ἐξέθρεψε χώρα}%
\end{groupe}%
%% v. 520-530
\begin{groupe}%
\typeVers{$2$ troch.}\chantChoeur{τήνδε τὴν θρασεῖαν οὕτω.}%
\end{groupe}%
%% v. 520-530
\begin{groupe}%
\typeVers{$2$ troch.}\chantChoeur{τάδε γὰρ εἰπεῖν τὴν πανοῦργον}%
\end{groupe}%
%% v. 520-530
\begin{groupe}%
\typeVers{$2$ troch.}\chantChoeur{κατὰ τὸ φανερὸν ὧδ᾽ ἀναιδῶς}%
\end{groupe}%
%% v. 526a
\begin{groupe}%
\decoupeVers{a}\typeVers{$2$ troch.}\chantChoeur{οὐκ ἂν ᾠόμην ἐν ἡμῖν\apparat{526}{ὑμῖν R : \lecon{ἡμῖν} Bonini}}%
\end{groupe}%
%% v. 526b
\begin{groupe}%
\ligne{526}\decoupeVers{b}\typeVers{$2$ troch\cat}\chantChoeur{οὐδὲ τολμῆσαί ποτ᾽ ἄν.}%
\end{groupe}%
%% v. 527
\begin{groupe}%
\typeVers{$2$ troch.}\chantChoeur{ἀλλὰ πᾶν γένοιτ᾽ ἂν ἤδη·%\testimonia{527-30}{Suidas α 1073}
\apparat{527}{ἀλλ᾽ ἅπαν R : \lecon{ἀλλὰ πᾶν} Cobet}%
\commentaire{527}{ἀλλ᾽ ἅπαν γένοιτ᾽ ἂν ἤδη·}{Le lieu-commun du \latin{omnia mutantur}\marqueur{1} avait bien-sûr cours aussi chez les Grecs. 
Brève recension : \\
Hérodote (V, 9) :%
\citComm{γένοιτο δ̓ ἂν πᾶν ἐν τῷ μακρῷ χρόνῳ.}%
Sophocle (\italique{Philoctète}, vers 305-306) : %
\citComm{πολλὰ γὰρ τάδε / ἐν τῷ μακρῷ γένοιτ᾽ ἂν ἀνθρώπων χρόνῳ.}%
Eupolis\marqueur{2} :%
\citComm{ἦ πολλά γ' ἐν μακρῷ χρόνῳ γίγνεται μεταλλαγῇ / (τῶν) πραγμάτων· μένει δὲ χρῆμ' οὐδὲν ἐν ταὐτῷ ῥυθμῷ.}%
Ménandre (fragment 685(466)\marqueur{3}) : 
\citComm{οὐκ ἔστ' ἄπιστον οὐδὲν ἐν θνητῷ βίῳ / οὐδ' ἂν γένοιτο.}
Et dans la pièce d'Aristophane d'antan, \italique{Lysistrata} (vers 258) : 
\citComm{Χορὸς γερόντων --- ἦ πόλλ᾽ ἄελπτ᾽ ἔνεστιν ἐν τῷ μακρῷ βίῳ φεῦ.}
Remarquons que dans tous ces extraits, le grec insiste sur la longueur de temps que prennent les changements pour s'opérer (ἐν τῷ μακρῷ χρόνῳ / ἐν θνητῷ βίῳ / ἐν τῷ μακρῷ βίῳ)... or c'est tout le contraire ici : ἤδη, \italique{aujourd'hui}, \italique{actuellement}. Ce contraste est certainement une allusion d'Aristophane aux temps mouvementés politiquement de la représentation des \italique{Thesmophories} -- où tout semble s'accélérer.}%
\debutNote{527}%
\note{1}{Cf. Tosi, \italique{Dictionnaire des sentences latines et grecques}, §26. La citation est d'Ovide (\italique{Métamorphoses}, XV, 165). Érasme (\italique{Adages}, 663) recense la formule équivalente de Térence : \latin{Omnium rerum vicissitudo est}.}%
\note{2}{Fragment 391(356) des \italique{Poetae Comici Graeci}, ed. Kassel \& Austin, V (1986).}%
\note{3}{\italique{Poetae Comici Graeci}, VI 2 (1998). (Le fragment est extrait de Stobée I, 8, 31.)}}%
\end{groupe}%
%% v. 520-5308
\begin{groupe}%
\typeVers{$2$ troch.}\chantChoeur{τὴν παροιμίαν δ᾽ ἐπαινῶ%\testimonia{528-30}{\Stobee III 13.29}
}%
\end{groupe}%
%% v. 529
\begin{groupe}%
\typeVers{$2$ troch.}\chantChoeur{τὴν παλαιάν· ὑπὸ λίθῳ γὰρ}%
\end{groupe}%
%% v. 530a
\begin{groupe}%
\decoupeVers{a}\typeVers{troch.}\chantChoeur{παντί που χρὴ}%
\end{groupe}%
%% v. 530b
\begin{groupe}%
\ligne{530}\decoupeVers{b}\typeVers{$2$ troch\cat}\chantChoeur{μὴ δάκῃ ῥητωρ ἀθρεῖν.
\commentaire{528-30}{τὴν παροιμίαν δ᾽ ἐπαινῶ / τὴν παλαιάν· ὑπὸ λίθῳ γὰρ / παντί που χρὴ / μὴ δάκῃ ῥητωρ ἀθρεῖν.}{Le proverbe – avec un scorpion à la place de l'orateur – était effectivement populaire puisqu'il nous est diversement parvenu dans les recueils de Proverbes. Ainsi des parémiographes Zénobie\marqueur{1} ou Diogénianus\marqueur{2} et d'Hésychius\marqueur{3} qui rapportent la version la plus courante : 
\citComm{ὑπὸ παντὶ λίθῳ σκόρπιος εὕδει: ἐπὶ τῶν κακοήθων λέγεται.}%
Macaire\marqueur{4} en précise le sens : \grec{ἐπὶ πολυπληθίᾳ τῶν πανταχοῦ ὄντων πονηρῶν ἀνθρώπων}, « rapport au grand nombre d'hommes qui sont méchants sous tout rapport ».
Le proverbe est attesté aussi par Praxilla\marqueur{5} et Sophocle\marqueur{6}.\\
Ici, la \italique{pierre} (λίθος) dont parle la coryphée est le piédestal (βῆμα) de la Pnyx sur lequel les orateurs -- Mnésiloque, en l'occurrence -- parlaient. (Aristophane use du mot dans le même sens aux vers 683 des \italique{Acharniens} ; 680 de la \italique{Paix} et 87 de l'\italique{Assemblée des femmes}.)
Hésychius (λ 999) :%
\citComm{λίθος· βῶλος. βωμός. καὶ βάσις. τὸ ἐν τῇ Ἀθηναίων ἐκκλησίᾳ βῆμα.}%
Enfin, la substitution au scorpion de l'orateur n'est pas gratuite : on parlait couramment de la \italique{morsure de la langue}\marqueur{7}.
Ainsi de la \italique{morsure du sycophante}, δῆγμα συκοφάντου, qui paraît proverbiale\marqueur{8}. Diogène le cynique eut cette réplique à l'ironie \italique{mordante} (rapportée par son homonyme laërtien) :
\citComm{ἐρωτηθεὶς τί τῶν θηρίων κάκιστα δάκνει ἔφη: τῶν μὲν ἀγρίων συκοφάντης, τῶν δὲ ἡμέρων κόλαξ.\\
Interrogé sur la bête dont la morsure est la plus mauvaise, [il répondit] que des bêtes sauvages c'était le sycophante, et des animaux domestiques, c'était le flatteur.}}%
\debutNote{528-30}%
\note{1}{\latin{Zenobii centuria} VI, 20 (tomus I). Je cite les proverbes d'après le recensement d'E. Leutsch et F. G. Schneidewin, \italique{Corpus paroemiographorum Graecorum}, en deux tomes. Signalons qu'une nouvelle édition critique des recensions de Zénobius et Diogénianus avec traduction en italien a paru en 2019 sous le titre \italique{I proverbi greci}, dir. Emanuelle Lelli.}%
\note{2}{\latin{Diogeniani centuria} VIII, 59 (tomus I).}%
\note{3}{Hésychius (υ 717). Photius et Suidas rapportent aussi ce proverbe : il est partout !}%
\note{4}{\latin{Macarii centuria} VIII, 69 (tomus II).}%
\note{5}{Le scholiaste attribue même le proverbe à Praxilla. Athénée (XV, 695d) donne un exemple de scolie (σκόλιον), sorte de chanson à boire, écrite sur ce proverbe ; comme Praxilla excella dans ce genre, il n'est pas impossible que la scolie fût d'elle, et que le scholiaste l'ait en tête lorsqu'il attribue un peu hâtivement la maternité du proverbe à Praxilla. Cette scolie correspond au fragment 903 des \italique{Poetae Melici Graeci}.}%
\note{6}{Fragment 37 de l'édition Radt.}%
\note{7}{Voir les sens figurés du verbe δάκνειν.}%
\note{8}{Voir le \italique{Ploutos}, vers 885, et le commentaire du scholiaste \italique{ad. loc.}}\\\passeLigne}%
\end{groupe}%
%% v. 531-532
\begin{groupe}%
\personnageGrec{Κορυφαῖος}%
\typeVers{$4$ia\cat}\tetramIambCat{\ligne{531}ἀλλ᾽ οὐ γάρ ἐστι τῶν ἀναισχύτων φύσει γυναικῶν%\testimonia{531-2}{Suidas α 1349 (sed ἄλλαι pro ἄρ᾽ εἰ)}
\apparat{531}{(Κορ.) add. von Velsen}%
\commentaire{531-73}{}{Tout ce passage est composé en tétramètres iambiques catalectiques ; mètre du conflit et des débats animés. Cf. aussi les \italique{Cavaliers} v. 335-366 ; \italique{Lysistrata} v. 350-381, 467-475.}\\
%
οὐδὲν κάκιον εἰς ἅπαντα πλὴν ἄρ᾽ εἰ γυναῖκες.\apparat{532}{\lecon{ἄρ᾽ εἰ} R : ἄλλαι Suid. α 1349}%
\commentaire{531-2}{ἀλλ᾽ οὐ γάρ ἐστι τῶν ἀναισχύτων φύσει γυναικῶν / οὐδὲν κάκιον εἰς ἅπαντα πλὴν ἄρ᾽ εἰ γυναῖκες.}{Oct. Feuillet, \italique{Dalila}, partie IV :
\citComm{\francais{Rien de plus semblable sous le soleil qu'une femme – et une femme.}}%
C'est la traduction presque textuelle de Ménandre (\italique{Sentences monostiques}, Γνῶμαι μονόστιχοι 109) :%
\citComm{Γυνὴ γυναικὸς πώποτ᾽ οὐδὲν διαφέρει.}%
(Note de Willems ; pour voir ce que Willems pense de la « manie des conjectures » visant à retoucher le texte d'Aristophane, on lira avec un plaisir non dissimulé la note entière n° 15, reproduite à la fin.)\\
Il serait piquant qu'Aristophane mette dans la bouche de la coryphée une citation d'Euripide lui-même pour dénigrer le comportement de certaines femmes ; le scholiaste n'en a malheureusement pas soufflé mot... on dispose néanmoins de citations d'Euripide très approchantes (fr. 494.27 Kannicht $=$ fr. 14.27 \italique{Mélanippe} Jouan-van Looy) : %
\citComm{τῆς μὲν κακῆς κάκι[ο]ν οὐδὲν γίγνεται / γυναικός.}%
}%
\\\passeLigne}%
\end{groupe}%
%% v. 533-539
\begin{groupe}%
%
\personnageGrec{Μίκα}%
\tetramIambCat{οὔτοι μὰ τὴν Ἄγλαυρον ὦ γυναῖκες εὖ φρονεῖτε,\apparat{533}{(γυνὴ τις) \Rpc : \lecon{(Μίκα)} Rogers \espApp Ἄγραυλον R : \lecon{Ἄγλαυρον} Daubuz}%
\commentaire{533}{μὰ τὴν Ἄγραυλον.}{Scholie d'un érudit :%
\citComm{κατὰ τῆς ᾿Αγραύλου ὦμνυον, κατὰ δὲ τῆς Πανδρόσου σπανιώτερον· κατὰ δὲ τῆς Ἕρσης οὐχ εὑρήχαμεν.\\
\francais{L'on jurait sur Agraule, plus rarement sur Pandrosos ; nous n'avons pas trouvé d'exemple sur Ersès.}} %
Aglaure (correction d'Agraule confirmée par l'épigraphie), Pandrose et Hersé étaient les trois filles de Kécrops, mythique fondateur et premier roi d'Athènes.}\\
\small{ἀλλ᾽ ἢ πεφάρμαχθ᾽ ἢ κακόν τι μέγα πεπόνθατ᾽ ἄλλο,}\\
ταύτην ἐῶσαι τὴν φθόρον τοιαῦτα περιυβρίζειν\\
ἡμᾶς ἁπάσας. εἰ μὲν οὖν τις ἔστιν· εἰ δὲ μή, ἡμεῖς%\testimonia{536-7}{(εἰ μὲν -- λαβοῦσαι) scholium\expo{bT} Homerus \italique{Il.} I 135-7 (ἔστι τις pro τις ἔστιν et αὐταί omisso)}
\\%\apparat{536}{\lecon{μή ἡμεῖς} R : μἠμεῖς Bentley}\\
αὐταί γε καὶ τὰ δουλάρια τέφραν ποθὲν λαβοῦσαι\apparat{537}{αὐτοί R : \lecon{αὐταί} Bonini \espApp \lecon{γε} R : τε Σ T ad Hom.\italique{Il}.I.135 %, Reiske \espApp πόθεν R : \lecon{ποθὲν} Küster
}\\
ταύτης ἀποψιλώσομεν τὸν χοῖρον, ἵνα διδαχθῇ%
%\commentaire{537-8}{τέφραν ποθὲν λαβοῦσαι / ταύτης ἀποψιλώσομεν τὸν χοῖρον}{.}
\\
γυνὴ γυναῖκας οὖσα μὴ κακῶς λέγειν τὸ λοιπόν.\\\passeLigne}%
\end{groupe}%
%% v. 540-543
\begin{groupe}%
\personnageGrec{Μνησίλοχος}%
\tetramIambCat{μὴ δῆτα τόν γε χοῖρον ὦ γυναῖκες. εἰ γὰρ οὔσης\\
παρρησίας κἀξὸν λέγειν ὅσαι πάρεσμεν ἀσταί,\\
εἶτ᾽ εἶπον ἅ ᾽γίγνωσκον ὑπὲρ Εὐριπίδου δίκαια,\\%\apparat{542}{ἅ γιγνωσκον R : \lecon{ἅ ᾽γίγνωσκον} Bentley}\\
διὰ τοῦτο τιλλομένην με δεῖ δοῦναι δίκην ὑφ᾽ ὑμῶν;\\\passeLigne}%
\end{groupe}%
%% v. 544-548
\begin{groupe}%
\personnageGrec{Μίκα}%
\tetramIambCat{οὐ γάρ σε δεῖ δοῦναι δίκην; ἥτις μόνη τέτληκας\\
\small{ὑπὲρ ἀνδρὸς ἀντειπεῖν, ὃς ἡμᾶς πολλὰ κακὰ δέδρακεν}\apparat{545}{ὡς R : \lecon{ὃς} Grynaeus}\\% (\italique{qui} Divus)}\\
ἐπίτηδες εὑρίσκων λόγους, ὅπου γυνὴ πονηρὰ\apparat{546}{ἐξεπίτηδες R : \lecon{ἐπίτηδες} Scaliger}\\
\small{ἐγένετο, Μελανίππας ποιῶν Φαίδρας τε· Πηνελόπην δὲ}\\
οὐπώποτ᾽ ἐποίησ᾽, ὅτι γυνὴ σώφρων ἔδοξεν εἶναι.\apparat{548}{ἐποίησεν R : \lecon{ἐποίησ᾽} Scaliger}\\\passeLigne}%
\end{groupe}%
%% v. 549-550
\begin{groupe}%
\personnageGrec{Μνησίλοχος}%
\tetramIambCat{\ligne{549}ἐγὼ γὰρ οἶδα ταἴτιον. Μίαν γὰρ οὐκ ἂν εἴποις\apparat{549}{ἂν ex ἂρ corr. \Rpc}%
\commentaire{549-50}{Μίαν γὰρ οὐκ ἂν εἴποις / τῶν νῦν γυναικῶν Πηνελόπην, Φαίδρας ἁπαξαπάσας.}{Le comique Euboulos paraît s'être inspiré de ce passage :%
\citComm{ὦ Ζεῦ πολυτίμητ', εἶτ' ἐγὼ κακῶς ποτε\\
ἐρῶ γυναῖκας ; νὴ Δί' ἀπολοίμην ἄρα\\
πάντων ἄριστον κτημάτων. Εἰ δ' ἐγένετο\\
κακὴ γυνὴ Μήδεια, Πηνελόπη δέ <γε>\\
μέγα πρᾶγμ'. ᾿Ερεῖ τις ὡς Κλυταιμνήστρα κακή·\\
Ἄλκηστιν ἀντέθηκα χρηστήν. ᾿Αλλ' ἴσως\\
Φαίδραν ἐρεῖ κακῶς τις · ἀλλὰ νὴ Δία\\
χρηστή τις ἦν μέντοι - τίς ; οἴμοι δείλαιος,\\
ταχέως γέ μ' αἱ χρησταὶ γυναῖκες ἐπέλιπον,\\
τῶν δ' αὖ πονηρῶν ἔτι λέγειν πολλὰς ἔχω.\\
\francais{Zeus très honoré, moi, dire du mal des femmes ? Ah! plutôt mourir ! Des femmes, le plus précieux de tous les biens ! Car si Médée s'est montrée méchante femme, Pénélope, elle, fut un vrai trésor. On dira que Clytemnestre fut mauvaise ? Je mets en regard la vertueuse Alceste. Mais peut-être médira-t-on de Phèdre ? Eh, par Zeus, vertueuse sans contredit fut.... qui donc ? Ah! malheureux ! les femmes vertueuses m'ont de suite fait faute, quand des mauvaises j'en ai encore un tas à citer.}}%
(Meineke, \italique{Fragm. Com. Gr}., t. III, p. 260. [$=$ fragment 115 (116. 117) de l'édition Kassel-Austin, \italique{Poet. Com. Gr.}]) (Note de Willems.)}
\\
\small{τῶν νῦν γυναικῶν Πηνελόπην, Φαίδρας δ᾽ ἁπαξαπάσας.}\apparat{550}{ἁπαξαπάσας R, M\expo{pc} : \lecon{δ᾽ ἁπαξαπάσας} M\expo{ac}}%
\commentaire{550}{Πηνελόπην, Φαίδρας.}{Si Phèdre est souvent citée par Aristophane comme (anti-)modèle d'impudicité (voir le vers 153 ; les \italique{Grenouilles} v. 1043 ; fr. 469(453) K.-A.(Kock)), Mélanippe apparaît seulement ici dans les œuvres conservées d'Aristophane.%
%\eff{développer sur Mélanippe ?}
}%
\\\passeLigne}%
\end{groupe}%
\endnumbering
\end{Leftside}
\begin{Rightside}
\beginnumbering
%% v. 509+ (part 2)
\begin{groupe}
\trimIambFr{\ligne{509}accoucher ». Car il avait donné des pieds contre la panse de la marmite.\testimonia{509}{(τὸ γὰρ ἦτρον -- ἐλάκτισεν) Suidas η 637 ; Zonaras 1012 ἦτρον \espTest (ἦτρον) Moeris η 14}
Tandis qu'il s'éloigne tout joyeux, elle débarrasse
la bouche de l'enfant qui se met à crier.\\
La vieille canaille qui l'a apporté\\
\ligne{513}court en souriant à l'époux, et lui dit :\\
« un lion, un lion t'est né, ton portrait moulé ;\testimonia{514-6}{Suid. λ 271 $\simeq$ Suid. π 2114}\testimonia{514}{Pollux IX 131 (λέων\expo{2} omisso et αὐτὸ ἔκμαγμα pro αὐτέκμαγμα) ; \Eustathe 1857.10, 462.40 ; cf. \AelDion ε 26}\\
outre tout le reste absolument, son membre\testimonia{515-6}{(πόσθιον -- κύτταρον) Suidas κ 2787}\\
est pareil au tien et tortu comme une pigne ».\testimonia{516}{(στρεβλὸν) Hesychius σ 1977 στρεβλὸν \espTest{(κύτταρον)} Hesychius κ 4747 ; cf. Photius κ 1270}\\
Ne sont-ce pas là de nos tours? Si fait, par Artémis,\\
c'en sont. Et nous nous fâchons contre Euripide,\\
« sans avoir enduré par-delà nos mérites ».;\\\passeLigne}%
\end{groupe}%
%% v. 520-530
\begin{groupe}%
\personnageFr{Chœur}%
\chantChoeurFr{Une chose me confond :}%
\end{groupe}%
%% v. 520-530
\begin{groupe}%
\chantChoeurFr{où a-t-on trouvé cette créature ?}%
\end{groupe}%
%% v. 520-530
\begin{groupe}%
\chantChoeurFr{Quel pays a formé}%
\end{groupe}%
%% v. 520-530
\begin{groupe}%
\chantChoeurFr{une telle effrontée ?}%
\end{groupe}%
%% v. 520-530
\begin{groupe}%
\chantChoeurFr{La coquine ! tenir ce langage}%
\end{groupe}%
%% v. 520-530
\begin{groupe}%
\chantChoeurFr{au grand jour, avec cette impudence ?}%
\end{groupe}%
%% v. 520-530
\begin{groupe}%
\chantChoeurFr{je n'aurais pas cru que parmi nous}%
\end{groupe}%
%% v. 520-530
\begin{groupe}%
\chantChoeurFr{elle l'eût jamais osé.}%
\end{groupe}%
%% v. 520-530
\begin{groupe}%
\chantChoeurFr{Mais maintenant tout arrive,}%
\end{groupe}%
%% v. 520-530
\begin{groupe}%
\chantChoeurFr{et j'approuve l'antique proverbe,\testimonia{527-30}{Suidas α 1073}}%
\end{groupe}%
%% v. 520-530
\begin{groupe}%
\chantChoeurFr{quasiment sous chaque pierre,\testimonia{528-30}{\Stobee III 13.29}%\commentaire{}{}{Euripide fait allusion à ce proverbe, Héraclès, 1002 : « Il n'y a pas de pierre sous laquelle ne se glisse un scorpion ».}
}%
\end{groupe}%
%% v. 520-530
\begin{groupe}%
\chantChoeurFr{il faut regarder}%
\end{groupe}%
%% v. 520-530
\begin{groupe}%
\chantChoeurFr{de crainte d'être mordu... par un orateur..\\\passeLigne}%
\end{groupe}%
%% v. 531-532
\begin{groupe}%
\personnageFr{Coryphée}%
\tetramIambCatFr{\ligne{531}Mais les femmes au naturel impudent,\testimonia{531-2}{Suidas α 1349 (sed ἄλλαι pro ἄρ᾽ εἰ)}
il n'est rien de pire à tous égards, si ce n'est des femmes.%\commentaire{531-2}{}{Cf. sur ce passage : Bull. cité, note 15, ci-après.}.
\\\passeLigne}%
\end{groupe}%
%% v. 533-539
\begin{groupe}%
\personnageFr{Mika}%
\tetramIambCatFr{Non, par Aglaure, ô femmes, vous n'êtes pas dans votre bon sens,
mais vous êtes maléficiées, ou vous avez pris quelque autre mal,
pour que vous laissiez cette peste nous braver toutes à outrance.
(Elle se tourne vers le bureau de l'assemblée.) Ainsi qu'on se décide...\testimonia{536-7}{(εἰ μὲν -- λαβοῦσαι) scholium\expo{bT} Homerus \italique{Il.} I 135-7 (ἔστι τις pro τις ἔστιν et αὐταί omisso)}
sinon nous-mêmes et nos esclaves\\
\ligne{537}irons chercher quelque part de la cendre,\\
et lui épilerons le bas-ventre, afin que, femme, elle apprenne
à ne plus médire des femmes à l'avenir.\\\passeLigne}%
\end{groupe}%
%% v. 540-543
\begin{groupe}%
\personnageFr{Mnésiloque}%
\tetramIambCatFr{Le bas-ventre ? Non pas, ô femmes. Car la franchise
étant de droit et nous toutes citoyennes ci-présentes ayant congé de parler,
si j'ai déduit les raisons que je savais en faveur d'Euripide,
\ligne{543}est-ce un motif pour m'épiler, afin de m'en faire porter la peine ?\\\passeLigne}%
\end{groupe}%
%% v. 544-548
\begin{groupe}%
\personnageFr{Mika}%
\tetramIambCatFr{\ligne{544}Quoi! tu n'en porterais pas la peine ? toi qui seule osas
répliquer en faveur d'un homme qui nous a tant malmenées,
trouvant à plaisir des sujets où une femme se montra mauvaise,
produisant des Mélanippes, des Phèdres ; mais jamais il ne produisit une Pénélope,
vu qu'elle passait pour sage.\\\passeLigne}%
\end{groupe}%
%% v. 549-550
\begin{groupe}%
\personnageFr{Mnésiloque}%
\tetramIambCatFr{\ligne{549}J'en sais bien la raison. Tu n'en citerais pas une seule,\\
parmi les femmes d'aujourd'hui, de Pénélope. Toutes \fin Phèdres, sans exception.\\\passeLigne}%
\end{groupe}%
\endnumbering
\end{Rightside}
\end{pages}
\Pages
\clearpage\ifodd\value{page}\hbox{}\newpage\fi
\enTeteG{ΕΠΕΙΣΟΔΙΟΝ Α' --- Codex Ravennas 429, f. 169r (Classense 171r) --- vv. 551–589}
\begin{adjustwidth}{-1.5cm}{-3.0cm}
\begin{adjustbox}{max totalheight=21.0cm, max width=16.5cm, keepaspectratio}
\begin{minipage}{16.5cm}
{\fontsize{7.2pt}{8.5pt}\selectfont
\noindent\textbf{Folio 14 --- Codex Ravennas 429, \texttt{f. 169r} (\texttt{Classense 171r}) : ΕΠΕΙΣΟΔΙΟΝ Α', vv. 551–589}\hfill \textit{Témoins :} $\text{R}^1, \Rpc, \Rdeux$ (\Mon)\par\vspace{2pt}
\hrule height 0.4pt\vspace{2.5pt}
}
{\fontsize{8.1pt}{9.5pt}\selectfont
\ligneRav{551}{\auDessus{γυ}{ν}}{ἀκούετ᾽ ὦ γυναῖκες οἷ᾽ εἴρηκεν ἡ πανοῦργος}{}
\ligneRav{}{—}{ἡμᾶς ἁπάσας αὖθις αὖ. : καὶ νὴ Δί᾽ οὐδέπω γε}{}
\ligneRav{553}{}{εἴρηχ᾽ ὅσα ξύνοιδ᾽· ἐπεὶ βούλεσθε \difference{πλείον᾽} εἴπω;}{\textbf{553} πλεῖον R : \lecon{πλείον᾽} Bentley}
\ligneRav{554}{\auDessus{γυ}{ν}}{ἀλλ᾽ οὐκ ἂν ἔτ᾽ ἔχοις· ὅσα γὰρ \difference{ᾔδησθ᾽} ἐξέχεας ἅπαντα·}{\textbf{554} ᾔδεις R : \lecon{ᾔδησθ᾽} Hirschig (Ar. \italique{Vesp}.146)}
\ligneRav{555}{μ\auDessus{ν}{Η}}{μὰ Δί᾽ οὐδὲ τήν γε μυριοστὴν μοῖραν ὧν ποιοῦμεν·}{\textbf{555} \lecon{add. γε} Faber}
\ligneRav{556}{}{\difference{ἐπεὶ τάδ᾽} \difference{οὐκ} εἴρηχ᾽, ὁρᾷς, ὡς \difference{στλεγγίδας} λαβοῦσαι}{\textbf{556} ἔπειτά γ᾽ R : \lecon{ἐπεὶ τάδ᾽} Dindorf \espApp οὐχ R : \lecon{οὐκ} \Mon \espApp τλεγγίδας R : \lecon{στλεγγίδας} \Mon}
\ligneRav{557}{—}{ἔπειτα σιφωνίζομεν \difference{τὸν κύσθον}. : ἐπιτριβείης·}{\textbf{557} τὸν σῖτον R, Σ, Suid. σ 512 : \lecon{τὸν κύσθον} Bibi cf. v.1114}
\ligneRav{558}{μ\auDessus{ν}{Η}}{\difference{ὥς τ᾽ αὖ τὰ} κρέ᾽ ἐξ Ἀπατουρίων ταῖς μαστροποῖς διδοῦσαι}{\textbf{558} ὥς τ᾽ αὐτὰ R : ὥς ταῦτα \Mon : \lecon{ὥς τ᾽ αὖ τὰ} Zanetti}
\ligneRav{558}{—}{ἔπειτα τὴν γαλῆν φαμεν– : τάλαιν᾽ ἐγώ· φλυαρεῖς·}{\textbf{558} ὥς τ᾽ αὐτὰ R : ὥς ταῦτα \Mon : \lecon{ὥς τ᾽ αὖ τὰ} Zanetti}
\ligneRav{560}{μ\auDessus{ν}{Η}}{οὐδ᾽ ὡς τὸν ἄνδρα τῷ πελέκει γυνὴ κατεσπόδησεν,}{\textbf{560} \lecon{τὸν ἄνδρα τῷ πελέκει γυνὴ} R : ἑτέρα τὸν ἄνδρα τῷ πελέκει Suid. κ 996}
\ligneRav{}{}{οὐκ εἶπον· οὐδ᾽ ὡς φαρμάκοις ἑτέρα τὸν ἄνδρ᾽ ἔμηνεν,}{}
\ligneRav{562}{—}{οὐδ᾽ ὡς ὑπὸ τῇ πυέλῳ κατώρυξέν ποτε— : ἐξόλοιο·}{\textbf{562} \lecon{ποτε} R : ποτ᾽ ed. pr.}
\ligneRav{563}{μ\auDessus{ν}{Η}}{Ἀχαρνικὴ τὸν πατέρα. : ταυτὶ δῆτ᾽ ἀνέκτ᾽ ἀκούειν;}{\textbf{563} \lecon{ἀχαρνικὴ} R : Ἁχαρνικὴ Dobree}
\ligneRav{}{μ\auDessus{ν}{Η}}{οὐδ᾽ ὡς σὺ τῆς δούλης τεκούσης ἄρρεν εἶτα σαυτῇ}{}
\ligneRav{565}{}{τοῦθ᾽ ὑπεβάλου, τὸ σὸν δὲ θυγάτριον παρῆκας αὐτῇ·}{}
\ligneRav{}{\auDessus{γυ}{ν}}{οὔ τοι μὰ τὼ θεὼ σὺ καταπροίξει λέγουσα ταυτί,}{}
\ligneRav{567}{—}{ἀλλ᾽ ἐκποκιῶ σου τὰς ποκάδας. : οὐ \difference{δὴ} μὰ Δία σύ γ᾽ ἅψει·}{\textbf{567} \lecon{ποκάδας} R, Suid. : πλοκ- Σ \espApp δὲ R : \lecon{δὴ} Bothe}
\ligneRav{}{\auDessus{γυ}{ν}}{καὶ μὴν ἰδού. : καὶ μὴν ἰδού. : λαβὲ θοἰμάτιον Φιλίστη·}{}
\ligneRav{569}{μ\auDessus{ν}{Η}}{\difference{πρόσθιγε} μόνον, κἀγώ σε νὴ τὴν Ἄρτεμιν— : τί δράσεις;}{\textbf{569} πρόσθες R : πρόσιθι Cobet : \lecon{πρόσθιγε} Willems}
\ligneRav{570}{μ\auDessus{ν}{Η}}{τὸν σησαμοῦνθ᾽ ὃν κατέφαγες, τοῦτον χεσεῖν ποιήσω·}{}
\ligneRav{571}{χορ/}{\difference{παύσασθε} λοιδορούμεναι· καὶ γὰρ γυνή τις ἡμῖν}{\textbf{571} (Χορ.) R : \lecon{(Κορ.)} von Velsen \espApp παύσασθαι R : \lecon{παύσασθε} Μυ2}
\ligneRav{}{}{ἐσπουδακυῖα προστρέχει. πρὶν οὖν ὁμοῦ γενέσθαι,}{}
\ligneRav{}{}{σιγᾶθ᾽, ἵν᾽ αὐτῆς κοσμίως πυθώμεθ᾽ ἅττα λέξει·}{}
\ligneRav{}{κλεισθ/}{φίλαι γυναῖκες ξυγγενεῖς τοὐμοῦ τρόπου·}{}
\ligneRav{575}{}{ὅτι μὲν φίλος εἴμ᾽ ὑμῖν, ἐπίδηλος ταῖς γνάθοις·}{}
\ligneRav{576}{}{γυναικομανῶ γὰρ προξενῶ θ᾽ ὑμῶν ἀεί·}{\textbf{576} \lecon{ὑμῶν} R : ὑμᾶς Suid. π 2541}
\ligneRav{}{}{καὶ νῦν ἀκούσας πρᾶγμα περὶ ὑμῶν μέγα}{}
\ligneRav{}{}{ὀλίγῳ τι πρότερον κατ᾽ ἀγορὰν λαλούμενον,}{}
\ligneRav{}{}{ἥκω φράσων τοῦτ᾽ ἀγγελῶν θ᾽ ὑμῖν, ἵνα}{}
\ligneRav{580}{}{\difference{σκοπῆτε καὶ τηρῆτε} \difference{μὴ καὶ} προσπέσῃ}{\textbf{580} σκοπεἴτε καὶ τηρεἴτε R : \lecon{σκοπῆτε καὶ τηρῆτε} Biset \espApp καὶ μὴ R : \lecon{μὴ καὶ} Porson (\italique{Aristophanica} p. 216) \espApp \textbf{581} ἡμῖν R : \lecon{ὑμῖν} \Mon \espApp ἀφράκτοις R : \lecon{ἀφάρκτοις} Dindorf (1830)}
\ligneRav{581}{}{\difference{ὑμῖν} \difference{ἀφάρκτοις} πρᾶγμα δεινὸν καὶ μέγα·}{\textbf{581} ἡμῖν R : \lecon{ὑμῖν} \Mon{} \espApp ἀφράκτοις R : \lecon{ἀφάρκτοις} Dindorf (1830, métathèse attique ancienne $-\rho\kappa-$)}
\ligneRav{}{χορ/}{τί δ᾽ ἔστιν ὦ παῖ; παῖδα γάρ σ᾽ εἰκὸς καλεῖν,}{}
\ligneRav{}{}{ἕως ἂν οὕτως τὰς γνάθους ψιλὰς ἔχῃς·}{}
\ligneRav{584}{κλεισθ/}{Εὐριπίδην \difference{φάσ᾽} ἄνδρα κηδεστήν τινα}{\textbf{584} ἔφας᾽ R : \lecon{φάσ᾽} Zanetti}
\ligneRav{585}{}{\difference{αὑτοῦ} γέροντα δεῦρ᾽ ἀναπέμψαι τήμερον·}{\textbf{585} αὐτου R : \lecon{αὑτοῦ} Biset}
\ligneRav{}{χορ/}{πρὸς ποῖον ἔργον ἢ τίνος γνώμης χάριν;}{}
\ligneRav{}{κλεισθ/}{ἵν᾽ ἅττα βουλεύοισθε καὶ μέλλοιτε δρᾶν,}{}
\ligneRav{}{}{ἐκεῖνος εἴη τῶν λόγων κατάσκοπος·}{}
\ligneRav{}{χορ/}{καὶ πῶς λέληθεν ἐν γυναιξὶν ὢν ἀνήρ;}{}
}
\vspace{3pt}\noindent\rule{4.5cm}{0.4pt}\par\vspace{1.5pt}
{\fontsize{7.0pt}{8.3pt}\selectfont
\noindent\textbf{Loci paralleli (vv. 551–589) ---} \textbf{554} (ᾔδησθ᾽) \EtymGen (A) ᾔδεισθα \espTest \textbf{557} (σιφωνίζομεν τὸν σῖτον) Suidas σ 512 \espTest (σιφωνίζομεν) Pollux VI 19 \espTest \textbf{560} Suidas κ 996 (ἑτέρα pro γυνὴ et κατεσπόδησε vel κατεσπόδωσεν) \espTest \textbf{566} Suidas κ 730 (sed καταπροίξῃ pro καταπροίξει) \espTest \textbf{567} (ἐκποκιῶ σου τὰς ποκάδας) Suidas ε 588 (ἐκπιῶ cod. A) ; Zonaras 672 ἐκποκιῶ \espTest (ἐκποκιῶ) Heschius ε 1637 ἐκποκιῶ ; \EtymMagnum 323.25 \espTest (ποκάδας) Hesychius π 2750 ποκάδες $=$ Photius π 1010 π\{λ\}οκάδες \espTest \textbf{570} (χεσεῖν) \EpimAlphab ε 201 \espTest \textbf{571-2} (πρὶν οὖν -- πυθώμεθ᾽) Suidas ο 292 (κοσμίως omisso) \espTest \textbf{574} (ξυγγενεῖς τοὐμοῦ τρόπου) Suidas ξ 83 \espTest \textbf{576} Suidas π 2541 προξενῶ θ' ὑμᾶς ἀεί\par
}
\vspace{2.5pt}\noindent\rule{4.5cm}{0.4pt}\par\vspace{1.5pt}
{\fontsize{6.9pt}{8.2pt}\selectfont
\noindent\textbf{Suite (en commentaire) de l'apparat critique \& notes paléographiques (\texttt{f. 169r}) ---} \textbf{vv. 556–557 (locus classicus de l'édition Louis Hauseux) :} --- \textbf{v. 557 :} τὸν σῖτον R, $\Sigma$, Souda $\sigma\ 512$ : \lecon{τὸν κύσθον} \lecon{Bibi} (cf. v. 1114). --- \textbf{v. 567 :} \lecon{ποκάδας} R, Souda (πλοκάδας $\Sigma$) \espApp δὲ R : \lecon{δὴ} Bothe. --- \textbf{v. 569 :} πρόσθες R (πρόσιθι Cobet) : \lecon{πρόσθιγε} Alphonse Willems (adopté par L. Hauseux). --- \textbf{v. 571 :} (Χορ.) R : \lecon{(Κορ.)} von Velsen \espApp παύσασθαι R (iotacisme $-\alpha\iota/-\varepsilon$) : \lecon{παύσασθε} \Mon$^{2}$. --- \textbf{v. 581 :} ἡμῖν R : \lecon{ὑμῖν} \Mon{} \espApp ἀφράκτοις R : \lecon{ἀφάρκτοις} Dindorf (1830, métathèse attique ancienne $-\rho\kappa-$).\par
}
\end{minipage}
\end{adjustbox}
\end{adjustwidth}
\newpage
\enTeteD{Codex Ravennas 429 (Biblioteca Classense) --- Fac-similé du f. 169r (vv. 551–589)}
\begin{adjustwidth}{-1.5cm}{-3.0cm}
\begin{center}
\includegraphics[width=16.5cm,height=20.5cm,keepaspectratio]{/workspaces/Thesmophories/dossier_ravennas_thesmophories/folios_individuels_ff162v-177r/14_f169r_classense171r_vv0551-0589.jpg}\par
\vspace{2pt}
{\scriptsize\textbf{Codex Ravennas 429, f. 169r (Classense 171r)} --- ΕΠΕΙΣΟΔΙΟΝ Α', \textbf{vv. 551–589} ($\text{R}^1, \Rpc, \Rdeux$ (\Mon)).}
\end{center}
\end{adjustwidth}
\newpage
\clearpage\ifodd\value{page}\hbox{}\newpage\fi
\enTeteG{ΘΕΣΜΟΦΟΡΙΑΖΟΥΣΑΙ --- Texte grec établi (f. 169r, vv. 551–589)}
\enTeteD{LES THESMOPHORIES --- Traduction française (f. 169r, vv. 551–589)}
\begin{pages}
\begin{Leftside}
\beginnumbering
%% v. 551-552a
\begin{groupe}%
%
\personnageGrec{Μίκα}%
\tetramIambCat{ἀκούετ᾽ ὦ γυναῖκες οἷ᾽ εἴρηκεν ἡ πανοῦργος\\
ἡμᾶς ἁπάσας αὖθις αὖ.\\\passeLigne}%
\end{groupe}%
%% v. 552b-553
\begin{groupe}%
\personnageGrec{Μνησίλοχος}%
\tetramIambCat{\ligne{552}\decalage{ἡμᾶς ἁπάσας αὖθις αὖ.}καὶ νὴ Δί᾽ οὐδέπω γε\\
εἴρηχ᾽ ὅσα ξύνοιδ᾽· ἐπεὶ βούλεσθε πλείον᾽ εἴπω;\apparat{553}{πλεῖον R : \lecon{πλείον᾽} Bentley}\\\passeLigne}%
\end{groupe}%
%% v. 554
\begin{groupe}%
\personnageGrec{Μίκα}%
\tetramIambCat{\small{ἀλλ᾽ οὐκ ἂν ἔτ᾽ ἔχοις· ὅσα γὰρ ᾔδησθ᾽ ἐξέχεας ἅπαντα.}%\testimonia{554}{(ᾔδησθ᾽) \EtymGen (A) ᾔδεισθα}
\apparat{554}{ᾔδεις R : \lecon{ᾔδησθ᾽} Hirschig (Ar. \italique{Vesp}.146)}\\\passeLigne}%
\end{groupe}%
%% v. 555-557a
\begin{groupe}%
\personnageGrec{Μνησίλοχος}%
\tetramIambCat{μὰ Δί᾽ οὐδὲ τήν γε μυριοστὴν μοῖραν ὧν ποιοῦμεν.\apparat{555}{%οὐδὲ τὴν R : %οὐδέπω τὴν Dobree : 
%\lecon{οὐδὲ τήν γε} Faber 
\lecon{add. γε} Faber \ligneApp{556} ἔπειτά γ᾽ R : \lecon{ἐπεὶ τάδ᾽} Dindorf %(1825) 
\espApp οὐχ R : \lecon{οὐκ} \Mon \espApp τλεγγίδας R : \lecon{στλεγγίδας} \Mon \ligneApp{557} τὸν σῖτον R, Σ, Suid. σ 512 : \lecon{τὸν κύσθον} Bibi cf. v.1114}\\% \ligneApp{557} τὸν σῖτον R, Σ, Suid. % σ 512 
%: \lecon{τὸν κύσθον} Bibi cf. v.1114}\\
ἐπεὶ τάδ᾽ οὐκ εἴρηχ᾽, ὁρᾷς, ὡς στλεγγίδας λαβοῦσαι
\commentaire{556-7}{ὁρᾷς ὡς τλεγγίδας λαβοῦσαι / ἔπειτα σιφωνίζομεν τὸν σῖτον.}{Voilà un trait sur les mœurs féminines qui n'a pas été compris des différents exégètes ; c'est donc à nous que revient la tâche – ô combien délicate ! – d'expliquer ce qui répugne aux oreilles. Muses, venez m'aider... mais vous êtes pucelles ; Muses, ne bougez donc. Ah ! que l'on se sent seul ; je ne peux même pas lancer en avant l'honorable professeur Jeffrey Henderson, le plus fin limier des obscénités d'Aristophane et dont le groin fureteur parvient même à les débusquer là où elles ne sont pas (celle-ci, du moins, lui aura échappé.) Mais la recherche est un plaisir solitaire. Attelons-nous y donc ! Pris littéralement, le texte a de quoi surprendre : « tu vois comme nous saisissons les strigiles, et qu'ensuite nous siphonnons le pain. (sic) » Si vous êtes pris de l'envie de corriger, avec Küster, τὸν σῖτον en τὸν οἶνον (l'on siphonnera plutôt du vin que du pain), renseignez-vous sur ce qu'est un strigile (une sorte de spatule ou plaque rectangulaire tout en longueur, légèrement recourbée pour servir de racloir lors de la toilette) ; j'offre l'édition d'Aristophane de Küster à qui parviendra à se saouler en buvant avec ce genre de coupe. (Il faudra d'ailleurs un jour que l'on m'explique l'intérêt d'introduire des émendations si c'est pour laisser un texte tout aussi obscur – et retouché en prime.) Qu'à cela ne tienne, ce n'est pas ce genre de peccadille qui doit arrêter Austin \& Olson : Aristote a dit un jour (\italique{Tropiques},145a24) qu'on pouvait se servir d'un strigile pour puiser de l'eau (hé ! ce sont des choses qui arrivent lorsqu'on est au bain ; A.\&O. auraient pu citer la phrase suivante qui leur était adressée~: Ἀλλ᾽ ὅμως εἴ τις ὁρίσαιτο τὴν στλεγγίδα ὄργανον πρὸς τὸ ἀρύειν, ἡμάρτηκεν, « Mais cependant, celui qui définirait le strigile comme instrument pour puiser se trompe »). Non, l'hypothèse d'ivrognesses s'enivrant à grandes rasades de strigiles est à écarter, sauf à m'expliquer : 1° comment boire au strigile \ \ \ 2° pourquoi aller chercher dans la salle de bain un instrument qui n'est pas censé en bouger \ \ \ 3° et quand bien même, qu'y aurait-il de déshonorant dans cette pratique ? \ \ \ Il faut chercher ailleurs, car les tentations des bains ne sont pas d'ordre biberonnesque. Aristophane vitupère ailleurs\marqueur{1} contre cette propension des Athéniens décadents à se délicater : les bains chauds sont un exemple parmi d'autres de ces mauvaises habitudes qui amollissent les cœurs et durcissent les vices. Pour comprendre exactement où il veut en venir, c'est Euripide lui-même qui nous donnera la clef de la salle de bain : dans son drame satyrique du \italique{Cyclope}, le héros aux-mille-tours-lubriques, se désole sur « son cher siphon » (v. 339, τὸν σίφωνα τὸν φίλον) qui depuis trop longtemps vit dans la solitude – solitude vis-à-vis du beau sexe ; le verbe χηρεύω servant aussi à désigner l'état de veuvage. Il parle évidemment de son πέος~; une fois attesté ce sens pour σίφων, celui du verbe σιφωνίζω s'éclaire d'une lumière nouvelle (le verbe est d'autant plus amusant qu'il joue sur le rapprochement, cher à Aristophane, entre βινεῖν et πίνειν).\\
La présence étonnante de τὸν σῖτον s'explique alors aisément : la leçon originale était τὸν κύσθον $=$ \latin{pudenda muliebria}, très proche phonétiquement de τὸν σῖτον. 
Iotacisme aidant, il s'agit d'une simple \italique{métathèse} du `s' (avec chute du `k').
Il est remarquable de noter que la même métathèse se produit au vers 1114 : le rustre archer Scythe, dans l'optique de prouver que Mnésiloque n'est pas une femme, conseille d'aller y voir de plus près :%
\citComm{σκέψαι †τὸ σκύτο†.}%
Scaliger \italique{à la grâce du mètre} -- et surtout à la grâce du sens -- avait rétabli :%
\citComm{σκέψαι τὸ κύστο ($=$ τὸν κύσθον dans le parler scythe).}%
Il est à croire qu'un scribe des \italique{Thesmophories} confondait avec un autre le mot grec pour désigner le « plus bel apanage » du corps féminin, qui tentait tendrement tant Brassens \& Villon : il s'y sera trompé par deux fois dans notre pièce (\latin{bis repetita...}).
J'arrête mes explications ici, de peur qu'un Calliclès de passage ne m'assène :%
\citComm{Καλλίκλης-Φιλόλογος --- οὐκ αἰσχύνῃ εἰς τοιαῦτα ἄγων, ὦ Πέπωκες, τοὺς λόγους;\marqueur{2}}%
%La présence étonnante de τὸν σῖτον s'explique alors aisément : la leçon originale était τὴν μᾶζαν, qu'on traduit faute de mieux par \italique{la galette}, et glosée τὸν σῖτον, \italique{le pain} par un érudit n'ayant pas tout saisi... car la μᾶζα peut désigner aussi le « plus bel apanage » du corps féminin, qui tentait tendrement tant Brassens \& Villon (cf. l'ingénieuse – et crue ! – explication du vers 835 des \italique{Acharniens} par Jeffrey Henderson\marqueur{2} – oreilles chastes s'abstenir ! – et, si la démonstration ne suffit pas, on pourra finir de se convaincre avec l'ordurier Hipponax\marqueur{3}). 
%Un scribe aura alors pris le Pirée pour un homme, l'attribut de la femme pour quelque denrée peu appétissante\marqueur{4} et une glose intempestive pour la leçon originale\marqueur{5}. 
%(Moi qui m'étais lancé dans ce commentaire avec la ferme intention de disperser la horde libidineuse des commentateurs qui collent aux basques d'Aristophane et maculent\marqueur{3} notre poète d'une horrible réputation de pornographe, il me déplaît de leur donner du grain à moudre. Las !.. Philologus semble moins à l'affût -- ou moins au fait -- des travers féminins. Aussi l'avons-nous suppléé en ce passage -- une fois n'est pas coutume !)\\% -- pour un autre passage d'Aristophane
%Plutus 1036 Γραῦς -- διὰ δακτυλίου μὲν οὖν ἔμεγ᾽ ἂν διελκύσαις.
%Χρεμύλος -- εἰ τυγχάνοι γ᾽ ὁ δακτύλιος ὢν τηλία.
%Διὰ δακτυλίου δεῖ σε ἑλκυσθῆναι Shakespeare (Souda online ; français Willems ; italien i proverbi greci ; dia koskinou la véritable expression ?)
Profitons-en pour remarquer à quel point les griefs des femmes envers Euripide ne sont pas sérieux dans la bouche d'Aristophane (si Euripide est misogyne, c'est bien là le moindre de ses défauts) ; car si les maris athéniens, s'en revenant d'une pièce d'Euripide, regardaient d'un œil soupçonneux leur femme tenant une marmite, avec quels yeux vont-ils dorénavant observer leur femme se saisir d'un strigile ou de tout instrument de forme analogue ?! Quoi qu'il en soit, comme pour l'histoire de l'amant qui s'échappe (v. 501), Mnésiloque n'invente rien mais reprend des thèmes certainement populaires du Landerneau athénien : le détournement à fin lubrique du strigile se retrouve illustrée, ce me semble, sur une coupe d'Onésimos\marqueur{4}, antérieure à la pièce d'Aristophane. Pour finir – était-il besoin de préciser ?.. – précisons tout-de-même : Mnésiloque mime en lançant son « ὁρᾷς » à Mika : « Ainsi font, font, font, les petites midinettes. Un siphon, fond, fond, on se siphonne toutes à fond.»}%
\debutNote{556-7}%
\note{1}{C'est un point de la querelle du Juste et de l'Injuste des \italique{Nuées} (v. 1044-1046) :
\citComm{Ἄδικος Λόγος — σε θερμῷ φησι λοῦσθαι πρῶτον οὐκ ἐάσειν.\\
καίτοι τίνα γνώμην ἔχων ψέγεις τὰ θερμὰ λουτρά;\\
Δίκαιος Λόγος — ὁτιὴ κάκιστόν ἐστι καὶ δειλὸν ποιεῖ τὸν ἄνδρα.\\
\francais{Raisonnement Injuste — il ne te permettra pas, dit-il, de te baigner dans l'eau chaude.\\
(Au Juste :) Mais d'après quel principe improuves-tu les bains chauds? \\
Raisonnement Juste — Parce qu'ils sont très pernicieux et qu'ils rendent l'homme lâche.}}%
Platon, d'accord sur ce point avec Aristophane comme sur beaucoup d'autres, nomme les bains chauds γεροντικὰ λουτρά, « des bains de vieillards » (\italique{Lois}, VI, 761c). Érasme en a fait un adage (2061): ἁπαλοὶ θερμολουσίαις, \italique{amollis par les bains chauds}.}%
%\note{2}{p. 113 de sa \italique{Maculate Muse}.}%
%\note{3}{fragment 28 de l'édition Kassel-Austin ; se reporter aux notes du \nieme{6} fragment pour l'éclaircissement.\eff{Et chez Masson ?}}%
%\note{4}{La μᾶζα étant un mets tout-à-fait commun et fait à la va-vite quand le σῖτος s'achetait chez un boulanger. Pour en savoir davantage sur la μᾶζα, l'on se reportera avec intérêt à la note essentielle de Willems (t. I, p. 96-99), reproduite à la fin.}%
%\note{5}{Confusion facilitée par la scansion identique.}%
%\note{2}{Iotacisme aidant, il s'agit d'une simple \italique{métathèse}.}%
\note{2}{Platon, \italique{Gorgias} 494e.}%
%\note{3}{Le chef de cette meute, le sus-nommé Jeffrey Henderson, ne s'en cache même pas : il affiche clairement ses `ambitions' jusque dans le titre de son exégèse porcine d'Aristophane : \italique{Maculate Muse}. Le livre fait autorité et cependant une bonne moitié -- sans exagération aucune -- n'est que pure élucubrication ! J'en viens à me demander si Hipponax ou Archiloque, qui traînent derrière eux une réputation plus ordurière encore que celle d'Aristophane, ne seraient pas eux aussi victimes du même genre de calomniateurs -- l'affreux Philologus ! ὦ βορβοροτάραξι Φιλόλογε !}%
\note{3}{Bâle, Antikenmuseum BS 440. ARV 326, 86 \italique{bis}. \italique{Para} 359. Elle est reproduite par John Boardman, fig. 233 in \italique{Les vases athéniens à figures rouges}, \italique{La période archaïque} ; première édition sous le titre \italique{Athenian Red Figure Vases}, \italique{The Archaic Period}, 1975. (Le lecteur pourra la retrouver aux annexes de ce livre.) En parcourant le livre de Boardmann, cette coupe a retenu mon attention : elle me paraît parfaitement illustrer l'anecdote de Mnésiloque. (Reconnaissons tout d'abord que l'interprétation est sujette à caution : j'ai montré cette coupe à différents amis et ai eu droit à autant d'interprétations qu'il y avait d'amis... interprétations qui en disaient d'ailleurs plus long sur leur propre psychologie que sur les mœurs des Athéniens ; profitons-en donc pour suggérer à l'éventuel psychanalyste qui se serait égaré à lire ces lignes d'abandonner fissa ses infantiles tests de Rorschach et de passer aux coupes athéniennes du \nieme{VI} siècle : dix vases athéniens d'Onésime seraient aussi révélateurs des obsessions ou perversions des supposés psychotiques que les taches d'encre du regretté Hermann, arraché trop tôt à l'affection des siens... Sans parler du gain en valeur artistique : incommensurable ! Refermons cette parenthèses et retournons à notre coupe :) qu'y ai-je donc vu ? À droite, un satyre (sa calvitie naissante est caractéristique des représentations de satyres de cette époque ; c'est entre autre sur ce détail de la laideur physique de Socrate, visible sur ses bustes, qu'Alcibiade rapproche selon moi Socrate du satyre Silène), un satyre donc, en train de se strigiler le fondement, pointe (du doigt) le Mont de Vénus d'une femme nue, lui désignant par là l'usage impropre du strigile auquel il lui arrive de céder. La clef qui me manque pour être tout-à-fait sûr de la bonne compréhension de cette scène est l'identification de l'objet dans la main de la femme : je crois y reconnaître un strigile enroulé dans du tissu. (plus recommandé pour atteindre le plaisir que la lame de bronze brute.) D'aucun y a vu le pagne du jeune homme à gauche. J'interprète cette scène comme un résumé des vices féminins de l'imaginaire grec : la femme est incapable de se contenir, elle cède à tous les plaisirs : le plaisir de la vigne (la femme est adossée à une outre) ; le plaisir de la chair (l'adultère est suggéré avec l'amant à gauche qui se rhabille) et -- toujours insatiable ! -- le plaisir des dieux. (nue, la voilà habillée pour l'hiver !..)
Une autre hypothèse d'un mien ami quant à cette scène vaut d'être rapportée : il l'avait interprétée quelque peu différemment : tout en restant dans le domaine d'Onan, il proposait de garder τὸν σῖτον et de comprendre : \italique{le pain} [d'argile] servant à récupérer les peaux mortes \& autres hormones raclées avec l'huile à l'aide du strigile (les στλεγγίσματα et autre γλοιός), raclures qui étaient ensuite parfois monnayées (notamment lorsqu'elle provenaient de virils athlètes)... et pourquoi pas par quelque vicelarde pour s'en foutre plein les zones érogènes ? Son explication avait l'avantage indéniable de conserver la leçon du manuscrit... Mais je crains que son imagination ne se soit quelque peu emballée : je n'ai pas trouvé d'exemple d'acception du mot σῖτος pour désigner ce «~pain d'argile~». Et si Pline nous confirme un usage médical des στλεγγίσματα, l'usage détourné, qui lui a été suggéré par certains papiers de « recherche », est un pur fantasme des modernes scholiastes.}\\
ἔπειτα σιφωνίζομεν τὸν κύσθον.%\testimonia{557}{(σιφωνίζομεν τὸν σῖτον) Suidas σ 512 \espTest (σιφωνίζομεν) Pollux VI 19}
%\apparat{557}{τὸν σῖτον R, Σ, Suid. σ 512 : \lecon{τὸν κύσθον} Bibi cf. v.1114}
\\\passeLigne}%
\end{groupe}%
%% v. 557b
\begin{groupe}%
\personnageGrec{Μίκα}%
\tetramIambCat{\ligne{557}\decalage{ἔπειτα σιφωνίζομεν τὸν κύσθον.}ἐπιτριβείης.\\\passeLigne}%
\end{groupe}%
%% v. 558-559a
\begin{groupe}%
%
\personnageGrec{Μνησίλοχος}%
\tetramIambCat{ὥς τ᾽ αὖ τὰ κρέ᾽ ἐξ Ἀπατουρίων ταῖς μαστροποῖς διδοῦσαι\\
ἔπειτα τὴν γαλῆν φαμεν–\apparat{558}{ὥς τ᾽ αὐτὰ R : ὥς ταῦτα \Mon : \lecon{ὥς τ᾽ αὖ τὰ} Zanetti}
\commentaire{558-9}{ὥς τ᾽ αὐτὰ κρέ᾽ ἐξ Ἀπατουρίων ταῖς μαστροποῖς διδοῦσαι / ἔπειτα τὴν γαλῆν φαμεν.}{Si Mnésiloque avait commencé un discours tout en nuance, il a bien vite perdu de sa prudence et se met à surpasser le maître Euripide dans la diffamation du genre féminin. Le ton monte et Mnésiloque en est maintenant à dénoncer les pratiques les plus honteuses des Athéniennes les plus dissolues.\\
Nous ne nous attarderons pas sur ce que pouvaient bien procurer les procureuses en échange de ces détournements de viandes\marqueur{1}.
Plus intéressante est l'excuse alléguée pour expliquer cette soudaine disparition carnée. La \italique{team 1\expo{er}} degré des commentateurs d'Aristophane n'a pas manqué de faire remarquer que la belette était un animal chapardeur\marqueur{2}. Certes ! mais ce pauvre animal a surtout \italique{bon dos} : on l'accuse bien facilement ; ainsi aux vers 1150-1152 de la \italique{Paix} :%
\citComm{ἦν δὲ καὶ πυός τις ἔνδον καὶ λαγῷα τέτταρα,\\
εἴ τι μὴ 'ξήνεγκεν αὐτῶν ἡ γαλῆ τῆς ἑσπέρας·\\
ἐψόφει γοῦν ἔνδον οὐκ οἶδ᾽ ἅττα κἀκυδοιδόπα\\
Il y avait aussi au logis de l'amouille et quatre morceaux de lièvre ; 
à moins que la belette n'en ait volé hier soir ; 
car il se passait là-dedans je ne sais quel bruit et quel remuement.}%
L'allusion au remue-ménage n'est pas limpide mais il est clair que le sens ne peut être pris littéralement : qui, entendant une belette le dévaliser, ne s'armerait aussitôt d'un bâton pour châtier la voleuse ? (Noter d'ailleurs ici aussi la présence de l'article : ἡ γαλῆ.) 
Aristophane joue certainement sur un détournement d'une locution proverbiale : la belette attique étant notre \italique{bouc-émissaire} à qui l'on fait \italique{porter le chapeau}. 
Conformité du grec et du français : la belette joue le même rôle que notre chat, animal qui n'a pas encore été introduit dans l'Athènes du \nieme{V} siècle : il chasse les rongeurs. Or, que lit-on chez La Fontaine ?%
\citComm{Le sommeil quitta son logis, /
Il eut pour hôte les soucis, /
Les soupçons, les alarmes vaines. / 
Tout le jour il avait l'œil au guet ; \textbf{et la nuit, /
Si quelque chat faisait du bruit, /
Le chat prenait l'argent.}\marqueur{3}}%
Ce n'est pas moi qui m'abuse : remplacez l'argent par la nourriture et les deux passages disent exactement la même chose.
Voici ce que l'on trouve dans le \italique{Dictionnaire de la langue verte} d'Alfred Delvau :%
\citComm{« C'est le chat ! » : Expression de l'argot du peuple, qui souligne ironiquement un doute, une dénégation. Ainsi, quelqu'un disant : Ce n'est pas moi qui ai fait cela. – Non ! c'est le chat, lui répondra-t-on.}%
Un autre passage de la \italique{Paix} montre encore mieux ce sens imagé : aux vers 793-795 :%
\citComm{καὶ γὰρ ἔφασχ᾽ ὁ πατὴρ ὃ παρ᾽ ἐλπίδας /
εἶχε τὸ δρᾶμα γαλῆν τῆς /
ἑσπέρας ἀπάγξαι.\\
Aussi bien leur père prétend-il qu'un drame qu'il pensait tenir contre toute attente fut étranglé le soir par une belette.}%
Juste avant ce passage, Aristophane faisait un honnête résumé des bienfaits qu'il a apportés à la Comédie, il en vient ensuite à étriller les nullités de ses poétereaux de contemporains. Nous en sommes à Carcinos \& Fils qui font des pieds et des mains pour pouvoir danser – c'est-à-dire obtenir la subvention d'un chœur.
Mais une fois le chœur accordé – à rebours de ce qui était attendu ! (παρ᾽ ἐλπίδας) [du bon goût ?] – Carcinos-père fait volte-face et prétend que son drame « fut étranglé le soir par une belette » (c'est encore le soir que la belette commet son forfait). 
Le scholiaste nous est ici précieux puisqu'il précise que Carcinos voulait représenter une pièce intitulée les \italique{Souris} (Μῦες\marqueur{4}) ; encore que la chronique du temps ne nous ait pas conservé l'affaire, il paraît aisé de reconstituer peu ou prou ce qui s'est passé. Les poètes devaient présenter à l'archonte éponyme leur projet. 
Pour les Dionysies, l'archonte sélectionnait trois tétralogies qui recevaient chacune le soutien financier d'un riche Athénien : le chorège. 
Ce n'est qu'alors que le poète se mettait à écrire. Et c'est à cette étape que les forces auront probablement manqué à notre Carcinos… ce dernier croyant s'en tirer avec un trait d'esprit : s'il ne peut pas présenter ses \italique{Souris}, c'est que la belette les a étranglées. 
Le mot ne peut être \italique{bon} que s'il existait effectivement une expression pour faire endosser la responsabilité d'un méfait à cette belette proverbiale.}%
\debutNote{558-9}%
\note{1}{Contentons-nous de renvoyer à la magistrale note d'Alphonse Willems sur la fête des Apaturies, t. II, pp. 118-120 ; le lecteur y découvrira l'importance démocratique de la κρεανομία, l'institution de la distribution de viande. Sur les débordements que de telles orgies pouvaient entraîner, l'on pourra – un exemple parmi d'autres – essayer de comprendre le jeu de mots d'Aristophane sur l'ἀνάρρυσις (au vers 890 de la \italique{Paix}), le second jour des Apaturies.}%
%\note{1}{Voir par exemple les vers 363-364 des \italique{Guêpes}. Lucien cite aussi l'adage ἁρπακτικώτεροι τῶν γαλῶν (\italique{Le pêcheur}, 34).}%
\note{2}{Voir par exemple les vers 363-364 des \italique{Guêpes}. Lucien, \italique{Piscator}, 34 : ἁρπακτικώτεροι τῶν γαλῶν.}%
\note{3}{Il s'agit – vous l'aurez bien sûr reconnue ; sinon, cessez immédiatement votre lecture et ne revenez que lorsque vous la connaîtrez par cœur ! – de la seconde fable du huitième livre : \italique{Le Savetier et le Financier}.}%
\note{4}{Probablement le drame satyrique. Μῦς étant un nom propre, comme chez Hérodote (VIII, 133), dans la \italique{Souda} (ο 680) ou chez Athénée (XI, 782b), les Μῦες désignant les enfants de ce Μῦς. Je m'inscris en faux contre l'interprétation de Paul Mazon (\italique{La} Paix\italique{, texte grec publié avec une introduction, des notes critiques et explicatives} (1904), v. 793) reprise par Taillardat (§792, p.463 n.5) qui pensent que Μῦες était le surnom de ces débauchés de Carcinites. L'on s'expliquerait difficilement la confusion du scholiaste (présent sur tous les bons manuscrits et repris par Suidas κ 396). Objection plus forte : cela n'expliquerait pas le texte de la \italique{Paix} : la belette (γαλῆ) est accusé d'avoir étranglé nuitamment non les choreutes mais la pièce elle-même (τὸ δρᾶμα).}\\\passeLigne}%
\end{groupe}%
%% v. 559b
\begin{groupe}%
\personnageGrec{Μίκα}%
\tetramIambCat{\ligne{559}\decalage{ἔπειτα τὴν γαλῆν φαμεν–}τάλαιν᾽ ἐγώ· φλυαρεῖς.\\\passeLigne}%
\end{groupe}%
%% v. 560-562a
\begin{groupe}%
\personnageGrec{Μνησίλοχος}%
\tetramIambCat{οὐδ᾽ ὡς τὸν ἄνδρα τῷ πελέκει γυνὴ κατεσπόδησεν,%\testimonia{560}{Suidas κ 996 (ἑτέρα pro γυνὴ et κατεσπόδησε vel κατεσπόδωσεν)}
\apparat{560}{\lecon{τὸν ἄνδρα τῷ πελέκει γυνὴ} R : ἑτέρα τὸν ἄνδρα τῷ πελέκει Suid. κ 996}\\
\small{οὐκ εἶπον· οὐδ᾽ ὡς φαρμάκοις ἑτέρα τὸν ἄνδρ᾽ ἔμηνεν,}\\
οὐδ᾽ ὡς ὑπὸ τῇ πυέλῳ κατώρυξέν ποτε—%\apparat{562}{\lecon{ποτε} R : ποτ᾽ ed. pr.}
\\\passeLigne}%
\end{groupe}%
%% v. 562b
\begin{groupe}%
\personnageGrec{Μίκα}%
\tetramIambCat{\ligne{562}\decalage{οὐδ᾽ ὡς ὑπὸ τῇ πυέλῳ κατώρυξέν ποτε—}ἐξόλοιο.\\\passeLigne}%
\end{groupe}%
%% v. 563a
\begin{groupe}%
\personnageGrec{Μνησίλοχος}%
\tetramIambCat{Ἀχαρνικὴ τὸν πατέρα.\apparat{563}{\lecon{ἀχαρνικὴ} R : Ἁχαρνικὴ Dobree}%
%\commentaire{563}{Ἀχαρνικὴ}{\eff{Une charbonnière ? qui aura charbonné son mari père ? Ou alors : après la hache et le poison, reste la corde (kat)ἀρτῶ à la la place de katwrucen Τὰ τρία τῶν εἰς θάνατον .}Ἀχαρνικῇ πύλῃ (épigraphie)
%κατώρης· κάτω ῥέπων <- (kat -ὄρνυμι) l'aura-t-elle précipité au pied de la porte d'acharne ? (Pris littéralement, le texte indique seulement la façon dont la femme a tenté de se débarrasser du corps et non comment elle l'a tuée -- ce qui constitue le réquisitoire de Mnésiloque ; de plus, si l'affaire était connue de tous, pourquoi préciser ici que c'est une acharnienne qui a fait e coup ? Cédons donc donc à la philologienne manie de vouloir à tout prix retoucher le texte -- au moins, contentons- nous de circonscrire cette retouche au commentaire et non au texte présenté au lecteur) "Je nai pas dit : ni comment, par un poison, une autre femme affolit son mari, ni comment au pied de la porte}
\\\passeLigne}%
\end{groupe}%
%% v. 563b
\begin{groupe}%
\personnageGrec{Μίκα}%
\tetramIambCat{\ligne{563}\decalage{Ἀχαρνικὴ τὸν πατέρα.}ταυτὶ δῆτ᾽ ἀνέκτ᾽ ἀκούειν;\\\passeLigne}%
\end{groupe}%
%% v. 564-565
\begin{groupe}%
\personnageGrec{Μνησίλοχος}%
\tetramIambCat{οὐδ᾽ ὡς σὺ τῆς δούλης τεκούσης ἄρρεν εἶτα σαυτῇ\\
τοῦθ᾽ ὑπεβάλου, τὸ σὸν δὲ θυγάτριον παρῆκας αὐτῇ.\\\passeLigne}%
\end{groupe}%
%% v. 566-567a
\begin{groupe}%
\personnageGrec{Μίκα}%
\tetramIambCat{οὔ τοι μὰ τὼ θεὼ σὺ καταπροίξει λέγουσα ταυτί,%\testimonia{566}{Suidas κ 730 (sed καταπροίξῃ pro καταπροίξει)}
\\
ἀλλ᾽ ἐκποκιῶ σου τὰς ποκάδας.%\testimonia{567}{(ἐκποκιῶ σου τὰς ποκάδας) Suidas ε 588 (ἐκπιῶ cod. A) ; Zonaras 672 ἐκποκιῶ \espTest (ἐκποκιῶ) Heschius ε 1637 ἐκποκιῶ ; \EtymMagnum 323.25 \espTest (ποκάδας) Hesychius π 2750 ποκάδες $=$ Photius π 1010 π\{λ\}οκάδες}
\apparat{567}{\lecon{ποκάδας} R, Suid. : πλοκ- Σ \espApp δὲ R : \lecon{δὴ} Bothe}\\\passeLigne}%
\end{groupe}%
%% v. 567b
\begin{groupe}%
\personnageGrec{Μνησίλοχος}%
\tetramIambCat{\ligne{567}\decalage{ἀλλ᾽ ἐκποκιῶ σου τὰς ποκάδας.}οὐ δὴ μὰ Δία σύ γ᾽ ἅψει.\\\passeLigne}%
\end{groupe}%
%% v. 568a
\begin{groupe}%
\personnageGrec{Μίκα}%
\tetramIambCat{καὶ μὴν ἰδού.\\\passeLigne}%
\end{groupe}%
%% v. 568b
\begin{groupe}%
\personnageGrec{Μνησίλοχος}%
\tetramIambCat{\ligne{568}\decalage{καὶ μὴν ἰδού.}καὶ μὴν ἰδού.\\\passeLigne}%
\end{groupe}%
%% v. 568c
\begin{groupe}%
\personnageGrec{Μίκα}%
\tetramIambCat{\ligne{568}\decalage{καὶ μὴν ἰδού.καὶ μὴν ἰδού.}λαβὲ θοἰμάτιον Φιλίστη.
\commentaire{568}{λαβὲ θοἰμάτιον Φιλίστη.}{Mettre bas les manteaux était de mise avant d'en venir aux mains. Ainsi au vers 408 des \italique{Guêpes} ; autre exemple : le fragment d'Hipponax\marqueur{1} contre Boupalos\marqueur{2} : Boupalos, sculpteur de Chios, avait fait d'Hipponax un portrait où il exagérait encore la laideur du modèle. Pour se venger, le poète écrivit contre lui des vers mordants et pleins de menaces. Suidas (β 452) nous a conservé celui-ci, qui paraît avoir passé en proverbe : %
\citComm{λάβετε μευ ταἰμάτια [θοιμάτιον codd.], κόψω Βουπάλου τον οφθαλμόν.\\
\francais{Prenez mon manteau, que je fasse sauter un œil à Boupalos.}}%
}%
\debutNote{568}%
\note{1}{Fr. 121 de l'édition Degani : \italique{Hipponactis testimonia et fragmenta}, Teubner, 1991 ; fr. 120 chez Olivier Masson, voir aussi son introduction p. 13.}%
\note{2}{Il est question de ce Boupalos au vers 361 de \italique{Lysistrata} -- le fragment sus-cité d'Hipponax nous a d'ailleurs été transmis par le scholiaste à ce passage ; je reprends la note de Willems de ce passage, p. 351 du second tome.}\\\passeLigne}%
\end{groupe}%
%% v. 569a
\begin{groupe}%
\personnageGrec{Μνησίλοχος}%
\tetramIambCat{πρόσθιγε μόνον, κἀγώ σε νὴ τὴν Ἄρτεμιν—\apparat{569}{πρόσθες R : πρόσιθι Cobet : \lecon{πρόσθιγε} Willems}%
\commentaire{569}{πρόσθες μόνον.}{Willems substitue πρόσθιγε (du verbe προσθιγγάνω, \italique{toucher}) à πρόσθες qu'il soupçonne de ne pas être le mot propre. Il s'en explique succinctement dans la \nieme{16} note reproduite à la fin. D'aucuns ont essayé de justifier la leçon du texte en avançant une tournure elliptique : πρόσθες [τὴν χεῖρα] (expression qui se trouve effectivement au vers 1669 des \italique{Phéniciennes} d'Euripide\marqueur{1}) mais tous les exemples proposés comportent χεῖρα, qui semble indispensable au sens. Cobet avait donc proposé de lire πρόσιθι μόνον, \italique{approche seulement} ; ce que l'on dirait effectivement dans nos langues modernes mais le grec use d'une image différente : il dit \italique{jette(-moi) seulement ta main} (cf. les exemples déjà cités) ou bien \italique{touche(-moi) seulement} (cf. les exemples proposés par Willems auxquels on peut ajouter \italique{Lysistrata} v. 365, qui est l'exacte correspondance de son émendation : ἅψαι μόνον).}%
\debutNote{569}%
\note{1}{L'on cite d'autres exemples équivalents : le vers 933 des \italique{Nuées}: κλαύσει, τὴν χεῖρ᾽ ἢν ἐπιβάλλῃς ; le fragment 309(277) (éd. K.-A. = \italique{Zonarae lexikon} p.1168) : τὴν χεῖρα μὴ ᾿πίβαλλε, μὴ κλάων καθῇ ; \italique{Lysistrata}, vers 435-436 \& 439-440 : τὴν χεῖρά μοι ... προσοίσει ... κλαύσεται \& εἴ ... μόνον τὴν χεῖρ᾽ ἐπιβαλεῖς, ἐπιχεσεῖ πατούμενος.}%
\\\passeLigne}%
\end{groupe}%
%% v. 569b
\begin{groupe}%
\personnageGrec{Μίκα}%
\tetramIambCat{\ligne{569}\decalage{πρόσθιγε μόνον, κἀγώ σε νὴ τὴν Ἄρτεμιν—}τί δράσεις;\\\passeLigne}%
\end{groupe}%
%% v. 570
\begin{groupe}%
\personnageGrec{Μνησίλοχος}%
\tetramIambCat{τὸν σησαμοῦνθ᾽ ὃν κατέφαγες, τοῦτον χεσεῖν ποιήσω.%\testimonia{570}{(χεσεῖν) \EpimAlphab ε 201}%
\commentaire{570}{τὸν σησαμοῦνθ᾽ ὃν κατέφαγες, τοῦτον χεσεῖν ποιήσω.}{«~Le gâteau de sésame que tu as engoulé, je vais te le faire chier.~» Mnésiloque aurait pu ajouter : que tu viens d'engouler \italique{en cachette} ; car ce \grec{σησαμοῦς}\marqueur{1} a été avalé en plein milieu des Thesmophories, le jour du \italique{jeûne}, \grec{νηστεία}.
Athénée (XIV, 647a) nous rapporte que les Syracusaines préparaient des gâteaux semblables de miel et de sésame, les \italique{mylles} pour leurs Thesmophories~:%
\citComm{ΜΥΛΛΟΙ. Ἡρακλείδης ὁ Συρακόσιος ἐν τῷ περὶ Θεσμῶν ἐν Συρακούσαις φησὶ τοῖς Παντελείοις τῶν Θεσμοφορίων ἐκ σησάμου καὶ μέλιτος κατασκευάζεσθαι ἐφήβαια γυναικεῖα, ἃ καλεῖσθαι κατὰ πᾶσαν Σικελίαν μυλλοὺς καὶ περιφέρεσθαι ταῖς θεαῖς.\\
\francais{Mylloi. Héraclide de Syracuse rapporte ce qui suit, dans son \italique{Traité des Rites religieux} : « Le jour des grandes Thesmophories, on fait à Syracuse avec de la farine, du miel et du sésame, la figure de la partie sexuelle d'une femme, que l'on porte en pompe, pour l'offrir aux déesses. On l'appelle \italique{mylle} par toute la Sicile. »}}%
C'est de peur que Mica risque de relâcher son gâteau. Chez Aristophane la peur se manifeste presque toujours physiquement par des troubles du bas-ventre\marqueur{2}. On dira en français \italique{avoir les pétoches}. % -- littéralement.
\latin{Pulchre, bene, recte.} Mais l'essentiel nous échappe encore : Mnésiloque ne promet pas à Mica un petit coup en cas d'agression de sa part, mais une riposte non graduée : une pluie de horions à lui assainir le côlon ! Car il faut comprendre : «~tu vas en chier [de peur] \italique{malgré} le gâteau de sésame par toi engueulé~».
Ce genre de gâteaux avait la vertu d'être astringent\marqueur{3}. C'est ce qu'il ressort d'un passage des \italique{Nuées} (vers 506-508) : Strepsiade hésite à rentrer dans l'antre de Socrate, qu'il compare aimablement à celui de Trophonius : %
\citComm{\ \ \ \ \ \ \ \ \ \ \ \ \ \ \ \ \ \ \ \ \ \ \ \ \ \ \ \ \ \ εἰς τὼ χεῖρέ νυν\\
δός μοι μελιτοῦτταν πρότερον, ὡς δέδοικ᾽ ἐγὼ\\
εἴσω καταβαίνων ὥσπερ εἰς Τροφωνίου.\\
Mets-moi d'abord en main un gâteau au miel ; car j'ai peur de descendre, comme si je pénétrais dans l'antre de Trophonios.}%
Il faut bien-sûr prononcer d'un trait `εἰς Τροφωνίου' pour faire entendre στρόφος, la \italique{colique}\marqueur{4}.\\
%()\\
C'est aussi ce qui explique le bon mot de table qu'Athénée (I, 4d) rapporte :%
\citComm{ἀμίδα δότω τις ἢ πλακοῦντά τις δότω\\
Qu'on m'apporte un pot de chambre ou bien un gâteau !}%
Au contraire de tous ces gâteaux de miel, la συρμαία (sorte de rave noire que prisaient les Égyptiens pour leur régime ; voir le commentaire \italique{infra} aux vers 855-857) possède des vertus purgatives. 
Cette opposition est vraisemblablement la raison de la confrontation dans la pièce perdue des \italique{Danaïdes}\marqueur{5} entre deux vendeurs : celui de syrméa et celui d'ἐγκρίδες\marqueur{6} :%
\citComm{ἐν δ᾿ Ἀριστοφάνης Δαναΐσι συρμαιοπῶλαι καὶ ἐγκριδοπῶλαι.}%
%Profitons-en pour nous pencher sur un problème d'établissement de texte : un passage des \italique{Deipnosophistes} d'Athénée (XIV, 645e), fragment n° 269(256) d'Aristophane dans l'édition de Kassel \& Austin des \italique{Poetae Comici Graeci}. Athénée y définit ce que sont que les ἐγκρίδες :%
%\citComm{ΕΓΚΡΙΔΕΣ. πεμμάτιον ἑψόμενον ἐν ἐλαίῳ καὶ μετὰ τοῦτο μελιτούμενον.\\
%μνημονεύει αὐτῶν Στησίχορος διὰ τούτων· χόνδρον τε καὶ ἐγκρίδας ἄλλα τε πέμματα καὶ μέλι χλωρόν.\\
%μνημονεύει αὐτῶν καὶ Ἐπίχαρμος καὶ ἐν τοῖς Ἐγχειρογάστορσι Νικοφῶν.\\
%Ἀριστοφάνης δʼ ἐν Δαναίσιν καὶ ποιητήν [sic] φησιν αὐτῶν εἶναι ἐν τούτοις· μήτʼ αρμα [re-sic] εἶναι ἐγκριδοπώλην.\\
%ENCRIDES. C'était un petit gâteau qu'on faisait frire dans l'huile, après quoi on le frottait de miel.\\
%Stésichore en parle ainsi : « Du gruau, des \italique{encrides}, et autres pâtisseries, et du miel tout récent. »\\
%Il en est aussi fait mention dans Épicharme et dans les \italique{Mercenaires} de Nicophon.\\
%Aristophane dit dans les \italique{Danaïdes} $[$ ... $]$}%
%Faut-il voir dans le titre de la pièce de Nicophon, les \italique{Mercenaires}, ἐγχειρογάστορες, un jeu de mots qui ferait entendre des problèmes gastriques, raison de la présence desdits ἐγκρίδες, remèdes supposés ? C'est possible, mais peu nous importe : concentrons-nous sur Aristophane !\\
%Ce fragment des \italique{Danaïdes} est à rapprocher de Pollux (VII 199 $=$ fr. 276(265) K.-A.) : %
%\citComm{ἐν δ᾿ Ἀριστοφάνης Δαναΐσι συρμαιοπῶλαι καὶ ἐγκριδοπῶλαι.\\
%Aristophane dans ses \italique{Danaïdes} [emploie les mots de] "vendeurs de syrméa" (syrméa $=$ purgatif, voir le commentaire au vers 857) et de "vendeurs d'encrides".}%
%\eff{Une fois tous ces éléments recoupés, il devient aisé de saisir l'idée du premier fragment : le "μήτʼ αρμα" inintelligible est en fait un κάθαρμα, un \italique{purgatif}. Vous me direz qu'il y a loin de μήτʼ αρμα à κάθαρμα ? détrompez-vous : on trouve de nombreuses abréviations dans les manuscrits ; les prépositions sont bien-sûr parmi les plus abrégées et κατά $=$ κ`  quand μετά $=$ μ`. Il n'y a donc, à la base, qu'un kappa mélu en mu\marqueur{5}.
%Supposer\marqueur{6} que dans le manuscrit copié les kappa étaient semblables aux mu permet de corriger l'autre aberration du fragment : καὶ ποιητήν est une mauvaise transcription de μαιοπώλην du mot συρμαιοπώλην (le συρ ayant sauté par haplographie, comme on dit, précédé qu'il est du σιν de Δαναίσιν).\\
%On obtient :%
%\citComm{Ἀριστοφάνης δʼ ἐν Δαναίσιν· συρμαιοπώλην φησιν αὐτῶν εἶναι ἐν τούτοις κάθαρμα εἶναι (ἐγκριδοπώλην).}%
%À la lettre, ce n'est toujours pas grec. Mais le sens général se laisse dégager : il s'agit d'une querelle entre les différents vendeurs ; celui d'encrides (ἐγκριδοπώλης) -- à moins que ce ne soit le vendeur de syrméa (συρμαιοπώλης) qui plastronne -- traite le vendeur de syrméa de purgatif (κάθαρμα). (Avec sans doute des jeu de mots sur la polysémie.)}%
}%
\debutNote{570}%
\note{1}{Suidas (π 3192) nous apprend que c'est une sorte de gâteau fait à partir de miel bouilli et de sésame : Πυραμοῦς: εἶδος πλακοῦντος ἐκ μέλιτος ἑφθοῦ καὶ πυροῦ πεφρυγμένου, ὡς σησαμοῦς ὁ διὰ σησαμῆς. ταῦτα δὲ ἐτίθεσαν ἆθλα τοῖς διαγρυπνηταῖς ἐν τοῖς συμποσίοις, καὶ ὁ διαγρυπνήσας μέχρι τῆς ἔω ταῦτα ἐλάμβανεν ἔπαθλον. ἡμέτερος ὁ πυραμοῦς.}%
\note{2}{On en ramasse à la pelle de ces exemples de déjections consécutives à une frayeur ; Taillardat (§309) en cite bon nombre : \italique{Assemblée des femmes}, v. 1060 ; \italique{Acharniens}, v. 350 ; et d'autres, mais la liste n'est toujours pas exhaustive : on pourrait encore la compléter avec les \italique{Guêpes}, v. 627 ; \italique{Oiseaux}, v. 66 ; etc. etc. Signalons la même promesse de relâchement des intestins après correction dans \italique{Lysistrata}, v. 440 : \grec{ἐπιχεσεῖ πατούμενος}.}%
\note{3}{Hypothèse, certes corroborée par les textes d'Aristophane que nous allons voir mais point testée sur mon propre appareil digestif. Est-ce digne d'un disciple de Willems, ce grand amateur de \italique{Realia} ?}%
\note{4}{En passant, Euripide use du même jeu de mots scatologique : au vers 300 de son \italique{Ion} : σὺν ἀνδρί $[$Créuse dit à Ion qu'elle est venue avec son mari, mais... ] σηκοὺς δ᾽ ἐκστρέφει Τροφωνίου... que le mari, Xuthus, était d'ailleurs véritablement allé consulter. Le tour était proverbial et Érasme l'a recensé (\italique{Adages} n° 677)}%
\note{5}{Pollux (VII 199 $=$ fr. 276(265) Kassel-Austin). Cf. aussi le passage -- corrompu -- d'Athénée (XIV, 645e) où il est question de l'ἐγκριδοπώλης dans les \italique{Danaïdes}.}%
\note{6}{Athénée (XIV, 645e) : ΕΓΚΡΙΔΕΣ. πεμμάτιον ἑψόμενον ἐν ἐλαίῳ καὶ μετὰ τοῦτο μελιτούμενον.}%
%\note{5}{Pour tout ce qui touche à la paléographie, je vais chercher les informations principalement dans le livre d'Édouard Thomson, \italique{Handbook of Greek and Latin Palaeography} (1906). Les abréviations et autres contractions en sont le septième chapitre (pp. 86 sqq.). On trouvera un utile tableau de l'évolution des cursives grecques à la page 149. Les abréviations comme κατά $=$ κ` ou μετά $=$ μ` ainsi que des kappa très semblables aux mu laissent à penser que le manuscrit mal interprété était contemporain d'Athénée : \nieme{II} siècle ap. J.-C.}%
%\note{6}{Supposition gratuite. Il faudrait qu'un fin connaisseur des \italique{Deipnosophistes} confirme la présence d'erreurs de ce genre ou bien infirme l'hypothèse. Pour une introduction à la tradition manuscrite d'Athénée, l'on se reportera avec intérêt à l'article de l'un de ces fins connaisseurs : Benoît Louyest, "L'épitomé du \italique{Banquet des sophistes} d'Athénée", in \italique{Rursus} 8 (2012).}%
\\\passeLigne}%
\end{groupe}%
%% v. 571-573
\begin{groupe}%
\personnageGrec{Κορυφαῖος}%
\tetramIambCat{παύσασθε λοιδορούμεναι· καὶ γὰρ γυνή τις ἡμῖν\apparat{571}{(Χορ.) R : \lecon{(Κορ.)} von Velsen \espApp παύσασθαι R : \lecon{παύσασθε} Μυ2}\\
ἐσπουδακυῖα προστρέχει. πρὶν οὖν ὁμοῦ γενέσθαι,%\testimonia{571-2}{(πρὶν οὖν -- πυθώμεθ᾽) Suidas ο 292 (κοσμίως omisso)}
\\
σιγᾶθ᾽, ἵν᾽ αὐτῆς κοσμίως πυθώμεθ᾽ ἅττα λέξει.\\\passeLigne}%
\end{groupe}%
%% v. 574-581
\begin{groupe}%
\personnageGrec{Κλεισθένης}%
\typeVers{trim. iamb.}\trimIamb{φίλαι γυναῖκες ξυγγενεῖς τοὐμοῦ τρόπου·%\testimonia{574}{(ξυγγενεῖς τοὐμοῦ τρόπου) Suidas ξ 83}
\\
ὅτι μὲν φίλος εἴμ᾽ ὑμῖν, ἐπίδηλος ταῖς γνάθοις·\\
γυναικομανῶ γὰρ προξενῶ θ᾽ ὑμῶν ἀεί.%\testimonia{576}{Suidas  π 2541 προξενῶ θ' ὑμᾶς ἀεί}
\apparat{576}{\lecon{ὑμῶν} R : ὑμᾶς Suid. π 2541}%
\commentaire{575-6}{}{Sur la conservation de ces deux vers que d'aucuns auraient voulu biffer, voir la note de Willems n° 17.% ; cette note est surtout instructive pour ses remarques de bas de page à propos de mesquins vols d'émendations que certains parmi les plus connus hellénistes se permettaient, pillant sans vergogne leurs confrères ne jouissant pas de la même notoriété. Ô le joli monde de ces esprits supérieurs versés dans la recherche.
}\\
καὶ νῦν ἀκούσας πρᾶγμα περὶ ὑμῶν μέγα\\
ὀλίγῳ τι πρότερον κατ᾽ ἀγορὰν λαλούμενον,
\commentaire{578}{ὀλίγῳ τι πρότερον κατ᾽ ἀγορὰν λαλούμενον.}{Clisthène rapporte les cancans de l'Agora, qui font grand cas de ce beau-père d'Euripide introduit chez les femmes pou plaider sa cause. Ce n'est certainement pas Euripide qui est à l'origine de la fuite et ce pour des raisons évidentes : \ \ \ 1° Si Aristophane, l'avait imaginé ainsi, il nous l'aurait fait comprendre directement ou par quelque allusion \ \ \ 2° Euripide n'a aucun intérêt à voir son complice être démasqué, qui ne pourra donc plus plaider sa cause \ \ \ 3° Euripide a beaucoup de défauts, mais pas celui d'être ingrat au point d'envoyer à une mort certaine son dévoué parent. \ \ \ Savoir quel domestique d'Agathon, quel voisin, ayant aperçu Mnésiloque attifé en femme, sera allé colporter la nouvelle sur l'Agora... voilà bien des détails sans importance sur lesquels Aristophane n'a pas daigné s'arrêter ; nous aurions fait de même si les professeurs parmi les plus admis ne s'étaient lancés dans ce genre d'extravagantes et vaines conjectures.}\\
ἥκω φράσων τοῦτ᾽ ἀγγελῶν θ᾽ ὑμῖν, ἵνα
\commentaire{579-87}{ἥκω φράσων...}{Passées les retrouvailles, Clisthène prend un ton paratragique -- l'heure est grave ! Noter à ce propos les redoublements\marqueur{1} : φράσων/ἀγγελῶν, σκοπῆτε/τηρῆτε, δεινὸν/μέγα, βουλεύοισθε/μέλλοιτε.}%
\debutNote{579-87}%
\note{1}{Redoublements propres au style eschyléen et raillés par l'Euripide des \italique{Grenouilles} (vv. 1153 sqq.).}\\
σκοπῆτε καὶ τηρῆτε μὴ καὶ προσπέσῃ\apparat{580}{σκοπεἴτε καὶ τηρεἴτε R : \lecon{σκοπῆτε καὶ τηρῆτε} Biset \espApp καὶ μὴ R : \lecon{μὴ καὶ} Porson (\italique{Aristophanica} p. 216) \espApp \textbf{581} ἡμῖν R : \lecon{ὑμῖν} \Mon \espApp ἀφράκτοις R : \lecon{ἀφάρκτοις} Dindorf (1830)}\\
ὑμῖν ἀφάρκτοις πρᾶγμα δεινὸν καὶ μέγα.\\\passeLigne}%
\end{groupe}%
%% v. 582-583
\begin{groupe}%
\personnageGrec{Κορυφαῖος}%
\trimIamb{τί δ᾽ ἔστιν ὦ παῖ; παῖδα γάρ σ᾽ εἰκὸς καλεῖν,\\
ἕως ἂν οὕτως τὰς γνάθους ψιλὰς ἔχῃς.%
\commentaire{582-3}{τί δ᾽ ἔστιν ὦ παῖ; παῖδα γάρ σ᾽ εἰκὸς καλεῖν, /	ἕως ἂν οὕτως τὰς γνάθους ψιλὰς ἔχῃς.}{Pique de la coryphée analogue à celle du chœur des \italique{Guêpes}, aux vers 1297-1298 : %
\citComm{Χορός – τί δ᾽ ἔστιν ὦ παῖ; παῖδα γάρ, κἂν ᾖ γέρων,\\
καλεῖν δίκαιον ὅστις ἂν πληγὰς λάβῃ.}%
}%
\\\passeLigne}%
\end{groupe}%
%% v. 584-585
\begin{groupe}%
\personnageGrec{Κλεισθένης}%
\trimIamb{Εὐριπίδην φάσ᾽ ἄνδρα κηδεστήν τινα\apparat{584}{ἔφας᾽ R : \lecon{φάσ᾽} Zanetti}\\
αὑτοῦ γέροντα δεῦρ᾽ ἀναπέμψαι τήμερον.\apparat{585}{αὐτου R : \lecon{αὑτοῦ} Biset}%
\commentaire{585}{ἀναπέμψαι.}{Le suffixe `ἀνα' s'explique par le fait que le Thesmophorion se trouvait probablement sur la colline de la Pnyx. (Cf. le vers 658.)}\\\passeLigne}%
\end{groupe}%
%% v. 586
\begin{groupe}%
\personnageGrec{Κορυφαῖος}%
\trimIamb{πρὸς ποῖον ἔργον ἢ τίνος γνώμης χάριν;%\commentaire{586}{πρὸς ποῖον ἔργον ἢ τίνος γνώμης χάριν;}{Vers paratragique.}
\\\passeLigne}%
\end{groupe}%
%% v. 587-588
\begin{groupe}%
\personnageGrec{Κλεισθένης}%
\trimIamb{ἵν᾽ ἅττα βουλεύοισθε καὶ μέλλοιτε δρᾶν,\\
%
ἐκεῖνος εἴη τῶν λόγων κατάσκοπος.\\\passeLigne}%
\end{groupe}%
%% v. 589
\begin{groupe}%
\personnageGrec{Κορυφαῖος}%
\trimIamb{καὶ πῶς λέληθεν ἐν γυναιξὶν ὢν ἀνήρ;\\\passeLigne}%
\end{groupe}%
\endnumbering
\end{Leftside}
\begin{Rightside}
\beginnumbering
%% v. 551-552a
\begin{groupe}%
\personnageFr{Mika}%
\tetramIambCatFr{\ligne{551}Vous entendez, femmes, comment la coquine\\
nous arrange toutes, une fois de plus.\\\passeLigne}%
\end{groupe}%
%% v. 552b-553
\begin{groupe}%
\personnageFr{Mnésiloque}%
\tetramIambCatFr{\ligne{552}\espace{2.8cm}Et, par Zeus, je n'ai pas encore\\
dit tout ce que je sais. En voulez-vous davantage ?\\\passeLigne}%
\end{groupe}%
%% v. 554
\begin{groupe}%
\personnageFr{Mika}%
\tetramIambCatFr{Tu ne pourrais. Ce que tu savais, tu l'as épanché.\testimonia{554}{(ᾔδησθ᾽) \EtymGen (A) ᾔδεισθα}\\\passeLigne}%
\end{groupe}%
%% v. 555-557a
\begin{groupe}%
\personnageFr{Mnésiloque}%
\tetramIambCatFr{Non, par Zeus, pas le dix-millième de ce que nous faisons.
Car enfin, je n'ai pas dit – tu vois ? – comment nous prenons des strigiles
pour ensuite nous siphonner \fin le con...\testimonia{557}{(σιφωνίζομεν τὸν σῖτον) Suidas σ 512 \espTest (σιφωνίζομεν) Pollux VI 19}\\\passeLigne}%
\end{groupe}%
%% v. 557b
\begin{groupe}%
\personnageFr{Mika}%
\tetramIambCatFr{\ligne{557}\espace{4.5cm}La peste soit de toi.\\\passeLigne}%
\end{groupe}%
%% v. 558-559a
\begin{groupe}%
\personnageFr{Mnésiloque}%
\tetramIambCatFr{... que les viandes des Apaturies% \commentaire{558}{τὰ κρέ᾽ ἐξ Ἀπατουρίων}{Cf. note sur la Paix, v. 893, p. 112, ci-devant}
, nous les distribuons à nos procureuses,
puis que nous accusons la belette...\\\passeLigne}%
\end{groupe}%
%% v. 559b
\begin{groupe}%
\personnageFr{Mika}%
\tetramIambCatFr{\ligne{559}\espace{3.9cm}Merci de moi ! Tu radotes.\\\passeLigne}%
\end{groupe}%
%% v. 560-562a
\begin{groupe}%
\personnageFr{Mnésiloque}%
\tetramIambCatFr{... ni que telle autre occit son époux d'un coup de hache,\testimonia{560}{Suidas κ 996 (ἑτέρα pro γυνὴ et κατεσπόδησε vel κατεσπόδωσεν)}
je ne l'ai pas dit ; ni qu'une autre rendit fou son mari par des drogues,
ni que naguère sous la \fin baignoire...\\\passeLigne}%
\end{groupe}%
%% v. 562b
\begin{groupe}%
\personnageFr{Mika}%
\tetramIambCatFr{\ligne{562}\espace{6cm}Crève...\\\passeLigne}%
\end{groupe}%
%% v. 563a
\begin{groupe}%
\personnageFr{Mnésiloque}%
\tetramIambCatFr{... cette Acharnienne enfouit son père.\\\passeLigne}%
\end{groupe}%
%% v. 563b
\begin{groupe}%
\personnageFr{Mika}%
\tetramIambCatFr{\ligne{563}\espace{3.5cm}Supporterons-nous ce langage ?\\\passeLigne}%
\end{groupe}%
%% v. 564-565
\begin{groupe}%
\personnageFr{Mnésiloque}%
\tetramIambCatFr{Ni que toi, quand ta servante mit au jour un garçon,\\
tu substituas l'enfant, et livras ta petite fille.\\\passeLigne}%
\end{groupe}%
%% v. 566-567a
\begin{groupe}%
\personnageFr{Mika}%
\tetramIambCatFr{Non, par les deux déesses, tu ne l'auras pas dit gratuitement.\testimonia{566}{Suidas κ 730 (sed καταπροίξῃ pro καταπροίξει)}
Je t'arracherai la toison.\\\passeLigne}%
\end{groupe}%
%% v. 567b
\begin{groupe}%
\personnageFr{Mnésiloque}%
\tetramIambCatFr{\ligne{567}\espace{3.3cm}Non, par Zeus, tu ne me toucheras pas.\testimonia{567}{(ἐκποκιῶ σου τὰς ποκάδας) Suidas ε 588 (ἐκπιῶ cod. A) ; Zonaras 672 ἐκποκιῶ \espTest (ἐκποκιῶ) Heschius ε 1637 ἐκποκιῶ ; \EtymMagnum 323.25 \espTest (ποκάδας) Hesychius π 2750 ποκάδες $=$ Photius π 1010 π\{λ\}οκάδες}\\\passeLigne}%
\end{groupe}%
%% v. 568a
\begin{groupe}%
\personnageFr{Mika}%
\tetramIambCatFr{Eh bien, tiens !\\\passeLigne}%
\end{groupe}%
%% v. 568b
\begin{groupe}%
\personnageFr{Mnésiloque}%
\tetramIambCatFr{\ligne{568}\espace{2.1cm}Eh bien, tiens !\\\passeLigne}%
\end{groupe}%
%% v. 568c
\begin{groupe}%
\personnageFr{Mika}%
\tetramIambCatFr{\ligne{568}\espace{4.3cm}Prends mon manteau, Philista.\\\passeLigne}%
\end{groupe}%
%% v. 569a
\begin{groupe}%
\personnageFr{Mnésiloque}%
\tetramIambCatFr{Atteins-moi seulement, et par Artémis...\\\passeLigne}%
\end{groupe}%
%% v. 569b
\begin{groupe}%
\personnageFr{Mika}%
\tetramIambCatFr{\ligne{569}\espace{6.2cm}Que feras-tu ?\\\passeLigne}%
\end{groupe}%
%% v. 570
\begin{groupe}%
\personnageFr{Mnésiloque}%
\tetramIambCatFr{Le gâteau de sésame que tu as dévoré, je te le ferai \fin évacuer.\testimonia{570}{(χεσεῖν) \EpimAlphab ε 201}\\\passeLigne}%
\end{groupe}%
%% v. 571-573
\begin{groupe}%
\personnageFr{Coryphée}%
\tetramIambCatFr{Trêve d'invectives. Car une femme\testimonia{571-2}{(πρὶν οὖν -- πυθώμεθ᾽) Suidas ο 292 (κοσμίως omisso)}
accourt vers nous en toute hâte. Avant qu'elle soit proche,
faites silence. Écoutons en bonne contenance ce qu'elle va nous dire.\\\passeLigne}%
\end{groupe}%
%% v. 574-581
\begin{groupe}%
\personnageFr{Clisthène}%
\trimIambFr{Femmes aimées, qui m'êtes unies par la parenté des mœurs,\testimonia{574}{(ξυγγενεῖς τοὐμοῦ τρόπου) Suidas ξ 83}
que je vous suis ami, mes joues le prouvent assez.
Oui, j'aime les femmes à la rage, et toujours je m'emploie pour vous.\testimonia{576}{Suidas  π 2541 προξενῶ θ' ὑμᾶς ἀεί}
Aussi, ayant appris sur votre sujet une chose capitale,
dont on jasait tout à l'heure sur l'agora,
je suis venu la raconter et vous donner avis
de veiller et d'être en garde, de crainte
qu'une grave et fâcheuse affaire ne vous tombe sur les bras au dépourvu.\\\passeLigne}%
\end{groupe}%
%% v. 582-583
\begin{groupe}%
\personnageFr{Coryphée}%
\trimIambFr{Qu'y a-t-il mon enfant ? Car il est naturel que je t'appelle enfant,
tant que tu auras ainsi les joues lisses.\\\passeLigne}%
\end{groupe}%
%% v. 584-585
\begin{groupe}%
\personnageFr{Clisthène}%
\trimIambFr{On dit qu'Euripide a envoyé ici même, aujourd'hui,\\
un de ses parents, un vieillard.\\\passeLigne}%
\end{groupe}%
%% v. 586
\begin{groupe}%
\personnageFr{Coryphée}%
\trimIambFr{Pourquoi faire, dans quel dessein ?\\\passeLigne}%
\end{groupe}%
%% v. 587-588
\begin{groupe}%
\personnageFr{Clisthène}%
\trimIambFr{À cette fin que, quoi que vous agitiez ou décidiez de faire,
il puisse épier vos discours.\\\passeLigne}%
\end{groupe}%
%% v. 589
\begin{groupe}%
\personnageFr{Coryphée}%
\trimIambFr{Et comment s'est-il dissimulé au milieu de femmes ? Un homme !\\\passeLigne}%
\end{groupe}%
\endnumbering
\end{Rightside}
\end{pages}
\Pages
\clearpage\ifodd\value{page}\hbox{}\newpage\fi
\enTeteG{ΕΠΕΙΣΟΔΙΟΝ Β' --- Codex Ravennas 429, f. 169v (Classense 171v) --- vv. 590–634}
\begin{adjustwidth}{-1.5cm}{-3.0cm}
\begin{adjustbox}{max totalheight=21.0cm, max width=16.5cm, keepaspectratio}
\begin{minipage}{16.5cm}
{\fontsize{7.2pt}{8.5pt}\selectfont
\noindent\textbf{Folio 15 --- Codex Ravennas 429, \texttt{f. 169v} (\texttt{Classense 171v}) : ΕΠΕΙΣΟΔΙΟΝ Β', vv. 590–634}\hfill \textit{Témoins :} $\text{R}^1, \Rpc, \Rdeux$ (\Mon)\par\vspace{2pt}
\hrule height 0.4pt\vspace{2.5pt}
}
{\fontsize{8.1pt}{9.5pt}\selectfont
\ligneRav{590}{κλεισθ/}{ἄφευσεν αὐτὸν \difference{κἀπέτιλ᾽} Εὐριπίδης}{\textbf{590} \lecon{ἄφευσεν} R : ἀφηῦσεν Meineke \espApp καπέτιλλ᾽ R : \lecon{κἀπέτιλ᾽} Bekker}
\ligneRav{}{}{καὶ τἄλλ᾽ ἅπανθ᾽ ὥσπερ γυναῖκ᾽ ἐσκεύασεν·}{}
\ligneRav{}{μ\auDessus{ν}{Η}}{πείθεσθε τούτῳ ταῦτα; τίς δ᾽ οὕτως ἀνὴρ}{}
\ligneRav{}{}{ἠλίθιος ὅστις τιλλόμενος ἠνείχετο;}{}
\ligneRav{594}{}{οὐκ \difference{οἴομαι 'γωγ᾽} ὦ πολυτιμήτω θεώ·}{\textbf{594} οἴομ᾽ ἔγωγ᾽ R : \lecon{οἴομαι 'γωγ᾽} anonymus Parisinus}
\ligneRav{595}{κλεισθ/}{ληρεῖς· ἐγὼ γὰρ οὐκ ἂν ἦλθον ἀγγελῶν,}{}
\ligneRav{596}{}{εἰ μὴ \difference{᾽πεπύσμην} \difference{ταῦτα} τῶν σάφ᾽ εἰδότων·}{\textbf{596} πεπύσμην R, \PsiUn : \lecon{᾽πεπύσμην} Faber \espApp ταυτὶ R : \lecon{ταῦτα} Bentley}
\ligneRav{}{χορ/}{τὸ πρᾶγμα τουτὶ δεινὸν εἰσαγγέλλεται·}{}
\ligneRav{598}{}{ἀλλ᾽ ὦ γυναῖκες οὐκ \difference{ἐλινύειν} ἐχρῆν,}{\textbf{598} ἐλιννύειν R : \lecon{ἐλινύειν} Bourdin}
\ligneRav{}{}{ἀλλὰ σκοπεῖν τὸν ἄνδρα καὶ ζητεῖν ὅπου}{}
\ligneRav{600}{}{λέληθεν \difference{ἡμᾶς} κρυπτὸς ἐγκαθήμενος·}{\textbf{600} ὑμᾶς R : \lecon{ἡμᾶς} \Mon}
\ligneRav{601}{}{καὶ σὺ \difference{ξυνέξευρ᾽} αὐτόν, ὡς ἂν τὴν χάριν}{\textbf{601} συνέξευρ᾽ R, Suid. π 2540 codd. AG : \lecon{ξυνέξευρ᾽} Suid. codd. FVM}
\ligneRav{602}{}{ταύτην τε κἀκείνην ἔχῃς ὦ πρόξενε·}{\textbf{602} \lecon{ταύτην} R : αὐτὴν Suid.}
\ligneRav{}{κλεισθ/}{φέρ᾽ ἴδω· τίς ἡ πρώτη σύ; : ποῖ τις τρέψεται;}{}
\ligneRav{604}{κλεισθ/}{ζητητέαι γάρ ἐστε. : κακοδαίμων ἐγώ·}{\textbf{604} \lecon{ζητητέαι γάρ ἐστε} R : ζητητέον τἄρ᾽ ἐστί Maas \espApp 604 ante 603 transp. Maas, post 607 collocare malui}
\ligneRav{605}{\auDessus{γυ}{ν}}{ἔμ᾽ \difference{ἥτις} \difference{εἴμ᾽ ἤρου}; Κλεωνύμου γυνή·}{\textbf{605} εἰτις ἤρου; R : \lecon{ἥτις} (Zanetti) \lecon{εἴμ᾽ ἤρου;} Bentley}
\ligneRav{606}{κλεισθ/}{γιγνώσκεθ᾽ ὑμεῖς ἥτις ἔσθ᾽ \difference{ἡδὶ} γυνή;}{\textbf{606} ἥδε R : ἥδ᾽ ἡ \Mon : \lecon{ἡδὶ} Grynaeus}
\ligneRav{}{χορ/}{γιγνώσκομεν δῆτ᾽. ἀλλὰ τὰς ἄλλας ἄθρει·}{}
\ligneRav{}{κλεισθ/}{ἡδὶ δὲ δὴ τίς ἐστιν ἡ τὸ παιδίον}{}
\ligneRav{609}{—}{ἔχουσα; : \difference{τίτθη νὴ Δί᾽} ἐμή. : διοίχομαι·}{\textbf{609} τίτθην η δί᾽ R : \lecon{τίτθη νὴ Δί᾽} Ellebodius}
\ligneRav{610}{κλεισθ/}{αὕτη σὺ ποῖ στρέφει; μέν᾽ αὐτοῦ. τί τὸ κακόν;}{\textbf{610} τί τὸ κακόν; continuat Clistheni R : Parenti tribuit (cum versu sequenti) \Rpc}
\ligneRav{611}{μ\auDessus{ν}{Η}}{ἔασον οὐρῆσαί μ᾽· ἀναίσχυντός \difference{τις εἶ}·}{\textbf{611} εἶ R : \lecon{τις εἶ} Küster}
\ligneRav{612}{κλεισθ/}{σὺ δ᾽ οὖν ποίει τοῦτ᾽. \difference{ἀναμενῶ} γὰρ ἐνθάδε·}{\textbf{612} ἀναμένω R : \lecon{ἀναμενῶ} Grynaeus}
\ligneRav{}{χορ/}{ἀνάμενε δῆτα καὶ σκόπει γ᾽ αὐτὴν σφόδρα·}{}
\ligneRav{}{}{μόνην γὰρ αὐτὴν ὦνερ οὐ γιγνώσκομεν·}{}
\ligneRav{615}{κλεισθ/}{\difference{πολύν} γε χρόνον οὐρεῖς σύ. : νὴ Δί᾽ ὦ μέλε·}{\textbf{615} πολύ R : \lecon{πολύν} Gelenius}
\ligneRav{}{}{στραγγουριῶ γάρ· ἐχθὲς ἔφαγον κάρδαμα·}{}
\ligneRav{}{κλεισθ/}{τί καρδαμίζεις; οὐ βαδιεῖ δεῦρ᾽ ὡς ἐμέ;}{}
\ligneRav{}{μ\auDessus{ν}{Η}}{τί δῆτά μ᾽ ἕλκεις ἀσθενοῦσαν; : εἰπέ μοι,}{}
\ligneRav{619}{—}{τίς \difference{ἔστ᾽} ἀνήρ σοι; : τὸν ἐμὸν ἄνδρα πυνθάνει;}{\textbf{619} ἔσθ᾽ R : \lecon{ἔστ᾽} \Mon}
\ligneRav{620}{}{τὸν δεῖνα γιγνώσκεις, τὸν ἐκ Κοθωκιδῶν·}{\textbf{620} add. signum interrogationis Canini}
\ligneRav{621}{κλεισθ/}{τὸν δεῖνα; ποῖον; : ἔσθ᾽ ὁ δεῖν᾽, ὃς καί ποτε}{\textbf{621} (Κλεισθ.) continuat R : \lecon{(Μνη.)} \Mon}
\ligneRav{}{—}{τὸν δεῖνα τὸν τοῦ δεῖνα— : ληρεῖν μοι δοκεῖς·}{}
\ligneRav{}{—}{ἀνῆλθες ἤδη δεῦρο πρότερον; : νὴ Δία}{}
\ligneRav{624}{—}{\difference{ὁσέτη} γε. : καὶ τίς \difference{σοὐστὶ} συσκηνήτρια;}{\textbf{624} ὅς᾽ ἔτη R : \lecon{ὁσέτη} G. H. Schaefer ad Gregor. Cor. p.65 \espApp σοὐστὶν R : \lecon{σοὐστὶ} Dindorf (σοῦ ᾽στὶ iam Biset)}
\ligneRav{625}{μ\auDessus{ν}{Η}}{ἡ δεῖν᾽ \difference{ἔμοιγ᾽}. : οἴμοι τάλας. οὐδὲν λέγεις·}{\textbf{625} ἔμοιγε R : \lecon{ἔμοιγ᾽} edd. pr.}
\ligneRav{}{\auDessus{γυ}{ν}}{ἄπελθ᾽. ἐγὼ γὰρ βασανιῶ ταύτην καλῶς}{}
\ligneRav{}{}{ἐκ τῶν ἱερῶν τῶν πέρυσι· σὺ δ᾽ ἀπόστηθί μοι,}{}
\ligneRav{}{}{ἵνα μὴ 'πακούσῃς ὢν ἀνήρ. σὺ δ᾽ εἰπέ μοι}{}
\ligneRav{}{}{ὅ τι πρῶτον ἡμῖν τῶν ἱερῶν ἐδείκνυτο·}{}
\ligneRav{630}{μ\auDessus{ν}{Η}}{τί δ᾽ \expo{μέντοι} πρῶτον ἦν; : ἐπίνομεν·}{\textbf{630} \lecon{μέντοι} \Rpc in mg., Suid. π 2563 : om. \Rac, \Mon}
\ligneRav{631}{\auDessus{γυ}{ν}}{τί δὲ \difference{μετὰ} τοῦτο δεύτερον; : προὐπίνομεν·}{\textbf{631} \lecon{δὲ} R, Suid. π 2563 : δαὶ Bentley \espApp με R : \lecon{μετὰ} Suid.}
\ligneRav{632}{\auDessus{γυ}{ν}}{ταυτὶ μὲν ἤκουσάς τινος· \difference{τί δ᾽ ἦν τρίτον};}{\textbf{632} τί δὲ τρίτον R : τί δ᾽ αὖ Bentley : τί δαὶ τρίτον Elmsley : \lecon{τί δ᾽ ἦν τρίτον} Teubner}
\ligneRav{}{μ\auDessus{ν}{Η}}{σκάφιον Ξένυλλ᾽ ᾔτησεν· οὐ γὰρ ἦν ἀμίς·}{}
\ligneRav{634}{\auDessus{γυ}{ν}}{ξυνέπινε \difference{δεῦρ᾽ ὦ} (δευρ ὦ \Rdeux) Ξενύλλ᾽· οὐκ ἦν ἀμίς·}{\textbf{634} δεῦρο \Rac, \Mon : \lecon{δευρ ὦ} \Rdeux \espApp Κλεισόσθενες R : \lecon{Κλείσθενες} Bonini}
}
\vspace{3pt}\noindent\rule{4.5cm}{0.4pt}\par\vspace{1.5pt}
{\fontsize{7.0pt}{8.3pt}\selectfont
\noindent\textbf{Loci paralleli (vv. 590–634) ---} \textbf{598} (ἐλινύειν) Hesychius ε 2095 ; cf. Moeris ε 61 (ἐλινύων $=$ ἀναπαυόμενος) $\sim$ Hesychius ε 2097 \espTest \textbf{598} (ἐλινύειν) Hesychius ε 2095, 2097 ; cf. Moeris ε 61 (ἐλινύων $=$ ἀναπαυόμενος)% $\sim$ Hesychius ε 2097 \espTest \textbf{601-2} Suidas π 2540 \espTest \textbf{616} Suidas σ 1163 (χθὲς γὰρ pro γάρ ἐχθὲς) \espTest \textbf{620} (Κοθωκιδῶν) Photius κ 859 Κοθωκίδαι (ap Scholium) $\sim$ Hesychius κ 3218 \espTest \textbf{630-1} Suidas π 2563 (sed προπίνομεν pro προὐ- cod. V) \espTest \textbf{633} Pollux X 45 ἐν Πολυίδῳ (sic) Ἀριετοφάνης· σκάφιον Ξένυλλ᾽ ἤτησεν· οὐ γὰρ ἦν ἁμίς \espTest Photius σ 286 σκάφιον\par
}
\vspace{2.5pt}\noindent\rule{4.5cm}{0.4pt}\par\vspace{1.5pt}
{\fontsize{6.9pt}{8.2pt}\selectfont
\noindent\textbf{Suite (en commentaire) de l'apparat critique \& notes paléographiques (\texttt{f. 169v}) ---} \textbf{Témoin papyrologique :} \textbf{\PsiUn{} (\italique{PSI} XI 1194)} couvre les \textbf{vv. 594–596}. --- \textbf{Supplément marginal au v. 630 :} La particule \lecon{μέντοι} (confirmée par la Souda $\pi\ 2563$) avait été omise dans la colonne par \Rac{} (et reste \textbf{omise dans \Mon{}}) ; elle a été rétablie dans la marge par \Rpc{}. --- \textbf{Dittographie syllabique au v. 634 :} En bas de page, $\text{R}^1$ écrit Κλεισόσθενες (dittographie de $-\sigma\text{ο}-$ dans Κλείσθενες, corrigée par Bonini), tandis que \Rdeux{} corrige δεῦρο (\Rac{}, \Mon{}) en \lecon{δευρ ὦ}. --- \textbf{v. 609 :} τίτθην η δί' R (médivision de $\nu$ final) : \lecon{τίτθη νὴ Δί'} Ellebodius. --- \textbf{v. 630 :} \lecon{μέντοι} \Rpc{} in mg., Souda $\pi\ 2563$ : om. $\text{R}^{\text{ac}}, \text{M}$.\par
}
\end{minipage}
\end{adjustbox}
\end{adjustwidth}
\newpage
\enTeteD{Codex Ravennas 429 (Biblioteca Classense) --- Fac-similé du f. 169v (vv. 590–634)}
\begin{adjustwidth}{-1.5cm}{-3.0cm}
\begin{center}
\includegraphics[width=16.5cm,height=20.5cm,keepaspectratio]{/workspaces/Thesmophories/dossier_ravennas_thesmophories/folios_individuels_ff162v-177r/15_f169v_classense171v_vv0590-0634.jpg}\par
\vspace{2pt}
{\scriptsize\textbf{Codex Ravennas 429, f. 169v (Classense 171v)} --- ΕΠΕΙΣΟΔΙΟΝ Β', \textbf{vv. 590–634} ($\text{R}^1, \Rpc, \Rdeux$ (\Mon)).}
\end{center}
\end{adjustwidth}
\newpage
\clearpage\ifodd\value{page}\hbox{}\newpage\fi
\enTeteG{ΘΕΣΜΟΦΟΡΙΑΖΟΥΣΑΙ --- Texte grec établi (f. 169v, vv. 590–634)}
\enTeteD{LES THESMOPHORIES --- Traduction française (f. 169v, vv. 590–634)}
\begin{pages}
\begin{Leftside}
\beginnumbering
%% v. 590-591
\begin{groupe}%
\personnageGrec{Κλεισθένης}%
\trimIamb{ἄφευσεν αὐτὸν κἀπέτιλ᾽ Εὐριπίδης\apparat{590}{\lecon{ἄφευσεν} R : ἀφηῦσεν Meineke \espApp καπέτιλλ᾽ R : \lecon{κἀπέτιλ᾽} Bekker}\\
καὶ τἄλλ᾽ ἅπανθ᾽ ὥσπερ γυναῖκ᾽ ἐσκεύασεν.\\\passeLigne}%
\end{groupe}%
%% v. 592-594
\begin{groupe}%
\personnageGrec{Μνησίλοχος}%
\trimIamb{πείθεσθε τούτῳ ταῦτα; τίς δ᾽ οὕτως ἀνὴρ\\
ἠλίθιος ὅστις τιλλόμενος ἠνείχετο;
\commentaire{593}{ἠλίθιος ὅστις τιλλόμενος ἠνείχετο;}{Willems, chevalier servant du \italique{Ravennas}, défend encore une fois la leçon du manuscrit dans sa \nieme{18} note, à retrouver à la fin.}\\
οὐκ οἴομαι 'γωγ᾽ ὦ πολυτιμήτω θεώ.\apparat{594}{οἴομ᾽ ἔγωγ᾽ R : \lecon{οἴομαι 'γωγ᾽} anonymus Parisinus}\\\passeLigne}%
\end{groupe}%
%% v. 595-596
\begin{groupe}%
\personnageGrec{Κλεισθένης}%
\trimIamb{ληρεῖς· ἐγὼ γὰρ οὐκ ἂν ἦλθον ἀγγελῶν,\\
εἰ μὴ ᾽πεπύσμην ταῦτα τῶν σάφ᾽ εἰδότων.\apparat{596}{πεπύσμην R, \PsiUn : \lecon{᾽πεπύσμην} Faber \espApp ταυτὶ R : \lecon{ταῦτα} Bentley}\\\passeLigne}%
\end{groupe}%
%% v. 597-602
\begin{groupe}%
\personnageGrec{Κορυφαῖος}%
\trimIamb{τὸ πρᾶγμα τουτὶ δεινὸν εἰσαγγέλλεται.\\
ἀλλ᾽ ὦ γυναῖκες οὐκ ἐλινύειν ἐχρῆν,%\testimonia{598}{(ἐλινύειν) Hesychius ε 2095 ; cf. Moeris ε 61 (ἐλινύων $=$ ἀναπαυόμενος) $\sim$ Hesychius ε 2097}
\apparat{598}{ἐλιννύειν R : \lecon{ἐλινύειν} Bourdin}\\%, cf. Moeris \italique{Gram.} ε 61, Hesychius ε 2097}\\
ἀλλὰ σκοπεῖν τὸν ἄνδρα καὶ ζητεῖν ὅπου\\
λέληθεν ἡμᾶς κρυπτὸς ἐγκαθήμενος.\apparat{600}{ὑμᾶς R : \lecon{ἡμᾶς} \Mon}\\
καὶ σὺ ξυνέξευρ᾽ αὐτόν, ὡς ἂν τὴν χάριν%\testimonia{601-2}{Suidas π 2540}
\apparat{601}{συνέξευρ᾽ R, Suid. π 2540 codd. AG : \lecon{ξυνέξευρ᾽} Suid. codd. FVM \ligneApp{602} \lecon{ταύτην} R : αὐτὴν Suid.}\\
ταύτην τε κἀκείνην ἔχῃς ὦ πρόξενε.%
\commentaire{602}{ὦ πρόξενε.}{Les proxènes ou hôtes publics, agents officiels dont les fonctions consistaient principalement à recevoir et à héberger les délégués de l'État dont ils exerçaient la proxénie. (Note de Willems.)}\\\passeLigne}%
\end{groupe}%
%% v. 603a
\begin{groupe}%
\personnageGrec{Κλεισθένης}%
\trimIamb{φέρ᾽ ἴδω· τίς ἡ πρώτη σύ;\\\passeLigne}%
\end{groupe}%
%% v. 603b
\begin{groupe}%
\personnageGrec{Μνησίλοχος}%
\trimIamb{\ligne{603}\decalage{φέρ᾽ ἴδω· τίς ἡ πρώτη σύ;}ποῖ τις τρέψεται;\\\passeLigne}%
\end{groupe}%
%% v. 604a
\begin{groupe}%
%
\personnageGrec{Κλεισθένης}%
\trimIamb{ζητητέαι γάρ ἐστε.%\apparat{604}{\lecon{ζητητέαι γάρ ἐστε} R : ζητητέον τἄρ᾽ ἐστί Maas \espApp 604 ante 603 transp. Maas, post 607 collocare malui}
\\\passeLigne}%
\end{groupe}%
%% v. 604b
\begin{groupe}%
\personnageGrec{Μνησίλοχος}%
\trimIamb{\ligne{604}\decalage{ζητητέαι γάρ ἐστε.}κακοδαίμων ἐγώ.\\\passeLigne}%
\end{groupe}%
%% v. 605
\begin{groupe}%
\personnageGrec{Μίκα}%
\trimIamb{ἔμ᾽ ἥτις εἴμ᾽ ἤρου; Κλεωνύμου γυνή.\apparat{605}{εἰτις ἤρου; R : \lecon{ἥτις} (Zanetti) \lecon{εἴμ᾽ ἤρου;} Bentley}% : ἥτις ἤρου; τοῦ Biset : alii alia}%
\commentaire{605}{Κλεωνύμου γυνή.}{L'évocation de Cléonyme est l'occasion de dresse un bref résumé des railleries à l'encontre de ce personnage parmi les plus plaisantés de l'Ancienne comédie. Cet Athénien nous est autrement connu par trois décrets qu'il proposa à la \grec{βουλή}\marqueur{1} l'an 426/5 ainsi qu'une proposition de fixer à 1000 drachmes la récompense contre des informations touchant à l'affaire de la parodie des Mystères d'Éleusis\marqueur{2} (contemporaine de celle des Hermocopides, l'an 415). Cléonyme est victime d'un tir de barrage d'Aristophane portant à la fois :\\
— au moral : sa lâcheté (il lui était arrivé un jour de jeter son bouclier pour mieux fuir. C'est surtout là-dessus que le poète ne tarit point : \italique{Cavaliers} v. 1372, \italique{Nuées} v. 353, \italique{Guêpes} v. 19, 592, \italique{Paix} v. 446, 672, 1186 (non nommé), 1275, 1295, \italique{Oiseaux} v. 290, 1475, 1481) ; sa duplicité (\italique{Acharniens} v. 88, il est traité de parjure v. 399 des \italique{Nuées}, de flatteur au v. 592 des \italique{Guêpes} (il est appelé Κολακώνυμος), de sycophante au v. 1479 des \italique{Oiseaux}) ; ses manières brusques de grand dadais (\italique{Acharniens} v. 844) ; ses mœurs peu ragoûtantes (accusé d'être un \grec{καταπύγων} au v. 675 des \italique{Nuées}) ;\\
— au physique : une laideur repoussante (v. 822 des \italique{Guêpes}, ) ; son androgynie (Socrate, dans une intention de parler exactement, emploie le féminin \grec{Κλεωνύμη} au v. 680 des \italique{Nuées}) ; une taille excessive (\italique{Acharniens} v. 88, \italique{Oiseaux}, v. 1476) ; sa gloutonnerie (\italique{Cavaliers} v. 1290, \italique{Oiseaux} v. 289).\\
Jouant sur les contrastes, Aristophane nous présente une femme de Cléonyme \italique{petite} de nom (c'est le sens de \grec{Μίκα}), virile et hardie dans les débats.\\
%Cette note sur l'occasion de tordre le cou à l'une des traditions philologiques les plus tenaces et qui laissait déjà en son temps Willems devant des abîmes de perplexité : pourquoi sont-ce les mêmes éminents professeurs qui se lancent dans des conjectures les plus extravagantes lorsqu'ils sont confrontés à d'obscurs passages allusifs et qui, face aux passages les plus clairs et dont le sens 
Revenons sur le grief le plus reproché à Cléonyme : avoir jeté son bouclier pour mieux prendre ses jambes à son cou, sa \grec{ῥίψασπις}. Il n'y a rien à y redire : le crime était théoriquement puni de mort\marqueur{3} et l'on comprend aisément que le poète ait exposé à l'opprobre publique le couard qui avait encore l'outrecuidance de venir palabrer après son forfait. Exégèse par trop simple pour nos éminents professeurs – dont le délassement consiste à extravaguer sur les passages les plus obscurs et remettre systématiquement en doute les passages les plus clairs\marqueur{4}. On ne la fait pas à ces \italique{viri docti}, comme disait Fritzsche : le dernier en date\marqueur{5} argue de ce que Cléonyme serait mort – condamné à mort – bien avant les dernière traces archéologiques de sa vie démocratique s'il avait effectivement commis ce qu'Aristophane lui reproche... avec le même raisonnement, les historiens des prochains millénaires déduiront que l'\italique{Élysée}, \italique{Matignon} et le \italique{palais Bourbon} ne pouvaient désigner que des prisons au vu des canailles hors-la-loi qui les fréquentaient ; prisons flanquées, semblera-t-il, d'une maison de santé pour ceux dont la panse menacait de crever à tout instant : le \italique{palais du Luxembourg}\marqueur{6}. Mais il n'y a pas besoin d'extrapoler : le même raisonnement intenable entraînait déjà de refuser le vers 1166 des \italique{Acharniens}\marqueur{7} : comment un bandit (un \italique{Oreste}, nom générique ; un \italique{Cartouche} ou \italique{Mandrin} chez nous) aurait-il pu sévir la nuit à Athènes alors que le banditisme était aussi puni de mort ?..
Certes, « \anglais{in interpreting the personal jokes in Aristophanes one must always be wary of the too serious or too literal explanation of the text} »\marqueur{8}, mais l'\italique{on ne détruit que ce que l'on remplace}, et pour remplacer une explication qui coule de source comme celle de nos passages, il vaut mieux présenter de solides arguments. Or M. Storey ne me fera pas accroire sans preuve\marqueur{9} que le motif pour lequel Aristophane poursuivit inlassablement de sa verve pendant des années Cléonyme fût, non pas une désertion devant l'ennemi doublée d'un abandon d'armes, mais une banale brigue pour échapper au service militaire\marqueur{10}.
}%
\debutNote{605}%
\note{1}{\italique{Inscriptiones Graecae}, I\expo{3}, 61.34, 68.5, 69.3-4.}%
\note{2}{Andocide, I-\italique{\grec{Περὶ τῶν μυστηρίων}}, 27. La démagogie de Cléonyme se voit elle-même débordée par celle de Pisandre qui renchérit à 10000 drachmes !..}%
\note{3}{Andocide, I-\italique{\grec{Περὶ τῶν μυστηρίων}}, 74.}%
\note{4}{Willems, tome I, p. 108, s'en plaignait déjà en son temps.}%
\note{5}{I. C. Storey, «~The `Blameless Shield' of Kleonymus~», in \italique{Rheinisches Museum für Philologie}, n°132 (1989), p. 247-261.}%
\note{6}{%«~Drôles de noms pour des prisons~!~» objecteront les plus prudents ; ce à quoi les plus érudits répondront que l'\italique{Élysée} fait référence ironiquement au lieu que devait rejoindre son unique résident après avoir reçu les douze balles qui devaient immanquablement être promises à ce ``héros jupitérien'', multirécidiviste de la haute-trahison. % ! « Mate ou gnons ! », tel devait être le mot d'ordre à ``Matignon'' %-- noter la eff{demander à Alexandre + exemples} -- 
%qui n'accueillit que de serviles et soumises serpillières. 
``Palais'' est à prendre dans son acception judiciaire ; suite logique des \italique{palais} de justice. % ; ``Bourbeux'' devait être le nom d'origine du premier quand on y constate l'entassement de la valetaille parlementaire et l'ambiance de basse-cour qui y régnait~; le second porte évidemment le nom de son plus célèbre prisonnier~: Rosa Luxembourg (et les jardins du même nom – attestés parfois au singulier – étaient probablement la cour dans laquelle les prisonniers allaient se délassant quand la folie régnant à l'intérieur se faisait trop oppressante).
}%
\note{7}{voir Willems, t. I p. 106 sq.}%
\note{8}{«~The `Blameless Shield' of Kleonymus~», p. 248.}%
\note{9}{Car je tiens pour nulle et non avenue sa citation du passage des \italique{Cavaliers} (v. 1369-1372) qui \italique{prouve} simplement que M. Storey n'a pas saisi l'image et le jeu de mots d'Aristophane sur \grec{κατὰ σπουδὰς} (mis à dessein pour faire entendre κατ᾿ ἀσπίδας en lieu et place de l'expression figée `κατὰ σπουδήν' qu'on trouve chez Thucydide, Xénophon, Élien, Plutarque, etc.). %Les autres passages sont encore moins probants : \italique{Paix} v. 1172-1187 et \italique{Guêpes} v. 1112-1121. 
Au reste, si M. Storey convient %-- comme il ne fait pas de doute et comme le scholiaste l'avait déjà remarqué -- 
que les vers 827-829 des \italique{Thesmophories} visent Cléonyme, je lui saurais gré d'expliquer le tour `ἀπὸ τῶν ὤμων' à l'aune de sa brillante exégèse.}%
\note{10}{Si banale qu'on en trouve à la pelle parmi nos % n'en verrait pas le bout de la liste si l'on s'avisait de recenser ceux qui y ont recouru parmi les 
modernes brigands. Citons seulement l'exemple le plus caricatural : M. Ciotti-le-bien-nommé, alors qu'il n'était encore que vil attaché-parlementaire fit jouer ses relations pour échapper au service militaire... maintenant que son derrière est confortablement installé au palais bourbeux, bien loin des fangeuses casernes, Chiotti n'a de cesse de haranguer le veau français pour un rétablissement dudit service militaire.% \italique{Ubi quae fuerant vitia, mores sunt.}%
}%
\\\passeLigne}%
\end{groupe}%
%% v. 606
\begin{groupe}%
\personnageGrec{Κλεισθένης}%
\trimIamb{γιγνώσκεθ᾽ ὑμεῖς ἥτις ἔσθ᾽ ἡδὶ γυνή;\apparat{606}{ἥδε R : ἥδ᾽ ἡ \Mon : \lecon{ἡδὶ} Grynaeus}\\\passeLigne}%
\end{groupe}%
%% v. 607
\begin{groupe}%
\personnageGrec{Κορυφαῖος}%
\trimIamb{γιγνώσκομεν δῆτ᾽. ἀλλὰ τὰς ἄλλας ἄθρει.\\\passeLigne}%
\end{groupe}%
%% v. 608-609a
\begin{groupe}%
\personnageGrec{Κλεισθένης}%
\trimIamb{ἡδὶ δὲ δὴ τίς ἐστιν ἡ τὸ παιδίον\\
ἔχουσα;\\\passeLigne}%
\end{groupe}%
%% v. 609b
\begin{groupe}%
\personnageGrec{Μίκα}%
\trimIamb{\ligne{609}\decalage{ἔχουσα;}τίτθη νὴ Δί᾽ ἐμή.\apparat{609}{τίτθην η δί᾽ R : \lecon{τίτθη νὴ Δί᾽} Ellebodius}%
\commentaire{609}{τίτθη.}{Remarquons en passant qu'une nourrice à Athènes n'était pas nécessairement une esclave. Cf. Démosth., LVII, 35. (Note de Willems.)}\\\passeLigne}%
\end{groupe}%
%% v. 609c
\begin{groupe}%
\personnageGrec{Μνησίλοχος}%
\trimIamb{\ligne{609}\decalage{ἔχουσα;τίτθη νὴ Δί᾽ ἐμή.}διοίχομαι.\\\passeLigne}%
\end{groupe}%
%% v. 610
\begin{groupe}%
\personnageGrec{Κλεισθένης}%
\trimIamb{αὕτη σὺ ποῖ στρέφει; μέν᾽ αὐτοῦ. τί τὸ κακόν;%\apparat{610}{τί τὸ κακόν; continuat Clistheni R : Parenti tribuit (cum versu sequenti) \Rpc}
\\\passeLigne}%
\end{groupe}%
%% v. 611a
\begin{groupe}%
\personnageGrec{Μνησίλοχος}%
\trimIamb{ἔασον οὐρῆσαί μ᾽· ἀναίσχυντός τις εἶ.\apparat{611}{εἶ R : \lecon{τις εἶ} Küster}\\\passeLigne}%
\end{groupe}%
%% v. 611b-612
\begin{groupe}%
\personnageGrec{Κλεισθένης}%
\trimIamb{%\ligne{611}\decalage{ἔασον οὐρῆσαί μ᾽.}\\
%\commentaire{611}{ἀναίσχυντός εἶ.}{J'attribue à Clisthène la réplique (qui traitera encore Mnésiloque au vers 638 d'\grec{ἀναίσχυντός}).}%
σὺ δ᾽ οὖν ποίει τοῦτ᾽. ἀναμενῶ γὰρ ἐνθάδε.\apparat{612}{ἀναμένω R : \lecon{ἀναμενῶ} Grynaeus}\\\passeLigne}%
\end{groupe}%
%% v. 613-614
\begin{groupe}%
\personnageGrec{Κορυφαῖος}%
\trimIamb{ἀνάμενε δῆτα καὶ σκόπει γ᾽ αὐτὴν σφόδρα·\\
μόνην γὰρ αὐτὴν ὦνερ οὐ γιγνώσκομεν.\\\passeLigne}%
\end{groupe}%
%% v. 615a
\begin{groupe}%
\personnageGrec{Κλεισθένης}%
\trimIamb{πολύν γε χρόνον οὐρεῖς σύ.\apparat{615}{πολύ R : \lecon{πολύν} Gelenius}\\\passeLigne}%
\end{groupe}%
%% v. 615b-616
\begin{groupe}%
\personnageGrec{Μνησίλοχος}%
\trimIamb{\ligne{615}\decalage{πολύν γε χρόνον οὐρεῖς σύ.}νὴ Δί᾽ ὦ μέλε·\\
στραγγουριῶ γάρ· ἐχθὲς ἔφαγον κάρδαμα.%\testimonia{616}{Suidas σ 1163 (χθὲς γὰρ pro γάρ ἐχθὲς)}%
\commentaire{616}{κάρδαμον}{%Le κάρδαμον est identifié au \italique{cresson alénois} par les différents dictionnaires de grec\marqueur{1}. Alphonse Willems, porté aux \italique{realia}, avait déjà formulé en son temps une objection à cette identification : Mnésiloque justifie le long temps qu'il met à uriner au motif d'avoir mangé du cardame la veille. Hippocrate (\italique{Du Régime}, II, 54) confirme que le cardame provoque la dysurie. Or, de ce que j'ai pu savoir grâce à l'internet -- le temple de la culture pour un néophyte en herbes comme moi --, le cresson, loin d'être dysurique, est au contraire diurétique !
%Une anecdote de Xénophon confirme les dires d'Hippocrate : chez les Perses, où il passait pour indécent de cracher, se moucher ou s'écarter pour faire de l'eau (\italique{Cyropédie}, I, 2, 16), on ne donnait aux enfants à leur repas du midi que du pain, avec du cardame pour tout assaisonnement (\italique{Cyrop}., I, 2, 8 et 11). La consommation de cresson paraît donc impropre aux mœurs perses. 
%Il me semble que le cardame grec n'était peut-être pas un cresson mais plutôt une crucifère cousine : une variété de moutarde ? Quelques indices me le font penser. 
%D'abord, le cardame est très souvent rapproché d'autres noms de moutarde : ainsi dans le passage déjà cité d'Hippocrate, le κάρδαμον est qualifié de θερμόν, tout comme la plante suivante : νᾶπυ. 
%Autre exemple : Athénée (IX, 366d) cite ce vers de Nicandre : Κάρδαμον, ἄρρινόν τε, μελάμφυλλόν τε σίνηπυ. Plus importante encore est cette glose d'Hésychius : \citComm{ἐρύσιμον· κάρδαμον ἄγριον}
%L'érysimon est identifié à notre \italique{sisymbre à cornes} (\latin{Sisymbrium polyceratium}) ; ce qui se recoupe avec mes cyber-recherches qui m'ont appris que le sisymbre était un succédané sauvage de la moutarde.
%
%\eff{Érasme 754}
Les botanistes s'accordent, je crois, à reconnaître dans le κάρδαμον, non le \italique{cresson}, comme on l'a cru jusqu'en ces derniers temps, mais une autre crucifère, connue des savants sous le nom de \latin{erucaria aleppica}. Suivant Fraas, \italique{Synopsis florae classicae}, les Grecs lui donnent encore aujourd'hui le même nom. Hippocrate dit positivement que le \grec{κάρδαμον} provoque la dysurie (\italique{Du Régime}, II, 54). On cite à ce sujet une singulière anecdote, recueillie par Xénophon. Chez les Perses, où il passait pour indécent de cracher, se moucher ou s'écarter pour faire de l'eau, on ne donnait aux enfants à leur repas du midi que du pain, avec du cardame pour tout assaisonnement. (\italique{Cyrop.}, I, 2, 8, 11 et 16) Mettez cela d'accord avec les traités de thérapeutique où le cresson est donné pour un sialagogue. (Note de Willems. Les dictionnaires  \LSJ , Bailly et Pape désignent le \grec{κάρδαμον} respectivement comme « \anglais{nose-smart}, \latin{Lepidium sativum} »,\italique{cresson alénois} ; « cresson » et « \allemand{eine Art Kresse mit bitterem Kraut} ».%
Pour aller dans le sens de Willems, ajoutons que le cresson, outre ses vertus sialagogiques, est loin d'être dysurique : il a des vertus au contraire diurétiques. Cf. le \italique{Dictionnaire des plantes médicinales et vénéneuses de France} de Paul-Victor Fournier, p. 731. Mais la raison de sa consommation chez les Perses est peut-être à chercher ailleurs : cette plante était réputée donner du courage.)}%
%\debutNote{616}%
%\note{1}{Les dictionnaires  \LSJ , Bailly et Pape désignent le \grec{κάρδαμον} respectivement comme « \anglais{nose-smart}, \latin{Lepidium sativum} » (\italique{cresson alénois}), « cresson » et « \allemand{eine Art Kresse mit bitterem Kraut} ».%
%Pour aller dans le sens de Willems, ajoutons que le cresson, outre ses vertus sialagogiques, est loin d'être dysurique : il a aussi des vertus au contraire diurétiques. Cf. le \italique{Dictionnaire des plantes médicinales et vénéneuses de France} de Paul-Victor Fournier, p. 731. Mais la raison de sa consommation chez les Perses est peut-être à chercher ailleurs : cette plante était réputée donner du courage.}%
\\\passeLigne}%
\end{groupe}%
%% v. 617
\begin{groupe}%
\personnageGrec{Κλεισθένης}%
\trimIamb{τί καρδαμίζεις; οὐ βαδιεῖ δεῦρ᾽ ὡς ἐμέ;\\\passeLigne}%
\end{groupe}%
%% v. 618a
\begin{groupe}%
\personnageGrec{Μνησίλοχος}%
\trimIamb{τί δῆτά μ᾽ ἕλκεις ἀσθενοῦσαν;\\\passeLigne}%
\end{groupe}%
%% v. 618b-619a
\begin{groupe}%
\personnageGrec{Κλεισθένης}%
\trimIamb{\ligne{618}\decalage{τί δῆτά μ᾽ ἕλκεις ἀσθενοῦσαν;}εἰπέ μοι,\\
τίς ἔστ᾽ ἀνήρ σοι;\apparat{619}{ἔσθ᾽ R : \lecon{ἔστ᾽} \Mon}\\\passeLigne}%
\end{groupe}%
%% v. 619b-620
\begin{groupe}%
\personnageGrec{Μνησίλοχος}%
\trimIamb{\ligne{619}\decalage{τίς ἔστ᾽ ἀνήρ σοι;}τὸν ἐμὸν ἄνδρα πυνθάνει;\\
τὸν δεῖνα γιγνώσκεις, τὸν ἐκ Κοθωκιδῶν.%\testimonia{620}{(Κοθωκιδῶν) Photius κ 859 Κοθωκίδαι (ap Scholium) $\sim$ Hesychius  κ 3218}%\apparat{620}{add. signum interrogationis Canini}%
\commentaire{619}{τὸν ἐκ Κοθωκιδῶν;}{« Du $[$dème$]$ des Cothôcides ». Ce dème faisait partie de la tribu des \grec{Οἰνηίς} au nord-ouest de la ville\marqueur{1}. Y a-t-il une raison particulière pour que Mnésiloque évoque ce dème ? certainement... mais cette raison nous reste obscure et aucune explication satisfaisante n'a jusqu'ici été avancée.}%
\debutNote{619}%
\note{1}{L'on pourra consulter le travail de John S. Traill, \italique{The Political Organization of Attica}, \italique{Hesperia} supplement XIV (1975). Voir notamment les pages 49, 62 et 68 ainsi que les cartes d'Athènes à la fin.}\\\passeLigne}%
\end{groupe}%
%% v. 621a
\begin{groupe}%
%
\personnageGrec{Κλεισθένης}%
\trimIamb{τὸν δεῖνα; ποῖον;\\\passeLigne}%
\end{groupe}%
%% v. 621b-622a
\begin{groupe}%
\personnageGrec{Μνησίλοχος}%
\trimIamb{\ligne{621}\decalage{τὸν δεῖνα; ποῖον;}ἔσθ᾽ ὁ δεῖν᾽, ὃς καί ποτε\apparat{621}{(Κλεισθ.) continuat R : \lecon{(Μνη.)} \Mon}\\
τὸν δεῖνα τὸν τοῦ δεῖνα—\\\passeLigne}%
\end{groupe}%
%% v. 622b-623a
\begin{groupe}%
\personnageGrec{Κλεισθένης}%
\trimIamb{\ligne{622}\decalage{τὸν δεῖνα τὸν τοῦ δεῖνα—}ληρεῖν μοι δοκεῖς.\\
ἀνῆλθες ἤδη δεῦρο πρότερον;\\\passeLigne}%
\end{groupe}%
%% v. 623b-624a
\begin{groupe}%
\personnageGrec{Μνησίλοχος}%
\trimIamb{\ligne{623}\decalage{ἀνῆλθες ἤδη δεῦρο πρότερον;}νὴ Δία\\
ὁσέτη γε.\apparat{624}{ὅς᾽ ἔτη R : \lecon{ὁσέτη} G. H. Schaefer ad Gregor. Cor. p.65 \espApp σοὐστὶν R : \lecon{σοὐστὶ} Dindorf (σοῦ ᾽στὶ iam Biset)}\\\passeLigne}%
\end{groupe}%
%% v. 624b
\begin{groupe}%
\personnageGrec{Κλεισθένης}%
\trimIamb{\ligne{624}\decalage{ὁσέτη γε.}καὶ τίς σοὐστὶ συσκηνήτρια;%
\commentaire{624}{συσκηνήτρια.}{Il faut croire d'après ce passage et le vers 658 que lors des Thesmophories, des tentes pour deux femmes étaient installées à proximité du Pnyx.}\\\passeLigne}%
\end{groupe}%
%% v. 625a
\begin{groupe}%
\personnageGrec{Μνησίλοχος}%
\trimIamb{ἡ δεῖν᾽ ἔμοιγ᾽.\apparat{625}{ἔμοιγε R : \lecon{ἔμοιγ᾽} edd. pr.}\\\passeLigne}%
\end{groupe}%
%% v. 625b
\begin{groupe}%
\personnageGrec{Κλεισθένης}%
\trimIamb{\ligne{625}\decalage{ἡ δεῖν᾽ ἔμοιγ᾽.}οἴμοι τάλας. οὐδὲν λέγεις.\\\passeLigne}%
\end{groupe}%
%% v. 626-629
\begin{groupe}%
\personnageGrec{Μίκα}%
\trimIamb{ἄπελθ᾽. ἐγὼ γὰρ βασανιῶ ταύτην καλῶς\\
ἐκ τῶν ἱερῶν τῶν πέρυσι· σὺ δ᾽ ἀπόστηθί μοι,\\
ἵνα μὴ 'πακούσῃς ὢν ἀνήρ. σὺ δ᾽ εἰπέ μοι\\
ὅ τι πρῶτον ἡμῖν τῶν ἱερῶν ἐδείκνυτο.\\\passeLigne}%
\end{groupe}%
%% v. 630
\begin{groupe}%
%
\personnageGrec{Μνησίλοχος}%
\trimIamb{φέρ᾽ ἴδω, τί μέντοι πρῶτον ἦν; ἐπίνομεν.%\testimonia{630-1}{Suidas π 2563 (sed προπίνομεν pro προὐ- cod. V)}
\apparat{630}{\lecon{μέντοι} \Rpc in mg., Suid. π 2563 : om. \Rac, \Mon}\\\passeLigne}%
\end{groupe}%
%% v. 631a
\begin{groupe}%
\personnageGrec{Μίκα}%
\trimIamb{τί δὲ μετὰ τοῦτο δεύτερον;\apparat{631}{\lecon{δὲ} R, Suid. π 2563 : δαὶ Bentley \espApp με R : \lecon{μετὰ} Suid.}\\\passeLigne}%
\end{groupe}%
%% v. 631b
\begin{groupe}%
\personnageGrec{Μνησίλοχος}%
\trimIamb{\ligne{631}\decalage{τί δὲ μετὰ τοῦτο δεύτερον;}προὐπίνομεν.\\\passeLigne}%
\end{groupe}%
%% v. 632
\begin{groupe}%
\personnageGrec{Μίκα}%
\trimIamb{ταυτὶ μὲν ἤκουσάς τινος· τί δ᾽ ἦν τρίτον;\apparat{632}{τί δὲ τρίτον R : τί δ᾽ αὖ Bentley : τί δαὶ τρίτον Elmsley : \lecon{τί δ᾽ ἦν τρίτον} Teubner}\\\passeLigne}%
\end{groupe}%
%% v. 633
\begin{groupe}%
\personnageGrec{Μνησίλοχος}%
\trimIamb{σκάφιον Ξένυλλ᾽ ᾔτησεν· οὐ γὰρ ἦν ἀμίς.%\testimonia{633}{Pollux X 45 ἐν Πολυίδῳ (sic) Ἀριετοφάνης· σκάφιον Ξένυλλ᾽ ἤτησεν· οὐ γὰρ ἦν ἁμίς \espTest Photius σ 286 σκάφιον}%
\commentaire{633}{σκάφιον Ξένυλλ᾽ ᾔτησεν· οὐ γὰρ ἦν ἀμίς.}{Acculé, c'est à ce vers que Mnésiloque fait son hara-kiri : il lance une dernière boutade en guise de baroud d'honneur ! L'explication communément admise laisse songeur quant à l'humour des commentateurs d'Aristophane : Mnésiloque emploierait selon eux un terme impropre (l'\italique{amis} serait une sorte de pot de chambre à usage exclusivement masculin en raison d'un col étroit... mais sa forme varie et le terme pourrait très bien aussi désigner ce genre d'ustensile pour les femmes\marqueur{1}). Quoi qu'il en soit, c'est un peu \italique{gros} : tout rustre qu'il soit, Mnésiloque ne peut pas commettre une telle bévue ; je dirais même plus~: dans ce domaine, il est certainement tout aussi savant que l'\italique{Athénienne de la rue}. 
Voici une explication bien meilleure : le \grec{σκάφιον} désignait à l'origine une \italique{coupe} (en forme de \italique{barque} selon l'étymologie). Il n'acquiert le sens quasi-exclusif d'\italique{urinoir} qu'en passant au latin \latin{scaphium}. À l'époque d'Aristophane, les deux sens devaient encore co-exister. La première partie du vers « σκάφιον Ξένυλλ᾽ ᾔτησεν » était donc interprétée par Mica favorablement~: Xénylla a demandé une coupe $[$pour faire une libation$]$... la chute est rude et Mnésiloque dans la seconde partie accuse les femmes d'un usage détourné des coupes consacrées qui passe la mesure : cet individu n'a rien des manières d'une femme : après avoir uriné publiquement, voilà qu'il prononce maintenant des \italique{insanités} (οὐδὲν ὑγιὲς v. 636).\\
%Ce n'est d'ailleurs pas autrement que le comprenait le scholiaste :
%\citComm{μὴ εὑρίσχων τί εἴπῃ ὁ κηδεστής φησιν ὅτι ἀντὶ ἀμίδος σκάφιον ᾔτησεν ἣ Ξένυλλα. ἀμὶς δέ ἐστιν οὐροδόχον ἀγγεῖον.\\
%\francais{Ne trouvant pas quoi répliquer, le Parent dit qu'en place d'un amis, c'est un scaphion qu'a demandé Xénylla. Un amis est un récipient pour recueillir l'urine.}}
Voici pour le sens propre -- si l'on peut dire. Mais, comme de coutume, Aristophane joue à différents niveaux : la réplique est aussi un clin d'œil au proverbe grec bien connu\marqueur{2} :%
\citComm{Ἄμας ἀπῄτουν, οἱ δ' ἀπηρνοῦντο σκάφας\\
Je demandais des bêches, ils me refusent des seaux.}%
(Jeux de consonances entre ἄμη/ἀμίς et σκάφη/σκάφιον.)
Apprécions à sa juste valeur l'à-propos de l'allusion au proverbe : il repose sur un quiproquo dû à la polysémie d'ἄμη -- comme est polysémique le σκάφιον dont parle Mnésiloque. De plus, le proverbe évoque ceux qui fournissent des réponses sans rapport avec la question posée -- ce qui est aussi le cas de notre pauvre Mnésiloque, bien embarrassé.
ἄμη et σκάφη pouvaient en effet être tous deux synonymes de \italique{seau}\marqueur{3}... mais c'était des \italique{houes} (autre sens de ἄμη) que demandait le paysan du proverbe...
}%
\debutNote{633}%
\note{1}{Pour toutes les remarques de cette note, je renvoie aux articles \latin{amis} et \latin{scaphium} du \italique{Dictionnaire des antiquités grecques et romaines} de Daremberg \& Saglio.}%
\note{2}{Il a été relevé par Suidas en α 1504 et α 1574. Érasme y consacre une entrée dans ses \italique{Adages} (au n°1149) et l'évoque aussi au n°1880. Maintes fois recensé par les parémiographes (par exemple Zénobius I, 83), il était connu de Plutarque (\italique{Du bavardage}, OM, 6, 35, 512e) et le scholiaste de la \italique{Paix} d'Aristophane le cite encore à propos de la chanson de Datis le Lydien aux vers 289-291 et de l'équivoque en grec que suggère la forme barbaristique χαίρομαι -- au passage : Aristophane joue sur le stéréotype du débauché « Lydien ($=$ Datis) à midi », Λυδὸς ἐν μεσημβρίᾳ (rapporté par Athénée, XII, 515d-517a et recensé par Érasme dans ses \italique{Adages}, n° 1594). Enfin, le scholiaste de Théocrite (\italique{Idylles}, IV, 10) cite aussi de ce proverbe. Pour finir, Aristophane y fait de nouveau allusion au fr. 927 Kassel-Austin : ἄγροικός εἰμι· τὴν σκάφην σκάφην λέγω.}%
\note{3}{Comme dans le proverbe pour évoquer abondance de biens : ἄμαις καὶ σκάφαις ἀρύσασθαι. Abondance de biens nuit, soit dit en passant, à la santé du corps comme celle de l'esprit. %(contrairement à un lieu commun par trop répandu dont 
Bloy n'a, à ma connaissance, étrangement pas fait l'exégèse de ce lieu commun -- Flaubert aurait parlé d'une idée reçue. % (c'est-à-dire bourgeoise). 
Les latins et leur solide bon sens savaient que \latin{Copia parit fastidium} et \latin{Neque laus in copia, neque culpa in penuria constitit}.)}\\\passeLigne}%
\end{groupe}%
%% v. 634-635a (part 1)
\begin{groupe}\personnageGrec{Μίκα}
\trimIamb{οὐδὲν λέγεις. δεῦρ᾽ ἐλθὲ δεῦρ᾽ ὦ Κλείσθενες·\apparat{634}{δεῦρο \Rac, \Mon : \lecon{δευρ ὦ} \Rdeux \espApp Κλεισόσθενες R : \lecon{Κλείσθενες} Bonini \ligneApp{635} ἀνὴρ R : \lecon{ἁνὴρ} Bentley}\\\passeLigne}%
\end{groupe}%
\endnumbering
\end{Leftside}
\begin{Rightside}
\beginnumbering
%% v. 590-591
\begin{groupe}%
\personnageFr{Clisthène}%
\trimIambFr{Euripide l'a flambé et épilé ;\\
et pour tout le reste l'a équipé en femme.\\\passeLigne}%
\end{groupe}%
%% v. 592-594
\begin{groupe}%
\personnageFr{Mnésiloque}%
\trimIambFr{Vous en fiez-vous à lui? Quel homme est assez
niais pour s'être laissé épiler.
Je n'en crois rien, pour ma part, ô déesses très honorées.\\\passeLigne}%
\end{groupe}%
%% v. 595-596
\begin{groupe}%
\personnageFr{Clisthène}%
\trimIambFr{Tu veux rire. Je ne serais pas venu l'annoncer\\
si je ne le tenais de ceux qui le savent bien.\\\passeLigne}%
\end{groupe}%
%% v. 597-602
\begin{groupe}%
\personnageFr{Coryphée}%
\trimIambFr{C'est une grave affaire qu'on nous annonce là.\\
Or ça, femmes, il s'agit, non de se morfondre,\testimonia{598}{(ἐλινύειν) Hesychius ε 2095, 2097 ; cf. Moeris ε 61 (ἐλινύων $=$ ἀναπαυόμενος)% $\sim$ Hesychius ε 2097
}\\
mais de guetter l'homme, de chercher où\\
il se dérobe à nous et se tient caché.\\
Toi, aide-nous à le trouver, que nous t'ayons\testimonia{601-2}{Suidas π 2540}\\
encore cette autre obligation, ô notre patron.\\\passeLigne}%
\end{groupe}%
%% v. 603a
\begin{groupe}%
\personnageFr{Clisthène}%
\trimIambFr{Voyons, toi d'abord, qui es-tu ?\\\passeLigne}%
\end{groupe}%
%% v. 603b
\begin{groupe}%
\personnageFr{Mnésiloque}%
\trimIambFr{\ligne{603}\espace{4cm}(À part :) Où se tourner ?\\\passeLigne}%
\end{groupe}%
%% v. 604a
\begin{groupe}%
\personnageFr{Clisthène}%
\trimIambFr{... Car on va vous passer en revue.\\\passeLigne}%
\end{groupe}%
%% v. 604b
\begin{groupe}%
\personnageFr{Mnésiloque}%
\trimIambFr{\ligne{604}\espace{2.5cm}(À part :) Malheureux moi !\\\passeLigne}%
\end{groupe}%
%% v. 605
\begin{groupe}%
\personnageFr{Mika}%
\trimIambFr{Tu me demandais qui je suis ? la femme de Cléonyme.\\\passeLigne}%
\end{groupe}%
%% v. 606
\begin{groupe}%
\personnageFr{Clisthène}%
\trimIambFr{Vous savez, vous autres, qui est cette femme ?\\\passeLigne}%
\end{groupe}%
%% v. 607
\begin{groupe}%
\personnageFr{Coryphée}%
\trimIambFr{Nous le savons. Examine les autres.\\\passeLigne}%
\end{groupe}%
%% v. 608-609a
\begin{groupe}%
\personnageFr{Clisthène}%
\trimIambFr{Et celle-là, qui est-elle ? Celle qui porte\\
l'enfant ?\\\passeLigne}%
\end{groupe}%
%% v. 609b
\begin{groupe}%
\personnageFr{Mika}%
\trimIambFr{\ligne{609}\espace{1.3cm}C'est ma nourrice, par Zeus.\\\passeLigne}%
\end{groupe}%
%% v. 609c
\begin{groupe}%
\personnageFr{Mnésiloque}%
\trimIambFr{\ligne{609}\espace{3.8cm}Je suis perdu. (Il va pour sortir.)\\\passeLigne}%
\end{groupe}%
%% v. 610
\begin{groupe}%
\personnageFr{Clisthène}%
\trimIambFr{Holà; toi, où vas-tu ? Reste ici. (Mnésiloque s'accroupit.) Qu'est-ce qui t'arrive ?\\\passeLigne}%
\end{groupe}%
%% v. 611a
\begin{groupe}%
\personnageFr{Mnésiloque}%
\trimIambFr{Laisse-moi lâcher l'eau. Tu es un impudent.\\\passeLigne}%
\end{groupe}%
%% v. 611b-612
\begin{groupe}%
\personnageFr{Clisthène}%
\trimIambFr{%\ligne{611}\espace{2.9cm}Tu es une impudente.\\
Eh bien, fais. J'attendrai ici.\\\passeLigne}%
\end{groupe}%
%% v. 613-614
\begin{groupe}%
\personnageFr{Coryphée}%
\trimIambFr{Oui, attends, et surveille-la bien.\\
Elle est la seule que nous ne connaissions pas.\\\passeLigne}%
\end{groupe}%
%% v. 615a
\begin{groupe}%
\personnageFr{Clisthène}%
\trimIambFr{Tu es bien longue à faire de l'eau.\\\passeLigne}%
\end{groupe}%
%% v. 615b-616
\begin{groupe}%
\personnageFr{Mnésiloque}%
\trimIambFr{\ligne{615}\espace{4.2cm}Par Zeus, l'ami,
c'est que j'ai une rétention : hier j'ai mangé du cardame.\testimonia{616}{Suidas σ 1163 (χθὲς γὰρ pro γάρ ἐχθὲς)}\\\passeLigne}%
\end{groupe}%
%% v. 617
\begin{groupe}%
\personnageFr{Clisthène}%
\trimIambFr{Que parles-tu de cardame ? Veux-tu venir ici, près de moi ? (Il la saisit)\\\passeLigne}%
\end{groupe}%
%% v. 618a
\begin{groupe}%
\personnageFr{Mnésiloque}%
\trimIambFr{Pourquoi me tirer ainsi ? une valétudinaire !\\\passeLigne}%
\end{groupe}%
%% v. 618b-619a
\begin{groupe}%
\personnageFr{Clisthène}%
\trimIambFr{\espace{4.5cm}Réponds-moi.\\
Quel est ton mari ?\\\passeLigne}%
\end{groupe}%
%% v. 619b-620
\begin{groupe}%
\personnageFr{Mnésiloque}%
\trimIambFr{\espace{2.7cm}Mon mari, dis-tu ?\\
Connais-tu... chose, du dème de Cothocides.\testimonia{620}{(Κοθωκιδῶν) Photius κ 859 Κοθωκίδαι (ap Scholium) $\sim$ Hesychius  κ 3218}\\\passeLigne}%
\end{groupe}%
%% v. 621a
\begin{groupe}%
\personnageFr{Clisthène}%
\trimIambFr{Chose ? Qui ça ?\\\passeLigne}%
\end{groupe}%
%% v. 621b-622a
\begin{groupe}%
\personnageFr{Mnésiloque}%
\trimIambFr{\espace{2cm}C'est chose, qui naguère...\\
avec chose, fils de chose...\\\passeLigne}%
\end{groupe}%
%% v. 622b-623a
\begin{groupe}%
\personnageFr{Clisthène}%
\trimIambFr{\espace{3.7cm}Tu m'as l'air de baliverner.\\
Es-tu déjà venue ici ?\\\passeLigne}%
\end{groupe}%
%% v. 623b-624a
\begin{groupe}%
\personnageFr{Mnésiloque}%
\trimIambFr{\espace{4.5cm}Oui, par Zeus,\\
tous les ans.\\\passeLigne}%
\end{groupe}%
%% v. 624b
\begin{groupe}%
\personnageFr{Clisthène}%
\trimIambFr{\espace{1.5cm}Et quelle est ta compagne de tente ?\\\passeLigne}%
\end{groupe}%
%% v. 625a
\begin{groupe}%
\personnageFr{Mnésiloque}%
\trimIambFr{La mienne, c'est chose...\\\passeLigne}%
\end{groupe}%
%% v. 625b
\begin{groupe}%
\personnageFr{Clisthène}%
\trimIambFr{\espace{2.1cm}Misère de moi. Tu ne dis rien.\\\passeLigne}%
\end{groupe}%
%% v. 626-629
\begin{groupe}%
\personnageFr{Mika}%
\trimIambFr{Retire-toi. C'est moi qui vais l'éplucher comme il faut
sur les cérémonies de l'an dernier. — Tiens-toi à distance !
tu es un homme, tu ne dois rien entendre. — Toi, dis-moi
lequel des rites sacrés nous fut révélé le premier.\\\passeLigne}%
\end{groupe}%
%% v. 630
\begin{groupe}%
\personnageFr{Mnésiloque}%
\trimIambFr{Voyons. Que fût-ce encore que le premier? ... Nous \fin bûmes.\testimonia{630-1}{Suidas π 2563 (sed προπίνομεν pro προὐ- cod. V)}\\\passeLigne}%
\end{groupe}%
%% v. 631a
\begin{groupe}%
\personnageFr{Mika}%
\trimIambFr{Et après ? le second ?\\\passeLigne}%
\end{groupe}%
%% v. 631b
\begin{groupe}%
\personnageFr{Mnésiloque}%
\trimIambFr{\espace{4cm}Nous bûmes à nos santés.\\\passeLigne}%
\end{groupe}%
%% v. 632
\begin{groupe}%
\personnageFr{Mika}%
\trimIambFr{Tu l'as appris de quelqu'une. Et le troisième ?\\\passeLigne}%
\end{groupe}%
%% v. 633
\begin{groupe}%
\personnageFr{Mnésiloque}%
\trimIambFr{Xenylla demanda une nef, car il n'y avait pas de pot \fin  de chambre.\testimonia{633}{Pollux X 45 ἐν Πολυίδῳ (sic) Ἀριετοφάνης· σκάφιον Ξένυλλ᾽ ἤτησεν· οὐ γὰρ ἦν ἁμίς \espTest Photius σ 286 σκάφιον}\\\passeLigne}%
\end{groupe}%
%% v. 634-635a (part 1)
\begin{groupe}\personnageFr{Mika}
\trimIambFr{Cela ne veut rien dire. Ici, viens ici, ô Clisthène.\\\passeLigne}%
\end{groupe}%
\endnumbering
\end{Rightside}
\end{pages}
\Pages
\end{document}
